\documentclass[11pt]{amsart}
\usepackage[margin=1.1in]{geometry}
\usepackage{amsmath,amssymb,amsthm,mathtools,mathrsfs}
\usepackage{enumitem}
\usepackage{booktabs,longtable,array}
\usepackage{xcolor}
\usepackage[colorlinks=true,linkcolor=blue!50!black,
  citecolor=green!40!black,urlcolor=blue!60!black,
  hypertexnames=false]{hyperref}

\setcounter{tocdepth}{1}

\theoremstyle{plain}
\newtheorem{theorem}{Theorem}[section]
\newtheorem{proposition}[theorem]{Proposition}
\newtheorem{conjecture}[theorem]{Conjecture}
\theoremstyle{definition}
\newtheorem{definition}[theorem]{Definition}
\newtheorem{question}[theorem]{Question}
\theoremstyle{remark}
\newtheorem{remark}[theorem]{Remark}
\numberwithin{equation}{section}

\newcommand{\taylor}[1]{\textcolor{red}{[Taylor: #1]} }
\newcommand{\todo}[1]{\textcolor{orange}{[todo: #1]}}
\newcommand{\agentreply}[1]{\textcolor{gray}{[agent: #1]}}

\newcommand{\arxivlink}[2]{\cite{#1}\,%
  \href{https://arxiv.org/abs/#2}{\textsf{arXiv:#2}}}
\newcommand{\ayoubpage}{\cite{AyoubHomepage}\,%
  \href{https://user.math.uzh.ch/ayoub/}{\textsf{user.math.uzh.ch/ayoub}}}
\newcommand{\stackstag}[1]{\cite[Tag~#1]{StacksProject}\,%
  \href{https://stacks.math.columbia.edu/tag/#1}{\textsf{#1}}}
\expandafter\def\expandafter\UrlBreaks\expandafter{\UrlBreaks\do\-}
\newcommand{\ghlink}[2]{\cite{#1}\,%
  \href{https://github.com/#2}{\nolinkurl{github.com/#2}}}

\newcommand{\C}{\mathbb{C}}
\newcommand{\R}{\mathbb{R}}
\newcommand{\Q}{\mathbb{Q}}
\newcommand{\Z}{\mathbb{Z}}
\newcommand{\N}{\mathbb{N}}
\newcommand{\F}{\mathbb{F}}
\newcommand{\Cp}{\mathbb{C}_{+}}
\newcommand{\Xbar}{\overline{X}}
\newcommand{\GR}{\Gamma_{\mathbb{R}}}
\newcommand{\GC}{\Gamma_{\mathbb{C}}}
\newcommand{\cO}{\mathcal{O}}
\newcommand{\cH}{\mathscr{H}}
\newcommand{\cA}{\mathcal{A}}
\newcommand{\cC}{\mathcal{C}}
\newcommand{\Sch}{\mathrm{Sch}}
\newcommand{\Schbar}{\overline{\mathrm{Sch}}}
\newcommand{\SmProj}{\mathrm{SmProj}}
\newcommand{\Den}{\mathrm{Den}}
\newcommand{\Flow}{\mathrm{Flow}}
\newcommand{\Mods}[1]{\mathrm{Mod}_{#1}}
\DeclareMathOperator{\Tr}{Tr}
\DeclareMathOperator{\tr}{tr}
\DeclareMathOperator{\Spec}{Spec}
\DeclareMathOperator{\spec}{spec}
\DeclareMathOperator{\divisor}{div}
\DeclareMathOperator{\Ima}{Im}
\DeclareMathOperator{\Rea}{Re}
\DeclareMathOperator{\diag}{diag}
\DeclareMathOperator{\ord}{ord}
\newcommand{\Sha}{\operatorname{Sha}}
\newcommand{\detreg}{{\det}_{\infty}}
\renewcommand{\Re}{\Rea}
\renewcommand{\Im}{\Ima}

\title[Hilbert--P\'olya operators and Deninger cohomology]{The Approach
to the Riemann Hypothesis via Hilbert--P\'olya Operators and Deninger
Cohomology Theories}

\date{21 November 2026}

\begin{document}







\begin{abstract}
These notes formulate one Deninger cohomology theory with a flow,
products, duality, and a compatible positive star.  A realization
whose spectral divisor is the completed zeta function would imply
the Riemann hypothesis by Deninger's formal argument after Serre.
We propose two conjectural comparisons with that theory:
Flach--Morin Weil--Arakelov cohomology on eligible objects as its
derived integer-twist fibres, and a multiplicity-preserving spectral
probe of the trivial norm-one sector of CCM adelic cyclic homology
as its degree-one part over number fields.
Neither comparison nor the common positive theory is constructed here.
Finite-field examples calibrate the axioms, while counterexamples
show that integer-twist fibres alone miss noninteger eigenvalues,
scalar trace spaces can lose zero multiplicities, and a naive
arithmetic tensor square need not have discrete spectral blocks.
The use of AI assistance is described in Section~\ref{sec:ai}.
This manuscript has not been checked by a human. Readers use it at
their own risk and are responsible for independently verifying any
claim, proof, computation, or citation they rely on.
\end{abstract}


\author{Printed by Taylor Dupuy \\
AI GENERATED}
\maketitle
\tableofcontents


\section{Introduction and main results}\label{sec:overview}

\subsection{The problem and the comparison proposal}
\label{subsec:search-findings}

The Riemann hypothesis asks whether every nontrivial zero of the
Riemann zeta function $\zeta(s)$ has real part $1/2$. Its analogue for a
smooth projective curve over a finite field is a theorem of Weil.
There the zeta function is determined by the action of Frobenius on
cohomology, and the bound on its zeros follows from positivity for
correspondences on the square of the curve. These notes ask whether
a comparable cohomology and positivity statement can be formulated for
$\Spec\Z$ and its arithmetic relatives.

The finite-field precedent separates two mathematical tasks. A
cohomology theory expresses the zeta function through linear operators
and supplies products and duality. A compatible positive form then
controls the eigenvalues of those operators. In the higher-dimensional
analogue, the Weil conjectures describe the expected zeta-function
properties, while the standard conjectures propose geometric
positivity sufficient for the cohomological argument. Deninger's
program replaces Frobenius over a finite field by a real flow on a
hypothetical arithmetic cohomology. The infinitesimal generator of that
flow is to have the zeta zeros as eigenvalues, and a compatible positive
form would force them onto the critical line.

The geometric input is part of the problem. The arithmetic curve
$\Spec\Z$ has the primes as its closed points; its Arakelov
compactification also includes the infinite place. An absolute or
$\F_1$ geometry seeks a setting in which this curve has a useful square,
playing the role of the surface in Weil's proof. Connes--Consani's
arithmetic and scaling sites supply concrete characteristic-one
constructions connected to this question: the arithmetic site's
counting distribution gives the completed zeta function, and periodic
orbits of the scaling site carry a Riemann--Roch theorem
\cite[Theorems~2.6--2.7]{ConnesConsani2014ArithmeticSite},
\cite[Theorem~5.17]{ConnesConsani2017ScalingGeometry}.
The global cohomology and intersection positivity needed to complete
this analogy remain construction problems in the sources considered
here.



A Deninger cohomology theory (Definition~\ref{def:deninger-theory}) is
a functor from compactified schemes to Deninger structures: a graded
flow with cup product and duality, whose positive star is a property,
together with the zeta divisor as an axiom on the functor.
Conjectures~\ref{conj:arithmetic-realization}
and~\ref{conj:arithmetic-positivity} ask for one such theory on the
common arithmetic and finite-field input category, and their full
versions add compact supports and Galois descent.  Its positive form
gives
the conditional zero-line theorem, Theorem~\ref{thm:rh}; for number
fields this is Deninger's argument after Serre.  The axiom list does
not assert that such a global theory has been constructed.

The comparison proposal in Conjecture~\ref{conj:fm-ccm-common} gives
two ways to test the same hypothetical theory.  The derived fibres
at integer twists should recover Flach--Morin cohomology on the
objects where it is defined.  Over a number field, a spectral probe
of the trivial norm-one sector of CCM adelic cyclic homology should
recover its degree-one zero spectrum with multiplicities.  This
extraction requires finite-dimensional generalized eigenspaces
with discrete support; its topological and trace comparison remains
open.  The existing sources establish a finite-field fixed-twist
comparison under a finite-generation hypothesis
(Theorem~\ref{thm:imp-fm20}) and an adelic spectral trace formula
\cite[Theorem~4.16]{ConnesConsaniMarcolli2009Weil}; they do not
construct these two comparisons to one positive arithmetic theory.

The manuscript first explains the cohomological precedents and the
finite-field logarithmic suspension, then proves the formal
implications and the negative tests.  In particular, integer-twist
fibres do not determine the full generator
(Proposition~\ref{prop:R3}), and scalar spectral models can lose
zero multiplicities (Proposition~\ref{prop:R9}).  The proposed
arithmetic-square K\"unneth and Hodge-index conditions require
additional product data; their naive strong forms fail in the
examples of Section~\ref{sec:kunneth}.  Sections~\ref{sec:candidates}
and~\ref{sec:imported} compare the candidate constructions with the
single axiom list and state the exact source hypotheses.






\subsection{A common spectral formalism}
\label{subsec:uniform-formalism}

This subsection concerns the values of a cohomology theory: spectral
cores with a flow, the Deninger axioms for a single datum, Deninger
structures, and the purity implication. The source category of
compactified schemes is fixed in Section~\ref{subsec:uniform-instances}.


Hermitian forms are linear in their first variable. For a nonzero
meromorphic function $Z$ on $\C$, write
$\divisor Z=\sum_\rho\ord_\rho(Z)[\rho]$, with poles having negative
order. The definitions of spectral core and positive infinitesimal
polarization below isolate the linear algebra used in the main theorem.

\begin{definition}\label{def:spectral-core}
Let $V$ be a complex vector space and $\Theta:V\to V$ a linear map.
\begin{enumerate}[label=(\alph*)]
\item For $\rho\in\C$ the \emph{generalized eigenspace}
  $V_{\Theta\sim\rho}$ is $\bigcup_{k\ge1}\ker(\Theta-\rho)^k$.
\item $(V,\Theta)$ is a \emph{spectral core} if $V$ is the algebraic
  direct sum of the $V_{\Theta\sim\rho}$, $\rho\in\C$, each of them is
  finite-dimensional, and the set
  $\spec\Theta=\{\rho:V_{\Theta\sim\rho}\ne0\}$ has no finite
  accumulation point.
\item The \emph{multiplicity} of $\rho$ is
  $m(\rho)=\dim V_{\Theta\sim\rho}$.
\end{enumerate}
\end{definition}

The cohomology axioms use a flow in place of Frobenius and a unit
adapted to the arithmetic base. The following definitions fix these
objects before the axioms are stated.

\begin{definition}\label{def:flow}
Let $V$ be a complex vector space which is the algebraic direct sum of
finite-dimensional subspaces $V_\alpha$, and let $\Theta:V\to V$ be a
linear map preserving each $V_\alpha$.
\begin{enumerate}[label=(\alph*)]
\item The \emph{flow} generated by $\Theta$ is the family of linear maps
  $\phi^t=e^{t\Theta}$, $t\in\R$, defined blockwise on each $V_\alpha$;
  $\Theta$ is its \emph{generator}.
\item For $n\in\Z$ the \emph{Tate twist} $V(n)$ is $V$ with generator
  $\Theta-n$ (flow $e^{-nt}\phi^t$).
\item $\C(-d)$ denotes the line $\C$ with generator multiplication by
  $d$.
\end{enumerate}
\end{definition}

\begin{definition}\label{def:unit}
Let $q>1$ be real.
\begin{enumerate}[label=(\alph*)]
\item The \emph{arithmetic unit} is the spectral core
  $\mathbb 1=(\C,0)$ with its multiplication.
\item The \emph{$q$-periodic unit} is the spectral core
  $\mathbb 1_q=\bigoplus_{n\in\Z}\C v^n$ (algebraic direct sum) with
  $\Theta v^n=(2\pi in/\log q)v^n$ and multiplication
  $v^nv^m=v^{n+m}$.
\item For a spectral core $(U,\Theta_U)$ and $k\in\Z$, $U(-k)$ is $U$
  with generator $\Theta_U+k$, as in Definition~\ref{def:flow}.
\end{enumerate}
\end{definition}

Over a number field the unit is $\mathbb 1$; over a fixed finite field
it is $\mathbb 1_q$, whose periodic spectrum records all logarithms
of Frobenius. The latter is Deninger's finite-Fourier-series model
\cite[pp.~167--168 and Proposition~3.1]{Deninger1998Analogies}.

\begin{definition}\label{def:aw}
Let $d\ge0$ be an integer.
\begin{enumerate}[label=(\alph*)]
\item A \emph{Weil datum of type $(d,0)$} is a tuple
  $(V^\bullet,\Theta,\cup,\Tr,*)$ consisting of spectral cores
  $(V^j,\Theta_j)$, $0\le j\le2d$, bilinear \emph{cup products}
  $\cup:V^i\times V^k\to V^{i+k}$ (zero if $i+k>2d$), a linear
  \emph{trace} $\Tr:V^{2d}\to\C$, and antilinear \emph{stars}
  $*_j:V^j\to V^{2d-j}$.
\item Its \emph{duality pairings} are
  $B_j(u,v)=\Tr(u\cup v)$ on $V^j\times V^{2d-j}$, and its \emph{star
  forms} are $h_j(u,v)=\Tr(u\cup*_jv)$ on $V^j$.
\item The \emph{spectral divisor} of $(V^j,\Theta_j)$ is
  $\operatorname{Sp}_j=\sum_\rho m_j(\rho)[\rho]$, where $m_j$ is the
  multiplicity function of Definition~\ref{def:spectral-core}.
\end{enumerate}
\end{definition}

\begin{definition}[Deninger axioms]
\label{def:weil-cohomology}
For an individual Weil datum $D$ of type $(d,0)$, a nonzero meromorphic
function $Z$ on $\C$, and a unit spectral core $U$ with multiplication,
we say that $D$ \emph{satisfies AW$k$ for $(Z,U)$} when condition AW$k$
below holds.
\begin{enumerate}[label=\textbf{AW\arabic*},ref=AW\arabic*,
  leftmargin=*]
\item\label{ax:AW1} \emph{(Grading, flow, and endpoints.)} The cup
  product is associative and graded-commutative with a unit $1\in V^0$,
  and $\Theta$ is a derivation:
  $\Theta_{i+k}(u\cup v)=\Theta_iu\cup v+u\cup\Theta_kv$. There are
  isomorphisms of spectral cores $V^0\cong U$, compatible with
  the multiplications and sending $1$ to $1$, and $V^{2d}\cong U(-d)$.
\item\label{ax:AW2} \emph{(Spectral realization of the divisor.)} The
  divisor identity
  \begin{equation}\label{eq:aw-div}
  \divisor Z=\sum_{j=0}^{2d}(-1)^{j+1}\operatorname{Sp}_j
  \tag{AW-div}
  \end{equation}
  holds.
\item\label{ax:AW3} \emph{(Duality.)} The trace is flow-equivariant into
  $\C(-d)$, that is, $\Tr(\Theta_{2d}x)=d\Tr(x)$, and for each $j$
  and $\rho$ the pairing $B_j$ restricts to a perfect pairing
  $V^j_{\Theta\sim\rho}\times V^{2d-j}_{\Theta\sim d-\rho}\to\C$.
\item\label{ax:AW4} \emph{(Positive compatible star.)} For each $j$,
  $\Theta_{2d-j}*_j=*_j\Theta_j+(d-j)*_j$, and $h_j$ is a positive
  definite Hermitian form on $V^j$.
\end{enumerate}
\end{definition}

Conditions \ref{ax:AW1}--\ref{ax:AW3} do not involve the stars, and we
use them also for tuples $(V^\bullet,\Theta,\cup,\Tr)$ without stars.
These are conditions on a single datum. A cohomology theory assigns
such data functorially to the objects of a source category, fixed in
Section~\ref{subsec:uniform-instances}.

It is natural to expect the same internal algebra axioms to apply to a
fixed positive-weight coefficient. The following obstruction shows why
the coefficient theory requires different product targets.

\begin{proposition}[Fixed-weight internal algebra obstruction]
\label{prop:weighted-algebra-obstruction}
Let $d\ge0$, $w\in\R$, and let $V^\bullet$ be a graded complex
algebra of spectral cores in degrees $0,\ldots,2d$. Suppose $\Theta$
is a derivation, $\Tr(\Theta x)=(d+w)\Tr(x)$, and antilinear maps
$*_j:V^j\to V^{2d-j}$ satisfy
\[
\Theta_{2d-j}*_j=*_j\Theta_j+(d-j)*_j,
\qquad h_j(u,v)=\Tr(u\cup *_jv).
\]
If every $h_j$ is positive definite Hermitian and $w\ne0$, then
every $V^j$ is zero. If the algebra is nonzero and unital, its unit
already forces $w=0$.
\end{proposition}

\begin{proof}
The derivation, trace, and star laws give
\[
h_j(\Theta u,v)+h_j(u,\Theta v)=(w+j)h_j(u,v).
\]
On each finite invariant generalized eigenspace,
$\Theta_j-(w+j)/2$ is skew-Hermitian and hence diagonalizable.
Thus $V^j$ is spanned by eigenvectors whose eigenvalues have real
part $(w+j)/2$. By the derivation law, a nonzero product of
eigenvectors in degrees $i,k$ has eigenvalue equal to the sum of
their eigenvalues. Its real part would be both
$w+(i+k)/2$ and $(w+i+k)/2$, forcing $w=0$. If $w\ne0$, all
products vanish, so all $h_j$ vanish; positivity gives $V^j=0$.
In the unital case $\Theta1=2\Theta1$, hence $\Theta1=0$.
The displayed identity in degree zero gives $w h_0(1,1)=0$.
\end{proof}

Deleting only the internal unit does not repair the positive-weight
axioms. Coefficient products instead take values in tensor-product
coefficients, and their pairings use the dual coefficient.
It is also natural to expect that moving the weight between the trace
and the star, or letting the unit carry it, avoids the obstruction; in
fact neither does. With trace law $\Tr(\Theta x)=\tau\Tr(x)$ and star
law $\Theta_{2d-j}*_j=*_j\Theta_j+\kappa_j*_j$, the computation in the
proof gives
$h_j(\Theta u,v)+h_j(u,\Theta v)=(\tau-\kappa_j)h_j(u,v)$, so only
the difference $\tau-\kappa_j$ matters. Keeping the intended weight
$w+j$ forces $\tau-\kappa_j=w+j$, and the unit argument applies
unchanged. A derivation of a unital algebra kills its unit, so the
unit cannot carry a nonzero weight without changing the derivation
rule of \ref{ax:AW1}.


These are the core cohomological requirements used for the RH
implication: algebra and flow, zeta realization, duality, and compatible
positivity. The endpoint normalization is Deninger's
\cite[(3), p.~168]{Deninger1998Analogies}; the trace, duality, and star
conditions follow his formulation and that of Flach--Morin
\cite[pp.~4--5]{Deninger2010HilbertPolya},
\cite[Conjecture~2.3(iii)--(iv), p.~522]
{FlachMorin2020DeningerWeilArakelov}.
Pullbacks are not required to preserve traces or stars, or to be
isometries. Absolute squares need not satisfy these axioms and may have
larger cohomology spaces (Theorem~\ref{thm:kun-square}).

The additional Hilbert-space and trace-formula requirements
\ref{ax:AW5}--\ref{ax:AW6} are recorded separately in
Definition~\ref{def:aw-analytic}. The regularized-determinant condition
AW-det is stronger than the divisor axiom \ref{ax:AW2}
(Definition~\ref{def:aw-det}). None is added to the hypotheses of the
purity implication below. The next definitions isolate that
implication's linear algebra and allow an arbitrary real weight offset;
the internal Weil datum above is restricted to weight zero.

\begin{definition}\label{def:inf-pol}
Let $V$ be a complex vector space, $\Theta:V\to V$ a linear map, and
$j\in\R$. A \emph{positive infinitesimal polarization of weight $j$} on
$(V,\Theta)$ is a positive definite Hermitian form $h$ on $V$ such that
\begin{equation}\label{eq:inf-pol}
h(\Theta u,v)+h(u,\Theta v)=j\,h(u,v)\qquad(u,v\in V).
\end{equation}
\end{definition}


\begin{definition}[Graded spectral realization]
\label{def:graded-realization}
Let $d\ge0$ be an integer, $w\in\R$, and $Z\not\equiv0$ a meromorphic
function on $\C$.
\begin{enumerate}[label=(\alph*)]
\item A \emph{graded spectral realization of type $(d,w)$}
of $Z$ consists of spectral cores $(V^j,\Theta_j)$, $0\le j\le2d$,
with
\[
 \divisor Z=\sum_{j=0}^{2d}(-1)^{j+1}
       \sum_{\rho\in\C}\dim V^j_{\Theta\sim\rho}[\rho].
\]
\item Such a realization is \emph{polarized} if each
$(V^j,\Theta_j)$ carries a positive
infinitesimal polarization of weight $w+j$.
\end{enumerate}
\end{definition}

The parameter $w$ records the intended weights; it imposes no spectral
restriction until a polarization is required. This definition keeps
only the hypotheses needed for zero location. The
datum of Definition~\ref{def:aw} adds cup products, duality, and a star;
its axioms \ref{ax:AW1}--\ref{ax:AW4} imply a polarized realization
of type $(d,0)$ by Proposition~\ref{prop:polarization-identity}.
Weighted coefficients use this spectral definition directly, without
an internal unit, cup product, or trace. A divisor identity is
weaker than a regularized determinant identity and contains no
normalization of the latter.

A Weil cohomology theory is a functor whose values carry the structure
the standard conjectures speak about. The same organization applies
here. The target category records the algebra, flow, and duality; the
positive star is a property, as for polarizable Hodge structures; and
the zeta realization is an axiom on the functor, because it refers to
the prescribed function of the source object.

\begin{definition}[Flow spaces]\label{def:flow-category}
$\Flow$ is the category whose objects are spectral cores $(V,\Theta)$
(Definition~\ref{def:spectral-core}) and whose morphisms
$f:(V,\Theta_V)\to(W,\Theta_W)$ are $\C$-linear maps with
$f\Theta_V=\Theta_Wf$.
\end{definition}

For $n\in\Z$ the twist $V(n)$ has generator $\Theta-n$. The arithmetic
unit $\mathbb 1$ and the periodic units $\mathbb 1_q$ of
Definition~\ref{def:unit} are objects of $\Flow$ with
multiplications.

\begin{proposition}[Exactness and additivity of divisors]
\label{prop:flow-exact}
Let $f:(V,\Theta_V)\to(W,\Theta_W)$ be a morphism in $\Flow$, and let
$\cdots\to V'^j\to V^j\to V''^j\to V'^{j+1}\to\cdots$ be a long exact
sequence in $\Flow$ with $V'^j=V^j=V''^j=0$ for $|j|$ large.
\begin{enumerate}[label=(\alph*)]
\item\label{item:flow-kernels} $\ker f$ and $\operatorname{coker}f$,
  with the induced generators, are spectral cores, and for each
  $\rho\in\C$ the functor $V\mapsto V_{\Theta\sim\rho}$ is exact on
  $\Flow$.
\item\label{item:flow-additive}
  $\sum_j(-1)^j\operatorname{Sp}(V^j)
   =\sum_j(-1)^j\operatorname{Sp}(V'^j)
   +\sum_j(-1)^j\operatorname{Sp}(V''^j)$.
\end{enumerate}
\end{proposition}

\begin{proof}[Proof of \eqref{item:flow-kernels}]
A linear map commuting with the generators preserves generalized
eigenspaces, so $f$ is the direct sum of its restrictions
$V_{\Theta\sim\rho}\to W_{\Theta\sim\rho}$. Kernels and cokernels are
computed blockwise. Each block is finite-dimensional, and the spectrum
of $\ker f$ or $\operatorname{coker}f$ lies in that of $V$ or $W$, so
it has no finite accumulation point.
\end{proof}

\begin{proof}[Proof of \eqref{item:flow-additive}]
By \eqref{item:flow-kernels} the restriction of the sequence to the
$\rho$-blocks is a bounded exact sequence of finite-dimensional
spaces, so the alternating sum of its dimensions vanishes. Comparing
coefficients of $[\rho]$ gives the identity.
\end{proof}

It is natural to expect $\Flow$ to be symmetric monoidal under
$V\otimes W$ with generator $\Theta\otimes1+1\otimes\Theta$, as the
target of a Weil cohomology theory is. In fact the algebraic tensor
product of two spectral cores need not be a spectral core
(Theorem~\ref{thm:kun-square}), so no tensor product is part of the
structure. The relative tensor product over $\mathbb 1_q$ is the
correct product on the finite-field objects
(Proposition~\ref{prop:kun-Fq-C} and Theorem~\ref{thm:kun-Fq}).

\begin{definition}[Deninger structures]
\label{def:deninger-structure}
Let $d\ge0$ be an integer and let $U$ be $\mathbb 1$ or some
$\mathbb 1_q$.
\begin{enumerate}[label=(\alph*)]
\item A \emph{Deninger structure of dimension $d$ with unit $U$} is a
  tuple $D=(V^\bullet,\Theta,\cup,\Tr)$ of spectral cores
  $(V^j,\Theta_j)$ for $0\le j\le2d$, bilinear products
  $\cup:V^i\times V^k\to V^{i+k}$, zero if $i+k>2d$, and a linear
  trace $\Tr:V^{2d}\to\C$, satisfying \ref{ax:AW1} and \ref{ax:AW3}
  for $U$.
\item A \emph{morphism} of Deninger structures is a graded unital
  algebra homomorphism commuting with the generators.
\item $\Den_d(U)$ is the category of Deninger structures of dimension
  $d$ with unit $U$. The category $\Den^\natural$ has all Deninger
  structures as objects and all morphisms of (b), including maps
  between different pairs $(d,U)$. Thus $\Den_d(U)$ is a full
  subcategory of $\Den^\natural$.
\end{enumerate}
\end{definition}

Explicitly, the product is associative, graded-commutative, and
unital, $\Theta$ is a derivation, $V^0\cong U$ compatibly with the
multiplications and units, $V^{2d}\cong U(-d)$, the trace is
flow-equivariant into $\C(-d)$, and the duality pairings are perfect
on generalized eigenspaces. The divisor condition \ref{ax:AW2} is
omitted because it refers to an external function, and the star of
\ref{ax:AW4} because its existence is a property. Morphisms need not
preserve traces.

\begin{proposition}[Disjoint-union obstruction]
\label{prop:den-disjoint-obstruction}
Let $\mathcal C$ be a category with an assignment of a pair
$(d_X,U_X)$ to each object $X$. If $\mathcal C$ has a morphism
$f:X\to Y$ with $(d_X,U_X)\ne(d_Y,U_Y)$, there is no functor
$H:\mathcal C^{\mathrm{op}}\to\coprod_{d,U}\Den_d(U)$ for which
$H(Z)\in\Den_{d_Z}(U_Z)$ for every object $Z$.
\end{proposition}

\begin{proof}
The image of $f$ would be a morphism $H(Y)\to H(X)$ between distinct
components of the disjoint union, but that category has no such
morphism.
\end{proof}

\begin{definition}[Polarizable Deninger structures]
\label{def:polarizable}
Let $D=(V^\bullet,\Theta,\cup,\Tr)$ be a Deninger structure of
dimension $d$ with unit $U$.
\begin{enumerate}[label=(\alph*)]
\item A \emph{polarization} of $D$ is a family of antilinear maps
  $*_j:V^j\to V^{2d-j}$ such that the Weil datum
  $(V^\bullet,\Theta,\cup,\Tr,*)$ satisfies \ref{ax:AW4}.
\item $D$ is \emph{polarizable} if it has a polarization.
\item $\Den^{\mathrm{pol}}_d(U)$ is the full subcategory of
  $\Den_d(U)$ of polarizable Deninger structures.
\end{enumerate}
\end{definition}

A polarized Deninger structure is a Weil datum satisfying
\ref{ax:AW1}, \ref{ax:AW3}, and \ref{ax:AW4}. When it also satisfies
\ref{ax:AW2} for a function $Z$, it is a polarized graded spectral
realization of $Z$ of type $(d,0)$
(Proposition~\ref{prop:polarization-identity}), and
Theorem~\ref{thm:main-general} below applies.

It is natural to expect morphisms of polarizable structures to
preserve the forms $h_j$. In fact, under the compatibility of traces
with cycle classes expected in
\cite[(2)--(5), pp.~4--5]{Deninger2025BSDDedekind}, the pullback along
the map of Arakelov compactifications induced by a finite extension
$L/K$ of number fields multiplies the degree-one form by $[L:K]$
\cite[(8), p.~5]{Deninger2025BSDDedekind}. Polarizability is therefore
a property of objects, as for polarizable Hodge structures, and
$\Den^{\mathrm{pol}}_d(U)$ is a full subcategory.

\begin{definition}[Deninger modules]\label{def:deninger-module}
Let $D=(V^\bullet,\Theta,\cup,\Tr)$ be a Deninger structure.
\begin{enumerate}[label=(\alph*)]
\item A \emph{Deninger module over $D$} is a family
  $(M^j,\Theta_j)_{j\in\Z}$ of spectral cores, zero for $|j|$ large,
  with an associative graded action $V^i\times M^k\to M^{i+k}$ of the
  algebra $V^\bullet$ for which $\Theta$ is a derivation.
\item $\Mods D$ is the category of Deninger modules over $D$, with
  graded $V^\bullet$-linear maps commuting with the generators as
  morphisms.
\end{enumerate}
\end{definition}

Deninger modules carry compactly supported cohomology of non-proper
objects and coefficient cohomology of weight $w\ne0$. No unit, trace,
or duality is required. A fixed positive-weight coefficient cannot
carry the unital structure of Definition~\ref{def:deninger-structure}
(Proposition~\ref{prop:weighted-algebra-obstruction}); its products
and dualities take values in tensor-product and dual coefficients, as
in Remark~\ref{rem:A-motivic}.

\begin{theorem}[General cohomological purity implication]

\label{thm:main-general}
Let $d\ge0$ be an integer, $w\in\R$, and $Z\not\equiv0$ meromorphic
on $\C$. Suppose that $Z$ has a polarized graded spectral realization
of type $(d,w)$. Then each $\Theta_j$ is diagonalizable, every one of
its eigenvalues has real part $(w+j)/2$, and
\[
 \ord_\rho Z=
 \begin{cases}
 (-1)^{j+1}\dim\ker(\Theta_j-\rho),
       &2\Re\rho-w=j\in\{0,\ldots,2d\},\\
 0,    &2\Re\rho-w\notin\{0,\ldots,2d\}.
 \end{cases}
\]
Thus zeros occur only on the odd-degree lines, poles only on the
even-degree lines, and the divisor of $Z$ determines each degreewise
spectral divisor without cancellation between degrees.
\end{theorem}

\begin{proof}
If $\Theta_jv=\rho v$ and $v\ne0$, the polarization identity gives
$(\rho+\bar\rho-w-j)h_j(v,v)=0$. Positivity therefore gives
$\Re\rho=(w+j)/2$. On every finite-dimensional invariant subspace,
$\Theta_j-(w+j)/2$ is skew-Hermitian and hence diagonalizable.
Applying this to every generalized eigenspace proves diagonalizability
on the algebraic direct sum. The lines $(w+j)/2$ are distinct, so at
each point the divisor identity has at most one nonzero summand.
\end{proof}

The theorem concerns a single function and its realization; its proof
uses no source geometry.

This is Deninger's infinitesimal version of Serre's argument

\cite[(7.11)--(7.13), pp.~32--33]{Deninger1994Motivic}, stated for
arbitrary real weight offset and without unnecessary geometric
hypotheses. Theorem~\ref{thm:A-schemes} gives its cup-product and
star formulation for $w=0$. Higher-dimensional arithmetic zeta
functions are not expected to have just one critical line:
Proposition~\ref{prop:P1Z} proves that already
$\zeta_{\mathbb{P}^1_\Z}(s)=\zeta(s)\zeta(s-1)$ has zeros in two
disjoint strips.

\begin{definition}\label{def:zeta}
Let $\zeta$ be the Riemann zeta function and let $K$ be a number field
with $r_1$ real places, $r_2$ complex places, discriminant $D_K$, and
Dedekind zeta function $\zeta_K$.
\begin{enumerate}[label=(\alph*)]
\item The \emph{archimedean gamma factors} are
  $\GR(s)=\pi^{-s/2}\Gamma(s/2)$ and $\GC(s)=2(2\pi)^{-s}\Gamma(s)$.
\item The \emph{completed zeta function} is
  $\Lambda(s)=\GR(s)\zeta(s)$.
\item The \emph{Riemann $\xi$-function} is
  $\xi(s)=\tfrac12 s(s-1)\Lambda(s)$.
\item The \emph{Riemann $\Xi$-function} is $F(z)=\Xi(z)=\xi(\tfrac12-iz)$.
\item The \emph{completed Dedekind zeta function} is
  $\Lambda_K(s)=|D_K|^{s/2}\GR(s)^{r_1}\GC(s)^{r_2}\zeta_K(s)$.
\item The \emph{Dedekind $\xi$-function} is
  $\xi_K(s)=s(s-1)\Lambda_K(s)$, and its zeros are the
  \emph{nontrivial zeros} of $\zeta_K$.
\end{enumerate}
\end{definition}

\begin{proposition}[Positive realization of the zeta zeros]
\label{prop:rh-positive-zero}
Let $\xi$ be the completed Riemann zeta function of
Definition~\ref{def:zeta}, and let $(V,\Theta)$ be a spectral core.
Suppose that, for every $\rho\in\C$,
\[
\dim V_{\Theta\sim\rho}=\ord_\rho\xi,
\]
and let $h$ be a positive infinitesimal polarization of weight one
on $(V,\Theta)$ (Definition~\ref{def:inf-pol}). Then RH holds.
\end{proposition}

\begin{proof}
For each zero $\rho$ of $\xi$, the multiplicity hypothesis gives a
nonzero finite-dimensional generalized $\rho$-eigenspace, which
contains an eigenvector $v\ne0$. Substituting $u=v$ into
\eqref{eq:inf-pol} yields
$(2\Re\rho-1)h(v,v)=0$. Since $h(v,v)>0$, one has
$\Re\rho=\tfrac12$.
\end{proof}

The multiplicity hypothesis is the exact zero-spectrum requirement;
positivity alone does not identify any eigenvalue with a zeta zero.
The form must satisfy \eqref{eq:inf-pol} for the same generator
\(\Theta\); an arbitrary positive inner product is insufficient
(Proposition~\ref{prop:content}\ref{item:ct-pos}).

\subsection{Geometric instances and the uniform conjectures}
\label{subsec:uniform-instances}

\begin{definition}[Arakelov compactification]
\label{def:arakelov-pair}
Let $X$ be a regular projective flat scheme over $\Spec\Z$. Its
\emph{Arakelov compactification} is the arithmetic pair
$\Xbar=(X,X_\infty)$, where $X_\infty$ consists of $X(\C)$ with
its complex-conjugation action $F_\infty$ and a K\"ahler form
$\omega$ satisfying $F_\infty^*\omega=-\omega$.
The dimension of $\Xbar$ is the Krull dimension of $X$.
\end{definition}

This is the convention of
\cite[Notation~2.2, pp.~520--522]{FlachMorin2020DeningerWeilArakelov}.
We use morphisms induced by morphisms of the underlying arithmetic
schemes and their conjugation-equivariant maps at infinity; these
morphisms are not required to preserve the chosen metrics. In particular, for a number field
$K$, the object $\Xbar_K=(\Spec\cO_K,X_{K,\infty})$ has dimension one.
The pair is not an ordinary scheme, and no product over an absolute
base is included in this convention. The broader framework of
\cite[Conjecture~2.3, p.~522]{FlachMorin2020DeningerWeilArakelov}
also considers separated finite-type arithmetic schemes with compact
supports. Positivity below concerns the projective regular objects.

\begin{definition}[Compactified schemes]
\label{def:compactified-schemes}\label{def:projective-inputs}
\begin{enumerate}[label=(\alph*)]
\item A \emph{compactified scheme} $\Xbar$ is a pair
  $(X,X_\infty)$, where $X$ is a separated scheme of finite type over
  $\Spec\Z$ and $X_\infty$ is the set $X(\C)$ with its complex conjugation
  $F_\infty$, together with, when $X$ is projective, flat, and regular
  over $\Spec\Z$, a K\"ahler form $\omega$ on $X(\C)$ with
  $F_\infty^*\omega=-\omega$.
\item $\Schbar$ is the category of compactified schemes. A morphism
  $\Xbar\to\overline Y$ is a morphism of schemes $X\to Y$ together
  with the induced $F_\infty$-equivariant map $X(\C)\to Y(\C)$; the
  K\"ahler forms are not required to be preserved.
\item $\Schbar^{\mathrm{reg}}$ is the full subcategory of objects with
  $X$ regular of pure dimension.
\item $\Schbar^{\mathrm{proj}}$ is the full subcategory of objects with
  $X$ projective, flat, and regular over $\Spec\Z$, or with $X$ smooth
  and projective over a finite field.
\item\label{item:schbar-one} $\Schbar_1$ is the full subcategory of
  $\Schbar^{\mathrm{proj}}$
  whose objects are the finite disjoint unions of the
  $\Xbar_K=(\Spec\cO_K,X_{K,\infty})$ for number fields $K$.
\item For a prime power $q$, $\SmProj_{\F_q}$ is the full subcategory
  of $\Schbar^{\mathrm{proj}}$ of smooth projective varieties over
  $\F_q$, and $\SmProj_{\F}$ is the full subcategory of smooth
  projective varieties over finite fields.
\item For $\Xbar\in\Schbar^{\mathrm{proj}}$, put $d_X=\dim X$ and
  $w_X=0$. The \emph{prescribed function} $Z_X$ is the completed zeta
  function $\zeta(X,s)\zeta(X_\infty,s)$ of
  \cite[pp.~520--521]{FlachMorin2020DeningerWeilArakelov} if $X$ is
  flat over $\Spec\Z$, and $\zeta_X(s)$ if $X$ is a variety over a
  finite field.
\item For connected $\Xbar\in\Schbar^{\mathrm{proj}}$, the \emph{unit}
  $U_X$ is $\mathbb 1$ if $X$ is flat over $\Spec\Z$, and
  $\mathbb 1_{q}$ if $X$ is a variety over a finite field with
  $H^0(X,\cO_X)\cong\F_{q}$.
\end{enumerate}
\end{definition}

For a variety over a finite field $X(\C)$ is empty, so its
compactified scheme is $(X,\varnothing)$. The objects $\Xbar_K$ are
finite flat over $\Spec\Z$; curves over finite fields are not objects
of $\Schbar_1$. The bar in $\Schbar$ indicates the compactification,
as it does in $\Xbar$, while $\Sch$ denotes the category of schemes.
The connected objects of $\Schbar^{\mathrm{proj}}$ are the projective
regular class of
\cite[Notation~2.2, pp.~521--522]{FlachMorin2020DeningerWeilArakelov},
with the morphisms made explicit, together with the smooth projective
varieties over finite fields.

\begin{remark}[Prescribed data and ambient geometries]
\label{rem:source-data}
The data $d_X$, $w_X$, $Z_X$, and $U_X$ of
Definition~\ref{def:compactified-schemes} are attached to the objects
of $\Schbar^{\mathrm{proj}}$. The meromorphic continuation of every
$Z_X$ is part of Conjecture~\ref{conj:arithmetic-realization} below.
The weight offset is zero for every object, as
Proposition~\ref{prop:weighted-algebra-obstruction} forces for unital
data; coefficients of nonzero weight use graded spectral realizations
(Definition~\ref{def:graded-realization}) and Deninger modules
(Definition~\ref{def:deninger-module}). A proposed geometry $\mathcal G$
over an absolute point enters through a comparison functor
$\Schbar\to\mathcal G$ together with its cohomological comparisons.
Agreement of zeta functions alone is not such a comparison: a
degreewise comparison must intertwine the generators and any pairings
it transports. Products in $\mathcal G$ need not carry Deninger
structures, since the tensor square of the Weil datum of an arithmetic
curve is not a spectral core (Theorem~\ref{thm:kun-square}). In
particular a Connes--Consani model for $\Spec\Z$ has the same target
function $\Lambda$ as the Arakelov model; the arithmetic and scaling
sites are semiringed topoi, not alternative definitions of an Arakelov
pair \cite[\S1, pp.~1--2]{ConnesConsani2015ScalingSite}.
\end{remark}

\begin{definition}[Deninger cohomology theory]
\label{def:deninger-theory}
For connected $\Xbar\in\Schbar^{\mathrm{proj}}$ let $d_X$, $Z_X$, and
$U_X$ be as in Definition~\ref{def:compactified-schemes}.
\begin{enumerate}[label=(\alph*)]
\item\label{item:den-core} A \emph{Deninger cohomology theory} is a
  functor $H:(\Schbar^{\mathrm{proj}})^{\mathrm{op}}\to\Den^\natural$ with
  $H(\Xbar)\in\Den_{d_X}(U_X)$ for every connected $\Xbar$, satisfying
  \ref{ax:F1} and \ref{ax:F2} below.
\item\label{item:den-full} A \emph{full Deninger cohomology theory} is
  a Deninger cohomology theory $H$ together with an extension $H_c$ to
  $\Schbar$, contravariant for proper morphisms and covariant for open
  immersions, whose value on $\Xbar$ is a Deninger module over
  $H(\Xbar_\Q)$ or over $H(\Spec\F_q)$, satisfying \ref{ax:F3} and
  \ref{ax:F4} below.
\item\label{item:den-pol} A Deninger cohomology theory, full or not,
  is \emph{polarized} if $H(\Xbar)$ is polarizable
  (Definition~\ref{def:polarizable}) for every connected
  $\Xbar\in\Schbar^{\mathrm{proj}}$.
\end{enumerate}
The axioms are the following.
\begin{enumerate}[label=\textbf{F\arabic*},ref=F\arabic*,leftmargin=*]
\item\label{ax:F1} \emph{(Endpoints.)} For connected $\Xbar$, the
  isomorphisms $H^0(\Xbar)\cong U_X$ and
  $H^{2d_X}(\Xbar)\cong U_X(-d_X)$ of
  Definition~\ref{def:deninger-structure} hold.
\item\label{ax:F2} \emph{(Zeta realization.)} For connected
  $\Xbar\in\Schbar^{\mathrm{proj}}$,
  $\divisor Z_X=\sum_j(-1)^{j+1}\operatorname{Sp}(H^j(\Xbar))$.
\item\label{ax:F3} \emph{(Localization.)} For a closed subscheme
  $Y\subset X$ with open complement $W$, there is a long exact sequence
  $\cdots\to H^j_c(W)\to H^j_c(X)\to H^j_c(Y)\to H^{j+1}_c(W)
  \to\cdots$ in $\Flow$.
\item\label{ax:F4} \emph{(Galois descent.)} For a finite Galois
  extension $L/K$ of number fields with group $G$, pullback along
  $\Xbar_L\to\Xbar_K$ induces
  $H^\bullet(\Xbar_K)\cong H^\bullet(\Xbar_L)^G$ compatibly with the
  generators and products.
\end{enumerate}
\end{definition}

The structure map $\mathbb P^1_{\F_q}\to\Spec\F_q$ is a morphism in
$\Schbar^{\mathrm{proj}}$ between connected objects of dimensions
$1$ and $0$. Proposition~\ref{prop:den-disjoint-obstruction} therefore
rules out the disjoint union of the $\Den_d(U)$ as the target of such a
theory. The cross-component morphisms in $\Den^\natural$ remove this
definitional obstruction.

A full Deninger cohomology theory is in particular a Deninger
cohomology theory. Axiom~\ref{ax:F1} restates the endpoint clause of
\ref{ax:AW1}, and \ref{ax:F2} is \ref{ax:AW2} for $(Z_X,U_X)$. Hence
the value $H(\Xbar)$ on a connected object is a datum without stars
satisfying \ref{ax:AW1}--\ref{ax:AW3} for $(Z_X,U_X)$, pullbacks
preserve cup products and commute with the generators, and a
polarization of $H(\Xbar)$ supplies stars satisfying \ref{ax:AW4}.
By Theorem~\ref{thm:main-general}, a polarized theory places the zeros
of each odd-degree determinant on $\Re s=j/2$; this is
\cite[Conjecture~2.3(v), p.~522]{FlachMorin2020DeningerWeilArakelov}.
Axiom~\ref{ax:F4} is the descent condition of
Definition~\ref{def:descent}. The compactly supported version of
\ref{ax:F2} requires archimedean factors for non-proper objects, which
are not specified here.

\begin{proposition}[Relative structure]
\label{prop:relative-structure}
Let $H$ be a Deninger cohomology theory, $K$ a number field, and
$\pi:\Xbar\to\Xbar_K$ a morphism in $\Schbar^{\mathrm{proj}}$ with
$\Xbar$ connected. Then $H(\pi)$ makes $H(\Xbar)$ a graded
$H(\Xbar_K)$-algebra, and $\Theta$ is a derivation for this action.
\end{proposition}

\begin{proof}
By functoriality $H(\pi):H(\Xbar_K)\to H(\Xbar)$ is a morphism in
$\Den^\natural$, a graded unital algebra homomorphism commuting with the
generators. Composing it with the cup product of $H(\Xbar)$ gives the
action, and $\Theta$ is a derivation of that cup product by
\ref{ax:AW1}.
\end{proof}

The proposition records what ``relative over $\cO_K$'' means for a
theory on $\Schbar^{\mathrm{proj}}$; it is a consequence of
functoriality and not an additional axiom.

\begin{proposition}[Finite-field calibration]
\label{prop:finite-field-calibration}
Let $H$ be a Deninger cohomology theory, and let $X$ be a smooth
projective geometrically connected variety of dimension $d$ over
$\F_q$. Let $\ell\nmid q$ be a prime, fix a field isomorphism
$\overline\Q_\ell\cong\C$, let
$H^j_\ell(X)=H^j(X_{\overline\F_q},\Q_\ell)\otimes\C$ with geometric
Frobenius $\mathrm F$, and let
$(\widehat H^j_\ell(X),\Theta_j)$ be its logarithmic suspension: the algebraic
direct sum, over all $\rho\in\C$ with $q^\rho$ an eigenvalue of
$\mathrm F$, of copies of the generalized eigenspace
$H^j_\ell(X)_{(q^\rho)}$, on which $\Theta_j$ is $\rho+N/\log q$ with
$N=\log(q^{-\rho}\mathrm F)$ nilpotent
(Definition~\ref{def:log-susp}).
\begin{enumerate}[label=(\roman*)]
\item\label{item:cal-curve} If $d=1$, then
  $\operatorname{Sp}(H^1(X))=\operatorname{Sp}(\widehat H^1_\ell(X))$.
\item\label{item:cal-pol} If $H$ is polarized, then for every $j$,
  $\operatorname{Sp}(H^j(X))=\operatorname{Sp}(\widehat H^j_\ell(X))$, the
  generator of $H^j(X)$ is diagonalizable, and $H^j(X)\cong\widehat
  H^j_\ell(X)$ in $\Flow$ if and only if $\mathrm F$ is semisimple on
  $H^j_\ell(X)$.
\end{enumerate}
\end{proposition}

\begin{proof}[Proof of \eqref{item:cal-curve}]
By Theorem~\ref{thm:B}\ref{item:B-det},
$\divisor\zeta_X=\sum_j(-1)^{j+1}\operatorname{Sp}(\widehat H^j_\ell(X))$.
Here $U_X=\mathbb 1_q$. The suspension of $(H^0_\ell(X),\mathrm F)=(\C,1)$ is
$\mathbb 1_q$, and that of $(H^2_\ell(X),\mathrm F)=(\C,q)$ has eigenvalues
$1+2\pi in/\log q$, $n\in\Z$, each simple, so it is $\mathbb 1_q(-1)$.
By \ref{ax:F1}, $H^0(X)\cong\mathbb 1_q$ and $H^2(X)\cong\mathbb
1_q(-1)$, so the degree-zero and degree-two terms agree on the two
sides, and \ref{ax:F2} gives
$\operatorname{Sp}(H^1(X))=\operatorname{Sp}(\widehat H^1_\ell(X))$.
\end{proof}

\begin{proof}[Proof of \eqref{item:cal-pol}]
A polarization of $H(X)$ and \ref{ax:F2} make
$(H^j(X),\Theta_j)_j$ a polarized graded spectral realization of
$\zeta_X$ of type $(d,0)$
(Proposition~\ref{prop:polarization-identity}). By
Theorem~\ref{thm:main-general} each generator is diagonalizable, and
$\dim\ker(\Theta_j-\rho\mid H^j(X))=(-1)^{j+1}\ord_\rho\zeta_X$ when
$2\Re\rho=j$, while $H^j(X)$ has no other eigenvalues. On the other
side, $\widehat H^j_\ell(X)_{\Theta\sim\rho}$ is the generalized
$q^\rho$-eigenspace of $\mathrm F$ on $H^j_\ell(X)$. By Deligne's theorem
\cite[Th\'eor\`eme~(1.6), p.~276]{Deligne1974WeilI} every eigenvalue
$\alpha$ of $\mathrm F$ on $H^j_\ell(X)$ has $|\alpha|=q^{j/2}$, so
$\widehat H^j_\ell(X)$ has spectrum on $\Re\rho=j/2$, and in the identity of
Theorem~\ref{thm:B}\ref{item:B-det} only the degree $j=2\Re\rho$
contributes at $\rho$. Hence
$\dim\widehat H^j_\ell(X)_{\Theta\sim\rho}=(-1)^{j+1}\ord_\rho\zeta_X
=\dim H^j(X)_{\Theta\sim\rho}$ for all $j$ and $\rho$. Two spectral
cores with diagonalizable generators and the same multiplicities are
isomorphic in $\Flow$. The generator of $\widehat H^j_\ell(X)$ is
diagonalizable if and only if every nilpotent part $N$ vanishes, that
is, if and only if $\mathrm F$ is semisimple on $H^j_\ell(X)$.
\end{proof}

It is natural to expect the comparison with $\ell$-adic cohomology to
be an axiom, as in a Weil cohomology theory. In fact the spectral part
of it follows from \ref{ax:F1} and \ref{ax:F2} for curves, and from
polarization in every dimension, by
Proposition~\ref{prop:finite-field-calibration}. The isomorphism there
is obtained degree by degree from multiplicities; it is neither
natural in $X$ nor compatible with cup products, and it is an
isomorphism with $\widehat H^j_\ell(X)$ only when Frobenius is semisimple,
which is not known in general. A natural comparison compatible with
products remains a separate requirement, stated in
Conjecture~\ref{conj:arithmetic-realization} below.

\begin{proposition}[The divisor does not determine degreewise spectra]
\label{prop:calibration-cancel}
Let $d\ge2$, let $U$ be $\mathbb 1$ or some $\mathbb 1_q$, let
$D=(V^\bullet,\Theta,\cup,\Tr)\in\Den_d(U)$, and let $\rho\in\C$ and
$1\le i\le2d-2$. There is $D'=(V'^\bullet,\Theta,\cup,\Tr)$ in
$\Den_d(U)$ with
$\sum_j(-1)^{j+1}\operatorname{Sp}(V'^j)
=\sum_j(-1)^{j+1}\operatorname{Sp}(V^j)$ and
$\operatorname{Sp}(V'^i)=\operatorname{Sp}(V^i)+\operatorname{Sp}(W)$,
where $W$ is $U$ with generator $\Theta_U+\rho$.
\end{proposition}

\begin{proof}
Let $W'$ be $U$ with generator $\Theta_U+d-\rho$, and write
$w_u\in W$, $w'_u\in W'$ for the element $u\in U$. Fix the isomorphism
$\phi:U\cong V^0$ of \ref{ax:AW1}. The pairing $B_0$ restricted to
$V^0_{\Theta\sim0}\times V^{2d}_{\Theta\sim d}$ is perfect and
$V^0_{\Theta\sim0}=\C\cdot1$, so there is $t\in V^{2d}$ with
$\Theta t=dt$ and $\Tr(t)=1$. For $k\in\{i,i+1\}$ add a copy $W_k$ of
$W$ in degree $k$ and a copy $W'_k$ of $W'$ in degree $2d-k$; these
degrees lie in $[1,2d-1]$. Let $V^0$ act on $W_k$ and $W'_k$ through
$\phi^{-1}$ and the multiplication of $U$, set
$w_u\cup w'_{u'}=\phi(uu')\cup t$ for $w_u\in W_k$,
$w'_{u'}\in W'_k$, with the graded-commutative sign for the opposite
order, and let all other products involving the added summands
vanish. The trace is unchanged. The derivation rule holds because the
shifts $\rho$ and $d-\rho$ add to the eigenvalue $d$ of $t$, and
associativity holds because every product of an added element with an
element of positive degree vanishes or lands in $V^{2d}$, whose
products with elements of positive degree vanish. On
generalized eigenspaces, $W_\sigma=U_{\sigma-\rho}$ pairs with
$W'_{d-\sigma}=U_{\rho-\sigma}$ by $(u,u')\mapsto\Tr(\phi(uu')\cup
t)$, which is perfect because the multiplication pairs
$U_\alpha$ perfectly with $U_{-\alpha}$ into $U_0=\C\cdot1$. The
endpoints are unchanged. The added divisor is
$(-1)^{i+1}\operatorname{Sp}(W)+(-1)^{i+2}\operatorname{Sp}(W)
+(-1)^{2d-i+1}\operatorname{Sp}(W')+(-1)^{2d-i}\operatorname{Sp}(W')=0$.
\end{proof}

It is natural to expect \ref{ax:F2} to determine the spectrum of each
$H^j(X)$ for a variety $X$ over $\F_q$, as it does for curves. In fact
for $d\ge2$ it determines only the alternating sum: for a surface,
Proposition~\ref{prop:calibration-cancel} with $i=1$ adds the same
spectrum to degrees one and two, and the same dual spectrum to degrees
two and three, without changing the divisor. The structure $D'$ is not
polarizable, since a polarization would place the spectrum of $W$ on
two different lines (Theorem~\ref{thm:main-general}); this is why
part~\eqref{item:cal-pol} of
Proposition~\ref{prop:finite-field-calibration} needs polarization.

\begin{remark}[Terminology and scope]
\label{rem:deninger-terminology}
\emph{Deninger cohomology theory} means
Definition~\ref{def:deninger-theory}, and the \emph{Deninger axioms}
\ref{ax:AW1}--\ref{ax:AW4} of Definition~\ref{def:weil-cohomology} are
conditions on a single datum. This is our terminology, not a claim of
existence or uniqueness. Flach--Morin denote Deninger's conjectural
groups by $H^\bullet_{\mathrm{dyn},c}$ and
$H^\bullet_{\mathrm{dyn}}$; their separately constructed fixed-twist
groups are called \emph{Weil--Arakelov cohomology}
\cite[abstract, Conjecture~2.3, and \S3]%
{FlachMorin2020DeningerWeilArakelov}. We call these fixed-twist
groups \emph{Flach--Morin cohomology} in these notes. This is a local
shorthand for their Weil--Arakelov cohomology, not the name used in
the cited paper. Their proposed comparison
\cite[Conjecture~4.1, p.~527]{FlachMorin2020DeningerWeilArakelov}
makes them candidate fixed-twist data for a Deninger cohomology
theory; it does not make them an instance of
Definition~\ref{def:deninger-theory}. The divisor condition
\ref{ax:F2} determines alternating spectral multiplicities; it does
not specify the functorial maps, cup products, or positive stars. For
a dimension-one object with prescribed endpoints, it does fix the
degree-one spectral divisor. This is a spectral constraint, not a
uniqueness theorem for the cohomology theory.
\end{remark}

The common substitutions in Theorem~\ref{thm:main-general} are as follows.
The table names the required realization, not its existence in every row.
\begin{center}\small
\begin{tabular}{@{}p{0.28\textwidth}p{0.23\textwidth}p{0.41\textwidth}@{}}
\toprule
Input & $(d_X,w_X)$ and $Z_X$ & Realization and consequence\\
\midrule
Number field $K$ & $(1,0)$; $\Lambda_K$ &
Arakelov or a specified absolute model; zeros on $\Re s=1/2$.\\
Smooth projective $X/\F_q$ & $(\dim X,0)$; $\zeta_X$ &
Logarithmic suspension of Frobenius; degree-$j$ line $\Re s=j/2$.\\
Regular projective flat $X/\Z$ & $(\dim X,0)$; completed $\zeta_X$ &
Conjectural global realization; odd-degree zero lines and even-degree
pole lines.\\
Pure weight-$w$ coefficients over a number field & $(1,w)$;
completed $L(M,s)$ &
Requires a polarized coefficient spectral realization; middle line
$\Re s=(w+1)/2$.\\
\bottomrule
\end{tabular}
\end{center}

For finite fields the suspension uses \emph{all} logarithms of each
Frobenius eigenvalue: $q^\rho=\alpha$ and
$\Re\rho=\log|\alpha|/\log q$. The unit is the periodic object
$\mathbb 1_q$ of Definition~\ref{def:unit}, rather than the arithmetic
unit $(\C,0)$ (Proposition~\ref{prop:unit}). Theorem~\ref{thm:B}
supplies the structural comparison; full positive polarization also
requires Frobenius semisimplicity. A global function field is the
curve case of that row. The coefficient row is conditional on a
specified coefficient category and analytic continuation; it does not
extend the number-field theorem to all automorphic $L$-functions.
The arithmetic squares of Section~\ref{sec:kunneth} require the larger
flow-module formalism and are not additional spectral-core objects in
this table.


The following conjectures ask for a Deninger cohomology theory
(Definition~\ref{def:deninger-theory}), in a core and a full version,
together with a finite-field comparison. Existence, zeta
identification, and duality are separated from positivity. A
realization through another source geometry must supply an input
functor and these same cohomological structures.



\begin{conjecture}[Uniform cohomological realization]
\label{conj:arithmetic-realization}
The prescribed functions $Z_X$ on $\Schbar^{\mathrm{proj}}$ admit
meromorphic continuation to $\C$, and there is a Deninger cohomology
theory $H$ whose restriction to each $\SmProj_{\F_q}$ is naturally
isomorphic, with generators and products, to the logarithmic
suspension of a chosen complex realization of $\ell$-adic cohomology,
$\ell\nmid q$.
\end{conjecture}

\begin{conjecture}[Full uniform cohomological realization]
\label{conj:arithmetic-realization-full}
The statement of Conjecture~\ref{conj:arithmetic-realization} holds
with $H$ a full Deninger cohomology theory.
\end{conjecture}

For $\Xbar=\Xbar_K$ the endpoint data are $(H^0,\Theta_0)=(\C,0)$ and
$(H^2,\Theta_2)=(\C,1)$, by \ref{ax:F1} with $U_X=\mathbb 1$. The
natural comparison with $\ell$-adic cohomology is part of the
conjecture rather than of Definition~\ref{def:deninger-theory}; its
spectral content is a consequence of the axioms in the cases of
Proposition~\ref{prop:finite-field-calibration}.



The cohomology, trace formula, duality, and determinant expectations
underlying Conjecture~\ref{conj:arithmetic-realization} are
\cite[Conjecture~2, p.~3]{Deninger2010HilbertPolya} and
\cite[Conjecture~2.3(i)--(iii), p.~522]
{FlachMorin2020DeningerWeilArakelov}. The completed function includes
Serre's archimedean factors of the generic fibre, as in
Section~\ref{subsec:schemes}; multiplication by a nowhere-vanishing
entire function
leaves its divisor unchanged.
Definition~\ref{def:deninger-structure} records the cup-product and
duality requirements; Section~\ref{sec:axioms} develops them and the
stronger regularized-determinant condition.
 An analytic
construction of those determinants, identified with the Euler and
gamma factors on their convergence half-plane, would also supply the
meromorphic continuation asserted here.

\begin{conjecture}[Uniform positivity]
\label{conj:arithmetic-positivity}
There is a Deninger cohomology theory $H$ as in
Conjecture~\ref{conj:arithmetic-realization} that is polarized.
\end{conjecture}

\begin{conjecture}[Full uniform positivity]
\label{conj:arithmetic-positivity-full}
There is a full Deninger cohomology theory $H$ as in
Conjecture~\ref{conj:arithmetic-realization-full} that is polarized.
\end{conjecture}

A full Deninger cohomology theory is in particular a Deninger
cohomology theory, so
Conjecture~\ref{conj:arithmetic-realization-full} implies
Conjecture~\ref{conj:arithmetic-realization}, and
Conjecture~\ref{conj:arithmetic-positivity-full} implies
Conjecture~\ref{conj:arithmetic-positivity}. Each positivity
conjecture implies the corresponding realization conjecture.


These conjectures concern the trivial coefficient and have $w_X=0$.
Nontrivial coefficients require additional spectral realizations,
external products and duality, and arithmetic Abel--Jacobi comparisons.


Together Conjectures~\ref{conj:arithmetic-realization}
and~\ref{conj:arithmetic-positivity} assert the existence of a
polarized Deninger cohomology theory with the finite-field comparison,
and their full versions assert a polarized full one.
The positive Hodge-star clause is the one in
\cite[Conjecture~2.3(iv), p.~522]{FlachMorin2020DeningerWeilArakelov}
used by Theorem~\ref{thm:main-general}. The positive form must be
compatible with the same generator that realizes the zeta divisor.
It plays the role of the positivity supplied by the Hodge standard
conjecture, together with the Lefschetz formalism; it is not a claim
that Grothendieck's standard conjectures already apply to an
unconstructed arithmetic cohomology.

\begin{theorem}[Riemann hypothesis for number fields]
\label{thm:main-number-fields}
Let $K$ be a number field. If $\Lambda_K$ has a polarized graded
spectral realization of type $(1,0)$ with endpoint data
$(V^0,\Theta_0)=(\C,0)$ and $(V^2,\Theta_2)=(\C,1)$, every
nontrivial zero of $\zeta_K$ has real part $1/2$.
If such realizations exist for all number fields,
every zero in the critical strip of every finite-order Hecke
$L$-function, and every zero or pole there of every Artin
$L$-function, has real part $1/2$.

\end{theorem}

\begin{proof}
The endpoint data and the divisor identity identify the spectrum in
degree one with the divisor of $s(s-1)\Lambda_K(s)$, where
$\Lambda_K$ is the completed Dedekind zeta function
(Definition~\ref{def:zeta}). Apply Theorem~\ref{thm:main-general} with
$d=1$, $w=0$. The Hecke and Artin conclusions follow by the Artin
formalism, with the full proof in Theorem~\ref{thm:A}.
\end{proof}

Conjectures~\ref{conj:arithmetic-realization}
and~\ref{conj:arithmetic-positivity} imply this theorem's hypotheses
on the Arakelov objects. An absolute or characteristic-one model
with the specified realization gives the identical implication.
For $K=\Q$ this is the Riemann hypothesis for $\zeta$.
 No simplicity
of zeros or holomorphy of nontrivial Artin $L$-functions is asserted.
The existence of a positive spectral model for a single completed
Dedekind zeta function is equivalent to its Riemann hypothesis
(Proposition~\ref{prop:content}); an objectwise construction
from the zeros alone does not supply the asserted functorial theory
or an independent proof of its positivity.

The finite-field counterpart, with its full hypotheses and proof, is
Theorem~\ref{thm:B}. Its logarithmically suspended cohomology satisfies
the spectral and duality axioms; the positive form exists precisely
when Frobenius is semisimple and pure. Purity is Deligne's theorem,
whereas general semisimplicity is not asserted. For curves, the Rosati
involution supplies the required positivity
(Proposition~\ref{prop:rosati-EF}).

\subsection{Base, source, and existing constructions}
\label{subsec:deninger-structures}

This subsection compares Definition~\ref{def:deninger-theory} with
the Weil cohomology formalism and with the existing constructions, and
explains why the source is $\Schbar$ rather than a category of schemes
over a fixed base. The definitions of $\Flow$, Deninger structures,
polarizability, and Deninger modules are in
Section~\ref{subsec:uniform-formalism}, and the definition of a
Deninger cohomology theory is in
Section~\ref{subsec:uniform-instances}.

It is natural to expect \ref{ax:F2} to be the analogue of the
Lefschetz trace formula axiom of a Weil cohomology theory. In fact
Weil cohomology theories have no such axiom: they are defined on smooth
projective varieties over an algebraically closed field by finiteness,
Poincar\'e duality, K\"unneth, and cycle-class axioms
\cite[\S1.2, pp.~362--363]{Kleiman1968Cycles}, and the trace formula
is a consequence \cite[Proposition~1.3.6, p.~366]{Kleiman1968Cycles}.
Axiom~\ref{ax:F2} instead refers to a function given independently of
$H$, because no point-counting primitive is available on the
cohomological side; its dynamical form is AW6$'$
(Theorem~\ref{thm:aw6-div}). Together with
Proposition~\ref{prop:flow-exact}\eqref{item:flow-additive}, axiom
\ref{ax:F3} propagates divisor identities along stratifications, as
$\zeta(X,s)=\zeta(W,s)\zeta(Y,s)$.

It is natural to expect a single class of admissible objects on which
every axiom is imposed. In fact
\cite[Conjecture~2.3, p.~522]{FlachMorin2020DeningerWeilArakelov}
imposes its parts on three classes: the grading and determinant
formula on all arithmetic schemes and their compactifications with
compact supports, the trace and perfect pairing for regular schemes of
pure dimension, and the Hodge star, positivity, and the location of
zeros for projective regular ones. The filtration
$\Schbar\supset\Schbar^{\mathrm{reg}}\supset\Schbar^{\mathrm{proj}}$
of Definition~\ref{def:compactified-schemes} records these classes.
Projective regular objects play the role of smooth projective
varieties.

\begin{proposition}[Schemes over $\cO_K$ form a non-full subcategory]
\label{prop:slice-not-full}
Let $K$ be a number field. The forgetful functor from schemes over
$\Spec\cO_K$ to schemes is faithful and not full if
$\operatorname{Aut}(K)\ne1$. The automorphism group of $\Spec\cO_K$
over $\Spec\cO_K$ is trivial, while its automorphism group as a
scheme is $\operatorname{Aut}(K)$.
\end{proposition}

\begin{proof}
Faithfulness holds because a morphism over $\Spec\cO_K$ is in
particular a morphism of schemes. The scheme automorphisms of
$\Spec\cO_K$ are the $\Spec\sigma$ for ring automorphisms $\sigma$ of
$\cO_K$, which are the restrictions of the elements of
$\operatorname{Aut}(K)$. A morphism over $\Spec\cO_K$ commutes with
the identity structure map, so the only automorphism over
$\Spec\cO_K$ is the identity.
\end{proof}

It is natural to expect a separate theory for each number field $K$,
defined on schemes over $\Spec\cO_K$. In fact such a family loses the
morphisms that are not $\cO_K$-linear, by
Proposition~\ref{prop:slice-not-full}. These include the Galois
actions used in Deninger's formulation over all number fields
\cite[Prediction~2.13, p.~8]{Deninger2025BSDDedekind} and in the
obstruction to real coefficients
\cite[2.9, p.~4, and Theorem~2.13, p.~5]{Deninger2022NoWeil}. Since
$\Spec\Z$ is terminal among schemes, every scheme is a scheme over
$\Spec\Z$, and Definition~\ref{def:deninger-theory} fixes no base. A
theory relative to $\cO_K$ is the restriction of $H$ with the algebra
structure of Proposition~\ref{prop:relative-structure}.

\begin{proposition}[Zeta functions do not depend on a base]
\label{prop:zeta-intrinsic}
Let $X$ be a separated scheme of finite type over $\Spec\Z$, and let
$X_0$ be its set of closed points. Then
$\zeta(X,s)=\prod_{x\in X_0}(1-N(x)^{-s})^{-1}$ is the same whether
$X$ is regarded as a scheme over $\Spec\Z$ or over $\Spec\cO_K$ for a
structure morphism $X\to\Spec\cO_K$.
\end{proposition}

\begin{proof}
The closed points of $X$ and the cardinalities $N(x)$ of their residue
fields are determined by the scheme $X$.
\end{proof}

\begin{proposition}[Weil restriction changes the zeta function]
\label{prop:weil-restriction}
Let $K$ be a number field of degree $n\ge2$. Then
$\operatorname{Res}_{\cO_K/\Z}\mathbb A^1_{\cO_K}\cong\mathbb A^n_\Z$,
$\zeta(\mathbb A^1_{\cO_K},s)=\zeta_K(s-1)$, and
$\zeta(\mathbb A^n_\Z,s)=\zeta(s-n)$, and these two functions are
different.
\end{proposition}

\begin{proof}
The Weil restriction represents $T\mapsto\cO_K\otimes_\Z\Gamma(T,\cO_T)$,
and $\cO_K$ is free of rank $n$ over $\Z$, so it is $\mathbb A^n_\Z$.
For a scheme $Y$ of finite type over $\Z$, the fibre of
$\mathbb A^1_Y$ over a closed point $y$ is $\mathbb A^1_{\kappa(y)}$,
with zeta function $(1-N(y)^{1-s})^{-1}$; hence
$\zeta(\mathbb A^1_Y,s)=\zeta(Y,s-1)$, and likewise
$\zeta(\mathbb A^n_\Z,s)=\zeta(s-n)$. The first function has its
rightmost pole at $s=2$, and the second at $s=n+1\ge3$.
\end{proof}

It is natural to expect a theory relative to $\cO_K$ to be carried to
$\Spec\Z$ by Weil restriction. In fact, by
Proposition~\ref{prop:weil-restriction}, Weil restriction changes the
prescribed function, while by Proposition~\ref{prop:zeta-intrinsic}
the forgetful functor preserves it. The zeta-preserving passage from
schemes over $\cO_K$ to $\Schbar$ is the forgetful functor.

\begin{proposition}[The suspension does not depend on the base field]
\label{prop:suspension-base}
Let $X$ be a smooth projective geometrically connected variety over
$\F_{q^n}$, and let $X'$ be the same scheme regarded as a variety over
$\F_q$. Let $\Sigma_q(X')$ be the logarithmic suspension of
$H^\bullet(X'_{\overline{\F}_q},\Q_\ell)\otimes\C$ formed with the
logarithms base $q$ of the eigenvalues of the geometric
$q$-Frobenius, and $\Sigma_{q^n}(X)$ the suspension of
$H^\bullet(X_{\overline{\F}_q},\Q_\ell)\otimes\C$ formed with the
logarithms base $q^n$ of the eigenvalues of the geometric
$q^n$-Frobenius. Then $\Sigma_q(X')\cong\Sigma_{q^n}(X)$ as graded
spectral cores, and $\Sigma_q(X')^0\cong\mathbb 1_{q^n}$.
\end{proposition}

\begin{proof}
Since $\F_{q^n}\otimes_{\F_q}\overline{\F}_q\cong\overline{\F}_q^{\,n}$,
the Galois module $H^j(X'_{\overline{\F}_q})$ is induced from
$H^j(X_{\overline{\F}_q})$ along $G_{\F_{q^n}}\subset G_{\F_q}$. If
$\beta$ is an eigenvalue of $\mathrm{Fr}_{q^n}$ with multiplicity $m$,
the characteristic polynomial of $\mathrm{Fr}_q$ on the induced module
contains the factor $(T^n-\beta)^m$, so the eigenvalues of
$\mathrm{Fr}_q$ over $\beta$ are the $n$ roots $\beta^{1/n}\zeta_n^k$,
each with multiplicity $m$. Their logarithms base $q$ are the numbers
$(\log\beta+2\pi i(k+nl))/(n\log q)$ for $0\le k<n$ and $l\in\Z$,
that is the numbers $(\log\beta+2\pi ij)/\log(q^n)$ for $j\in\Z$,
which are the logarithms base $q^n$ of $\beta$, with the same
multiplicities. In degree zero $\beta=1$, and the eigenvalues
$2\pi ij/\log(q^n)$ are those of $\mathbb 1_{q^n}$.
\end{proof}

It is natural to expect the base of a relative theory, $\Spec\cO_K$ or
$\Spec\Z$, to play the role of the base field $\F_q$. In fact the base
field enters the axioms only through the unit $U_X$ of \ref{ax:F1},
the cohomology of the absolute point, and by
Proposition~\ref{prop:suspension-base} even that unit is determined by
the variety and not by a chosen base field. The cohomology of $\Xbar_K$
has $H^0=\C$, $H^1$ infinite-dimensional, and $H^2=\C(-1)$
\cite[(3), p.~168]{Deninger1998Analogies}: this is the shape of the
cohomology of a curve over $\F_q$, not of a unit. Schemes over
$\cO_K$ are the analogue of varieties over a curve, and
\cite[p.~172]{Deninger1998Analogies} describes the relative structure
of Proposition~\ref{prop:relative-structure} as a module structure
over the cohomology of the base.

It is natural to expect the Flach--Morin and Connes--Consani--Marcolli
constructions to stand to Definition~\ref{def:deninger-theory} as
$\ell$-adic and crystalline cohomology stand to the Weil cohomology
axioms, which both satisfy
\cite[\S1.2, pp.~362--363, and Examples~1.2.5, p.~364]%
{Kleiman1968Cycles}. In fact neither is an instance. The Flach--Morin
complexes are real complexes at a fixed integer twist without a flow;
their relation to the flow is
\cite[Conjecture~4.1, p.~527, and Remark~4.15, p.~535]%
{FlachMorin2020DeningerWeilArakelov}, as recorded in
Remark~\ref{rem:deninger-terminology}. Their pullbacks are constructed
for equidimensional flat maps of regular proper arithmetic schemes
\cite[\S6.6, p.~1522]{FlachMorin2018WeilEtale}, and the
compactified complexes depend on the K\"ahler form
\cite[\S3.2, p.~523]{FlachMorin2020DeningerWeilArakelov}. The adelic
cyclic homology of
\cite[Definition~4.10, p.~21]{ConnesConsaniMarcolli2009Weil} is a
degree-one space attached to a global field, without schemes or cup
products, and its functoriality for Galois extensions passes through a
non-canonical covering \cite[\S9, pp.~50--54]%
{ConnesConsaniMarcolli2009Weil}. Over number fields no instance of
Definition~\ref{def:deninger-theory} is known. Over finite fields
the logarithmic suspension of $\ell$-adic cohomology satisfies the
corresponding spectral and duality axioms (Theorem~\ref{thm:B}).
What the existing constructions share is the
degree-one spectral datum of the objects of $\Schbar_1$ together with
the finite-field suspension.

\subsection{The absolute-base philosophy and the search outcomes}
\label{subsec:f1-search}

The $\F_1$ philosophy asks what geometry would let arithmetic use the
same square-and-correspondences argument as a curve over a finite
field. For a smooth projective curve $C/\F_q$, the surface
$C\times_{\F_q}C$ carries the diagonal and the graphs of powers of
Frobenius. Their intersections encode point counts, and the positivity
of correspondences, expressed by the Hodge index theorem or the Rosati
involution, supplies Weil's bound
(Section~\ref{subsec:weil-curves}). The closed points of $\Spec\Z$
are the primes, with Euler factors $(1-p^{-s})^{-1}$. The analogy calls
for a geometric setting in which the arithmetic curve has an equally
useful square and correspondences.

An absolute base is meant to make that arithmetic square possible.
The ordinary product
\[
 \Spec\Z\times_{\Spec\Z}\Spec\Z=\Spec\Z
\]
does not produce a surface. The notation
$\Spec\Z\times_{\Spec\F_1}\Spec\Z$ expresses the desired replacement
in a proposed geometry over a base below $\Spec\Z$; it is not a fibre
product of ordinary schemes over a one-element field. Manin's absolute
motives and tensor products of zeta functions make this motivation
explicit \cite[p.~121, p.~132, and \S5.2, p.~158]{Manin1995Lectures}.
Adding the infinite places to obtain $\overline{\Spec\Z}$ addresses a
second requirement, the completion of the arithmetic curve and its zeta
function. It does not itself specify this absolute base or its products.
A uniform treatment must say both what the compactified objects are
and how their proposed geometry supplies products.

\begin{remark}[There is no single $\F_1$ geometry]
\label{rem:f1-not-unique}
It is natural to speak of ``the'' category of schemes over $\F_1$. In
fact several such categories have been constructed for different
purposes, and no single category is known to which all of them are
equivalent. Among them are Deitmar's monoid schemes
\cite{Deitmar2004SchemesF1}, Soul\'e's varieties
\cite{Soule2003Varietes}, the relative schemes of Toën--Vaquié
\cite{ToenVaquie2005SpecZ}, Durov's generalized rings
\cite{Durov2007Arakelov}, Borger's $\Lambda$-schemes
\cite{Borger2009Lambda}, Lorscheid's blue schemes
\cite{Lorscheid2011Blueprints}, the glued schemes of Connes and
Consani \cite{ConnesConsani2009SchemesF1}, and their arithmetic and
scaling sites
\cite{ConnesConsani2014ArithmeticSite,ConnesConsani2015ScalingSite};
surveys are \cite{LopezPenaLorscheid2009Mapping} and
\cite[p.~1]{Lorscheid2018Everyone}. Some comparisons are theorems:
Deitmar's monoid schemes are equivalent to the Toën--Vaquié schemes
relative to sets \cite[Theorem~36]{Vezzani2010Deitmar}. Others
are not equivalences. The arithmetic square also behaves differently:
in Durov's theory $\Z\otimes_{\F_1}\Z=\Z$
\cite[0.5.19, p.~36]{Durov2007Arakelov}, and in $\Lambda$-geometry the
product $\Spec\Z\times_{\Spec\F_1}\Spec\Z$ is $\Spec\Z$
(Proposition~\ref{prop:borger-square}), while the arithmetic site has
a square of a different kind
\cite[\S6, p.~22]{ConnesConsani2016ArithmeticGeometry}. A cohomological
formulation therefore cannot take its source to be ``the'' $\F_1$
category; an absolute geometry enters through a chosen comparison
functor from $\Schbar$, as in
Section~\ref{subsec:uniform-instances}.
\end{remark}

\begin{proposition}[Products over $\F_1$ in $\Lambda$-geometry]
\label{prop:borger-square}
In the $\Lambda$-algebraic geometry of \cite{Borger2009Lambda}, let
$\F_1$ denote $\Z$ with its unique $\Lambda$-structure. Then
$\Spec\Z$ is a terminal $\Lambda$-scheme, and
$\Spec\Z\times_{\Spec\F_1}\Spec\Z=\Spec\Z$ in the category of
$\Lambda$-schemes.
\end{proposition}

\begin{proof}
By \cite[p.~2]{Borger2009Lambda}, $\Z$ with its $\Lambda$-structure is
the initial object in the category of $\Lambda$-rings, so its spectrum
is terminal among $\Lambda$-schemes. A product with a terminal object
is the other factor.
\end{proof}

It is natural to expect the source of a Deninger cohomology theory to
be a category over $\Spec\F_1$ containing the arithmetic square. In
fact such a square cannot be an object on which the theory takes
values in spectral cores (Theorem~\ref{thm:kun-square}), in
$\Lambda$-geometry the square collapses
(Proposition~\ref{prop:borger-square}), and by
Remark~\ref{rem:f1-not-unique} there is no single $\F_1$ category to
choose. The source of Definition~\ref{def:deninger-theory} is
therefore $\Schbar^{\mathrm{proj}}$, and a geometry containing the
square is an ambient category reached by a comparison functor.

The search found precise failures of several expected properties of
such a square. An untwisted tensor formula over $\Spec\Z$ is
incompatible with the spectral normalization of the arithmetic curve:
it would identify its degree-one flow module with two copies of itself
(Proposition~\ref{prop:kun-over-Z}). General derived tensor formulas
over arbitrary bases also fail in the finite-field examples of
Propositions~\ref{prop:kun-EM} and~\ref{prop:EM-cochain}.
These results delimit particular K\"unneth formulas; they do not rule
out an absolute product with a different base and comparison theorem.

Even an absolute square cannot inherit every analytic property of the
curve if it retains the expected cross products. Under the spectral
realization, duality, pullback, cup-product, and multiplicative-trace
hypotheses of Theorem~\ref{thm:kun-square}, the cross-product map on
degree-one tensors is injective. For a square, the pairs of zeros
$\rho$ and $1-\rho$ then give infinitely many independent vectors with
flow eigenvalue $1$. Thus its degree-two flow module fails the finite
multiplicity requirement for the spectral cores used here. Under the
positive Hilbert-space realization, the corresponding tensor-sum
operator has no compact resolvent
(Proposition~\ref{prop:kun-AW6}). The infinite multiplicity needs no
extra hypothesis on zero spacings; the stronger accumulation claim in
Theorem~\ref{thm:kun-square} has its own stated hypothesis. Deninger and
Manin had already identified the sum-of-zeros difficulty
\cite[(7.32), pp.~41--42]{Deninger1994Motivic},
\cite[p.~132]{Manin1995Lectures}.

Restricting tensor pairs to imaginary parts of the same sign does not
supply the missing cohomology ring. This restriction, Kurokawa's cone,
excludes the tensors pairing nonreal $\rho$ with $1-\rho$. An algebra
containing the pullbacks from both factors must contain all their
cross products, including these excluded tensors
(Proposition~\ref{prop:kun-cone}). The periodic-orbit model exhibits
another difficulty: the tensor sum for distinct prime periods
$\log p$ and $\log p'$ has dense imaginary spectrum, whereas a repeated
prime period gives infinite multiplicities
(Proposition~\ref{prop:kun-orbits}). The corresponding diagonal flows
have, respectively, no closed orbits or nonisolated families of them
(Proposition~\ref{prop:kun-torus}); a product trace formula cannot simply
retain the isolated-orbit description of the factors.

Connes--Consani supply concrete characteristic-one geometry in which to
pursue the absolute-base idea. Characteristic one here refers to
idempotent addition, as in Boolean and tropical semirings, rather than
to an ordinary field. Their arithmetic site is a semiringed topos, and
their scaling site extends its scalars from the Boolean semifield to
the tropical semifield. The arithmetic site has a constructed square
and Frobenius correspondences
\cite[\S6]{ConnesConsani2016ArithmeticGeometry}. In the version of the
site used in \cite[Theorems~2.6--2.7]{ConnesConsani2014ArithmeticSite},
points over the maximal compact subsemiring $[0,1]_{\max}$ of the tropical
semifield
recover an adelic quotient, and the counting distribution gives the
completed Riemann zeta function
\cite[Theorems~2.6--2.7]{ConnesConsani2014ArithmeticSite}.
Thus the search has actual geometric objects and zeta identities to
compare, not only a proposed notation for an absolute square.

The available Riemann--Roch theorems identify a further part of the
analogy that has been realized. On the scaling site's periodic orbit
at a prime $p$, whose length is $\log p$, Connes--Consani construct
divisors and prove a Riemann--Roch formula
(Theorem~\ref{thm:imp-cc16}). Their later Riemann--Roch theorem on
$\overline{\Spec\Z}$ is also a construction in dimension one
\cite[Theorem~1.1, p.~2]{ConnesConsani2023RiemannRochSpecZ}.
Neither statement is the global intersection inequality on a square
needed for the RH strategy. The checked sources propose that further
step through a Riemann--Roch and intersection-theoretic framework
\cite[\S3, pp.~16--18, and p.~22]{ConnesConsani2018RiemannRoch}.

The search did not find a comparison identifying these constructions
with a single cohomology of the compactified arithmetic curve that
has the required global zeta spectrum, duality, and positive star.
Connes--Consani's archimedean cyclic cohomology realizes gamma factors
(Theorem~\ref{thm:imp-cc14}); it does not by that theorem realize the
nontrivial zeros. Ayoub's existing realization theories and the known
Weil--Arakelov constructions likewise supply specific comparison
results, whose scope is listed in Table~\ref{tab:theory-status}.
The missing global identification and positivity theorem are open
requirements in this synthesis. Their absence from the checked corpus
is a search outcome, not a nonexistence theorem for $\F_1$ cohomology.

The finite-field comparison also produces a constructive correction.
After replacing Frobenius by its full logarithmic suspension, even the
finite-field product formula requires a relative tensor product over
the periodic unit $\mathbb 1_q$, the suspended cohomology of the base.
Tensoring instead over $\C$ repeats every spectral value infinitely
(Proposition~\ref{prop:kun-Fq-C}); tensoring over $\mathbb 1_q$ gives
the actual K\"unneth isomorphism of Theorem~\ref{thm:kun-Fq}.
The choice of base and tensor product is therefore part of the
geometric proposal, including in the case where the cohomology already
exists.

The surviving proposal separates the curve's zeta realization from the
larger space needed for its square. Conjecture~\ref{conj:kunneth}
requires a multiplicative surface datum and a K\"unneth isomorphism,
with a Hodge-index condition on the tensors pairing eigenvalues
$\rho$ and $1-\rho$. It imposes no finite-multiplicity, determinant, or
compact-resolvent condition on that surface. With its stated duality
and star compatibilities, the Hodge-index clause gives purity and
hence RH by Theorem~\ref{thm:kun-hodge}. This is the role analogous to
a standard conjecture: the required positivity is a geometric
statement about correspondences on a square. A construction realizing
these data in the arithmetic or scaling site, together with the
comparison identifying its intersection form, remains to be supplied.

In the uniform formulation below, compactified arithmetic schemes,
finite-field varieties, and a specified absolute geometry can provide
inputs to the same spectral implication. The objects of
$\Schbar^{\mathrm{proj}}$, on which a zeta determinant is required, do
not include absolute squares. This distinction
is forced by the tensor obstruction above
and permits a larger target for the surface cohomology. The
$\F_1$ part of the search therefore contributes both a geometric
construction problem and restrictions on its possible solution:
realize the required cohomology on an appropriate absolute square,
compare its pairings with Weil positivity, and prove the compatible
positive structure. The manuscript reports those restrictions together with
the constructions that are already available.

\subsection{Comparison with the Weil and standard conjectures}

Over a finite field, $\ell$-adic cohomology is already constructed;
the Lefschetz and Hodge standard conjectures seek algebraic
correspondences and positivity on it. Their conjunction gives the
Riemann hypothesis by Serre's argument
\cite[Theorem~5-6, p.~19]{Kleiman1994Standard}. In the arithmetic
setting both the global spectral cohomology and its positive structure
are conjectural. Ayoub's Weil theory, Connes--Consani's cyclic homology
and arithmetic sites, and Weil--Arakelov complexes are actual
constructions with different coefficient rings, source categories, and
comparison theorems. Table~\ref{tab:theory-status} states those
differences and the remaining conjectures separately.

The analogy between Theorem~A and Theorem~B is exact for positivity and
duality, and not symmetric in one place. Over $\F_q$ the cohomology is
constructed independently of the zeta function, and rationality and
\ref{ax:AW1}--\ref{ax:AW3} are theorems about it, rationality being
derived from the Lefschetz trace formula for all powers of Frobenius
(Theorem~\ref{thm:ff-lefschetz}\ref{item:fl-converse}; in these notes
rationality itself is imported, Theorem~\ref{thm:imp-deligne}, and the
Lefschetz formula is read off from it); over a number
field no object is
known that satisfies \ref{ax:AW1} with an arithmetically defined cup
product, and the divisor identity \ref{ax:AW2} is imposed as an axiom.
Given \ref{ax:AW1}, \ref{ax:AW2} is equivalent to the trace-formula
axiom AW6$'$, whose geometric side is the explicit formula of
$\zeta_K$ (Definition~\ref{def:aw6prime},
Theorem~\ref{thm:aw6-div}): imposing AW6$'$ instead of \ref{ax:AW2}
changes the form of the assumption, not its strength. Test functions
supported in a bounded interval $(0,L)$ of positive times do not
determine the spectrum (Theorem~\ref{thm:window}); every tail of times
does. The determinant
normalization, exactly one over $\F_q$ (Theorem~\ref{thm:ff-det}), is
computed over every number field: it is $2^{-n/2}|D_K|^{-s/2}$
relative to $\Lambda_K$ and exactly one relative to Deninger's
completion (Theorem~\ref{thm:det-K}).
Table~\ref{tab:dictionary} is the dictionary between the
two settings.

{\small
\begin{longtable}{@{}>{\raggedright\arraybackslash}p{0.3\textwidth}
>{\raggedright\arraybackslash}p{0.3\textwidth}
>{\raggedright\arraybackslash}p{0.3\textwidth}@{}}
\caption{Dictionary between Theorem~B over $\F_q$ and Theorem~A over
number fields.}\label{tab:dictionary}\\
\toprule
Smooth projective $X/\F_q$ & Number field $K$, object $\Xbar_K$ &
Where\\
\midrule
\endfirsthead
\toprule
Smooth projective $X/\F_q$ & Number field $K$, object $\Xbar_K$ &
Where\\
\midrule
\endhead
\bottomrule
\endfoot
existence of $H^\bullet$: a theorem (Grothendieck), independent of
$Z(X,t)$ & existence of a Weil datum from arithmetic input: open &
Thm.~\ref{thm:imp-weil-coh}, Sec.~\ref{sec:open}\\
geometric Frobenius $\mathrm{F}$ on $H^j$ & flow $e^{t\Theta}$ with
generator $\Theta_j$ on a spectral core $V^j$; $\mathrm{F}=q^\Theta$ on
the suspension & Def.~\ref{def:flow}, \ref{def:log-susp}\\
eigenvalue $\alpha$ of $\mathrm{F}$ & the progression
$\Sigma_q(\alpha)$ of eigenvalues $\rho$ of $\Theta$ &
Def.~\ref{def:suspension}\\
$|\alpha|=q^{j/2}$ & $\Re\rho=j/2$ (for $K$: $\Re\rho=\tfrac12$) &
Prop.~\ref{prop:purity}\\
rationality: $Z(X,t)=\prod_j\det(1-\mathrm{F}t\mid H^j)^{(-1)^{j+1}}$,
a theorem derived from the Lefschetz trace formula &
\ref{ax:AW2}: $\divisor\Lambda_K=\sum_j(-1)^{j+1}\operatorname{Sp}_j$,
imposed as an axiom, equivalent given \ref{ax:AW1} to AW6$'$; AW-det,
with normalization term $2^{-n/2}|D_K|^{-s/2}$, and $1$ for
$\widehat\zeta_K$ &
Def.~\ref{def:aw}, \ref{def:aw-det}, Thm.~\ref{thm:aw6-div},
\ref{thm:det-K}\\
Poincar\'e duality, $\mathrm{F}=q^d$ on $H^{2d}$ & \ref{ax:AW3}:
duality pairing $\rho$ with $d-\rho$ & Rem.~\ref{rem:aw-comments}\\
unit $H^0$ with $\mathrm{F}=1$; periodic unit $\mathbb 1_q$ after
suspension & arithmetic unit $\mathbb 1=(\C,0)$; $D_0=s/2\pi$ &
Prop.~\ref{prop:unit}, Sec.~\ref{subsec:det-normalization}\\
positive Frobenius polarization $h(\mathrm{F}u,\mathrm{F}v)=q^jh(u,v)$
& positive infinitesimal polarization $h(\Theta u,v)+h(u,\Theta v)=
jh(u,v)$ (AW4-V) & Def.~\ref{def:frob-pol}, \ref{def:inf-pol}\\
Weil/Rosati positivity: $\Tr(\alpha\alpha^\dagger)>0$ on
$\operatorname{End}^0(J)$ with $\pi\pi^\dagger=q$, the multiplicative
algebra form AW4-E$_{\mathrm F}$ & positive trace involution on an
algebra containing $\Theta$ with $\Theta^\dagger=j-\Theta$
(AW4-E), which AW4-E$_{\mathrm F}$ implies on the suspension &
Prop.~\ref{prop:algebra-form}, \ref{prop:aw4-EF}, \ref{prop:rosati-EF},
\ref{prop:aw4-forms}\\
Castelnuovo positivity of correspondences on $C\times C$ (Hodge index)
& Hodge index clause on the algebraic classes of
$\Xbar_K\boxtimes\Xbar_K$: equivalent to purity with definite star
forms & Thm.~\ref{thm:kun-hodge}, Conj.~\ref{conj:kunneth}\\
K\"unneth: $H^\bullet(X\times Y)=H^\bullet(X)\otimes H^\bullet(Y)$;
after suspension, a tensor product over $\mathbb 1_q$ & conjectural;
the tensor square is not a spectral core, and there is no formula over
$\Spec\Z$ & Thm.~\ref{thm:kun-Fq}, \ref{thm:kun-square},
Prop.~\ref{prop:kun-over-Z}\\
suspended star, $\operatorname{tr}(u\cup*_jv)=h_j(u,v)$,
$*_j\mathrm{F}=q^{j-d}\mathrm{F}*_j$ (when $\mathrm{F}$ is semisimple
and pure) & \ref{ax:AW4}: positive flow-compatible star,
$\Theta_{2d-j}*_j=*_j\Theta_j+(d-j)*_j$ & Def.~\ref{def:susp-datum},
Prop.~\ref{prop:polarization-identity}\\
positive form exists $\Leftrightarrow$ $\mathrm{F}$ semisimple and pure
& package on one object $\Leftrightarrow$ RH for that object &
Prop.~\ref{prop:weil-converse}, \ref{prop:content}\\
Lefschetz trace formula for all $\mathrm F^n$, a theorem, from which
rationality follows; two-sided form & the trace-formula axiom AW6$'$
(explicit formula of $\zeta_K$), an axiom, equivalent to \ref{ax:AW2}
given \ref{ax:AW1}; one-sided bounded windows of times do not
suffice &
Def.~\ref{def:aw6prime}, Thm.~\ref{thm:ff-lefschetz},
\ref{thm:aw6-div}, \ref{thm:window}\\
RH: Weil (curves, via positivity), Deligne (all $X$, not via
positivity); the positivity axiom known where $\mathrm{F}$ is semisimple
& the axioms: open & Rem.~\ref{rem:B-status}, Sec.~\ref{sec:open}\\
\end{longtable}
}

\subsection{Other obstructions found by the search}

\emph{What fails.} Sections~\ref{subsec:witnesses}
and~\ref{sec:negative} record seventeen
findings (R1--R17, correcting the expectations O01--O23 of the register
in Section~\ref{subsec:neg-register})
about expectations that are natural when intuition is imported from the
finite-field case and that are false. The Weil formalism, perfect
duality, and a flow-compatible quaternionic star do not force purity
(Propositions~\ref{prop:R1}--\ref{prop:R2-star}); comparisons at integer
twists do not see a nonreal spectrum (Proposition~\ref{prop:R3});
skew-symmetry does not control a spectrum (Proposition~\ref{prop:R4});
real coefficients, general leafwise cohomology, and local determinant
formulas each fail for a precise reason
(Remarks~\ref{rem:R5-import} and~\ref{rem:R6},
Propositions~\ref{prop:R5} and~\ref{prop:R7}); scalar de Branges models
do not see zero multiplicities, and several positivity statements that
look sufficient are not (Propositions~\ref{prop:R9}--\ref{prop:R16} and
Remark~\ref{rem:R13}); and the literal formulas of
one version of a recent comparison produce a degenerate seminorm
(Proposition~\ref{prop:R17}).

Sixteen further corrected expectations
 (O24--O39),
listed in the register of Section~\ref{subsec:neg-register}, concern the
main theorem and its calibration: the axioms other than positivity can
be satisfied with any spectrum symmetric under $\rho\mapsto1-\rho$ with
even central multiplicity, and a flow-compatible form
that is not positive allows zeros off the line
(Proposition~\ref{prop:content}); Galois descent controls only the
rational orbits of isotypic divisors (Proposition~\ref{prop:descent});
the zeta function of $\mathbb{P}^1_\Z$ has zeros on two lines
(Proposition~\ref{prop:P1Z}), and without positivity the total zeta
function of an arithmetic surface does not determine its degreewise
spectra (Proposition~\ref{prop:no-positivity-cancel}); over $\F_q$, a
``formal curve'' with Poincar\'e duality, rational zeta function,
functional equation, and integral point counts violates purity
(Proposition~\ref{prop:formal-curve}), Weil's positivity is not a form
on $H^1$ (Proposition~\ref{prop:ff-N7}), a single branch of the
logarithm destroys the divisor (Proposition~\ref{prop:one-branch}), and
the unit object of the axioms must be the periodic cohomology of the
base (Proposition~\ref{prop:unit}).

Seventeen more
 (O40--O56) concern
the multiplicative form of Weil's positivity
(Proposition~\ref{prop:aw4-EF}), the parity of central orders
(Proposition~\ref{prop:N-central-order}), and the K\"unneth conjecture
(Section~\ref{sec:kunneth}).

Seventeen further
 (O57--O73) concern the
normalization of the determinant over a number field and the
trace-formula axiom (Sections~\ref{subsec:det-normalization}
and~\ref{subsec:trace-axiom}), and two more (O74--O75) the
Eilenberg--Moore comparison over a base
(Section~\ref{sec:kunneth}).

\subsection{The construction and positivity problems}

For a single number field, \ref{ax:AW1}--\ref{ax:AW3} can be
satisfied by assigning vector spaces to the known zeta zeros.
The full objectwise package exists if and only if RH holds for that
Dedekind zeta function (Proposition~\ref{prop:content}). Such a model
does not construct cohomology from arithmetic geometry or establish
comparisons with existing theories. The research problem is to
construct the spaces, products, and pairings independently of the zeros
and then prove that the same realization has the compatible positivity
of \ref{ax:AW4}.

The distinction persists in the other cases considered here. For one
Galois extension, the package with descent follows from the relevant
Riemann hypotheses together with Artin holomorphy
(Proposition~\ref{prop:class-converse}). Over a finite field the positive
Frobenius-compatible form requires both purity and semisimplicity.
For higher-dimensional arithmetic objects, the resulting zero and pole
lines depend on cohomological degree
(Theorem~\ref{thm:A-schemes}).
\subsection{The K\"unneth and Hodge-index conjecture}
\label{subsec:intro-kunneth}

Weil's argument uses correspondences on $C\times C$. To state its
arithmetic analogue, the following definitions allow a square with
larger cohomology spaces than those of the curve. The symbol
$\Xbar_K\boxtimes\Xbar_L$ denotes the intended arithmetic square;
it does not specify an ordinary fibre product over $\Spec\Z$.

\begin{definition}\label{def:aw4-alg}
A Weil datum $(V^\bullet,\Theta,\cup,\Tr,*)$ of type $(1,0)$ satisfies
the \emph{algebraic part of \ref{ax:AW4} in degree one} if $*_1$ is
bijective, $\Theta_1*_1=*_1\Theta_1$,
$B_1(*_1x,*_1y)=\overline{B_1(x,y)}$ for all $x,y\in V^1$, and the star
form $h_1$ is Hermitian.
\end{definition}

\begin{definition}\label{def:surface-datum}
Let $D_K=(V_K^\bullet,\Theta,\cup,\Tr_K,*)$ and
$D_L=(V_L^\bullet,\Theta,\cup,\Tr_L,*)$ be Weil data of type $(1,0)$ on
$\Xbar_K$ and $\Xbar_L$.
\begin{enumerate}[label=(\alph*)]
\item A \emph{surface datum} over $(D_K,D_L)$ consists of an associative
  graded-commutative unital graded algebra $H^\bullet=\bigoplus_{j=0}^4
  H^j$ over $\C$, a degree-preserving derivation $\Theta_S$ of it, a
  linear \emph{trace} $\Tr_S:H^4\to\C$, antilinear \emph{stars}
  $*_S:H^j\to H^{4-j}$, and graded unital algebra homomorphisms
  $\mathrm{pr}_1^*:V_K^\bullet\to H^\bullet$ and
  $\mathrm{pr}_2^*:V_L^\bullet\to H^\bullet$ with
  $\Theta_S\mathrm{pr}_i^*=\mathrm{pr}_i^*\Theta$.
\item Its \emph{K\"unneth map} is $\kappa:V_K^\bullet\otimes
  V_L^\bullet\to H^\bullet$, $\kappa(x\otimes y)=\mathrm{pr}_1^*x\cup
  \mathrm{pr}_2^*y$.
\item Its \emph{algebraic part} is $N_S=\kappa\bigl(\bigoplus_\rho
  V^1_{K,\Theta\sim\rho}\otimes V^1_{L,\Theta\sim1-\rho}\bigr)\subset
  H^2$.
\item Its \emph{star forms} are $h_S^j(u,v)=\Tr_S(u\cup*_Sv)$ on $H^j$.
\item It is \emph{multiplicative} if
  $\Tr_S(\mathrm{pr}_1^*x\cup\mathrm{pr}_2^*y)=\Tr_K(x)\Tr_L(y)$ for
  $x\in V_K^2$, $y\in V_L^2$, and
  $*_S\kappa(x\otimes y)=(-1)^{|x||y|}\kappa(*x\otimes*y)$ for
  homogeneous $x,y$.
\item It satisfies the \emph{Hodge index clause} if $h_S^2$ is positive
  definite on $N_S$.
\end{enumerate}
\end{definition}

The Hodge-index clause is positivity of the star form on $N_S$.
It is not positivity of the intersection pairing on all of $H^2$.
A geometric construction must compare these tensors and their form
with actual correspondences on the proposed arithmetic square.

\begin{definition}\label{def:tensor-datum}
Let $D=(V^\bullet,\Theta,\cup,\Tr,*)$ and
$\widetilde D=(\widetilde V^\bullet,\widetilde\Theta,\cup,\widetilde\Tr,
\tilde*)$ be Weil data of types $(d,0)$ and $(\tilde d,0)$.
\begin{enumerate}[label=(\alph*)]
\item The \emph{tensor datum} $D\otimes\widetilde D$ has
  $(V\otimes\widetilde V)^n=\bigoplus_{a+b=n}V^a\otimes\widetilde V^b$
  with the flow
  $\Theta\otimes1+1\otimes\widetilde\Theta$, the product
  $(x\otimes y)\cup(x'\otimes y')=(-1)^{|y||x'|}(x\cup x')\otimes
  (y\cup y')$, the trace $\Tr\otimes\widetilde\Tr$ on
  $V^{2d}\otimes \widetilde V^{2\tilde d}$, and the stars
  $*(x\otimes y)=(-1)^{|x||y|}*x\otimes\tilde*y$ for homogeneous $x,y$.
\item For $n\ge0$, the \emph{Tate datum} $T_n$ is the Weil datum of type
  $(n,0)$ with $T_n^{2i}=\C h^i$ for $0\le i\le n$ and $T_n^j=0$ for odd
  $j$, generator $\Theta h^i=ih^i$, product $h^ih^k=h^{i+k}$ (zero if
  $i+k>n$), trace $\Tr(ch^n)=c$, and stars $*(ch^i)=\bar ch^{n-i}$.
\end{enumerate}
\end{definition}

\begin{conjecture}[K\"unneth conjecture]\label{conj:kunneth}
There are, for every number field $K$, a Weil datum $D_K$ of type
$(1,0)$ on $\Xbar_K$ satisfying \ref{ax:AW1}, \ref{ax:AW2}, and
\ref{ax:AW3} for $(\Lambda_K,\mathbb 1)$ and the algebraic part of
\ref{ax:AW4} in degree one (Definition~\ref{def:aw4-alg}), and, for
every pair of number fields $(K,L)$, a multiplicative surface datum
$S_{K,L}$ over $(D_K,D_L)$ (Definition~\ref{def:surface-datum}) such
that:
\begin{enumerate}[label=(K\arabic*),ref=(K\arabic*)]
\item\label{item:K-iso} the K\"unneth map of $S_{K,L}$ is an isomorphism
  of graded algebras, for the Koszul product on $V_K^\bullet\otimes
  V_L^\bullet$ (Definition~\ref{def:tensor-datum});
\item\label{item:K-hodge} $S_{K,K}$ satisfies the Hodge index clause.
\end{enumerate}
\end{conjecture}

No spectral-core, determinant, or compact-resolvent condition is imposed
on $S_{K,L}$.

The curve data in this conjecture satisfy only the algebraic star
condition before the surface positivity is imposed. Under the stated
compatibilities, the Hodge-index clause gives purity and definiteness
of the degree-one star form on each generalized eigenspace, and hence
RH (Theorem~\ref{thm:kun-hodge}). The signs of that form may depend on
the eigenspace. No comparison identifying these data with the
cohomology of the arithmetic or scaling site is asserted here.

The stronger product expectations have explicit obstructions. Ordinary
products over $\Spec\Z$ do not satisfy the proposed untwisted
K\"unneth formula (Proposition~\ref{prop:kun-over-Z}); the general
derived comparison also fails without its additional hypotheses
(Propositions~\ref{prop:kun-EM} and~\ref{prop:EM-cochain}). On an
arithmetic self-square, duality forces infinite multiplicity, and a
positive Hilbert realization has no compact resolvent
(Theorem~\ref{thm:kun-square}, Proposition~\ref{prop:kun-AW6}).
The same-sign cone omits required cross products
(Proposition~\ref{prop:kun-cone}). Positive product results remain over
the finite-field periodic unit and for Tate data
(Theorem~\ref{thm:kun-Fq}, Proposition~\ref{prop:kun-tate}).

\emph{Other results.} Section~\ref{sec:debranges} reproduces
Lagarias's RH-equivalent Hermite--Biehler criterion, keeping
multiplicities explicit (Proposition~\ref{prop:hb}), and records what
canonical systems and the Weil form do and do not supply.

\emph{What is imported.} Section~\ref{sec:classical} states, with exact
source locators and hypotheses, the classical analytic facts used;
Sections~\ref{sec:finite-field} and~\ref{sec:class} state the
finite-field, Artin, and representation-theoretic facts where they are
used; and Section~\ref{sec:imported} states in the same way the
comparison theorems relevant to candidate theories
(Section~\ref{sec:candidates}) and the de Branges and Weil-positivity
results used in Section~\ref{sec:debranges}. None of them constructs
the object of Theorem~\ref{thm:A}. The provenance of every other
statement is recorded in Subsection~\ref{subsec:provenance}.

\emph{What is open.} Section~\ref{sec:open} lists the open questions,
including whether a trace formula tested on a two-sided bounded window
of times determines the spectrum (Question~\ref{q:trace-axiom});
Subsection~\ref{sec:lean} records the Lean formalization, claim by
claim.





\section{Prior work}\label{sec:prior}

These notes are an exposition of previous work. Their central
conditional statement, that a cohomology of the Arakelov
compactification of $\Spec\cO_K$ with a flow and a positive star
compatible with it would force the zeros of $\zeta_K$ onto the critical
line, is Deninger's formal argument, which he models on Serre's
K\"ahler analogue of the Weil conjectures; its finite-field model is
Weil's proof of the Riemann hypothesis for curves in the form given by
Serre and Kleiman. This section surveys that work and its neighbours,
with exact locators, and then records, statement by statement, where
the results of the later sections appear in the literature or where
their ideas come from. What the notes add is organization: one list of
axioms with exact equivalences among their forms, a calibration of the
axioms against the finite-field case, negative results and explicit
counterexamples that delimit what the axioms other than positivity can
give, and a Lean formalization of the linear-algebraic parts.

\subsection{Positivity and the Riemann hypothesis over finite
fields}\label{subsec:prior-ff}

Weil's proofs of the Riemann hypothesis for a curve $C$ over $\F_q$ rest
on one positivity statement: for a nonzero class $\xi$ of
correspondences on $C\times C$, the trace $\sigma(\xi\circ\xi')$ is
positive, and applied to $\xi$ built from the Frobenius graph it gives
$|\alpha|=q^{1/2}$ for the Frobenius roots. It is stated with the
deduction of the Riemann hypothesis in
\cite[p.~346 and p.~347]{Weil1941RHFunctionFields}; for the earlier
note \cite{Weil1940CRAS} and for the book \cite{Weil1948Courbes} see
the account in \cite[\S1]{Milne2016RHFF}; Milne locates the
positivity statement of the book as its Theorem~10, p.~54
\cite[p.~12]{Milne2016RHFF}. In the language of abelian varieties
the same positivity is that of the Rosati involution,
$\Tr(\alpha\alpha^\dagger)>0$ on $\operatorname{End}^0(J)$
(\cite{Weil1948VarietesAbeliennes}, as reported in
\cite[p.~23]{Milne2016RHFF}), \cite[\S21, Theorem~1, p.~192]%
{Mumford1974AV}, \cite[Theorem~14.3, p.~62]{MilneAV2008}, and the
Riemann hypothesis follows by applying it to Frobenius, for which
$\pi\pi^\dagger=q$ \cite[Theorems~1.27 and~1.30, pp.~21--22]%
{Milne2016RHFF}.

Serre proved the K\"ahler analogue: on a smooth complex projective
variety, an endomorphism that multiplies a hyperplane class by $q$ has
eigenvalues of absolute value $q^{r/2}$ on $H^r$, because after
normalization it is unitary for the positive definite Hermitian form
given by the Hodge--Riemann bilinear relations; his second theorem is
the higher-dimensional analogue of Castelnuovo's positivity
\cite[Th\'eor\`emes~1 and~2, pp.~392--394]{Serre1960Analogues}. This
``unitary operator'' pattern is the one transported by Grothendieck's
standard conjectures and by Deninger's program. Kleiman and
Grothendieck formulated the standard conjectures of Lefschetz type and
of Hodge type and showed that together they imply the Riemann
hypothesis with semisimple Frobenius by Serre's argument
\cite[pp.~359--361, 3.11 pp.~380--381, Lemma~4.6 p.~385,
Theorem~4.7 p.~386]{Kleiman1968Cycles},
\cite[Remark~(2), p.~197]{Grothendieck1969Standard},
\cite[Theorem~5-6, p.~19]{Kleiman1994Standard}; both conjectures remain
open in general. Deligne proved the Riemann hypothesis for smooth
projective varieties over $\F_q$ by another route, through Lefschetz
pencils, monodromy, and Rankin's method applied to tensor powers
\cite[Th\'eor\`eme~(1.6), p.~276, and \S3, pp.~283--284]{Deligne1974WeilI},
and extended it to weights of mixed sheaves
\cite[Th\'eor\`eme~1, p.~138]{Deligne1980WeilII}; no positivity on
algebraic cycles enters. The survey \cite{Milne2016RHFF} covers this
history. Serre also observed that a Weil cohomology cannot have
coefficients in $\R$, $\Q$, or $\Q_p$, because for a supersingular
elliptic curve $\operatorname{End}(E)\otimes\R$ is Hamilton's
quaternions \cite[(2.2) and fn.~25, p.~27]{Milne2016RHFF}; this is the
finite-field form of the obstruction of Section~\ref{subsec:prior-den}
below.

\subsection{Standard conjectures and Birch--Swinnerton-Dyer}
\label{subsec:prior-sc-bsd}

The comparison above suggests asking whether Grothendieck's standard
conjectures also imply Birch--Swinnerton-Dyer. The phrase ``standard
conjectures'' needs care here. Grothendieck's Lefschetz and Hodge standard
conjectures concern algebraic correspondences and positivity on algebraic
cycles of smooth projective varieties; they do not state a comparison
between rational points and the central zero of an elliptic $L$-function
\cite[pp.~193--197]{Grothendieck1969Standard}. The divisor Tate
conjecture on a surface over a finite field, and the conjectural
arithmetic spectral and Abel--Jacobi maps for an elliptic motive, are
different inputs. We record the supported implications and the exact
gaps. The examples below refute particular shortcuts; they do not prove
logical independence of BSD from all geometric standard conjectures.

Let $E/K$ be an elliptic curve over a global field, and let
$r=\operatorname{rank}E(K)$. Write
\[
 \Sha(E/K)=\ker\!\left(H^1(K,E)\longrightarrow
                  \prod_v H^1(K_v,E)\right).
\]
When $L(E/K,s)$
continues to $s=1$, put $r_{\mathrm{an}}=\ord_{s=1}L(E/K,s)$; here $L$
includes the finite Euler factors at bad places.
Rank BSD says $r_{\mathrm{an}}=r$. Full BSD also asserts finiteness of
$\Sha(E/K)$ and identifies the first nonzero Taylor coefficient with
period, regulator, local-component, and torsion factors. Over $\Q$,
Wiles gives the refined formula in terms of these factors
\cite[p.~2, Remark~1]{Wiles2000BSD}. Its complete finite-place
$L$-function is distinct from adjoining an archimedean factor in the
spectral argument below. The general number-field formula has its
own local measures.

Fix a prime $\ell\ne\operatorname{char}K$. The geometric divisor-class
map used below is the base-changed first Chern class
\[
 c_1:\operatorname{Pic}(E)\otimes\Q_\ell\longrightarrow
 H^2_{\mathrm{\acute et}}(E_{\overline K},\Q_\ell(1)).
\]

\begin{proposition}\label{prop:bsd-picard-gap}
Let $E/K$ be an elliptic curve with origin $O$. Under the degree and
$O$-splitting, $\operatorname{Pic}(E)\cong\Z\oplus E(K)$, while the
geometric class map $c_1$ records only the degree factor. Consequently
its kernel contains $E(K)\otimes_\Z\Q_\ell$.
\end{proposition}
\begin{proof}
The degree map is split by $n\mapsto\mathcal O_E(nO)$, and
$\operatorname{Pic}^0(E)(K)=E(K)$ by the Abel--Jacobi isomorphism for an
elliptic curve. On $E_{\overline K}$, a divisor class in $H^2$ is its
degree times the class of a point, so every degree-zero line bundle
has zero cohomological class. Tensoring with $\Q_\ell$ gives the stated
kernel.
\end{proof}

\begin{remark}\label{rem:bsd-implication-status}
The examples in this subsection refute specific formal deductions from
positivity, divisor data, or prime-primary information. They are not
elliptic curves violating BSD, and they do not establish logical
nonimplication from Grothendieck's standard conjectures imposed on all
varieties. No such all-varieties implication is proved here.
\end{remark}

The rational points sit in the homologically trivial part, not in the
degree-one divisor class measured by the curve's standard-conjecture
positivity. Over a function field, they can instead be represented by
horizontal divisor classes on a fibred surface; that surface is the
object in Theorem~\ref{thm:bsd-function-field} below.

For the number-field spectral route, let $V_{E,(1)}$ denote a finite
generalized $1$-eigenspace of a proposed operator $\Theta$ for the
elliptic motive. The relevant spectral point is $1$, whereas the
weight-zero zeta object of $\Spec\Z$ has central point $1/2$.
Deninger's proposed arithmetic Abel--Jacobi map compares
$E(K)\otimes\C$ with the central eigenspace
\cite[p.~8 and Example~5, p.~10]{Deninger2010HilbertPolya}.
For the conditional number-field statement below, write
$\Lambda_E(s)=g(s)L(E/K,s)$ for a proposed completed factor;
the proposition specifies the required behavior of $g$ at $1$.

\begin{proposition}\label{prop:bsd-positive-gap}
For every $r\geq0$ and $k>0$ there are a positive Hermitian space
$V$, an operator $\Theta$, and an isometric injection
$\C^r\hookrightarrow\ker(\Theta-1)$ satisfying
$\Theta+\Theta^*=2$ but with
$\dim\ker(\Theta-1)=r+k$.
\end{proposition}
\begin{proof}
Take $V=\C^{r+k}$ with its usual Hermitian metric,
$\Theta=1$, and include the first $r$ coordinates. The cokernel
has dimension $k$.
\end{proof}

This is a linear-algebra counterexample to a positivity-only rank
deduction, not an elliptic curve violating BSD. A positive height form
does not make the arithmetic Abel--Jacobi map surjective.

Consider the Jordan matrix
$J=\bigl(\begin{smallmatrix}1&1\\0&1\end{smallmatrix}\bigr)$.

\begin{proposition}\label{prop:bsd-jordan-gap}
For the matrix $J$ just displayed,
$\ord_{s=1}\det(s-J)=2$, whereas $\dim\ker(J-1)=1$.
\end{proposition}
\begin{proof}
The determinant is $(s-1)^2$, and $J-1$ has rank one.
\end{proof}

Thus determinant multiplicity cannot be replaced by eigenspace
dimension without semisimplicity or an equivalent block comparison.

\begin{proposition}[Conditional rank BSD]\label{prop:bsd-spectral-rank}
Let $E/K$ be an elliptic curve over a number field. Assume that:
$L(E/K,s)$ is meromorphic near $1$; the factor $g$ in
$\Lambda_E(s)$ is holomorphic and nonzero at $1$;
a spectral realization has a finite-dimensional generalized
$1$-eigenspace $V_{E,(1)}$ and the exact order identity
$\ord_{s=1}\Lambda_E(s)=\dim V_{E,(1)}$;
on this block a positive Hermitian form satisfies
$\Theta+\Theta^*=2$; and an arithmetic Abel--Jacobi map is an
isomorphism
\[
  E(K)\otimes_\Z\C\xrightarrow{\sim}
  \ker(\Theta-1\mid V_{E,(1)}).
\]
Then $\ord_{s=1}L(E/K,s)=\operatorname{rank}E(K)$.
\end{proposition}
\begin{proof}
The factor $g$ is a unit at $1$, so it does not change the order.
On the generalized block, $A=\Theta-1$ is nilpotent. The positive
form and $\Theta+\Theta^*=2$ make $A$ skew-Hermitian, hence
diagonalizable; thus $A=0$. The assumed order identity now gives
\[
  \ord_{s=1}L(E/K,s)=\dim V_{E,(1)}
  =\dim\ker(\Theta-1)=\dim(E(K)\otimes\C)=r.
\]
\end{proof}

The exact order identity includes any noncancellation requirement for
an alternating spectral determinant. This is the finite-dimensional
form of Deninger's conditional argument
\cite[p.~8 and Example~5, p.~10]{Deninger2010HilbertPolya}; it does not
construct the spectral realization or the Abel--Jacobi isomorphism.
His printed elliptic example concerns $E/\Q$; the statement above
for general $K$ is the same conditional linear-algebra argument.

Deninger's coefficient cohomology and its determinant formula are
proposals, and he presents the arithmetic Abel--Jacobi isomorphism
as a further condition for rank BSD
\cite[pp.~6--8 and Example~5, p.~10]{Deninger2010HilbertPolya}.
The classical isomorphism $E(K)\cong\operatorname{Pic}^0(E)(K)$ used
in Proposition~\ref{prop:bsd-picard-gap} is a different map.
The N\'eron--Tate height theorem gives a nondegenerate height pairing
on $E(K)\otimes\Q$, but it neither constructs the proposed spectral
target nor proves that an arithmetic map onto it is surjective.
Once the spectral order identity is assumed, the dimension assertion
in this last comparison is precisely the missing rank equality.
Conjectures~\ref{conj:arithmetic-realization} and
\ref{conj:arithmetic-positivity} concern trivial coefficients and
do not include the elliptic realization or this comparison.

Proposition~\ref{prop:formal-curve}\ref{item:fc-impure}--\ref{item:fc-forms}
and Remark~\ref{rem:B-status} illustrate why duality and a formal
Frobenius relation do not supply the positive form automatically.
In particular, the formal curve is not the cohomology of a curve.

There is a different established implication over a function field.
Let $C/\F_q$ be smooth, projective, and geometrically connected, and
let $K=\F_q(C)$. Let $X/\F_q$ be a smooth projective surface with a
proper flat morphism $\pi:X\to C$, a zero section, and smooth
geometrically connected generic fibre $E/K$ of genus one. Closed
fibres of $\pi$ may be singular. The divisor Tate conjecture for $X$
may be written in pole-order form
\[
  -\ord_{s=1}\zeta(X,s)=\operatorname{rank}\operatorname{NS}(X),
\]
with divisors understood over $\F_q$.
Its equivalence with the divisor cycle-map form is an arithmetic
theorem \cite[Theorem~6.1.2]{Ulmer2012CurvesJacobians}.

Let $\mathcal E/C$ be the N\'eron model of $E$, let
$t=\#E(K)_{\mathrm{tors}}$, and put
$c_v=\#\Phi_v(k_v)$ for its component group at $v$.
Let $R_{\log}$ be the determinant of the logarithmic N\'eron--Tate
height pairing on a basis of $E(K)$ modulo torsion, with $R_{\log}=1$
in rank zero. In this convention, the full function-field coefficient
formula is
\begin{equation}\label{eq:bsd-function-field}
 \lim_{s\to1}\frac{L(E/K,s)}{(s-1)^r}
 =\frac{\#\Sha(E/K)\,R_{\log}\prod_v c_v}{t^2}
    q^{\chi(C,\operatorname{Lie}\mathcal E)}.
\end{equation}
The height takes values in $\Q\log q$; a height divided by $\log q$
requires multiplying its rank-$r$ regulator by $(\log q)^r$.
The exponent is the coherent Euler characteristic $h^0-h^1$.
This is the elliptic specialization of
\cite[Conjecture~1.3, (1.10) and \S1.4]%
{LichtenbaumRamachandranSuzuki2022ArtinTateBSD}; its discriminant
includes the torsion denominator.

\begin{theorem}[Tate to function-field BSD]
\label{thm:bsd-function-field}
Let $C/\F_q$ be smooth, projective, and geometrically connected, and
let $K=\F_q(C)$. Let $\pi:X\to C$ be a proper flat morphism from a
smooth projective surface $X/\F_q$, with a zero section and smooth
geometrically
connected generic elliptic fibre $E/K$. If the divisor Tate
pole-order conjecture holds for $X$, then rank BSD, finiteness of
$\Sha(E/K)$, and \eqref{eq:bsd-function-field} hold for $E/K$.
\end{theorem}
\begin{proof}
The Tate--Shioda comparison identifies the defect between the pole
order of $\zeta(X,s)$ and $\operatorname{rank}\operatorname{NS}(X)$
with $\ord_{s=1}L(E/K,s)-\operatorname{rank}E(K)$; equivalently, the
rank equalities are equivalent
\cite[Proposition~2.4]{LichtenbaumRamachandranSuzuki2022ArtinTateBSD}.
The Artin--Tate--Milne theorem upgrades the Tate pole-order assertion
to the Artin--Tate special-value conjecture and finiteness of the
Brauer group of $X$
\cite[Theorem~6.1.4]{Ulmer2012CurvesJacobians}.
The Artin--Grothendieck comparison identifies
Brauer finiteness with $\Sha(E/K)$ finiteness. Kato--Trihan's
characteristic-$p$ theorem then yields the complete BSD statement,
including the characteristic-prime part. These imported implications
are assembled, under the stated surface hypotheses, in
\cite[Theorem~1.1 and \S1.1]%
{LichtenbaumRamachandranSuzuki2022ArtinTateBSD}; see also
\cite[p.~541]{KatoTrihan2003BSD}.
\end{proof}

Thus the divisor Tate conjecture is imposed on the \emph{surface}
$X$, while BSD concerns its generic elliptic curve over $K$.
This implication uses arithmetic comparison theorems in addition to
positivity or Frobenius purity. It is not a deduction of number-field
BSD from Grothendieck's standard conjectures. In particular, the
pole-order equality is not obtained by matching individual fixed
classes with points; the Tate--Shioda comparison subtracts the same
vertical and base contributions on both sides.

\begin{proposition}\label{prop:bsd-coefficient-gap}
Let $f$ be meromorphic near $s=1$ with nonzero leading coefficient,
and let $c>0$ with $c\ne1$. Then $f$ and $cf$ have the same divisor
but different leading coefficients. The Euclidean line $\R$ has
lattices $\Z$ and $2\Z$ with the same real span and metric but Gram
determinants $1$ and $4$.
\end{proposition}
\begin{proof}
Multiplying by $c$ scales the first nonzero Taylor coefficient
without creating a zero or pole. The displayed lattices have bases
$1$ and $2$, whose squared lengths are $1$ and $4$.
\end{proof}

The first example concerns divisor-only data: a specified Euler
product fixes its own normalization. The second shows that a complex
metric alone does not name an integral lattice or regulator.
Tensoring an arithmetic group with $\Q$ also discards finite torsion.

\begin{proposition}\label{prop:bsd-primary-gap}
Let $G=\bigoplus_{\ell\ \mathrm{prime}}\Z/\ell\Z$. For every prime
$\ell$, the $\ell$-primary subgroup $G[\ell^\infty]$ is finite,
but $G$ is infinite.
\end{proposition}
\begin{proof}
The $\ell$-primary subgroup is its one summand of order $\ell$;
there are infinitely many nonzero summands.
\end{proof}

This group example shows why prime-primary finiteness alone is not
the whole-group finiteness in BSD. It is not an example of an
elliptic curve with infinite $\Sha$.

Alternating cancellation can lose still more graded information, as
Proposition~\ref{prop:no-positivity-cancel} shows in a Weil datum.
For a number-field elliptic curve, the rational period--regulator
comparison and the integral comparisons at each prime in the
Tamagawa-number conjecture supply additional conjectural input for
the full number-field BSD formula; rank equality alone supplies
none of those factors
\cite[\S\S1--3, Conjectures~1--3 and \S5.4]{Flach2004ETNC}.
The elementary examples do not prove that all standard conjectures
are logically insufficient for BSD. They identify precisely the
missing comparisons in the direct positivity argument.

\subsection{The explicit formula and Weil's
criterion}\label{subsec:prior-weil}

Weil generalized Riemann's explicit formula to Hecke $L$-functions and
to curves over finite fields, wrote its geometric side as a sum of
local terms, and observed that the Riemann hypothesis is equivalent to
the positivity of the resulting Hermitian functional on test functions
of the form $g*\tilde g$ \cite{Weil1952Explicit}, as reported in
\cite[\S V, pp.~8--9]{Bombieri2000Clay}. He later extended the
criterion to Artin--Hecke $L$-functions, where positivity is equivalent
to the Riemann hypothesis together with Artin's conjecture, and called
the analogy with the function-field case, where positivity follows from
the index theorem on $C\times C$, the most serious argument he knew in
favour of these conjectures \cite[\S16, p.~16]{Weil1972Explicit}.
Yoshida and Bombieri studied the functional on test functions of
bounded support: positivity holds unconditionally when the support is
short enough that no prime power enters, and the Riemann hypothesis is
equivalent to nondegeneracy on every support
\cite[Proposition~1, p.~285, Theorems~1 and~2, pp.~310 and~321]%
{Yoshida1992Hermitian}, \cite[Theorems~1 and~12, pp.~191 and~226--227]%
{Bombieri2000Lincei}, \cite[Theorem, p.~2]{Burnol2000Explicites}.
Connes and Consani gave an operator-theoretic reason for the archimedean
positivity in the prime-free window
\cite[Theorem~1, p.~3]{ConnesConsani2020WeilPositivity}.

\subsection{Deninger's program}\label{subsec:prior-den}

Deninger showed that Serre's archimedean local factors are
zeta-regularized determinants of a derivation on an
infinite-dimensional ``archimedean cohomology''
\cite[Theorem~(4.1), p.~252]{Deninger1991Gamma}, and that the local
factors at the finite primes are regularized determinants of a
derivation whose eigenvalues are all the values $\rho$ with $q^\rho$ a
Frobenius eigenvalue
\cite[Proposition~(2.3), p.~142, and Theorem~(2.7), p.~144]%
{Deninger1992Local}. In \S3 of the same paper, labelled speculative, he
proposed an arithmetic site whose cohomology with its flow would give
the completed zeta function as an alternating product of regularized
determinants, with $\xi_{\mathrm D}(s)=\detreg((s-\Theta)/2\pi\mid
H^1)$ \cite[(3.1)--(3.2), pp.~146--147]{Deninger1992Local}, where
$\xi_{\mathrm D}(s)=2^{-1/2}(2\pi)^{-2}s(s-1)\pi^{-s/2}\Gamma(s/2)
\zeta(s)$ is Deninger's normalization (it differs from $\xi$ of
Definition~\ref{def:zeta} by the constant recorded after
Theorem~\ref{thm:imp-deninger-product}). The resulting prediction
$\prod_\rho(s-\rho)/2\pi=\xi_{\mathrm D}(s)$ was proved by Deninger for
$\Re s>1$ \cite[Theorem~3.3, p.~147]{Deninger1992Local}, extended to
the cut plane by Soul\'e (in a letter reported in
\cite[Remark~3.4(1), p.~147]{Deninger1992Local}; see
\cite{SchroeterSoule1994Deninger}, as reported in
\cite[p.~168]{Deninger1998Analogies}), and generalized to Hecke
$L$-series of number fields by Illies (Diplomarbeit, M\"unster 1992, as
reported in \cite[Theorem~7.18, p.~35]{Deninger1994Motivic}); an earlier
related formula for every number field, up to a constant and assuming
the Riemann hypothesis, is Kurokawa's
\cite[p.~24]{Kurokawa1988Parabolic}. In Section~7 of
\cite{Deninger1994Motivic}, again explicitly speculative, Deninger
described a hypothetical site with an $\R$-action, a Lefschetz trace
formula equivalent to Weil's explicit formula, the determinant formula,
Poincar\'e duality, and the argument that a flow-equivariant Hodge star
making $\Tr(f\cup*f')$ positive definite would force the eigenvalues
onto the critical lines, ``by adapting Serre's argument''
\cite[(7.2)--(7.13), pp.~29--33, and (7.19), p.~36]%
{Deninger1994Motivic}. This is the earliest statement of the
implication positivity $\Rightarrow$ Riemann hypothesis among the
sources cited here; it is restated in
\cite[pp.~168--169 and p.~172]{Deninger1998Analogies},
\cite[pp.~4--5]{Deninger2010HilbertPolya}, and
\cite[2.7, p.~3]{Deninger2022NoWeil}. The same paper already
contemplates a K\"unneth formula for $\Spec\Z\boxtimes\Spec\Z$ and
notes that the resulting regularized product over sums of zeros does
not make sense \cite[(7.32), pp.~41--42]{Deninger1994Motivic}.

Deninger's ICM address \cite{Deninger1998Analogies} rewrites the
explicit formula as a distributional Lefschetz formula, conjectures the
corresponding trace formula for flows on foliated spaces, and derives
the expected shape of ``$\Spec\Z$'': three-dimensional, with closed
orbits of length $\log p$ for the primes and one fixed point for the
infinite place \cite[(5)--(6), pp.~170--171, (20)--(22),
pp.~179--180, and p.~181]{Deninger1998Analogies}. The dynamical
analogue of the positivity argument is a theorem for Riemannian
foliations \cite[Theorem~2.1, p.~8]{Deninger2001NumberTheory}, and for
an ordinary elliptic curve over $\F_q$ the explicit formula is realized
as a transversal index formula on a solenoid
\cite{Deninger2002Explicit}; Deninger and Singhof showed that the
leafwise Hodge decomposition needed in general fails for some
foliations \cite[Proposition~1.2, p.~214]{DeningerSinghof2001Counterexample}.
Deninger's note on the Hilbert--P\'olya strategy restates the
determinant conjecture and the positivity argument and relates the
positive form on one eigenspace to the Beilinson--Bloch height pairing
\cite[Conjecture~2, p.~3, pp.~4--5, and pp.~7--9]%
{Deninger2010HilbertPolya}. Flach and Morin codified the expected
properties as a list of axioms, including an antilinear Hodge star with
positive form, and related them to unconditionally defined
Weil--Arakelov complexes \cite[Conjecture~2.3, p.~522, Conjecture~4.1,
and Theorem~4.3, p.~527]{FlachMorin2020DeningerWeilArakelov}. Finally
Deninger proved that no functorial cohomology with real coefficients
and an endomorphism with the expected Galois multiplicities can exist
for rings of integers, because of symplectic Artin representations of
root number $-1$, crediting the finite-field form of the obstruction
to Serre \cite[\S1, p.~1, and Theorem~2.13, p.~5]{Deninger2022NoWeil}.
None of these works constructs the full global dynamical cohomology
with the spectral realization and positive star required here.


\subsection{Products, the absolute point, and the adele class
space}\label{subsec:prior-cc}

\begingroup\emergencystretch=1em
Manin advocated Deninger's formalism together with Kurokawa's tensor
product of zeta functions, whose zeros are the sums of zeros, as
heuristics for ``absolute motives'' and Descartes powers of $\Spec\Z$
over an absolute point, and recorded the difficulty that for $\zeta\otimes
\zeta$ the zeros $\rho$ and $1-\rho$ make the construction ill-defined
\cite[p.~121, p.~132, and \S5.2, p.~158]{Manin1995Lectures}; Kurokawa's
definition restricts to zeros with imaginary parts of equal sign
\cite[p.~219]{Kurokawa1992MultipleZeta}. Connes realized the critical
zeros as an absorption spectrum on the adele class space and showed,
for global fields of positive characteristic, that a global trace
formula is equivalent to the Riemann hypothesis for all
Gr\"ossencharakter $L$-functions; the global formula is left open
\cite[Theorem~1, p.~13, Theorem~5, p.~42, and p.~3]%
{Connes1998TraceFormula}. Connes and Consani constructed the arithmetic
site and the scaling site, whose points recover the adele class space
modulo $\widehat\Z^\times$ and whose periodic orbits carry a
Riemann--Roch formula
\cite[Theorem~2.6]{ConnesConsani2014ArithmeticSite},
\cite[Theorem~1.1, p.~2]{ConnesConsani2016ArithmeticGeometry},
\cite[Theorems~5.1 and~6.7]{ConnesConsani2015ScalingSite},
\cite[Theorems~3.9 and~5.17]{ConnesConsani2017ScalingGeometry}, and
proposed to derive the Riemann hypothesis from a Riemann--Roch theorem
and an intersection theory on the square of the scaling site, which
remain to be found \cite[\S3, pp.~16--18, and p.~22]%
{ConnesConsani2018RiemannRoch}, \cite[\S4.3.2, pp.~20--21]%
{Connes2016EssayRH}; see also the dictionary of
\cite[\S2.4, p.~12]{ConnesConsaniMarcolli2009Weil} and the
one-dimensional Riemann--Roch theorem on $\overline{\Spec\Z}$ of
\cite[Theorem~1.1, p.~2]{ConnesConsani2023RiemannRochSpecZ}.
\par\endgroup

\subsection{De Branges spaces}\label{subsec:prior-db}

De Branges proposed to derive the Riemann hypothesis from positivity
conditions on the Hilbert space of entire functions with structure
function $E(z)=\xi(1-iz)$ and on its reciprocal space
\cite{DeBranges1986RHHilbert}, \cite{DeBranges1992Euler}, as reported in
\cite[Theorem~2, p.~4]{ConreyLi2000Positivity}; Conrey and Li showed by
high-precision numerical evaluation that these conditions fail, and
report a proof by Sarnak, without numerical calculation, for the
reciprocal space \cite[(3.2) and~(3.4), pp.~6--7, and Remark, p.~9]%
{ConreyLi2000Positivity}. Lagarias proved that
$E(z)=\xi(\tfrac12-iz)+\xi'(\tfrac12-iz)$ is a de Branges function if
and only if the Riemann hypothesis holds
\cite[Theorem~1, p.~9]{Lagarias2006HilbertEntire}, and that shifted
versions are de Branges functions unconditionally
\cite[Theorem~2.1, p.~8]{Lagarias2005Differenced}. Suzuki constructed
canonical systems with positive Hamiltonians for shifted $\xi$
unconditionally \cite[Theorem~2.3, p.~8]{Suzuki2012Canonical},
\cite[Theorem~2.4, p.~9]{Suzuki2016Hamiltonians}, and, assuming the
Riemann hypothesis, identified the completion of $C_c^\infty(\R)$ for
Weil's form with a model space of a de Branges space
\cite[Theorem~1.1, p.~2]{Suzuki2023WeilHilbert}; a later paper
corrects the even extension of the kernels used there
\cite[p.~27]{Suzuki2026WeilScrew}.

\subsection{Surveys and formalizations}\label{subsec:prior-surveys}

Surveys of this circle of ideas include Bombieri's and Sarnak's
problem descriptions \cite{Bombieri2000Clay}, \cite{Sarnak2004Clay},
Conrey's \cite{Conrey2003Notices}, Lagarias's \cite{Lagarias2006Survey},
Connes's \cite{Connes2000Frontiers}, \cite{Connes2016EssayRH}, the
physically oriented \cite{BerryKeating1999SIAM} and
\cite{SchumayerHutchinson2011Physics}, Milne's history of the finite-field
case \cite{Milne2016RHFF}, and Deninger's own surveys
\cite{Deninger1998Analogies}, \cite{Deninger2001NumberTheory},
\cite{Deninger2008Foliated}.

\begin{sloppypar}
Formalizations related to these notes include the following.
\begin{itemize}
\item The statement of the Riemann hypothesis in Lean, due to
  D.~Loeffler: \ghlink{Mathlib2026RiemannHypothesis}%
  {leanprover-community/mathlib4}.
\item An independent Lean formalization of the Hilbert--P\'olya
  reduction, contemporaneous with the one of Subsection~\ref{sec:lean}:
  \ghlink{Csillag2026HilbertPolyaLean}%
  {frostynumeral/hilbert-polya-lean-verification}.
\item Li's criterion: \ghlink{Bulka2026LiCriterionLean}%
  {nicholasbulka/li-criterion-rh-equivalence-lean}.
\item Weil's explicit formula for $\zeta$:
  \ghlink{Anthropic2026Zeta23}{anthropics/formal-math}.
\item The Riemann hypothesis for hyperelliptic curves over finite
  fields, by the method of Stepanov and Bombieri and without cohomology:
  \ghlink{MathInc2026RHCurves}{math-inc/RiemannHypothesisCurves}.
\end{itemize}
\end{sloppypar}
The overlaps with the formalization of these notes are listed in
Subsection~\ref{sec:lean}.



\section{Conventions and analytic preliminaries}\label{sec:conventions}

Throughout, Hermitian forms are linear in the first argument and
antilinear in the second. For a meromorphic function $f\not\equiv0$ on
$\C$ and $\rho\in\C$ we write $\ord_\rho f\in\Z$ for its order at
$\rho$; for an entire $f\not\equiv0$ its \emph{divisor} is the formal
sum $\divisor f=\sum_\rho \ord_\rho(f)[\rho]$.

Definition~\ref{def:zeta} is given in the introduction.

For $K=\Q$ one has $D_\Q=1$, $r_1=1$, $r_2=0$, so
$\Lambda_\Q=\Lambda$ and $\xi_\Q=2\xi$. The classical analytic facts
used below are collected, with sources, in
Section~\ref{sec:classical}. By Theorem~\ref{thm:imp-riemann}, $F$ is an
even real entire function of order one whose zeros, counted with
multiplicity, correspond to the nontrivial zeros of $\zeta$ under
$\rho=\tfrac12+i\tau\leftrightarrow z=-\tau$; the Riemann
hypothesis (RH) is the statement that all zeros of $F$ are real.

\begin{proposition}\label{prop:F0}
Let $F$ be as in Definition~\ref{def:zeta}. Then $F(0)=\xi(\tfrac12)>0$.
\end{proposition}

\begin{proof}
For real $0<s<1$ the alternating series
$\eta(s)=\sum_{n\ge1}(-1)^{n-1}n^{-s}$ converges, and grouping its terms
in consecutive pairs shows $\eta(s)>0$. On $\Re s>1$ one has
$\eta(s)=(1-2^{1-s})\zeta(s)$, and both sides are analytic on
$\Re s>0$, $s\ne1$ (the left by Dirichlet's test and locally uniform
convergence of the paired series), so the identity persists there. For
$0<s<1$ we have $1-2^{1-s}<0$, hence $\zeta(s)<0$. Since
$\Gamma(\tfrac14)>0$, $\Lambda(\tfrac12)<0$ and
$\xi(\tfrac12)=\tfrac12\cdot\tfrac12\cdot(-\tfrac12)\Lambda(\tfrac12)>0$.
\end{proof}

\begin{definition}\label{def:hb}
Let $\Cp=\{z\in\C:\Im z>0\}$.
\begin{enumerate}[label=(\alph*)]
\item The \emph{reflection} of a function $E$ defined on a set
  $\Omega\subset\C$ with $\overline\Omega=\Omega$ is the function
  $E^\#(z)=\overline{E(\bar z)}$ on $\Omega$.
\item An entire function $E$ satisfies the \emph{Hermite--Biehler (HB)
  inequality} if $|E^\#(z)|<|E(z)|$ for every $z\in\Cp$. (Common real
  zeros of $E$ and $E^\#$ are allowed; when a convention forbids them we
  say so.)
\item For an entire function $G$, its \emph{HB companion} is the entire
  function $E_G=G+iG'$.
\item The \emph{Hardy space} $H^2(\Cp)$ is the space of holomorphic $G$
  on $\Cp$ with $\sup_{y>0}\int_\R|G(x+iy)|^2\,dx<\infty$.
\end{enumerate}
\end{definition}

Definition~\ref{def:flow} is stated in the introduction.


Definition~\ref{def:inf-pol} is given in the introduction.

The identity \eqref{eq:inf-pol} is Deninger's
\cite[(7.13), p.~33]{Deninger1994Motivic}.
Equivalently, $\Theta-j/2$ is skew-symmetric for $h$. On a
finite-dimensional $\Theta$-stable block, \eqref{eq:inf-pol} integrates
to $h(\phi^tu,\phi^tv)=e^{jt}h(u,v)$: the flow is a similitude of
weight $j$.

Definition~\ref{def:spectral-core} is given in the introduction.

On a spectral core the flow of Definition~\ref{def:flow} is defined
blockwise on the generalized eigenspaces. Spectral cores are the
algebraic part of the spaces in
\cite[Definition~2.1(i), p.~521]{FlachMorin2020DeningerWeilArakelov}
and \cite[\S1]{Deninger1994Motivic}; the condition of no finite
accumulation point is added here. The Tate twist of
Definition~\ref{def:flow} is \cite[p.~168]{Deninger1998Analogies} and
\cite[p.~4]{Deninger2010HilbertPolya}.

The following linear algebra is used throughout, from the finite-field
calibration on.

\begin{proposition}[{Serre; Deninger
\cite[(7.11)--(7.13), pp.~32--33]{Deninger1994Motivic}}]%
\label{prop:purity}
Let $V$ be a complex vector space, $\Theta:V\to V$ a linear map,
$j\in\R$, and $h$ a positive infinitesimal polarization of weight $j$
on $(V,\Theta)$ (Definition~\ref{def:inf-pol}). Then:
\begin{enumerate}[label=(\roman*)]
\item\label{item:pur-real} every eigenvalue $\rho$ of $\Theta$ satisfies
  $\Re\rho=j/2$;
\item\label{item:pur-orth} eigenvectors of $\Theta$ with distinct
  eigenvalues are $h$-orthogonal;
\item\label{item:pur-ss} $\Theta$ is diagonalizable on every
  finite-dimensional $\Theta$-stable subspace of $V$; in particular, if
  $(V,\Theta)$ is a spectral core then $V_{\Theta\sim\rho}=
  \ker(\Theta-\rho)$ for every $\rho$.
\end{enumerate}
\end{proposition}

Throughout the proof put $A=\Theta-j/2$; by \eqref{eq:inf-pol},
$h(Au,v)=-h(u,Av)$.

\begin{proof}[Proof of \ref{item:pur-real}]
If $\Theta v=\rho v$ with $v\ne0$, then \eqref{eq:inf-pol} with $u=v$
gives $(\rho+\bar\rho-j)h(v,v)=0$, and $h(v,v)>0$.
\end{proof}

\begin{proof}[Proof of \ref{item:pur-orth}]
Let $\Theta u=\rho u$ and $\Theta v=\sigma v$ with $\rho\ne\sigma$. By
\ref{item:pur-real}, $\rho=j/2+i\tau$ and $\sigma=j/2+i\tau'$ with
$\tau\ne\tau'$ real. Then \eqref{eq:inf-pol} gives
$(\rho+\bar\sigma-j)h(u,v)=i(\tau-\tau')h(u,v)=0$.
\end{proof}

\begin{proof}[Proof of \ref{item:pur-ss}]
Let $W$ be finite-dimensional and $\Theta$-stable. If $\Theta|_W$ were
not diagonalizable, some eigenvalue $\rho$ would admit $w\in W$ with
$v=(\Theta-\rho)w\ne0$ and $(\Theta-\rho)v=0$. By
\ref{item:pur-real}, $\rho-j/2=i\tau$ with $\tau\in\R$, so
$\Theta-\rho=A-i\tau$, and $A-i\tau$ is again skew-symmetric because
$i\tau$ is purely imaginary. Hence
$h(v,v)=h((A-i\tau)w,v)=-h(w,(A-i\tau)v)=0$, contradicting $v\ne0$.
For a spectral core apply this to $W=V_{\Theta\sim\rho}$.
\end{proof}

Proposition~\ref{prop:purity} excludes Jordan blocks, not repeated
eigenvalues: it says nothing about simplicity.
It is the linear algebra of Serre's K\"ahler analogue of the Weil
conjectures \cite[Th\'eor\`eme~1, pp.~392--393]{Serre1960Analogues}, in
the infinitesimal form in which Deninger uses it
\cite[(7.11)--(7.13), pp.~32--33]{Deninger1994Motivic},
\cite[p.~169]{Deninger1998Analogies}; the semisimplicity
clause~\ref{item:pur-ss} is \cite[2.7, p.~3]{Deninger2022NoWeil}.


\subsection{Classical inputs}\label{sec:classical}

The following classical results are used without proof. Each is quoted
from a source in which its statement was checked; where the source
states the result without proving it, this is said.

\begin{theorem}[Riemann; Hadamard; de la Vall\'ee Poussin; von
Mangoldt]\label{thm:imp-riemann}
Let $\xi$ and $F$ be as in Definition~\ref{def:zeta}. Then:
\begin{enumerate}[label=(\roman*)]
\item\label{item:riem-entire} $\xi$ extends to an entire function with
  $\xi(s)=\xi(1-s)$, and all zeros of $\xi$ lie in $0<\Re s<1$
  \cite[Theorem~17.17]{Sutherland2017Lecture17};
\item\label{item:riem-order} $\xi$ is an entire function of order one,
  and the number $N(T)$ of its zeros $\rho$ with $0<\Im\rho\le T$,
  counted with multiplicity, satisfies $N(T)=O(T\log T)$ as
  $T\to\infty$ \cite[pp.~12--14]{PerezMarco2018NotesRH},
  \cite[p.~4]{Riemann1859Wilkins}.
\end{enumerate}
\end{theorem}

The source for~\ref{item:riem-entire} contains a proof; the sources
for~\ref{item:riem-order} state it (in the second, heuristically) with
references to the standard proofs. The absence of zeros of $\xi$ on
$\Re s=1$ in~\ref{item:riem-entire} is the nonvanishing of $\zeta$ on
that line (Hadamard, de la Vall\'ee Poussin). Since $\GR$ has no zeros
and its only poles are at $0,-2,-4,\dots$, the zeros of $\xi$ are
exactly the zeros of $\zeta$ in $0<\Re s<1$, with multiplicity: the
\emph{nontrivial zeros}. Because $\xi$ is real on $\R$,
$\xi(\bar s)=\overline{\xi(s)}$. Consequently $F$ is an entire function
of order one, even (from $\xi(s)=\xi(1-s)$), and real
($F(\bar z)=\overline{F(z)}$); its zeros lie in the strip
$|\Im z|<\tfrac12$ and correspond to the nontrivial zeros under
$\rho=\tfrac12+i\tau\leftrightarrow z=-\tau$; and the number of zeros
of $F$ in $|z|\le R$ is $O(R\log R)$, because such a zero corresponds
to a zero $\rho$ with $|\Im\rho|\le R$ and the zeros are symmetric
under $\rho\mapsto\bar\rho$.

\begin{theorem}[Hecke]\label{thm:imp-hecke}
Let $K$ be a number field and $\Lambda_K$ as in
Definition~\ref{def:zeta}. Then $\Lambda_K$ extends to a meromorphic
function on $\C$, holomorphic except for simple poles at $s=0$ and
$s=1$, and $\Lambda_K(s)=\Lambda_K(1-s)$
\cite[Remark~19.13]{Sutherland2017Lecture19}.
\end{theorem}

The source states this result, attributing it to Hecke, without proof,
and normalizes $\GC(s)=(2\pi)^{-s}\Gamma(s)$; the two normalizations
differ by the nonzero constant factor $2^{r_2}$, which changes neither
the poles nor the functional equation.

\begin{proposition}\label{prop:xiK}
Let $K$ be a number field and $\xi_K$ as in Definition~\ref{def:zeta}.
Then $\xi_K$ is entire, $\xi_K(s)=\xi_K(1-s)$, every zero of $\xi_K$
lies in the strip $0\le\Re s\le1$, $\xi_K(0)\xi_K(1)\ne0$, and for
$0\le\Re s\le1$, $s\notin\{0,1\}$, the order of $\xi_K$ at $s$ equals
the order of $\zeta_K$ at $s$. In particular every zero of $\zeta_K$
with $0<\Re s<1$ is a zero of $\xi_K$ of the same order.
\end{proposition}

\begin{proof}
By Theorem~\ref{thm:imp-hecke}, $s(s-1)\Lambda_K(s)$ has removable
singularities at $0$ and $1$, so $\xi_K$ is entire, and
$\xi_K(1-s)=(1-s)(-s)\Lambda_K(1-s)=\xi_K(s)$. A simple pole has
nonzero residue, so $\xi_K(1)=\operatorname{Res}_{s=1}\Lambda_K\ne0$,
and $\xi_K(0)=\xi_K(1)$. For $\Re s>1$ the Euler product
$\zeta_K(s)=\prod_{\mathfrak p}(1-N\mathfrak p^{-s})^{-1}$ converges
absolutely, so $\zeta_K(s)\ne0$; the factors $|D_K|^{s/2}$, $\GR(s)$,
$\GC(s)$ have no zeros; hence $\xi_K(s)\ne0$ for $\Re s>1$, and by the
functional equation for $\Re s<0$. For $0\le\Re s\le1$ and
$s\notin\{0,1\}$ the factors $s(s-1)|D_K|^{s/2}\GR(s)^{r_1}
\GC(s)^{r_2}$ are holomorphic and nonzero ($\Gamma$ has no zeros and
its poles are at $0,-1,-2,\dots$), so $\xi_K$ and $\zeta_K$ have the
same order there.
\end{proof}

For $K=\Q$ this recovers the correspondence between zeros of
$\xi_\Q=2\xi$ and nontrivial zeros of $\zeta$.

\begin{definition}\label{def:hasse-weil-zeta}
Let $q$ be a prime power.
\begin{enumerate}[label=(\alph*)]
\item The \emph{zeta function} of a variety $X$ over $\F_q$ is the power
  series
  \[
  Z(X,t)=\exp\Bigl(\sum_{n\ge1}\#X(\F_{q^n})\,\frac{t^n}n\Bigr),
  \]
  and we write $\zeta_X(s)=Z(X,q^{-s})$.
\item For an elliptic curve $E$ over $\F_q$, the \emph{Frobenius
  roots} of $E$ are the complex numbers $\alpha,\beta$ (with
  multiplicity) such that $\det(1-F^*t\mid H^1)=(1-\alpha t)
  (1-\beta t)$, where $\ell\nmid q$ is a prime, $F^*$ is geometric
  Frobenius on $H^1(E_{\overline{\F}_q},\Q_\ell)$ (as in
  Theorem~\ref{thm:imp-deligne} below), computed after embedding
  $\overline\Q_\ell$ in $\C$ (by Theorem~\ref{thm:imp-deligne} this
  polynomial has integer coefficients, so $\alpha,\beta$ depend on
  neither $\ell$ nor the embedding).
\end{enumerate}
\end{definition}

\begin{theorem}[{\cite[(1.5.4) and Th\'eor\`eme~(1.6), p.~276]%
{Deligne1974WeilI}}]\label{thm:imp-deligne}
Let $X_0$ be a smooth projective variety over $\F_q$, $X$ its base
change to $\overline{\F}_q$, $\ell\nmid q$ a prime, and $F^*$ the
geometric Frobenius acting on $H^i(X,\Q_\ell)$. Then
$Z(X_0,t)=\prod_i\det(1-F^*t\mid H^i(X,\Q_\ell))^{(-1)^{i+1}}$; each
$\det(1-F^*t\mid H^i(X,\Q_\ell))$ has integer coefficients independent
of $\ell$; and every complex conjugate of every root $\alpha$ of its
reciprocal polynomial satisfies $|\alpha|=q^{i/2}$.
\end{theorem}

These are the statements conjectured by Weil \cite[pp.~507--508]%
{Weil1949Numbers}. For an elliptic curve $E$ over $\F_q$ we also use
the standard facts that $H^0$, $H^1$, $H^2$ of $E_{\overline{\F}_q}$
have dimensions $1,2,1$ and that $F^*$ acts on $H^0$ by $1$ and on
$H^2$ by $q$; they are not part of the statement quoted above. With
them, Theorem~\ref{thm:imp-deligne} gives
\begin{equation}\label{eq:elliptic-zeta}
Z(E,t)=\frac{(1-\alpha t)(1-\beta t)}{(1-t)(1-qt)},\qquad
|\alpha|=|\beta|=q^{1/2},
\end{equation}
where $\alpha,\beta$ are the Frobenius roots of $E$
(Definition~\ref{def:hasse-weil-zeta}).

\subsection{The Weil form}\label{subsec:weil-form}

By Theorem~\ref{thm:imp-riemann} and the discussion after it, every
zero of $F$ lies in the strip $|\Im z|<\tfrac12$ and the number $N(R)$
of zeros of $F$ in $|z|\le R$, counted with multiplicity, satisfies
$N(R)=O(R\log R)$; by Proposition~\ref{prop:F0}, $F(0)\ne0$. If
$(r_n)$ lists the zeros of $F$ with multiplicity, then
\begin{equation}\label{eq:sum-r2}
\sum_n|r_n|^{-2}<\infty .
\end{equation}
Indeed, $\delta=\min_n|r_n|>0$ because the zeros are discrete and
$F(0)\ne0$, and grouping the zeros into the dyadic shells
$2^k\delta\le|r|<2^{k+1}\delta$ bounds the sum by
$\sum_{k\ge0}N(2^{k+1}\delta)(2^k\delta)^{-2}$, which converges since
$N(R)=O(R\log R)$.

\begin{definition}\label{def:weil-form}
Let $F$ be as in Definition~\ref{def:zeta}, let $\Gamma$ be the set of
distinct zeros of $F$, and for $\gamma\in\Gamma$ let
$m_\gamma=\ord_\gamma F\ge1$.
\begin{enumerate}[label=(\alph*)]
\item The \emph{Fourier transform} of a compactly supported
  $\psi\in L^1(\R)$ is the entire function
  $\widehat\psi(z)=\int_\R\psi(t)e^{izt}\,dt$, $z\in\C$.
\item The \emph{adjoint} of a function $\psi$ on $\R$ is
  $\widetilde\psi(t)=\overline{\psi(-t)}$.
\item The \emph{Weil distribution} is the linear functional
  $W(f)=\sum_{\gamma\in\Gamma}m_\gamma\widehat f(\gamma)$ on
  $C_c^\infty(\R)$; the series converges absolutely by
  Proposition~\ref{prop:weil-form}.
\item The \emph{Weil form} is $Q_W(\psi)=W(\psi*\widetilde\psi)$ for
  $\psi\in C_c^\infty(\R)$.
\end{enumerate}
\end{definition}

The Fourier inversion formula for this convention carries the factor
$1/2\pi$. The sign of the exponent does not affect any statement about
the location of zeros, because $F$ is even. The form $Q_W$ is Weil's
quadratic functional \cite[\S16, p.~16]{Weil1972Explicit}; see
\cite[\S V, p.~9]{Bombieri2000Clay} and
\cite[Theorem~1, p.~191]{Bombieri2000Lincei}.

\begin{proposition}\label{prop:weil-form}
Let $F$, $\Gamma$, $m_\gamma$, $W$, and $Q_W$ be as in
Definition~\ref{def:weil-form}, and let $f,\psi\in C_c^\infty(\R)$.
Then:
\begin{enumerate}[label=(\roman*)]
\item\label{item:wf-transform} $\psi*\widetilde\psi\in C_c^\infty(\R)$
  and $(\psi*\widetilde\psi)\widehat{\ }(z)
  =\widehat\psi(z)\,\overline{\widehat\psi(\bar z)}$ for all $z\in\C$;
\item\label{item:wf-conv} the series
  $\sum_{\gamma}m_\gamma\widehat f(\gamma)$ converges absolutely, and
  \[
  Q_W(\psi)=\sum_{\gamma\in\Gamma}m_\gamma\,\widehat\psi(\gamma)\,
    \overline{\widehat\psi(\bar\gamma)},
  \]
  an absolutely convergent series whose sum is real;
\item\label{item:wf-rh} if every zero of $F$ is real, then
  $Q_W(\psi)=\sum_{\gamma}m_\gamma|\widehat\psi(\gamma)|^2\ge0$.
\end{enumerate}
\end{proposition}

\begin{proof}[Proof of \ref{item:wf-transform}]
Convolution of compactly supported smooth functions is compactly
supported and smooth. Since $e^{iz(s+u)}=e^{izs}e^{izu}$, Fubini's
theorem gives $(g*h)\widehat{\ }=\widehat g\,\widehat h$ for compactly
supported integrable $g,h$. Moreover, substituting $u=-t$,
\[
\widehat{\widetilde\psi}(z)=\int_\R\overline{\psi(u)}e^{-izu}\,du
=\overline{\int_\R\psi(u)e^{i\bar zu}\,du}
=\overline{\widehat\psi(\bar z)},
\]
because $\overline{e^{-izu}}=e^{i\bar zu}$ for real $u$.
\end{proof}

\begin{proof}[Proof of \ref{item:wf-conv}]
Let $\operatorname{supp}f\subset[-a,a]$. For $|\Im z|\le\tfrac12$ and
$t\in[-a,a]$ one has $|e^{izt}|\le e^{a/2}$. Integrating by parts $k$
times, $(-iz)^k\widehat f(z)=\widehat{f^{(k)}}(z)$, so
$|z|^k|\widehat f(z)|\le e^{a/2}\|f^{(k)}\|_1$. Using $k=0$ and $k=2$
and $(1+|z|)^2\le2(1+|z|^2)$,
\[
|\widehat f(z)|\le C_f(1+|z|)^{-2}\quad(|\Im z|\le\tfrac12),\qquad
C_f=2e^{a/2}\bigl(\|f\|_1+\|f''\|_1\bigr).
\]
All zeros satisfy $|\Im\gamma|<\tfrac12$, so
$\sum_\gamma m_\gamma|\widehat f(\gamma)|\le C_f\sum_n|r_n|^{-2}<\infty$
by \eqref{eq:sum-r2}. Applying this to
$f=\psi*\widetilde\psi$ and using \ref{item:wf-transform} gives the
displayed formula. Since $F$ is real, $\gamma\mapsto\bar\gamma$ permutes
$\Gamma$ with $m_{\bar\gamma}=m_\gamma$, and the term indexed by
$\bar\gamma$ is the complex conjugate of the term indexed by $\gamma$;
by absolute convergence the sum is real.
\end{proof}

\begin{proof}[Proof of \ref{item:wf-rh}]
If $\gamma$ is real then $\bar\gamma=\gamma$ and the $\gamma$-term is
$m_\gamma|\widehat\psi(\gamma)|^2$.
\end{proof}

\section{Finite-field calibration}\label{sec:finite-field}

This section isolates the linear algebra by which positivity forces the
Riemann hypothesis for varieties over $\F_q$, in the two forms in which
it occurs: a positive form on the cohomology for which Frobenius is a
similitude, and a positive involution on an algebra of operators
containing Frobenius. Weil's theorem for curves is an instance of the
second form, not of the first. The section then shows how either form
becomes, after logarithmic suspension of Frobenius, exactly the
hypothesis of the purity statement used over number fields
(Proposition~\ref{prop:purity}). Section~\ref{subsec:thmB} places the
result in the frame of the axioms of Section~\ref{sec:axioms}. In
applications $q$ is a prime power; the linear algebra only uses that
$q>1$ is real.

It is natural to expect that a finite-dimensional graded cohomology
with a Frobenius, perfect Poincar\'e duality, a rational zeta function
with functional equation, and integral point counts already forces
$|\alpha|=q^{j/2}$, and that any nondegenerate Frobenius-compatible form
does; in fact none of this suffices, and positivity is the whole
content (Propositions~\ref{prop:formal-curve}
and~\ref{prop:weil-converse}).

\subsection{Frobenius polarizations}\label{subsec:frob-pol}

\begin{definition}\label{def:frob-pol}
Let $q>1$ be real, $j\in\R$, $V$ a finite-dimensional complex vector
space, and $\mathrm{F}:V\to V$ a linear map. A \emph{positive Frobenius
polarization of weight $j$} on $(V,\mathrm{F})$ is a positive definite
Hermitian form $h$ on $V$ with
$h(\mathrm{F}u,\mathrm{F}v)=q^jh(u,v)$ for all $u,v\in V$.
\end{definition}

This is the vector-space form of the hypothesis of Serre's argument
\cite[Th\'eor\`eme~1 and its proof, pp.~392--393]{Serre1960Analogues}.

\begin{proposition}[{Serre \cite[p.~393]{Serre1960Analogues}; Kleiman
\cite[Lemma~4.6, p.~385]{Kleiman1968Cycles}}]\label{prop:weil-linear}
Let $q>1$ be real, $j\in\R$, $V$ a finite-dimensional complex vector
space, $\mathrm{F}:V\to V$ linear, and $h$ a positive Frobenius
polarization of weight $j$ on $(V,\mathrm{F})$. Then $\mathrm{F}$ is
diagonalizable in an $h$-orthonormal basis, distinct eigenspaces of
$\mathrm{F}$ are $h$-orthogonal, and every eigenvalue $\alpha$ of
$\mathrm{F}$ satisfies $|\alpha|=q^{j/2}$.
\end{proposition}

\begin{proof}
Put $U=q^{-j/2}\mathrm{F}$. Then $h(Uu,Uv)=h(u,v)$, so $U$ is an
isometry of the finite-dimensional inner-product space $(V,h)$, hence
injective and therefore unitary. A unitary operator on a
finite-dimensional inner-product space is diagonalizable in an
orthonormal basis with eigenvalues of modulus one (spectral theorem for
normal operators), and its eigenspaces are the spans of the basis
vectors with a given eigenvalue, hence mutually orthogonal. Hence
$\mathrm{F}=q^{j/2}U$ is diagonalizable in the same basis, with the
same eigenspaces, and its eigenvalues have modulus $q^{j/2}$.
\end{proof}

\begin{proposition}\label{prop:weil-converse}
Let $q>1$ be real, $j\in\R$, $V$ a finite-dimensional complex vector
space, and $\mathrm{F}:V\to V$ linear.
\begin{enumerate}[label=(\roman*)]
\item\label{item:wc-exist} $(V,\mathrm{F})$ admits a positive Frobenius
  polarization of weight $j$ if and only if $\mathrm{F}$ is
  diagonalizable and all its eigenvalues have modulus $q^{j/2}$.
\item\label{item:wc-all} In that case the positive Frobenius
  polarizations of weight $j$ are exactly the orthogonal sums
  $\bigoplus_\alpha h_\alpha$ of arbitrary positive definite Hermitian
  forms $h_\alpha$ on the eigenspaces $V_\alpha$ of $\mathrm{F}$.
\end{enumerate}
\end{proposition}

\begin{proof}[Proof of \ref{item:wc-exist}]
One direction is Proposition~\ref{prop:weil-linear}. Conversely, the
form described in~\ref{item:wc-all} is a positive Frobenius
polarization: for $u\in V_\alpha$, $v\in V_\beta$,
$h(\mathrm{F}u,\mathrm{F}v)=\alpha\bar\beta h(u,v)$, which is $0=q^j\cdot
0$ if $\alpha\ne\beta$ and $|\alpha|^2h(u,v)=q^jh(u,v)$ if
$\alpha=\beta$.
\end{proof}

\begin{proof}[Proof of \ref{item:wc-all}]
The computation just made shows that every such orthogonal sum is a
positive Frobenius polarization. Conversely, a positive Frobenius
polarization makes distinct eigenspaces orthogonal
(Proposition~\ref{prop:weil-linear}), so it is the orthogonal sum of its
restrictions to the eigenspaces.
\end{proof}

On a single finite-dimensional space, the existence of a positive
Frobenius polarization is therefore \emph{equivalent} to purity
together with semisimplicity, and when one exists it is far from
unique. The content of the positivity hypothesis lies entirely in
producing a form from geometry; the same phenomenon over $\Spec\Z$ is
Proposition~\ref{prop:content}.

It is natural to expect that purity of the eigenvalues is all a
positive Frobenius polarization encodes; in fact it also encodes
semisimplicity, and purity alone does not produce one.

\begin{proposition}\label{prop:N-jordan}
Let $q>1$ be real, $j\in\R$, $V=\C^2$, and
$\mathrm{F}=\begin{pmatrix}q^{j/2}&1\\0&q^{j/2}\end{pmatrix}$. Then
every eigenvalue of $\mathrm{F}$ has absolute value $q^{j/2}$, and
$(V,\mathrm{F})$ admits no positive Frobenius polarization of weight
$j$.
\end{proposition}

\begin{proof}
The only eigenvalue is $q^{j/2}$, and $\mathrm{F}$ is not
diagonalizable; Proposition~\ref{prop:weil-linear} excludes a positive
Frobenius polarization.
\end{proof}

For a nonsemisimple Frobenius one must therefore specify separately how
its nilpotent part is transported to a generator
(Definition~\ref{def:log-susp}); a positive realization does not follow
from absolute values.

\begin{definition}\label{def:formal-curve}
The \emph{formal curve} is the graded $\C$-algebra $H^\bullet$ with
$H^0=\C\cdot1$, $H^1=\C e_1\oplus\C e_2$, $H^2=\C\eta$, unit $1$,
products $e_1\cup e_2=-e_2\cup e_1=\eta$ and all other products of
positive-degree elements zero, trace $\Tr(c\eta)=c$, and the graded
linear map $\mathrm{F}$ equal to $1$ on $H^0$, to $\diag(3,1)$ on $H^1$
in the basis $e_1,e_2$, and to $3$ on $H^2$; put $q=3$. Its
\emph{Lefschetz numbers} are $N_n=\sum_i(-1)^i\Tr(\mathrm{F}^n\mid H^i)$
for $n\ge1$, and its \emph{zeta function} is
$Z(t)=\prod_i\det(1-\mathrm{F}t\mid H^i)^{(-1)^{i+1}}$.
\end{definition}

\begin{proposition}\label{prop:formal-curve}
Let $H^\bullet$, $\mathrm{F}$, $q=3$, $N_n$, and $Z(t)$ be as in
Definition~\ref{def:formal-curve}.
\begin{enumerate}[label=(\roman*)]
\item\label{item:fc-weil} $\mathrm{F}$ is an automorphism of the graded
  algebra; the pairing $B(u,v)=\Tr(u\cup v)$ on $H^1$ is alternating
  and perfect with $B(\mathrm{F}u,\mathrm{F}v)=qB(u,v)$; $N_n=0\in\Z$
  for all $n$; $Z(t)=P_1(t)/((1-t)(1-qt))$ with $P_1(t)=(1-t)(1-3t)
  \in\Z[t]$, and $q\,t^2P_1(1/(qt))=P_1(t)$.
\item\label{item:fc-impure} The eigenvalues $3$ and $1$ of $\mathrm{F}$
  on $H^1$ do not have modulus $q^{1/2}$.
\item\label{item:fc-forms} The Hermitian forms $h$ on $H^1$ with
  $h(\mathrm{F}u,\mathrm{F}v)=q\,h(u,v)$ are exactly those with
  $h(e_1,e_1)=h(e_2,e_2)=0$; the nonzero ones are nondegenerate of
  signature $(1,1)$, and none is positive definite.
\item\label{item:fc-alg} On the subalgebra $H^0\oplus H^2$, which
  contains the classes of the unit and of a point, the pairing
  $(a\cdot1+b\eta,a'\cdot1+b'\eta)\mapsto ab'+ba'$ is nondegenerate and
  $\mathrm{F}$-compatible of weight one, and $\mathrm{F}$ is semisimple
  with eigenvalues $1$ and $q$ there.
\end{enumerate}
\end{proposition}

\begin{proof}[Proof of \ref{item:fc-weil}]
$\mathrm{F}(e_1\cup e_2)=3\eta=\mathrm{F}e_1\cup\mathrm{F}e_2$, and the
other products are compatible trivially. $B(e_1,e_2)=1=-B(e_2,e_1)$ and
$B(e_i,e_i)=0$, so $B$ is alternating and perfect, and
$B(\mathrm{F}e_1,\mathrm{F}e_2)=3=qB(e_1,e_2)$. Next
$N_n=1-(3^n+1)+3^n=0$. The formula for $Z$ is the definition, and
$3t^2P_1(1/(3t))=3t^2(1-\tfrac1{3t})(1-\tfrac1t)=(3t-1)(t-1)=P_1(t)$.
\end{proof}

\begin{proof}[Proof of \ref{item:fc-impure}]
$|3|=3$ and $|1|=1$, while $q^{1/2}=\sqrt3$.
\end{proof}

\begin{proof}[Proof of \ref{item:fc-forms}]
Write $h_{kl}=h(e_k,e_l)$. The condition reads
$f_kf_lh_{kl}=3h_{kl}$ with $(f_1,f_2)=(3,1)$, that is $9h_{11}=3h_{11}$,
$h_{22}=3h_{22}$, and $3h_{12}=3h_{12}$. So $h_{11}=h_{22}=0$ and
$h_{12}=\overline{h_{21}}$ is arbitrary; the Gram matrix
$\bigl(\begin{smallmatrix}0&c\\\bar c&0\end{smallmatrix}\bigr)$ has
eigenvalues $\pm|c|$.
\end{proof}

\begin{proof}[Proof of \ref{item:fc-alg}]
The pairing is $(u,v)\mapsto\Tr(u\cup v)$ restricted to $H^0\oplus H^2$
(with $\Tr$ extended by zero on $H^0$), its Gram matrix in the basis
$1,\eta$ is $\bigl(\begin{smallmatrix}0&1\\1&0\end{smallmatrix}\bigr)$,
and $\Tr(\mathrm{F}u\cup\mathrm{F}v)=q\Tr(u\cup v)$ because
$\mathrm{F}$ is multiplicative and is $q$ on $H^2$.
\end{proof}


The formal curve is linear algebra, not the cohomology of a curve:
$N_n=0$ for all $n$ is impossible for a curve over $\F_3$ by the
Hasse--Weil bound itself. Item~\ref{item:fc-alg} shows that
nondegeneracy and Frobenius compatibility of the intersection pairing on
the algebraic classes of $X$ itself (for a curve, the classes of degree
$0$ and $2$) say nothing about $H^1$; since $\mathrm{F}$ has eigenvalues
$1$ and $q$ there, no Frobenius-compatible pairing on these classes can
be positive, so positivity has to be sought elsewhere. The positivity used in
Weil's proof
lives instead on the algebra of correspondences
(Section~\ref{subsec:weil-curves}), and in Grothendieck's formulation it
is the Hodge standard conjecture on $X\times X$ together with the
Lefschetz-type standard conjecture
\cite[Introduction, pp.~359--360]{Kleiman1968Cycles},
\cite[Definition~3.6, p.~39, and Theorems~3.15 and~3.17, p.~42]%
{Milne2016RHFF}. The finite-field form of the quaternionic block of
Proposition~\ref{prop:R2}, with $\mathrm{F}=q^\Theta$, is
Proposition~\ref{prop:R2-frob}.

\subsection{The algebra form of positivity}\label{subsec:algebra-form}

\begin{definition}\label{def:pos-involution}
Let $k\in\{\Q,\R,\C\}$ and $A$ a finite-dimensional unital
$k$-algebra. A \emph{positive trace involution} on $A$ is a pair
$(\tau,\dagger)$ of a $k$-linear functional $\tau:A\to k$ with
$\tau(ab)=\tau(ba)$ and a map $a\mapsto a^\dagger$ with
$(ab)^\dagger=b^\dagger a^\dagger$, $a^{\dagger\dagger}=a$, and
$\tau(aa^\dagger)>0$ for $a\ne0$ (a real number), where $\dagger$ is
$k$-linear if $k\ne\C$; if $k=\C$, $\dagger$ is antilinear and
moreover $\tau(a^\dagger)=\overline{\tau(a)}$.
\end{definition}

For $k=\C$ it enters the arithmetic axioms
(Definition~\ref{def:aw4-forms}). For the algebra of correspondences
with the transpose and its trace, positivity is
\cite[3.11 and Lemma~4.6]{Kleiman1968Cycles}.

\begin{proposition}[{Kleiman \cite[Lemma~4.6, p.~385]{Kleiman1968Cycles}}]%
\label{prop:algebra-form}
Let $k\in\{\Q,\R,\C\}$, $A$ a finite-dimensional unital $k$-algebra
with a positive trace involution $(\tau,\dagger)$, $q>1$ real, $j\in\R$,
and $g\in A$ with $gg^\dagger=q^j$.
\begin{enumerate}[label=(\roman*)]
\item\label{item:af-alg} The minimal polynomial of $g$ over $k$ is
  separable, and each of its complex roots has modulus $q^{j/2}$.
\item\label{item:af-rep} If $V$ is a finite-dimensional complex vector
  space and $\iota:A\otimes_k\C\to\operatorname{End}_\C(V)$ an injective
  unital $\C$-algebra homomorphism, then $\iota(g)$ is diagonalizable and all
  its eigenvalues have modulus $q^{j/2}$.
\end{enumerate}
\end{proposition}

\begin{proof}[Proof of \ref{item:af-alg}]
In the finite-dimensional algebra $A$, $gg^\dagger=q^j$ makes $g$
invertible with inverse $q^{-j}g^\dagger$, so also $g^\dagger g=q^j$.
We construct a positive definite Hermitian form $h$ on
$A_\C=A\otimes_k\C$ with $h(ga,gb)=q^jh(a,b)$, where $g$ acts by left
multiplication $L_g$. If $k=\C$, put $h(a,b)=\tau(ab^\dagger)$: it is
linear in $a$, antilinear in $b$, Hermitian because
$\overline{h(b,a)}=\tau((ba^\dagger)^\dagger)=\tau(ab^\dagger)$, positive
definite by hypothesis, and $h(ga,gb)=\tau(gab^\dagger g^\dagger)=
\tau(ab^\dagger g^\dagger g)=q^jh(a,b)$. If $k\ne\C$, let
$S(a,b)=\tfrac12(\tau(ab^\dagger)+\tau(ba^\dagger))$ on
$A_\R=A\otimes_k\R$. It is symmetric, satisfies $S(ga,gb)=q^jS(a,b)$ by
the same computation, and is positive definite: its Gram matrix in a
$k$-basis of $A$ has entries in $k$ and is positive on nonzero vectors
with coordinates in $k$, hence positive semidefinite over $\R$; a real
null vector would lie in the kernel of this matrix, which is spanned by
vectors with coordinates in $k$, and there are none. Let $h$ be the
Hermitian extension of $S$ to $A_\C=A_\R\otimes_\R\C$. In either case,
Proposition~\ref{prop:weil-linear} applied to $(A_\C,L_g,h)$ shows that
$L_g$ is diagonalizable with eigenvalues of modulus $q^{j/2}$, so its
minimal polynomial is separable with roots of that modulus. Since $A$ is
unital, $p(L_g)=0$ if and only if $p(g)=0$ (evaluate at $1$, and
conversely $p(L_g)a=p(g)a$), for polynomials $p$ over $\C$; and the
minimal polynomial of $g$ over $\C$ in $A_\C$ equals its minimal
polynomial over $k$ in $A$.
\end{proof}

\begin{proof}[Proof of \ref{item:af-rep}]
Since $\iota$ is an injective unital algebra homomorphism,
$p(\iota(g))=0$ if
and only if $p(g)=0$ in $A_\C$, so $\iota(g)$ has the minimal polynomial
of $g$, which by~\ref{item:af-alg} is separable with roots of modulus
$q^{j/2}$.
\end{proof}

Proposition~\ref{prop:algebra-form} is Kleiman's form of the positivity
argument \cite[Lemma~4.6, p.~385]{Kleiman1968Cycles}, which is stated
for the normalized Frobenius with $g'\circ g=\mathrm{id}$ and includes
the separability of~\ref{item:af-alg}; the form $gg^\dagger=q^j$ is
\cite[Proposition~3.10]{Milne2016RHFF}. With $g$ replaced by a
generator $\theta$ and $gg^\dagger=q^j$ by $\theta+\theta^\dagger=j$,
the same argument (with Proposition~\ref{prop:purity} in place of
Proposition~\ref{prop:weil-linear}) gives
Proposition~\ref{prop:aw4-forms}\ref{item:aw4e-pure}, which in this
additive form is Deninger's identity
\cite[(7.13), p.~33]{Deninger1994Motivic}; the two are
related by $g=q^\theta$ (Proposition~\ref{prop:susp-converse}).

\subsection{Weil's theorem for curves}\label{subsec:weil-curves}

\begin{theorem}[Grothendieck; Deligne]\label{thm:imp-weil-coh}
Let $X$ be a smooth projective geometrically connected variety of
dimension $d$ over $\F_q$, $\ell\nmid q$ a prime,
$H^j=H^j(X_{\overline\F_q},\Q_\ell)$, and $\mathrm{F}$ the geometric
Frobenius. Then $H^j$ is finite-dimensional and vanishes outside
$0\le j\le2d$; the cup product makes $H^\bullet$ an associative
graded-commutative algebra with unit, and $\mathrm{F}$ is an algebra
automorphism; $H^0=\Q_\ell\cdot1$ with $\mathrm{F}=1$, and $H^{2d}$ is
one-dimensional with $\mathrm{F}=q^d$; there is a linear isomorphism
$\operatorname{tr}:H^{2d}\to\Q_\ell$, and the pairings
$(u,v)\mapsto\operatorname{tr}(u\cup v)$ on $H^j\times H^{2d-j}$ are
perfect \cite[\S3, pp.~10--11]{Kleiman1994Standard},
\cite[Theorem~24.1, p.~144, \S25, p.~147, and proof of Theorem~27.12,
p.~160]{MilneLEC2013} (the duality theorem there is stated for torsion
coefficients, from which the statement for $\Q_\ell$ follows by passage
to the limit). If
$X$ is a curve, $\mathrm{F}$ acts semisimply on $H^1$ (Weil;
\cite[p.~10]{Kleiman1994Standard}).
\end{theorem}

Here $\mathrm{F}$ is the pullback by the $q$-Frobenius endomorphism, an
algebra endomorphism by functoriality, and it is injective because
$\operatorname{tr}(\mathrm{F}u\cup\mathrm{F}v)=q^d\operatorname{tr}
(u\cup v)$ and the pairings are perfect; the trace is identified with
$\Q_\ell$ after choosing a generator of $\Z_\ell(1)$, which does not
affect anything below.


\begin{theorem}[Weil; {\cite{MilneAV2008}}]\label{thm:imp-rosati}
Let $C$ be a smooth projective geometrically connected curve of genus
$g\ge1$ over $\F_q$, $J$ its Jacobian, $\pi\in\operatorname{End}(J)$ the
$q$-Frobenius endomorphism, $\lambda$ the canonical principal
polarization, $\operatorname{End}^0(J)=\operatorname{End}(J)\otimes\Q$,
$\alpha\mapsto\alpha^\dagger=\lambda^{-1}\alpha^\vee\lambda$ the Rosati
involution, $\ell\nmid q$ a prime, and $V_\ell J$ the rational
$\ell$-adic Tate module. Then:
\begin{enumerate}[label=(\roman*)]
\item\label{item:ro-char} for $\alpha\in\operatorname{End}(J)$ the
  characteristic polynomial $P_\alpha$ of $\alpha$ on $V_\ell J$ has
  integer coefficients independent of $\ell$, and the map
  $\operatorname{End}(J)\otimes\Q_\ell\to\operatorname{End}(V_\ell J)$
  is injective \cite[Theorem~10.9, p.~46, Theorem~10.15, p.~49, and
  Proposition~10.20, p.~50]{MilneAV2008};
\item\label{item:ro-pos} the bilinear form $(\alpha,\beta)\mapsto
  \Tr(\alpha\circ\beta^\dagger)$ on $\operatorname{End}^0(J)$, with
  $\Tr$ the trace on $V_\ell J$, is positive definite
  \cite[Theorem~14.3, p.~62]{MilneAV2008};
\item\label{item:ro-pi} $\pi\pi^\dagger=q$
  \cite[Chapter~II, Lemma~1.2, p.~76]{MilneAV2008};
\item\label{item:ro-count} $\#C(\F_q)=1-\sum_ia_i+q$, where $a_1,\dots,
  a_{2g}$ are the roots of $P_\pi$
  \cite[Chapter~III, Theorem~11.1, p.~118]{MilneAV2008}.
\end{enumerate}
\end{theorem}

The source proves these statements except that it defers the
calculation behind~\ref{item:ro-pos} to other work; a proof of the
positivity by intersection theory is also given in
\cite[Theorem~1.27, p.~21]{Milne2016RHFF}.
The positivity~\ref{item:ro-pos} is due to Weil
\cite{Weil1948VarietesAbeliennes}, as reported in
\cite[p.~23]{Milne2016RHFF}.

\begin{proposition}[Weil]\label{prop:weil-curves}
Let $C$, $g\ge1$, $J$, $\pi$, and $\ell$ be as in
Theorem~\ref{thm:imp-rosati}, and let $\iota:\overline\Q_\ell\to\C$ be
a field isomorphism. Then $\pi$ acts semisimply on
$V_\ell J\otimes_{\Q_\ell,\iota}\C$, every root $a$ of $P_\pi$ satisfies
$|\gamma(a)|=q^{1/2}$ for every field embedding $\gamma:\Q(a)\to\C$,
and $|\#C(\F_{q^n})-q^n-1|\le2g\,q^{n/2}$ for all $n\ge1$.
\end{proposition}

\begin{proof}
By Theorem~\ref{thm:imp-rosati}\ref{item:ro-char}, $\tau(\alpha)=
\Tr(\alpha\mid V_\ell J)$ is a $\Q$-valued linear functional on
$A=\operatorname{End}^0(J)$ with $\tau(ab)=\tau(ba)$, and
by~\ref{item:ro-pos} it forms with the Rosati involution a positive
trace involution over $k=\Q$ (Definition~\ref{def:pos-involution}).
By~\ref{item:ro-pi}, $\pi\pi^\dagger=q$. By
Proposition~\ref{prop:algebra-form}\ref{item:af-alg} with $j=1$, the
minimal polynomial $m$ of $\pi$ over $\Q$ is separable with all complex
roots of modulus $q^{1/2}$. Since $\operatorname{End}^0(J)$ acts on
$V_\ell J$ through an algebra homomorphism, $m(\pi)=0$ on $V_\ell J$, so $\pi$
is semisimple there (its minimal
polynomial divides the separable $m$), and every root of $P_\pi$ is a
root of $m$; the conjugates of a root of $m$ are roots of $m$. For the
last assertion apply~\ref{item:ro-count} to $C$ over $\F_{q^n}$, whose
Frobenius endomorphism is $\pi^n$ with characteristic polynomial roots
$a_i^n$ on $V_\ell J$.
\end{proof}

The same instance can be run on the surface $C\times C$. Positivity of
the trace form $X\mapsto\Tr(X\circ X')$ on correspondences modulo the
trivial ones expresses the positivity of the Rosati involution, the
transpose $X'$ corresponding to $\dagger$ and the graph of Frobenius to
$\pi$ \cite[fn.~7, p.~5]{Milne2016RHFF}; on $C\times C$ it follows from
the Hodge index theorem \cite[Theorem~V.1.9, p.~364]{Hartshorne1977},
and \cite[Exercises~V.1.9 and~V.1.10, p.~368]{Hartshorne1977} derive
from it the inequality $|N-q-1|\le2g\sqrt q$ directly. Both routes prove
$\tau(aa^\dagger)>0$ for the same algebra and then apply
Proposition~\ref{prop:algebra-form}; they differ only in how positivity
is proved.

It is natural to expect that Weil's positivity is a canonical positive
Frobenius polarization on $H^1(C)$, so that Weil's theorem is an
instance of Proposition~\ref{prop:weil-linear}; in fact the positive
form that Weil constructs lives on the endomorphism algebra
(Proposition~\ref{prop:weil-curves}), the complex space
$H^1(C)\otimes_\iota\C$ carries no complex conjugation induced by its
$\ell$-adic structure, and the Frobenius polarizations on it form the
family of Proposition~\ref{prop:weil-converse}\ref{item:wc-all}, none of
which is singled out.

\begin{proposition}\label{prop:ff-N7}
Let $C$ be a smooth projective geometrically connected curve of genus
$g\ge1$ over $\F_q$, $\ell\nmid q$ a prime, $\iota:\overline\Q_\ell\to\C$
a field isomorphism, $H=H^1(C_{\overline{\F}_q},\Q_\ell)$, and
$V=H\otimes_{\Q_\ell,\iota}\C$ with geometric Frobenius
$\mathrm{F}$.
\begin{enumerate}[label=(\roman*)]
\item\label{item:n7-conj} The rule $v\otimes z\mapsto v\otimes\bar z$
  does not define a map $V\to V$.
\item\label{item:n7-exist} $(V,\mathrm{F})$ admits positive Frobenius
  polarizations of weight one, and they are exactly the orthogonal sums
  of arbitrary positive definite forms on the eigenspaces of
  $\mathrm{F}$.
\end{enumerate}
\end{proposition}

\begin{proof}[Proof of \ref{item:n7-conj}]
By Hensel's lemma $\Q_\ell$ contains a square root $x$ of $-m$ for some
positive integer $m$ (take $m=7$ if $\ell=2$, and $m$ with $-m$ a
nonzero square modulo $\ell$ otherwise). Since $\iota$ fixes $\Q$,
$\iota(x)^2=-m$, so $\iota(x)$ is purely imaginary and nonzero. For
$v\ne0$ in $H$, the vector $xv\otimes1=v\otimes\iota(x)$ would be sent
by the rule both to $xv\otimes1$ and to $v\otimes\overline{\iota(x)}=
-v\otimes\iota(x)=-(xv\otimes1)$, which differ.
\end{proof}

\begin{proof}[Proof of \ref{item:n7-exist}]
By Theorem~\ref{thm:imp-weil-coh}, $\mathrm{F}$ is semisimple on $V$,
and by Theorem~\ref{thm:imp-deligne} its eigenvalues have modulus
$q^{1/2}$. Apply Proposition~\ref{prop:weil-converse}.
\end{proof}

Thus the positivity that Weil proves is an instance of the algebra form
(Proposition~\ref{prop:algebra-form}), with $A=\operatorname{End}^0(J)$,
the Rosati involution, and $g=\pi$: a multiplicative relation
$\pi\pi^\dagger=q$, which Section~\ref{subsec:mult-form} recasts over
$\C$ as the multiplicative algebra form AW4-E$_{\mathrm F}$
(Proposition~\ref{prop:rosati-EF}) and converts into the additive
relation of the arithmetic axioms (Proposition~\ref{prop:aw4-EF}). A
polarization in the vector-space form on $H^1$ is obtained only
afterwards, from semisimplicity and purity, and is not canonical. The
passage from $V_\ell J$ to $H^1$ goes through the identification of
$H^1(C_{\overline\F_q},\Q_\ell)$ with the dual of
$\Q_\ell\otimes T_\ell J$ \cite[Example~3.8, p.~40]{Milne2016RHFF};
its compatibility with Frobenius is not re-derived here. The proof
of~\ref{item:n7-exist} uses the imported semisimplicity and purity
statements; the proof of Theorem~\ref{thm:B}\ref{item:B-curves} uses
Proposition~\ref{prop:rosati-EF} and this identification instead. The
arithmetic axioms carry both the vector-space and the algebra forms of
positivity (Definition~\ref{def:aw4-forms}).

\begin{remark}\label{rem:ff-status}
Let $X$ be a smooth projective variety over $\F_q$, $\ell\nmid q$ a
prime, and $\iota:\overline{\Q}_\ell\to\C$ a field isomorphism. By
Theorem~\ref{thm:imp-deligne}, the eigenvalues of geometric Frobenius on
$H^j(X_{\overline{\F}_q},\Q_\ell)\otimes_{\Q_\ell,\iota}\C$ all have
absolute value $q^{j/2}$; Deligne's proof does not produce a positive
Frobenius polarization. By Proposition~\ref{prop:weil-converse}, such a
polarization exists if and only if Frobenius also acts semisimply there,
and semisimplicity of Frobenius on $H^j$ is known for curves and abelian
varieties (Weil) and a few other varieties, is implied by the standard
conjectures, and is unknown in general
\cite[p.~10]{Kleiman1994Standard},
\cite[Aside~3.18]{Milne2016RHFF}. For curves,
Weil's positivity of the Rosati involution gives semisimplicity and
purity of $\pi$ on $V_\ell J$ (Proposition~\ref{prop:weil-curves}); via
the identification of $H^1$ with the dual of the Tate module
(\cite[Example~3.8, p.~40]{Milne2016RHFF}, Frobenius compatibility not
re-derived here) this yields semisimplicity and purity of Frobenius on
$H^1$, and Proposition~\ref{prop:weil-converse} then supplies the
non-canonical polarizations of Proposition~\ref{prop:ff-N7}.
\end{remark}

\subsection{Logarithmic suspension}\label{subsec:suspension}

\begin{definition}\label{def:suspension}
Let $q>1$ be real and $\alpha\in\C^\times$. The \emph{logarithmic
suspension} of $\alpha$ is the set
$\Sigma_q(\alpha)=\{\rho\in\C: q^\rho=\alpha\}$, where
$q^\rho=e^{\rho\log q}$ with $\log q\in\R$.
\end{definition}

Explicitly, if $\ell$ is any complex logarithm of $\alpha$, then
$\Sigma_q(\alpha)=\{(\ell+2\pi in)/\log q:n\in\Z\}$; changing $\ell$
reindexes this set. The set $\Sigma_q(\alpha)$ is the spectrum of
Deninger's construction
\cite[Corollary~(2.6), p.~8, and (2.8), p.~10]{Deninger1994Motivic};
see also \cite[Remark~7.2]{Hesselholt2016THHZeta}.

\begin{definition}\label{def:log-susp}
Let $q>1$ be real, $V$ a finite-dimensional complex vector space, and
$\mathrm{F}:V\to V$ invertible, with generalized eigenspaces
$V_{(\alpha)}$. On $V_{(\alpha)}$ let
$N_\alpha=\log(\alpha^{-1}\mathrm{F})=\sum_{k\ge1}(-1)^{k+1}
(\alpha^{-1}\mathrm{F}-1)^k/k$, a nilpotent map (the sum is finite).
The \emph{logarithmic suspension} of $(V,\mathrm{F})$ is the pair
$(\widehat V,\Theta)$ with
\[
\widehat V=\bigoplus_{\rho\in\C,\ q^\rho\in\spec\mathrm{F}}
\widehat V_\rho,\qquad\widehat V_\rho=V_{(q^\rho)},\qquad
\Theta|_{\widehat V_\rho}=\rho+\tfrac1{\log q}N_{q^\rho}
\]
(algebraic direct sum, one copy of $V_{(q^\rho)}$ for each $\rho$).
\begin{enumerate}[label=(\alph*)]
\item If $h$ is a Hermitian form on $V$, the \emph{suspended form}
  $\widehat h$ is the orthogonal sum of the restrictions of $h$ to the
  summands $\widehat V_\rho$.
\item If $\widehat h$ is a Hermitian form on $\widehat V$ and
  $(\rho_\alpha)$ is a choice of one $\rho_\alpha\in\Sigma_q(\alpha)$
  for each eigenvalue $\alpha$ of $\mathrm{F}$, the \emph{restricted
  form} on $V=\bigoplus_\alpha V_{(\alpha)}$ is the orthogonal sum of
  the restrictions of $\widehat h$ to the summands
  $\widehat V_{\rho_\alpha}=V_{(\alpha)}$.
\end{enumerate}
\end{definition}

The index set is the union of the sets $\Sigma_q(\alpha)$ over the
eigenvalues $\alpha$ of $\mathrm{F}$; no branch of the logarithm is
chosen. On each summand $q^\Theta=q^\rho\exp(N_{q^\rho})=\mathrm{F}$.
Definition~\ref{def:log-susp} is Deninger's functor from regular
singular connections on $\mathbb{G}_m$ with monodromy $\mathrm{F}$
\cite[\S2, pp.~6--8]{Deninger1994Motivic}, applied to $\ell$-adic
cohomology as in \cite[\S4, p.~16]{Deninger1994Motivic}, written out
explicitly; the name is the one used in these notes.

\begin{proposition}[{(i)--(ii) Deninger
\cite[Corollary~(2.8), p.~10, and~(4.1), p.~16]{Deninger1994Motivic}}]%
\label{prop:suspension}
Let $q>1$ be real, $V$ a finite-dimensional complex vector space,
$\mathrm{F}:V\to V$ invertible, and $(\widehat V,\Theta)$ its
logarithmic suspension.
\begin{enumerate}[label=(\roman*)]
\item\label{item:susp-zeros} for each eigenvalue $\alpha$, the zeros of
  the entire function $s\mapsto1-\alpha q^{-s}$ are exactly the points
  of $\Sigma_q(\alpha)$, each simple; consequently the zeros of
  $\det(1-\mathrm{F}q^{-s})$, counted with multiplicity, are exactly the
  eigenvalues of $\Theta$ on $\widehat V$, counted with algebraic
  multiplicity;
\item\label{item:susp-core} $(\widehat V,\Theta)$ is a spectral core
  (Definition~\ref{def:spectral-core}) with $\widehat
  V_{\Theta\sim\rho}=\widehat V_\rho$;
\item\label{item:susp-pol} if $h$ is a positive Frobenius polarization
  of weight $j\in\R$ on $(V,\mathrm{F})$, then $\widehat h$ is a positive
  infinitesimal polarization of weight $j$ on $(\widehat V,\Theta)$
  (Definition~\ref{def:inf-pol});
\item\label{item:susp-line} in that case every eigenvalue of $\Theta$
  has real part $j/2$ and $\Theta$ is diagonalizable.
\end{enumerate}
\end{proposition}

\begin{proof}[Proof of \ref{item:susp-zeros}]
$1-\alpha q^{-s}=0$ if and only if $q^s=\alpha$, that is,
$s\in\Sigma_q(\alpha)$. The derivative $\alpha(\log q)q^{-s}$ never
vanishes, so each zero is simple. Since
$\det(1-\mathrm{F}q^{-s})=\prod_\alpha(1-\alpha q^{-s})^{\dim
V_{(\alpha)}}$, its zero at $\rho$ has order $\dim V_{(q^\rho)}$, which
is the algebraic multiplicity of $\rho$ as an eigenvalue of $\Theta$,
because $\Theta-\rho$ is nilpotent on $\widehat V_\rho$ and $\Theta$ has
no other eigenvalue $\rho$ on the other summands.
\end{proof}

\begin{proof}[Proof of \ref{item:susp-core}]
On $\widehat V_\rho$, $\Theta-\rho=N_{q^\rho}/\log q$ is nilpotent, so
$\widehat V_\rho$ is the generalized $\rho$-eigenspace, of finite
dimension. The set of eigenvalues is a finite union of arithmetic
progressions $\Sigma_q(\alpha)$ with nonzero step $2\pi i/\log q$, so it
has no finite accumulation point.
\end{proof}

\begin{proof}[Proof of \ref{item:susp-line} and \ref{item:susp-pol}]
By Proposition~\ref{prop:weil-linear}, $\mathrm{F}$ is diagonalizable
with $|\alpha|=q^{j/2}$, so every $N_\alpha=0$, $\Theta=\rho$ on
$\widehat V_\rho$, and $q^{\Re\rho}=|q^\rho|=q^{j/2}$ gives
$\Re\rho=j/2$. For $u\in\widehat V_\rho$, $v\in\widehat V_{\rho'}$,
$\widehat h(\Theta u,v)+\widehat h(u,\Theta v)=(\rho+\bar\rho')
\widehat h(u,v)$, which is $0$ if $\rho\ne\rho'$ (distinct summands are
orthogonal) and $2\Re\rho\,\widehat h(u,v)=j\,\widehat h(u,v)$ if
$\rho=\rho'$.
\end{proof}

Thus a positive Frobenius polarization of weight $j$ suspends to a
positive infinitesimal polarization of weight $j$, the hypothesis of
Proposition~\ref{prop:purity}, on a spectral core whose spectrum is the
full periodic zero divisor of $\det(1-\mathrm{F}q^{-s})$. The converse
holds as well: the infinitesimal identity is the logarithm of
$\mathrm{F}^*\mathrm{F}=q^j$.

\begin{proposition}\label{prop:susp-converse}
Let $q>1$ and $j$ be real, $V$ a finite-dimensional complex vector
space, $\mathrm{F}:V\to V$ invertible, $(\widehat V,\Theta)$ its
logarithmic suspension, $\widehat h$ a positive infinitesimal
polarization of weight $j$ on $(\widehat V,\Theta)$, and $h$ the
restricted form of $\widehat h$ for some choice $(\rho_\alpha)$
(Definition~\ref{def:log-susp}). Then $h$ is a positive Frobenius
polarization of weight $j$ on $(V,\mathrm{F})$.
\end{proposition}

\begin{proof}
By Proposition~\ref{prop:purity}, $\Theta$ is diagonalizable with
eigenvalues of real part $j/2$; so every $N_\alpha=0$ and
$\mathrm{F}=\alpha=q^{\rho_\alpha}$ on $V_{(\alpha)}$, with
$|\alpha|=q^{j/2}$. Then $h(\mathrm{F}u,\mathrm{F}v)=\alpha\bar\beta
h(u,v)$ for $u\in V_{(\alpha)}$, $v\in V_{(\beta)}$, which is
$q^jh(u,v)$ when $\alpha=\beta$ and $0=q^j\cdot0$ otherwise.
\end{proof}

\subsection{The multiplicative algebra form}\label{subsec:mult-form}

It is natural to expect that the positivity Weil proves for curves is
literally an instance of an additive relation $\theta^\dagger=j-\theta$
for the generator $\theta$ of a flow, the form in which positivity
enters the arithmetic axioms (Section~\ref{sec:axioms}); in fact the
Rosati involution supplies the multiplicative relation
$\pi\pi^\dagger=q$ on the endomorphism algebra of the Jacobian, and the
additive relation is obtained from it by functional calculus inside the
same algebra, once the multiplicative relation has itself forced
purity. This subsection makes the passage exact.

\begin{definition}\label{def:aw4-EF}
Let $q>1$ and $j$ be real, $V$ a finite-dimensional complex vector
space, and $\mathrm{F}:V\to V$ a linear map.
\begin{enumerate}[label=(\alph*)]
\item $(V,\mathrm{F})$ satisfies the \emph{multiplicative algebra form
  AW4-E$_{\mathrm F}$ of weight $j$} if there are a subalgebra
  $A\subset\operatorname{End}_\C(V)$ containing $\mathrm{id}_V$ and
  $\mathrm{F}$ and a positive trace involution $(\tau,\dagger)$ on $A$
  over $k=\C$ (Definition~\ref{def:pos-involution}) with
  $\mathrm{F}\mathrm{F}^\dagger=q^j$.
\item If $\mathrm{F}$ is invertible, a \emph{sheet} of its logarithmic
  suspension $(\widehat V,\Theta)$ (Definition~\ref{def:log-susp}) is a
  map $\sigma$ from the set of eigenvalues of $\mathrm{F}$ to $\C$ with
  $\sigma(\alpha)\in\Sigma_q(\alpha)$ for every eigenvalue $\alpha$.
\item The \emph{sheet space} of a sheet $\sigma$ is
  $\widehat V_\sigma=\bigoplus_\alpha\widehat V_{\sigma(\alpha)}\subset
  \widehat V$, and its \emph{sheet map} $\iota_\sigma:V\to\widehat
  V_\sigma$ sends $v\in V_{(\alpha)}$ to its copy in
  $\widehat V_{\sigma(\alpha)}=V_{(\alpha)}$.
\item The \emph{interpolation polynomial} of a sheet $\sigma$ is the
  polynomial $p_\sigma\in\C[X]$ of degree less than the number of
  eigenvalues of $\mathrm{F}$ with $p_\sigma(\alpha)=\sigma(\alpha)$ for
  every eigenvalue $\alpha$.
\end{enumerate}
\end{definition}

Distinct eigenvalues $\alpha$ have distinct $\sigma(\alpha)$, because
$q^{\sigma(\alpha)}=\alpha$; so the sum in~(c) is direct, $\iota_\sigma$
is an isomorphism, and $p_\sigma$ exists and is unique by Lagrange
interpolation. For $m\in\Z$ the map $\sigma+2\pi im/\log q$ is again a
sheet, and $\widehat V$ is the direct sum of the sheet spaces of the
sheets $\sigma+2\pi im/\log q$, $m\in\Z$: every $\rho$ with
$q^\rho=\alpha$ equals $\sigma(\alpha)+2\pi im/\log q$ for exactly one
$m$. For a subalgebra $A\subset\operatorname{End}_\C(V)$ we write
$A^t=\{a^t:a\in A\}\subset\operatorname{End}_\C(V^*)$ for the algebra of
transposes acting on the dual space $V^*$.

\begin{proposition}\label{prop:aw4-EF}
Let $q>1$ and $j$ be real, $V$ a finite-dimensional complex vector
space, $\mathrm{F}:V\to V$ linear, and suppose that $(V,\mathrm{F})$
satisfies AW4-E$_{\mathrm F}$ of weight $j$
(Definition~\ref{def:aw4-EF}) with the algebra $A$ and the positive
trace involution $(\tau,\dagger)$. Let $(\widehat V,\Theta)$ be the
logarithmic suspension of $(V,\mathrm{F})$.
\begin{enumerate}[label=(\roman*)]
\item\label{item:EF-pure} $\mathrm{F}$ is invertible,
  $\mathrm{F}^{-1}=q^{-j}\mathrm{F}^\dagger$ lies in $A$, $\mathrm{F}$
  is diagonalizable, and every eigenvalue $\alpha$ of $\mathrm{F}$
  satisfies $|\alpha|=q^{j/2}$.
\item\label{item:EF-sheet} For every sheet $\sigma$, the element
  $\theta_\sigma=p_\sigma(\mathrm{F})$ of $A$ satisfies
  $\Theta\iota_\sigma=\iota_\sigma\theta_\sigma$ and
  $\theta_\sigma^\dagger=j-\theta_\sigma$.
\item\label{item:EF-add} Every vector of $\widehat V$ lies in a
  finite-dimensional $\Theta$-stable subspace $\widehat V'$ which carries
  a subalgebra $A'\subset\operatorname{End}_\C(\widehat V')$ containing
  $\mathrm{id}_{\widehat V'}$ and $\theta'=\Theta|_{\widehat V'}$, and a
  positive trace involution $(\tau',\dagger')$ on $A'$ over $\C$ with
  $\theta'^{\dagger'}=j-\theta'$.
\item\label{item:EF-dual} The dual $(V^*,\mathrm{F}^t)$ satisfies
  AW4-E$_{\mathrm F}$ of weight $j$ with the algebra $A^t$, the trace
  $a^t\mapsto\tau(a)$, and the involution $a^t\mapsto(a^\dagger)^t$.
\end{enumerate}
\end{proposition}

\begin{proof}[Proof of \ref{item:EF-pure}]
Since $\mathrm{F}\mathrm{F}^\dagger=q^j\,\mathrm{id}_V$ with $q^j\ne0$
and $V$ is finite-dimensional, $\mathrm{F}$ is invertible with inverse
$q^{-j}\mathrm{F}^\dagger\in A$. The inclusion
$A=A\otimes_\C\C\to\operatorname{End}_\C(V)$ is an injective unital
homomorphism, so Proposition~\ref{prop:algebra-form}\ref{item:af-rep}
with $k=\C$ and $g=\mathrm{F}$ shows that $\mathrm{F}$ is
diagonalizable with eigenvalues of modulus $q^{j/2}$.
\end{proof}

\begin{proof}[Proof of \ref{item:EF-sheet}]
By~\ref{item:EF-pure}, $V_{(\alpha)}=\ker(\mathrm{F}-\alpha)$, so the
nilpotent map $N_\alpha$ of Definition~\ref{def:log-susp} vanishes and
$\Theta=\sigma(\alpha)$ on $\widehat V_{\sigma(\alpha)}$, while
$p_\sigma(\mathrm{F})$ acts on $V_{(\alpha)}$ by
$p_\sigma(\alpha)=\sigma(\alpha)$; hence
$\Theta\iota_\sigma=\iota_\sigma\theta_\sigma$. The element
$\theta_\sigma$ lies in $A$ because $A$ contains $\mathrm{id}_V$ and
$\mathrm{F}$. The involution fixes $\mathrm{id}_V$: from
$(ab)^\dagger=b^\dagger a^\dagger$ with $a=\mathrm{id}_V$ and
$b=\mathrm{id}_V^\dagger$ one gets $\mathrm{id}_V=\mathrm{id}_V
\mathrm{id}_V^\dagger=\mathrm{id}_V^\dagger$, using
$b^\dagger=\mathrm{id}_V$. Since $\dagger$ is additive, antilinear, and
reverses products, $p(\mathrm{F})^\dagger=\bar p(\mathrm{F}^\dagger)$
for every $p\in\C[X]$, where $\bar p$ has the conjugate coefficients.
By~\ref{item:EF-pure}, $\mathrm{F}^\dagger=q^j\mathrm{F}^{-1}$ acts on
$V_{(\alpha)}$ by $q^j\alpha^{-1}=\bar\alpha$, because
$|\alpha|^2=q^j$. So $\theta_\sigma^\dagger$ acts on $V_{(\alpha)}$ by
$\bar p_\sigma(\bar\alpha)=\overline{\sigma(\alpha)}$. Every
$\rho\in\Sigma_q(\alpha)$ has $q^{\Re\rho}=|\alpha|=q^{j/2}$, so
$\Re\sigma(\alpha)=j/2$ and $\overline{\sigma(\alpha)}=j-\sigma(\alpha)$.
Hence $\theta_\sigma^\dagger$ and $j-\theta_\sigma$ agree on every
$V_{(\alpha)}$, hence on $V$; both lie in
$A\subset\operatorname{End}_\C(V)$, so they are equal.
\end{proof}

\begin{proof}[Proof of \ref{item:EF-add}]
Fix a sheet $\sigma_0$ and put $\sigma_m=\sigma_0+2\pi im/\log q$ for
$m\in\Z$. A vector of $\widehat V$ has finitely many nonzero components,
so it lies in $\widehat V'=\bigoplus_{|m|\le n}\widehat V_{\sigma_m}$ for
some $n$; this space is finite-dimensional and $\Theta$-stable. Let
$A'$ be the set of maps acting on each $\widehat V_{\sigma_m}$,
$|m|\le n$, as $\iota_{\sigma_m}a_m\iota_{\sigma_m}^{-1}$ for some
$a_m\in A$, a subalgebra of $\operatorname{End}_\C(\widehat V')$
isomorphic to the product of $2n+1$ copies of $A$; put
$\tau'=\sum_m\tau(a_m)$ and let $\dagger'$ act componentwise. Every
identity in Definition~\ref{def:pos-involution} holds componentwise,
and $\tau'(xx^{\dagger'})=\sum_m\tau(a_ma_m^\dagger)>0$ unless every
$a_m$ vanishes, so $(\tau',\dagger')$ is a positive trace involution.
By~\ref{item:EF-sheet}, $\theta'$ is the element of $A'$ with
components $\theta_{\sigma_m}$, and $\theta'^{\dagger'}$ has components
$j-\theta_{\sigma_m}$, that is, $\theta'^{\dagger'}=j-\theta'$.
\end{proof}

\begin{proof}[Proof of \ref{item:EF-dual}]
Transposition $a\mapsto a^t$ is a bijective linear map
$A\to A^t$ reversing products, and $(\mathrm{F}^t)^{-1}=
(\mathrm{F}^{-1})^t$. Transporting $\tau$ and $\dagger$ along it
preserves every identity in Definition~\ref{def:pos-involution}: for
instance $(a^tb^t)^\dagger=((ba)^\dagger)^t=(a^\dagger b^\dagger)^t=
(b^\dagger)^t(a^\dagger)^t$, and
$\tau(a^t(a^\dagger)^t)=\tau(a^\dagger a)=\tau(aa^\dagger)>0$ for
$a\ne0$. Finally $\mathrm{F}^\dagger\mathrm{F}=q^j$ because
$\mathrm{F}^\dagger=q^j\mathrm{F}^{-1}$, so
$\mathrm{F}^t(\mathrm{F}^t)^\dagger=(\mathrm{F}^\dagger\mathrm{F})^t=
q^j$.
\end{proof}

Item~\ref{item:EF-add} is the condition AW4-E of
Definition~\ref{def:aw4-forms} in Section~\ref{sec:axioms}, read on
$(\widehat V,\Theta)$ with weight $j$. The multiplicative form is
canonical, while the additive form depends on the choice of sheets,
exactly as the flow depends on the logarithm of Frobenius
(Proposition~\ref{prop:one-branch} below); the passage between them
uses the purity supplied by the multiplicative form itself
(item~\ref{item:EF-pure}).

\begin{proposition}\label{prop:rosati-EF}
Let $C$, $g\ge1$, $J$, $\pi$, $\ell$, and the Rosati involution
$a\mapsto a^\dagger$ on $\operatorname{End}^0(J)$ be as in
Theorem~\ref{thm:imp-rosati}, let $\iota:\overline\Q_\ell\to\C$ be a
field isomorphism, put $V=V_\ell J\otimes_{\Q_\ell,\iota}\C$, let
$\mathrm{F}:V\to V$ be the map induced by $\pi$, and let $A\subset
\operatorname{End}_\C(V)$ be the image of the homomorphism
$\operatorname{End}^0(J)\otimes_\Q\C\to\operatorname{End}_\C(V)$
induced by the action on $V_\ell J$.
\begin{enumerate}[label=(\roman*)]
\item\label{item:rEF-inj} This homomorphism is injective.
\item\label{item:rEF-inv} There is a unique antilinear map $\dagger$
  on $A$ sending the image of $a\otimes z$ to the image of
  $a^\dagger\otimes\bar z$ ($a\in\operatorname{End}^0(J)$, $z\in\C$),
  and with $\tau$ the trace of endomorphisms of $V$, the pair
  $(\tau,\dagger)$ is a positive trace involution on $A$ over $\C$.
\item\label{item:rEF-EF} $(V,\mathrm{F})$ satisfies AW4-E$_{\mathrm F}$
  of weight one with this algebra and involution.
\end{enumerate}
\end{proposition}

\begin{proof}[Proof of \ref{item:rEF-inj}]
By Theorem~\ref{thm:imp-rosati}\ref{item:ro-char}, the map
$\operatorname{End}^0(J)\otimes_\Q\Q_\ell\to\operatorname{End}_{\Q_\ell}
(V_\ell J)$ is injective. Tensoring over $\Q_\ell$ with $\C$ (through
$\iota$) preserves injectivity, and
$\operatorname{End}_{\Q_\ell}(V_\ell J)\otimes_{\Q_\ell}\C=
\operatorname{End}_\C(V)$.
\end{proof}

\begin{proof}[Proof of \ref{item:rEF-inv}]
The map $(a,z)\mapsto a^\dagger\otimes\bar z$ is $\Q$-bilinear, so it
induces an additive, antilinear map on
$\operatorname{End}^0(J)\otimes_\Q\C$, which by~\ref{item:rEF-inj}
defines $\dagger$ on $A$; it reverses products and is an involution
because the Rosati involution does, as one checks on pure tensors. The
trace $\tau$ satisfies $\tau(xy)=\tau(yx)$. For
$a\in\operatorname{End}^0(J)$ let $\Tr(a)$ be its trace on $V_\ell J$;
it is rational by Theorem~\ref{thm:imp-rosati}\ref{item:ro-char}, and
$\tau(a\otimes z)=z\Tr(a)$. The Rosati involution is the adjoint for
the nondegenerate alternating Weil pairing $e^\lambda$ on $V_\ell J$,
$e^\lambda(aa_1,a_2)=e^\lambda(a_1,a^\dagger a_2)$
\cite[\S14, p.~61]{MilneAV2008}, so $a^\dagger$ has the characteristic
polynomial of $a$ and $\Tr(a^\dagger)=\Tr(a)$; hence
$\tau(x^\dagger)=\overline{\tau(x)}$ on pure tensors and, by
additivity, on $A$. For positivity, let $(a_i)_{1\le i\le n}$ be a
$\Q$-basis of $\operatorname{End}^0(J)$ and $G_{ik}=\Tr(a_ia_k^\dagger)
\in\Q$. The matrix $G$ is symmetric, since
$\Tr(a_ka_i^\dagger)=\Tr((a_ia_k^\dagger)^\dagger)$, and positive
definite on $\Q^n$ by Theorem~\ref{thm:imp-rosati}\ref{item:ro-pos};
as in the proof of
Proposition~\ref{prop:algebra-form}\ref{item:af-alg}, it is then
positive definite on $\R^n$. Every $x\in A$ is the image of
$\sum_ia_i\otimes z_i$ for a unique $z\in\C^n$, and
$\tau(xx^\dagger)=\sum_{i,k}z_i\bar z_kG_{ik}=u^TGu+v^TGv$ for
$z=u+iv$ with $u,v\in\R^n$, which is positive for $z\ne0$.
\end{proof}

\begin{proof}[Proof of \ref{item:rEF-EF}]
$A$ contains the identity and $\mathrm{F}$, the image of
$\pi\otimes1$, and $\mathrm{F}\mathrm{F}^\dagger$ is the image of
$\pi\pi^\dagger\otimes1=q\otimes1$ by
Theorem~\ref{thm:imp-rosati}\ref{item:ro-pi}.
\end{proof}

Thus Weil's positivity is an instance of the multiplicative algebra form,
and by Proposition~\ref{prop:aw4-EF} it gives, inside the same algebra,
the semisimplicity and purity of $\pi$ on $V_\ell J\otimes_\iota\C$
(the complex form of Proposition~\ref{prop:weil-curves}) and the
additive algebra form on the logarithmic suspension of
$(V_\ell J\otimes_\iota\C,\pi)$. The passage to $H^1$ is the one import
that remains. It consists of three facts, none re-derived here:
$H_1(C_{\mathrm{et}},\Q_\ell)=\Q_\ell\otimes T_\ell J$
\cite[Example~3.8, p.~40]{Milne2016RHFF}; the canonical isomorphism
$H^1(J_{\overline\F_q},\Z_\ell)\cong\operatorname{Hom}(T_\ell J,\Z_\ell)$
and its compatibility with the Galois action
\cite[Chapter~I, Theorem~12.1(a), p.~55, and Remark~12.5,
p.~56]{MilneAV2008}; and the isomorphism $H^1(J_{\overline\F_q},\Z_\ell)
\to H^1(C_{\overline\F_q},\Z_\ell)$ induced by an Abel--Jacobi map
\cite[Chapter~III, Corollary~9.6, p.~114]{MilneAV2008}, which does not
depend on the base point, because translations act trivially on
$H^1(J)$ (the addition map pulls $x\in H^1$ back to $x\otimes1+
1\otimes x$ \cite[Chapter~I, (3), p.~55]{MilneAV2008}), and hence
commutes with the pullbacks by the $q$-power
Frobenius endomorphisms of $C$ and $J$, the latter being $\pi$. The
naturality of these isomorphisms, which turns pullback by $\pi$ into
the transpose of the action of $\pi$ on $T_\ell J$, is also taken from
the source. So $H^1(C_{\overline\F_q},\Q_\ell)$ is the dual of
$V_\ell J$ with Frobenius the transpose of $\pi$; in genus $0$ both are
zero. Granting this,
Proposition~\ref{prop:aw4-EF}\ref{item:EF-dual} carries the
multiplicative form from $V_\ell J$ to $H^1$ (proof of
Theorem~\ref{thm:B}\ref{item:B-curves}).


\section{Analytic and positivity conditions on Deninger structures}
\label{sec:axioms}

We formulate axioms, following Deninger's conjectural determinant
formula \cite[Conjecture~2, p.~3]{Deninger2010HilbertPolya}, his further
expectations of a trace, a cup-product pairing, and a positive star
\cite[pp.~4--5]{Deninger2010HilbertPolya}, and the reformulation by
Flach and Morin \cite[Definition~2.1, Conjecture~2.3]%
{FlachMorin2020DeningerWeilArakelov}, in a form
sufficient for an exact implication to the Riemann hypothesis. This
section describes the stronger cup-product, duality, and star package
over the uniform spectral formalism of
Section~\ref{subsec:uniform-formalism}. Its global arithmetic
construction with compatible positivity is a target; the suspended
finite-field structures are constructed in Theorem~\ref{thm:B}.
The values of a Deninger cohomology theory
(Definition~\ref{def:deninger-theory}) are Deninger structures
(Definition~\ref{def:deninger-structure}); the conditions below are
imposed on a single such structure, polarized or not, and on its
Weil datum.


For a number field $K$ we write $\Xbar_K=(\Spec\cO_K,X_{K,\infty})$ for
the Arakelov compactification of $\Spec\cO_K$: an arithmetic pair, not
a scheme. Its dimension is the Krull dimension of $\cO_K$, namely one.

The \emph{arithmetic test category} is the full subcategory
$\Schbar_1$ of $\Schbar^{\mathrm{proj}}$
(Definition~\ref{def:compactified-schemes}\ref{item:schbar-one}), whose
objects are the finite disjoint unions of the $\Xbar_K$.

Compare the Arakelov compactifications of
\cite[Notation~2.2]{FlachMorin2020DeningerWeilArakelov} and
\cite[\S2]{Deninger2022NoWeil}.

\begin{proposition}[Number fields and arithmetic curves]
\label{prop:A1-number-fields}
Let $\mathcal N$ be the category of number fields with field
homomorphisms as morphisms. The assignment
$K\mapsto\Xbar_K$, $(\alpha:K\to L)\mapsto\Spec(\alpha|_{\cO_K})$,
is an equivalence from $\mathcal N^{\mathrm{op}}$ onto the full
subcategory of connected objects of $\Schbar_1$.
\end{proposition}

\begin{proof}
A ring homomorphism $\cO_K\to\cO_L$ extends uniquely to $K\to L$, and
a field homomorphism $K\to L$ carries integral elements to integral
elements, hence $\cO_K$ into $\cO_L$. Thus
$\operatorname{Hom}_{\mathcal N}(K,L)=\operatorname{Hom}(\cO_K,\cO_L)
=\operatorname{Hom}(\Spec\cO_L,\Spec\cO_K)$. The map on infinite
places is induced by the scheme morphism, so it adds no choice. Every
connected object of $\Schbar_1$ is some $\Xbar_K$.
\end{proof}

In particular the morphisms $\Xbar_L\to\Xbar_K$ between connected
objects are exactly those induced by field embeddings
$K\hookrightarrow L$, with the corresponding finite map
$\Spec\cO_L\to\Spec\cO_K$ and map of infinite places, and the
automorphisms of $\Xbar_K$ form $\operatorname{Aut}(K)$.

A contravariant theory on $\Schbar_1$ already contains $\Xbar_\Q$, and
all
number fields are needed for the Galois-functoriality test of
Remark~\ref{rem:R5-import}. For the implication to RH alone, the single
object $\Xbar_\Q$ suffices. Larger categories appear in
\cite[Conjecture~2, p.~3]{Deninger2010HilbertPolya} (separated
$\Z$-schemes of finite type) and in
\cite[Conjecture~2.3, p.~522]{FlachMorin2020DeningerWeilArakelov}
(arithmetic schemes, with cohomology with compact supports, together
with Arakelov compactifications of projective flat regular schemes);
they extend the minimal number-field package and motivate the uniform
source category of Definition~\ref{def:projective-inputs}.


It is natural to expect such a larger theory to carry a K\"unneth
formula, as Weil cohomologies over a field do; in fact its natural
strong forms are inconsistent with the axioms below. No K\"unneth
formula is imposed in the axioms, and the test category $\Schbar_1$ is not
presented as a scheme category over a ``field with one element''.
Section~\ref{sec:kunneth} states the K\"unneth conjecture in the form
that survives, records its refutations (no formula over $\Spec\Z$, even
in derived form; no spectral core, no compact resolvent, and no
determinant on the tensor square of an arithmetic curve), and proves the
Hodge index statement on the arithmetic surface that corresponds to the
positivity in Weil's proof.

Spectral cores are defined in the introduction
(Definition~\ref{def:spectral-core});
 by
Proposition~\ref{prop:suspension}\ref{item:susp-core}, a logarithmic
suspension is one.

\begin{definition}\label{def:regdet}
Let $(V,\Theta)$ be a spectral core with multiplicities $m$.
\begin{enumerate}[label=(\alph*)]
\item A \emph{spectral cut} on a nonempty open set $U\subset\C$ is a
  choice, for each $\rho\in\spec\Theta$, of a holomorphic branch
  $L_\rho$ of $\log((s-\rho)/2\pi)$ on $U$ (so $U\cap\spec\Theta=
  \emptyset$). For $s\in U$ and $w\in\C$ put
  $((s-\rho)/2\pi)^{-w}=e^{-wL_\rho(s)}$.
\item The \emph{spectral zeta function} for this cut is
  \[
  Z(w;s)=\sum_{\rho\in\spec\Theta}m(\rho)
    \Bigl(\frac{s-\rho}{2\pi}\Bigr)^{-w}\qquad(s\in U),
  \]
  wherever the series converges absolutely.
\item $(V,\Theta)$ is \emph{zeta-regularizable} for the cut if, for
  every $s\in U$, the series converges for $\Re w$ large; there is
  $\varepsilon=\varepsilon(s)>0$ such that $Z(w;s)$ continues to a
  meromorphic function of $w$ on the half-plane $\Re w>-\varepsilon$
  that is regular at $w=0$; the function $s\mapsto\partial_wZ(0;s)$ is
  holomorphic on $U$; and $s\mapsto\exp(-\partial_wZ(0;s))$ continues to
  an entire function of $s$.
\item The \emph{regularized determinant} for the cut is that entire
  function, $\detreg\bigl((s-\Theta)/2\pi\bigr)=
  \exp(-\partial_wZ(0;s))$.
\end{enumerate}
\end{definition}

Two half-planes $\Re w>-\varepsilon$ and $\Re w>-\varepsilon'$ are
nested, and each contains $w=0$ and the half-plane of absolute
convergence. By the identity theorem the continuation in (c) is
therefore unique, so $\partial_wZ(0;s)$ and $\detreg$ depend only on
the spectral cut.

Whenever a regularized determinant is used below, a spectral cut is
fixed and recorded as part of the data. If $V$ is
finite-dimensional, then for every spectral cut
$\partial_wZ(0;s)=-\sum_\rho m(\rho)L_\rho(s)$, so
$\detreg((s-\Theta)/2\pi)=\prod_\rho((s-\rho)/2\pi)^{m(\rho)}$
independently of the cut. In particular, if $V=\C$ and $\Theta=\rho_0$,
then $\detreg((s-\Theta)/2\pi)=(s-\rho_0)/2\pi$. For infinite spectra
the determinant depends on the cut; the one cut used below is Deninger's
\cite[(1.1), p.~3]{Deninger1994Motivic}.

\begin{definition}\label{def:principal-cut}
Let $(V,\Theta)$ be a spectral core whose spectrum lies in a vertical
strip $\{a\le\Re s\le b\}$, and let $\operatorname{Log}$ be the
principal branch of the logarithm on $\C\setminus(-\infty,0]$.
\begin{enumerate}[label=(\alph*)]
\item The \emph{principal cut} of $(V,\Theta)$ is the spectral cut on
  $U_{\mathrm{pr}}=\C\setminus\bigcup_{\rho\in\spec\Theta}
  \bigl(\rho+(-\infty,0]\bigr)$ given by
  $L_\rho(s)=\operatorname{Log}((s-\rho)/2\pi)$.
\item The \emph{reflected cut} of $(V,\Theta)$ is the spectral cut on
  $U_{\mathrm{rf}}=\C\setminus\bigcup_{\rho\in\spec\Theta}
  \bigl(\rho+[0,\infty)\bigr)$ given by the branch $L_\rho(s)$ of
  $\log((s-\rho)/2\pi)$ with imaginary part in $[0,2\pi)$.
\end{enumerate}
\end{definition}

The sets $U_{\mathrm{pr}}$ and $U_{\mathrm{rf}}$ are open and contain
the half-planes $\Re s>b$ and $\Re s<a$ respectively: only finitely many
eigenvalues have imaginary part in a given bounded interval, so each
union of rays is locally finite, hence closed. The principal cut is
Deninger's convention $-\pi<\arg\le\pi$, restricted to
$U_{\mathrm{pr}}$; there it agrees with his pointwise definition.
The reflected cut computes Deninger's $\detreg(-\Theta)$
\cite[(1.3)--(1.7), pp.~4--6]{Deninger1994Motivic}; the regularized
determinant with the principal cut is also
\cite[Definition~2.1, p.~521]{FlachMorin2020DeningerWeilArakelov}.

The internal unital axioms apply in weight zero. For compactified
number rings $\Xbar_K$, $d=1$; the finite-field calibration uses the
axioms with $d=\dim X$, $w=0$, and a periodic unit object; arithmetic
schemes use them with $d\ge2$. Pure coefficients of weight $w$ use
graded spectral realizations of type $(1,w)$, as in
Remark~\ref{rem:A-motivic}; their products have separate targets.

Definition~\ref{def:unit} is stated in the introduction.


The $q$-periodic unit is Deninger's space $R_p$ of finite Fourier series
on $\R/(\log p)\Z$ with $\Theta=d/dy$, complexified, for $q=p$
\cite[pp.~167--168 and Proposition~3.1]{Deninger1998Analogies}.

Definition~\ref{def:aw} and the core axioms
\ref{ax:AW1}--\ref{ax:AW4} of
Definition~\ref{def:weil-cohomology} are stated in the introduction.

\begin{definition}[Additional analytic conditions]
\label{def:aw-analytic}
Let a Weil datum of type $(d,0)$ be as in Definition~\ref{def:aw}.
It satisfies AW5 or AW6 when the corresponding condition below holds.
\begin{enumerate}[label=\textbf{AW\arabic*},ref=AW\arabic*,
  leftmargin=*,start=5]
\item\label{ax:AW5} \emph{(Hilbert realization.)} Each $h_j$ is
  positive definite; let $\cH^j$ be the completion of $(V^j,h_j)$. The
  normalized flow
  $U_j(t)=e^{-jt/2}\phi^t$ extends to a strongly continuous unitary
  group on $\cH^j$ whose generator $A_j$ is skew-adjoint and has $V^j$
  as a core; put $T_j=-iA_j$, so that $T_j$ is self-adjoint and the
  closure of $\Theta_j$ is $j/2+iT_j$.
\item\label{ax:AW6} \emph{(No loss on completion; schematic.)} $T_j$
  has pure point spectrum equal to
  $\{(\rho-j/2)/i:\rho\in\spec\Theta_j\}$, with the same finite
  multiplicities and no finite accumulation; its counting function
  satisfies $N_j(R)=O(R^M)$ for some $M$; and for test functions $f$ in
  a class $\mathcal S_X$ the trace $\tr f(T_j)$ equals the
  explicit-formula distribution of the object, with an archimedean
  regularization.
\end{enumerate}
\end{definition}


The endpoint clause of \ref{ax:AW1} is Deninger's normalization of
$H^0$ and $H^2$ \cite[(3), p.~168]{Deninger1998Analogies}; the duality
of \ref{ax:AW3} and the star law of \ref{ax:AW4} are
\cite[pp.~168--169 and p.~172]{Deninger1998Analogies},
\cite[pp.~4--5]{Deninger2010HilbertPolya}, and
\cite[Conjecture~2.3(iii)--(iv), p.~522]%
{FlachMorin2020DeningerWeilArakelov}; the Hilbert completion of
\ref{ax:AW5} is \cite[p.~33]{Deninger1994Motivic}.
Positive-weight coefficients use separate spectral realizations, not
these internal AW axioms. Condition \ref{ax:AW6} is schematic as
stated. For $\Xbar_K$ its trace-formula clause is made precise by the
axiom AW6$'$ of Definition~\ref{def:aw6prime} below, and its first
clause follows from \ref{ax:AW5}
(Proposition~\ref{prop:aw5-contains}); Theorems~\ref{thm:rh}
and~\ref{thm:A} do not use \ref{ax:AW5} or \ref{ax:AW6}. The divisor
identity \ref{ax:AW2} is
imposed. Over $\F_q$ rationality is derived from the Lefschetz trace
formula; over $\Xbar_K$ the trace-formula axiom AW6$'$ is, given
\ref{ax:AW1}, equivalent to \ref{ax:AW2}
(Theorem~\ref{thm:aw6-div}), so replacing \ref{ax:AW2} by
AW6$'$ changes the form of the assumption, not its strength.
 The two further forms of the positivity axiom, and the
determinant form of \ref{ax:AW2}, are separate definitions.

\begin{definition}\label{def:aw4-forms}
Let $(V^\bullet,\Theta,\cup,\Tr,*)$ be a Weil datum of type $(d,0)$ and
$0\le j\le2d$.
\begin{enumerate}[label=(\alph*)]
\item The datum satisfies \emph{AW4-V in degree $j$} if
  $(V^j,\Theta_j)$ admits a positive infinitesimal polarization of
  weight $j$ (Definition~\ref{def:inf-pol}).
\item The datum satisfies \emph{AW4-E in degree $j$} if every vector of
  $V^j$ lies in a finite-dimensional $\Theta_j$-stable subspace $V'$
  which carries a subalgebra $A_{V'}\subset\operatorname{End}_\C(V')$
  containing $\mathrm{id}_{V'}$ and $\theta_{V'}=\Theta_j|_{V'}$, and a
  positive trace involution $(\tau_{V'},\dagger)$ on $A_{V'}$ over $k=\C$
  (Definition~\ref{def:pos-involution}) with
  $\theta_{V'}^\dagger=j-\theta_{V'}$.
\end{enumerate}
\end{definition}

\begin{definition}\label{def:aw-det}
Let $(V^j,\Theta_j)_{0\le j\le2d}$ be graded spectral cores, and let
$Z\not\equiv0$ be meromorphic on $\C$. They satisfy \emph{AW-det for
$Z$} if, for each $j$, a spectral cut for $(V^j,\Theta_j)$ is
fixed as part of the data, $(V^j,\Theta_j)$ is zeta-regularizable for
it (Definition~\ref{def:regdet}), the entire function
$D_j(s)=\detreg((s-\Theta_j)/2\pi)$ satisfies
$\divisor D_j=\operatorname{Sp}_j$, and there are recorded constants
$a,b\in\C$, the \emph{normalization term} $e^{a+bs}$, with
\begin{equation}\label{eq:aw-det}
e^{a+bs}\,Z(s)=\prod_{j=0}^{2d}D_j(s)^{(-1)^{j+1}}.
\tag{AW-det}
\end{equation}
\end{definition}

Taking divisors in \eqref{eq:aw-det} gives the spectral divisor
identity of Definition~\ref{def:graded-realization}. For an
unweighted Weil datum this is \ref{ax:AW2}. The converse fails:
the divisor identity asks nothing of regularization. All
zero-location theorems below use \ref{ax:AW2}, or AW6$'$, which is
equivalent to it given \ref{ax:AW1} (Theorem~\ref{thm:aw6-div}), and
never AW-det, so no value of the
normalization term enters them; AW-det records the stronger,
determinant-level expectation of Deninger
\cite[Conjecture~2, p.~3]{Deninger2010HilbertPolya} and Flach--Morin
\cite[Conjecture~2.3]{FlachMorin2020DeningerWeilArakelov}. The
normalization term for $\Xbar_K$ is computed in
Theorem~\ref{thm:det-K}.

For a number field $K$ the endpoint normalization is
$V^0=\C\cdot1$ with $\Theta_0=0$ and $V^2\cong\C(-1)$.

A Weil datum on $\Xbar_K$ satisfying \ref{ax:AW1}--\ref{ax:AW4} for
$(\Lambda_K,\mathbb 1)$ is a polarized Deninger structure in
$\Den^{\mathrm{pol}}_1(\mathbb 1)$, with a chosen polarization,
satisfying \ref{ax:AW2} for $\Lambda_K$. A Deninger cohomology theory
(Definition~\ref{def:deninger-theory}) restricts to a functor on
$\Schbar_1$ whose value on each $\Xbar_K$ is such a structure
without the polarization.
The trace-formula clause of \ref{ax:AW6} is made precise for $\Xbar_K$
by the axiom AW6$'$ below. It tests the flow against smooth compactly
supported functions of the time $t$ and compares the result with the
geometric side of the explicit formula.

\begin{definition}\label{def:test-classes}
Let $(V,\Theta)$ be a spectral core with multiplicity function $m$.
\begin{enumerate}[label=(\alph*)]
\item The \emph{transform} of $\varphi\in C_c^\infty(\R)$ is the entire
  function $\Phi(s)=\int_\R\varphi(t)e^{st}\,dt$.
\item The \emph{two-sided class} is $\mathcal S=C_c^\infty(\R)$, the
  \emph{one-sided class} is $\mathcal S^+=C_c^\infty((0,\infty))$, for
  $T_0\in\R$ the \emph{tail class} is
  $\mathcal S_{>T_0}=C_c^\infty((T_0,\infty))$, and for $L>0$ the
  \emph{window classes} are $\mathcal S_{(0,L)}=C_c^\infty((0,L))$ and
  $\mathcal S_{(-L,L)}=C_c^\infty((-L,L))$.
\item $(V,\Theta)$ satisfies the \emph{growth condition} (G) if
  $\spec\Theta$ lies in a vertical strip $\{a\le\Re s\le b\}$ and
  $\sum_{|\rho|\le R}m(\rho)=O(R^M)$ as $R\to\infty$ for some $M$.
\item For $\varphi\in\mathcal S$ such that the series converges
  absolutely, the \emph{distributional trace} of the flow of
  $(V,\Theta)$ at $\varphi$ is
  $\Tr_{\mathrm{dis}}(V)(\varphi)=\sum_\rho m(\rho)\Phi(\rho)$.
\end{enumerate}
\end{definition}

Under $x=e^t$, the function $x\mapsto\varphi(\log x)$ lies in
$C_c^\infty(\R_{>0})$ and $\Phi$ is its Mellin transform
$\int_0^\infty\varphi(\log x)x^{s-1}\,dx$; the distributional trace is
Deninger's \cite[(4) and the display after it, p.~170]
{Deninger1998Analogies}.
The source writes the sum over $\rho$ without an explicit weight; the
weight $m(\rho)$ is the standard reading, and it is the one that the
proof of Theorem~\ref{thm:imp-explicit} establishes.

\begin{definition}\label{def:geom-side}
Let $K$ be a number field with $r_1$ real and $r_2$ complex places and
discriminant $D_K$, and let $\varphi\in\mathcal S$.
\begin{enumerate}[label=(\alph*)]
\item The \emph{archimedean terms} at $\varphi$ are
  \[
  \begin{aligned}
  W_\R(\varphi)&=\log\pi\cdot\varphi(0)+\int_0^\infty\Bigl[
  \frac{\varphi(t)+e^{-t}\varphi(-t)}{1-e^{-2t}}
  -\frac{e^{-2t}\varphi(0)}t\Bigr]dt,\\
  W_\C(\varphi)&=2\log(2\pi)\cdot\varphi(0)+\int_0^\infty\Bigl[
  \frac{\varphi(t)+e^{-t}\varphi(-t)}{1-e^{-t}}
  -\frac{2e^{-t}\varphi(0)}t\Bigr]dt.
  \end{aligned}
  \]
\item The \emph{geometric side} of the explicit formula of $K$ at
  $\varphi$ is
  \[
  \begin{aligned}
  G_K(\varphi)={}&-\log|D_K|\,\varphi(0)+r_1W_\R(\varphi)
  +r_2W_\C(\varphi)\\
  &+\sum_{\mathfrak p}\log N\mathfrak p\sum_{k\ge1}
  \bigl(\varphi(k\log N\mathfrak p)
  +N\mathfrak p^{-k}\varphi(-k\log N\mathfrak p)\bigr),
  \end{aligned}
  \]
  the sum running over the nonzero prime ideals of $\cO_K$.
\end{enumerate}
\end{definition}

The prime sum is finite because $\varphi$ has compact support; both
integrands are bounded near $t=0$, where the two terms of size
$\varphi(0)/t$ cancel, and are integrable at infinity. The geometric
side involves only the prime ideals, the archimedean places, and the
discriminant of $K$, not its zeros. For $\varphi\in\mathcal S^+$ it
reduces to
\begin{equation}\label{eq:geom-one-sided}
G_K(\varphi)=\sum_{\mathfrak p}\log N\mathfrak p\sum_{k\ge1}
\varphi(k\log N\mathfrak p)+r_1\int_0^\infty
\frac{\varphi(t)\,dt}{1-e^{-2t}}+r_2\int_0^\infty
\frac{\varphi(t)\,dt}{1-e^{-t}}.
\end{equation}

\begin{definition}\label{def:aw6prime}
Let $K$ be a number field, let $(V^\bullet,\Theta,\cup,\Tr,*)$ be a
Weil datum of type $(1,0)$ on $\Xbar_K$, and let
$\mathcal S'\subset\mathcal S$.
\begin{enumerate}[label=(\alph*)]
\item The datum satisfies \emph{AW6$'$ for $(\Lambda_K,\mathbb 1)$ on
  $\mathcal S'$} if each $(V^j,\Theta_j)$ satisfies (G) and, for every
  $\varphi\in\mathcal S'$, the three distributional traces converge
  absolutely and
  \begin{equation}\label{eq:aw6prime}
  \sum_{j=0}^2(-1)^j\Tr_{\mathrm{dis}}(V^j)(\varphi)=G_K(\varphi).
  \tag{AW6$'$}
  \end{equation}
  \emph{AW6$'$} without qualification means $\mathcal S'=\mathcal S$.
\item Suppose that each star form $h_j$ is positive definite and that
  the datum satisfies \ref{ax:AW5}, with unitary groups $U_j$ on the
  completions $\cH^j$. The datum satisfies \emph{AW6$'$ in Hilbert form
  on $\mathcal S'$} if each $(V^j,\Theta_j)$ satisfies (G) and, for
  every $\varphi\in\mathcal S'$, the operators
  $A^{(j)}_\varphi=\int_\R\varphi(t)e^{jt/2}U_j(t)\,dt$ (strong
  integrals) are of trace class and
  $\sum_{j=0}^2(-1)^j\tr A^{(j)}_\varphi=G_K(\varphi)$.
\end{enumerate}
\end{definition}

In the functional calculus of the self-adjoint operator $T_j$ of
\ref{ax:AW5}, $A^{(j)}_\varphi=f_j(T_j)$ with
$f_j(\tau)=\Phi(\tfrac j2+i\tau)$. So a trace ``$\tr f(T_j)$'' of the
schematic \ref{ax:AW6} must use a test function depending on the degree
$j$, and the terms $\Phi(0)$ and $\Phi(1)$ of the explicit formula are
the traces in degrees $0$ and $2$. The test function must also be the
transform of a compactly supported $\varphi$, not an arbitrary Schwartz
function of $\tau$ (Proposition~\ref{prop:schwartz}). The form (a) uses
only the spectral cores, with no Hilbert space, positivity, or star; the
form (b) presupposes \ref{ax:AW5}, which presupposes that $h_j$ is
positive definite and then yields AW4-V
(Proposition~\ref{prop:aw5-contains}).
\begin{remark}\label{rem:aw-comments}
The following comments record the scope of the axioms. The
quaternionic refinement $*_1^2=-1$ in degree one
\cite[2.7, p.~3]{Deninger2022NoWeil} is compatible with
\ref{ax:AW4} and is used only in Proposition~\ref{prop:root-parity}. By
\ref{ax:AW1} and \ref{ax:AW3}, for eigenvectors $u\in V^j$, $v\in
V^{2d-j}$ with eigenvalues $\rho,\sigma$ one has
$d\Tr(u\cup v)=\Tr(\Theta(u\cup v))=(\rho+\sigma)\Tr(u\cup v)$, so
the pairing is supported on $\sigma=d-\rho$; perfectness then gives
$m_j(\rho)=m_{2d-j}(d-\rho)$, the spectral shadow of the functional
equation. Spectral duality determines the reflected divisor, not the
nonvanishing factors of a functional equation; conductor and epsilon
data are extra normalization data. Axioms \ref{ax:AW5}--\ref{ax:AW6}
are not used in Theorems~\ref{thm:rh} and~\ref{thm:A};
\ref{ax:AW5} presupposes that $h_j$ is positive definite, and with its
unitary flow it yields AW4-V (Proposition~\ref{prop:aw5-contains});
both make the operator and
trace-formula comparisons of Sections~\ref{sec:debranges}
and~\ref{sec:open} meaningful. Finally, the data must be constructed
from arithmetic input: defining a vector for each known zero is a formal
encoding, not the desired construction
(Propositions~\ref{prop:converse} and~\ref{prop:trivial-model}).
\end{remark}

It is natural to expect that, as over a finite field, the zeta
function of a higher-dimensional arithmetic object has all its zeros on
one line, and that zeros and poles coming from different degrees may
cancel in the alternating product; in fact under \ref{ax:AW4} each
degree has its own line, the lines are disjoint, and no cancellation
occurs (Theorem~\ref{thm:A-schemes},
Propositions~\ref{prop:P1Z} and~\ref{prop:no-positivity-cancel}).

\begin{remark}\label{rem:aw-analogy}
Over $\F_q$ the zeta function is a rational function of $q^{-s}$ and
purity is read off from characteristic polynomials of Frobenius. Over a
number field there is no single $q$, and $\Lambda_K$ is not such a
rational function: \ref{ax:AW2} replaces the cohomological expression
of $Z(X,t)$ by the divisor of a flow, \ref{ax:AW3} replaces Poincar\'e
duality for Frobenius by duality for a flow, and \ref{ax:AW4} replaces
the positive Frobenius polarization of Definition~\ref{def:frob-pol} by
an infinitesimal one, exactly as in the suspension of
Proposition~\ref{prop:suspension}. Table~\ref{tab:dictionary} in
Section~\ref{sec:overview} lists the correspondence. For a pure
motive of weight $w$, the expected critical line of its $L$-function
is $\Re s=(w+1)/2$; this is the separate graded spectral convention.
\end{remark}

\section{Purity and the conditional Riemann hypothesis}
\label{sec:purity}

The purity statement, Proposition~\ref{prop:purity}, is stated in
Section~\ref{sec:conventions}, before its first use.

\begin{proposition}[{\cite[p.~5]{Deninger2010HilbertPolya}}]%
\label{prop:polarization-identity}
Let $d\ge0$ be an integer, and let $(V^\bullet,\Theta,\cup,
\Tr,*)$ be a Weil datum of type $(d,0)$ (Definition~\ref{def:aw})
satisfying the derivation rule of \ref{ax:AW1}, the trace equivariance
of \ref{ax:AW3}, and the star law of \ref{ax:AW4}. Then for each $j$ the
star form $h_j(u,v)=\Tr(u\cup*_jv)$ satisfies
$h_j(\Theta_ju,v)+h_j(u,\Theta_jv)=j\,h_j(u,v)$ for all
$u,v\in V^j$.
\end{proposition}

\begin{proof}
The star law gives $*_j(\Theta_jv)=\Theta_{2d-j}(*_jv)-(d-j)*_jv$, and
$\Tr(u\cup*_j\Theta_jv)=h_j(u,\Theta_jv)$ by definition. Hence, using
the derivation rule and then trace equivariance,
\begin{align*}
h_j(\Theta_ju,v)+h_j(u,\Theta_jv)
&=\Tr(\Theta_ju\cup*_jv)+\Tr(u\cup\Theta_{2d-j}*_jv)
  -(d-j)\Tr(u\cup*_jv)\\
&=\Tr\bigl(\Theta_{2d}(u\cup*_jv)\bigr)-(d-j)h_j(u,v)
 =dh_j(u,v)-(d-j)h_j(u,v).\qedhere
\end{align*}
\end{proof}

The computation is Deninger's \cite[p.~5]{Deninger2010HilbertPolya},
\cite[(7.12)--(7.13), p.~33]{Deninger1994Motivic},
\cite[2.7, p.~3]{Deninger2022NoWeil}.
Thus under \ref{ax:AW4} the star form $h_j$ is a positive infinitesimal
polarization of weight $j$ on $(V^j,\Theta_j)$: the star form of the
positivity axiom implies AW4-V in every degree. The algebra form does
too.

\begin{proposition}\label{prop:aw4-forms}
Let $(V^\bullet,\Theta,\cup,\Tr,*)$ be a Weil datum of type $(d,0)$,
and let $0\le j\le2d$.
\begin{enumerate}[label=(\roman*)]
\item\label{item:aw4e-pure} If the datum satisfies AW4-E in degree
  $j$ (Definition~\ref{def:aw4-forms}), then $\Theta_j$ is
  diagonalizable on every generalized eigenspace, every eigenvalue
  $\rho$ of $\Theta_j$ satisfies $\Re\rho=j/2$, and the datum
  satisfies AW4-V in degree $j$.
\item\label{item:aw4v-e} If the datum satisfies AW4-V in degree $j$,
  then it satisfies AW4-E in degree $j$.
\item\label{item:aw4v-star} If the datum satisfies the derivation rule
  of \ref{ax:AW1}, satisfies \ref{ax:AW3}, and satisfies AW4-V in every
  degree, then there are antilinear maps $*_j':V^j\to V^{2d-j}$ such
  that the datum with stars $*'_j$ satisfies \ref{ax:AW4}, with star
  forms equal to the given polarizations.
\end{enumerate}
\end{proposition}

\begin{proof}[Proof of \ref{item:aw4e-pure}]
Let $V'$, $A=A_{V'}$, $\theta=\theta_{V'}$, and $(\tau,\dagger)$ be as in
Definition~\ref{def:aw4-forms}. Put $\langle a,b\rangle=\tau(ab^\dagger)$
on $A$. It is linear in $a$ and antilinear in $b$, Hermitian because
$\overline{\langle b,a\rangle}=\tau((ba^\dagger)^\dagger)=
\tau(ab^\dagger)$, and positive definite by hypothesis. Let
$L_\theta(a)=\theta a$. Then
$\langle\theta a,b\rangle=\tau(\theta ab^\dagger)=\tau(ab^\dagger\theta)
=\tau(a(\theta^\dagger b)^\dagger)=\langle a,\theta^\dagger b\rangle$,
so the adjoint of $L_\theta$ is $L_{\theta^\dagger}=j-L_\theta$,
and $\langle\,,\rangle$ is a positive infinitesimal polarization of
weight $j$ on $(A,L_\theta)$. By
Proposition~\ref{prop:purity}\ref{item:pur-ss} and~\ref{item:pur-real}
(with $V=A$, which is finite-dimensional), $L_\theta$ is diagonalizable
with eigenvalues of real part $j/2$, so its minimal polynomial $m$
is separable with roots on $\Re s=j/2$. Since $A$ is unital,
$p(L_\theta)=0$ if and only if $p(\theta)=0$ (evaluate at $1$, and
conversely $p(L_\theta)a=p(\theta)a$), so $m$ is the minimal polynomial
of $\theta$ in $A$, hence in $\operatorname{End}_\C(V')$. Therefore
$\theta$ is diagonalizable on $V'$ with eigenvalues of real part
$j/2$. Let $v\in V^j_{\Theta\sim\rho}$ be nonzero, with
$(\Theta_j-\rho)^kv=0$, and choose $V'\ni v$ as above. Then $\rho$ is an
eigenvalue of $\theta$, so $\Re\rho=j/2$, and since $\theta$ is
diagonalizable on $V'$, $(\theta-\rho)^kv=0$ forces $(\theta-\rho)v=0$.
So $V^j=\bigoplus_\rho\ker(\Theta_j-\rho)$ with
$\Re\rho=j/2$. Let $h$ be the orthogonal sum of arbitrary positive
definite Hermitian forms on the finite-dimensional eigenspaces. For
eigenvectors $u,v$ with eigenvalues $\rho,\rho'$,
$h(\Theta u,v)+h(u,\Theta v)=(\rho+\bar\rho')h(u,v)$, which is $0$ if
$\rho\ne\rho'$ and $jh(u,v)$ if $\rho=\rho'$; so $h$ is a positive
infinitesimal polarization of weight $j$.
\end{proof}

\begin{proof}[Proof of \ref{item:aw4v-e}]
Let $h$ be a positive infinitesimal polarization of weight $j$ on
$(V^j,\Theta_j)$. Every vector of $V^j$ lies in a finite sum $V'$ of
generalized eigenspaces, which is finite-dimensional and
$\Theta_j$-stable. Take $A_{V'}=\operatorname{End}_\C(V')$, $\tau_{V'}$ the
trace, and $a^\dagger$ the $h|_{V'}$-adjoint of $a$. Then $(\tau_{V'},\dagger)$
is a positive trace involution: in an $h$-orthonormal basis
$a^\dagger$ is the conjugate transpose and $\tau_{V'}(aa^\dagger)$ is the
sum of the squared moduli of the entries of $a$. Identity
\eqref{eq:inf-pol} says $\theta_{V'}^\dagger=j-\theta_{V'}$.
\end{proof}

\begin{proof}[Proof of \ref{item:aw4v-star}]
Fix $j$ and a positive infinitesimal polarization $h$ of weight
$j$ on $V^j$. By Proposition~\ref{prop:purity}, $V^j$ is the direct
sum of the eigenspaces $\ker(\Theta_j-\rho)$, which are $h$-orthogonal.
For $v$ in the $\rho$-eigenspace the linear form $u\mapsto h(u,v)$
vanishes on the other eigenspaces, and by \ref{ax:AW3} the pairing $B_j$
restricts to a perfect pairing of $V^j_{\Theta\sim\rho}$ with
$V^{2d-j}_{\Theta\sim d-\rho}$ and vanishes between
$V^j_{\Theta\sim\rho'}$ and $V^{2d-j}_{\Theta\sim d-\rho}$ for
$\rho'\ne\rho$ (Remark~\ref{rem:aw-comments}). So there is a unique
$*'_jv\in V^{2d-j}_{\Theta\sim d-\rho}$ with $B_j(u,*'_jv)=h(u,v)$ for
all $u\in V^j$; the map $*'_j$ is antilinear, and its star form is
$h$. The derivation rule and the trace equivariance give
$B_j(\Theta u,x)+B_j(u,\Theta x)=dB_j(u,x)$, so
\[
B_j(u,\Theta*'_jv)=d\,h(u,v)-h(\Theta u,v)
=(d-j)h(u,v)+h(u,\Theta v)=B_j\bigl(u,(d-j)*'_jv+*'_j\Theta v\bigr)
\]
for all $u$, using \eqref{eq:inf-pol}; since $B_j$ is perfect, this is
the star law.
\end{proof}

Thus, as bare existence statements, AW4-V and AW4-E in degree $j$ are
equivalent, and both are equivalent to ``$\Theta_j$ is semisimple with
all eigenvalues on $\Re s=j/2$''; in presence of the derivation rule
and \ref{ax:AW3}, AW4-V in every degree is equivalent to \ref{ax:AW4}
(Proposition~\ref{prop:polarization-identity} and
\ref{item:aw4v-star}). The two forms differ in what a
construction must supply: a form on the whole of $V^j$, or a positive
involution on an algebra of operators containing the flow generator.
Over $\F_q$ the second is what is actually proved for curves
(Proposition~\ref{prop:weil-curves} and
Proposition~\ref{prop:ff-N7}).

\begin{theorem}[{Deninger's formal argument, after Serre;
\cite[(7.11)--(7.13), pp.~32--33]{Deninger1994Motivic},
\cite[p.~169]{Deninger1998Analogies},
\cite[pp.~4--5]{Deninger2010HilbertPolya},
\cite[2.7, p.~3]{Deninger2022NoWeil}}]\label{thm:rh}
Let $K$ be a number field and suppose that the single object
$\Xbar_K$ carries a Weil datum of type $(1,0)$ satisfying
\ref{ax:AW1}, \ref{ax:AW2}, and \ref{ax:AW3} for
$(\Lambda_K,\mathbb 1)$, and satisfying either \ref{ax:AW4}, or AW4-V
in degree one, or AW4-E in degree one. Then:
\begin{enumerate}[label=(\roman*)]
\item\label{item:rh-divisor} $\operatorname{Sp}_1=\divisor\xi_K$, where
  $\xi_K(s)=s(s-1)\Lambda_K(s)$;
\item\label{item:rh-line} every zero of $\xi_K$ has real part
  $\tfrac12$;
\item\label{item:rh-ss} $\Theta_1$ is diagonalizable, with
  $\dim\ker(\Theta_1-\rho)=\ord_\rho\xi_K$ for every $\rho\in\C$.
\end{enumerate}
\end{theorem}

\begin{proof}[Proof of \ref{item:rh-divisor}]
By \ref{ax:AW1} with $U=\mathbb 1$, $(V^0,\Theta_0)=(\C,0)$ and
$(V^2,\Theta_2)\cong(\C,1)$, so $\operatorname{Sp}_0=[0]$ and
$\operatorname{Sp}_2=[1]$. By Theorem~\ref{thm:imp-hecke}, $\Lambda_K$
has simple poles at $0$ and $1$ and no others, so
$\divisor\xi_K=\divisor\Lambda_K+[0]+[1]$. With $d=1$,
\eqref{eq:aw-div} reads $\divisor\Lambda_K=-[0]+\operatorname{Sp}_1-[1]$.
\end{proof}

\begin{proof}[Proof of \ref{item:rh-line} and \ref{item:rh-ss}]
Under \ref{ax:AW4}, Proposition~\ref{prop:polarization-identity} (with
$d=1$, $w=0$, $j=1$) shows that $h_1$ is a positive infinitesimal
polarization of weight one; under AW4-E in degree one,
Proposition~\ref{prop:aw4-forms}\ref{item:aw4e-pure} provides one; under
AW4-V it is given. Proposition~\ref{prop:purity} then gives
$\Re\rho=\tfrac12$ for every eigenvalue $\rho$ of $\Theta_1$ and
$V^1_{\Theta\sim\rho}=\ker(\Theta_1-\rho)$. By~\ref{item:rh-divisor} the
zeros of $\xi_K$, with multiplicity, are the eigenvalues of $\Theta_1$,
with multiplicity $m_1(\rho)=\dim\ker(\Theta_1-\rho)$.
\end{proof}

By Proposition~\ref{prop:xiK}, every zero of $\zeta_K$ with
$0<\Re s<1$ is a zero of $\xi_K$ of the same order, so the conclusion
contains the Riemann hypothesis for $\zeta_K$; for $K=\Q$, where
$\xi_\Q=2\xi$, it is the Riemann hypothesis. The argument does not
presuppose that zeros are simple, uses no regularized determinant and
neither \ref{ax:AW5} nor \ref{ax:AW6}, and uses from \ref{ax:AW3} only
the trace equivariance. In particular the normalization term of
\eqref{eq:aw-det} never enters: \ref{ax:AW2} in the form
\eqref{eq:aw-div} identifies the zeros with the degree-one spectrum
independently of any determinant convention. The statement is
conditional: no structure satisfying \ref{ax:AW1}--\ref{ax:AW4} on
$\Xbar_\Q$ is known.

Theorem~\ref{thm:rh} is Deninger's formal implication
\cite[(7.11)--(7.13), pp.~32--33]{Deninger1994Motivic},
\cite[pp.~168--169]{Deninger1998Analogies},
\cite[pp.~4--5]{Deninger2010HilbertPolya}, including the semisimplicity
in~\ref{item:rh-ss} \cite[2.7, p.~3]{Deninger2022NoWeil}, modelled on
Serre's argument \cite{Serre1960Analogues}; in Flach--Morin's
formulation it is \cite[Conjecture~2.3(iv)--(v), p.~522]%
{FlachMorin2020DeningerWeilArakelov}. The statement here differs only
in bookkeeping: \ref{ax:AW2} is imposed in divisor form, the poles at
$0$ and $1$ are carried by the endpoint clause of \ref{ax:AW1} rather
than by distinct eigenvalues in degrees $0$, $1$, $2$
\cite[p.~168]{Deninger1998Analogies}, and \ref{ax:AW4} may be replaced
by either of its two other forms.

\begin{definition}\label{def:diag-op}
Let $\cH$ be a Hilbert space with an orthonormal basis $(e_k)$ and let
$(\tau_k)$ be real numbers. The \emph{diagonal operator}
$T=\diag(\tau_k)$ has domain
$\mathcal D(T)=\{\sum c_ke_k:\sum|\tau_kc_k|^2<\infty\}$ and acts by
$T(\sum c_ke_k)=\sum\tau_kc_ke_k$.
\end{definition}

\begin{proposition}\label{prop:hilbert}
Let $(V,\Theta)$ be a spectral core, $j\in\R$, and $h$ a positive
infinitesimal polarization of weight $j$ on $(V,\Theta)$. Let $\cH$
be the Hilbert-space completion of $(V,h)$. Then there is an orthonormal
basis $(e_k)$ of $\cH$ contained in $V$ with $\Theta e_k=(j/2+i\tau_k)
e_k$, $\tau_k\in\R$, and the diagonal operator $T=\diag(\tau_k)$ of
Definition~\ref{def:diag-op} is self-adjoint, satisfies $\Theta=j/2+iT$
on $V$, and has pure point spectrum $\{\tau_k\}$ with finite
multiplicities and no finite accumulation.
\end{proposition}

\begin{proof}
By Proposition~\ref{prop:purity}, each $V_{\Theta\sim\rho}$ is an
eigenspace with $\rho=j/2+i\tau$, and distinct eigenspaces are
orthogonal. Choosing an orthonormal basis of each (finite-dimensional)
eigenspace gives an orthonormal family $(e_k)$ spanning $V$, hence an
orthonormal basis of $\cH$. The operator $T$ is symmetric on its domain.
If $y=\sum y_ke_k\in\mathcal D(T^*)$ with $T^*y=z$, then testing against
$x=e_k$ gives $z_k=\tau_ky_k$, so $\sum|\tau_ky_k|^2=\|z\|^2<\infty$,
$y\in\mathcal D(T)$, and $Ty=z$; hence $T=T^*$. The set
$\{\tau_k\}=\{(\rho-j/2)/i:\rho\in\spec\Theta\}$ has no finite
accumulation point because $\spec\Theta$ has none, so it is closed and
discrete. The spectrum of a diagonal operator with real entries is the
closure of the set of entries, which is therefore $\{\tau_k\}$, and each
eigenspace is spanned by the finitely many $e_k$ with that entry.
\end{proof}

Proposition~\ref{prop:hilbert} produces the diagonal completion, the
Hilbert space envisaged in \cite[p.~33]{Deninger1994Motivic} and
\cite[p.~169]{Deninger1998Analogies}. Axiom
\ref{ax:AW5} asks more: that a geometrically defined completion, with a
geometrically defined unitary flow, coincides with it, with a
skew-adjoint generator. A self-adjoint realization cannot be inferred
from skew-symmetry of $\Theta_j-j/2$ on a dense domain alone
(Proposition~\ref{prop:R4}).

\begin{proposition}\label{prop:converse}
The following are equivalent.
\begin{enumerate}[label=(\roman*)]
\item\label{item:conv-rh} The Riemann hypothesis holds.
\item\label{item:conv-model} There is a spectral core $(V,\Theta)$ with
  a positive infinitesimal polarization of weight one such that
  $\spec\Theta$ with
  multiplicities is the multiset of nontrivial zeros of $\zeta$.
\end{enumerate}
\end{proposition}

\begin{proof}[Proof of \ref{item:conv-rh}$\Rightarrow$\ref{item:conv-model}]
Let $V=\bigoplus_\rho\C^{m(\rho)}$ over the distinct nontrivial zeros
$\rho$, with $m(\rho)$ the order of the zero, $\Theta=\rho$ on the
$\rho$-summand, and $h$ the orthogonal sum of the standard forms. The
zeros are discrete (they are zeros of the entire function $\xi$), so
$(V,\Theta)$ is a spectral core. For $u,v$ in the $\rho$-summand,
$h(\Theta u,v)+h(u,\Theta v)=(\rho+\bar\rho)h(u,v)=h(u,v)$ because
$\Re\rho=\tfrac12$; distinct summands are orthogonal on both sides.
\end{proof}

\begin{proof}[Proof of \ref{item:conv-model}$\Rightarrow$\ref{item:conv-rh}]
Proposition~\ref{prop:purity}\ref{item:pur-real}.
\end{proof}

\begin{remark}\label{rem:reformulation}
Proposition~\ref{prop:converse} shows that the bare existence of a
positive spectral model is a reformulation of RH. The model in its
proof supplies no cup product, trace, functoriality, local comparison, or
trace formula, and it does not establish the determinant identity
\eqref{eq:aw-det} for the diagonal operator. The research problem is to
derive the spectrum and the positivity from an independently defined
arithmetic object; compare Propositions~\ref{prop:weil-linear}
and~\ref{prop:weil-converse} over $\F_q$. The idea of realizing the
zeros as a spectrum goes back to Hilbert and P\'olya; see
\cite[p.~3]{Connes1998TraceFormula}, \cite[p.~147]{Deninger1992Local},
and \cite[p.~5]{Sarnak2004Clay}.
\end{remark}

\subsection{Positivity cannot be weakened}\label{subsec:witnesses}

The first two negative findings (R1 and R2 of
Section~\ref{sec:negative}) are stated here, before the main theorem,
because Section~\ref{subsec:content} uses them: they show that
Proposition~\ref{prop:purity} fails as soon as its form is only
nondegenerate, or positive on part of the space.

It is natural to expect that duality for a flow, together with a
Hermitian form compatible with it, forces purity, and that positivity of
such a form on one spectral piece (for instance of a height pairing on
one eigenspace) propagates to all of them; in fact neither holds without
positivity on the whole space.

\begin{definition}\label{def:R1-forms}
Let $j\in\R$ and let $V=\C^2$ with standard basis $e_1,e_2$.
\begin{enumerate}[label=(\alph*)]
\item For $\rho\in\C$, the \emph{dual pair} is
  $\Theta^{\mathrm{dual}}_\rho=\diag(\rho,\,j-\bar\rho)$ with the
  form $h^{\mathrm{dual}}(u,v)=u_1\bar v_2+u_2\bar v_1$.
\item For real $\tau_1\ne\tau_2$, the \emph{split pair} is
  $\Theta^{\mathrm{split}}=\diag(j/2+i\tau_1,\,j/2+i\tau_2)$ with the
  form $h^{\mathrm{split}}(u,v)=u_1\bar v_1-u_2\bar v_2$.
\end{enumerate}
\end{definition}

\begin{proposition}\label{prop:R1}
Let $j\in\R$ and let $V=\C^2$ with standard basis $e_1,e_2$.
\begin{enumerate}[label=(\roman*)]
\item\label{item:R1-dual} For $\rho\in\C$, $h^{\mathrm{dual}}$ is a
  nondegenerate Hermitian form satisfying \eqref{eq:inf-pol} for
  $\Theta=\Theta^{\mathrm{dual}}_\rho$, and
  $\Theta^{\mathrm{dual}}_\rho$ has the eigenvalue $\rho$; in
  particular $\Re\rho\ne j/2$ is possible.
\item\label{item:R1-partial} For real $\tau_1\ne\tau_2$,
  $h^{\mathrm{split}}$ is a nondegenerate Hermitian form satisfying
  \eqref{eq:inf-pol} for $\Theta=\Theta^{\mathrm{split}}$, it is
  positive definite on the eigenline
  $\ker(\Theta^{\mathrm{split}}-j/2-i\tau_1)$, and
  $h^{\mathrm{split}}(e_2,e_2)=-1$.
\end{enumerate}
\end{proposition}

\begin{proof}[Proof of \ref{item:R1-dual}]
Put $\Theta=\Theta^{\mathrm{dual}}_\rho$ and $h=h^{\mathrm{dual}}$, and
write $h(u,v)=u^{T}H\bar v$ with $H=\bigl(\begin{smallmatrix}0&1\\1&0
\end{smallmatrix}\bigr)$; $H$ is real symmetric and invertible, so $h$
is Hermitian and nondegenerate. Identity \eqref{eq:inf-pol} is
equivalent to $\Theta^TH+H\bar\Theta=jH$. Both sides vanish on the
diagonal, and the off-diagonal entries of the left side are
$\rho+\overline{(j-\bar\rho)}=j$ and $(j-\bar\rho)+\bar\rho=j$.
\end{proof}

\begin{proof}[Proof of \ref{item:R1-partial}]
Put $\Theta=\Theta^{\mathrm{split}}$ and $h=h^{\mathrm{split}}$. Here
$H=\diag(1,-1)$, and $\Theta^TH+H\bar\Theta$ is diagonal with
entries $2\Re(j/2+i\tau_k)(\pm1)=j(\pm1)$, so \eqref{eq:inf-pol}
holds. The eigenline for $j/2+i\tau_1$ is $\C e_1$, on which
$h(ce_1,ce_1)=|c|^2$; and $h(e_2,e_2)=-1$.
\end{proof}

In \ref{item:R1-dual} both eigenlines are $h$-isotropic, as they must be
by the computation in the proof of
Proposition~\ref{prop:purity}\ref{item:pur-real}; positivity, not
nondegeneracy, is what excludes this. Item \ref{item:R1-partial} is the
orthogonal sum of a positive and a negative line. Deninger's comparison
with heights \cite[pp.~7--9]{Deninger2010HilbertPolya} is therefore
motivation and a compatibility test, not a construction of the full
positive form required by \ref{ax:AW4}. Proposition~\ref{prop:R2} below
shows that the gap persists in presence of a cup product, a trace, and a
quaternionic star.

It is natural to expect that a quaternionic star, perfect weight-one
duality, and the star law of \ref{ax:AW4} together force eigenvalues onto
the critical line; in fact all of this structure is compatible with
eigenvalues off the line, the form $h_1$ then being indefinite.

\begin{definition}\label{def:R2-block}
Let $0<a<\tfrac12$, $T>0$, and $\lambda=\tfrac12+a+iT$. Let $V=\C^4$ with
standard basis $e_1,e_2,f_1,f_2$ (in this order), let $I_2$ be the
$2\times2$ identity matrix, and put
\[
\Theta=\diag(\lambda,\bar\lambda,1-\lambda,1-\bar\lambda),\quad
\Omega=\begin{pmatrix}0&I_2\\-I_2&0\end{pmatrix},\quad
R=\begin{pmatrix}0&-1\\1&0\end{pmatrix},\quad
S=\begin{pmatrix}R&0\\0&R\end{pmatrix}.
\]
Let $B(v,w)=v^T\Omega w$ and let $J:V\to V$ be the antilinear map
$Jv=S\bar v$. The \emph{quaternionic block} is the graded algebra
$V^0=\C\cdot1$, $V^1=V$, $V^2=\C\eta$ with unit $1$, products
$v\cup w=B(v,w)\eta$ for $v,w\in V^1$ and all other products of
positive-degree elements zero, generators $\Theta_0=0$,
$\Theta_1=\Theta$, $\Theta_2=1$, trace $\Tr(c\eta)=c$, and antilinear
stars $*_0(c\cdot1)=\bar c\eta$, $*_1=J$, $*_2(c\eta)=\bar c\cdot1$.
Put $h_j(u,v)=\Tr(u\cup*_jv)$. We also use the antilinear map
$J_+v=-\Omega\bar v$ on $V$, the matrix
$\Pi=\bigl(\begin{smallmatrix}0&I_2\\I_2&0\end{smallmatrix}\bigr)$, the
characteristic polynomial $p(s)=\det(s-\Theta)$, and the variable
$w=s-\tfrac12$.
\end{definition}

\begin{proposition}\label{prop:R2}
Let $0<a<\tfrac12$ and $T>0$, and let the quaternionic block of
Definition~\ref{def:R2-block} be given. Then:
\begin{enumerate}[label=(\roman*)]
\item\label{item:R2-axioms} the block, regarded as a Weil datum of type
  $(1,0)$ (Definition~\ref{def:aw}), satisfies \ref{ax:AW1} for $U=\mathbb 1$,
  \ref{ax:AW3}, and the star law of \ref{ax:AW4} in degrees $0,1,2$;
  moreover $J^2=-1$, and $h_0,h_2$ are positive definite;
\item\label{item:R2-form} $h_1(v,w)=v^TH\bar w$ with
  $H=\Omega S=\bigl(\begin{smallmatrix}0&R\\-R&0\end{smallmatrix}\bigr)$,
  $h_1$ is Hermitian and nondegenerate of signature $(2,2)$, and
  $h_1(e_1+f_2,e_1+f_2)=-2$;
\item\label{item:R2-off} every eigenvalue of $\Theta_1$ has real part
  $\tfrac12+a$ or $\tfrac12-a$, hence none lies on $\Re s=\tfrac12$.
\end{enumerate}
\end{proposition}

\begin{proof}[Proof of \ref{item:R2-axioms}]
Since $\Theta$ is diagonal and the only nonzero entries of $\Omega$ are
at positions $(1,3),(2,4)$ (value $1$) and $(3,1),(4,2)$ (value $-1$),
the identity $\Theta^T\Omega+\Omega\Theta=\Omega$ reduces to
$\lambda+(1-\lambda)=1$ and $\bar\lambda+(1-\bar\lambda)=1$. Hence
$B(\Theta v,w)+B(v,\Theta w)=B(v,w)$, which is the derivation rule in
degrees $1+1$; in all other bidegrees the derivation rule reads $0=0$ or
involves the unit. The form $B$ is alternating because $\Omega^T=-\Omega$,
so the product is graded-commutative; it is associative because every
product of three positive-degree elements lies in degree at least $3$,
where the algebra vanishes. The normalizations $V^0=\C\cdot1$,
$\Theta_0=0$, $V^2=\C$, $\Theta_2=1$ hold by definition. For
\ref{ax:AW3}, $\Tr(\Theta_2x)=\Tr(x)$; the pairing $V^0\times V^2\to\C$
is $(c,c'\eta)\mapsto cc'$; and on $V^1$ the eigenlines are $\C e_1,\C
e_2,\C f_1,\C f_2$ with eigenvalues $\lambda,\bar\lambda,1-\lambda,
1-\bar\lambda$, and $B(e_1,f_1)=B(e_2,f_2)=1$, so the eigenline for
$\rho$ is paired perfectly with that for $1-\rho$. For the star law in
degree one, $J\Theta=\Theta J$ amounts to $\Theta S=S\bar\Theta$; since
$S$ interchanges the first two and the last two coordinates up to sign,
this reduces to $\lambda=\overline{\bar\lambda}$ and
$1-\lambda=\overline{1-\bar\lambda}$. In degree zero,
$\Theta_2*_0=*_0\Theta_0+*_0$ holds since both sides send $c$ to $\bar
c\eta$; in degree two, $\Theta_0*_2=0=*_2\Theta_2-*_2$. Since $S$ is real
and $S^2=-I$, $J^2v=S\overline{S\bar v}=S^2v=-v$. Finally
$h_0(c,c')=h_2(c\eta,c'\eta)=c\bar c'$.
\end{proof}

\begin{proof}[Proof of \ref{item:R2-form}]
$h_1(v,w)=B(v,Jw)=v^T\Omega S\bar w$, and a direct multiplication gives
$\Omega S=\bigl(\begin{smallmatrix}0&R\\-R&0\end{smallmatrix}\bigr)$.
Since $R^T=-R$, $H$ is real symmetric, so $h_1$ is Hermitian. Moreover
$H^2=\diag(-R^2,-R^2)=I$ and $\tr H=0$, so the eigenvalues of $H$ are
$1$ and $-1$, each with multiplicity two. With $v=e_1+f_2=(1,0,0,1)$,
$v^THv=2H_{14}=2R_{12}=-2$.
\end{proof}

\begin{proof}[Proof of \ref{item:R2-off}]
The eigenvalues are $\tfrac12\pm a\pm iT$ and $0<a$.
\end{proof}

By Proposition~\ref{prop:polarization-identity}, $h_1$ satisfies
\eqref{eq:inf-pol} with $j=1$; Proposition~\ref{prop:purity} does not
apply because $h_1$ is indefinite, and indeed every eigenline is
$h_1$-isotropic ($H$ has zero diagonal). Axiom \ref{ax:AW2} is not
addressed: the block is a witness against deducing purity from
\ref{ax:AW1}, \ref{ax:AW3}, and the algebraic part of \ref{ax:AW4}, not a
model of all axioms.

It is natural to expect that, if the given star is not positive, one may
choose a positive star afterwards and keep the flow compatibility; in
fact positivity and compatibility with the flow must hold for the same
star.

\begin{proposition}\label{prop:R2-star}
Let $0<a<\tfrac12$ and $T>0$, and let $V,\Theta,\Omega,B,J_+,\Pi$ be
as in Definition~\ref{def:R2-block}. Then:
\begin{enumerate}[label=(\roman*)]
\item\label{item:R2s-pos} $J_+^2=-1$ and $B(v,J_+w)=v^T\bar w$, a
  positive definite Hermitian form;
\item\label{item:R2s-defect} $(\Theta J_+-J_+\Theta)v=-2a\,\Pi\bar v$
  for all $v\in V$; in particular $\Theta J_+\ne J_+\Theta$;
\item\label{item:R2s-none} there is no antilinear map $J':V\to V$ with
  $\Theta J'=J'\Theta$ for which $(v,w)\mapsto B(v,J'w)$ is a positive
  definite Hermitian form.
\end{enumerate}
\end{proposition}

\begin{proof}[Proof of \ref{item:R2s-pos}]
$J_+^2v=-\Omega\overline{(-\Omega\bar v)}=\Omega^2v=-v$, and
$B(v,J_+w)=-v^T\Omega^2\bar w=v^T\bar w$.
\end{proof}

\begin{proof}[Proof of \ref{item:R2s-defect}]
$\Theta J_+v-J_+\Theta v=-(\Theta\Omega-\Omega\bar\Theta)\bar v$. The
matrix $\Theta\Omega-\Omega\bar\Theta$ is supported on the nonzero
positions of $\Omega$, with entries
$\lambda-(1-\bar\lambda)$ at $(1,3)$,
$\bar\lambda-(1-\lambda)$ at $(2,4)$,
$-(1-\lambda)+\bar\lambda$ at $(3,1)$, and
$-(1-\bar\lambda)+\lambda$ at $(4,2)$; each equals
$\lambda+\bar\lambda-1=2a$. Hence $\Theta\Omega-\Omega\bar\Theta=2a\Pi$,
which is nonzero since $a>0$.
\end{proof}

\begin{proof}[Proof of \ref{item:R2s-none}]
Suppose $J'$ is such a map and put $h'(v,w)=B(v,J'w)$. Using
$\Theta J'=J'\Theta$ and $\Theta^T\Omega+\Omega\Theta=\Omega$ (proof of
Proposition~\ref{prop:R2}\ref{item:R2-axioms}),
\[
h'(\Theta v,w)+h'(v,\Theta w)=B(\Theta v,J'w)+B(v,\Theta J'w)
=B(v,J'w)=h'(v,w).
\]
Proposition~\ref{prop:purity}\ref{item:pur-real} with $j=1$ then gives
$\Re\rho=\tfrac12$ for every eigenvalue, contradicting
Proposition~\ref{prop:R2}\ref{item:R2-off}.
\end{proof}

The commutation defect in \ref{item:R2s-defect} has operator norm $2a$,
twice the distance of the eigenvalues from the critical line.

\section{The class of zeta functions and the main theorem}
\label{sec:class}

This section states the main theorem, for the class
\[
\cC_{\mathrm{Ded}}=\{\Lambda_K: K\text{ a number field}\}
\]
of completed Dedekind zeta functions, which contains
$\Lambda=\Lambda_\Q$. Its objects $\Xbar_K$ are the arithmetic curves,
so the finite-field theorem it mirrors is the Riemann hypothesis for
curves over $\F_q$ (Theorem~\ref{thm:B}); its analytic inputs
(continuation, functional equation, gamma factors) are theorems; and
the Artin formalism expresses finite-order Hecke and Artin $L$-functions
through it, so the theorem also covers their zeros in the critical strip.
The
section then isolates the axiom that carries the content
(Section~\ref{subsec:content}), treats Galois descent
(Section~\ref{subsec:descent}), places the finite-field theorem in the
same frame (Section~\ref{subsec:thmB}), and records the statement for
higher-dimensional arithmetic schemes (Section~\ref{subsec:schemes}).

\subsection{Artin and Hecke \texorpdfstring{$L$}{L}-functions}
\label{subsec:artin}

To avoid a clash with the zeros $\rho$ and the constant $\pi$, Galois
representations are denoted $\sigma$.

\begin{definition}\label{def:artin}
Let $k$ be a number field, $L/k$ a finite Galois extension with group
$G$, and $\sigma:G\to\mathrm{GL}(V_\sigma)$ a representation on a
finite-dimensional complex vector space.
\begin{enumerate}[label=(\alph*)]
\item The \emph{Artin $L$-function} $L(\sigma,s)$ is the Euler product
  $\prod_{\mathfrak p}\det\bigl(1-\sigma(\mathrm{Frob}_{\mathfrak P})
  N\mathfrak p^{-s}\mid V_\sigma^{I_{\mathfrak P}}\bigr)^{-1}$ over the
  finite primes $\mathfrak p$ of $k$, where $\mathfrak P$ is a prime of
  $L$ above $\mathfrak p$ with inertia group $I_{\mathfrak P}$, for
  $\Re s>1$ \cite[pp.~9 and~11--12]{Cogdell2014ArtinLFunctions}.
\item The \emph{completed Artin $L$-function} is
  $\Lambda(\sigma,s)=L_\infty(\sigma,s)L(\sigma,s)$, where
  $L_\infty(\sigma,s)$ is the product over the archimedean places $v$ of
  $k$ of $\GR(s)^{n_+}\GR(s+1)^{n_-}$ if $v$ is real and $\sigma$
  restricted to a decomposition group at $v$ is $n_+\cdot1\oplus n_-\cdot
  \mathrm{sgn}$, and of $\GC(s)^{\dim\sigma}$ if $v$ is complex
  \cite[p.~13]{Cogdell2014ArtinLFunctions}.
\item The \emph{dual} and the \emph{conjugate} of $\sigma$ are
  $\sigma^\vee$ and $\bar\sigma$; an irreducible $\sigma$ is
  \emph{self-dual} if $\sigma\cong\sigma^\vee$, and \emph{defined over
  $\R$} if $V_\sigma\cong V_0\otimes_\R\C$ for a real representation
  $V_0$ of $G$.
\item For a modulus $\mathfrak m$ of $k$ and a character $\chi$ of the
  ray class group of $\mathfrak m$, the \emph{finite-order Hecke
  $L$-function} is the Weber $L$-series
  $L(\chi,s)=\prod_{\mathfrak p\nmid\mathfrak m}
  (1-\chi(\mathfrak p)N\mathfrak p^{-s})^{-1}$
  \cite[p.~8]{Milne2020CFT}.
\end{enumerate}
\end{definition}

For $\sigma=1_G$, $\Lambda(1_G,s)=\GR(s)^{r_1}\GC(s)^{r_2}\zeta_k(s)
=|D_k|^{-s/2}\Lambda_k(s)$, so $\Lambda(1_G,s)$ and $\Lambda_k$ have
the same divisor.

\begin{theorem}[Artin; Hecke; Brauer]\label{thm:imp-artin}
Let $k$, $L/k$, and $G$ be as in Definition~\ref{def:artin}.
\begin{enumerate}[label=(\roman*)]
\item\label{item:art-formal} The functions $\Lambda(\sigma,s)$ are
  additive in $\sigma$, unchanged by inflation from a quotient
  $\mathrm{Gal}(L'/k)$ of $G$, and inductive: for $H\subset G$ and a
  representation $\chi$ of $H=\mathrm{Gal}(L/L^H)$,
  $\Lambda(\mathrm{Ind}_H^G\chi,s)=\Lambda(\chi,s)$, the right side
  formed over $L^H$ \cite[pp.~9, 12, and~13]{Cogdell2014ArtinLFunctions}.
\item\label{item:art-fe} Each $\Lambda(\sigma,s)$ extends
  meromorphically to $\C$ and satisfies
  $\Lambda(\sigma,s)=\varepsilon(\sigma,s)\Lambda(\sigma^\vee,1-s)$,
  where $\varepsilon(\sigma,s)=W(\sigma)A_\sigma^{-(s-1/2)}$ with
  $|W(\sigma)|=1$ and a real number $A_\sigma>0$
  \cite[pp.~10--11, 13, and~15]{Cogdell2014ArtinLFunctions}.
\item\label{item:art-brauer} (Brauer) $L(\sigma,s)=\prod_i
  L(\psi_i,s)^{n_i}$ with $n_i\in\Z$, where $\psi_i$ are one-dimensional
  characters of subgroups $H_i$ of $G$ and $L(\psi_i,s)$ is formed over
  $L^{H_i}$ \cite[pp.~10--11, (2)]{Cogdell2014ArtinLFunctions},
  \cite[Theorem~2, p.~7]{Kramar2005ArtinBrauer}.
\item\label{item:art-weber} (Artin reciprocity) The Artin
  $L$-function of a one-dimensional representation of $\mathrm{Gal}(L/k)$
  is, up to finitely many Euler factors, a finite-order Hecke
  $L$-function of $k$ \cite[p.~10]{Milne2020CFT}.
\item\label{item:art-hecke} (Weber, Hecke) A finite-order Hecke
  $L$-function $L(\chi,s)$ with $\chi$ nontrivial extends to an entire
  function \cite[p.~8]{Milne2020CFT}; and if $M$ is the class field of a
  ray class group of $k$, then up to finitely many Euler factors
  $\zeta_M(s)=\prod_\chi L(\chi,s)$, the product over the characters
  $\chi$ of that ray class group
  \cite[p.~4]{Cogdell2014ArtinLFunctions}.
\end{enumerate}
\end{theorem}

The sources cited state these results; the proofs are in the works of
Artin, Hecke, and Brauer to which they refer, except that
\cite{Kramar2005ArtinBrauer} proves the group-theoretic part
of~\ref{item:art-brauer}. An Euler factor at a finite prime has zeros
and poles only on $\Re s=0$, so the finitely many factors
in~\ref{item:art-weber} and~\ref{item:art-hecke} do not affect zeros or
poles in $0<\Re s<1$. Two consequences are used repeatedly. First,
by~\ref{item:art-formal} and the decomposition of $\mathrm{Ind}_H^G1$
into irreducibles,
\begin{equation}\label{eq:artin-factor}
\divisor\Lambda_{L^H}=\sum_{\sigma}\langle\sigma,
\mathrm{Ind}_H^G1\rangle\divisor\Lambda(\sigma,s),
\end{equation}
the sum over the irreducible representations of $G$, where
$\langle\,,\rangle$ is the multiplicity pairing of characters. Second,
by~\ref{item:art-fe}, $\ord_{1-\rho}\Lambda(\sigma^\vee,s)=
\ord_\rho\Lambda(\sigma,s)$ for all $\rho$. If $\sigma$ is self-dual,
applying~\ref{item:art-fe} twice gives $W(\sigma)^2=1$, so
$W(\sigma)=\pm1$, and the function $t\mapsto A_\sigma^{t/2}
\Lambda(\sigma,\tfrac12+t)$ is even if $W(\sigma)=1$ and odd if
$W(\sigma)=-1$.

\begin{theorem}[{Frobenius--Schur; Fr\"ohlich--Queyrut; Serre}]
\label{thm:imp-fs}
Let $G$ be a finite group and $\sigma$ an irreducible complex
representation of $G$.
\begin{enumerate}[label=(\roman*)]
\item\label{item:fs-form} If $\sigma$ is self-dual, then up to scaling
  there is exactly one nonzero $G$-invariant bilinear form on
  $V_\sigma$; it is nondegenerate and either symmetric or alternating,
  and it is symmetric if and only if $\sigma$ is defined over $\R$
  \cite[Theorem~2.1, p.~2]{ConradFrobeniusSchur}.
\item\label{item:fs-fq} If $G=\mathrm{Gal}(L/k)$ for number fields and
  $\sigma$ is defined over $\R$, then $W(\sigma)=1$ (Fr\"ohlich and
  Queyrut, as quoted in \cite[p.~6]{Deninger2022NoWeil}).
\item\label{item:fs-even} If $\sigma$ is self-dual and not defined over
  $\R$, it occurs with even multiplicity in the complexification of
  every real representation of $G$
  \cite[\S13.2]{Serre1977LinearReps}, as quoted in
  \cite[p.~5]{Deninger2022NoWeil}.
\item\label{item:fs-omar} There is a Galois extension $L/\Q$ with group
  the quaternion group of order eight whose two-dimensional irreducible
  representation $\sigma$ has $W(\sigma)=-1$ (constructed by Armitage,
  as reported in \cite[\S4, p.~29]{Cougnard1983BasesNormales}; examples
  due to Omar are quoted in \cite[p.~5]{Deninger2022NoWeil}).
\end{enumerate}
\end{theorem}

The two-dimensional irreducible representation of the quaternion
group has character with values $2,-2,0$, which are rational, and is
self-dual and not defined over $\R$ \cite[Example~1.1, p.~1]%
{ConradFrobeniusSchur}.

\subsection{The main theorem}\label{subsec:thmA}

\begin{theorem}[{Theorem~A; (i) and~(iv) after Deninger, (ii) and~(iii)
by the Artin formalism}]\label{thm:A}
Suppose that for every number field $K$ the object $\Xbar_K$ carries a
Weil datum of type $(1,0)$ satisfying \ref{ax:AW1}, \ref{ax:AW2},
\ref{ax:AW3}, and \ref{ax:AW4} for $(\Lambda_K,\mathbb 1)$. Then:
\begin{enumerate}[label=(\roman*)]
\item\label{item:A-dedekind} for every number field $K$, every zero of
  $\xi_K$, that is every nontrivial zero of $\zeta_K$, has real part
  $\tfrac12$; in particular the Riemann hypothesis holds;
\item\label{item:A-hecke} every zero in $0<\Re s<1$ of every
  finite-order Hecke $L$-function of every number field has real part
  $\tfrac12$;
\item\label{item:A-artin} every zero and every pole in $0<\Re s<1$ of
  every Artin $L$-function $L(\sigma,s)$, for every finite Galois
  extension of number fields and every representation $\sigma$, has
  real part $\tfrac12$;
\item\label{item:A-ss} for every $K$, $\Theta_1$ is diagonalizable on
  $V^1_K$, with $\dim\ker(\Theta_1-\rho)=\ord_\rho\zeta_K$ for every
  nontrivial zero $\rho$ of $\zeta_K$ and $\ker(\Theta_1-\rho)=0$
  otherwise, and $\Theta_1-\tfrac12$ is skew-Hermitian for $h_1$.
\end{enumerate}
\end{theorem}

The hypothesis is used object by object: \ref{item:A-dedekind}
and~\ref{item:A-ss} for a given $K$ need the axioms only on $\Xbar_K$,
and \ref{item:A-hecke}--\ref{item:A-artin} only on the objects $\Xbar_M$
for the fields $M$ cut out by the characters involved; no functoriality
is used. By Theorem~\ref{thm:rh}, \ref{ax:AW4} may be replaced by AW4-V
or AW4-E in degree one; the conclusions remain valid, except that in
the last clause of~\ref{item:A-ss} the form $h_1$ is replaced by the
polarization provided by that hypothesis.

\begin{proof}[Proof of \ref{item:A-dedekind} and \ref{item:A-ss}]
Theorem~\ref{thm:rh} gives both, except for the orders and the last
clause. For $\rho$ with $0<\Re\rho<1$ the order of $\xi_K$ equals that
of $\zeta_K$ (Proposition~\ref{prop:xiK}), and by
Theorem~\ref{thm:rh}\ref{item:rh-line} every zero of $\xi_K$ lies in
that strip. By Proposition~\ref{prop:polarization-identity}, $h_1$ is a
positive infinitesimal polarization of weight one, which says that
$\Theta_1-\tfrac12$ is skew-Hermitian for $h_1$.
\end{proof}

\begin{proof}[Proof of \ref{item:A-hecke}]
Let $\chi$ be a character of the ray class group of a modulus of $k$,
$M$ the corresponding class field, and $\rho$ a zero of $L(\chi,s)$
with $0<\Re\rho<1$. By
Theorem~\ref{thm:imp-artin}\ref{item:art-hecke}, up to finitely many
Euler factors, which have no zeros or poles in the strip,
$\zeta_M(s)=\prod_{\chi'}L(\chi',s)$. Each factor with $\chi'\ne1$ is
entire, and the factor with $\chi'=1$ agrees with $\zeta_k$ up to
finitely many Euler factors, so it is holomorphic in the strip by
Theorem~\ref{thm:imp-hecke}. Hence $\ord_\rho\zeta_M\ge\ord_\rho
L(\chi,s)>0$; by Proposition~\ref{prop:xiK}, $\rho$ is a zero of
$\xi_M$, and $\Re\rho=\tfrac12$ by~\ref{item:A-dedekind} for $M$.
\end{proof}

\begin{proof}[Proof of \ref{item:A-artin}]
Let $\sigma$ be a representation of $G=\mathrm{Gal}(L/k)$. By
Theorem~\ref{thm:imp-artin}\ref{item:art-brauer}, $L(\sigma,s)=
\prod_iL(\psi_i,s)^{n_i}$ with $n_i\in\Z$, and
by Theorem~\ref{thm:imp-artin}\ref{item:art-weber} each
$L(\psi_i,s)$ is, up to finitely many
Euler factors, a finite-order Hecke $L$-function of $L^{H_i}$. As in
the previous proof each $L(\psi_i,s)$ is holomorphic on $0<\Re s<1$. If
$\rho$ in that strip is a zero or pole of $L(\sigma,s)$, then
$0\ne\ord_\rho L(\sigma,s)=\sum_in_i\ord_\rho L(\psi_i,s)$, so some
$L(\psi_i,s)$ vanishes at $\rho$, and $\Re\rho=\tfrac12$
by~\ref{item:A-hecke}.
\end{proof}

Items~\ref{item:A-dedekind} and~\ref{item:A-ss} are Deninger's
implication (Theorem~\ref{thm:rh} and the discussion after it); items
\ref{item:A-hecke} and~\ref{item:A-artin} are the classical
consequence of the Riemann hypothesis for all $\zeta_K$ through Brauer
induction, Artin reciprocity, and Hecke's theory
(Theorem~\ref{thm:imp-artin}).
Nothing is gained for zero location by enlarging the class to Hecke or
Artin $L$-functions: $\cC_{\mathrm{Ded}}$ already controls them through
the Artin formalism. What functoriality adds is multiplicity information
and Artin's holomorphy conjecture for the rational orbits of characters
(Proposition~\ref{prop:descent}).
Subsection~\ref{sec:lean} records which parts of Theorem~\ref{thm:A} are
formalized.
\subsection{Which axiom carries the content}\label{subsec:content}

It is natural to expect that a Weil datum with the correct spectrum
and perfect duality is already a nontrivial arithmetic statement, and
that a flow-compatible Hermitian form, or a positive one, is what forces
purity; in fact \ref{ax:AW1}--\ref{ax:AW3} can be satisfied with any
symmetric spectrum, each half of \ref{ax:AW4} can be satisfied
separately with zeros off the line, and the full package on one object
is equivalent to the Riemann hypothesis for that object. The same is
true over $\F_q$: a positive Frobenius-compatible form exists if and
only if Frobenius is semisimple and pure
(Proposition~\ref{prop:weil-converse}), and the content of Weil's
theorem lies in constructing the form from geometry.

\begin{definition}\label{def:trivial-model}
Let $\Sigma=\sum_\rho m(\rho)[\rho]$ be a divisor on $\C$ with
$m(\rho)\ge0$, no finite accumulation point of its support, invariant
under $\rho\mapsto1-\rho$, and with $m(\tfrac12)$ even. Choose
\begin{enumerate}[label=(\alph*)]
\item for each unordered pair $\{\rho,1-\rho\}$ with $\rho\ne\tfrac12$
  in the support, bases $(e_{\rho,k})_{k\le m(\rho)}$ and
  $(e_{1-\rho,k})_{k\le m(\rho)}$, one of them designated first;
\item a basis $(e_{1/2,k})_{k\le m(1/2)}$ grouped into consecutive pairs.
\end{enumerate}
The \emph{trivial model} with spectrum $\Sigma$ is the Weil datum of
type $(1,0)$ with $V^0=\C\cdot1$, $V^2=\C\eta$, $V^1$ the algebraic
direct sum of the lines $\C e_{\rho,k}$, generators $\Theta_0=0$,
$\Theta_2=1$, $\Theta_1e_{\rho,k}=\rho e_{\rho,k}$, the alternating form
$B$ on $V^1$ with $B(e,f)=1=-B(f,e)$ for each pair $(e,f)$ (the first
member of a pair in (a) being the designated vector, the pairs in (b)
in their given order) and $B=0$ on all other pairs of basis vectors,
products $u\cup v=B(u,v)\eta$ on $V^1$ and the unit products otherwise,
trace $\Tr(c\eta)=c$, and stars $*_0(c)=\bar c\eta$,
$*_2(c\eta)=\bar c$, and $*_1$ the antilinear map with $*_1e=f$,
$*_1f=-e$ on each pair $(e,f)$ if $\Sigma$ is supported on
$\Re s=\tfrac12$, and $*_1=0$ otherwise.
\end{definition}

\begin{proposition}\label{prop:trivial-model}
Let $\Sigma$ be as in Definition~\ref{def:trivial-model} and consider
its trivial model.
\begin{enumerate}[label=(\roman*)]
\item\label{item:tm-aw13} It satisfies \ref{ax:AW1} for $U=\mathbb 1$
  and \ref{ax:AW3}, and $\operatorname{Sp}_1=\Sigma$.
\item\label{item:tm-aw4} If $\Sigma$ is supported on $\Re s=\tfrac12$,
  it satisfies \ref{ax:AW4}, and $*_1^2=-1$.
\item\label{item:tm-iff} Some Weil datum with $\operatorname{Sp}_1=
  \Sigma$ satisfies \ref{ax:AW1} for $U=\mathbb 1$, \ref{ax:AW3}, and
  \ref{ax:AW4} if and only if $\Sigma$ is supported on $\Re s=\tfrac12$.
\end{enumerate}
\end{proposition}

\begin{proof}[Proof of \ref{item:tm-aw13}]
The spectrum of $\Theta_1$ is $\Sigma$ by construction and has no finite
accumulation point, so each $(V^j,\Theta_j)$ is a spectral core. The
product is associative (every product of three positive-degree elements
lies in degree at least three, where it vanishes) and
graded-commutative ($B$ is alternating). The derivation rule is
$0=0$ in the bidegrees involving $V^0$ trivially, and in bidegree
$(1,1)$ it reads $B(\Theta u,v)+B(u,\Theta v)=B(u,v)$, which holds on
basis vectors because $B(e_{\rho,k},e_{\rho',l})\ne0$ forces
$\rho+\rho'=1$. The endpoint conditions hold by construction. For
\ref{ax:AW3}, $\Tr(\Theta_2x)=\Tr(x)$; the pairing $V^0\times V^2$ is
$(c,c'\eta)\mapsto cc'$; and $B$ pairs the span of the
$e_{\rho,k}$ perfectly with the span of the $e_{1-\rho,k}$.
\end{proof}

\begin{proof}[Proof of \ref{item:tm-aw4}]
On the line, $1-\rho=\bar\rho$, so $*_1$ maps the $\rho$-eigenspace to
the $\bar\rho$-eigenspace and commutes with $\Theta_1$ (for an
antilinear map this is $\Theta_1*_1e=\bar\rho\,*_1e=*_1(\rho e)$), which
is the star law with $d=j=1$. The laws in degrees $0$ and $2$ are
$\Theta_2*_0=*_0\Theta_0+*_0$ and $\Theta_0*_2=*_2\Theta_2-*_2$, both
immediate. On a pair $(e,f)$, $h_1(e,e)=B(e,f)=1$,
$h_1(f,f)=B(f,-e)=1$, $h_1(e,f)=B(e,-e)=0$, and $h_1(f,e)=B(f,f)=0$;
distinct pairs are $h_1$-orthogonal. So $h_1$ is the standard positive
form in the basis, and $h_0(c,c')=h_2(c\eta,c'\eta)=c\bar c'$. Finally
$*_1^2e=*_1f=-e$ and $*_1^2f=-*_1e=-f$.
\end{proof}

\begin{proof}[Proof of \ref{item:tm-iff}]
If $\Sigma$ lies on the line, use \ref{item:tm-aw13}
and~\ref{item:tm-aw4}. Conversely, by
Proposition~\ref{prop:polarization-identity} the form $h_1$ is a
positive infinitesimal polarization of weight one, and
Proposition~\ref{prop:purity}\ref{item:pur-real} puts every eigenvalue
on the line.
\end{proof}

\begin{proposition}\label{prop:content}
Let $K$ be a number field and $\Sigma_K=\divisor\xi_K$.
\begin{enumerate}[label=(\roman*)]
\item\label{item:ct-free} The trivial model with spectrum $\Sigma_K$
  satisfies \ref{ax:AW1}, \ref{ax:AW2}, and \ref{ax:AW3} for
  $(\Lambda_K,\mathbb 1)$, unconditionally.
\item\label{item:ct-pos} The trivial model with spectrum $\Sigma_K$
  carries a positive definite Hermitian form on $V^1$ (not of the form
  $h_1$ and not satisfying \eqref{eq:inf-pol} unless $\Sigma_K$ lies on
  the line).
\item\label{item:ct-indef} For every divisor $\Sigma$ as in
  Definition~\ref{def:trivial-model} which is supported in
  $0<\Re s<1$ and invariant under $\rho\mapsto\bar\rho$ there is a Weil
  datum of type $(1,0)$ satisfying \ref{ax:AW1} for $U=\mathbb 1$,
  \ref{ax:AW3}, and the star law of \ref{ax:AW4}, with
  $\operatorname{Sp}_1=\Sigma$ and with star form $h_1$ Hermitian,
  nondegenerate, and satisfying \eqref{eq:inf-pol} with weight one;
  $h_1$ is indefinite as soon as $\Sigma$ has a point off the line.
\item\label{item:ct-iff} Some Weil datum on $\Xbar_K$ satisfies
  \ref{ax:AW1}--\ref{ax:AW4} for $(\Lambda_K,\mathbb 1)$ if and only if
  every zero of $\xi_K$ has real part $\tfrac12$.
\end{enumerate}
\end{proposition}

\begin{proof}[Proof of \ref{item:ct-free}]
$\Sigma_K$ is invariant under $\rho\mapsto1-\rho$ by
Proposition~\ref{prop:xiK}, its support has no finite accumulation
point, and $m(\tfrac12)=\ord_{1/2}\xi_K$ is even because
$t\mapsto\xi_K(\tfrac12+t)$ is an even function. By
Proposition~\ref{prop:trivial-model}\ref{item:tm-aw13} the trivial
model satisfies \ref{ax:AW1} and \ref{ax:AW3} with
$\operatorname{Sp}_1=\Sigma_K$, and then \eqref{eq:aw-div} holds exactly
as in the proof of Theorem~\ref{thm:rh}\ref{item:rh-divisor}, read
backwards.
\end{proof}

\begin{proof}[Proof of \ref{item:ct-pos}]
Take the form for which the basis vectors $e_{\rho,k}$ are orthonormal.
If it satisfied \eqref{eq:inf-pol}, Proposition~\ref{prop:purity}
would put $\Sigma_K$ on the line.
\end{proof}

\begin{proof}[Proof of \ref{item:ct-indef}]
The support of $\Sigma$ splits into the points on the line, the real
pairs $\{\rho,1-\rho\}$ with $\rho\ne\tfrac12$, and the quartets
$\{\lambda,\bar\lambda,1-\lambda,1-\bar\lambda\}$ with $\lambda$
nonreal and $\Re\lambda\ne\tfrac12$; the multiplicity is constant on
each pair and quartet. Take $V^0=\C\cdot1$, $V^2=\C\eta$ and let $V^1$
be the algebraic direct sum of the following blocks, with products
$u\cup v=B(u,v)\eta$ for the direct sum $B$ of the blocks' forms, and
star the direct sum of the blocks' stars. For the part on the line, the
trivial model of Definition~\ref{def:trivial-model}, whose star form is
positive (Proposition~\ref{prop:trivial-model}\ref{item:tm-aw4}). For
each real pair and each unit of multiplicity, the plane $\C e\oplus\C f$
with $\Theta=\diag(\rho,1-\rho)$, $B(e,f)=1=-B(f,e)$, $B(e,e)=
B(f,f)=0$, and $*(ae+bf)=-\bar ae+\bar bf$; then $*$ commutes with the
real diagonal $\Theta$, and $h_1(ae+bf,a'e+b'f)=a\bar b'+b\bar a'$ is
the dual-pair form of Definition~\ref{def:R1-forms} with $j=1$, of
signature $(1,1)$. For each quartet and each unit of multiplicity,
choose $\lambda$ in it with $\Re\lambda>\tfrac12$ and $\Im\lambda>0$
and take the block $V^1$ of Definition~\ref{def:R2-block} with
$a=\Re\lambda-\tfrac12\in(0,\tfrac12)$ and $T=\Im\lambda$, with its
form $B$ and star $J$; by Proposition~\ref{prop:R2} it satisfies the
axioms listed and its star form has signature $(2,2)$. The direct sum
satisfies \ref{ax:AW1} and \ref{ax:AW3} blockwise as in the proof of
Proposition~\ref{prop:trivial-model}\ref{item:tm-aw13} (the form $B$
pairs the $\rho$- and $(1-\rho)$-eigenspaces of each block), and the
star law blockwise. By Proposition~\ref{prop:polarization-identity},
$h_1$ satisfies \eqref{eq:inf-pol}. It is the orthogonal sum of the
blocks' forms, hence nondegenerate, and indefinite if some block is off
the line.
\end{proof}

\begin{proof}[Proof of \ref{item:ct-iff}]
If every zero of $\xi_K$ lies on the line, the trivial model with
spectrum $\Sigma_K$ satisfies \ref{ax:AW1}--\ref{ax:AW3} by
\ref{item:ct-free} and \ref{ax:AW4} by
Proposition~\ref{prop:trivial-model}\ref{item:tm-aw4}. The converse is
Theorem~\ref{thm:rh}.
\end{proof}

\begin{remark}\label{rem:content}
Proposition~\ref{prop:content} is the exact analogue, for one object,
of Propositions~\ref{prop:weil-linear} and~\ref{prop:weil-converse}
over $\F_q$: positivity \emph{and} flow compatibility of one and the
same star is the whole content of the axioms, and dropping either half
destroys the implication (\ref{item:ct-pos} and~\ref{item:ct-indef};
for a positive star that fails to commute with the flow see also
Proposition~\ref{prop:R2-star}). The axioms are therefore to be judged
not by the implication, which is formal, but by whether the functor,
cup product, trace, descent, and trace formula can be constructed
independently of the zeros (Section~\ref{sec:open}). The same
collapse happens for any single meromorphic function $G$ with simple
poles at $0$ and $1$ and no other poles, with $\divisor G$ invariant
under $\rho\mapsto1-\rho$ and even central order, for instance a
completed member of the Selberg class that is self-dual (so that
$\divisor G$ is invariant under $\rho\mapsto1-\rho$), has root number
$+1$ (so that the central order is even), and has a simple pole at
$s=1$:
Proposition~\ref{prop:trivial-model}\ref{item:tm-iff}, applied to the
divisor of the entire function $s(s-1)G(s)$, shows that ``a cohomology
for $G$'' of type $(1,0)$ exists if and only if the zeros of $G$ lie on
the line. (For entire $G$ the endpoint blocks of \ref{ax:AW1} would
have to be removed, and for poles of higher order they do not
suffice.) An analytic class supplies no object for such a
cohomology to be a functor of, no cup product with an independent
meaning, no descent, and no geometric side of a trace formula; this is
why the headline class is geometric.
\end{remark}

It is natural to expect that self-duality and a simple pole at $s=1$
suffice for the collapse of Remark~\ref{rem:content}, so that for such
a function a ``cohomology'' of type $(1,0)$ exists if and only if its
zeros lie on the line; in fact the even central order is a separate
condition, and an odd central order excludes every Weil datum with the
axioms other than positivity, whatever the location of the zeros.

\begin{proposition}[{after \cite[2.9, p.~4, and Theorem~2.13, p.~5]%
{Deninger2022NoWeil}}]\label{prop:N-central-order}
Let $\sigma$ be the two-dimensional irreducible representation of the
quaternion group of order eight, realized as the Galois group of an
extension $L/\Q$ with $W(\sigma)=-1$ as in
Theorem~\ref{thm:imp-fs}\ref{item:fs-omar}, let $\Lambda(\sigma,s)$
be its completed Artin $L$-function over $\Q$, and put
$G(s)=\Lambda(s)\Lambda(\sigma,s)$.
\begin{enumerate}[label=(\roman*)]
\item\label{item:Nc-sym} $\divisor G$ is invariant under
  $\rho\mapsto1-\rho$.
\item\label{item:Nc-odd} $\ord_{1/2}G$ is odd.
\item\label{item:Nc-none} No Weil datum of type $(1,0)$ satisfies
  \ref{ax:AW1} for $U=\mathbb 1$, \ref{ax:AW2} for $(G,\mathbb 1)$, and
  \ref{ax:AW3}.
\end{enumerate}
\end{proposition}

\begin{proof}[Proof of \ref{item:Nc-sym}]
The representation $\sigma$ is self-dual (the remark after
Theorem~\ref{thm:imp-fs}), so by
Theorem~\ref{thm:imp-artin}\ref{item:art-fe},
$\ord_{1-\rho}\Lambda(\sigma,s)=\ord_\rho\Lambda(\sigma,s)$; and
$\Lambda(s)=\Lambda(1-s)$ by Theorem~\ref{thm:imp-riemann}.
\end{proof}

\begin{proof}[Proof of \ref{item:Nc-odd}]
$\Lambda(\tfrac12)\ne0$ by Proposition~\ref{prop:F0}. Since $\sigma$
is self-dual with $W(\sigma)=-1$, the meromorphic function
$t\mapsto A_\sigma^{t/2}\Lambda(\sigma,\tfrac12+t)$ is odd (the
consequence of Theorem~\ref{thm:imp-artin}\ref{item:art-fe} recorded
after it), and it is not identically zero because the Euler product
does not vanish for $\Re s>1$. An odd meromorphic function
$ct^m+O(t^{m+1})$ with $c\ne0$ has $(-1)^m=-1$, so
$\ord_{1/2}\Lambda(\sigma,s)$ is odd, and so is $\ord_{1/2}G$.
\end{proof}

\begin{proof}[Proof of \ref{item:Nc-none}]
By \ref{ax:AW2} with $d=1$, $\divisor G=-\operatorname{Sp}_0+
\operatorname{Sp}_1-\operatorname{Sp}_2$, and $\operatorname{Sp}_0=[0]$,
$\operatorname{Sp}_2=[1]$ by \ref{ax:AW1}; so the multiplicity
$m_1(\tfrac12)$ equals $\ord_{1/2}G$, which is odd
by~\ref{item:Nc-odd}. By graded commutativity in \ref{ax:AW1}, the
pairing $B_1(u,v)=\Tr(u\cup v)$ on $V^1$ is alternating, and by
\ref{ax:AW3} it restricts to a perfect pairing of
$V^1_{\Theta\sim1/2}$ with itself. A perfect alternating form exists
only in even dimension, a contradiction.
\end{proof}

Proposition~\ref{prop:trivial-model}\ref{item:tm-iff} therefore does
not apply to $G$; the hypothesis ``even central order'' of
Definition~\ref{def:trivial-model} is a root-number condition, the
same parity that Proposition~\ref{prop:real-coeff} below reads off the
central space.
\subsection{Galois descent}\label{subsec:descent}

For a finite Galois extension $L/K$ with group $G$, each $g\in G$
induces an automorphism of $\Xbar_L$ over $\Xbar_K$, so a contravariant
assignment on $\Schbar_1$ makes $G$ act on $V^\bullet_L$ (by pullback along
$g^{-1}$), commuting with $\Theta$ and $\cup$. Every generalized
eigenspace of $\Theta_1$ on $V^1_L$ is then a finite-dimensional
representation of $G$. The character values of a finite group are
algebraic numbers.

\begin{definition}\label{def:descent}
Let $X\mapsto(H^\bullet(\Xbar),\Theta,\cup,\Tr,*)$ be a contravariant
assignment of Weil data of type $(1,0)$ on a full subcategory of
$\Schbar_1$.
\begin{enumerate}[label=(\alph*)]
\item The assignment satisfies \emph{Galois descent} if, for every
  finite Galois extension $L'/K'$ of fields whose objects lie in the
  subcategory, with group $G'$, pullback along $\Xbar_{L'}\to\Xbar_{K'}$
  induces an isomorphism $H^\bullet(\Xbar_{K'})\cong H^\bullet(\Xbar_{L'})^{G'}$
  compatible with $\Theta$ and $\cup$.
\end{enumerate}
Let moreover $L/K$ be a finite Galois extension of number fields with
group $G$ such that $\Xbar_L$ lies in the subcategory.
\begin{enumerate}[label=(\alph*),resume]
\item For an irreducible representation $\sigma$ of $G$, the space
  $M_\sigma=\operatorname{Hom}_G(V_\sigma,H^1(\Xbar_L))$, with $\Theta$ acting
  by composition, is the \emph{isotypic multiplicity space} of
  $\sigma$.
\item The \emph{central space} is the generalized eigenspace
  $V_{1/2}=H^1(\Xbar_L)_{\Theta\sim1/2}$, a representation of $G$.
\item The \emph{rational orbit} $O(\sigma)$ of an irreducible
  representation $\sigma$ of $G$ is the set of irreducible
  representations $\tau$ of $G$ whose character is
  $\gamma\circ\chi_\sigma$ for some field automorphism $\gamma$ of $\C$,
  and $\chi_{O(\sigma)}=\sum_{\tau\in O(\sigma)}\chi_\tau$.
\end{enumerate}
\end{definition}

Each $M_\sigma$ is a spectral core, and $H^1(\Xbar_L)\cong\bigoplus_\sigma
V_\sigma\otimes M_\sigma$ compatibly with $\Theta$
\cite[p.~4]{Deninger2022NoWeil}. The set $O(\sigma)$ is finite, and
$\chi_{O(\sigma)}$ takes rational values; if $\chi_\sigma$ is
rational-valued, $O(\sigma)=\{\sigma\}$.
The functoriality expectations behind Definition~\ref{def:descent} are
\cite[p.~4]{Deninger2022NoWeil} and
\cite[\S\S2.8--2.10]{Deninger2025BSDDedekind}.

It is natural to expect that functoriality in the field, added to the
axioms on each object, contributes transcendental content beyond zero
location; in fact, without positivity, it contributes Artin's
holomorphy conjecture for the rational-valued characters and the
corresponding isotypic multiplicities, and no information on the
location of zeros.

\begin{proposition}[{Artin; \cite[\S13.1, Theorems~29--30,
pp.~102--103]{Serre1977LinearReps}}]\label{prop:perm-span}
Let $G$ be a finite group and $\sigma$ an irreducible representation of
$G$. Then $\chi_{O(\sigma)}$ is a $\Q$-linear combination of the
permutation characters $\chi_{\mathrm{Ind}_H^G1}$, $H$ running over the
cyclic subgroups of $G$.
\end{proposition}

\begin{proof}
Call $g,g'\in G$ \emph{cyclically equivalent} if $\langle g\rangle$ and
$\langle g'\rangle$ are conjugate. The function $\chi_{O(\sigma)}$ is
constant on cyclic equivalence classes: it is a class function, and if
$g$ has order $n$ and $k$ is prime to $n$, the eigenvalues of
$\tau(g^k)$ are the $k$-th powers of those of $\tau(g)$, so
$\chi_\tau(g^k)=\gamma_k(\chi_\tau(g))$ for an automorphism $\gamma_k$
of $\C$ sending every $n$-th root of unity $\zeta$ to $\zeta^k$; since
$\gamma_k$ permutes $O(\sigma)$, $\chi_{O(\sigma)}(g^k)=
\chi_{O(\sigma)}(g)$. The permutation character
$\chi_{\mathrm{Ind}_H^G1}(g)=\#\{xH\in G/H:x^{-1}gx\in H\}$ is also
constant on these classes, and is nonzero exactly when $\langle
g\rangle$ is conjugate to a subgroup of $H$. Let $H_1,\dots,H_r$
represent the conjugacy classes of cyclic subgroups, ordered by
nondecreasing order, and $g_i$ a generator of $H_i$. The matrix
$(\chi_{\mathrm{Ind}_{H_i}^G1}(g_k))_{i,k}$ has integer entries, and an
entry with $k>i$ vanishes, since then $H_k$ is not conjugate to a
subgroup of $H_i$; the diagonal entries are at least one. So the
matrix is invertible over $\Q$, the $r$ permutation characters span
the $r$-dimensional space of rational-valued functions on the cyclic
equivalence classes, and $\chi_{O(\sigma)}$ lies in this space.
\end{proof}

\begin{proposition}\label{prop:descent}
Let $L/K$ be a finite Galois extension of number fields with group $G$,
and suppose a contravariant assignment on the full subcategory of
$\Schbar_1$ with objects $\Xbar_{L^H}$, $H\subset G$, satisfies Galois
descent and, on each $\Xbar_{L^H}$,
\ref{ax:AW1} and \ref{ax:AW2} for $(\Lambda_{L^H},\mathbb 1)$, and let
$M_\sigma$ and $O(\sigma)$ be as in Definition~\ref{def:descent}. Then
for every irreducible representation $\sigma$ of $G$, with completed
Artin $L$-functions $\Lambda(\tau,s)$ of $L/K$
(Definition~\ref{def:artin}),
\[
\sum_{\tau\in O(\sigma)}\operatorname{Sp}(\Theta\mid M_\tau)=
\divisor\prod_{\tau\in O(\sigma)}\Lambda(\tau,s)
+\delta_{\sigma,1}\bigl([0]+[1]\bigr).
\]
Consequently:
\begin{enumerate}[label=(\roman*)]
\item\label{item:desc-mult} if $\chi_\sigma$ is rational-valued, then
  for every $\rho$ (with $\rho\notin\{0,1\}$ if $\sigma=1$) the
  algebraic multiplicity of $\rho$ on $M_\sigma$ is
  $\ord_\rho\Lambda(\sigma,s)$;
\item\label{item:desc-artin} for $\sigma\ne1$ the product
  $\prod_{\tau\in O(\sigma)}\Lambda(\tau,s)$ is entire; in particular
  $\Lambda(\sigma,s)$ is entire whenever $\chi_\sigma$ is
  rational-valued and $\sigma\ne1$.
\end{enumerate}
\end{proposition}

\begin{proof}
Let $\Gamma\subset G$. By descent for $L/L^\Gamma$ and the decomposition of
$H^1(\Xbar_L)$,
\[
H^1(\Xbar_{L^\Gamma})\cong H^1(\Xbar_L)^\Gamma\cong
\bigoplus_\sigma V_\sigma^\Gamma\otimes M_\sigma,
\]
and $\dim V_\sigma^\Gamma=\langle\sigma,\mathrm{Ind}_\Gamma^G1\rangle$ by
Frobenius reciprocity, so
\[
\operatorname{Sp}(\Theta_1\mid H^1(\Xbar_{L^\Gamma}))=\sum_\sigma\langle\sigma,
\mathrm{Ind}_\Gamma^G1\rangle\operatorname{Sp}(\Theta\mid M_\sigma).
\]
The proof of Theorem~\ref{thm:rh}\ref{item:rh-divisor}, which uses
only \ref{ax:AW1} and \ref{ax:AW2}, gives
$\operatorname{Sp}(\Theta_1\mid H^1(\Xbar_{L^\Gamma}))
=\divisor\Lambda_{L^\Gamma}
+[0]+[1]$.
By \eqref{eq:artin-factor} and $\langle1,\mathrm{Ind}_\Gamma^G1\rangle=1$,
\[
\divisor\Lambda_{L^\Gamma}+[0]+[1]=\sum_\sigma\langle\sigma,
\mathrm{Ind}_\Gamma^G1\rangle\bigl(\divisor\Lambda(\sigma,s)+
\delta_{\sigma,1}([0]+[1])\bigr).
\]
Put $\Delta(\tau)=\operatorname{Sp}(\Theta\mid M_\tau)-\divisor
\Lambda(\tau,s)-\delta_{\tau,1}([0]+[1])$, a divisor, and for a virtual
character $\chi$ put $\Phi(\chi)=\sum_\tau\langle\tau,\chi\rangle
\Delta(\tau)$, a $\Z$-linear function. We have shown
$\Phi(\chi_{\mathrm{Ind}_\Gamma^G1})=0$ for every subgroup $\Gamma$. By
Proposition~\ref{prop:perm-span} there is an integer $n\ge1$ with
$n\chi_{O(\sigma)}$ an integral combination of permutation characters,
so $0=\Phi(n\chi_{O(\sigma)})=n\sum_{\tau\in O(\sigma)}\Delta(\tau)$,
and $\sum_{\tau\in O(\sigma)}\Delta(\tau)=0$ because the group of
divisors is torsion-free. This is the displayed identity (note that
$O(1)=\{1\}$). If $\chi_\sigma$ is rational-valued, $O(\sigma)=
\{\sigma\}$, which gives~\ref{item:desc-mult}. For~\ref{item:desc-artin},
the left side of the displayed identity is effective.
\end{proof}

Proposition~\ref{prop:descent} corrects the expectation that descent
forces every isotypic divisor: the argument controls only sums over
rational orbits, since the objects $\Xbar_{L^H}$ see only permutation
characters, which are rational-valued. With descent,
\ref{ax:AW1}--\ref{ax:AW2} imply the holomorphy of these orbit products
without any positivity; positivity is needed only for the location of
zeros. The proposition also shows that \ref{ax:AW2} cannot be weakened
to an inclusion of the zeros in the spectrum without losing
\ref{item:desc-artin}. Conversely, at the level of one finite Galois
extension of $\Q$, the Riemann hypothesis and Artin's conjecture for its
Artin $L$-functions produce the package with descent.

\begin{definition}\label{def:AL}
For a finite Galois extension $L/\Q$, let $\cA_L$ be the full
subcategory of $\Schbar_1$ whose objects are the $\Xbar_K$ with $K\subset
L$.
\end{definition}

\begin{proposition}[{compare \cite[Theorem~2.15]{Deninger2025BSDDedekind}}]%
\label{prop:class-converse}
Let $L/\Q$ be a finite Galois extension with group $G$, and let
$\Lambda(\sigma,s)$ be the completed Artin $L$-functions over $\Q$ of
the irreducible representations $\sigma$ of $G$. Suppose that for every
irreducible $\sigma\ne1$ the function $\Lambda(\sigma,s)$ is entire and
all its zeros have real part $\tfrac12$, and that all zeros of $\xi$
have real part $\tfrac12$. Then there is a contravariant assignment on
$\cA_L$ of Weil data of type $(1,0)$, satisfying Galois descent, whose
value on each $\Xbar_K$ satisfies \ref{ax:AW1}--\ref{ax:AW4} for
$(\Lambda_K,\mathbb 1)$, with pullbacks commuting with $\Theta$ and
$\cup$.
\end{proposition}

\begin{proof}
For irreducible $\sigma$ put $\Sigma_\sigma=\divisor\Lambda(\sigma,s)$
if $\sigma\ne1$ and $\Sigma_1=\divisor\xi$; by hypothesis these are
effective divisors supported on $\Re s=\tfrac12$. Let $M_\sigma$ be the
algebraic direct sum of lines $\C e$, one for each point of
$\Sigma_\sigma$ counted with multiplicity, with $\Theta e=\rho e$ on
the lines over $\rho$, and write $M_{\sigma,\rho}$ for the span of the
lines over $\rho$. Put $H^1(\Xbar_L)=\bigoplus_\sigma V_\sigma\otimes M_\sigma$
with $G$ acting on the first factors and $\Theta$ on the second;
$H^0(\Xbar_L)=\C\cdot1$ and $H^2(\Xbar_L)=\C\eta$ with trivial action,
$\Theta_0=0$, $\Theta_2=1$, $\Tr(c\eta)=c$. This is a spectral core,
since only finitely many $\sigma$ occur.

\emph{Positive form.} Let $h$ be the orthogonal sum of tensor products
of $G$-invariant inner products on the $V_\sigma$ with the standard
forms on the $M_\sigma$. It is positive, $G$-invariant, and satisfies
\eqref{eq:inf-pol} with weight one, because all eigenvalues lie on
$\Re s=\tfrac12$ and distinct eigenspaces are orthogonal.

\emph{Alternating duality.} By
Theorem~\ref{thm:imp-artin}\ref{item:art-fe} (and $\xi(s)=\xi(1-s)$ for
$\sigma=1$), $\dim M_{\sigma^\vee,1-\rho}=\dim M_{\sigma,\rho}$. The
involution $(\sigma,\rho)\mapsto(\sigma^\vee,1-\rho)$ acts on the set of
summands $V_\sigma\otimes M_{\sigma,\rho}$. For a pair of distinct
summands exchanged by it, choose an isomorphism
$M_{\sigma^\vee,1-\rho}\cong M_{\sigma,\rho}^*$, let $\beta$ be the
resulting $G$-invariant perfect pairing of the two summands (built from
the canonical pairing of $V_\sigma$ with $V_{\sigma^\vee}\cong
V_\sigma^*$), and put $B(x,y)=\beta(x,y)=-B(y,x)$ for $x$ in the first
summand and $y$ in the second. A summand fixed by the involution has
$\sigma\cong\sigma^\vee$ and $\rho=\tfrac12$. By
Theorem~\ref{thm:imp-fs}\ref{item:fs-form}, $V_\sigma$ carries a
nondegenerate $G$-invariant bilinear form $b_\sigma$ that is either
alternating, or symmetric with $\sigma$ defined over $\R$. If it is
alternating, let $c$ be a nondegenerate symmetric form on
$M_{\sigma,1/2}$. If it is symmetric and $\sigma\ne1$, then
$W(\sigma)=1$ by Theorem~\ref{thm:imp-fs}\ref{item:fs-fq}, so
$t\mapsto A_\sigma^{t/2}\Lambda(\sigma,\tfrac12+t)$ is even
(consequence of Theorem~\ref{thm:imp-artin}\ref{item:art-fe}), hence
$\dim M_{\sigma,1/2}=\ord_{1/2}\Lambda(\sigma,s)$ is even, and we let
$c$ be a nondegenerate alternating form on $M_{\sigma,1/2}$ (for
$\sigma=1$ use $\xi(s)=\xi(1-s)$ instead). Put $B=b_\sigma\otimes
c$ on the summand, which is then alternating and $G$-invariant. Let $B$
vanish on all other pairs of summands. Then $B$ is alternating,
$G$-invariant, supported on pairs of eigenvalues with sum one, and
pairs $V^1_{\Theta\sim\rho}$ perfectly with $V^1_{\Theta\sim1-\rho}$.
Put $u\cup v=B(u,v)\eta$ on $H^1(\Xbar_L)$ and the unit products otherwise.

\emph{Star.} Since $B$ is perfect on each pair of eigenspaces with
eigenvalues $\rho$ and $1-\rho=\bar\rho$, there is a unique antilinear
map $*_1$ with $B(u,*_1v)=h(u,v)$ for all $u,v$; it maps the
$\rho$-eigenspace to the $\bar\rho$-eigenspace, so
$\Theta_1*_1=*_1\Theta_1$, it is bijective because $h$ is
nondegenerate, and it is $G$-equivariant because $B$ and $h$ are
$G$-invariant. With $*_0(c)=\bar c\eta$ and $*_2(c\eta)=\bar c$, the
datum on $\Xbar_L$ satisfies \ref{ax:AW1}, \ref{ax:AW3}, and
\ref{ax:AW4}, as in the proof of
Proposition~\ref{prop:trivial-model}.

\emph{The functor.} For $K\subset L$ put
$V^\bullet_K=(V^\bullet_L)^{\mathrm{Gal}(L/K)}$ with the restricted
structure. For an embedding $K\hookrightarrow K'$ of subfields of $L$,
extend it to $g\in G$ and let the pullback be $v\mapsto gv$; it does
not depend on the extension, and it maps
$(V_L)^{\mathrm{Gal}(L/K)}$ into $(V_L)^{\mathrm{Gal}(L/K')}$ because
$g\,\mathrm{Gal}(L/K)g^{-1}=\mathrm{Gal}(L/gK)\supset
\mathrm{Gal}(L/K')$. This is contravariant, commutes with $\Theta$ and
$\cup$, and satisfies descent because invariants of invariants are
invariants. On each $V^1_K$ the form $h$ stays positive, $*_1$ preserves
invariants, and $B(u,*_1u)=h(u,u)>0$ shows that $B$ stays perfect on
each pair of eigenspaces, so \ref{ax:AW1}, \ref{ax:AW3}, and
\ref{ax:AW4} hold on $\Xbar_K$. Finally, with $H=\mathrm{Gal}(L/K)$,
$\operatorname{Sp}_1(V^1_K)=\sum_\sigma\langle\sigma,
\mathrm{Ind}_H^G1\rangle\Sigma_\sigma=\divisor\Lambda_K+[0]+[1]$ by
\eqref{eq:artin-factor}, which is \eqref{eq:aw-div} for $\Lambda_K$ as
in the proof of Theorem~\ref{thm:rh}\ref{item:rh-divisor}.
\end{proof}

It is natural to expect that a cohomology with real coefficients,
complexified, could satisfy the package with descent; in fact Galois
descent turns Deninger's no-go theorem into an inconsistency.

\begin{proposition}\label{prop:real-coeff}
Let $L/\Q$ be a finite Galois extension with group $G$ and $\sigma$ an
irreducible representation of $G$ with rational-valued character, which
is self-dual, not defined over $\R$, and has $W(\sigma)=-1$. Suppose a
contravariant assignment on the full subcategory $\cA_L$ of $\Schbar_1$
(Definition~\ref{def:AL}) satisfies Galois descent and, on each object,
\ref{ax:AW1} and
\ref{ax:AW2}. Then $\sigma$ occurs with odd multiplicity in the central space
$V_{1/2}$ of Definition~\ref{def:descent}; consequently
$V_{1/2}$ is not the complexification of a real representation of $G$.
Such $L$ and $\sigma$ exist.
\end{proposition}

\begin{proof}
The multiplicity of $\sigma$ in $V_{1/2}$ is the dimension of the
generalized $\tfrac12$-eigenspace of $M_\sigma$, which is
$\ord_{1/2}\Lambda(\sigma,s)$ by
Proposition~\ref{prop:descent}\ref{item:desc-mult} with $K=\Q$. Since
$\sigma$ is self-dual and $W(\sigma)=-1$, the function
$t\mapsto A_\sigma^{t/2}\Lambda(\sigma,\tfrac12+t)$ is odd and not
identically zero, so its order at $t=0$ is odd. If $V_{1/2}$ were the
complexification of a real representation, $\sigma$ would occur in it
with even multiplicity by Theorem~\ref{thm:imp-fs}\ref{item:fs-even}.
The quaternion extension of
Theorem~\ref{thm:imp-fs}\ref{item:fs-omar} with its two-dimensional
representation satisfies the hypotheses (see the remark after that
theorem).
\end{proof}

This is the argument of Deninger's no-go theorem
(Theorem~\ref{thm:imp-deninger-nogo}), whose hypothesis~2.12
\cite{Deninger2022NoWeil} is here supplied by descent for the one
rational character that the argument needs. The obstruction needs the
class to contain the quaternion fields; for the single object $\Xbar_\Q$
it says nothing (compare Remark~\ref{rem:R5-import}).

\begin{proposition}[{after \cite[2.7--2.8, pp.~3--4]{Deninger2022NoWeil}}]%
\label{prop:root-parity}
Let $L/K$ be a finite Galois extension of number fields with group $G$,
and suppose that a contravariant assignment on the full subcategory
of $\Schbar_1$ with objects $\Xbar_{L^H}$, $H\subset G$, satisfies Galois
descent and, on each
object, \ref{ax:AW1} and \ref{ax:AW2}, and that the star $*_1$ on
$V^1_L$ commutes with $G$ and with $\Theta_1$ on the central space
$V_{1/2}$ of Definition~\ref{def:descent} and satisfies $*_1^2=-1$
there. Let
$\sigma$ be an irreducible representation of $G$ defined over $\R$.
Then $\sigma$ occurs in $V_{1/2}$ with even multiplicity; if moreover
$\chi_\sigma$ is rational-valued, $\ord_{1/2}L(\sigma,s)$ is even.
\end{proposition}

\begin{proof}
The space $V_{1/2}$ is $G$-stable and $*_1$-stable (an antilinear map
commuting with $\Theta_1$ maps the $\rho$-generalized eigenspace to the
$\bar\rho$ one). Write $V_\sigma=V_0\otimes_\R\C$ and let $j_\sigma$ be
the $G$-equivariant complex conjugation of this real form, so
$j_\sigma^2=1$. The $\sigma$-isotypic part of $V_{1/2}$ is
$V_\sigma\otimes N$ with $N$ the generalized $\tfrac12$-eigenspace of
$M_\sigma$, and $*_1$ maps it to the $\bar\sigma$-isotypic part, which
is the same part since $\bar\sigma\cong\sigma$. Let $\kappa$ be complex
conjugation on $N$ in some basis. Then $(j_\sigma\otimes\kappa)^{-1}*_1$
is linear and $G$-equivariant on $V_\sigma\otimes N$, hence equals
$1\otimes A$ for a linear $A$ by Schur's lemma, and
$*_1=j_\sigma\otimes j_N$ with $j_N=\kappa A$ antilinear. Then
$-1=*_1^2=j_\sigma^2\otimes j_N^2=1\otimes j_N^2$, so $j_N^2=-1$, a
quaternionic structure on $N$, and $\dim N$ is even. If $\chi_\sigma$
is rational-valued, $\dim N=\ord_{1/2}\Lambda(\sigma,s)=
\ord_{1/2}L(\sigma,s)$ by Proposition~\ref{prop:descent}%
\ref{item:desc-mult}, the gamma factors being holomorphic and nonzero
at $\tfrac12$.
\end{proof}

For $\sigma$ self-dual and not defined over $\R$ the same argument,
with an antilinear $j_\sigma$ of square $-1$, gives a real structure on
$N$ and no parity constraint. The prediction is consistent with the
functional equation: representations defined over $\R$ have
$W(\sigma)=1$ (Theorem~\ref{thm:imp-fs}\ref{item:fs-fq}) and hence even
central order, while the others may have $W(\sigma)=-1$ and odd central
order (Proposition~\ref{prop:real-coeff}). Root numbers therefore enter
the package through the parity of central orders, never through the
location of zeros.
\subsection{The finite-field instance}\label{subsec:thmB}

This subsection places Weil's theorem in the frame of
Definition~\ref{def:aw}: for a smooth projective geometrically connected
variety over $\F_q$, the logarithmic suspension of $\ell$-adic
cohomology is a Weil datum of type $(d,0)$ which satisfies
\ref{ax:AW1}--\ref{ax:AW3} by Grothendieck's formalism, whatever its
stars, and whose positivity axiom holds exactly when Frobenius is
semisimple and pure. The determinant identity holds with normalization
term $1$. We first record the algebra needed to suspend cup products
and stars.

\begin{proposition}\label{prop:susp-product}
Let $X$, $\ell$, $H^\bullet$, $\mathrm{F}$, and $\operatorname{tr}$ be as
in Theorem~\ref{thm:imp-weil-coh}, let $\iota:\overline\Q_\ell\to\C$ be
a field isomorphism, and write $H^j$ also for
$H^j\otimes_{\Q_\ell,\iota}\C$, with the
induced $\mathrm{F}$, cup product, and trace.
\begin{enumerate}[label=(\roman*)]
\item\label{item:sp-mult} The cup product maps
  $H^i_{(\alpha)}\times H^k_{(\beta)}$ into $H^{i+k}_{(\alpha\beta)}$.
\item\label{item:sp-der} On each $H^j_{(\alpha)}$ let $N=\log(\alpha^{-1}
  \mathrm{F})$ as in Definition~\ref{def:log-susp}. Then
  $N(u\cup v)=Nu\cup v+u\cup Nv$ for all $u,v$.
\item\label{item:sp-star} If $h_j$ is a positive Frobenius polarization
  of weight $j$ on $(H^j,\mathrm{F})$, there is a unique antilinear map
  $*_j:H^j\to H^{2d-j}$ with $\operatorname{tr}(u\cup*_jv)=h_j(u,v)$
  for all $u,v$, and it maps $H^j_{(q^\rho)}$ into
  $H^{2d-j}_{(q^{d-\rho})}$ for every $\rho$ with $q^\rho$ an
  eigenvalue.
\end{enumerate}
\end{proposition}

\begin{proof}[Proof of \ref{item:sp-mult}]
By Theorem~\ref{thm:imp-weil-coh}, $\mathrm{F}$ is an automorphism of
the graded algebra $H^\bullet$. If $(\mathrm{F}\otimes\mathrm{F}-
\alpha\beta)^n(u\otimes v)=0$, then $(\mathrm{F}-\alpha\beta)^n(u\cup
v)=0$; and $\mathrm{F}\otimes\mathrm{F}$ has the single eigenvalue
$\alpha\beta$ on $H_{(\alpha)}\otimes H_{(\beta)}$.
\end{proof}

\begin{proof}[Proof of \ref{item:sp-der}]
By~\ref{item:sp-mult} the semisimple part $\mathrm{F}_s$, equal to
$\alpha$ on $H_{(\alpha)}$, is multiplicative, hence so is the
unipotent part $\mathrm{F}_u=\mathrm{F}_s^{-1}\mathrm{F}$, and
$N=\log\mathrm{F}_u$ on each $H_{(\alpha)}$. For integers $t\ge0$,
$\exp(tN)=\mathrm{F}_u^t$ is multiplicative; both sides of
$\exp(tN)(u\cup v)=\exp(tN)u\cup\exp(tN)v$ are polynomials in $t$
because $N$ is nilpotent, so the identity holds for all real $t$, and
differentiating at $t=0$ gives the claim.
\end{proof}

\begin{proof}[Proof of \ref{item:sp-star}]
The Poincar\'e pairing is perfect, so $*_j$ exists, is unique, and is
antilinear. By Proposition~\ref{prop:weil-linear}, $\mathrm{F}$ is
semisimple on $H^j$ with eigenvalues of modulus $q^{j/2}$. From
$h_j(\mathrm{F}u,\mathrm{F}v)=q^jh_j(u,v)$ and
$\operatorname{tr}(\mathrm{F}u\cup\mathrm{F}w)=q^d\operatorname{tr}
(u\cup w)$ we get $\operatorname{tr}(\mathrm{F}u\cup*_j\mathrm{F}v)=
q^{j-d}\operatorname{tr}(\mathrm{F}u\cup\mathrm{F}*_jv)$ for all $u$,
so $*_j\mathrm{F}=q^{j-d}\mathrm{F}*_j$. For an eigenvector $v$ with
eigenvalue $\alpha=q^\rho$, $\Re\rho=j/2$, this gives
$\mathrm{F}(*_jv)=q^{d-j}\bar\alpha\,*_jv=q^{d-\rho}\,*_jv$, because
$q^{\bar\rho}=\bar\alpha$ and $\bar\rho=j-\rho$.
\end{proof}

\begin{definition}\label{def:susp-datum}
Let $X$, $\ell$, $\iota$, $H^\bullet$, $\mathrm{F}$, and
$\operatorname{tr}$ be as in Proposition~\ref{prop:susp-product}, and
let $(\widehat H^j,\Theta_j)$ be the logarithmic suspension of
$(H^j,\mathrm{F})$ (Definition~\ref{def:log-susp}).
\begin{enumerate}[label=(\alph*)]
\item The \emph{suspended datum} of $(X,\ell,\iota)$ is the Weil datum
  $(\widehat H^\bullet,\Theta,\cup,\Tr,0)$ of type $(d,0)$ with zero
  stars, where for $x$ the copy of $u\in H^i_{(q^\rho)}$ in
  $\widehat H^i_\rho$ and $y$ the copy of $v\in H^k_{(q^{\rho'})}$ in
  $\widehat H^k_{\rho'}$, the product $x\cup y$ is the copy of $u\cup
  v$ in $\widehat H^{i+k}_{\rho+\rho'}$ if $q^{\rho+\rho'}$ is an
  eigenvalue of $\mathrm{F}$ on $H^{i+k}$ and $0$ otherwise, and
  $\Tr(x)=\operatorname{tr}(u)$ if $x$ is the copy of $u\in H^{2d}$ in
  the summand $\widehat H^{2d}_d$ and $\Tr(x)=0$ on the other summands.
\item If positive Frobenius polarizations $h_j$ of weight $j$ on
  $(H^j,\mathrm{F})$ are given for all $j$, the \emph{suspended stars}
  are the antilinear maps $\widehat*_j$ sending the copy of $v\in
  H^j_{(q^\rho)}$ in $\widehat H^j_\rho$ to the copy of $*_jv$ in
  $\widehat H^{2d-j}_{d-\rho}$, with $*_j$ as in
  Proposition~\ref{prop:susp-product}\ref{item:sp-star}.
\end{enumerate}
\end{definition}

By Proposition~\ref{prop:susp-product}, the product lands in the
indicated summand, the suspended stars are well defined, and
$\widehat H^{2d}_d$ is a summand because $\mathrm{F}=q^d$ on $H^{2d}$.
Axioms \ref{ax:AW1}--\ref{ax:AW3} and AW-det do not involve the stars.
The suspended datum is Deninger's cohomology
$H^\bullet(X/L)=D(H^\bullet(X_{\overline\F_q},\Q_\ell)\otimes\C)$,
which inherits products and duality from $\ell$-adic cohomology because
$D$ is an exact tensor functor
\cite[\S4, pp.~16--17]{Deninger1994Motivic}.

The regularized determinants of suspended Frobenius are computed first.
Recall the Hurwitz zeta function $\zeta_H(w,a)=\sum_{n\ge0}(n+a)^{-w}$
for $a\in\C\setminus(-\infty,0]$ and $\Re w>1$, with principal powers.

\begin{theorem}[{Lerch; \cite[proof of (2.7), p.~9]%
{Deninger1994Motivic}}]
\label{thm:imp-hurwitz}
For $a\in\C\setminus(-\infty,0]$, $\zeta_H(w,a)$ continues
meromorphically in $w$ to $\C$ with a single simple pole at $w=1$, and
$\zeta_H(0,a)=\tfrac12-a$ and $\exp(\partial_w\zeta_H(0,a))=
\Gamma(a)/\sqrt{2\pi}$.
\end{theorem}

The source states these classical facts without proof.

It is natural to expect that the identity
$\detreg((s-\Theta)/2\pi)=1-\alpha q^{-s}$ holds only up to an
unspecified normalization and that the factor $2\pi$ must be tuned; in
fact, for the principal cut, the constant is exactly $1$, and passing
to the reflected cut changes the determinant by a factor $e^{a+bs}$
with nonzero $b$.

\begin{theorem}[Deninger; Hesselholt]\label{thm:ff-det}
Let $q>1$ be real, $V$ a finite-dimensional complex vector space,
$\mathrm{F}:V\to V$ invertible, and $(\widehat V,\Theta)$ its
logarithmic suspension. Then $(\widehat V,\Theta)$ is zeta-regularizable
for its principal cut (Definition~\ref{def:principal-cut}), and
\[
\detreg\bigl((s-\Theta)/2\pi\bigr)=\det(1-\mathrm{F}q^{-s}\mid V)
=\prod_\alpha(1-\alpha q^{-s})^{\dim V_{(\alpha)}}.
\]
Moreover $Z(0;s)=0$ for every $s$ in the domain of the cut, so the same
identity holds with $2\pi$ replaced by any $\delta>0$.
\end{theorem}

\begin{proof}
The spectrum is the union of the progressions $\Sigma_q(\alpha)$, with
multiplicity $\dim V_{(\alpha)}$ on each point of $\Sigma_q(\alpha)$
(Proposition~\ref{prop:suspension}\ref{item:susp-zeros}), and the
spectral zeta function is the corresponding weighted sum over $\alpha$;
it suffices to treat one progression with multiplicity one. Fix a
logarithm $\lambda$ of $\alpha$, put $\rho_n=(\lambda+2\pi in)/\log q$
for $n\in\Z$, and for $s\in\C$ put $\sigma=s\log q-\lambda$ and
$z=i\sigma/2\pi$. Then
\[
t_n=\frac{s-\rho_n}{2\pi}=\frac{-i}{\log q}(n+z),\qquad
t_{-m}=\frac{i}{\log q}(m-z)\quad(m\ge1).
\]
Let $R=\{z'\in\C:\Im z'>0\}\cup\{z'\in\C:0<\Re z'<1\}$. Replacing
$\lambda$ by $\lambda+2\pi in_0$ replaces $z$ by $z+n_0$ and only
reindexes the spectrum. The point $s$ lies in the domain
$U_{\mathrm{pr}}$ of the principal cut exactly when no $t_n$ lies in
$(-\infty,0]$, that is,
unless $\Re z\in\Z$ and $\Im z\le0$; so for $s\in U_{\mathrm{pr}}$ there
is $n_0$ with $z+n_0\in R$, and such $s$ form an open set. Assume
$z\in R$. For $n\ge0$ the principal argument of $n+z$ lies in
$(-\pi/2,\pi]$, so the principal argument of $t_n$ is
$\arg(n+z)-\pi/2$ and $t_n^{-w}=(\log q)^we^{i\pi w/2}(n+z)^{-w}$; for
$m\ge1$ the principal argument of $m-z$ lies in $(-\pi,\pi/2]$, so
$t_{-m}^{-w}=(\log q)^we^{-i\pi w/2}(m-z)^{-w}$. Moreover $z$ and
$1-z$ lie outside $(-\infty,0]$. Hence, for $\Re w>1$,
\[
Z(w;s)=\sum_{n\in\Z}t_n^{-w}=(\log q)^w\Bigl[e^{i\pi w/2}\zeta_H(w,z)
+e^{-i\pi w/2}\zeta_H(w,1-z)\Bigr],
\]
which by Theorem~\ref{thm:imp-hurwitz} continues meromorphically in
$w$, is regular at $w=0$, depends holomorphically on $s$ near $w=0$,
and has $Z(0;s)=(\tfrac12-z)+(\tfrac12-(1-z))=0$. Differentiating at
$w=0$, the factor $(\log q)^w$ contributes $\log\log q\cdot Z(0;s)=0$,
and
\[
\partial_wZ(0;s)=\tfrac{i\pi}2\bigl[(\tfrac12-z)-(z-\tfrac12)\bigr]
+\partial_w\zeta_H(0,z)+\partial_w\zeta_H(0,1-z),
\]
so, using $\Gamma(z)\Gamma(1-z)=\pi/\sin\pi z$,
\[
e^{-\partial_wZ(0;s)}=e^{-i\pi/2}e^{i\pi z}\frac{2\pi}{\Gamma(z)
\Gamma(1-z)}=-ie^{i\pi z}\cdot2\sin\pi z=1-e^{2\pi iz}=1-e^{-\sigma}
=1-\alpha q^{-s}.
\]
The right side does not depend on $n_0$, so
$e^{-\partial_wZ(0;s)}=1-\alpha q^{-s}$ on all of $U_{\mathrm{pr}}$,
and it continues to that entire function. With $\delta$ in place of
$2\pi$, the factor $(\log q)^w$ becomes $(\delta\log q/2\pi)^w$, which
again contributes $0$ because $Z(0;s)=0$.
\end{proof}

Theorem~\ref{thm:ff-det} is \cite[(2.7), pp.~8--9, and (2.8), p.~10]%
{Deninger1994Motivic} and \cite[Proposition~7.1; p.~21 of the arXiv version]%
{Hesselholt2016THHZeta}; the proof above follows the former.

\begin{proposition}[{\cite[(2.10), p.~11, and~(4.1), p.~16]%
{Deninger1994Motivic}}]\label{prop:ff-cut}
Let $q>1$, $\alpha\in\C^\times$, and let $(\widehat V,\Theta)$ be the
logarithmic suspension of $(\C,\alpha)$. Then $(\widehat V,\Theta)$ is
zeta-regularizable for its reflected cut
(Definition~\ref{def:principal-cut}), and
\[
\detreg\bigl((s-\Theta)/2\pi\bigr)=1-\alpha^{-1}q^{s}=
-\alpha^{-1}q^{s}\,(1-\alpha q^{-s}).
\]
\end{proposition}

\begin{proof}
For $w\notin[0,\infty)$ the argument of $w$ in $[0,2\pi)$ is
$\operatorname{Arg}(-w)+\pi$. Hence, for $s\in U_{\mathrm{rf}}$,
$((s-\rho)/2\pi)^{-w}=e^{-i\pi w}((\rho-s)/2\pi)^{-w}$ with the
principal power on the right. The points $-\rho$, $\rho\in
\Sigma_q(\alpha)$, form $\Sigma_q(\alpha^{-1})$, and
$(\rho-s)/2\pi=((-s)-(-\rho))/2\pi$, so for $s\in U_{\mathrm{rf}}$ the
spectral zeta function of the reflected cut is
$Z'(w;s)=e^{-i\pi w}Z^-(w;-s)$,
where $Z^-$ is the principal spectral zeta function of the suspension of
$(\C,\alpha^{-1})$, and $-s$ lies in its principal domain. By
Theorem~\ref{thm:ff-det}, $Z^-(0;-s)=0$ and
$e^{-\partial_wZ^-(0;-s)}=1-\alpha^{-1}q^{s}$, so
$-\partial_wZ'(0;s)=i\pi Z^-(0;-s)-\partial_wZ^-(0;-s)$ gives the
claim.
\end{proof}

The two answers have the same divisor and differ by the factor
$-\alpha^{-1}q^s=e^{a+bs}$ with $b=\log q\ne0$: the normalization term of
Definition~\ref{def:aw-det} is a genuine convention and is recorded as
data.

It is natural to expect that the suspension can be done with a single
logarithm of each eigenvalue; in fact a single branch destroys the
divisor.

\begin{proposition}\label{prop:one-branch}
Let $q>1$, $\alpha\in\C^\times$, and $\rho_0\in\Sigma_q(\alpha)$, and let
$(V,\Theta)=(\C,\rho_0)$. Then $(V,\Theta)$ is a spectral core whose
spectral divisor $[\rho_0]$ differs from the divisor
$\sum_{\rho\in\Sigma_q(\alpha)}[\rho]$ of $1-\alpha q^{-s}$, and
$\detreg((s-\Theta)/2\pi)=(s-\rho_0)/2\pi$ for every spectral cut, a
polynomial in $s$ and not a function of $q^{-s}$.
\end{proposition}

\begin{proof}
$\Sigma_q(\alpha)$ is infinite, and the determinant of a
one-dimensional core was computed after Definition~\ref{def:regdet}.
\end{proof}

The result also depends on the choice of $\rho_0$, which is defined only
modulo $2\pi i/\log q$; the whole progression is needed, and with it
Theorem~\ref{thm:ff-det} recovers the Euler factor exactly.

It is natural to expect that the finite-field zeta function satisfies
the axioms verbatim, with $V^0=\C$ and $V^{2d}=\C(-d)$; in fact the
unit object must be the cohomology of the base, which over $\F_q$ is
periodic.

\begin{proposition}[{Deninger \cite[pp.~167--168 and
Proposition~3.1]{Deninger1998Analogies}}]\label{prop:unit}
Let $X$ be a smooth projective geometrically connected variety over
$\F_q$, $\ell\nmid q$ a prime, $\iota:\overline\Q_\ell\to\C$ a field
isomorphism, and $\widehat H^0$ the degree-zero part of the suspended
datum (Definition~\ref{def:susp-datum}). Then $\widehat H^0$, with its
product, is isomorphic to the $q$-periodic unit $\mathbb 1_q$ of
Definition~\ref{def:unit}; its spectrum is $(2\pi i/\log q)\Z$; its
principal-cut determinant is $1-q^{-s}$, not $s/2\pi$; and it is not
isomorphic to the arithmetic unit $\mathbb 1$.
\end{proposition}

\begin{proof}
By Theorem~\ref{thm:imp-weil-coh}, $H^0=\C\cdot1$ with $\mathrm{F}=1$,
so $\widehat H^0$ has one line $\C1_\rho$ for each $\rho\in\Sigma_q(1)=
(2\pi i/\log q)\Z$, with $\Theta=\rho$ there, and by
Definition~\ref{def:susp-datum} $1_\rho\cup1_{\rho'}=1_{\rho+\rho'}$.
Sending $1_{2\pi in/\log q}$ to $v^n$ is the isomorphism. The
determinant is Theorem~\ref{thm:ff-det} with $\alpha=1$. The spectrum
of $\mathbb 1$ is $\{0\}$.
\end{proof}

\begin{theorem}[{Theorem~B; (i)--(iii) Deninger
\cite[\S4, Proposition~(4.1), and pp.~16--17]{Deninger1994Motivic}}]%
\label{thm:B}
Let $X$ be a smooth projective geometrically connected variety of
dimension $d$ over $\F_q$, $\ell\nmid q$ a prime, $\iota:\overline
\Q_\ell\to\C$ a field isomorphism, $H^j$ short for
$H^j(X_{\overline\F_q},\Q_\ell)
\otimes_{\Q_\ell,\iota}\C$ with geometric Frobenius $\mathrm{F}$, and
$(\widehat H^\bullet,\Theta,\cup,\Tr,0)$ the suspended datum
(Definition~\ref{def:susp-datum}). Then:
\begin{enumerate}[label=(\roman*)]
\item\label{item:B-aw1} it satisfies \ref{ax:AW1} for
  $(\zeta_X,\mathbb 1_q)$;
\item\label{item:B-det} it satisfies AW-det for $\zeta_X$ with the
  principal cuts and normalization term $1$, with
  $D_j(s)=\det(1-\mathrm{F}q^{-s}\mid H^j)$, and hence \ref{ax:AW2};
\item\label{item:B-aw3} it satisfies \ref{ax:AW3} with $w=0$;
\item\label{item:B-aw4} for each $j$, the following are equivalent: the
  datum satisfies AW4-V in degree $j$; it satisfies AW4-E in degree
  $j$; $\mathrm{F}$ is semisimple on $H^j$ and $|\alpha|=q^{j/2}$ for
  every eigenvalue $\alpha$. If they hold for every $j$, the datum
  with the suspended stars of Definition~\ref{def:susp-datum} satisfies
  \ref{ax:AW4};
\item\label{item:B-purity} if the datum satisfies AW4-V in degree $j$,
  every eigenvalue $\alpha$ of $\mathrm{F}$ on $H^j$ satisfies
  $|\alpha|=q^{j/2}$;
\item\label{item:B-curves} if $X$ is a curve, positive Frobenius
  polarizations exist in every degree, and the datum with the
  suspended stars satisfies \ref{ax:AW1}--\ref{ax:AW4}.
\end{enumerate}
\end{theorem}

\begin{proof}[Proof of \ref{item:B-aw1}]
By Proposition~\ref{prop:susp-product}\ref{item:sp-mult} and
Theorem~\ref{thm:imp-weil-coh}, the product is associative and
graded-commutative with unit the copy of $1$ in $\widehat H^0_0$. Since
$\Theta=\rho+N/\log q$ on $\widehat H_\rho$, indices add under the
product, and $N$ is a derivation
(Proposition~\ref{prop:susp-product}\ref{item:sp-der}), $\Theta$ is a
derivation. By Theorem~\ref{thm:imp-weil-coh}, $H^0=\C$ with
$\mathrm{F}=1$, so $\widehat H^0\cong\mathbb 1_q$ by
Proposition~\ref{prop:unit}; and $H^{2d}$ is a line with
$\mathrm{F}=q^d$, so $\widehat H^{2d}$ has one line over each point of
$\Sigma_q(q^d)=d+(2\pi i/\log q)\Z$, which is $\mathbb 1_q(-d)$.
\end{proof}

\begin{proof}[Proof of \ref{item:B-det}]
By Theorem~\ref{thm:ff-det}, $D_j(s)=\det(1-\mathrm{F}q^{-s}\mid H^j)$,
whose divisor is $\operatorname{Sp}_j$ by
Proposition~\ref{prop:suspension}\ref{item:susp-zeros}. By
Theorem~\ref{thm:imp-deligne},
$\zeta_X(s)=\prod_j\det(1-\mathrm{F}q^{-s}\mid H^j)^{(-1)^{j+1}}$,
which is \eqref{eq:aw-det} with $a=b=0$.
\end{proof}

\begin{proof}[Proof of \ref{item:B-aw3}]
$\Tr$ vanishes off $\widehat H^{2d}_d$, where $\Theta=d$, so
$\Tr(\Theta x)=d\Tr(x)$. For $x\in\widehat H^j_\rho$ and
$y\in\widehat H^{2d-j}_{\rho'}$, $x\cup y$ lies in the summand
$\rho+\rho'$, so $\Tr(x\cup y)$ vanishes unless $\rho'=d-\rho$, and
then equals $\operatorname{tr}(u\cup v)$. The Poincar\'e pairing
satisfies $\operatorname{tr}(\mathrm{F}u\cup\mathrm{F}v)=
q^d\operatorname{tr}(u\cup v)$, so it vanishes on the image of
$\mathrm{F}\otimes\mathrm{F}-q^d$, which contains
$H^j_{(\alpha)}\otimes H^{2d-j}_{(\beta)}$ when $\alpha\beta\ne q^d$.
Since the pairing is perfect (Theorem~\ref{thm:imp-weil-coh}), its
restriction to $H^j_{(q^\rho)}\times H^{2d-j}_{(q^{d-\rho})}$ is
perfect, which by Proposition~\ref{prop:suspension}\ref{item:susp-core}
is the pairing of $\widehat H^j_{\Theta\sim\rho}$ with
$\widehat H^{2d-j}_{\Theta\sim d-\rho}$.
\end{proof}

\begin{proof}[Proof of \ref{item:B-aw4} and \ref{item:B-purity}]
If AW4-V holds in degree $j$, Proposition~\ref{prop:susp-converse}
gives a positive Frobenius polarization of weight $j$ on
$(H^j,\mathrm{F})$, so $\mathrm{F}$ is semisimple and pure by
Proposition~\ref{prop:weil-linear}; this also
proves~\ref{item:B-purity}. If $\mathrm{F}$ is semisimple and pure,
Proposition~\ref{prop:weil-converse} gives a positive Frobenius
polarization and Proposition~\ref{prop:suspension}\ref{item:susp-pol}
gives AW4-V. AW4-V and AW4-E are equivalent by
Proposition~\ref{prop:aw4-forms}. Suppose now that positive Frobenius
polarizations $h_j$ exist for all $j$, and let $*_j$ be as in
Proposition~\ref{prop:susp-product}\ref{item:sp-star}, so that
$\widehat*_j$ maps $\widehat H^j_\rho$ to $\widehat H^{2d-j}_{d-\rho}$
and every eigenvalue $\rho$ has $\Re\rho=j/2$. On $\widehat H^j_\rho$ we have
$\Theta_{2d-j}\widehat*_j=
(d-\rho)\widehat*_j$ and $\widehat*_j\Theta_j=\bar\rho\,\widehat*_j$,
and $d-\rho=\bar\rho+(d-j)$: this is the star law. Finally, for $x$ the
copy of $u$ over $\rho$ and $y$ the copy of $v$ over $\rho'$,
$x\cup\widehat*_jy$ lies over $\rho+d-\rho'$, so $\Tr(x\cup
\widehat*_jy)$ is $h_j(u,v)$ if $\rho=\rho'$ and $0$ otherwise: the
star form is the suspended form $\widehat h_j$, which is positive
definite.
\end{proof}

\begin{proof}[Proof of \ref{item:B-curves}]
For $j=0,2$, $\mathrm{F}$ is $1$ and $q$ on a line. For $j=1$, if $C$
has genus $0$ then $H^1=0$ by the identification recalled after
Proposition~\ref{prop:rosati-EF}. Otherwise, by
Proposition~\ref{prop:rosati-EF}, $(V_\ell J\otimes_\iota\C,\pi)$
satisfies AW4-E$_{\mathrm F}$ of weight one; under the identification
of $H^1$ with the dual of $V_\ell J\otimes_\iota\C$, with
$\mathrm{F}$ the transpose of $\pi$ (the imported identification
recalled after Proposition~\ref{prop:rosati-EF}),
Proposition~\ref{prop:aw4-EF}\ref{item:EF-dual} shows that
$(H^1,\mathrm{F})$ satisfies AW4-E$_{\mathrm F}$ of weight one. By
Proposition~\ref{prop:aw4-EF}\ref{item:EF-pure}, $\mathrm{F}$ is
semisimple on $H^1$ with eigenvalues of modulus $q^{1/2}$, and by
Proposition~\ref{prop:aw4-EF}\ref{item:EF-add} the datum satisfies
AW4-E in degree one directly. Positive Frobenius polarizations then
exist in every degree by Proposition~\ref{prop:weil-converse}; apply
\ref{item:B-aw4} and~\ref{item:B-aw1}--\ref{item:B-aw3}.
\end{proof}

Thus for curves the suspended datum satisfies all four axioms, and the
finite-field Riemann hypothesis is the conclusion of the purity argument
applied to it. Weil's positivity of the Rosati involution is an instance
of the multiplicative algebra form AW4-E$_{\mathrm F}$
(Proposition~\ref{prop:rosati-EF}); the additive algebra form AW4-E on
the suspension follows from it inside the same algebra
(Proposition~\ref{prop:aw4-EF}), and the star form \ref{ax:AW4} from
AW4-E through~\ref{item:B-aw4}, with the non-canonical polarizations of
Proposition~\ref{prop:weil-converse}. The only import in this chain is
the Frobenius compatibility of the identification of $H^1$ with the
dual of the Tate module. For a
general $X$ the statement is weaker: \ref{ax:AW1}--\ref{ax:AW3} are
theorems, the conclusion $|\alpha|=q^{j/2}$ is Deligne's theorem, but
the hypothesis AW4 in degree $j$ is \emph{equivalent} to Deligne's
theorem together with the semisimplicity of Frobenius on $H^j$, which is
open in general (Remark~\ref{rem:ff-status}).

It is natural to expect that, since Deligne proved $|\alpha|=q^{j/2}$,
the finite-field instance of the positivity axiom is a theorem for
every smooth projective variety; in fact it is known only where
Frobenius is known to be semisimple.

\begin{remark}\label{rem:B-status}
By Theorem~\ref{thm:B}\ref{item:B-aw4} and Remark~\ref{rem:ff-status},
the status of the axioms for a smooth projective $X/\F_q$ is: AW1--AW3
and AW-det hold for all $X$; AW4 in degree $j$ holds if and only if
Frobenius is semisimple on $H^j$ (purity being Deligne's theorem), which
is known for curves and abelian varieties and implied by the standard
conjectures \cite[p.~10]{Kleiman1994Standard},
\cite[Aside~3.18, p.~42]{Milne2016RHFF}. Deligne's proof of purity does not
pass through a positive form. What Weil proved for curves, and what the
standard conjectures would give in general, is positivity of an
involution on an algebra of correspondences
(Proposition~\ref{prop:algebra-form}), never a positive form on $H^j$
itself: the Weil form of the standard conjectures is positive definite
on some summands of $H^j$ and negative definite on others
\cite[Definition~3.6, p.~39, Theorem~3.15, and fn.~36, p.~42]%
{Milne2016RHFF}. The arithmetic package
therefore mirrors the conjectural route of Serre, Grothendieck, and
Kleiman, not the proved route of Deligne.
\end{remark}

\begin{remark}\label{rem:B-iota}
Positivity is imposed on the complex side, for one field isomorphism
$\iota$ at a time, and gives $|\iota(\alpha)|=q^{j/2}$ for that $\iota$.
Weil's Riemann hypothesis asks this for every complex conjugate of every
eigenvalue; it follows from the statement for one $\iota$ once the
characteristic polynomial is known to have rational coefficients
(Theorem~\ref{thm:imp-deligne}), and the Rosati argument proves it for
every embedding at once (Proposition~\ref{prop:weil-curves}). Compared
with the finite-field formalism, the arithmetic axioms lose rationality
and integrality in $q^{-s}$, finite dimensionality, and Frobenius as an
operator (only the flow $e^{t\Theta}$ remains, with
$\mathrm{F}=e^{\Theta\log q}$), and they gain the archimedean places,
whose gamma factors, together with $\zeta_K$, produce the poles of
$\Lambda_K$ at $0$ and $1$ that the endpoint blocks account for in the
proof of Theorem~\ref{thm:rh}. The integer-twist blindness of
Proposition~\ref{prop:R3} is visible on the suspension: $\ker(\Theta-n)$
on $\widehat H^{2n}$ is the $\mathrm{F}=q^n$ eigenspace, so cones of
integer twists never see $\Im\rho\ne0$, although the determinants of
Theorem~\ref{thm:ff-det} contain the whole spectrum.
\end{remark}
\subsection{Higher-dimensional arithmetic schemes}\label{subsec:schemes}

For a regular projective flat arithmetic scheme $X$ of dimension $d$
the Hasse--Weil zeta function
$\zeta_X(s)=\prod_x(1-N(x)^{-s})^{-1}$, the product over the closed
points $x$ with residue field of order $N(x)$, converges for $\Re s>d$;
its completion $\widehat\zeta_X$ by Serre's gamma factors, its
meromorphic continuation, and its functional equation are conjectural in
general \cite[Conjecture~1, p.~2]{Deninger2010HilbertPolya}. A Weil
datum with \ref{ax:AW2} for such an object presupposes the
continuation, just as the existence of a Weil cohomology over $\F_q$
yields rationality.

It is natural to expect that all nontrivial zeros of the zeta function
of a higher-dimensional arithmetic scheme lie on one line; in fact
already for $\mathbb{P}^1_\Z$ they lie in two disjoint strips.

\begin{proposition}[{compare \cite[Conjecture~1(c), p.~2]%
{Deninger2010HilbertPolya}}]\label{prop:P1Z}
Let $X=\mathbb{P}^1_\Z$, so $d=2$. Then $\zeta_X(s)=\zeta(s)\zeta(s-1)$.
The zeros of $\zeta_X$ in $0<\Re s<2$ are exactly the points of
$Z\cup(Z+1)$, where $Z$ is the set of nontrivial zeros of $\zeta$; the
set $Z$ is nonempty, $Z$ lies in $0<\Re s<1$, and $Z+1$ in $1<\Re s<2$.
In particular no vertical line contains all zeros of $\zeta_X$ in
$0<\Re s<2$.
\end{proposition}

\begin{proof}
The closed points of $\mathbb{P}^1_\Z$ lie over primes $p$, and grouping
them by $p$ gives $\zeta_X(s)=\prod_pZ(\mathbb{P}^1_{\F_p},p^{-s})$ for
$\Re s>2$. Since $\#\mathbb{P}^1(\F_{p^n})=p^n+1$,
$Z(\mathbb{P}^1_{\F_p},t)=1/((1-t)(1-pt))$, so
$\zeta_X(s)=\prod_p(1-p^{-s})^{-1}(1-p^{1-s})^{-1}=\zeta(s)\zeta(s-1)$.
In $0<\Re s<2$ the zeros of $\zeta(s)$ are the points of $Z$ (in
$0<\Re s<1$) together with no others, since $\zeta$ has no zeros on
$\Re s\ge1$ (Euler product for $\Re s>1$; Theorem~\ref{thm:imp-riemann}
for $\Re s=1$); the zeros of $\zeta(s-1)$ there are the points $\rho+1$
with $\rho$ a zero of $\zeta$ in $-1<\Re\rho<1$, and $\zeta$ has no zeros
with $-1<\Re\rho\le0$ by the functional equation $\xi(s)=\xi(1-s)$ and
the preceding sentence (the trivial zeros are at $-2,-4,\dots$, and at
$\rho=0$ one uses $\xi(0)\ne0$, Proposition~\ref{prop:xiK}). So the
zeros are $Z\cup(Z+1)$. Finally $Z\ne\emptyset$: otherwise $\xi$, an
entire function of order one (Theorem~\ref{thm:imp-riemann}), would have
no zeros and would be $e^{a+bs}$ by Hadamard's theorem
\cite[Theorem~1.4, p.~2]{Dupuy2017Hadamard}; then $\xi(s)=\xi(1-s)$
forces $b=0$, so $\xi$ would be constant, whereas $\xi(s)\to\infty$ as
$s\to+\infty$ along $\R$, because $\zeta(s)\ge1$ for real $s>1$ and
$\pi^{-s/2}\Gamma(s/2)\to\infty$.
\end{proof}

\begin{theorem}[{(i) after Deninger \cite[p.~172]{Deninger1998Analogies},
\cite[p.~5]{Deninger2010HilbertPolya}}]\label{thm:A-schemes}
Let $d\ge1$, let $Z\not\equiv0$ be meromorphic on $\C$, and let a Weil
datum of type $(d,0)$ satisfy the derivation rule of \ref{ax:AW1},
\ref{ax:AW2} for $Z$, the trace equivariance of \ref{ax:AW3}, and
\ref{ax:AW4}. Then:
\begin{enumerate}[label=(\roman*)]
\item\label{item:As-line} every eigenvalue $\rho$ of $\Theta_j$ has
  $\Re\rho=j/2$, and each $\Theta_j$ is diagonalizable;
\item\label{item:As-ord} for every $\rho$ with $j=2\Re\rho\in
  \{0,1,\dots,2d\}$, $\ord_\rho Z=(-1)^{j+1}m_j(\rho)$, and $Z$ has no
  zero or pole $\rho$ with $2\Re\rho\notin\{0,1,\dots,2d\}$;
\item\label{item:As-rh} hence every zero of $Z$ has real part in
  $\{\tfrac12,\tfrac32,\dots,\tfrac{2d-1}2\}$, every pole has real part
  in $\{0,1,\dots,d\}$, the zeros are exactly the odd-degree spectra and
  the poles exactly the even-degree spectra, and $Z$ determines every
  spectral divisor $\operatorname{Sp}_j$.
\end{enumerate}
\end{theorem}

\begin{proof}
By Proposition~\ref{prop:polarization-identity} with $w=0$, $h_j$ is a
positive infinitesimal polarization of weight $j$, and
Proposition~\ref{prop:purity} gives~\ref{item:As-line}. By
\ref{item:As-line}, the supports of the $\operatorname{Sp}_j$ lie on the
pairwise disjoint lines $\Re s=j/2$, so in \eqref{eq:aw-div} the
coefficient of $[\rho]$ is the single term with $j=2\Re\rho$, which
gives~\ref{item:As-ord}; and~\ref{item:As-rh} follows, the sign
$(-1)^{j+1}$ being $+1$ for odd $j$ and $-1$ for even $j$.
\end{proof}

Poles are part of the statement: for an arithmetic surface whose
generic fibre is an elliptic curve $E$ over $\Q$, $\zeta_X(s)=\zeta(s)
\zeta(s-1)/L(E,s)$ up to finitely many Euler factors, and the zeros of
$L(E,s)$ on $\Re s=1$ are poles of $\zeta_X$ coming from $V^2$. The
theorem uses no regularized determinant. Without positivity its last
clause fails.

It is natural to expect that the determinant identity for the total
zeta function of a higher-dimensional object determines the spectra in
each degree, and that zeros and poles of different degrees may cancel
in the alternating product regardless of positivity; in fact both
depend on \ref{ax:AW4}: without it the degreewise spectra are not
determined, and with it no cancellation occurs
(Theorem~\ref{thm:A-schemes}).

\begin{definition}\label{def:adjoin-pair}
Let $(V^\bullet,\Theta,\cup,\Tr,*)$ be a Weil datum of type $(2,0)$
with unit $1\in V^0=\C\cdot1$ for the cup product and $\Theta_0=0$,
let $\eta\in V^4$ satisfy
$\Theta_4\eta=2\eta$ and $\Tr\eta=1$, and let $\rho\in\C$. The
\emph{adjunction of the pair $\{\rho,2-\rho\}$} is the tuple obtained by
adding to $V^1$ and $V^3$ lines $\C e$ and $\C f$, and to $V^2$ lines
$\C g$ and $\C g'$, with $\Theta e=\rho e$, $\Theta f=(2-\rho)f$,
$\Theta g=\rho g$, $\Theta g'=(2-\rho)g'$; with products $1\cup x=x\cup
1=x$ for the new vectors $x$, $e\cup f=\eta=-f\cup e$,
$g\cup g'=\eta=g'\cup g$, and all other products involving a new vector
equal to zero; and with the stars extended by zero on the new lines.
\end{definition}

\begin{proposition}\label{prop:no-positivity-cancel}
Let $Z\not\equiv0$ be meromorphic on $\C$, let a Weil datum
$(V^\bullet,\Theta,\cup,\Tr,*)$ of type $(2,0)$ satisfy \ref{ax:AW1} for
$U=\mathbb 1$, \ref{ax:AW2} for $Z$, and \ref{ax:AW3}, and let
$\rho\in\C$. Then there is $\eta\in V^4$ with
$\Theta_4\eta=2\eta$ and $\Tr\eta=1$, and the adjunction of the pair
$\{\rho,2-\rho\}$ (Definition~\ref{def:adjoin-pair}) satisfies
\ref{ax:AW1} for $U=\mathbb 1$, \ref{ax:AW2} for the same $Z$, and
\ref{ax:AW3}, while its spectral divisors $\operatorname{Sp}_1$,
$\operatorname{Sp}_2$, and $\operatorname{Sp}_3$ differ from those of the
given datum.
\end{proposition}

\begin{proof}
By \ref{ax:AW1}, $V^4\cong\C(-2)$ is a line with $\Theta_4=2$, and by
\ref{ax:AW3} the pairing $V^0\times V^4\to\C$, $(c,x)\mapsto\Tr(c\cup
x)$, is perfect, so $\Tr$ is nonzero on $V^4$ and $\eta$ exists. Every
product of two new vectors that is not declared zero lands in $V^4$,
and every product of three elements of positive degree involving a new
vector vanishes (it lies in degree above $4$ or involves a vanishing
product), so the product is associative; the signs are those required
by graded commutativity in degrees $1+3$ and $2+2$. The derivation rule
holds on the new products because the eigenvalues of the factors add
up to $2$, and on products with the unit trivially. The endpoint
conditions are unchanged. For \ref{ax:AW3} the new lines pair perfectly,
with eigenvalues $\rho\leftrightarrow2-\rho$, in degrees
$1\leftrightarrow3$ and $2\leftrightarrow2$, and trivially with the old
ones. For \ref{ax:AW2}, the right side of \eqref{eq:aw-div} changes by
$[\rho]-([\rho]+[2-\rho])+[2-\rho]=0$. The new lines raise
multiplicities in degrees $1$, $2$, and $3$, so $\operatorname{Sp}_1$,
$\operatorname{Sp}_2$, and $\operatorname{Sp}_3$ change.
\end{proof}

With \ref{ax:AW4} the adjunction is impossible unless $\Re\rho=1$ and
$\Re(2-\rho)=1$, in which case the new lines in $V^1$ and $V^3$ would
violate Theorem~\ref{thm:A-schemes}\ref{item:As-line}. The same
phenomenon for trace formulas: a Lefschetz or explicit-formula identity
for $X$ determines only the alternating sum $\sum_j(-1)^jm_j(\rho)$,
which for $d=1$ is the degree-one divisor, and for $d\ge2$ becomes a
degreewise statement only in combination with \ref{ax:AW4}.

\begin{remark}\label{rem:A-motivic}
Let $M$ be a specified pure coefficient of weight $w$ over a number
field, with a specified meromorphic completed $L(M,s)$.
If this function admits a polarized graded spectral realization
of type $(1,w)$ in Definition~\ref{def:graded-realization},
Theorem~\ref{thm:main-general} gives
$\Re\rho=(w+\nu)/2$ in degree $\nu$. Thus its zeros lie on
$\Re s=(w+1)/2$, and its poles can lie only on
$\Re s=w/2$ and $\Re s=w/2+1$.
This is the coefficient weight calculation of
\cite[(7.11)--(7.13), pp.~32--33]{Deninger1994Motivic} and
\cite[pp.~6--7]{Deninger2010HilbertPolya}.
It uses a spectral coefficient realization, not an internal
unital Weil datum of positive weight. In Deninger's coefficient
formalism the duality and star maps have dual or complementary
coefficient targets. Such a coefficient realization is additional
to the unweighted conjectures on $\Schbar^{\mathrm{proj}}$. The trivial
coefficient has $w=0$ and recovers the number-field case.
An arithmetic Abel--Jacobi comparison is a further hypothesis in
the conditional BSD rank argument.
\end{remark}

\subsection{The normalization of the determinant}
\label{subsec:det-normalization}

Theorem~\ref{thm:A} uses \ref{ax:AW2} in the divisor form, so it is
independent of any regularization. For AW-det on $\Xbar_K$ the endpoint
blocks give $D_0=s/2\pi$ and $D_2=(s-1)/2\pi$, and \eqref{eq:aw-det}
reads $D_1(s)=e^{a+bs}\xi_K(s)/4\pi^2$. Over $\F_q$ the normalization
term is exactly $1$ for the principal cut (Theorem~\ref{thm:B}); over a
number field it is determined by the spectrum alone and is computed in
Theorem~\ref{thm:det-K}.

\begin{theorem}[{Deninger; Soul\'e; compare Kurokawa;
\cite[Theorem~3.3, (3.3.1), p.~147]{Deninger1992Local},
\cite[pp.~164, 167, and p.~168]{Deninger1998Analogies}}]
\label{thm:imp-deninger-product}
Let $\rho$ run over the nontrivial zeros of $\zeta$, with multiplicity.
The zeta-regularized product $\prod_\rho(s-\rho)/2\pi$, formed with the
arguments in $(-\pi,\pi]$, equals
$\tfrac{s}{2\pi}\tfrac{s-1}{2\pi}\,2^{-1/2}\pi^{-s/2}\Gamma(s/2)
\zeta(s)$.
\end{theorem}

The first source proves this formula for $\Re s>1$
\cite[Theorem~3.3, p.~147]{Deninger1992Local}, and reports its
extension to all of $U_{\mathrm{pr}}$ (the source's set $\Omega$) in a
letter of Soul\'e \cite[Remark~3.4(1), p.~147]{Deninger1992Local};
for $\Re s>1$ the first source takes the arguments in $(-\pi/2,\pi/2)$,
which agrees there with the principal branch, and the convention
$(-\pi,\pi]$ on $U_{\mathrm{pr}}$ is that of
\cite[p.~167]{Deninger1998Analogies};
the second restates it, attributing it to Deninger and to Schr\"oter and
Soul\'e \cite{SchroeterSoule1994Deninger}, and notes that a related
formula had been observed earlier by Kurokawa
\cite[p.~24]{Kurokawa1988Parabolic}, for every number field, up to a
constant and assuming the Riemann hypothesis.
With $\xi$ as in Definition~\ref{def:zeta}, the right side is
$2^{-1/2}\xi(s)/2\pi^2$. Theorem~\ref{thm:det-K} below reproves it and
extends it to every number field.

It is natural to expect that for $K=\Q$ the determinant of the
nontrivial zeros is exactly $\xi/2\pi^2$, that is, that AW-det holds
with normalization term $1$ relative to $\Lambda$; in fact, by
Theorem~\ref{thm:imp-deninger-product}, the normalization term
relative to $\Lambda$ is the constant $2^{-1/2}$, and it is $1$
relative to Deninger's completion (Theorem~\ref{thm:det-K}). The
zero-location theorems use neither value.

The rest of this subsection computes the normalization term for every
number field. It is stated with the gamma factors that zeta
regularization singles out.

\begin{definition}\label{def:deninger-gamma}
Let $K$ be a number field with $r_1$ real and $r_2$ complex places,
discriminant $D_K$, and degree $n=[K:\Q]=r_1+2r_2$.
\begin{enumerate}[label=(\alph*)]
\item \emph{Deninger's archimedean factors} are
  $\Gamma^{\mathrm D}_{\R}(s)=2^{-1/2}\pi^{-s/2}\Gamma(s/2)$ and
  $\Gamma^{\mathrm D}_{\C}(s)=(2\pi)^{-s}\Gamma(s)$.
\item \emph{Deninger's completion} of $\zeta_K$ is
  $\widehat\zeta_K(s)=\Gamma^{\mathrm D}_{\R}(s)^{r_1}
  \Gamma^{\mathrm D}_{\C}(s)^{r_2}\zeta_K(s)$.
\item Let $(V,\Theta)$ be a spectral core with spectrum
  $\divisor\xi_K$, counted with multiplicity $m(\rho)=\ord_\rho\xi_K$,
  and let $s\in\C$ be such that $\Im s\ne\Im\rho$ for every zero $\rho$
  of $\xi_K$. The \emph{half-plane pieces} of its spectral zeta function
  for the principal cut are
  \[
  Z_+(w;s)=\sum_{\Im\rho>\Im s}m(\rho)
    \Bigl(\frac{s-\rho}{2\pi}\Bigr)^{-w},\qquad
  Z_-(w;s)=\sum_{\Im\rho<\Im s}m(\rho)
    \Bigl(\frac{s-\rho}{2\pi}\Bigr)^{-w},
  \]
  with principal powers, wherever the series converge absolutely; so
  $Z_+$ collects the zeros above $s$.
\item The \emph{Stirling constants} of $K$ are
  \[
  A=\tfrac n2,\quad B=\tfrac12\log|D_K|-\tfrac n2(1+\log2\pi),\quad
  C=2-\tfrac{r_1+r_2}2,\quad
  E=\tfrac{r_1}2\log4\pi+\tfrac{r_2}2\log8\pi.
  \]
\item The \emph{standard branch} of $\log\xi_K$ on $\Re z>1$ is
  \[
  \begin{aligned}
  f(z)={}&\log z+\log(z-1)+\tfrac z2\log|D_K|
  +r_1\bigl(-\tfrac z2\log\pi+\log\Gamma(\tfrac z2)\bigr)\\
  &+r_2\bigl(\log2-z\log2\pi+\log\Gamma(z)\bigr)
  +\sum_{\mathfrak p}\sum_{k\ge1}\frac{N\mathfrak p^{-kz}}k,
  \end{aligned}
  \]
  with principal logarithms, $\log\Gamma$ the branch on
  $\C\setminus(-\infty,0]$ that is real on $(0,\infty)$, and the sum
  over the nonzero prime ideals.
\end{enumerate}
\end{definition}

Deninger's factors are those of \cite[p.~247]{Deninger1991Gamma} and
Flach--Morin \cite[(1), p.~521]{FlachMorin2020DeningerWeilArakelov};
they are $\Gamma^{\mathrm D}_{\R}=2^{-1/2}\GR$ and
$\Gamma^{\mathrm D}_{\C}=\tfrac12\GC$, so
\begin{equation}\label{eq:zeta-hat}
\widehat\zeta_K(s)=2^{-n/2}|D_K|^{-s/2}\Lambda_K(s).
\end{equation}
The set of $s$ in Definition~\ref{def:deninger-gamma}(c) is
$U_{\mathrm{pr}}\cap U_{\mathrm{rf}}$ for the spectrum $\divisor\xi_K$
(Definition~\ref{def:principal-cut}).

The proofs use three classical inputs.

\begin{theorem}[{Stirling; Hadamard;
\cite[Lemma~3.1, p.~24]{Hindry2010ZetaL},
\cite[p.~3]{HasanalizadeShenWong2021CountingZeros},
\cite[Theorem~1.4, p.~2, specialized to order one, and Lemma~4.10,
p.~11]{Dupuy2017Hadamard}}]
\label{thm:imp-stirling}
Let $K$ be a number field.
\begin{enumerate}[label=(\roman*)]
\item\label{item:st-stirling} With $\log\Gamma$ as in
  Definition~\ref{def:deninger-gamma}(e),
  $\log\Gamma(z)=(z-\tfrac12)\log z-z+\tfrac12\log2\pi+R(z)$ with
  $|R(z)|\le c_\delta/|z|$ for $|z|\ge1$ and $|\arg z|\le\pi-\delta$,
  for each $\delta>0$.
\item\label{item:st-order} $\xi_K$ is an entire function of order one.
\item\label{item:st-hadamard} An entire function of order one with
  zeros $\rho\ne0$ counted with multiplicity satisfies
  $\sum_\rho|\rho|^{-2}<\infty$ and is
  $e^{a+bs}s^{k}\prod_\rho(1-s/\rho)e^{s/\rho}$ for some $a,b\in\C$
  and $k\ge0$.
\end{enumerate}
\end{theorem}

The first source states~\ref{item:st-stirling} ``when $s$ stays away
from negative reals'', which is the sector condition; the second
recalls~\ref{item:st-order}; the third proves~\ref{item:st-hadamard}
in the form of Hadamard's factorization theorem. The count of zeros of
$\zeta_K$ \cite[Theorem~1.1, p.~2]{HasanalizadeShenWong2021CountingZeros}
gives, for the number $N_K(T)$ of zeros $\rho$ of $\xi_K$ with
$|\Im\rho|\le T$, the bounds $N_K(T)=O(T\log T)$ and
$N_K(T+1)-N_K(T)=O(\log T)$.

By~\ref{item:st-order} and~\ref{item:st-hadamard}, the function
\begin{equation}\label{eq:G-K}
G_{\xi}(z)=-(\log\xi_K)''(z)=\sum_\rho\frac{m(\rho)}{(z-\rho)^2}
\qquad(z\notin\divisor\xi_K),
\end{equation}
with $m(\rho)=\ord_\rho\xi_K$, is given by an absolutely and locally
uniformly convergent series.

\begin{proposition}\label{prop:xiK-asymp}
Let $K$ be a number field with Stirling constants $A,B,C,E$ and standard
branch $f$ of $\log\xi_K$ (Definition~\ref{def:deninger-gamma}(d),(e)),
and let $s\in\C$. Then, for $z=s+x$ as $x\to+\infty$,
\[
f(z)=Az\log z+Bz+C\log z+E+O(|z|^{-1}),\qquad
f'(z)=A\log z+A+B+\frac Cz+O(|z|^{-2}),
\]
and $G_\xi(z)=-f''(z)=-A/z+C/z^2+O(|z|^{-3})$.
\end{proposition}

\begin{proof}
As $x\to\infty$, $\arg z\to0$, so Theorem~\ref{thm:imp-stirling}
\ref{item:st-stirling} applies to $z$ and $z/2$; the remainder $R$ is
holomorphic on the sector, and Cauchy's estimate on the disc of radius
$|z|/4$ about $z$, which lies in a slightly larger sector, gives
$R'(z)=O(|z|^{-2})$ and $R''(z)=O(|z|^{-3})$. Expanding,
$-\tfrac z2\log\pi+\log\Gamma(\tfrac z2)=\tfrac z2\log z
-\tfrac z2(1+\log2\pi)-\tfrac12\log z+\tfrac12\log4\pi+O(|z|^{-1})$ and
$\log2-z\log2\pi+\log\Gamma(z)=z\log z-z(1+\log2\pi)-\tfrac12\log z
+\tfrac12\log8\pi+O(|z|^{-1})$, with derivatives of the remainders
$O(|z|^{-2})$ and $O(|z|^{-3})$; and $\log z+\log(z-1)=2\log z
+O(|z|^{-1})$, likewise. Each rational prime $p$ lies below at most
$n$ prime ideals, each of norm a power $p^{f}$ with $1\le f\le n$, so
the prime sum and its first two derivatives are bounded by a constant
multiple of $\sum_{m\ge2}(1+\log m)^2m^{-\Re z}$, which decays
exponentially in $\Re z$ and is therefore $O(|z|^{-3})$. Adding
up gives the expansion of $f$ with the stated $A$, $B$, $C$, $E$, and
differentiating it termwise gives those of $f'$ and $f''$.
\end{proof}

\begin{proposition}\label{prop:beta}
Let $a\in\C\setminus(-\infty,0]$ and $0<\Re w<2$. Then, with the
principal power,
\[
a^{-w}=P(w)\int_0^\infty x^{1-w}(x+a)^{-2}\,dx,\qquad
P(w)=\frac{\sin\pi w}{\pi(1-w)}.
\]
The function $P$ is entire, $P(1)=1$, and
$(2\pi)^wP(w)=w+(1+\log2\pi)w^2+O(w^3)$ at $w=0$.
\end{proposition}

\begin{proof}
For $a>0$ the substitution $x=at$ gives $a^{-w}B(2-w,w)=
a^{-w}\Gamma(2-w)\Gamma(w)=a^{-w}(1-w)\pi/\sin\pi w$. For $a$ in a
compact subset of $\C\setminus(-\infty,0]$ one has $|x+a|\ge c(1+x)$
for $x\ge0$, so the integral is holomorphic in $a$ and the identity
persists. The function $P$ has a removable singularity at $w=1$ with
value $1$, and $P(w)=w+w^2+O(w^3)$.
\end{proof}

\begin{theorem}[{(ii) Soul\'e; Illies;
\cite[Theorem~7.18(ii), p.~35]{Deninger1994Motivic}}]\label{thm:det-K}
Let $K$ be a number field with $n=[K:\Q]$, $r_1$ real and $r_2$
complex places, and discriminant $D_K$, and let $(V,\Theta)$ be a
spectral core with spectrum $\divisor\xi_K$, counted with multiplicity.
\begin{enumerate}[label=(\roman*)]
\item\label{item:dK-zeta} For $s\in U_{\mathrm{pr}}$ the spectral
  zeta function $Z(w;s)$ of the principal cut converges absolutely for
  $\Re w>1$ and continues holomorphically to
  $\{\Re w>-1\}\setminus\{1\}$, with a simple pole at $w=1$ of residue
  $-\pi n$, and
  $Z(0;s)=\tfrac n2s+2-\tfrac{r_1+r_2}2$.
\item\label{item:dK-det} $(V,\Theta)$ is zeta-regularizable for its
  principal cut, and
  \[
  D_1(s)=\detreg\Bigl(\frac{s-\Theta}{2\pi}\Bigr)
  =2^{-n/2}(2\pi)^{-2}|D_K|^{-s/2}\xi_K(s)
  =\frac s{2\pi}\cdot\frac{s-1}{2\pi}\cdot\widehat\zeta_K(s).
  \]
\item\label{item:dK-norm} Every Weil datum of type $(1,0)$ on $\Xbar_K$
  satisfying \ref{ax:AW1} and \ref{ax:AW2} for $(\Lambda_K,\mathbb 1)$
  satisfies AW-det for $\Lambda_K$ with the principal cuts and
  normalization term
  \[
  e^{a_K+b_Ks}=2^{-n/2}|D_K|^{-s/2},\qquad a_K=-\tfrac n2\log2,\quad
  b_K=-\tfrac12\log|D_K|,
  \]
  and satisfies AW-det for $\widehat\zeta_K$ with normalization term
  $1$.
\item\label{item:dK-fe} $D_1(1-s)=|D_K|^{s-1/2}D_1(s)$.
\end{enumerate}
\end{theorem}

\begin{proof}[Proof of \ref{item:dK-zeta}]
Let $s\in U_{\mathrm{pr}}$, put $g_s(x)=G_\xi(s+x)$ for $x\ge0$ (the
ray $s+[0,\infty)$ misses the zeros), and let $1<u=\Re w<2$. For a zero
$\rho$ put $a=s-\rho\notin(-\infty,0]$. For $x\ge2|a|$ one has
$|x+a|\ge x/2$, so $\int_{2|a|}^\infty x^{1-u}|x+a|^{-2}dx\le
\tfrac4u(2|a|)^{-u}$. Since $0\le\Re\rho\le1$, all but finitely many
zeros satisfy $|\Im a|\ge|a|/4$, and then
$\int_0^{2|a|}x^{1-u}|x+a|^{-2}dx\le16|a|^{-2}(2|a|)^{2-u}/(2-u)$.
Hence $\sum_\rho m(\rho)\int_0^\infty x^{1-u}|x+s-\rho|^{-2}dx$ is
bounded by a constant times $\sum_\rho m(\rho)|s-\rho|^{-u}$, which is
finite because $N_K(T)=O(T\log T)$. The same bound gives absolute
convergence of $Z(w;s)$ for $\Re w>1$. Since dividing by $2\pi$ does
not change arguments, $((s-\rho)/2\pi)^{-w}=(2\pi)^w(s-\rho)^{-w}$, and
Proposition~\ref{prop:beta}, Fubini's theorem, and~\eqref{eq:G-K} give,
for $1<\Re w<2$,
\[
Z(w;s)=(2\pi)^wP(w)\,I(w),\qquad I(w)=\int_0^\infty x^{1-w}g_s(x)\,dx.
\]
By Proposition~\ref{prop:xiK-asymp}, $g_s(x)=c_1x^{-1}+c_2x^{-2}
+O(x^{-3})$ as $x\to\infty$ with $c_1=-A$ and $c_2=As+C$, and $g_s$ is
continuous on $[0,\infty)$. Hence, for $1<\Re w<2$,
\[
I(w)=\int_0^1x^{1-w}g_s\,dx+\int_1^\infty x^{1-w}\Bigl(g_s-\frac{c_1}x
-\frac{c_2}{x^2}\Bigr)dx+\frac{c_1}{w-1}+\frac{c_2}w,
\]
where the first integral is holomorphic for $\Re w<2$ and the second for
$\Re w>-1$. So $I$ continues meromorphically to $-1<\Re w<2$ with
simple poles at $w=1$ and $w=0$, of residues $c_1$ and $c_2$. The
factor $(2\pi)^wP(w)$ is entire, vanishes at $w=0$, and equals $2\pi$ at
$w=1$ (Proposition~\ref{prop:beta}). Together with the absolutely
convergent series for $\Re w>1$, this gives the continuation of $Z$,
the residue $2\pi c_1=-\pi n$ at $w=1$, and $Z(0;s)=c_2$. For later use
we record, from the expansions at $w=0$,
\begin{equation}\label{eq:dZ0}
\partial_wZ(0;s)=(1+\log2\pi)c_2+I_0,\qquad
I_0=\lim_{X\to\infty}\Bigl[\int_0^Xx\,g_s(x)\,dx-c_1X-c_2\log X\Bigr],
\end{equation}
where $I_0$ is the constant term of $I$ at $w=0$: it equals
$\int_0^1xg_s\,dx+\int_1^\infty x(g_s-c_1x^{-1}-c_2x^{-2})\,dx-c_1$,
which is the limit displayed.
\end{proof}

\begin{proof}[Proof of \ref{item:dK-det}]
Extend the standard branch $f$
(Definition~\ref{def:deninger-gamma}(e)) to
$U_{\mathrm{pr}}$ by $f(s)=f(s+X)-\int_0^X(\xi_K'/\xi_K)(s+x)\,dx$ for
any $X$ with $\Re s+X>1$. The value does not depend on $X$; since the
segment $s+[0,X]$ stays away from the discrete set of zeros for $s$ in a
small disc, $f$ is holomorphic on $U_{\mathrm{pr}}$, with
$f'=\xi_K'/\xi_K$, $f''=-G_\xi$, and $e^f=\xi_K$. Integrating by parts,
$\int_0^Xx\,g_s(x)\,dx=-Xf'(s+X)+f(s+X)-f(s)$. With $z=s+X$,
Proposition~\ref{prop:xiK-asymp} gives
\[
-Xf'(z)+f(z)=As\log z+Bs+C\log z+E-C-AX+o(1),
\]
using $z-X=s$ and $X/z\to1$. Subtracting $c_1X+c_2\log X=-AX+(As+C)\log X$
and letting $X\to\infty$ gives $I_0=Bs+E-C-f(s)$, so by~\eqref{eq:dZ0}
\[
\partial_wZ(0;s)=(1+\log2\pi)(As+C)+Bs+E-C-f(s)
=\tfrac s2\log|D_K|+\tfrac n2\log2+2\log2\pi-f(s),
\]
because $A(1+\log2\pi)+B=\tfrac12\log|D_K|$ and $C\log2\pi+E=
2\log2\pi+(\tfrac{r_1}2+r_2)\log2$. The right side is holomorphic on
$U_{\mathrm{pr}}$, and $\exp(-\partial_wZ(0;s))=2^{-n/2}(2\pi)^{-2}
|D_K|^{-s/2}\xi_K(s)$ is entire. This is zeta-regularizability
(Definition~\ref{def:regdet}) and the first formula; the second is
\eqref{eq:zeta-hat} with $\xi_K=s(s-1)\Lambda_K$.
\end{proof}

\begin{proof}[Proof of \ref{item:dK-norm}]
By \ref{ax:AW1}, $V^0\cong\mathbb 1$ and $V^2\cong\mathbb 1(-1)$ are
lines with generators $0$ and $1$, so $D_0=s/2\pi$ and $D_2=(s-1)/2\pi$
for every cut, and $\operatorname{Sp}_0=[0]$, $\operatorname{Sp}_2=[1]$.
By \ref{ax:AW2}, $\operatorname{Sp}_1=\divisor\Lambda_K+[0]+[1]=
\divisor\xi_K$, so \ref{item:dK-det} applies to $V^1$, each $D_j$ has
divisor $\operatorname{Sp}_j$, and
$D_1/(D_0D_2)=4\pi^2D_1/(s(s-1))=2^{-n/2}|D_K|^{-s/2}\Lambda_K
=\widehat\zeta_K$.
\end{proof}

\begin{proof}[Proof of \ref{item:dK-fe}]
By \ref{item:dK-det} and $\xi_K(1-s)=\xi_K(s)$
(Proposition~\ref{prop:xiK}).
\end{proof}

For $K=\Q$, Theorem~\ref{thm:det-K}\ref{item:dK-det} is
Theorem~\ref{thm:imp-deninger-product}. For general $K$, the formula
$D_1=\frac s{2\pi}\frac{s-1}{2\pi}\widehat\zeta_K$ is the case of the
trivial character of \cite[Theorem~7.18(ii), p.~35]{Deninger1994Motivic},
which the source attributes to Soul\'e and to the Diplomarbeit of
Illies (M\"unster, 1992). It is
natural to expect that computing the normalization needs the general
theory of regularized products; in fact the proof above uses only
Hadamard's factorization, Stirling's formula, and the zero count
(Theorem~\ref{thm:imp-stirling} and the remarks after it), and a beta
integral (Proposition~\ref{prop:beta}). The
constant $a_K$ is the regularized determinant of the archimedean
spectra: the spectra $\{-2\nu\}_{\nu\ge0}$ and $\{-\nu\}_{\nu\ge0}$ have
regularized products $\prod_\nu\bigl(\frac\nu\pi+\frac s{2\pi}\bigr)=
\Gamma^{\mathrm D}_{\R}(s)^{-1}$ and
$\prod_\nu\bigl(\frac\nu{2\pi}+\frac s{2\pi}\bigr)=
\Gamma^{\mathrm D}_{\C}(s)^{-1}$
\cite[Proposition~2.1, p.~248]{Deninger1991Gamma}, and
$a_K=r_1\log(\Gamma^{\mathrm D}_{\R}/\GR)+r_2\log(\Gamma^{\mathrm D}_{\C}
/\GC)$.

These notes keep the completion $\Lambda_K$ of
Definition~\ref{def:zeta}. With it, the normalization term of AW-det on
$\Xbar_K$ is the explicit function $2^{-n/2}|D_K|^{-s/2}$ of
Theorem~\ref{thm:det-K}\ref{item:dK-norm}; with Deninger's completion
$\widehat\zeta_K$ it is $1$, exactly as over $\F_q$
(Theorem~\ref{thm:B}\ref{item:B-det}).

It is natural to expect that for every number field the regularized
determinant of the nontrivial zeros reproduces the completion
$\Lambda_K$ of Definition~\ref{def:zeta}, that is,
$D_1=\xi_K/4\pi^2$ with normalization term $1$ relative to $\Lambda_K$
(for $K=\Q$, $D_1=\xi/2\pi^2$); in fact the normalization term relative
to $\Lambda_K$ is $2^{-n/2}|D_K|^{-s/2}$, which is not identically $1$
for any $K$, and the normalization term is $1$ exactly relative to
$\widehat\zeta_K$, which has Deninger's gamma factors and no
discriminant factor (Theorem~\ref{thm:det-K}).

\begin{theorem}[{after \cite[Lemma~(1.7), p.~6, and (7.19.1),
p.~36]{Deninger1994Motivic}}]\label{thm:det-K-rf}
Let $K$ and $(V,\Theta)$ be as in Theorem~\ref{thm:det-K}. Then
$(V,\Theta)$ is zeta-regularizable for its reflected cut
(Definition~\ref{def:principal-cut}); for $s\in U_{\mathrm{rf}}$ its
spectral zeta function is $Z_{\mathrm{rf}}(w;s)=e^{-i\pi w}Z(w;1-s)$,
with $Z$ that of the principal cut; and its regularized determinant is
\[
D_{\mathrm{rf}}(s)=i^{r_2}e^{-i\pi ns/2}\,2^{-n/2}(2\pi)^{-2}
|D_K|^{(s-1)/2}\,\xi_K(s)
=e^{a'+b's}D_1(s),
\]
with $a'=\tfrac{i\pi r_2}2-\tfrac12\log|D_K|$ and
$b'=\log|D_K|-\tfrac{i\pi n}2$.
\end{theorem}

\begin{proof}
For $u\notin[0,\infty)$ the argument of $u$ in $[0,2\pi)$ is
$\operatorname{Arg}(-u)+\pi$. Hence, for $s\in U_{\mathrm{rf}}$,
$((s-\rho)/2\pi)^{-w}=e^{-i\pi w}(((1-s)-(1-\rho))/2\pi)^{-w}$ with the
reflected power on the left and the principal power on the right. As
$\rho\mapsto1-\rho$ permutes the zeros with multiplicity and
$1-s\in U_{\mathrm{pr}}$ exactly when $s\in U_{\mathrm{rf}}$, summing
gives $Z_{\mathrm{rf}}(w;s)=e^{-i\pi w}Z(w;1-s)$, which is regular at
$w=0$ by Theorem~\ref{thm:det-K}\ref{item:dK-zeta}. Then
$-\partial_wZ_{\mathrm{rf}}(0;s)=i\pi Z(0;1-s)-\partial_wZ(0;1-s)$, so
$D_{\mathrm{rf}}(s)=e^{i\pi Z(0;1-s)}D_1(1-s)$. By
Theorem~\ref{thm:det-K}\ref{item:dK-zeta},
$e^{i\pi Z(0;1-s)}=e^{i\pi(n-r_1-r_2)/2}e^{2\pi i}e^{-i\pi ns/2}
=i^{r_2}e^{-i\pi ns/2}$, and by \ref{item:dK-det} and
\ref{item:dK-fe} of the same theorem
$D_1(1-s)=2^{-n/2}(2\pi)^{-2}|D_K|^{(s-1)/2}\xi_K(s)$. The quotient
$D_{\mathrm{rf}}/D_1=i^{r_2}|D_K|^{s-1/2}e^{-i\pi ns/2}$ is
$e^{a'+b's}$.
\end{proof}

The cut dependence has the shape of Proposition~\ref{prop:ff-cut}: the
real part of the slope $b'$ is the logarithm of the conductor, $\log
|D_K|$ in place of $\log q$, and the phase $e^{-i\pi ns/2}$ has no
finite-field counterpart because the archimedean part of the spectrum
is not periodic.

It is natural to expect that the half-plane pieces $Z_\pm$ of the
spectral zeta function are regular at $w=0$, so that $D_1(1-s)$ and
$D_1(s)$ can be compared piece by piece through the phases
$e^{\pm i\pi w}$; in fact each piece has a pole at $w=0$ for every
$s\ne\tfrac12$, with the residue that the main term
$n\log T$ of the counting function predicts, and the poles cancel in
the sum. It is then natural to expect that the normalization over a
number field cannot be determined without a finer regularization of
the pieces; in fact the full spectral zeta function is regular at $w=0$
and determines it (Theorem~\ref{thm:det-K}).

\begin{proposition}[{after \cite[(1.3)--(1.7), pp.~4--6, and
Theorem~(7.20), p.~37]{Deninger1994Motivic}}]\label{prop:Zpm}
Let $K$ and $(V,\Theta)$ be as in Theorem~\ref{thm:det-K}, and let
$s\in\C$ with $\Im s\ne\Im\rho$ for every zero $\rho$ of $\xi_K$. Then
the half-plane pieces $Z_\pm(w;s)$ of
Definition~\ref{def:deninger-gamma}(c) converge absolutely for
$\Re w>1$, $Z_++Z_-=Z$, and each continues meromorphically to
$\Re w>-1$ with poles only at $w=0$ and $w=1$:
\[
\operatorname{Res}_{w=0}Z_\pm(w;s)=\mp\frac{i\,n\,(s-\frac12)}{2\pi},
\qquad
Z_\pm(w;s)=\pm\frac{i\,n}{(w-1)^2}+O\bigl((w-1)^{-1}\bigr)
\quad(w\to1).
\]
In particular neither piece is regular at $w=0$ unless $s=\tfrac12$.
With the opposite labelling (the piece over the zeros below $s$ called
$Z_+$) the signs of both residues change.
\end{proposition}

\begin{proof}
Absolute convergence and $Z_++Z_-=Z$ follow from
Theorem~\ref{thm:det-K}\ref{item:dK-zeta}, since no zero has
$\Im\rho=\Im s$. For a zero above $s$, $s-\rho$ lies in the lower half
plane and its argument in $[0,2\pi)$ is its principal argument plus
$2\pi$; for a zero below $s$ the two arguments agree. Hence, for
$\Re w>1$, the reflected spectral zeta function is
$Z_{\mathrm{rf}}(w;s)=e^{-2\pi iw}Z_+(w;s)+Z_-(w;s)$, and by
Theorem~\ref{thm:det-K-rf}
\[
(1-e^{-2\pi iw})\,Z_+(w;s)=Z(w;s)-e^{-i\pi w}Z(w;1-s).
\]
Both $s$ and $1-s$ lie in $U_{\mathrm{pr}}$, so by
Theorem~\ref{thm:det-K}\ref{item:dK-zeta} the right side, $N(w)$, is
holomorphic on $\{\Re w>-1\}\setminus\{1\}$. The factor
$1-e^{-2\pi iw}$ vanishes on $\Re w>-1$ only at $w=0,1,2,\dots$, simply;
for $\Re w>1$, $Z_+$ is holomorphic, so the only possible poles are at
$0$ and $1$. At $w=0$, $N(0)=Z(0;s)-Z(0;1-s)=n(s-\tfrac12)$ and
$1-e^{-2\pi iw}=2\pi iw+O(w^2)$, so the residue is
$n(s-\frac12)/2\pi i$. At $w=1$, $N(w)=-2\pi n/(w-1)+O(1)$ because
$e^{-i\pi w}=-1+O(w-1)$ and both $Z(\cdot;s)$ and $Z(\cdot;1-s)$ have
residue $-\pi n$, while $1-e^{-2\pi iw}=2\pi i(w-1)+O((w-1)^2)$; so
$Z_+=in/(w-1)^2+O((w-1)^{-1})$. The statements for $Z_-=Z-Z_+$ follow,
$Z$ being regular at $0$ and having a simple pole at $1$.
\end{proof}

It is natural to expect that zeta-regularizability of $\divisor\xi_K$
follows from its counting function alone, which has an explicit main
term and error $O(\log T)$; in fact an error of that size can create a
pole at $w=0$, and the proof of Theorem~\ref{thm:det-K} uses the
Hadamard product and the Stirling asymptotics of $\xi_K$ along
horizontal rays, not only the count. It is equally natural to expect
that for the zeros of $\xi_K$ the fluctuation of the count contributes
to the residues of $Z_\pm$ at $w=0$; in fact it contributes nothing,
the residues being those predicted by the main term
(Proposition~\ref{prop:Zpm}).

\begin{proposition}[{compare \cite[Remark~(3.4)(2), p.~147]%
{Deninger1992Local}}]\label{prop:N-count}
Let $K$ be a number field, let $\Sigma=\divisor\xi_K$ and
\[
\Sigma'=\Sigma+\sum_{k\ge1}\bigl([\tfrac12+ie^k]+[\tfrac12-ie^k]\bigr).
\]
\begin{enumerate}[label=(\roman*)]
\item\label{item:Nc-core} $\Sigma'$ is the spectrum, with multiplicity,
  of a spectral core; it is invariant under $\rho\mapsto1-\rho$ and
  $\rho\mapsto\bar\rho$; and its number of points with
  $0<\Im\rho\le T$ exceeds that of $\Sigma$ by $\lfloor\log T\rfloor$
  for $T\ge1$.
\item\label{item:Nc-pole} For every $s$ in the principal domain
  $U_{\mathrm{pr}}$ of $\Sigma'$, the spectral zeta function of
  $\Sigma'$ for the principal cut continues meromorphically to
  $\Re w>-1$ and has a simple pole at $w=0$ with residue $2$. So
  $\Sigma'$ is not zeta-regularizable for its principal cut.
\end{enumerate}
\end{proposition}

\begin{proof}[Proof of \ref{item:Nc-core}]
The added points are discrete, symmetric under both maps, and
$\#\{k\ge1:e^k\le T\}=\lfloor\log T\rfloor$; take the direct sum of
the spectral cores with spectra $\Sigma$ and the added points.
\end{proof}

\begin{proof}[Proof of \ref{item:Nc-pole}]
The spectral zeta function of $\Sigma'$ is that of $\Sigma$, which is
regular at $w=0$ by Theorem~\ref{thm:det-K}\ref{item:dK-zeta}, plus
$Y(w)=\sum_k\bigl[((s-\frac12-ie^k)/2\pi)^{-w}
+((s-\frac12+ie^k)/2\pi)^{-w}\bigr]$. Put
$\varepsilon_k=(s-\tfrac12)e^{-k}$ and let $k_0$ satisfy
$|\varepsilon_k|<\tfrac12$ for $k\ge k_0$. For such $k$,
$s-\tfrac12\mp ie^k=e^k(\mp i)(1\pm i\varepsilon_k)$, and the principal
logarithms add because $|\operatorname{Arg}(1\pm i\varepsilon_k)|<\pi/6$;
so the $k$-th term is
$(2\pi)^we^{-kw}\bigl[e^{i\pi w/2}(1+i\varepsilon_k)^{-w}
+e^{-i\pi w/2}(1-i\varepsilon_k)^{-w}\bigr]$. Since
$|(1\pm i\varepsilon)^{-w}-1|\le c|w||\varepsilon|$ for $w$ bounded and
$|\varepsilon|<\tfrac12$, the sum over $k\ge k_0$ equals
\[
\frac{2(2\pi)^w\cos(\pi w/2)\,e^{-k_0w}}{1-e^{-w}}+h(w),
\]
where $h$ is a sum of terms $O(|w|e^{-k(\Re w+1)})$, holomorphic on
$\Re w>-1$. The first term has a simple pole at $w=0$ with residue $2$,
and the finitely many terms with $k<k_0$ are entire. So $Y$, and with it
the spectral zeta function of $\Sigma'$, has a simple pole at $w=0$ of
residue $2$, and Definition~\ref{def:regdet}(c) fails.
\end{proof}


\subsection{The trace-formula axiom}\label{subsec:trace-axiom}

Over $\F_q$ the cohomology is constructed independently of the zeta
function, the Lefschetz trace formula for all powers of Frobenius is a
theorem about it, and rationality with the correct multiplicities is a
consequence of that formula. This subsection compares the divisor axiom
\ref{ax:AW2} on $\Xbar_K$ with the trace-formula axiom AW6$'$ of
Definition~\ref{def:aw6prime}, and calibrates the comparison against
$\F_q$. The input on the number-field side is the explicit formula.

\begin{theorem}[{Weil; \cite[Theorem~4.1, (17), p.~15]%
{Deninger2001NumberTheory}, \cite[(1), p.~1]{Deninger2002Explicit}}]
\label{thm:imp-explicit}
Let $K$ be a number field and $\varphi\in C_c^\infty(\R)$, with
transform $\Phi$ and geometric side $G_K(\varphi)$
(Definitions~\ref{def:test-classes} and~\ref{def:geom-side}). Then
\[
\Phi(0)-\sum_\rho\ord_\rho(\xi_K)\,\Phi(\rho)+\Phi(1)=G_K(\varphi),
\]
the sum running over the distinct zeros $\rho$ of $\xi_K$ and
converging absolutely.
\end{theorem}

The first source states this identity for every $K$, with the
archimedean terms given explicitly for $\varphi$ supported in
$(0,\infty)$ or in $(-\infty,0)$, where they agree with
Definition~\ref{def:geom-side}; it notes that they can be written on
all of $\R$ in different ways and refers to Barner
\cite{Barner1981Explicit} for a proof. The explicit formula in this
generality is due to Weil \cite{Weil1952Explicit}, as reported in
\cite[\S V, p.~8]{Bombieri2000Clay}. The
second states it for $K=\Q$ with the archimedean term $W_\R$ of
Definition~\ref{def:geom-side} on all of $\R$. For general $K$ the form
of $W_\R$ and $W_\C$ at $t=0$ is the one produced by the contour
argument, which we recall.

\begin{proof}[Derivation of the archimedean terms]
Put $\lambda=\Lambda_K'/\Lambda_K=\xi_K'/\xi_K-1/s-1/(s-1)$ and fix
$c>1$. Integration by parts gives $|\Phi(s)|\le C_k(1+|s|)^{-k}$ on
every vertical strip, for every $k$. By
Theorem~\ref{thm:imp-stirling}\ref{item:st-order}
and~\ref{item:st-hadamard},
$\xi_K'/\xi_K(s)=b+\sum_\rho(1/(s-\rho)+1/\rho)$ with
$\sum|\rho|^{-2}<\infty$, the zeros repeated with multiplicity. Since
$N_K(T+2)-N_K(T-2)=O(\log T)$, each interval $[j,j+1]$ contains a
point $T_j$ at distance at least $1/(C\log j)$ from the imaginary part
of every zero. For $s=\sigma+iT_j$ with $1-c\le\sigma\le c$, the zeros
with $|\rho|>2|s|$ contribute $O(|s|\sum|\rho|^{-2})$, and the
$O(|s|\log|s|)$ others contribute at most
$|s|\cdot C\log j/\min|\rho|$ each; so $|\lambda(s)|=O(T_j^3)$, and
likewise at $-T_j$ because the zeros are symmetric under conjugation.
The residue theorem on the rectangles with vertices $1-c\pm iT_j$,
$c\pm iT_j$ therefore gives, in the limit, since the residue of
$\Phi\lambda$ at $0$ and $1$ is $-\Phi(0)$, $-\Phi(1)$ and at $\rho$ is
$\ord_\rho(\xi_K)\Phi(\rho)$,
\[
\frac1{2\pi i}\Bigl[\int_{(c)}-\int_{(1-c)}\Bigr]\Phi(s)\lambda(s)\,ds
=-\Phi(0)+\sum_\rho\ord_\rho(\xi_K)\Phi(\rho)-\Phi(1).
\]
On $\Re s=1-c$ substitute $s=1-u$ and use $\lambda(1-u)=-\lambda(u)$,
which follows from $\Lambda_K(1-s)=\Lambda_K(s)$
(Theorem~\ref{thm:imp-hecke}); this gives
\[
\Phi(0)-\sum_\rho\ord_\rho(\xi_K)\Phi(\rho)+\Phi(1)
=-\frac1{2\pi i}\int_{(c)}\bigl[\Phi(s)+\Phi(1-s)\bigr]
\lambda(s)\,ds.
\]
For $\Re s>1$, $\lambda(s)=\tfrac12\log|D_K|+r_1\GR'/\GR(s)
+r_2\GC'/\GC(s)-\sum_{\mathfrak p}\sum_{k\ge1}\log N\mathfrak p\,
N\mathfrak p^{-ks}$, and Fourier inversion on $\Re s=c$ gives
$\frac1{2\pi i}\int_{(c)}\Phi(s)e^{-sT}ds=\varphi(T)$ and
$\frac1{2\pi i}\int_{(c)}\Phi(1-s)e^{-sT}ds=e^{-T}\varphi(-T)$. The
constant term gives $-\log|D_K|\,\varphi(0)$ and the prime terms give
the prime sum of $G_K$ (the interchange is justified by
$|\Phi(c+iy)|\le C(1+|y|)^{-2}$). For the archimedean terms insert
Gauss's integral
$\psi(s)=\Gamma'/\Gamma(s)=\int_0^\infty\bigl(e^{-u}/u
-e^{-su}/(1-e^{-u})\bigr)du$ ($\Re s>0$): with $u=2t$,
$\GR'/\GR(s)=-\tfrac12\log\pi+\int_0^\infty\bigl(e^{-2t}/2t
-e^{-st}/(1-e^{-2t})\bigr)dt$, and
$\GC'/\GC(s)=-\log2\pi+\psi(s)$. Interchanging the $s$- and
$t$-integrals (Fubini: the $t$-integral of the absolute value grows at
most linearly in $|s|$ on $\Re s=c$, and $\Phi$ decays faster than any
power there), the term $-e^{-st}/(1-e^{-2t})$ produces
$(\varphi(t)+e^{-t}\varphi(-t))/(1-e^{-2t})$, the term $e^{-2t}/2t$
produces $-e^{-2t}\varphi(0)/t$, and the constant produces
$\log\pi\cdot\varphi(0)$; this is $W_\R$, and $W_\C$ follows in the
same way.
\end{proof}

The derivation uses Gauss's integral for $\psi$, a classical
consequence of the product formula for $\Gamma$, which is not
re-derived here. For $\varphi\in\mathcal S^+$ the theorem is the
one-sided explicit formula of \cite[(18), p.~15]%
{Deninger2001NumberTheory} with~\eqref{eq:geom-one-sided}, and no
choice at $t=0$ enters.

\begin{proposition}\label{prop:exp-unique}
Let $\Sigma\subset\{a\le\Re s\le b\}$ be a discrete set with
$\#\{\rho\in\Sigma:|\rho|\le R\}=O(R^B)$, and let $c:\Sigma\to\C$
satisfy $|c_\rho|\le C(1+|\rho|)^A$.
\begin{enumerate}[label=(\roman*)]
\item\label{item:eu-conv} For every $\varphi\in C_c^\infty(\R)$ the
  series $\sum_\rho c_\rho\Phi(\rho)$ converges absolutely, and
  $\varphi\mapsto\sum_\rho c_\rho\Phi(\rho)$ is a distribution $T_c$ on
  $\R$.
\item\label{item:eu-tail} If, for some $T_0\in\R$, $T_c(\varphi)=0$
  for every $\varphi$ in $C_c^\infty((T_0,\infty))$, or for every
  $\varphi$ in $C_c^\infty((-\infty,T_0))$, then $c=0$.
\end{enumerate}
\end{proposition}

\begin{proof}[Proof of \ref{item:eu-conv}]
Fix a real $s_0>b$ and an integer $N>A+B$. Partial summation against
the counting function gives $\sum_\rho|c_\rho||\rho-s_0|^{-N}<\infty$.
Put $g(t)=\sum_\rho c_\rho(\rho-s_0)^{-N}e^{\rho t}$; since
$|e^{\rho t}|\le e^{\max(|a|,|b|)|t|}$, the series converges uniformly
on compact sets, so $g$ is continuous. For $\varphi\in C_c^\infty(\R)$,
$N$ integrations by parts give $\int e^{\rho t}(-\frac d{dt}-s_0)^N
\varphi\,dt=(\rho-s_0)^N\Phi(\rho)$, hence
$|\Phi(\rho)|\le C_\varphi|\rho-s_0|^{-N}$ on the strip, which gives
absolute convergence, and $\sum_\rho c_\rho\Phi(\rho)=\int g\,
(-\frac d{dt}-s_0)^N\varphi\,dt$, which is continuous in $\varphi$.
So $T_c=(\frac d{dt}-s_0)^Ng$ in the sense of distributions.
\end{proof}

\begin{proof}[Proof of \ref{item:eu-tail}]
Suppose $T_c=0$ on $(T_0,\infty)$. Then $(\frac d{dt}-s_0)^Ng=0$ there,
so $g(t)=e^{s_0t}P(t)$ on $(T_0,\infty)$ for a polynomial $P$ of degree
less than $N$ (the distributional solutions of a linear differential
equation with constant coefficients are classical). For $\Re s>s_0$
compute $\mathcal L(s)=\int_{T_0}^\infty g(t)e^{-st}\,dt$ in two ways.
Termwise, which is justified by $\sum_\rho|c_\rho||\rho-s_0|^{-N}
\int_{T_0}^\infty e^{(b-\Re s)t}dt<\infty$,
\[
\mathcal L(s)=M(s):=\sum_\rho c_\rho(\rho-s_0)^{-N}
\frac{e^{(\rho-s)T_0}}{s-\rho}.
\]
From the polynomial form, $\mathcal L(s)=e^{-(s-s_0)T_0}Q(1/(s-s_0))$
with $Q$ a polynomial, holomorphic on $\C\setminus\{s_0\}$. The series
$M$ converges locally uniformly on the complement of
$\Sigma_c=\{\rho:c_\rho\ne0\}$, because on a compact set $D$ disjoint
from $\Sigma_c$ the factors $e^{(\rho-s)T_0}$ are bounded and
$|s-\rho|^{-1}\le C_D$; so $M$ is meromorphic on $\C$ with simple poles
exactly at the points of $\Sigma_c$, of residue
$c_\rho(\rho-s_0)^{-N}\ne0$. The set $\C\setminus(\Sigma_c\cup\{s_0\})$
is connected, and $M$ agrees there with
$e^{-(s-s_0)T_0}Q(1/(s-s_0))$ by the identity theorem, which has no
pole in $\Sigma_c$ (note $s_0\notin\Sigma_c$). So $\Sigma_c$ is empty.
The case $(-\infty,T_0)$ follows by $t\mapsto-t$, $\rho\mapsto-\rho$.
\end{proof}

It is natural to expect that an exponential sum over a discrete
spectrum determines its coefficients only if it vanishes for all
positive times; in fact vanishing on one tail $(T_0,\infty)$ suffices
(Proposition~\ref{prop:exp-unique}\ref{item:eu-tail}), while
vanishing on a bounded interval does not (Theorem~\ref{thm:window}).

The trace formula of Definition~\ref{def:aw6prime} is the form in which
Deninger expects the explicit formula to arise from a Lefschetz formula
\cite[(10), p.~172]{Deninger1998Analogies},
\cite[(19), p.~16]{Deninger2001NumberTheory}.

\begin{theorem}\label{thm:aw6-div}
Let $K$ be a number field and let a Weil datum of type $(1,0)$ on
$\Xbar_K$ satisfy \ref{ax:AW1} for $(\Lambda_K,\mathbb 1)$. The
following are equivalent:
\begin{enumerate}[label=(\roman*)]
\item\label{item:ad-two} the datum satisfies AW6$'$;
\item\label{item:ad-one} it satisfies AW6$'$ on $\mathcal S^+$;
\item\label{item:ad-tail} it satisfies AW6$'$ on $\mathcal S_{>T_0}$
  for some $T_0\in\R$;
\item\label{item:ad-div} the datum satisfies \ref{ax:AW2} for
  $(\Lambda_K,\mathbb 1)$, that is,
  $\operatorname{Sp}_1=\divisor\xi_K$.
\end{enumerate}
In particular, for $T_0\ge0$, the identity~\eqref{eq:aw6prime} for
test functions supported in $(T_0,\infty)$, whose geometric side
involves only the prime powers with $k\log N\mathfrak p>T_0$ and the
archimedean places, determines $\operatorname{Sp}_1$ with multiplicities.
\end{theorem}

\begin{proof}
\ref{item:ad-two}$\Rightarrow$\ref{item:ad-one}$\Rightarrow$%
\ref{item:ad-tail} is clear, since $\mathcal S^+=\mathcal S_{>0}$. For
\ref{item:ad-tail}$\Rightarrow$\ref{item:ad-div}: by \ref{ax:AW1},
$V^0\cong\mathbb 1$ and $V^2\cong\mathbb 1(-1)$, so
$\operatorname{Sp}_0=[0]$, $\operatorname{Sp}_2=[1]$,
and~\eqref{eq:aw6prime} reads $\Phi(0)-\sum_\rho m_1(\rho)\Phi(\rho)
+\Phi(1)=G_K(\varphi)$. Subtracting Theorem~\ref{thm:imp-explicit} gives
$\sum_\rho c_\rho\Phi(\rho)=0$ for $\varphi\in\mathcal S_{>T_0}$, with
$c=m_1-\ord(\xi_K)$ on $\Sigma=\spec\Theta_1\cup\divisor\xi_K$. By (G)
and $N_K(T)=O(T\log T)$, $\Sigma$ lies in a vertical strip with
polynomially bounded counting function, and $|c_\rho|$ is bounded by
the sum of the two counting functions at $|\rho|$. By
Proposition~\ref{prop:exp-unique}\ref{item:eu-tail}, $c=0$, that is,
$\operatorname{Sp}_1=\divisor\xi_K$; and $\divisor\Lambda_K=
\divisor\xi_K-[0]-[1]=-\operatorname{Sp}_0+\operatorname{Sp}_1
-\operatorname{Sp}_2$ is \ref{ax:AW2}. For
\ref{item:ad-div}$\Rightarrow$\ref{item:ad-two}: $V^0$ and $V^2$ are
lines, and $\spec\Theta_1=\divisor\xi_K$ lies in $0\le\Re s\le1$ with
counting function $O(R\log R)$, so (G) holds in every degree; the
traces converge absolutely by
Proposition~\ref{prop:exp-unique}\ref{item:eu-conv}, and the left side
of~\eqref{eq:aw6prime} is that of Theorem~\ref{thm:imp-explicit}.
\end{proof}

It is natural to expect that, as over $\F_q$, imposing a trace-formula
axiom in place of \ref{ax:AW2} would make the divisor identity a
theorem and the arithmetic axioms weaker; in fact, by
Theorem~\ref{thm:aw6-div}, given \ref{ax:AW1} the axiom AW6$'$ is
equivalent to \ref{ax:AW2}.
Imposing AW6$'$ instead of \ref{ax:AW2} changes the form of the
assumption, from ``the flow has the right spectrum'' to ``the flow has
the right trace formula'', not its strength. The geometric side of
AW6$'$ is built from the prime ideals, the archimedean places, and the
discriminant alone, but no object satisfying it is known.

It is natural to expect that, as over $\F_q$, where
$\log Z(X,t)=\sum_kN_kt^k/k$ is an identity of power series, the trace
formula integrates to the determinant identity; in fact AW6$'$
determines the divisor only, and the determinant identity with its
normalization term is a separate analytic statement about
$\divisor\xi_K$: by Theorems~\ref{thm:aw6-div}
and~\ref{thm:det-K}\ref{item:dK-norm}, a datum satisfying \ref{ax:AW1}
and AW6$'$ satisfies AW-det for $\Lambda_K$ with normalization term
$2^{-n/2}|D_K|^{-s/2}$.

For a datum of type $(d,0)$ with $d\ge2$ the same argument shows that a
trace formula determines only the alternating divisor
$\sum_j(-1)^{j+1}\operatorname{Sp}_j$; it does not determine the
degreewise spectra without positivity
(Proposition~\ref{prop:no-positivity-cancel},
Theorem~\ref{thm:A-schemes}). For $d=1$, \ref{ax:AW1} fixes degrees
$0$ and $2$, so no cancellation between degrees is possible.

The finite-field calibration is the following theorem.

\begin{definition}\label{def:lefschetz-numbers}
Let $X$ be a smooth projective geometrically connected variety over
$\F_q$, with $H^j$ and $\mathrm F$ as in Theorem~\ref{thm:B}; $\mathrm F$
is invertible on each $H^j$ (Theorem~\ref{thm:imp-weil-coh}). The
\emph{Lefschetz numbers} of $X$ are
$N_k=\sum_j(-1)^j\Tr(\mathrm F^k\mid H^j)$, $k\in\Z$.
\end{definition}

\begin{theorem}[{(i) Grothendieck; (ii) after
\cite[(2), p.~3]{Deninger2002Explicit},
\cite[(10), p.~172]{Deninger1998Analogies}}]\label{thm:ff-lefschetz}
Let $X$ be a smooth projective geometrically connected variety of
dimension $d$ over $\F_q$, with $H^j$, $\mathrm F$, and the suspended
datum $(\widehat H^\bullet,\Theta,\cup,\Tr,0)$ as in
Theorem~\ref{thm:B}, multiplicity functions $m_j$ of
$(\widehat H^j,\Theta_j)$, and Lefschetz numbers $N_k$
(Definition~\ref{def:lefschetz-numbers}).
\begin{enumerate}[label=(\roman*)]
\item\label{item:fl-counts} $N_k=\#X(\F_{q^k})$ for $k\ge1$,
  $N_0=\sum_j(-1)^j\dim H^j$, and $N_{-k}=q^{-dk}N_k$.
\item\label{item:fl-formula} For every $\varphi\in C_c^\infty(\R)$,
  \[
  \sum_{j=0}^{2d}(-1)^j\sum_\rho m_j(\rho)\Phi(\rho)
  =\log q\sum_{k\in\Z}N_k\,\varphi(k\log q),
  \]
  the left side converging absolutely.
\item\label{item:fl-converse} Let $(W^\bullet,\Theta',\cup,\Tr',*')$ be
  a Weil datum of type $(d,0)$ whose spectral cores satisfy (G), with
  spectral divisors $\operatorname{Sp}'_j$. If the identity
  of~\ref{item:fl-formula}, with its right side, holds for its
  multiplicity functions for every $\varphi\in C_c^\infty((T_0,\infty))$
  for some $T_0$, then
  $\sum_j(-1)^{j+1}\operatorname{Sp}'_j=\divisor\zeta_X$, that is, the
  datum satisfies \ref{ax:AW2} for $\zeta_X$.
\end{enumerate}
\end{theorem}

\begin{proof}[Proof of \ref{item:fl-counts}]
Take logarithms in the rationality statement of
Theorem~\ref{thm:imp-deligne}. Since
$\log\det(1-\mathrm Ft\mid H^j)^{-1}=\sum_k\Tr(\mathrm F^k\mid H^j)
t^k/k$, this gives $\sum_k\#X(\F_{q^k})t^k/k=\sum_kN_kt^k/k$. The value $N_0$ is
immediate. By Theorem~\ref{thm:imp-weil-coh}, the pairing of $H^j$
with $H^{2d-j}$ is perfect and satisfies
$\operatorname{tr}(\mathrm Fu\cup\mathrm Fv)=
q^d\operatorname{tr}(u\cup v)$. So $\mathrm F$ on $H^{2d-j}$ is $q^d$
times the adjoint of $\mathrm F^{-1}$ on $H^j$, and hence
$\Tr(\mathrm F^{-k}\mid H^j)=q^{-dk}\Tr(\mathrm F^k\mid H^{2d-j})$, and
$(-1)^j=(-1)^{2d-j}$.
\end{proof}

\begin{proof}[Proof of \ref{item:fl-formula}]
By Definition~\ref{def:log-susp} and
Proposition~\ref{prop:suspension}\ref{item:susp-core},
$\spec\Theta_j$ is the union of the progressions $\Sigma_q(\alpha)$
over the eigenvalues $\alpha$ of $\mathrm F$ on $H^j$, with
$m_j(\rho)=\dim H^j_{(\alpha)}$ for $\rho\in\Sigma_q(\alpha)$. Fix
$\alpha$ and $\rho_0\in\Sigma_q(\alpha)$, and put $L=\log q$ and
$\psi(t)=\varphi(t)e^{\rho_0t}$. Then $\Sigma_q(\alpha)=\{\rho_0
+2\pi in/L\}$, and the Poisson summation formula for the smooth
compactly supported function $\psi$ gives
$\sum_n\Phi(\rho_0+2\pi in/L)=\sum_n\int\psi(t)e^{2\pi int/L}dt
=L\sum_k\psi(kL)=L\sum_k\alpha^k\varphi(kL)$, since
$e^{\rho_0kL}=\alpha^k$. Summing over $\alpha$ with multiplicity and
over $j$ with signs gives the identity; absolute convergence holds
because $\Phi$ decays faster than any power on vertical strips.
\end{proof}

\begin{proof}[Proof of \ref{item:fl-converse}]
Subtract the identity~\ref{item:fl-formula} for the suspended datum. The
difference is $\sum_\rho c_\rho\Phi(\rho)=0$ for
$\varphi\in C_c^\infty((T_0,\infty))$, where $c_\rho$ is the difference
of the two alternating sums of multiplicities; both spectra lie in
vertical strips with polynomially bounded counting functions. By
Proposition~\ref{prop:exp-unique}\ref{item:eu-tail}, $c=0$, so
$\sum_j(-1)^{j+1}\operatorname{Sp}'_j$ is the alternating spectral
divisor of the suspended datum, which is $\divisor\zeta_X$ by
Theorem~\ref{thm:B}\ref{item:B-det}.
\end{proof}

So over $\F_q$ the Lefschetz formula for all powers of Frobenius from
some point on determines the alternating divisor of $\zeta_X$ by the
same uniqueness statement, Proposition~\ref{prop:exp-unique}, that
gives Theorem~\ref{thm:aw6-div}; the difference between the two
settings is that over $\F_q$ the formula is a theorem about an
independently constructed cohomology. The atom at $t=0$ of the right
side of~\ref{item:fl-formula} is $\log q\cdot N_0$, an Euler
characteristic; in the geometric side over $K$ its place is taken by
$-\log|D_K|$ and the archimedean constants.

It is natural to expect that a faithful analogue of the Lefschetz
formula must be one-sided, since only the numbers $N_k$ with $k\ge1$
count points; in fact $N_0$ and $N_{-k}$ are fixed by the Euler
characteristic and by duality, so the suspended datum satisfies the
two-sided identity of Theorem~\ref{thm:ff-lefschetz}, and given
\ref{ax:AW1} the one-sided and two-sided forms of AW6$'$ are
equivalent (Theorem~\ref{thm:aw6-div}).

The next statements concern the Hilbert form of AW6$'$.

\begin{proposition}\label{prop:aw5-contains}
Let a Weil datum of type $(d,0)$ have positive definite star form $h_j$
in degree $j$ and satisfy the condition of \ref{ax:AW5} for this $j$.
\begin{enumerate}[label=(\roman*)]
\item\label{item:a5-pol} $h_j$ is a positive infinitesimal polarization
  of weight $j$ on $(V^j,\Theta_j)$, so the datum satisfies AW4-V
  in degree $j$.
\item\label{item:a5-pure} $\Theta_j$ is diagonalizable on $V^j$, every
  eigenvalue has real part $j/2$, and distinct eigenspaces are
  $h_j$-orthogonal.
\item\label{item:a5-spec} $T_j$ is the diagonal operator of an
  orthonormal basis of $\cH^j$ consisting of eigenvectors of $\Theta_j$
  in $V^j$; its spectrum is pure point, equal to
  $\{(\rho-j/2)/i:\rho\in\spec\Theta_j\}$ with the multiplicities
  $m_j$, without finite accumulation point.
\item\label{item:a5-trace} If $(V^j,\Theta_j)$ satisfies (G), then for
  every $\varphi\in C_c^\infty(\R)$ the operator
  $\int_\R\varphi(t)e^{jt/2}U_j(t)\,dt$ is of trace class, with
  trace $\sum_\rho m_j(\rho)\Phi(\rho)$.
\end{enumerate}
\end{proposition}

\begin{proof}[Proof of \ref{item:a5-pol}]
Each $v\in V^j$ lies in a finite-dimensional $\Theta_j$-stable subspace,
on which $U_j(t)v=e^{-jt/2}e^{t\Theta_j}v$ is smooth in $t$.
Differentiating $h_j(U_j(t)v,U_j(t)v)=h_j(v,v)$ at $t=0$ gives
$h_j(\Theta v,v)+h_j(v,\Theta v)=jh_j(v,v)$. Both sides are the
diagonal values of Hermitian forms, so they agree as forms,
which is~\eqref{eq:inf-pol}.
\end{proof}

\begin{proof}[Proof of \ref{item:a5-pure}]
Proposition~\ref{prop:purity}, by~\ref{item:a5-pol}.
\end{proof}

\begin{proof}[Proof of \ref{item:a5-spec}]
By~\ref{item:a5-pol} and Proposition~\ref{prop:hilbert} there is an
orthonormal basis $(e_k)$ of $\cH^j$ in $V^j$ with $\Theta_je_k=
(j/2+i\tau_k)e_k$, and $D=\diag(\tau_k)$ is self-adjoint with the
stated spectrum. On $V^j$, which is the span of the $e_k$, the
generator $A_j$ of $U_j$ acts as $\Theta_j-j/2$, so $T_j=-iA_j$
agrees with $D$ there. Since $V^j$ is a core for $A_j$, $T_j$ is the
closure of its restriction to $V^j$; the closure of the restriction of
$D$ to finite linear combinations of the $e_k$ is $D$ (truncate the
expansion of a vector in its domain). So $T_j=D$.
\end{proof}

\begin{proof}[Proof of \ref{item:a5-trace}]
By~\ref{item:a5-spec}, $U_j(t)e_k=e^{it\tau_k}e_k$, so the operator
maps $e_k$ to $\Phi(j/2+i\tau_k)e_k=\Phi(\rho_k)e_k$. A bounded
operator that is diagonal in an orthonormal basis with absolutely
summable entries is of trace class with trace their sum, and the
entries are summable by Proposition~\ref{prop:exp-unique}
\ref{item:eu-conv} with $c=m_j$.
\end{proof}

It is natural to expect that the requirement of \ref{ax:AW6} that the
completion lose no spectrum is an axiom independent of \ref{ax:AW5}; in
fact it follows from \ref{ax:AW5}
(Proposition~\ref{prop:aw5-contains}\ref{item:a5-spec}), and what
\ref{ax:AW6} adds to \ref{ax:AW5} is exactly (G) and the trace formula,
that is, AW6$'$. It is also natural to expect \ref{ax:AW5} to be a
purely analytic completeness condition, so that ``\ref{ax:AW1},
\ref{ax:AW5}, and AW6$'$ imply \ref{ax:AW2}'' would be the analogue of
``Lefschetz implies rationality''; in fact \ref{ax:AW5} presupposes a
positive definite form $h_j$ and, with its unitary flow, yields the
positivity axiom in the form AW4-V
(Proposition~\ref{prop:aw5-contains}\ref{item:a5-pol}), so these
hypotheses imply the Riemann hypothesis for $\zeta_K$, while over $\F_q$
rationality needs no polarization. The analogue of ``Lefschetz implies
rationality'' is Theorem~\ref{thm:aw6-div}.

\begin{theorem}\label{thm:aw6-hilb}
Let $K$ be a number field and let a Weil datum of type $(1,0)$ on
$\Xbar_K$ satisfy \ref{ax:AW1} for $(\Lambda_K,\mathbb 1)$, have
positive definite star forms, and satisfy \ref{ax:AW5}.
\begin{enumerate}[label=(\roman*)]
\item\label{item:ah-forms} For every $\mathcal S'\subset\mathcal S$,
  the datum satisfies AW6$'$ in Hilbert form on $\mathcal S'$ if and
  only if it satisfies AW6$'$ on $\mathcal S'$.
\item\label{item:ah-rh} If it satisfies AW6$'$ on $\mathcal S$, on
  $\mathcal S^+$, or on a tail class, then
  $\operatorname{Sp}_1=\divisor\xi_K$, every zero of $\xi_K$ has real
  part $\tfrac12$, and $\dim\ker(\Theta_1-\rho)=\ord_\rho\xi_K$ for
  every $\rho$.
\item\label{item:ah-conv} If every zero of $\xi_K$ has real part
  $\tfrac12$, the trivial model with spectrum $\divisor\xi_K$
  (Definition~\ref{def:trivial-model}) has positive definite star forms
  and satisfies \ref{ax:AW1}, \ref{ax:AW3}, \ref{ax:AW4}, \ref{ax:AW5},
  and AW6$'$ in both forms.
\end{enumerate}
In particular, some Weil datum on $\Xbar_K$ with positive definite star
forms satisfies \ref{ax:AW1}, \ref{ax:AW5}, and AW6$'$ if and only if
every zero of $\xi_K$ has real part $\tfrac12$.
\end{theorem}

\begin{proof}[Proof of \ref{item:ah-forms}]
Both forms include (G), and by
Proposition~\ref{prop:aw5-contains}\ref{item:a5-trace} with $w=0$ the
traces $\tr A^{(j)}_\varphi$ are the distributional traces.
\end{proof}

\begin{proof}[Proof of \ref{item:ah-rh}]
By Theorem~\ref{thm:aw6-div}, $\operatorname{Sp}_1=\divisor\xi_K$; by
Proposition~\ref{prop:aw5-contains}\ref{item:a5-pure}, $\Theta_1$ is
diagonalizable with every eigenvalue on $\Re s=\tfrac12$.
\end{proof}

\begin{proof}[Proof of \ref{item:ah-conv}]
By Propositions~\ref{prop:trivial-model} and~\ref{prop:content}
\ref{item:ct-free}, the trivial model satisfies \ref{ax:AW1},
\ref{ax:AW2}, \ref{ax:AW3}, and \ref{ax:AW4}, and its star forms are the
standard positive forms in the given bases; $h_1$ is a positive
infinitesimal polarization of weight one because $\Theta_1$ is
diagonal in an $h_1$-orthonormal basis with eigenvalues of real part
$\tfrac12$. By Proposition~\ref{prop:hilbert}, $T=\diag(\tau_k)$ is
self-adjoint on $\cH^1$, $U_1(t)=e^{itT}$ is a strongly continuous
unitary group extending $e^{-t/2}e^{t\Theta_1}$, its generator $iT$ is
skew-adjoint, and $V^1$ is a core for it, as in the proof of
Proposition~\ref{prop:aw5-contains}\ref{item:a5-spec}. In degrees $0$
and $2$, $U_j(t)=1$ on a line. This is \ref{ax:AW5}. AW6$'$ follows
from Theorem~\ref{thm:aw6-div} and~\ref{item:ah-forms}.
\end{proof}

\begin{definition}\label{def:window-objects}
Let $g\in C_c^\infty(\R)$, $L>0$, $q>1$, and let $N\ge2$ be an integer.
\begin{enumerate}[label=(\alph*)]
\item The \emph{reflection} of $g$ is the function
  $g^\dagger(t)=e^{-t}\overline{g(-t)}$, and the transform of $g$ is
  $G(s)=\int_\R g(t)e^{st}\,dt$.
\item The \emph{half-shifted coset} is
  $C_L=\{\tfrac12+\tfrac{2\pi i}L(n+\tfrac12):n\in\Z\}$.
\item The \emph{phantom set} is
  $S_N=\{\sqrt q\,e^{i\pi(2m+1)/N}:0\le m<N\}$.
\end{enumerate}
\end{definition}

\begin{proposition}[{Weil; \cite[\S V, p.~9]{Bombieri2000Clay}}]%
\label{prop:aw6-weil-pos}
Let $K$ and the datum be as in Theorem~\ref{thm:aw6-hilb}, and suppose
it satisfies AW6$'$. Let $g\in C_c^\infty(\R)$, with reflection
$g^\dagger$ and transform $G$ (Definition~\ref{def:window-objects}).
Then the transform of $g*g^\dagger$ is $G(s)\overline{G(1-\bar s)}$,
and
\[
G_K(g*g^\dagger)=2\Re\bigl(G(0)\overline{G(1)}\bigr)
-\sum_\rho\ord_\rho(\xi_K)|G(\rho)|^2
\le2\Re\bigl(G(0)\overline{G(1)}\bigr).
\]
In particular $G_K(g*g^\dagger)\le0$ when $G(0)=G(1)=0$.
\end{proposition}

\begin{proof}
Substituting $u=-t$, $\int g^\dagger(t)e^{st}dt=\overline{\int
g(u)e^{(1-\bar s)u}du}=\overline{G(1-\bar s)}$, and convolution becomes
multiplication. By Theorem~\ref{thm:aw6-hilb}\ref{item:ah-rh}, every
$\rho$ in $\operatorname{Sp}_1=\divisor\xi_K$ satisfies
$1-\bar\rho=\rho$, so the transform at $\rho$ is $|G(\rho)|^2$; at $0$
and $1$ it is $G(0)\overline{G(1)}$ and $G(1)\overline{G(0)}$.
Insert these into~\eqref{eq:aw6prime}.
\end{proof}

For $K=\Q$ and $G(0)=G(1)=0$ this is Weil's positivity criterion in the
form of Theorem~\ref{thm:imp-weil-crit} and
Proposition~\ref{prop:weil-crit-notes}. The terms from degrees $0$ and
$2$ enter the alternating trace with sign $+$ and are not squares; the
two linear conditions $G(0)=G(1)=0$ remove them, which is the role of
the poles $0$ and $1$ in Weil's criterion.

\medskip
The remaining statements concern test classes that do not determine
the spectrum. It is natural to expect that the trace formula for test
functions supported in a bounded interval $(0,L)$, which involves the
prime ideals with $N\mathfrak p<e^L$ and the archimedean places,
determines the degree-one spectrum, since such test functions separate
the points of the spectrum; in fact it does not.

\begin{theorem}\label{thm:window}
Let $K$ be a number field, $L>0$, and $C_L$ the half-shifted coset
(Definition~\ref{def:window-objects}).
\begin{enumerate}[label=(\roman*)]
\item\label{item:w-poisson} For every $\varphi\in C_c^\infty(\R)$,
  $\sum_{\rho\in C_L}\Phi(\rho)=L\sum_{k\in\Z}(-1)^ke^{kL/2}\varphi(kL)$.
\item\label{item:w-datum} Let $\Sigma=\divisor\xi_K+\sum_{\rho\in
  C_L}[\rho]$. The trivial model with spectrum $\Sigma$ satisfies
  \ref{ax:AW1} for $(\Lambda_K,\mathbb 1)$, \ref{ax:AW3}, (G), and
  AW6$'$ on $\mathcal S_{(0,L)}$; its degree-one spectrum is
  $\Sigma\ne\divisor\xi_K$, so it does not satisfy \ref{ax:AW2}.
\item\label{item:w-rh} If every zero of $\xi_K$ has real part
  $\tfrac12$, that trivial model also has positive definite star forms
  and satisfies \ref{ax:AW4} and \ref{ax:AW5}.
\end{enumerate}
\end{theorem}

\begin{proof}[Proof of \ref{item:w-poisson}]
With $\psi(t)=\varphi(t)e^{t/2}e^{i\pi t/L}$, the Poisson summation
formula gives $\sum_n\Phi(\tfrac12+\tfrac{2\pi i}L(n+\tfrac12))
=\sum_n\int\psi(t)e^{2\pi int/L}dt=L\sum_k\psi(kL)$, and
$\psi(kL)=(-1)^ke^{kL/2}\varphi(kL)$.
\end{proof}

\begin{proof}[Proof of \ref{item:w-datum}]
The set $C_L$ is discrete, invariant under $\rho\mapsto1-\rho$ ($n\mapsto
-n-1$), and does not contain $\tfrac12$; so $\Sigma$ satisfies the
hypotheses of Definition~\ref{def:trivial-model}, the multiplicity of
$\tfrac12$ being $\ord_{1/2}\xi_K$, which is even (proof of
Proposition~\ref{prop:content}\ref{item:ct-free}).
Proposition~\ref{prop:trivial-model}\ref{item:tm-aw13} gives
\ref{ax:AW1}, \ref{ax:AW3}, and $\operatorname{Sp}_1=\Sigma$; (G) holds
because $\Sigma$ lies in $0\le\Re s\le1$ with counting function
$O(R\log R)$. For $\varphi\in\mathcal S_{(0,L)}$, by
Theorem~\ref{thm:imp-explicit} and~\ref{item:w-poisson} the left side
of~\eqref{eq:aw6prime} is $G_K(\varphi)-L\sum_k(-1)^ke^{kL/2}
\varphi(kL)=G_K(\varphi)$, since no point $kL$ lies in $(0,L)$.
\end{proof}

\begin{proof}[Proof of \ref{item:w-rh}]
$\Sigma$ then lies on $\Re s=\tfrac12$. Apply
Proposition~\ref{prop:trivial-model}\ref{item:tm-aw4}, and the proof
of Theorem~\ref{thm:aw6-hilb}\ref{item:ah-conv}, which uses only that
the spectrum lies on the line.
\end{proof}

It is natural to expect that a class of test functions determines the
spectrum as soon as it separates the points of the spectrum; in fact
separation is necessary but not sufficient. The window class
$\mathcal S_{(0,L)}$ separates points, since $e^{\rho t}\ne e^{\rho't}$
on every interval when $\rho\ne\rho'$, and still fails
(Theorem~\ref{thm:window}), while every tail class succeeds
(Proposition~\ref{prop:exp-unique}). The failure is a Poisson-summation
phenomenon of infinite spectra: an infinite exponential sum can vanish
on an interval. The arithmetic form of ``all powers of Frobenius'' is
therefore a class of test functions whose supports reach arbitrarily
far, not a class large enough to separate eigenvalues. The coset $C_L$
is the logarithmic suspension of phantom Frobenius eigenvalues, as the
next proposition shows.

It is natural to expect that finitely many Lefschetz identities
determine the characteristic polynomial of Frobenius; in fact they do
so only together with a bound on the dimension.

\begin{proposition}\label{prop:ff-phantom}
Let $q>1$, let $N\ge2$ be an integer, and let $S_N$ be the phantom set
(Definition~\ref{def:window-objects}).
\begin{enumerate}[label=(\roman*)]
\item\label{item:ph-sums} $\sum_{\alpha\in S_N}\alpha^k=0$ for
  $1\le k<N$ and $\sum_{\alpha\in S_N}\alpha^N=-Nq^{N/2}$; every
  $\alpha\in S_N$ has $|\alpha|=\sqrt q$, and $S_N$ is stable under
  $\alpha\mapsto q/\alpha=\bar\alpha$.
\item\label{item:ph-susp} The logarithmic suspension of
  $(\C^{S_N},\diag(S_N))$ has spectrum $C_{N\log q}$.
\end{enumerate}
\end{proposition}

\begin{proof}[Proof of \ref{item:ph-sums}]
$\sum_me^{i\pi(2m+1)k/N}=e^{i\pi k/N}\sum_{m=0}^{N-1}e^{2\pi imk/N}$
vanishes unless $N\mid k$, and equals $Ne^{i\pi}$ for $k=N$. The
modulus is clear, and $q/\alpha=\sqrt q\,e^{-i\pi(2m+1)/N}
=\sqrt q\,e^{i\pi(2(N-1-m)+1)/N}$.
\end{proof}

\begin{proof}[Proof of \ref{item:ph-susp}]
With $L=N\log q$, the progression $\Sigma_q(\alpha)$ of
$\alpha=\sqrt q\,e^{i\pi(2m+1)/N}$ is
$\{\tfrac12+\tfrac{2\pi i}L(m+\tfrac12+nN):n\in\Z\}$, and as $m$ runs
over $0,\dots,N-1$ and $n$ over $\Z$, $m+nN$ runs over $\Z$ once.
\end{proof}

So for a curve $C$ over $\F_q$, adding $\diag(S_N)$ to $\mathrm F$ on
$V^1$ leaves the numbers $N_k$ of Theorem~\ref{thm:ff-lefschetz} for
$1\le k<N$ unchanged, preserves purity and the symmetry
$\alpha\mapsto q/\alpha$, and changes the characteristic polynomial on
$V^1$. For a fixed curve the Newton identities recover that polynomial
from $N_1,\dots,N_{2g}$ once the genus $g$ is known; the arithmetic
spectrum is infinite, and no window suffices.

It is natural to expect that the two-sided window $(-L,L)$ gives no
more information than $(0,L)$, by the symmetry $\rho\mapsto1-\rho$; in
fact the atom of the geometric side at $t=0$ detects every spectrum
added on the critical line. For the datum of Theorem~\ref{thm:window},
\ref{item:w-poisson} shows that its spectral side differs from
$G_K(\varphi)$ by $L\varphi(0)$ on $\mathcal S_{(-L,L)}$.

\begin{proposition}\label{prop:window-two}
Let $K$ be a number field and let a Weil datum of type $(1,0)$ on
$\Xbar_K$ satisfy \ref{ax:AW1} for $(\Lambda_K,\mathbb 1)$, with
$(V^1,\Theta_1)$ satisfying (G), $\spec\Theta_1\subset\{\Re s=\tfrac12\}$,
and $m_1(\rho)\ge\ord_\rho\xi_K$ for every $\rho$. If the datum
satisfies AW6$'$ on $\mathcal S_{(-L,L)}$ for some $L>0$, then
$\operatorname{Sp}_1=\divisor\xi_K$.
\end{proposition}

\begin{proof}
Put $c=m_1-\ord(\xi_K)\ge0$; it is supported on the line. Subtracting
Theorem~\ref{thm:imp-explicit}, $\sum_\rho c_\rho\Phi(\rho)=0$ for
$\varphi\in\mathcal S_{(-L,L)}$. For $g\in C_c^\infty((-L/2,L/2))$ the
function $g*g^\dagger$ of Proposition~\ref{prop:aw6-weil-pos} lies in
$\mathcal S_{(-L,L)}$, and its transform at a point $\rho$ of the line
is $|G(\rho)|^2$, because $1-\bar\rho=\rho$. So
$\sum_\rho c_\rho|G(\rho)|^2=0$, and $c_\rho G(\rho)=0$ for all $\rho$
and $g$. For $\rho=\tfrac12+i\gamma$ take $g\ge0$, not identically
zero, supported in an interval about $0$ so short that
$\Re(e^{t/2}e^{i\gamma t})>0$ on it; then $\Re G(\rho)>0$. So $c=0$.
\end{proof}

The hypotheses of Proposition~\ref{prop:window-two} put the whole
spectrum, hence every zero of $\xi_K$, on the critical line; the
proposition is a uniqueness statement under that assumption. Whether a
spectrum that omits some zeros of $\xi_K$ can satisfy AW6$'$
on a two-sided window is open (Question~\ref{q:trace-axiom}).

It is natural to expect that any Schwartz function $f$ of the spectral
variable $\tau$ can serve as a test function, via ``$\tr f(T_1)$'';
in fact the prime sum of the geometric side then diverges in general.

\begin{proposition}\label{prop:schwartz}
Let $K$ be a number field, $g(t)=\exp(-(1+t^2)^{1/4})$, and
$\varphi(t)=e^{-t/2}g(t)$. Then $g$ is a Schwartz function, so
$f(\tau)=\int_\R\varphi(t)e^{(1/2+i\tau)t}\,dt=\int_\R g(t)e^{i\tau t}
\,dt$ is a Schwartz function of $\tau$; but
$\sum_{\mathfrak p}\log N\mathfrak p\,\varphi(\log N\mathfrak p)$
diverges.
\end{proposition}

\begin{proof}
Every derivative of $g$ is $g$ times a polynomial in derivatives of
$(1+t^2)^{1/4}$, which are bounded, and $g$ decays faster than every
power of $|t|$; so $g$ is a Schwartz function and so is its Fourier
transform $f$. For $t\ge1$, $(1+t^2)^{1/4}\le2^{1/4}t^{1/2}$, so
$\varphi(\log x)\ge x^{-1/2}\exp(-2^{1/4}(\log x)^{1/2})\ge x^{-3/4}$
for $x$ large. Hence the terms of the sum are eventually at least
$N\mathfrak p^{-1}$. Since $\zeta_K$ has a pole at $s=1$
(Theorem~\ref{thm:imp-hecke}), $\log\zeta_K(\sigma)=\sum_{\mathfrak p}
\sum_kN\mathfrak p^{-k\sigma}/k$ tends to infinity as $\sigma\to1^+$,
while the terms with $k\ge2$ stay bounded; so
$\sum_{\mathfrak p}N\mathfrak p^{-1}=\infty$.
\end{proof}

The admissible test functions are transforms of $\varphi$ for which the
prime sum converges and the spectral side can be evaluated at complex
$\rho$; $C_c^\infty(\R)$ is the natural class, used throughout.
\section{The K\"unneth conjecture}\label{sec:kunneth}

Weil's proof of the Riemann hypothesis for a curve $C$ over $\F_q$ takes
place on the surface $C\times C$: the Frobenius graph and the diagonal are
correspondences, and the positivity it uses is Castelnuovo's positivity
of the trace form on correspondences, a consequence of the Hodge index
theorem on $C\times C$ (Section~\ref{subsec:weil-curves}). A Weil
cohomology over a field has a K\"unneth formula, which is what places
these classes in the cohomology of the surface. 

This section asks what
the analogous statement for $\Xbar_K$ should be. It first fixes the
algebra of tensor products of flows and the data a cohomology of the
``arithmetic surface'' $\Xbar_K\boxtimes\Xbar_L$ would carry
(Section~\ref{subsec:kun-defs}); it then records the refutations, which
show that the natural strong forms of a K\"unneth formula are
inconsistent with the axioms of Section~\ref{sec:axioms}
(Section~\ref{subsec:kun-neg}); it states the conjecture in the form
that survives, with the two notions of compactness separated
(Section~\ref{subsec:kun-conj}); and it proves what holds: a K\"unneth
formula relative to the periodic unit over $\F_q$ and for Tate objects
and an Eilenberg--Moore formula over a base
(Section~\ref{subsec:kun-pos}), and a Hodge index statement on the
arithmetic surface which is equivalent to purity together with
definiteness of the star form, the exact counterpart of the positivity
in Weil's proof (Section~\ref{subsec:kun-hodge}). No surface is
constructed. A K\"unneth formula for the cohomology of
$\Spec\Z\boxtimes\Spec\Z$, with flow $\Theta\otimes1+1\otimes\Theta$,
is already contemplated by Deninger, who observes that the resulting
regularized product over the sums $\rho+\rho'$ of zeros does not make
sense and relates this to Kurokawa's restriction to zeros with
imaginary parts of equal sign \cite[(7.32), pp.~41--42]%
{Deninger1994Motivic}.

\subsection{Flow modules and surface data}\label{subsec:kun-defs}

\begin{definition}\label{def:flow-module}
Let $q>1$ be real.
\begin{enumerate}[label=(\alph*)]
\item A \emph{flow module} is a complex vector space $V$ with a linear
  map $\Theta_V:V\to V$; a morphism of flow modules is a linear map
  commuting with the maps $\Theta$.
\item The \emph{tensor product} of flow modules $(V,\Theta_V)$ and
  $(W,\Theta_W)$ is $(V\otimes_\C W,\Theta_V\otimes1+1\otimes\Theta_W)$.
\item A \emph{flow algebra} is a commutative unital $\C$-algebra $R$
  with a derivation $\Theta_R$; a \emph{flow $R$-module} is an
  $R$-module $M$ with a linear map $\Theta_M$ such that
  $\Theta_M(rm)=\Theta_R(r)m+r\Theta_M(m)$.
\item The \emph{relative tensor product} $M\otimes_RN$ of flow
  $R$-modules is the quotient of $M\otimes_\C N$ by the span of the
  elements $rm\otimes n-m\otimes rn$, with the map induced by
  $\Theta_M\otimes1+1\otimes\Theta_N$.
\item The \emph{periodicity action} on the logarithmic suspension
  $(\widehat V,\Theta)$ of an invertible linear map on a
  finite-dimensional space (Definition~\ref{def:log-susp}) is the action
  of the $q$-periodic unit $\mathbb 1_q=\C[v^{\pm1}]$
  (Definition~\ref{def:unit}) in which $v$ sends the copy of $u$ in
  $\widehat V_\rho$ to the copy of $u$ in $\widehat V_{\rho+2\pi
  i/\log q}$.
\end{enumerate}
\end{definition}

The rule $\Theta\otimes1+1\otimes\Theta$ is the one Deninger derives
from a hypothetical K\"unneth formula \cite[(7.32)]{Deninger1994Motivic};
compare Manin's additive tensor product \cite[(1.22)]{Manin1995Lectures}.

The map in~(d) is well defined: $\Theta_M\otimes1+1\otimes\Theta_N$ sends
$rm\otimes n-m\otimes rn$ to the sum of the three relations
$(\Theta_Rr)m\otimes n-m\otimes(\Theta_Rr)n$,
$r(\Theta_Mm)\otimes n-\Theta_Mm\otimes rn$, and
$rm\otimes\Theta_Nn-m\otimes r\Theta_Nn$. In~(e), $\mathbb 1_q$ is a flow
algebra ($\Theta v^n=(2\pi in/\log q)v^n$ is a derivation), the copy of
$u\in V_{(q^\rho)}$ has the same generalized eigenspace in both summands,
and $\Theta(v\cdot x)=v\cdot\Theta x+(2\pi i/\log q)v\cdot x$, so
$\widehat V$ is a flow $\mathbb 1_q$-module.

\begin{proposition}\label{prop:tensor-eigen}
Let $(V,\Theta_V)$ and $(W,\Theta_W)$ be spectral cores
(Definition~\ref{def:spectral-core}) with multiplicities $m_V$, $m_W$.
Then $V\otimes W$ is the direct sum of the finite-dimensional
$\Theta$-stable subspaces $V_{\Theta\sim\rho}\otimes W_{\Theta\sim\rho'}$,
on which $\Theta=\Theta_V\otimes1+1\otimes\Theta_W$ has the single
eigenvalue $\rho+\rho'$; for every $\lambda\in\C$,
\[
(V\otimes W)_{\Theta\sim\lambda}=\bigoplus_{\rho+\rho'=\lambda}
V_{\Theta\sim\rho}\otimes W_{\Theta\sim\rho'},\qquad
\dim(V\otimes W)_{\Theta\sim\lambda}=\sum_{\rho+\rho'=\lambda}
m_V(\rho)m_W(\rho').
\]
\end{proposition}

\begin{proof}
On $V_{\Theta\sim\rho}\otimes W_{\Theta\sim\rho'}$ the commuting maps
$(\Theta_V-\rho)\otimes1$ and $1\otimes(\Theta_W-\rho')$ are nilpotent,
so $\Theta-\rho-\rho'$ is nilpotent there, and on the blocks with
$\rho+\rho'\ne\lambda$ the map $\Theta-\lambda$ is invertible. A vector
is in the generalized $\lambda$-eigenspace if and only if each of its
block components is.
\end{proof}

The sums may be infinite; this is the phenomenon behind
Theorem~\ref{thm:kun-square}.

\begin{proposition}\label{prop:B-support}
Let $(V^\bullet,\Theta,\cup,\Tr,*)$ be a Weil datum of type $(d,0)$
satisfying the derivation rule of \ref{ax:AW1} and the trace
equivariance of \ref{ax:AW3}. Then $B_j$ vanishes on
$V^j_{\Theta\sim\rho}\times V^{2d-j}_{\Theta\sim\sigma}$ unless
$\rho+\sigma=d$.
\end{proposition}

\begin{proof}
By the derivation rule and the trace equivariance,
$B_j(\Theta u,v)+B_j(u,\Theta v)=dB_j(u,v)$. So the linear form on
$V^j_{\Theta\sim\rho}\otimes V^{2d-j}_{\Theta\sim\sigma}$ induced by
$B_j$ vanishes on the image of $\Theta\otimes1+1\otimes\Theta-d$.
That map has the single eigenvalue $\rho+\sigma-d$ on this
finite-dimensional space (proof of Proposition~\ref{prop:tensor-eigen}),
so it is surjective when $\rho+\sigma\ne d$.
\end{proof}

Definition~\ref{def:aw4-alg} is stated in the introduction.


The trivial model of Definition~\ref{def:trivial-model} with spectrum on
the line satisfies it: for a pair $(e,f)$,
$B_1(*_1(ae+bf),*_1(a'e+b'f))=\bar a\bar b'-\bar b\bar a'$. Positivity of
$h_1$ is not part of the condition.

Definition~\ref{def:surface-datum} is stated in the introduction.


The Hodge index clause is the positivity of the star form $h_S^2$ on
$N_S$, that is, the Castelnuovo (Hodge--Riemann) form of the index
theorem on the classes coming from $V^1_K\otimes V^1_L$; it is not a
statement about the signature of the duality pairing
$(u,v)\mapsto\Tr_S(u\cup v)$ on $H^2$, which is the arithmetic
counterpart of the intersection pairing. The notation
$\Xbar_K\boxtimes\Xbar_L$ for the object whose cohomology a
surface datum would be is only notation: no such object is defined here.
The algebraic part plays the role of the classes of correspondences: by
Proposition~\ref{prop:tensor-eigen}, $N_S$ lies in the generalized
eigenspace of $\Theta_S$ at~$1$, the analogue of the eigenvalue $q$ of
Frobenius on the algebraic classes in $H^2(C\times C)$.

Definition~\ref{def:tensor-datum} is stated in the introduction.


The Tate datum is the cohomology of $\mathbb{P}^n$ with its Tate
twists; its spectrum in degree $2i$ is $\{i\}$. When $K=L$ and the
spectra of both factors lie on the line, the following subspace is the
one singled out by Kurokawa's splitting of the zeros into the halves
$\Im\rho>0$ and $\Im\rho<0$, built into his definition of tensor
products of zeta functions \cite[p.~219]{Kurokawa1992MultipleZeta} and
proposed to define the tensor square of $\zeta$
\cite[\S5.2, p.~158]{Manin1995Lectures}.

\begin{definition}\label{def:kurokawa-cone}
Let $(V,\Theta)$ be a spectral core. The \emph{Kurokawa cone} of
$V\otimes V$ is the sum of the subspaces
$V_{\Theta\sim\rho}\otimes V_{\Theta\sim\rho'}$ over the pairs with
$\Im\rho\cdot\Im\rho'\ge0$.
\end{definition}

\subsection{Refutations}\label{subsec:kun-neg}

The findings of this subsection are stated first, before the positive
results they delimit.

It is natural to expect a K\"unneth formula for fibre products over
$\Spec\Z$, as for products over a field; in fact the scheme-theoretic
square of $\Spec\Z$ over itself is $\Spec\Z$, and a datum is never its
own tensor square.

\begin{proposition}\label{prop:kun-over-Z}
Let a Weil datum $(V^\bullet,\Theta,\cup,\Tr,*)$ of type $(1,0)$ on
$\Xbar_\Q$ satisfy \ref{ax:AW1} for $U=\mathbb 1$ and \ref{ax:AW2} for
$(\Lambda,\mathbb 1)$. Then there is no isomorphism of flow modules
between $V^1$ and the degree-one part $V^0\otimes V^1\oplus V^1\otimes
V^0$ of its tensor square.
\end{proposition}

\begin{proof}
By \ref{ax:AW1}, $V^0=\C$ with $\Theta_0=0$, so the degree-one part of
the tensor square is isomorphic to $V^1\oplus V^1$, with multiplicities
$2m_1(\rho)$. By the proof of Theorem~\ref{thm:rh}\ref{item:rh-divisor},
which uses only \ref{ax:AW1} and \ref{ax:AW2},
$\operatorname{Sp}_1=\divisor\xi_\Q$, and $\xi$ has a zero (proof of
Proposition~\ref{prop:P1Z}); at such a zero $2m_1(\rho)\ne m_1(\rho)$.
\end{proof}

It is natural to expect at least a derived K\"unneth formula over a
base, the cohomology of $X\times_SY$ being that of the derived tensor
product of the cohomologies of $X$ and $Y$ over that of $S$; in fact
this fails already in \'etale cohomology of curves over $\mathbb{P}^1$.

\begin{proposition}\label{prop:kun-EM}
Let $E$ be an elliptic curve over $\F_q$, $x:E\to\mathbb{P}^1_{\F_q}$ a
finite morphism of degree two, $\ell\nmid q$ a prime,
$R=H^\bullet(\mathbb{P}^1_{\overline\F_q},\Q_\ell)$, and
$M=H^\bullet(E_{\overline\F_q},\Q_\ell)$, a graded $R$-module through
$x^*$.
\begin{enumerate}[label=(\roman*)]
\item\label{item:EM-R} $R=\Q_\ell\oplus\Q_\ell h$ with $h$ of degree two
  and $h^2=0$; $H^1(E_{\overline\F_q},\Q_\ell)$ is a two-dimensional
  $R$-submodule of $M$ killed by $h$, and it is a direct summand of $M$.
\item\label{item:EM-tor} The graded module $\operatorname{Tor}^R_1(M,M)$
  has a nonzero component in internal degree four.
\item\label{item:EM-vanish} $H^j((E\times_{\mathbb{P}^1}E)_{\overline
  \F_q},\Q_\ell)=0$ for $j\ge3$.
\end{enumerate}
\end{proposition}

\begin{proof}[Proof of \ref{item:EM-R}]
By Theorem~\ref{thm:imp-weil-coh} for $\mathbb{P}^1$, $H^0=\Q_\ell$ and
$H^2$ is a line on which Frobenius is $q$. Since
$\#\mathbb{P}^1(\F_{q^n})=q^n+1$, $Z(\mathbb{P}^1,t)=1/((1-t)(1-qt))$, so by
Theorem~\ref{thm:imp-deligne} the characteristic polynomial of
Frobenius on $H^1$ is $1$; Frobenius is injective (remark after
Theorem~\ref{thm:imp-weil-coh}), so $H^1=0$. The product $h^2$ lies in
$H^4=0$. The cohomology of $E$ vanishes in degrees above two and
$H^1(E)$ has dimension two (the facts recalled after
Theorem~\ref{thm:imp-deligne}), so $h\cdot H^1(E)\subset H^3(E)=0$, and
$M=(H^0\oplus H^2)\oplus H^1$ as $R$-modules.
\end{proof}

\begin{proof}[Proof of \ref{item:EM-tor}]
By~\ref{item:EM-R}, $\operatorname{Tor}_1^R(M,M)$ contains
$\operatorname{Tor}_1^R(H^1,H^1)$ as a direct summand, and $H^1$ is the
sum of two copies of $R/hR$ placed in degree one. The complex
$R[-2]\xrightarrow{h}R\to R/hR\to0$, with $R[-2]$ the module $R$ with
degrees raised by two, extends to a graded free resolution of $R/hR$,
because the kernel of multiplication by $h$ on $R$ is $hR$.
After tensoring with $R/hR$ the map $h$ becomes zero, so
$\operatorname{Tor}_1^R(R/hR,R/hR)$ is a copy of $R/hR$ in degree two,
and $\operatorname{Tor}_1^R(H^1,H^1)$ is four copies of it in degree
$1+1+2=4$.
\end{proof}

\begin{proof}[Proof of \ref{item:EM-vanish}]
The fibre product is finite over $E$, hence of dimension one; its
reduced subscheme is a reduced separated scheme of finite type and of
dimension one, in general reducible (a variety in the sense of the
source), which has the same
\'etale site \cite[pp.~22--23]{MilneLEC2013}, and a variety of dimension one
over an algebraically closed field has cohomological dimension at most
two \cite[Theorem~15.1, p.~105]{MilneLEC2013}. Passing to the limit over
$\Z/\ell^n$ and tensoring with $\Q_\ell$ gives the claim.
\end{proof}

In the derived tensor product $M\otimes^{\mathbf L}_RM$, the component
of $\operatorname{Tor}_1^R(M,M)$ of internal degree four sits in total
degree $4-1=3$; so the cohomology of $E\times_{\mathbb{P}^1}E$ is not
that of the derived tensor product over $R$.

It is natural to expect that the derived tensor product over the
cohomology ring fails only because it forgets the cochains, and that
over the \'etale cochain algebras the formula for
$E\times_{\mathbb{P}^1}E$ is repaired; in fact the derived tensor
product of the cochain algebras has cohomology of the same dimensions
as the one over $R$, nonzero in every degree.

\begin{definition}\label{def:dg-tor}
Let $k$ be a field.
\begin{enumerate}[label=(\alph*)]
\item For a differential graded $k$-algebra $A$, cohomologically
  graded, a right differential graded $A$-module $B$, and a left one
  $B'$, the \emph{derived tensor product} $B\otimes^{\mathbf L}_AB'$ is
  the complex $B\otimes_A\operatorname{total}(P^\bullet)$ for a proper
  projective resolution $P^\bullet$ of $B'$
  \cite[pp.~238--239]{McCleary2001UsersGuide}; its cohomology is the
  differential $\operatorname{Tor}^A(B,B')$
  \cite[Definition~7.5, p.~240]{McCleary2001UsersGuide}. A graded
  algebra is regarded as a differential graded algebra with zero
  differential.
\item Let $E$, $x$, and $\ell$ be as in Proposition~\ref{prop:kun-EM},
  with $\ell\ne2$, and let $k$ be $\F_\ell$ or $\Q_\ell$. A \emph{dg
  model of $x^*$ over $k$} is a homomorphism $A\to B$ of differential
  graded $k$-algebras together with isomorphisms of graded algebras
  $H^\bullet(A)\cong H^\bullet(\mathbb{P}^1_{\overline\F_q},k)$ and
  $H^\bullet(B)\cong H^\bullet(E_{\overline\F_q},k)$ that carry the map
  induced on cohomology to $x^*$; $B$ is a right and a left $A$-module
  through $A\to B$.
\end{enumerate}
\end{definition}

The \'etale cochain algebras $C^\bullet(\mathbb{P}^1_{\overline\F_q},k)
\to C^\bullet(E_{\overline\F_q},k)$, with the pullback along $x$, are
meant as dg models in this sense; the statements below hold for every
dg model.

\begin{proposition}\label{prop:EM-bar}
Let $k$ be a field, $A\to B$ a homomorphism of differential graded
$k$-algebras, $R=H^\bullet(A)$, and $M=H^\bullet(B)$. Suppose that
$R=k\oplus kh$ with $h$ of degree two, that $M^j=0$ for
$j\notin\{0,1,2\}$, that $M^0=k\cdot1$, that $\dim M^1=2$, and that
$M^2$ is one-dimensional and spanned by the image of $h$. Then
$H^n(B\otimes^{\mathbf
L}_AB)$ has dimension $1$, $4$, $5$ for $n=0,1,2$ and dimension $4$ for
every $n\ge3$.
\end{proposition}

\begin{proof}
The image of $h$ in $M^2$ is $h\cdot1=1\cdot h$, and $h$ kills $M^1$
and $M^2$ from either side, since $M^3=M^4=0$; so the left and right
$R$-module structures on $M$ agree, the map $R\to M$, $r\mapsto r\cdot1$,
is injective with image $M^0\oplus M^2$, and $M=R\oplus M^1$ with
$M^1$ the sum of two copies of $R/hR$ placed in degree one. As in the
proof of Proposition~\ref{prop:kun-EM}\ref{item:EM-tor}, the kernel of
multiplication by $h$ on $R$ is $hR$, so
$\cdots\xrightarrow{h}R[-4]\xrightarrow{h}R[-2]\xrightarrow{h}R\to
R/hR\to0$ is a graded free resolution, and
$\operatorname{Tor}^R_i(R/hR,R/hR)$ is one-dimensional, in internal
degree $2i$, for every $i\ge0$. Since $\operatorname{Tor}^R_i$ vanishes
for $i\ge1$ when one argument is $R$, the module
$\operatorname{Tor}_0^R(M,M)=M\otimes_RM$ has dimensions $1,4,5$ in
internal degrees $0,1,2$ and vanishes in other degrees, and for $i\ge1$
the module $\operatorname{Tor}^R_i(M,M)=\operatorname{Tor}^R_i(M^1,M^1)$
is four-dimensional and concentrated in internal degree $2+2i$.

The filtration of $B\otimes_A\operatorname{total}(P^\bullet)$ by the
resolution degree, $\{0\}=F^1\subset F^0\subset F^{-1}\subset\cdots$,
is exhaustive, and its spectral sequence has
$E_2^{-i,j}=\operatorname{Tor}_i^R(M,M)^j$ (internal degree $j$)
\cite[Theorem~7.6, p.~241]{McCleary2001UsersGuide} and differentials
$d_r:E_r^{-i,j}\to E_r^{-i+r,j-r+1}$
\cite[Theorem~2.6, p.~33]{McCleary2001UsersGuide}. Since $F^1=\{0\}$
the filtration is weakly convergent \cite[p.~62]{McCleary2001UsersGuide},
so $E_\infty^{-i,j}$ is the graded piece of an exhaustive filtration of
$H^{j-i}(B\otimes^{\mathbf L}_AB)$
\cite[Theorem~3.2, p.~63]{McCleary2001UsersGuide}. The $E_2$-term is
nonzero only at $(0,0)$, $(0,1)$, $(0,2)$, and $(-i,2+2i)$, $i\ge1$. For
$r\ge2$, $d_r$ leaves the column $i=0$ for a column $-i+r>0$, where
$E_r$ is zero; from $(-i,2+2i)$ it reaches $(-(i-r),3+2i-r)$, which for
$i=r$ has internal degree $3+i>2$ and for $i-r\ge1$ would require
$3+2i-r=2+2(i-r)$, that is $r=-1$. So every $d_r$, $r\ge2$, vanishes,
$E_\infty=E_2$, and $\dim H^n(B\otimes^{\mathbf L}_AB)=
\sum_i\dim E_2^{-i,n+i}$. For $n\le2$ only $i=0$ contributes, and for
$n\ge3$ only $i=n-2$.
\end{proof}

\begin{proposition}\label{prop:EM-cochain}
Let $E$, $x$, and $\ell$ be as in Proposition~\ref{prop:kun-EM}, with
$\ell\ne2$, let $k$ be $\F_\ell$ or $\Q_\ell$, and let $A\to B$ be a
dg model of $x^*$ over $k$. Then $H^n(B\otimes^{\mathbf L}_AB)$ has
dimension $1$, $4$, $5$ for $n=0,1,2$ and dimension $4$ for every
$n\ge3$, whereas $H^n((E\times_{\mathbb{P}^1}E)_{\overline\F_q},k)=0$
for $n\ge3$. In particular $B\otimes^{\mathbf L}_AB$ is not
quasi-isomorphic to any complex whose cohomology is
$H^\bullet((E\times_{\mathbb{P}^1}E)_{\overline\F_q},k)$.
\end{proposition}

\begin{proof}
We verify the hypotheses of Proposition~\ref{prop:EM-bar}. For
$k=\Q_\ell$ the dimensions of $R$ and $M$ are those of
Proposition~\ref{prop:kun-EM}\ref{item:EM-R}. Fix compatible isomorphisms
$\mu_{\ell^n}\cong\Z/\ell^n$ over $\overline\F_q$. For $k=\F_\ell$, a
complete
connected nonsingular curve of genus $g$ over an algebraically closed
field has $H^j(\cdot,\mu_\ell)$ of dimensions $1,2g,1$ for $j=0,1,2$
and zero for $j>2$ \cite[Proposition~14.2, p.~90]{MilneLEC2013}. In
both cases $h^2\in H^4=0$, and it remains to show that
$x^*:H^2(\mathbb{P}^1_{\overline\F_q},k)\to H^2(E_{\overline\F_q},k)$ is
nonzero. For a curve $X$ as above and $n\ge1$, the Kummer sequence
identifies $H^2(X,\mu_{\ell^n})$ with $\operatorname{Pic}(X)/\ell^n$
(proof of \cite[Proposition~14.2, p.~90]{MilneLEC2013}), and the degree
identifies this with $\Z/\ell^n$, because $\operatorname{Pic}^0(X)$ is
$\ell$-divisible \cite[Proposition~14.1, p.~89]{MilneLEC2013}. The
Kummer sequence is functorial, so $x^*$ becomes the pullback of line
bundles, which multiplies degrees by two: the pullback of the divisor
of a closed point is its scheme-theoretic fibre, of length
$\deg x=2$ since $x$ is finite and flat. As $\ell$ is odd, $x^*$ is an
isomorphism on $H^2(\cdot,\mu_{\ell^n})$ for every $n$; take $n=1$ for
$\F_\ell$, and pass to the limit and tensor with $\Q_\ell$
\cite[p.~126]{MilneLEC2013} for $\Q_\ell$. The vanishing for $n\ge3$ is
Proposition~\ref{prop:kun-EM}\ref{item:EM-vanish}, whose proof gives it
for $\F_\ell$-coefficients before the passage to the limit.
\end{proof}

So the cohomology of the cochain-level derived tensor product has the
dimensions of the derived tensor product over $R$
(Proposition~\ref{prop:EM-bar}), and the failure of the formula for
$E\times_{\mathbb{P}^1}E$ is not a matter of the coefficient algebra.

\begin{remark}\label{rem:N-EM-scope}
It is natural to expect, from Propositions~\ref{prop:kun-EM}
and~\ref{prop:EM-cochain}, that no derived K\"unneth formula over a base
holds in any form. In fact the failure is one of the map $x$, not of
the base: for every cochain enhancement with $\F_\ell$-coefficients
(Definition~\ref{def:EM-setting}) the cochains of $Y\times_SZ$ are
the derived tensor product of those of $Y$ and $Z$ over those of $S$
whenever $Rf_*\F_\ell$ lies in the thick subcategory generated by the
constant sheaf (Proposition~\ref{prop:EM}), for instance for
$Y=\mathbb{P}^1\times F$ over the same base $\mathbb{P}^1$ and every
proper $Z\to\mathbb{P}^1$, $x$ included
(Proposition~\ref{prop:EM-product}), and $x$ violates this hypothesis
(Proposition~\ref{prop:EM-x}). Nor do
Propositions~\ref{prop:kun-over-Z}, \ref{prop:kun-EM},
and~\ref{prop:EM-cochain}, which concern the tensor product over a
point and over the cohomology or the cochains of a base, exclude a
formula of that kind for the arithmetic surface, which is a product
over $\Spec\F_1$ and not a fibre product over $\Spec\Z$, the square of
$\Spec\Z$ over itself being $\Spec\Z$.
\end{remark}

\begin{definition}\label{def:kun-notation}
Let $K$ be a number field and $D_K=(V_K^\bullet,\Theta,\cup,\Tr_K,*)$ a
Weil datum of type $(1,0)$ on $\Xbar_K$.
\begin{enumerate}[label=(\alph*)]
\item $Z(K)$ is the set of zeros of $\xi_K$.
\item $V_K=V_K^1$ is the degree-one part, with the pairing
  $B_K=B_1$ and the star form $h_1$, and $V_{K,\rho}=
  V_{K,\Theta\sim\rho}$ for $\rho\in\C$.
\item The \emph{space of algebraic tensors} is $N_K=\bigoplus_\rho
  V_{K,\rho}\otimes V_{K,1-\rho}\subset V_K\otimes V_K$.
\item The \emph{tensor star form} is the sesquilinear form
  $h_\otimes=h_1\otimes h_1$ on $V_K\otimes V_K$, $h_\otimes(x\otimes
  y,x'\otimes y')=h_1(x,x')h_1(y,y')$.
\item For $\rho$ with $\Re\rho=\tfrac12$, the \emph{endomorphism map}
  is $\Phi_\rho:V_{K,\rho}\otimes V_{K,1-\rho}\to
  \operatorname{End}(V_{K,\rho})$, $\Phi_\rho(x\otimes y)(z)=
  B_K(z,y)x$, and when $h_1$ is nondegenerate on $V_{K,\rho}$,
  $u\mapsto u^\star$ is the adjoint of $u\in\operatorname{End}
  (V_{K,\rho})$ with respect to $h_1|_{V_{K,\rho}}$.
\end{enumerate}
\end{definition}

The same notation is used for a second field $L$. When
$\Re\rho=\tfrac12$ the space $V_{K,1-\rho}$ is $V_{K,\bar\rho}$.

It is natural to expect that the tensor product of the cohomologies of
two arithmetic curves is again a spectral core, so that the axioms make
sense on the arithmetic surface; in fact duality forces an
infinite-dimensional generalized eigenspace, and under a mild hypothesis
on the zeros also a finite accumulation point of the spectrum.

\begin{theorem}\label{thm:kun-square}
Let $K,L$ be number fields, let $D_K$ and $D_L$ be Weil data of type
$(1,0)$ on $\Xbar_K$ and $\Xbar_L$ satisfying \ref{ax:AW1} for
$U=\mathbb 1$, \ref{ax:AW2}, and \ref{ax:AW3} for $(\Lambda_K,\mathbb 1)$
and $(\Lambda_L,\mathbb 1)$, and let $S$ be a surface datum over
$(D_K,D_L)$ with $\Tr_S(\mathrm{pr}_1^*x\cup\mathrm{pr}_2^*y)=
\Tr_K(x)\Tr_L(y)$ for $x\in V_K^2$, $y\in V_L^2$, and use the notation
of Definition~\ref{def:kun-notation}.
\begin{enumerate}[label=(\roman*)]
\item\label{item:ks-inj} $\kappa$ is injective on $V_K\otimes V_L$, and
  $\Theta_S\kappa=\kappa(\Theta\otimes1+1\otimes\Theta)$.
\item\label{item:ks-mult} The generalized eigenspace of $\Theta_S$ on
  $H^2$ at $1$ contains $\kappa\bigl(\bigoplus_{\rho\in Z(K)\cap Z(L)}
  V_{K,\rho}\otimes V_{L,1-\rho}\bigr)$, of dimension
  $\sum_{\rho\in Z(K)\cap Z(L)}\ord_\rho\xi_K\ord_\rho\xi_L$, which is
  infinite if $K=L$.
\item\label{item:ks-acc} If $K=L$ and the number of distinct imaginary
  parts in $(0,T]$ of zeros of $\xi_K$ is not $O(T)$ as $T\to\infty$,
  the set of eigenvalues of $\Theta_S$ on $H^2$ has a finite
  accumulation point, with imaginary part $0$ and real part in $[0,2]$.
\item\label{item:ks-core} If $Z(K)\cap Z(L)$ is infinite, in particular
  if $K=L$, then $(H^2,\Theta_S)$ is not a spectral core; so no surface
  datum over $(D_K,D_L)$ with multiplicative trace is a Weil datum of
  type $(2,0)$, and neither \ref{ax:AW2} nor AW-det can be imposed on
  it.
\end{enumerate}
\end{theorem}

\begin{proof}[Proof of \ref{item:ks-inj}]
$\Theta_S$ is a derivation commuting with the $\mathrm{pr}_i^*$, which
gives the intertwining. For $x,x'\in V_K$ and $y,y'\in V_L$, associativity
and graded commutativity give $\kappa(x\otimes y)\cup
\kappa(x'\otimes y')=-\mathrm{pr}_1^*(x\cup x')\cup\mathrm{pr}_2^*
(y\cup y')$, so $\Tr_S(\kappa(x\otimes y)\cup\kappa(x'\otimes y'))=
-B_{K}(x,x')B_{L}(y,y')$, with $B_K,B_L$ the degree-one duality
pairings. Let $z=\sum z_{\rho\rho'}$ with $z_{\rho\rho'}\in
V_{K,\rho}\otimes V_{L,\rho'}$ and $\kappa(z)=0$. For $u\in V_{K,1-\rho}\otimes
V_{L,1-\rho'}$, Proposition~\ref{prop:B-support} gives
$0=\Tr_S(\kappa(z)\cup\kappa(u))=-(B_K\otimes B_L)(z_{\rho\rho'},u)$.
By \ref{ax:AW3}, $B_K\otimes B_L$ is a perfect pairing of the
finite-dimensional spaces $V_{K,\rho}\otimes V_{L,\rho'}$ and
$V_{K,1-\rho}\otimes V_{L,1-\rho'}$, so $z_{\rho\rho'}=0$.
\end{proof}

\begin{proof}[Proof of \ref{item:ks-mult}]
By the proof of Theorem~\ref{thm:rh}\ref{item:rh-divisor}, the
multiplicity of $\rho$ on $V_K$ is $\ord_\rho\xi_K$, and likewise for
$V_L$; and $\ord_{1-\rho}\xi_L=\ord_\rho\xi_L$ by
Proposition~\ref{prop:xiK}. By Proposition~\ref{prop:tensor-eigen} the
blocks $V_{K,\rho}\otimes V_{L,1-\rho}$ lie in the generalized eigenspace of
$\Theta\otimes1+1\otimes\Theta$ at~$1$, and by~\ref{item:ks-inj} their
images are independent and lie in that of $\Theta_S$. The count of
zeros \cite[p.~1]{HasanalizadeShenWong2021CountingZeros}, of order
$(n_KT/\pi)\log T$ up to $T$ with $n_K=[K:\Q]$, shows that $Z(K)$ is
infinite.
\end{proof}

\begin{proof}[Proof of \ref{item:ks-acc}]
Let $\gamma_1<\gamma_2<\cdots$ be the distinct positive imaginary parts
of zeros of $\xi_K$. If $\gamma_{n+1}-\gamma_n\ge\delta>0$ for all
$n\ge n_0$, at most $n_0+T/\delta$ of them lie in $(0,T]$; so there are
$n_1<n_2<\cdots$ with $g_k=\gamma_{n_k+1}-\gamma_{n_k}$ positive and
strictly decreasing to~$0$. Choose zeros $\rho_k,\rho'_k$ with imaginary
parts $\gamma_{n_k}$, $\gamma_{n_k+1}$. Then $1-\rho'_k\in Z(K)$, and by
Proposition~\ref{prop:tensor-eigen} and~\ref{item:ks-inj},
$\lambda_k=\rho_k+1-\rho'_k$ is an eigenvalue of $\Theta_S$ on
$\kappa(V_{K,\rho_k}\otimes V_{K,1-\rho'_k})\subset H^2$. The real parts of
zeros lie in $[0,1]$ (Proposition~\ref{prop:xiK}), so
$\Re\lambda_k\in[0,2]$, while $\Im\lambda_k=-g_k$ are distinct and tend
to~$0$. A bounded infinite set of points has an accumulation point.
\end{proof}

\begin{proof}[Proof of \ref{item:ks-core}]
A spectral core has finite-dimensional generalized eigenspaces, and a
spectral cut, hence AW-det, is defined only for spectral cores;
\ref{ax:AW2} is a condition on a Weil datum, whose components are
spectral cores.
\end{proof}

The obstruction of~\ref{item:ks-mult} is the one Manin records for the
tensor square of $\zeta$, whose zeros $\rho$ and $1-\rho$ make the
eigenvalue~$1$ of infinite multiplicity
\cite[p.~132, difficulty~(i)]{Manin1995Lectures} and Deninger notes
for the K\"unneth cohomology of $\Spec\Z\boxtimes\Spec\Z$
\cite[(7.32), p.~42]{Deninger1994Motivic}; here it is forced by
the cohomological formalism (pullbacks, cup product, a multiplicative
trace, and duality), not by a choice of regularization. The hypothesis
of~\ref{item:ks-acc} holds, for instance, if a positive proportion of
the zeros are simple, since the count of zeros up to height $T$ is of
order $T\log T$; the multiplicity bound $O(\log T)$ for zeros alone only
gives a number of distinct imaginary parts at least of order $T$, which
does not exclude $O(T)$, and the
accumulation point is not established here without the hypothesis.

It is natural to expect that Kurokawa's cone, which does have finite
multiplicities for $\zeta$ in the known range, is the cohomology of the
arithmetic surface; in fact it is not closed under the cup products
that the surface must carry, and it contains none of the algebraic
classes attached to nonreal zeros.

\begin{proposition}\label{prop:kun-cone}
In the situation of Theorem~\ref{thm:kun-square} with $K=L$, let $W_c$
be the Kurokawa cone of $V_K\otimes V_K$
(Definition~\ref{def:kurokawa-cone}).
\begin{enumerate}[label=(\roman*)]
\item\label{item:kc-alg} Every subalgebra of $H^\bullet$ containing
  $\mathrm{pr}_1^*V_K$ and $\mathrm{pr}_2^*V_K$ contains $\kappa(V_K\otimes
  V_K)$, and $\kappa(W_c)\ne\kappa(V_K\otimes V_K)$.
\item\label{item:kc-classes} $W_c\cap\bigoplus_\rho V_{K,\rho}\otimes
  V_{K,1-\rho}=\bigoplus_{\rho\in Z(K)\cap\R}V_{K,\rho}\otimes V_{K,1-\rho}$,
  which is finite-dimensional.
\end{enumerate}
\end{proposition}

\begin{proof}[Proof of \ref{item:kc-alg}]
Such a subalgebra contains every product $\mathrm{pr}_1^*x\cup
\mathrm{pr}_2^*y=\kappa(x\otimes y)$. The function $\xi_K$ is real on
$\R$, so $Z(K)$ is stable under conjugation; it is infinite
(Theorem~\ref{thm:kun-square}\ref{item:ks-mult}) and its real points
form a discrete subset of $[0,1]$, so it contains some $\rho$ with
$\Im\rho>0$, and $V_{K,\rho}\otimes V_{K,\bar\rho}$ is a nonzero block
outside $W_c$; $\kappa$ is injective.
\end{proof}

\begin{proof}[Proof of \ref{item:kc-classes}]
$\Im(1-\rho)=-\Im\rho$, so $\Im\rho\cdot\Im(1-\rho)\ge0$ exactly when
$\rho$ is real, and $Z(K)\cap\R$ is finite.
\end{proof}

It is natural to expect that analytic compactness, the compact
resolvent of the self-adjoint generators required by \ref{ax:AW6},
survives the tensor product; in fact it survives for spectra bounded on
one side and fails for the symmetric spectra that duality produces.

\begin{definition}\label{def:tensor-diag}
Let $\cH_1,\cH_2$ be Hilbert spaces with orthonormal bases
$(e_i)_{i\in I}$ and $(f_k)_{k\in I'}$, and let $T_1=\diag(\tau_i)$ and
$T_2=\diag(\tau'_k)$ be real diagonal operators
(Definition~\ref{def:diag-op}).
\begin{enumerate}[label=(\alph*)]
\item The \emph{tensor sum} $T_1\boxplus T_2$ is the diagonal operator
  $\diag(\tau_i+\tau'_k)$ on the Hilbert tensor product
  $\cH_1\widehat\otimes\cH_2$ with orthonormal basis
  $(e_i\otimes f_k)$.
\item A self-adjoint operator $T$ has \emph{compact resolvent} if
  $(T-i)^{-1}$ is compact; its \emph{counting function} $N_T(R)$ is the
  number of eigenvalues of absolute value at most $R$, with multiplicity.
\end{enumerate}
\end{definition}

\begin{proposition}\label{prop:kun-diag}
Let $T_1$, $T_2$, and $T=T_1\boxplus T_2$ be as in
Definition~\ref{def:tensor-diag}.
\begin{enumerate}[label=(\roman*)]
\item\label{item:kd-sa} $T$ is self-adjoint and agrees with
  $T_1\otimes1+1\otimes T_2$ on the algebraic tensor product of the
  spans of the $e_i$ and the $f_k$.
\item\label{item:kd-cr} $T$ has compact resolvent if and only if, for
  every $R>0$, only finitely many pairs $(i,k)$ satisfy
  $|\tau_i+\tau'_k|\le R$.
\item\label{item:kd-semi} If $T_1$ and $T_2$ have compact resolvent and
  $\tau_i\ge-c_1$, $\tau'_k\ge-c_2$ for all $i,k$, with $c_1,c_2\ge0$,
  then $T$ has compact resolvent and
  $N_T(R)\le N_{T_1}(R+c)N_{T_2}(R+c)$ with $c=c_1+c_2$.
\item\label{item:kd-sym} If infinitely many pairs $(i,k)$ satisfy
  $\tau_i+\tau'_k=0$, then $0$ is an eigenvalue of $T$ of infinite
  multiplicity and $T$ does not have compact resolvent.
\item\label{item:kd-sharp} Both bounds $N_T(R)\le N_{T_1}(R+c_2)
  N_{T_2}(R+c_1)$ with $c_1,c_2\ge0$, and $N_T(R)\le N_{T_1}(R+c)
  N_{T_2}(R+c)$ with $c=c_1+c_2$ when some $c_i<0$, fail in general:
  with $\{\tau_i\}=\{-10\}\cup\{1,2,3,\dots\}$ and
  $\{\tau'_k\}=\{10,11,12,\dots\}$ (so $c_1=10$, $c_2=0$) the first
  fails at $R=0$, and with $\{\tau_i\}=\{10,11,\dots\}$ and
  $\{\tau'_k\}=\{0,1,2,\dots\}$ (so $c_1=-10$, $c_2=0$) the second
  fails at $R=10$.
\end{enumerate}
\end{proposition}

\begin{proof}[Proof of \ref{item:kd-sa}]
The proof of Proposition~\ref{prop:hilbert} that a diagonal operator
with real entries is self-adjoint applies verbatim, and both operators
send $e_i\otimes f_k$ to $(\tau_i+\tau'_k)e_i\otimes f_k$.
\end{proof}

\begin{proof}[Proof of \ref{item:kd-cr}]
$(T-i)^{-1}$ is diagonal with entries $c_{ik}=(\tau_i+\tau'_k-i)^{-1}$.
A diagonal operator $\diag(c_n)$ is compact if and only if, for every
$\varepsilon>0$, only finitely many $|c_n|$ are at least $\varepsilon$:
if so, it is the norm limit of its finite-rank truncations; if not,
infinitely many orthonormal $e_n$ have $\|Ce_n-Ce_m\|\ge\varepsilon
\sqrt2$, so their images have no convergent subsequence. Finally
$|c_{ik}|\ge\varepsilon$ if and only if
$|\tau_i+\tau'_k|^2\le\varepsilon^{-2}-1$.
\end{proof}

\begin{proof}[Proof of \ref{item:kd-semi}]
If $|\tau_i+\tau'_k|\le R$, then $\tau_i\le R-\tau'_k\le R+c_2$ and
$\tau_i\ge-c_1$, so $|\tau_i|\le R+c$ because $c_1,c_2\ge0$; likewise
$|\tau'_k|\le R+c$. The argument of~\ref{item:kd-cr}, applied to the
diagonal operators $T_1$ and $T_2$ themselves, shows that the numbers of
such $i$ and of such $k$ are at most $N_{T_1}(R+c)$ and $N_{T_2}(R+c)$,
both finite; apply~\ref{item:kd-cr} to $T$.
\end{proof}

\begin{proof}[Proof of \ref{item:kd-sym}]
The vectors $e_i\otimes f_k$ with $\tau_i+\tau'_k=0$ are orthonormal
eigenvectors with eigenvalue $0$; apply~\ref{item:kd-cr} with any $R$.
\end{proof}

\begin{proof}[Proof of \ref{item:kd-sharp}]
In the first example the pair $(-10,10)$ gives $N_T(0)\ge1$, while
$N_{T_1}(0)=0$ because no $\tau_i$ vanishes. In the second, the pair
$(10,0)$ gives $N_T(10)\ge1$, while $R+c=0$ and $N_{T_1}(0)=0$.
\end{proof}

\begin{proposition}\label{prop:kun-AW6}
Let $K$ be a number field and let a Weil datum of type $(1,0)$ on
$\Xbar_K$ satisfy \ref{ax:AW1}--\ref{ax:AW4} for $(\Lambda_K,\mathbb
1)$. Let $T_1$ be the diagonal operator of Proposition~\ref{prop:hilbert}
for $(V^1,\Theta_1,h_1)$ with $j=1$. Then $0$ is an eigenvalue of
infinite multiplicity of $T_1\boxplus T_1$; in particular
$T_1\boxplus T_1$ has no compact resolvent, and the requirement of
\ref{ax:AW6} that the spectrum have finite multiplicities fails for it.
\end{proposition}

\begin{proof}
By Theorem~\ref{thm:rh} every eigenvalue is
$\rho=\tfrac12+i\tau$ with $\rho\in Z(K)$ and $\tau$ an eigenvalue of
$T_1$. Since $\xi_K$ is real on $\R$ and $\xi_K(s)=\xi_K(1-s)$,
$\bar\rho=\tfrac12-i\tau$ is also a zero, so $-\tau$ is an eigenvalue of
$T_1$; there are infinitely many zeros
(Theorem~\ref{thm:kun-square}\ref{item:ks-mult}), hence infinitely many
pairs of basis vectors with $\tau_i+\tau_k=0$. Apply
Proposition~\ref{prop:kun-diag}\ref{item:kd-sym}.
\end{proof}

On $V^1\otimes V^1$ the closure of the flow generator is
$1+i(T_1\boxplus T_1)$, the weight-two normalization of \ref{ax:AW5};
the eigenvalue $0$ of $T_1\boxplus T_1$ is the eigenvalue $1$ of
Theorem~\ref{thm:kun-square}\ref{item:ks-mult}. The failure is not a
problem of domains (the operator is diagonal) but of symmetry: the pairs
$(\rho,1-\rho)$ are the shadow of duality on the compact object, and by
Proposition~\ref{prop:kun-diag}\ref{item:kd-semi} a spectrum bounded on
one side, such as the archimedean tower $\{-2n:n\ge0\}$ of
$\GR$, would keep a compact resolvent and a polynomial counting function.

It is natural to expect that, after logarithmic suspension over $\F_q$,
the suspension of a product is the tensor product of the suspensions;
in fact the tensor product over $\C$ has infinite multiplicities, and the
correct statement is the tensor product over the periodic unit
(Theorem~\ref{thm:kun-Fq}).

\begin{proposition}\label{prop:kun-Fq-C}
Let $q>1$, let $V,W$ be nonzero finite-dimensional complex vector spaces
with invertible linear maps $\mathrm{F},\mathrm{G}$, and let
$\widehat V,\widehat W$ be their logarithmic suspensions. Then every
generalized eigenspace of $\Theta\otimes1+1\otimes\Theta$ on
$\widehat V\otimes_\C\widehat W$ is either zero or infinite-dimensional.
\end{proposition}

\begin{proof}
If $\widehat V_{\rho_1}\otimes\widehat W_{\rho_2}$ is a nonzero block,
so is $\widehat V_{\rho_1+2\pi ik/\log q}\otimes
\widehat W_{\rho_2-2\pi ik/\log q}$ for every $k\in\Z$, and all of these
lie in the generalized eigenspace at $\rho_1+\rho_2$
(Proposition~\ref{prop:tensor-eigen}).
\end{proof}

It is natural to expect the absolute product of two closed orbits, the
building blocks of Deninger's picture in which the prime $p$
contributes the periodic unit of period $\log p$
\cite[pp.~167--168 and p.~181]{Deninger1998Analogies}, to be again a
well-behaved
spectral object; in fact it has infinite multiplicities for one prime,
dense spectrum for two, and Kurokawa's cone repairs the first only at
the price of unbounded multiplicities.

\begin{proposition}\label{prop:kun-orbits}
Let $p\ne p'$ be primes and write $\mathbb 1_p$ for the $p$-periodic
unit of Definition~\ref{def:unit}.
\begin{enumerate}[label=(\roman*)]
\item\label{item:ko-same} Every eigenvalue $2\pi ik/\log p$ of
  $\mathbb 1_p\otimes\mathbb 1_p$ has infinite multiplicity.
\item\label{item:ko-dense} $\mathbb 1_p\otimes\mathbb 1_{p'}$ has the
  eigenvalues $2\pi i(m/\log p+n/\log p')$, $m,n\in\Z$, each simple,
  and they are dense in $i\R$.
\item\label{item:ko-cone} On the Kurokawa cone of
  $\mathbb 1_p\otimes\mathbb 1_p$ the eigenvalue $2\pi ik/\log p$ has
  multiplicity $|k|+1$.
\end{enumerate}
\end{proposition}

\begin{proof}[Proof of \ref{item:ko-same}]
$\mathbb 1_p$ is the suspension of $(\C,1)$ (Proposition~\ref{prop:unit});
apply Proposition~\ref{prop:kun-Fq-C}.
\end{proof}

\begin{proof}[Proof of \ref{item:ko-dense}]
If $m/\log p+n/\log p'=m'/\log p+n'/\log p'$, then
$(m-m')\log p'=(n'-n)\log p$, that is $p'^{m-m'}=p^{n'-n}$, which by
unique factorization forces $m=m'$ and $n=n'$. For the same reason
$\theta=\log p/\log p'$ is irrational, so the subgroup $\Z+\theta\Z$ of
$\R$ is not cyclic and hence dense, and the eigenvalues are
$(2\pi i/\log p)(m+n\theta)$.
\end{proof}

\begin{proof}[Proof of \ref{item:ko-cone}]
The cone is spanned by the $v^m\otimes v^n$ with $mn\ge0$, and
$\#\{(m,n):m+n=k,\ mn\ge0\}=|k|+1$.
\end{proof}

It is natural to expect the trace formula of the arithmetic surface to
be the product of the explicit formulas of the factors, as a Lefschetz
formula summing over isolated closed orbits of the diagonal flow; in
fact the closed orbits of the diagonal flow on a product of two closed
orbits of lengths $\log p$ and $\log p'$ exist only for $p=p'$ and are
never isolated.

\begin{proposition}\label{prop:kun-torus}
Let $a,b>0$, let $\mathbb T=(\R/a\Z)\times(\R/b\Z)$ with the flow
$\phi^t(x,y)=(x+t,y+t)$, and let $p,p'$ be primes.
\begin{enumerate}[label=(\roman*)]
\item\label{item:kt-irr} If $a/b\notin\Q$, no point of $\mathbb T$ has
  a periodic orbit.
\item\label{item:kt-rat} If $a/b\in\Q$, every orbit is closed, and no
  closed orbit is isolated among closed orbits.
\item\label{item:kt-primes} $\log p/\log p'\in\Q$ if and only if
  $p=p'$.
\end{enumerate}
\end{proposition}

\begin{proof}[Proof of \ref{item:kt-irr} and \ref{item:kt-rat}]
$\phi^t(x,y)=(x,y)$ if and only if $t\in a\Z\cap b\Z$, which is nonzero
exactly when $a/b\in\Q$, and then does not depend on $(x,y)$. In that
case the orbit of $(x,y)$ meets the circle $\{x\}\times\R/b\Z$ in the
finite set of points $(x,y+ma)$, $m\in\Z$, so for all small $s\ne0$ the
orbit of $(x,y+s)$ is a different closed orbit, arbitrarily close.
\end{proof}

\begin{proof}[Proof of \ref{item:kt-primes}]
$m\log p=n\log p'$ with positive integers $m,n$ means $p^m=p'^n$.
\end{proof}

A Lefschetz formula that sums over isolated closed orbits therefore
does not apply to the product of two closed orbits with the diagonal
flow, and for $p\ne p'$ the spectral side of
Proposition~\ref{prop:kun-orbits}\ref{item:ko-dense} has no geometric
counterpart at all. A trace formula on the arithmetic surface, if there
is one, must be of the type of non-isolated fixed points.

It is natural to expect the archimedean factor of a tensor product of
representations to be obtained from the factors by adding spectra, as
the nonarchimedean zeros of Theorem~\ref{thm:kun-Fq} are; in fact it
depends multiplicatively on the action of complex conjugation, which
the flow does not see.

\begin{proposition}\label{prop:kun-arch}
Let $\chi$ be the nontrivial character of $\mathrm{Gal}(\Q(i)/\Q)$.
Then $L_\infty(\chi,s)=\GR(s+1)$ and
$L_\infty(\chi\otimes\chi,s)=\GR(s)$ (Definition~\ref{def:artin}); the
divisor of poles of $\GR(s)$ is $\sum_{n\ge0}[-2n]$, whereas the sum of
the pole divisor of $\GR(s+1)$ with itself under
$(\rho,\rho')\mapsto\rho+\rho'$ is $\sum_{k\ge0}(k+1)[-2-2k]$.
\end{proposition}

\begin{proof}
Complex conjugation acts on $\Q(i)$ nontrivially, so $\chi$ restricted
to the decomposition group at the real place is the sign character and
$\chi\otimes\chi$ is trivial; Definition~\ref{def:artin}(b) gives the
two factors. The poles of $\Gamma(s/2)$ are at $s=-2n$ and those of
$\Gamma((s+1)/2)$ at $s=-1-2n$, $n\ge0$, all simple, and
$\#\{(m,n):(-1-2m)+(-1-2n)=-2-2k\}=k+1$.
\end{proof}

The two divisors differ (the first contains $0$, the second does not):
the datum lost by the flow is the involution at infinity, a
multiplicative datum. A relative K\"unneth formula over the archimedean
point, if there is one, must keep that involution before taking
invariants.

It is natural to expect that the obstruction to real coefficients is an
accident of degree one (Remark~\ref{rem:R5-import}); in fact a K\"unneth
formula turns the coefficients into a tensor functor, and on the
finite-field calibration no tensor functor of the relevant kind with
real values exists.

\begin{remark}\label{rem:kun-real}
The underlying obstruction is Serre's (Proposition~\ref{prop:R5}).
Deninger observes that, for motives over finite fields, exact faithful
tensor functors into real vector spaces do not exist
\cite[p.~184]{Deninger1998Analogies}. A cohomology with a K\"unneth
formula, restricted to the finite-field objects, is a tensor functor;
so a K\"unneth-compatible theory which restricts there to an exact
faithful tensor functor has non-real coefficients. The quaternionic
refinement $*_1^2=-1$ is compatible with tensor products: on
$V^1\otimes V^1$ the square of $*\otimes*$ is $(-1)(-1)=1$.
\end{remark}

\subsection{The conjecture and compactness}\label{subsec:kun-conj}

The refutations leave one form of the conjecture. Two notions of
compactness enter it and must be kept apart. \emph{Geometric}
compactness is that of $\Xbar_K$, whose infinite places are included:
it is what makes the duality \ref{ax:AW3} and the symmetry
$\rho\leftrightarrow1-\rho$ of the spectrum hold, and it is exactly this
symmetry that produces the infinite multiplicities of
Theorem~\ref{thm:kun-square}. \emph{Analytic} compactness is the compact
resolvent of \ref{ax:AW6}; it survives products of spectra bounded on
one side and fails on the tensor square of every compact arithmetic
curve (Propositions~\ref{prop:kun-diag} and~\ref{prop:kun-AW6}). The two
requirements are in tension on products. One may either $(\alpha)$ keep the
K\"unneth isomorphism and the algebraic classes, and drop the
spectral-core, determinant, and compact-resolvent conditions on the
surface, or $(\beta)$ keep those conditions by passing to a one-sided subspace
such as Kurokawa's cone, which is not closed under cup products and
contains no algebraic classes with nonreal eigenvalues
(Proposition~\ref{prop:kun-cone}). Form $(\beta)$ is refuted; the conjecture
is stated in form $(\alpha)$.

Conjecture~\ref{conj:kunneth} is stated in the introduction.


A K\"unneth formula for an arithmetic cohomology of
$\Spec\Z\boxtimes\Spec\Z$, with flow $\Theta\otimes1+1\otimes\Theta$,
is contemplated in \cite[(7.32), pp.~41--42]{Deninger1994Motivic}, and
positivity on a square is the aim of the Riemann--Roch strategy of
Connes and Consani \cite[\S3, pp.~16--18]{ConnesConsani2018RiemannRoch};
Conjecture~\ref{conj:kunneth} is one precise form of these
expectations. Section~\ref{subsec:f1-search} explains their absolute-base
motivation and distinguishes the constructed sites from the required
cohomological and intersection-theoretic comparisons.


The conjecture is meant for data constructed from arithmetic input
independent of the zeros of the $\xi_K$. That requirement is not a
mathematical condition, and Proposition~\ref{prop:kun-collapse} below
shows that without it the conjecture is equivalent to the Riemann
hypothesis for every $\zeta_K$, exactly as for the axioms on one object
(Proposition~\ref{prop:content}). Its role is the one Grothendieck's
construction of $H^\bullet(C\times C)$ plays over $\F_q$: the surface
must exist before its positivity is proved. Candidate ambient
categories for such a construction include Deninger's dynamical
systems, in which the product carries the diagonal flow, and, for the
local analogue at $p$, Scholze's products of diamonds, which ``give good
meaning to the non-existent products $\Spec\Q_p\times_{\Spec\F_1}
\cdots$'' and relate a hypothetical $\Spec\Z\times\Spec\Z$ to a
geometric realization of the Galois group
\cite[\S6, pp.~13--14]{Scholze2018PadicGeometry}. Relative products are
allowed only over unit-like bases: over $\F_q$ the formula is a theorem
(Theorem~\ref{thm:kun-Fq}), over $\Spec\Z$ it is refuted
(Proposition~\ref{prop:kun-over-Z}), and over the archimedean point the
involution at infinity must be kept (Proposition~\ref{prop:kun-arch}).

{\small
\begin{longtable}{@{}>{\raggedright\arraybackslash}p{0.34\textwidth}
>{\raggedright\arraybackslash}p{0.34\textwidth}
>{\raggedright\arraybackslash}p{0.22\textwidth}@{}}
\caption{Which compactness each conclusion on the arithmetic surface
needs.}\label{tab:compactness}\\
\toprule
Conclusion & Compactness needed & Status\\
\midrule
\endfirsthead
\toprule
Conclusion & Compactness needed & Status\\
\midrule
\endhead
\bottomrule
\endfoot
K\"unneth map injective, intertwining the flows & duality \ref{ax:AW3}
on the factors (geometric) & Thm.~\ref{thm:kun-square}\ref{item:ks-inj}\\
algebraic classes $N_S$ and the Hodge index clause & geometric: the
classes pair $\rho$ with $1-\rho$ & Thm.~\ref{thm:kun-hodge}\\
spectral core, determinants, \ref{ax:AW2} on the surface & finite
multiplicities: contradicted by geometric compactness &
refuted, Thm.~\ref{thm:kun-square}\ref{item:ks-core}\\
compact resolvent (\ref{ax:AW6}) on the surface & analytic;
incompatible with the symmetric spectrum & refuted,
Prop.~\ref{prop:kun-AW6}\\
compact resolvent for one-sided spectra & analytic only & proved,
Prop.~\ref{prop:kun-diag}\ref{item:kd-semi}\\
\end{longtable}
}

Geometric compactness enters the divisor identity only through the
boundary. For an open object $X$ with compactification $\Xbar$ and
boundary $\partial X$ at the infinite places, one expects compactly
supported cohomology $H_c^\bullet(X)$ with a flow and a long exact
sequence relating it to $H^\bullet(\Xbar)$ and $H^\bullet(\partial X)$;
Deninger's picture assigns to the infinite place the tower
$\{-2n:n\ge0\}$ of poles of $\GR$
\cite[pp.~167--168]{Deninger1998Analogies}.
It is natural to expect that the divisor identity for the open object
decides the connecting maps of that sequence, for instance whether the
eigenvalue $0$ of $H^0(\Xbar)$ survives in $H^0_c(X)$; in fact the
alternating spectral divisor of any such sequence is independent of the
maps.

\begin{proposition}\label{prop:kun-open}
Let $0\to E^0\to F^0\to G^0\to E^1\to F^1\to G^1\to E^2\to F^2\to0$ be
an exact sequence of spectral cores and linear maps commuting with the
generators. Then
\[
\sum_{j=0}^2(-1)^j\operatorname{Sp}(E^j)=\sum_{j=0}^2(-1)^j
\operatorname{Sp}(F^j)-\sum_{j=0}^1(-1)^j\operatorname{Sp}(G^j).
\]
In particular, if $F^0=\C$ with generator $0$, and two such sequences
have the same spaces $F^j$ and $G^j$, with $F^0\to G^0$ injective in the
first and zero in the second, then they have the same alternating
divisor $\sum_j(-1)^j\operatorname{Sp}(E^j)$, while
$\operatorname{Sp}(E^0)$ is $0$ in the first case and $[0]$ in the
second.
\end{proposition}

\begin{proof}
A linear map commuting with the generators sends each generalized
eigenspace into the generalized eigenspace for the same eigenvalue, so
the sequence is the direct sum, over $\lambda\in\C$, of its restrictions
to the generalized $\lambda$-eigenspaces. These are exact sequences of
finite-dimensional spaces, so the alternating sum of their dimensions
vanishes; this is the displayed identity, coefficient by coefficient.
In the special case, $E^0$ is the kernel of $F^0\to G^0$.
\end{proof}

So the divisor identity for the compactly supported cohomology of
$\Spec\Z$, with its trivial zeros coming from the boundary tower,
follows from the identities for $\Xbar_\Q$ and $\partial X$ whatever the
connecting maps, and it does not determine whether $0$ lies in the
spectrum of $H^0_c$. If moreover the eigenvalue $0$ occurs in
$H^0(\partial X)$ with multiplicity one, as for the tower $\{-2n:n\ge0\}$,
then, since $\xi(0)\ne0$ so that $0$ does not occur in
$H^1(\Xbar_\Q)$, the multiplicity of $0$ in $H^1_c$ is $1$ minus the
rank of $H^0(\Xbar_\Q)\to H^0(\partial X)$, and the divisor identity
does not determine it either.

\subsection{What holds}\label{subsec:kun-pos}

Over a finite field the formula holds once the suspension is tensored
over the periodic unit, with determinant normalization exactly one.

\begin{theorem}\label{thm:kun-Fq}
Let $q>1$, let $V,W$ be finite-dimensional complex vector spaces with
invertible linear maps $\mathrm{F},\mathrm{G}$, and let $\widehat V$,
$\widehat W$, and $\widehat{V\otimes W}$ be the logarithmic suspensions of
$(V,\mathrm{F})$, $(W,\mathrm{G})$, and $(V\otimes W,\mathrm{F}\otimes
\mathrm{G})$, with the periodicity actions of
Definition~\ref{def:flow-module}(e).
\begin{enumerate}[label=(\roman*)]
\item\label{item:kF-iso} There is an isomorphism of flow modules
  $\widehat{V\otimes W}\to\widehat V\otimes_{\mathbb 1_q}\widehat W$
  sending the copy of $u\otimes w$, $u\in V_{(\alpha)}$,
  $w\in W_{(\beta)}$, in the summand indexed by $\rho$ (with
  $q^\rho=\alpha\beta$) to the class of the copy of $u$ over $\rho_1$
  tensored with the copy of $w$ over $\rho-\rho_1$, for any
  $\rho_1\in\Sigma_q(\alpha)$.
\item\label{item:kF-var} If $X,Y$ are smooth projective geometrically
  connected varieties over $\F_q$, $\ell\nmid q$, and
  $\iota:\overline\Q_\ell\to\C$, then in each degree $n$ the suspended
  datum of $X\times_{\F_q}Y$ is isomorphic, as a flow module, to
  $\bigoplus_{a+b=n}\widehat H^a_\ell(X)\otimes_{\mathbb 1_q}
  \widehat H^b_\ell(Y)$, and its principal-cut determinant is
  $\prod_{a+b=n}\det(1-\mathrm{F}\otimes\mathrm{G}\,q^{-s}\mid
  H^a_\ell(X)\otimes H^b_\ell(Y))$, with normalization term $1$.
\end{enumerate}
\end{theorem}

\begin{proof}[Proof of \ref{item:kF-iso}]
Write $c=2\pi i/\log q$. On $V_{(\alpha)}\otimes W_{(\beta)}$ the commuting maps
$\alpha^{-1}\mathrm{F}\otimes1$ and $1\otimes\beta^{-1}\mathrm{G}$
are unipotent, so $\mathrm{F}\otimes\mathrm{G}$ has the single
eigenvalue $\alpha\beta$ there, and the summand of
$\widehat{V\otimes W}$ over $\rho$ is
$\bigoplus_{\alpha\beta=q^\rho}V_{(\alpha)}\otimes W_{(\beta)}$. The
space $\widehat V\otimes_\C\widehat W$ is the direct sum of the blocks
$\widehat V_{\rho_1}\otimes\widehat W_{\rho_2}$, each a copy of
$V_{(q^{\rho_1})}\otimes W_{(q^{\rho_2})}$. The relations
$vx\otimes y-x\otimes vy$ identify the block $(\rho_1+c,\rho_2)$ with
the block $(\rho_1,\rho_2+c)$ through the identity of
$V_{(\alpha)}\otimes W_{(\beta)}$, and those for $v^{\pm n}$ follow from
those for $v$ (for instance $v^{-1}x\otimes y-x\otimes v^{-1}y$ is minus
the relation for $v$ applied to $v^{-1}x\otimes v^{-1}y$). So the
quotient is the direct sum, over the orbits of $(k c,-k c)$,
$k\in\Z$, on pairs $(\rho_1,\rho_2)$, of one copy of
$V_{(\alpha)}\otimes W_{(\beta)}$; and these orbits correspond
bijectively to the triples $(\rho,\alpha,\beta)$ with $\rho=\rho_1+
\rho_2$, $\alpha=q^{\rho_1}$, $\beta=q^{\rho_2}$. This is the stated map
and it is bijective. On $V_{(\alpha)}\otimes W_{(\beta)}$ the logarithm
of the unipotent part of $(\alpha\beta)^{-1}\mathrm{F}\otimes\mathrm{G}$
is $N_\alpha\otimes1+1\otimes N_\beta$ (the logarithm of a product of
commuting unipotent maps is the sum of the logarithms), so both
generators are $\rho+(N_\alpha\otimes1+1\otimes N_\beta)/\log q$.
\end{proof}

\begin{proof}[Proof of \ref{item:kF-var}]
The cup product induces an isomorphism $\bigoplus_{a+b=n}
H^a(X)\otimes H^b(Y)\to H^n(X\times Y)$ with $\Q_\ell$-coefficients
\cite[Theorem~22.4, p.~136]{MilneLEC2013} (the source states it with
$\Z_\ell$-coefficients and a Tor term, which vanishes after tensoring
with $\Q_\ell$). It is compatible with the Frobenius by naturality of
the K\"unneth map in $X$ and $Y$: the Frobenius of $X\times Y$ is the
product of the Frobenius endomorphisms, and pullbacks along it commute
with the projections and with cup products.
 Apply~\ref{item:kF-iso} and
Theorem~\ref{thm:B}\ref{item:B-det} for $X\times Y$.
\end{proof}

Theorem~\ref{thm:kun-Fq} is the spectral shadow of the strong
K\"unneth theorem for topological periodic cyclic homology: for smooth
and proper dg categories over a perfect field $k$ of characteristic
$p$, $\mathrm{TP}(X)\wedge^{\mathbf L}_{\mathrm{TP}(k)}\mathrm{TP}(Y)
\to\mathrm{TP}(X\otimes_kY)$ is an equivalence
\cite[Theorem~A, pp.~2--3]{BlumbergMandell2017StrongKunneth}, the
tensor product being taken over $\mathrm{TP}(k)$ and not over the
sphere; Hesselholt identifies the Hasse--Weil zeta function with a
regularized determinant on $\mathrm{TP}$
\cite[Theorem~A and footnote~1, p.~2]{Hesselholt2016THHZeta}. The tensor
product over $\mathbb 1_q$ plays the role of the tensor product over
$\mathrm{TP}(k)$; over $\C$ it has infinite multiplicities
(Proposition~\ref{prop:kun-Fq-C}).

Over a base of positive dimension a derived K\"unneth formula holds
under a hypothesis on the map, in the form of an \'etale
Eilenberg--Moore statement, and the double cover of
Proposition~\ref{prop:kun-EM} violates that hypothesis.

\begin{definition}\label{def:EM-setting}
Let $K$ be a separably closed field, $\ell$ a prime invertible in $K$,
and $\Lambda=\F_\ell$. Schemes are of finite type over $K$.
\begin{enumerate}[label=(\alph*)]
\item For a scheme $S$, $\langle\Lambda_S\rangle$ is the smallest
  strictly full triangulated subcategory of $D(S_{\mathrm{et}},\Lambda)$
  that contains the constant sheaf $\Lambda_S$ and is closed under
  direct summands.
\item A \emph{cochain enhancement} assigns to each scheme $T$ a
  differential graded $\Lambda$-algebra $C^\bullet(T)$ whose cohomology
  algebra is $H^\bullet(T,\Lambda)$ with the cup product, and an exact
  functor $R\Gamma_T$ from $D^+(T_{\mathrm{et}},\Lambda)$ to the derived
  category of differential graded $C^\bullet(T)$-modules which lifts
  $R\Gamma(T,-)$ and sends $\Lambda_T$ to $C^\bullet(T)$; to each
  morphism $f:Y\to S$ a homomorphism $f^*:C^\bullet(S)\to C^\bullet(Y)$
  inducing the pullback on cohomology, and natural isomorphisms
  $R\Gamma_S(Rf_*L)\cong R\Gamma_Y(L)$ of $C^\bullet(S)$-modules,
  $C^\bullet(Y)$-modules being $C^\bullet(S)$-modules through $f^*$; and
  to each $S$ a natural transformation
  $\mu_S:R\Gamma_S(L)\otimes^{\mathbf L}_{C^\bullet(S)}R\Gamma_S(L')\to
  R\Gamma_S(L\otimes^{\mathbf L}_\Lambda L')$ of exact functors in
  each variable, which for $L=\Lambda_S$ is the canonical
  identification. Derived tensor products are those of
  Definition~\ref{def:dg-tor}.
\end{enumerate}
\end{definition}

The \'etale cochain algebras with $\F_\ell$-coefficients of
Remark~\ref{rem:N-EM-scope} are meant to form a cochain enhancement in
this sense; Godement resolutions are the usual construction, and their
properties are not verified here. The statements below that involve
$C^\bullet$ assume a cochain enhancement and hold for every one.

\begin{proposition}\label{prop:EM-fibre}
In the setting of Definition~\ref{def:EM-setting}, let $f:Y\to S$ and
$g:Z\to S$ be proper, let $\tilde g:Y\times_SZ\to Y$ be the projection
and $h=f\circ\tilde g$. Then $Rh_*\Lambda\cong Rf_*\Lambda
\otimes^{\mathbf L}_\Lambda Rg_*\Lambda$ in $D(S_{\mathrm{et}},\Lambda)$.
\end{proposition}

\begin{proof}
Let $\tilde f:Y\times_SZ\to Z$ be the other projection. By proper base
change for $g$ along $f$ \stackstag{0F0C},
$f^{-1}Rg_*\Lambda\cong R\tilde g_*\tilde f^{-1}\Lambda=
R\tilde g_*\Lambda$. By the projection formula for the proper
morphism $f$ \stackstag{0F0G}, applied to $\Lambda_Y$ and
$Rg_*\Lambda$,
$Rf_*\Lambda\otimes^{\mathbf L}_\Lambda Rg_*\Lambda\cong
Rf_*(f^{-1}Rg_*\Lambda)\cong Rf_*R\tilde g_*\Lambda=Rh_*\Lambda$.
\end{proof}

\begin{definition}\label{def:EM-kappa}
In the setting of Proposition~\ref{prop:EM-fibre}, fix a cochain
enhancement. The \emph{Eilenberg--Moore comparison map}
\[
\kappa:C^\bullet(Y)\otimes^{\mathbf L}_{C^\bullet(S)}C^\bullet(Z)\to
C^\bullet(Y\times_SZ)
\]
is the composite of the identifications $C^\bullet(Y)\cong
R\Gamma_S(Rf_*\Lambda)$ and $C^\bullet(Z)\cong R\Gamma_S(Rg_*\Lambda)$,
of $\mu_S$ for the pair $(Rf_*\Lambda,Rg_*\Lambda)$, of $R\Gamma_S$
applied to the isomorphism of Proposition~\ref{prop:EM-fibre}, and of
$R\Gamma_S(Rh_*\Lambda)\cong C^\bullet(Y\times_SZ)$.
\end{definition}

\begin{proposition}[{after \cite[Proposition~1.1.9, p.~6, and
Corollary~1.1.10, p.~7]{Lurie2011DAGXIII}}]\label{prop:EM}
In the setting of Definition~\ref{def:EM-kappa}, if $Rf_*\Lambda$ lies
in $\langle\Lambda_S\rangle$, then $\kappa$ is an isomorphism in the
derived category of $C^\bullet(S)$-modules.
\end{proposition}

\begin{proof}
All maps in the definition of $\kappa$ other than $\mu_S$ are
isomorphisms, so it suffices that $\mu_S$ is an isomorphism for the
pair $(Rf_*\Lambda,Rg_*\Lambda)$. Fix $L'=Rg_*\Lambda$ and let
$\mathcal C$ be the class of $L$ in $D^+(S_{\mathrm{et}},\Lambda)$ for
which $\mu_S$ is an isomorphism at $(L,L')$. It contains $\Lambda_S$.
Both sides of $\mu_S$ are exact functors of $L$ and $\mu_S$ is a
natural transformation of exact functors, so a distinguished triangle
with two vertices in $\mathcal C$ gives a morphism of long exact
cohomology sequences which is an isomorphism at two of every three
places, and the third vertex lies in $\mathcal C$ by the five lemma;
$\mathcal C$ is closed under shifts, and under direct summands because
a retract of an isomorphism is an isomorphism. So $\mathcal C$
contains $\langle\Lambda_S\rangle$.
\end{proof}

\begin{proposition}\label{prop:EM-criterion}
In the setting of Definition~\ref{def:EM-kappa}, suppose that only
finitely many cohomology sheaves $R^if_*\Lambda$ are nonzero and that
each has a finite filtration whose graded pieces are finite direct sums
of copies of $\Lambda_S$. Then $Rf_*\Lambda$ lies in
$\langle\Lambda_S\rangle$, and $\kappa$ is an isomorphism for every
proper $g:Z\to S$.
\end{proposition}

\begin{proof}
A short exact sequence of sheaves is a distinguished triangle, so each
$R^if_*\Lambda$ lies in $\langle\Lambda_S\rangle$ by induction on the
length of the filtration. The truncation triangles
$\tau_{\le i-1}Rf_*\Lambda\to\tau_{\le i}Rf_*\Lambda\to
R^if_*\Lambda[-i]$ and induction on the number of nonzero cohomology
sheaves give the first claim; apply Proposition~\ref{prop:EM}.
\end{proof}

Proposition~\ref{prop:EM} is the \'etale form of the Eilenberg--Moore
theorem in Lurie's formulation, with the same thick-subcategory proof
\cite[Proposition~1.1.9, p.~6, and Corollary~1.1.10, p.~7]%
{Lurie2011DAGXIII}; the category $\langle\Lambda_S\rangle$ corresponds
to Lurie's nilpotent objects \cite[Definition~1.1.7, p.~6]%
{Lurie2011DAGXIII}, and the nilpotence condition goes back to Dwyer
\cite{Dwyer1974StrongConvergence}, as reported in
\cite[Theorem~8\textsuperscript{bis}.30]{McCleary2001UsersGuide}.

\begin{proposition}\label{prop:EM-product}
In the setting of Definition~\ref{def:EM-kappa}, let $K$ be
algebraically closed, $F$ a smooth projective connected curve over $K$,
$Y=S\times_KF$, and $f$ the first projection. Then $Rf_*\Lambda$ lies in
$\langle\Lambda_S\rangle$, and for every proper $g:Z\to S$ the map
$\kappa:C^\bullet(S\times_KF)\otimes^{\mathbf L}_{C^\bullet(S)}
C^\bullet(Z)\to C^\bullet(Z\times_KF)$ is an isomorphism.
\end{proposition}

\begin{proof}
Let $a:S\to\Spec K$ and $b:F\to\Spec K$ be the structure morphisms. By
proper base change for $b$ along $a$ \stackstag{0F0C},
$Rf_*\Lambda\cong a^{-1}Rb_*\Lambda$. Since $K$ is separably closed,
$Rb_*\Lambda$ is a complex of $\Lambda$-vector spaces, hence isomorphic
to $\bigoplus_jH^j(F,\Lambda)[-j]$, and $H^j(F,\Lambda)$ has dimension
$1,2g,1$ for $j=0,1,2$ and vanishes for $j>2$, with $g$ the genus of
$F$ \cite[Proposition~14.2, p.~90]{MilneLEC2013}. So $Rf_*\Lambda$ is
a finite direct sum of shifts of $\Lambda_S$, and
$Y\times_SZ=Z\times_KF$; apply Proposition~\ref{prop:EM-criterion}.
\end{proof}

It is natural to expect that the failure of
Propositions~\ref{prop:kun-EM} and~\ref{prop:EM-cochain} is a defect of
the base $\mathbb{P}^1$; in fact over the same base the formula holds for
$\mathbb{P}^1\times F$ and every proper $Z\to\mathbb{P}^1$
(Proposition~\ref{prop:EM-product}), and what fails is the hypothesis
of Proposition~\ref{prop:EM} for the map $x$.

\begin{proposition}\label{prop:EM-x}
Let $E$ and $x$ be as in Proposition~\ref{prop:kun-EM}, $\ell\nmid2q$ a
prime, $\Lambda=\F_\ell$, and consider $x$ over $K=\overline\F_q$. Then
$Rx_*\Lambda$ does not lie in $\langle\Lambda_{\mathbb{P}^1}\rangle$, and
for every cochain enhancement the map $\kappa$ for the pair $(x,x)$ is
not an isomorphism.
\end{proposition}

\begin{proof}
Write $S=\mathbb{P}^1_{\overline\F_q}$. Let $\mathcal D$ be the class of
objects of $D(S_{\mathrm{et}},\Lambda)$ each of whose cohomology
sheaves is isomorphic to $\Lambda_S^{\oplus r}$ for some $r\ge0$. Since
$S$ is connected, $\operatorname{Hom}(\Lambda_S^{\oplus a},
\Lambda_S^{\oplus b})$ is the space of $b\times a$ matrices over
$\Lambda$, so kernels, cokernels, and images of idempotents of such
sheaves are again of this form; and every extension of
$\Lambda_S^{\oplus a}$ by $\Lambda_S^{\oplus b}$ splits, because
$H^1(S,\Lambda)=0$ \cite[Proposition~14.2, p.~90]{MilneLEC2013}. The
long exact sequence of cohomology sheaves of a distinguished triangle
therefore shows that $\mathcal D$ is closed under cones, and it is
closed under shifts and direct summands; so
$\langle\Lambda_S\rangle\subset\mathcal D$. Since $x$ is finite,
$R^qx_*\Lambda=0$ for $q\ge1$ \stackstag{03QP}, so $Rx_*\Lambda=
x_*\Lambda$ and $H^1(S,x_*\Lambda)=H^1(E_{\overline\F_q},\Lambda)$,
which is two-dimensional \cite[Proposition~14.2, p.~90]{MilneLEC2013}.
If $x_*\Lambda$ were isomorphic to $\Lambda_S^{\oplus r}$, this group
would be $H^1(S,\Lambda)^{\oplus r}=0$. So $Rx_*\Lambda\notin
\mathcal D\supset\langle\Lambda_S\rangle$.

For the second claim, $x^*:C^\bullet(S)\to C^\bullet(E_{\overline\F_q})$
is a dg model of $x^*$ over $\F_\ell$ (Definition~\ref{def:dg-tor}),
so by Proposition~\ref{prop:EM-cochain} the source of $\kappa$ has
nonzero cohomology in every degree $n\ge3$, whereas its target has
none.
\end{proof}

Over $\F_q$ the tensor product of Theorem~\ref{thm:kun-Fq} is taken over
the periodic unit, which is the suspension of the cohomology of the
base $\Spec\F_q$; over a base with more cohomology the suspended
derived tensor product is taken over the suspension of that cohomology.

\begin{proposition}\label{prop:kun-susp-tor}
Let $q>1$, let $R$ be a finite-dimensional commutative $\C$-algebra with
an algebra automorphism $\mathrm F_R$, and let $M$, $N$ be
finite-dimensional $R$-modules with invertible linear maps
$\mathrm F_M$, $\mathrm F_N$ such that $\mathrm F_M(rm)=\mathrm F_R(r)
\mathrm F_M(m)$ and $\mathrm F_N(rn)=\mathrm F_R(r)\mathrm F_N(n)$.
Let $\widehat R$, $\widehat M$, $\widehat N$ be the logarithmic
suspensions.
\begin{enumerate}[label=(\roman*)]
\item\label{item:ksT-alg} Through the isomorphisms of
  Theorem~\ref{thm:kun-Fq}\ref{item:kF-iso}, the product of $R$ makes
  $\widehat R$ a flow algebra containing $\mathbb 1_q=\widehat{\C}$,
  and the actions of $R$ make $\widehat M$, $\widehat N$ flow
  $\widehat R$-modules.
\item\label{item:ksT-tor} For every $i\ge0$ there is an isomorphism of
  complex vector spaces
  $\operatorname{Tor}_i^{\widehat R}(\widehat M,\widehat N)\cong
  \widehat{\operatorname{Tor}_i^R(M,N)}$, where $\operatorname{Tor}_i^R$
  carries the automorphism induced by $\mathrm F_M$ and $\mathrm F_N$;
  for $i=0$ it is an isomorphism of flow modules
  $\widehat M\otimes_{\widehat R}\widehat N\cong\widehat{M\otimes_RN}$.
\end{enumerate}
\end{proposition}

\begin{proof}[Proof of \ref{item:ksT-alg}]
A linear map commuting with the automorphisms preserves generalized
eigenspaces and commutes with the maps $N_\alpha$ of
Definition~\ref{def:log-susp}, so it induces a morphism of flow
modules between suspensions. The isomorphisms of
Theorem~\ref{thm:kun-Fq}\ref{item:kF-iso} are, by their explicit
formula, natural in both factors and compatible with the associativity,
symmetry, and unit isomorphisms, the suspension of $(\C,1)$ being
$\mathbb 1_q$ (Proposition~\ref{prop:unit}): for the symmetry, choose
$\rho-\rho_1$ in place of $\rho_1$ for the exchanged factors; for the
unit, the copy of $1$ over $\rho-\rho_1$ is $v^k\cdot1$ with
$\rho-\rho_1=2\pi ik/\log q$. Composing their inverses
with the suspensions of the product $R\otimes R\to R$ and of the
actions gives the stated structures; the flow is a derivation because
the isomorphisms and the suspended maps are morphisms of flow modules.
\end{proof}

\begin{proof}[Proof of \ref{item:ksT-tor}]
Let $\bar R=R/\C\cdot1$ and $P_i=R\otimes\bar R^{\otimes i}\otimes M$
with the bar differential; the maps $r[r_1|\cdots|r_i]m\mapsto
1[r|r_1|\cdots|r_i]m$ form a contracting homotopy, so $P_\bullet\to M$ is
a resolution by the free modules $R\otimes V_i$,
$V_i=\bar R^{\otimes i}\otimes M$, and the maps induced by
$\mathrm F_R$ and $\mathrm F_M$ make it equivariant. Taking generalized
eigenspaces is exact, so suspension is exact and $\widehat P_\bullet\to
\widehat M$ is exact. Choosing one $\rho_\alpha\in\Sigma_q(\alpha)$ for
each eigenvalue $\alpha$ shows that $\widehat V$ is a free
$\mathbb 1_q$-module of rank $\dim V$, so
$\widehat{R\otimes V_i}\cong\widehat R\otimes_{\mathbb 1_q}\widehat V_i$
is a free $\widehat R$-module and $\widehat P_\bullet$ is a free
resolution of $\widehat M$. The maps $\widehat P\otimes_{\mathbb 1_q}
\widehat N\cong\widehat{P\otimes_\C N}\to\widehat{P\otimes_RN}$ are
$\widehat R$-balanced and natural in $P$, and for $P=R\otimes V$ they
induce the isomorphism $\widehat P\otimes_{\widehat R}\widehat N\cong
\widehat V\otimes_{\mathbb 1_q}\widehat N\cong\widehat{V\otimes N}=
\widehat{P\otimes_RN}$. So $\widehat P_\bullet\otimes_{\widehat R}
\widehat N\cong\widehat{P_\bullet\otimes_RN}$ as complexes, and
exactness of suspension gives the claim; for $i=0$ all maps are
morphisms of flow modules.
\end{proof}

For $R=\C$ with $\mathrm F_R=1$, the cohomology of the base
$\Spec\overline\F_q$, one has $\widehat R=\mathbb 1_q$,
$\operatorname{Tor}_i^R=0$ for $i\ge1$, and the case $i=0$ of
Proposition~\ref{prop:kun-susp-tor}\ref{item:ksT-tor} is
Theorem~\ref{thm:kun-Fq}\ref{item:kF-iso};
Theorem~\ref{thm:kun-Fq}\ref{item:kF-var} is its application to the
cohomology of two varieties over $\Spec\F_q$. For $R$ and $M$ the
cohomology of $\mathbb{P}^1_{\overline\F_q}$ and of $E_{\overline\F_q}$
in Proposition~\ref{prop:kun-EM}, tensored with $\C$ through $\iota$ and
with their Frobenius maps, $\operatorname{Tor}_i^R(M,M)$ is nonzero for
every $i$ (proof of Proposition~\ref{prop:EM-bar}; a graded free
resolution is a free resolution of the underlying modules), so
$\operatorname{Tor}_i^{\widehat R}(\widehat M,\widehat M)$ is nonzero
for every $i$, whereas over the point all higher
$\operatorname{Tor}_i^{\mathbb 1_q}$ vanish.

It is natural to expect that the positivity and duality axioms pass to
tensor products whenever the product is spectrally meaningful; this is
the case.

\begin{proposition}\label{prop:kun-tensor}
Let $D$ and $\widetilde D$ be Weil data of types $(d,0)$ and
$(\tilde d,0)$ satisfying the derivation rule of \ref{ax:AW1},
\ref{ax:AW3}, and \ref{ax:AW4}, and suppose that every
$(V^a\otimes\widetilde V^b,\Theta\otimes1+1\otimes\widetilde\Theta)$ is
a spectral core. Then the tensor datum $D\otimes\widetilde D$
(Definition~\ref{def:tensor-datum}) is a Weil datum of type
$(d+\tilde d,0)$ satisfying the derivation rule, \ref{ax:AW3}, and
\ref{ax:AW4}, with star forms $h_{D\otimes\widetilde D}=
\bigoplus_{a+b=n}h_a\otimes\tilde h_b$ in degree $n$; if $D$ and
$\widetilde D$ satisfy \ref{ax:AW1} for $U=\mathbb 1$, so does
$D\otimes\widetilde D$.
\end{proposition}

\begin{proof}
The Koszul product is associative, graded-commutative, and unital, and
$\Theta\otimes1+1\otimes\widetilde\Theta$ is a derivation of it. The
trace is supported on bidegree $(2d,2\tilde d)$, where the generator
acts on $V^{2d}\otimes\widetilde V^{2\tilde d}$ with
$\Tr(\Theta x)=d\Tr(x)$ and $\widetilde\Tr(\widetilde\Theta y)=
\tilde d\widetilde\Tr(y)$, which gives trace equivariance into
$\C(-d-\tilde d)$. The duality pairing is $B(x\otimes y,x'\otimes y')=
(-1)^{|y||x'|}B(x,x')\widetilde B(y,y')$ when $|x|+|x'|=2d$ and
$|y|+|y'|=2\tilde d$, and zero otherwise; by \ref{ax:AW3} for $D$ and
$\widetilde D$ and Proposition~\ref{prop:B-support}, it pairs the block
$V^a_{\Theta\sim\rho}\otimes\widetilde V^b_{\Theta\sim\rho'}$ perfectly
with $V^{2d-a}_{\Theta\sim d-\rho}\otimes\widetilde
V^{2\tilde d-b}_{\Theta\sim\tilde d-\rho'}$ and vanishes on the other
pairs of blocks, which by Proposition~\ref{prop:tensor-eigen} is
\ref{ax:AW3} for $D\otimes\widetilde D$. For $|x|=a$ and $|y|=b$,
\[
\Theta(*x\otimes\tilde*y)=(*\Theta x+(d-a)*x)\otimes\tilde*y
+*x\otimes(\tilde*\widetilde\Theta y+(\tilde d-b)\tilde*y),
\]
which with the constant sign $(-1)^{ab}$ is the star law in degree
$a+b$. The star form of $x\otimes y$ and $x'\otimes y'$ vanishes unless
both have bidegree $(a,b)$, and then equals
$(-1)^{ab}(-1)^{b(2d-a)}\Tr(x\cup*x')\widetilde\Tr(y\cup\tilde*y')=
h_a(x,x')\tilde h_b(y,y')$. A tensor product of positive definite
Hermitian forms is positive definite (in orthonormal bases the products
of basis vectors are orthonormal). The endpoint conditions are
$\C\otimes\C=\C$ and $\C(-d)\otimes\C(-\tilde d)=\C(-d-\tilde d)$.
\end{proof}

It is natural to expect the absolute product to be well behaved when one
factor has finite-dimensional cohomology; for the Tate objects this is
the case, and it is the K\"unneth reading of the identity
$\zeta_{\mathbb{P}^1_\Z}(s)=\zeta(s)\zeta(s-1)$
(Proposition~\ref{prop:P1Z}).

\begin{proposition}\label{prop:kun-tate}
Let $X$ be a scheme of finite type over $\Z$, $n\ge0$, and $\zeta_X$ the
Hasse--Weil zeta function of Section~\ref{subsec:schemes}.
\begin{enumerate}[label=(\roman*)]
\item\label{item:kT-zeta} $\zeta_{X\times_\Z\mathbb{P}^n_\Z}(s)=
  \prod_{i=0}^n\zeta_X(s-i)$ on the half-plane where the Euler products
  converge absolutely.
\item\label{item:kT-datum} If a Weil datum $D$ of type $(d,0)$ satisfies
  the derivation rule, \ref{ax:AW3}, and \ref{ax:AW4}, then so does
  $D\otimes T_n$ (Definition~\ref{def:tensor-datum}); if $D$ satisfies
  \ref{ax:AW1} for $U=\mathbb 1$, so does $D\otimes T_n$; if $D$
  satisfies \ref{ax:AW2} for $Z$, then $D\otimes T_n$ satisfies
  \ref{ax:AW2} for $Z_n(s)=\prod_{i=0}^nZ(s-i)$; and if $D$ satisfies
  AW-det for $Z$ with principal cuts and normalization term
  $e^{a+bs}$, then $D\otimes T_n$ satisfies AW-det for $Z_n$ with
  principal cuts and normalization term $\prod_{i=0}^ne^{a+b(s-i)}$.
\end{enumerate}
\end{proposition}

\begin{proof}[Proof of \ref{item:kT-zeta}]
Grouping closed points by residue characteristic, as in the proof of
Proposition~\ref{prop:P1Z}, $\zeta_X(s)=\prod_pZ(X_{\F_p},p^{-s})$.
Since $\#\mathbb{P}^n(\F_{p^m})=\sum_{i=0}^np^{mi}$,
$\#(X\times\mathbb{P}^n)(\F_{p^m})=\sum_ip^{mi}\#X(\F_{p^m})$, and
exponentiating gives $Z((X\times\mathbb{P}^n)_{\F_p},t)=\prod_iZ(X_{\F_p},
p^it)$; take the product over $p$ with $t=p^{-s}$.
\end{proof}

\begin{proof}[Proof of \ref{item:kT-datum}]
$T_n$ satisfies the derivation rule, \ref{ax:AW1} for $U=\mathbb 1$,
\ref{ax:AW3} ($h^i$ pairs with $h^{n-i}$, eigenvalues $i+(n-i)=n$), and
\ref{ax:AW4} ($\Theta*h^i=(n-i)*h^i=*\Theta h^i+(n-2i)*h^i$ and
$h_{2i}(h^i,h^i)=1$). The spaces $V^a\otimes\C h^i=V^a(-i)$ are spectral
cores, so Proposition~\ref{prop:kun-tensor} applies. In degree $k$ the
spectral divisor of $D\otimes T_n$ is $\sum_i\operatorname{Sp}_{k-2i}
(D)$ shifted by $i$, and $(-1)^{k+1}=(-1)^{k-2i+1}$, so the alternating
sum is $\sum_i$ of the divisor of $Z(s-i)$. For the principal cut, the
spectral zeta function of $V(-i)$ at $s$ is that of $V$ at $s-i$, so its
determinant is $D_j(s-i)$; spectral zeta functions of finite direct
sums add on the intersection of the domains, which contains a
right half-plane, so determinants multiply, and the alternating product
is $\prod_ie^{a+b(s-i)}Z(s-i)$.
\end{proof}

\subsection{The Hodge index clause}\label{subsec:kun-hodge}

It is natural to expect that positivity on the arithmetic surface is a
new condition, stronger or weaker than positivity on the curve; in fact,
on the algebraic tensors, it is equivalent to purity together with
definiteness of the star form on each eigenspace, and it is the
positivity of an explicit involution on
$\bigoplus_\rho\operatorname{End}(V_{K,\rho})$. It implies the algebra
form AW4-E of the positivity axiom, and for a suitable star it follows
from AW4-E. This is the exact counterpart of the positivity in Weil's
proof: there the trace form of correspondences on $C\times C$ is
positive (Castelnuovo; the Hodge index theorem), and this positivity is
equivalent to that of the Rosati involution
(Section~\ref{subsec:weil-curves}). The deduction of
$\Tr(u'\circ u)>0$ from the Hodge index conjecture on $X\times X$ is
Kleiman's \cite[3.11, pp.~380--381]{Kleiman1968Cycles}, after Serre
\cite[Th\'eor\`eme~2, p.~394]{Serre1960Analogues}; the positivity of
the Rosati involution is \cite[\S21, Theorem~1, p.~192]{Mumford1974AV}.
Theorem~\ref{thm:kun-hodge} is the arithmetic analogue of these
statements.

\begin{theorem}\label{thm:kun-hodge}
Let $K$ be a number field and let a Weil datum of type $(1,0)$ on
$\Xbar_K$ satisfy the derivation rule of \ref{ax:AW1}, \ref{ax:AW3}, and
the algebraic part of \ref{ax:AW4} in degree one
(Definition~\ref{def:aw4-alg}), with the notation of
Definition~\ref{def:kun-notation}.
\begin{enumerate}[label=(\roman*)]
\item\label{item:kh-iso} If $\Re\rho\ne\tfrac12$, then $h_1$ vanishes
  on $V_{K,\rho}\times V_{K,\rho}$ and $h_\otimes$ vanishes on
  $V_{K,\rho}\otimes V_{K,1-\rho}$.
\item\label{item:kh-sign} If $\Re\rho=\tfrac12$ and $V_{K,\rho}\ne0$,
  then $h_1$ is nondegenerate on $V_{K,\rho}$ and on $V_{K,1-\rho}$, of
  the same signature $(p,q)$, and $h_\otimes$ is nondegenerate on
  $V_{K,\rho}\otimes V_{K,1-\rho}$ of signature $(p^2+q^2,2pq)$.
\item\label{item:kh-iff} $h_\otimes$ is positive definite on $N_K$ if
  and only if every eigenvalue $\rho$ of $\Theta_1$ has
  $\Re\rho=\tfrac12$ and $h_1$ is definite, positive or negative, on
  each $V_{K,\rho}$.
\item\label{item:kh-pure} In that case $\Theta_1$ is diagonalizable,
  the datum satisfies AW4-V and AW4-E in degree one, and if it also
  satisfies \ref{ax:AW1} for $U=\mathbb 1$ and \ref{ax:AW2} for
  $(\Lambda_K,\mathbb 1)$, every zero of $\xi_K$ has real part
  $\tfrac12$.
\item\label{item:kh-surface} If $S$ is a multiplicative surface datum
  over the datum and itself, then $\kappa$ is injective on $N_K$ and
  $h_S^2(\kappa u,\kappa u')=h_\otimes(u,u')$ for $u,u'\in V_K\otimes
  V_K$; so $S$ satisfies the Hodge index clause if and only if
  $h_\otimes$ is positive definite on $N_K$.
\item\label{item:kh-rosati} If $\Re\rho=\tfrac12$, then $\Phi_\rho$ is a
  linear isomorphism, $h_\otimes(u,u')=\tr(\Phi_\rho(u)
  \Phi_\rho(u')^\star)$ for $u,u'\in V_{K,\rho}\otimes V_{K,1-\rho}$, and
  $\Phi_\rho(\sum_ie_i\otimes f_i)$ is the identity whenever
  $(e_i)$ is a basis of $V_{K,\rho}$ and $(f_i)$ a basis of
  $V_{K,1-\rho}$ with $B_K(e_k,f_i)=\delta_{ki}$.
\end{enumerate}
\end{theorem}

\begin{proof}[Proof of \ref{item:kh-iso}]
An antilinear map commuting with $\Theta_1$ sends $V_{K,\rho}$ into
$V_{K,\bar\rho}$, since $(\Theta_1-\bar\rho)^k*_1=
*_1(\Theta_1-\rho)^k$. So $h_1(x,y)=B_K(x,*_1y)$ with
$x\in V_{K,\rho}$ and $*_1y\in V_{K,\bar\rho}$, which vanishes by
Proposition~\ref{prop:B-support} unless $\bar\rho=1-\rho$, that is
$\Re\rho=\tfrac12$. For $x,x'\in V_{K,\rho}$,
$h_\otimes(x\otimes y,x'\otimes y')=h_1(x,x')h_1(y,y')=0$.
\end{proof}

\begin{proof}[Proof of \ref{item:kh-sign}]
Now $\bar\rho=1-\rho$. The bijection $*_1$ maps $V_{K,\rho}$ into
$V_{K,1-\rho}$ and $V_{K,1-\rho}$ into $V_{K,\rho}$, injectively, so it
maps $V_{K,\rho}$ onto $V_{K,1-\rho}$. By \ref{ax:AW3}, $B_K$ pairs
$V_{K,\rho}$ and $V_{K,1-\rho}$ perfectly, so $h_1(x,y)=B_K(x,*_1y)$ is
nondegenerate on $V_{K,\rho}$, and likewise on $V_{K,1-\rho}$. For
$a,b\in V_{K,\rho}$, $h_1(*_1a,*_1b)=B_K(*_1a,*_1^2b)=
\overline{B_K(a,*_1b)}=\overline{h_1(a,b)}$, so $*_1$ maps subspaces of
$V_{K,\rho}$ on which $h_1$ is positive (negative) definite to subspaces
of $V_{K,1-\rho}$ of the same dimension with the same property; the
signatures agree. In bases orthogonal for the two restrictions, the
products of basis vectors are $h_\otimes$-orthogonal with values the
products of the signs, which gives $(pp+qq,pq+qp)$.
\end{proof}

\begin{proof}[Proof of \ref{item:kh-iff}]
If some eigenvalue has $\Re\rho\ne\tfrac12$, then $V_{K,1-\rho}\ne0$ by
\ref{ax:AW3} and the nonzero block $V_{K,\rho}\otimes V_{K,1-\rho}$ of
$N_K$ is $h_\otimes$-null by~\ref{item:kh-iso}. If $\Re\rho=\tfrac12$
and $h_1$ is indefinite on $V_{K,\rho}$, then $2pq>0$
in~\ref{item:kh-sign}. Conversely, if all eigenvalues are on the line
and $h_1$ is definite on each $V_{K,\rho}$, each block is positive
definite by~\ref{item:kh-sign}, and distinct blocks are
$h_\otimes$-orthogonal: $h_1(V_{K,\rho},V_{K,\sigma})=
B_K(V_{K,\rho},*_1V_{K,\sigma})$ vanishes by
Proposition~\ref{prop:B-support} unless $\bar\sigma=1-\rho$, that is
$\sigma=\rho$.
\end{proof}

\begin{proof}[Proof of \ref{item:kh-pure}]
Let $\varepsilon_\rho=\pm1$ make $\varepsilon_\rho h_1$ positive on
$V_{K,\rho}$, and let $h'$ be the orthogonal sum of the forms
$\varepsilon_\rho h_1|_{V_{K,\rho}}$, a positive definite Hermitian form
on $V_K$. The computation proving
Proposition~\ref{prop:polarization-identity} for $j=1$ uses only the
derivation rule, the trace equivariance, and the star law in degree one,
so $h_1$ satisfies \eqref{eq:inf-pol} with weight one; since each
$V_{K,\rho}$ is $\Theta_1$-stable and distinct $V_{K,\rho}$ are $h_1$-
and $h'$-orthogonal, so does $h'$. Thus AW4-V holds in degree one,
AW4-E by Proposition~\ref{prop:aw4-forms}\ref{item:aw4v-e},
diagonalizability by Proposition~\ref{prop:purity}, and under
\ref{ax:AW1} and \ref{ax:AW2} Theorem~\ref{thm:rh} gives the last
assertion.
\end{proof}

\begin{proof}[Proof of \ref{item:kh-surface}]
For $x,x',y,y'\in V_K$, multiplicativity of the star, associativity, and
graded commutativity give $\kappa(x\otimes y)\cup*_S\kappa(x'\otimes
y')=-\kappa(x\otimes y)\cup\kappa(*x'\otimes*y')=\mathrm{pr}_1^*(x\cup
*x')\cup\mathrm{pr}_2^*(y\cup*y')$, and multiplicativity of the trace
gives $h_S^2(\kappa(x\otimes y),\kappa(x'\otimes y'))=
h_1(x,x')h_1(y,y')$. Injectivity is the proof of
Theorem~\ref{thm:kun-square}\ref{item:ks-inj}, which uses only the
derivation rule, \ref{ax:AW3}, and the multiplicativity of the trace.
\end{proof}

\begin{proof}[Proof of \ref{item:kh-rosati}]
$\Phi_\rho$ is injective because $B_K$ pairs $V_{K,\rho}$ and
$V_{K,1-\rho}$ perfectly, and both spaces have dimension $m(\rho)^2$.
Write $y=*_1y_0$ and $y'=*_1y'_0$ with $y_0,y'_0\in V_{K,\rho}$. Then
$\Phi_\rho(x\otimes y)=x\,h_1(\cdot,y_0)$, whose adjoint is
$y_0\,h_1(\cdot,x)$, so $\tr(\Phi_\rho(x\otimes y)\Phi_\rho(x'\otimes
y')^\star)=h_1(y'_0,y_0)h_1(x,x')$; and
$h_1(y,y')=\overline{h_1(y_0,y'_0)}=h_1(y'_0,y_0)$ by the computation
in~\ref{item:kh-sign}. The last assertion is
$\sum_iB_K(z,f_i)e_i=z$ for $z=\sum_kc_ke_k$.
\end{proof}

Thus on the algebraic tensors, identified by~\ref{item:kh-rosati} with
$\bigoplus_\rho\operatorname{End}(V_{K,\rho})$, the K\"unneth form is the
trace form $(u,u')\mapsto\tr(uu'^\star)$ of an involution, and the Hodge
index clause is the positivity of that involution: the arithmetic form
of Castelnuovo's positivity $\Tr(X\circ X')>0$ on correspondences of
$C\times C$ and of the Rosati positivity of
Theorem~\ref{thm:imp-rosati}\ref{item:ro-pos}. The identity of
$V_{K,\rho}$ for all $\rho$ at once, the component of the diagonal,
lies only in the product $\prod_\rho\operatorname{End}(V_{K,\rho})$,
not in $N_K$: the diagonal is not an algebraic class of the algebraic
tensor product, and its ``self-intersection'' $\sum_\rho m(\rho)$
diverges, the counterpart of $2g$ for a curve of infinite genus. The
operators $\int f(t)\phi^t\,dt$ generated by the flow also lie only in
that product. If all eigenvalues are on the line and $h_1$ is positive,
such an operator acts on $V_{K,\rho}$ by the scalar $\int
f(t)e^{t\rho}\,dt$, and the formal value of the trace form on it is
$\sum_\rho m(\rho)|\int f(t)e^{t\rho}\,dt|^2$, the shape of the Weil
form of Proposition~\ref{prop:weil-form}\ref{item:wf-rh}. Off the line
the restriction carries no information: by~\ref{item:kh-iso} the form
vanishes on the blocks of eigenvalues off the line, so the Hodge index
clause on these operators alone is not the Weil positivity criterion.

It is natural to expect the converse of~\ref{item:kh-pure}, that AW4-E
in degree one implies the Hodge index clause; in fact it fails for a
given star, because AW4-E constrains the spectrum only, while the Hodge
index clause also constrains the star.

\begin{definition}\label{def:split-datum}
Let $\rho=\tfrac12+i\gamma$ with $\gamma>0$. The \emph{split-signature
datum} at $\rho$ is the Weil datum of type $(1,0)$ with $V^0=\C\cdot1$,
$V^2=\C\eta$, $\Theta_0=0$, $\Theta_2=1$, and $V^1$ with basis
$e_1,e_2,f_1,f_2$, $\Theta_1e_i=\rho e_i$, $\Theta_1f_i=\bar\rho f_i$;
with the product $u\cup v=B(u,v)\eta$ on $V^1$ for the alternating form
with $B(e_i,f_k)=\delta_{ik}=-B(f_k,e_i)$ and $B(e_i,e_k)=B(f_i,f_k)=0$,
and the unit products otherwise; with the trace $\Tr(c\eta)=c$; and
with the antilinear stars $*_0(c)=\bar c\eta$, $*_2(c\eta)=\bar c$,
$*_1e_i=\varepsilon_if_i$, $*_1f_i=-\varepsilon_ie_i$, where
$(\varepsilon_1,\varepsilon_2)=(1,-1)$.
\end{definition}

For this datum the notation of Definition~\ref{def:kun-notation} is used
verbatim, with $V=V^1$, $V_\rho=\C e_1\oplus\C e_2$, and
$V_{\bar\rho}=\C f_1\oplus\C f_2$; the definition there uses only a Weil
datum of type $(1,0)$, not the field.

\begin{proposition}\label{prop:kun-hodge-strict}
Let $\rho=\tfrac12+i\gamma$ with $\gamma>0$, and consider the
split-signature datum at $\rho$ (Definition~\ref{def:split-datum}).
It satisfies \ref{ax:AW1} for $U=\mathbb 1$, \ref{ax:AW3}, the
algebraic part of \ref{ax:AW4} in degree one, AW4-V, and AW4-E in degree
one, and $h_\otimes$ is not positive definite on
$V_\rho\otimes V_{\bar\rho}$.
\end{proposition}

\begin{proof}
\ref{ax:AW1} and \ref{ax:AW3} hold as for the trivial model
(Proposition~\ref{prop:trivial-model}\ref{item:tm-aw13}), since $B$
pairs the $\rho$- and $\bar\rho=(1-\rho)$-eigenspaces perfectly and
vanishes otherwise. The map $*_1$ is bijective, sends $V_\rho$ to
$V_{\bar\rho}$ and back, so it commutes with $\Theta_1$; on basis
vectors $B(*_1e_i,*_1f_k)=-\varepsilon_i\varepsilon_kB(f_i,e_k)=
\delta_{ik}$, $B(*_1f_k,*_1e_i)=-\delta_{ik}$, and the other values
vanish, so $B(*_1x,*_1y)=\overline{B(x,y)}$ by antilinearity. The form
$h_1(x,y)=B(x,*_1y)$ has $h_1(e_i,e_k)=\varepsilon_k\delta_{ik}$,
$h_1(f_i,f_k)=\varepsilon_k\delta_{ik}$, and $h_1(e_i,f_k)=h_1(f_k,e_i)=0$,
so it is Hermitian, of signature $(1,1)$ on $V_\rho$. The spectrum
$\{\rho,\bar\rho\}$ lies on the line and $\Theta_1$ is diagonal; the
orthogonal sum of positive definite forms on the two eigenspaces
satisfies \eqref{eq:inf-pol} with weight one, since for eigenvectors
with eigenvalues $\sigma,\sigma'$ both sides are $(\sigma+\bar\sigma')
h(u,v)$, which is $h(u,v)$ if $\sigma=\sigma'$ and $0$ otherwise. So
AW4-V holds, and AW4-E by
Proposition~\ref{prop:aw4-forms}\ref{item:aw4v-e}. By
Theorem~\ref{thm:kun-hodge}\ref{item:kh-sign}, $h_\otimes$ has
signature $(2,2)$ on $V_\rho\otimes V_{\bar\rho}$.
\end{proof}

With the star $*'_1$ of Proposition~\ref{prop:aw4-forms}%
\ref{item:aw4v-star} built from a positive polarization, the same
datum satisfies \ref{ax:AW4}, and the tensor star form of $*'_1$ is
positive definite, being the tensor square of a positive form: the
Hodge index clause is a property of the star, not of the spectrum
alone.

It is natural to expect the conjecture, like any existence statement
about the surface, to have content beyond the Riemann hypothesis; in
fact, without the requirement that the surface be constructed
independently of the zeros, it collapses as the axioms on one object do.

\begin{proposition}\label{prop:kun-collapse}
Let $K$ and $L$ be number fields and $D_K$, $D_L$ Weil data of type
$(1,0)$ on $\Xbar_K$ and $\Xbar_L$ satisfying \ref{ax:AW1} for
$U=\mathbb 1$.
\begin{enumerate}[label=(\roman*)]
\item\label{item:kl-tensor} The tensor datum $D_K\otimes D_L$
  (Definition~\ref{def:tensor-datum}) with $\mathrm{pr}_1^*x=x\otimes1$
  and $\mathrm{pr}_2^*y=1\otimes y$ is a multiplicative surface datum
  over $(D_K,D_L)$ whose K\"unneth map is the identity.
\item\label{item:kl-iff} The existence of the data of
  Conjecture~\ref{conj:kunneth} satisfying~\ref{item:K-iso}
  and~\ref{item:K-hodge}, with no condition on how they are
  constructed, is equivalent to the statement that every zero of every
  $\xi_K$ has real part $\tfrac12$.
\end{enumerate}
\end{proposition}

\begin{proof}[Proof of \ref{item:kl-tensor}]
By \ref{ax:AW1} the products are associative, graded-commutative, and
unital, and $\Theta(1)=0$ because $\Theta$ is a derivation. The Koszul
tensor product of such algebras is again associative, graded-commutative,
and unital, $\Theta\otimes1+1\otimes\Theta$ is a derivation of it, the
maps $x\mapsto x\otimes1$ and $y\mapsto1\otimes y$ are unital algebra
homomorphisms commuting with the flows (using $\Theta(1)=0$),
$(x\otimes1)\cup(1\otimes y)=x\otimes y$, the trace
$\Tr_K\otimes\Tr_L$ is multiplicative, and the star rule is the
definition of the stars in Definition~\ref{def:tensor-datum}.
\end{proof}

\begin{proof}[Proof of \ref{item:kl-iff}]
If every $\xi_K$ has its zeros on the line, the trivial models with
spectra $\divisor\xi_K$ satisfy \ref{ax:AW1}--\ref{ax:AW4}
(Proposition~\ref{prop:content}), and in particular the algebraic part
of \ref{ax:AW4} in degree one (the computation after
Definition~\ref{def:aw4-alg}), with $h_1$ positive; take
$S_{K,L}=D_K\otimes D_L$ as in~\ref{item:kl-tensor}. Its K\"unneth map
is an isomorphism, and $h_\otimes=h_1\otimes h_1$ is positive definite,
so the Hodge index clause holds by
Theorem~\ref{thm:kun-hodge}\ref{item:kh-surface}. Conversely, the Hodge
index clause for $S_{K,K}$ and
Theorem~\ref{thm:kun-hodge}\ref{item:kh-surface}, \ref{item:kh-iff},
and~\ref{item:kh-pure} put every zero of $\xi_K$ on the line.
\end{proof}

Theorem~\ref{thm:kun-hodge} and Proposition~\ref{prop:kun-collapse} say
that, as existence statements, the Hodge index clause and the
positivity axiom on the curve carry the same content; for a given star
the clause is the stronger condition
(Proposition~\ref{prop:kun-hodge-strict}). What the arithmetic surface
would add is a place
where the positivity could be proved from geometry, as Weil proved
Castelnuovo's inequality by intersection theory on $C\times C$; the
statement to be proved is fixed by Theorem~\ref{thm:kun-hodge}.
\section{Negative results}\label{sec:negative}

This section records inferences that look plausible on the way to a
Deninger cohomology theory or a de Branges realization of the zeros,
but are not valid. Each finding is introduced by a sentence of the form
``it is natural to expect $X$; in fact $Y$'', followed by a numbered
statement. Unless the contrary is said explicitly, a finding refutes an
inference, not a theorem in the literature: the witnesses are
counterexamples to deductions from the listed properties, not to the
Weil conjectures, to the existence of any cohomology theory, or to any
published main theorem. Section~\ref{subsec:neg-register} maps the
findings to a register of the expectations they correct.

\subsection{Algebraic findings}\label{subsec:neg-algebraic}

The first two findings, R1 and R2 (Definitions~\ref{def:R1-forms}
and~\ref{def:R2-block}, Propositions~\ref{prop:R1}, \ref{prop:R2},
and~\ref{prop:R2-star}), are stated in Section~\ref{subsec:witnesses},
before the main theorem, whose Section~\ref{subsec:content} uses them.

It is natural to expect that over a finite field a Frobenius commuting
with a quaternionic structure and preserving a perfect alternating
duality of weight one is pure; in fact the block of
Definition~\ref{def:R2-block} exponentiates to a counterexample.

\begin{proposition}\label{prop:R2-frob}
Let $0<a<\tfrac12$, $T>0$, $q>1$, let $\Theta$, $\Omega$, and $J$ be as
in Definition~\ref{def:R2-block}, and let
$\mathrm{F}=q^\Theta=\diag(q^\lambda,q^{\bar\lambda},q^{1-\lambda},
q^{1-\bar\lambda})$. Then $\mathrm{F}^T\Omega\mathrm{F}=q\Omega$,
$\mathrm{F}J=J\mathrm{F}$, $J^2=-1$, and the eigenvalues of
$\mathrm{F}$ have moduli $q^{1/2+a}$ and $q^{1/2-a}$.
\end{proposition}

\begin{proof}
$\Theta$ is diagonal and $\Theta^T\Omega+\Omega\Theta=\Omega$ (proof of
Proposition~\ref{prop:R2}\ref{item:R2-axioms}), so the nonzero entries
of $\Omega$ sit at pairs of eigenvalues with sum one, and
$\mathrm{F}^T\Omega\mathrm{F}$ multiplies each of them by $q$. Since
$J\Theta=\Theta J$ and $J$ is antilinear, $J$ maps the $\mu$-eigenline
of $\Theta$ to the $\bar\mu$-eigenline, on which $\mathrm{F}$ is
$q^{\bar\mu}=\overline{q^\mu}$; so $J\mathrm{F}=\mathrm{F}J$. The
equality $J^2=-1$ is Proposition~\ref{prop:R2}\ref{item:R2-axioms}, and
the moduli are $q^{\Re\mu}$.
\end{proof}

It is natural to expect that comparing $V$ with its integer Tate twists
by cones sees every summand, and that adjoining an acyclic block leaves
determinant and special-value data unchanged; in fact integer-twist
cones are blind to summands with nonreal spectrum, while determinants
change by a nonconstant factor.

\begin{definition}\label{def:twist-cone}
Let $V$ be a finite-dimensional complex vector space, $\Theta:V\to V$
linear, and $n\in\Z$. The \emph{twist cone} of $(V,\Theta)$ at $n$ is the
two-term complex $C_n(V,\Theta)=[V\xrightarrow{\Theta-n}V]$ in degrees
$0$ and $1$.
\end{definition}

\begin{proposition}\label{prop:R3}
Let $0<a<\tfrac12$, $T>0$, and let $(V,\Theta)$, $p$, and $w$ be as in
Definition~\ref{def:R2-block}. Then:
\begin{enumerate}[label=(\roman*)]
\item\label{item:R3-acyclic} for every $n\in\Z$ the map $\Theta-n$ is
  invertible, so $C_n(V,\Theta)$ has zero cohomology;
\item\label{item:R3-poly} $p(s)$ equals
  \[
  ((w-a)^2+T^2)((w+a)^2+T^2)=w^4+2(T^2-a^2)w^2+(T^2+a^2)^2,
  \]
  satisfies $p(1-s)=p(s)$ and $p(s)>0$ for real $s$, and has degree
  four;
\item\label{item:R3-det} if $(V',\Theta')$ is a spectral core that is
  zeta-regularizable for a spectral cut on an open set $U$, then
  $(V'\oplus V,\Theta'\oplus\Theta)$ is a spectral core that is
  zeta-regularizable for a spectral cut on a nonempty open subset
  $U_0\subset U$ extending the given one, and, with the determinants
  taken for these two cuts,
  \[
  \detreg\bigl((s-\Theta'\oplus\Theta)/2\pi\bigr)
  =(2\pi)^{-4}p(s)\,\detreg\bigl((s-\Theta')/2\pi\bigr).
  \]
\end{enumerate}
\end{proposition}

\begin{proof}[Proof of \ref{item:R3-acyclic}]
The eigenvalues $\tfrac12\pm a\pm iT$ have nonzero imaginary part, so
none equals $n$.
\end{proof}

\begin{proof}[Proof of \ref{item:R3-poly}]
$(s-\lambda)(s-\bar\lambda)=(w-a)^2+T^2$ and
$(s-(1-\lambda))(s-(1-\bar\lambda))=(w+a)^2+T^2$. With
$A=w^2+a^2+T^2$ the product is $(A-2aw)(A+2aw)=A^2-4a^2w^2$, which
expands to the displayed quartic. The quartic is even in $w$, and
$s\mapsto1-s$ is $w\mapsto-w$. For real $s$ both factors are at least
$T^2>0$.
\end{proof}

\begin{proof}[Proof of \ref{item:R3-det}]
The sum is a spectral core: its generalized eigenspaces are those of
the summands added together, and its spectrum is the union of
$\spec\Theta'$ with four points. Let $Z'$ be the spectral zeta
function of $(V',\Theta')$ for the given cut on $U$. Removing from $U$
the four closed rays $\rho_k+(-\infty,0]$ through the eigenvalues
$\rho_k$ of $\Theta$ leaves a nonempty open set $U_0$. Extend the cut
to the sum by keeping the given branches on $\spec\Theta'$ and taking
the principal branch of $\log((s-\rho_k)/2\pi)$ on $U_0$ for each
$\rho_k\notin\spec\Theta'$; in either case the branch $L_{\rho_k}$
satisfies $e^{L_{\rho_k}(s)}=(s-\rho_k)/2\pi$. On $U_0$ the spectral
zeta function of the sum is
$Z'(w;s)+\sum_{k=1}^4e^{-wL_{\rho_k}(s)}$; the added terms are entire
in $w$, holomorphic in $s$, and have $w$-derivative $-L_{\rho_k}(s)$
at $w=0$. Hence
\[
\exp\bigl(-\partial_w Z_{V'\oplus V}(0;s)\bigr)
=\exp\bigl(-\partial_wZ'(0;s)\bigr)\prod_{k=1}^4\frac{s-\rho_k}{2\pi}
\]
on $U_0$, and the right side continues to the entire function
$(2\pi)^{-4}p(s)\detreg((s-\Theta')/2\pi)$.
\end{proof}

Thus the quasi-isomorphism classes of all integer-twist cones, and
their long exact sequences, are unchanged by adjoining $(V,\Theta)$,
whereas a degree-one determinant is multiplied by $(2\pi)^{-4}p(s)$, a
nonconstant factor which is positive on the real axis and symmetric
under $s\mapsto1-s$. A determinant-line trivialization carrying metric
or basis data can retain this factor even for an acyclic complex. The
finding concerns integer-cone data alone; it is not a counterexample to
the full package of Deninger and Flach--Morin
\cite[Conjecture~4.1]{FlachMorin2020DeningerWeilArakelov}, which also
specifies determinant identities and leading coefficients.

It is natural to expect that for a higher-dimensional variety all zeros
of its zeta function lie on a single critical line; in fact purity
holds degree by degree, and zeros from different degrees lie on
different lines.

\begin{proposition}\label{prop:N-lines}
Let $q$ be a prime power, $E$ an elliptic curve over $\F_q$ with
Frobenius roots $\alpha,\beta$, and $X=E\times\mathbb{P}^1$, with zeta
functions $Z(E,t)$, $Z(X,t)$, and $\zeta_X$ as in
Definition~\ref{def:hasse-weil-zeta}. Then:
\begin{enumerate}[label=(\roman*)]
\item\label{item:N-product} $Z(X,t)=Z(E,t)\,Z(E,qt)$;
\item\label{item:N-zeros} the zeros of $\zeta_X$ form the set
  $\Sigma_q(\alpha)\cup\Sigma_q(\beta)\cup
  \bigl(1+\Sigma_q(\alpha)\bigr)\cup\bigl(1+\Sigma_q(\beta)\bigr)$
  (Definition~\ref{def:suspension}); each of the lines
  $\Re s=\tfrac12$ and $\Re s=\tfrac32$ contains infinitely many of
  them, and every zero lies on one of these two lines;
\item\label{item:N-poles} the poles of $\zeta_X$ lie on the lines
  $\Re s\in\{0,1,2\}$.
\end{enumerate}
\end{proposition}

\begin{proof}[Proof of \ref{item:N-product}]
$\#\mathbb{P}^1(\F_{q^n})=q^n+1$, so
$\#X(\F_{q^n})=\#E(\F_{q^n})+q^n\#E(\F_{q^n})$. Exponentiating
$\sum_n\#E(\F_{q^n})t^n/n$ gives $Z(E,t)$, and exponentiating
$\sum_n\#E(\F_{q^n})(qt)^n/n$ gives $Z(E,qt)$.
\end{proof}

\begin{proof}[Proof of \ref{item:N-zeros} and \ref{item:N-poles}]
By \eqref{eq:elliptic-zeta} (the case of elliptic curves of
Theorem~\ref{thm:imp-deligne}, Section~\ref{sec:classical}),
$Z(E,t)=(1-\alpha t)(1-\beta t)/((1-t)(1-qt))$ with
$|\alpha|=|\beta|=q^{1/2}$. By \ref{item:N-product},
\[
\zeta_X(s)=\frac{(1-\alpha q^{-s})(1-\beta q^{-s})
(1-\alpha q^{1-s})(1-\beta q^{1-s})}
{(1-q^{-s})(1-q^{1-s})(1-q^{1-s})(1-q^{2-s})}.
\]
By Proposition~\ref{prop:suspension}\ref{item:susp-zeros},
$1-\alpha q^{-s}$ vanishes exactly on $\Sigma_q(\alpha)$, and
$1-\alpha q^{1-s}$ vanishes exactly when $s-1\in\Sigma_q(\alpha)$;
likewise for $\beta$. Points of $\Sigma_q(\alpha)$ and
$\Sigma_q(\beta)$ have real part $\tfrac12$ since
$|\alpha|=|\beta|=q^{1/2}$, and points of $1+\Sigma_q(\alpha)$ and
$1+\Sigma_q(\beta)$ real part $\tfrac32$; each of these sets is an
infinite arithmetic
progression. The denominator vanishes exactly where $q^{s}\in\{1,q,q^2\}$,
that is on $\Re s\in\{0,1,2\}$. Since numerator and denominator have no
common zero (their real parts differ), the zeros of $\zeta_X$ are those
of the numerator and its poles those of the denominator.
\end{proof}

The zeros on $\Re s=\tfrac12$ come from $H^1(X)$ (weight one) and those
on $\Re s=\tfrac32$ from $H^3(X)$ (weight three); this is the
degree-by-degree form of purity, the finite-field counterpart of
Theorem~\ref{thm:A-schemes}.


\subsection{Operator-theoretic and analytic findings}
\label{subsec:neg-operator}

It is natural to expect that a skew-symmetric operator has purely
imaginary spectrum, so that the algebraic purity argument controls the
whole spectrum of a completion; in fact a closed, densely defined
skew-symmetric operator can have every complex number in its spectrum
and no eigenvalues at all.

\begin{definition}\label{def:R4-D}
Let $H^1_0(0,1)$ be the space of absolutely continuous
$f:[0,1]\to\C$ with $f'\in L^2(0,1)$ and $f(0)=f(1)=0$. The
\emph{Dirichlet derivative} is the operator $Df=f'$ on $L^2(0,1)$ with
domain $H^1_0(0,1)$.
\end{definition}

\begin{proposition}\label{prop:R4}
Let $D$ be the Dirichlet derivative of Definition~\ref{def:R4-D}. Then:
\begin{enumerate}[label=(\roman*)]
\item\label{item:R4-closed} $D$ is densely defined and closed;
\item\label{item:R4-skew} $\langle Df,g\rangle=-\langle f,Dg\rangle$
  for all $f,g\in H^1_0(0,1)$;
\item\label{item:R4-noeig} $D$ has no eigenvalues;
\item\label{item:R4-range} for every $z\in\C$ the function
  $t\mapsto e^{zt}$ does not lie in the range of $D-z$; hence the
  spectrum of $D$ is $\C$;
\item\label{item:R4-notskew} $D$ is not skew-adjoint.
\end{enumerate}
\end{proposition}

\begin{proof}[Proof of \ref{item:R4-closed}]
$C_c^\infty(0,1)\subset H^1_0(0,1)$ is dense in $L^2(0,1)$. Let
$f_n\in H^1_0(0,1)$ with $f_n\to f$ and $f_n'\to g$ in $L^2$. Put
$G(x)=\int_0^xg$. Since $f_n(0)=0$, the Cauchy--Schwarz inequality gives
$|f_n(x)-G(x)|\le\|f_n'-g\|_2$ for $x\in[0,1]$, so $f_n\to G$
uniformly. Hence $f=G$ almost everywhere, $G$ is absolutely continuous
with $G'=g\in L^2$, $G(0)=0$, and $G(1)=\lim f_n(1)=0$. Thus $f\in
H^1_0(0,1)$ and $Df=g$.
\end{proof}

\begin{proof}[Proof of \ref{item:R4-skew}]
Integration by parts for absolutely continuous functions gives
$\int_0^1f'\bar g=[f\bar g]_0^1-\int_0^1f\bar g'=-\int_0^1f\bar g'$.
\end{proof}

\begin{proof}[Proof of \ref{item:R4-noeig}]
If $f'=zf$ with $f\in H^1_0(0,1)$, then $e^{-zt}f$ is absolutely
continuous with derivative zero almost everywhere, hence constant, equal
to its value $f(0)=0$ at $t=0$.
\end{proof}

\begin{proof}[Proof of \ref{item:R4-range}]
If $(D-z)f=g$ with $f\in H^1_0(0,1)$, then $(e^{-zt}f)'=e^{-zt}g$, and
integrating over $[0,1]$ gives
$\int_0^1e^{-zt}g(t)\,dt=e^{-z}f(1)-f(0)=0$. For $g(t)=e^{zt}$ the integral is
$\int_0^11\,dt=1$. So $D-z$ is not surjective for any $z$, and every
$z$ lies in the spectrum.
\end{proof}

\begin{proof}[Proof of \ref{item:R4-notskew}]
If $D^*=-D$, then $iD$ is self-adjoint, so $iD-i$ is surjective; then
$D-1=-i(iD-i)$ is surjective, contradicting \ref{item:R4-range} with
$z=1$.
\end{proof}

The eigenvector argument of Proposition~\ref{prop:purity} is algebraic
and remains valid; what fails is any claim about the full spectrum of a
completion obtained without a domain theorem. A self-adjoint
realization requires such a statement, which is why \ref{ax:AW5} asks
for a strongly continuous unitary group with skew-adjoint generator.

It is natural to expect that a scalar de Branges model built from $F$
retains the multiplicities of the zeros of $F$; in fact a real zero of
order $k$ becomes a common real zero of order $k-1$ of $E$ and $E^\#$,
and a cyclic self-adjoint operator cannot carry multiplicity in its
eigenspaces at all.

\begin{definition}\label{def:cyclic}
Let $A$ be a self-adjoint operator on a Hilbert space $\cH$. A vector
$v_0\in\cH$ is \emph{cyclic} for $A$ if the closed linear span of
$\{v_0\}\cup\{(A-z)^{-1}v_0:z\in\C\setminus\R\}$ is $\cH$.
\end{definition}

\begin{proposition}\label{prop:R9}
\leavevmode
\begin{enumerate}[label=(\roman*)]
\item\label{item:R9-order} Let $G$ be an entire function with $G^\#=G$
  and $G\not\equiv0$, with HB companion $E=E_G$
  (Definition~\ref{def:hb}). Then $E^\#=G-iG'$, and if $r\in\R$ is a
  zero of $G$ of order $k\ge1$, then $E$ and $E^\#$ vanish at $r$ to
  order exactly $k-1$ (in particular $E(r)E^\#(r)\ne0$ when $k=1$).
\item\label{item:R9-example} Let $G(z)=(z^2-1)^2$ and $E=E_G$. Then
  $E(z)=(z^2-1)(z^2-1+4iz)$ and $E^\#(z)=(z^2-1)(z^2-1-4iz)$, and $\pm1$
  are common simple zeros of $E$ and $E^\#$.
\item\label{item:R9-cyclic} Let $A$ be a self-adjoint operator on a
  Hilbert space $\cH$ with a cyclic vector $v_0$. Then
  $\dim\ker(A-\lambda)\le1$ for every $\lambda\in\C$.
\end{enumerate}
\end{proposition}

\begin{proof}[Proof of \ref{item:R9-order}]
Since $(G')^\#=(G^\#)'=G'$, we get $E^\#=G^\#-i(G')^\#=G-iG'$. Write
$G=(z-r)^kU$ with $U$ entire and $U(r)\ne0$. Then
$G'=(z-r)^{k-1}(kU+(z-r)U')$, so
\[
E=(z-r)^{k-1}\bigl((z-r)U+ikU+i(z-r)U'\bigr),
\]
and the bracket takes the value $ikU(r)\ne0$ at $r$. The same computation
with $-i$ in place of $i$ applies to $E^\#$.
\end{proof}

\begin{proof}[Proof of \ref{item:R9-example}]
$G'(z)=4z(z^2-1)$, so $E=(z^2-1)(z^2-1+4iz)$ and, by
\ref{item:R9-order}, $E^\#=(z^2-1)(z^2-1-4iz)$. The second factors take
the values $\pm4i$ at $z=\pm1$.
\end{proof}

\begin{proof}[Proof of \ref{item:R9-cyclic}]
Let $K=\ker(A-\lambda)$; we may assume $K\ne0$, so $\lambda\in\R$. Let
$P$ be the orthogonal projection onto $K$ and $R_z=(A-z)^{-1}$ for
$z\notin\R$. For $v\in K$, $(A-z)v=(\lambda-z)v$, so
$R_zv=(\lambda-z)^{-1}v\in K$. For $u\perp K$ and $v\in K$,
$\langle R_zu,v\rangle=\langle u,R_{\bar z}v\rangle=0$, using
$R_z^*=R_{\bar z}$; so $R_z$ preserves $K^\perp$ as well, and
$PR_z=R_zP$. Hence $PR_zv_0=R_zPv_0=(\lambda-z)^{-1}Pv_0\in\C Pv_0$,
and $P$ maps the span of $\{v_0\}\cup\{R_zv_0\}$ into $\C Pv_0$. By
continuity of $P$ and cyclicity, $K=P\cH\subset\C Pv_0$.
\end{proof}

For a positive measure $\mu=\sum_rm_r\delta_r$ on a discrete subset of
$\R$, multiplication by the variable on $L^2(\mu)$ has, at each atom $r$,
the eigenspace of functions supported on $\{r\}$, which is
one-dimensional whatever the weight $m_r$: the residue of $-F'/F$ at $r$
records $m_r$, the eigenspace dimension does not. In particular the zero
operator on $\C^2$ admits no cyclic vector. By \ref{item:R9-order}, the
function $E=E_F=F+iF'$ of Proposition~\ref{prop:hb} may have common real
zeros with $E^\#$ when $F$ has multiple real zeros; requiring $E$ to have
no real zeros would assume simplicity. A scalar de Branges realization
of arithmetic $H^1$ therefore needs structure beyond self-adjointness,
and a multiplicity-preserving identification requires additional data.

It is natural to expect that the semilocal de Branges spaces of
Connes--Consani--Moscovici may be identified with the global space
attached to $\Xi$ once both are known to be Hilbert spaces; in fact an
abstract isomorphism of Hilbert spaces, unitary or merely bounded with
bounded inverse, carries no spectral information.

\begin{definition}\label{def:ccm23-semilocal}
For a finite set $S$ of places of $\Q$ containing $\infty$ and
$\lambda>0$, we write $\mathcal B^S_\lambda$ for the Hilbert space of
entire functions constructed in
\cite[\S4.8]{ConnesConsaniMoscovici2024Prolate} from Sonin spaces,
with its $S$-dependent induced norm, and $\upsilon_S$ for the
isomorphism of hilbertian spaces (a linear homeomorphism) of
\cite[Proposition~4.8]{ConnesConsaniMoscovici2024Prolate}.
\end{definition}

\begin{proposition}\label{prop:R10}
Let $\cH_1,\cH_2$ be separable infinite-dimensional complex Hilbert
spaces and let $(\lambda_n)_{n\ge1}$, $(\mu_n)_{n\ge1}$ be real sequences
with $|\lambda_n|\to\infty$ and $|\mu_n|\to\infty$. Then:
\begin{enumerate}[label=(\roman*)]
\item\label{item:R10-iso} there is a unitary map $U:\cH_1\to\cH_2$;
\item\label{item:R10-ops} there are self-adjoint operators $A_1$ on
  $\cH_1$ and $A_2$ on $\cH_2$ with pure point spectrum such that
  $\dim\ker(A_1-x)=\#\{n:\lambda_n=x\}$ and
  $\dim\ker(A_2-x)=\#\{n:\mu_n=x\}$ for every $x\in\R$;
\item\label{item:R10-noneq} if $\#\{n:\lambda_n=x\}\ne\#\{n:\mu_n=x\}$
  for some $x$, there is no bounded linear bijection
  $W:\cH_1\to\cH_2$, unitary or not, with
  $W\mathcal D(A_1)=\mathcal D(A_2)$ and $WA_1=A_2W$ on
  $\mathcal D(A_1)$.
\end{enumerate}
\end{proposition}

\begin{proof}[Proof of \ref{item:R10-iso}]
Choose orthonormal bases $(e_n)_{n\ge1}$ of $\cH_1$ and $(f_n)_{n\ge1}$
of $\cH_2$ and let $Ue_n=f_n$.
\end{proof}

\begin{proof}[Proof of \ref{item:R10-ops}]
Let $A_1e_n=\lambda_ne_n$ with domain
$\{\sum c_ne_n:\sum|\lambda_nc_n|^2<\infty\}$, and similarly $A_2$. The
argument in the proof of Proposition~\ref{prop:hilbert} shows that these
diagonal operators are self-adjoint. A vector $\sum c_ne_n$ lies in
$\ker(A_1-x)$ if and only if $c_n(\lambda_n-x)=0$ for all $n$, so the
kernel is spanned by the $e_n$ with $\lambda_n=x$, finitely many since
$|\lambda_n|\to\infty$. The spectrum is the closure of $\{\lambda_n\}$,
which is this discrete set.
\end{proof}

\begin{proof}[Proof of \ref{item:R10-noneq}]
If $v\in\ker(A_1-x)$, then $Wv\in\mathcal D(A_2)$ and
$A_2Wv=WA_1v=xWv$. If $u\in\ker(A_2-x)$, then
$W^{-1}u\in\mathcal D(A_1)$ and $WA_1W^{-1}u=A_2u=xu=W(xW^{-1}u)$,
so $A_1W^{-1}u=xW^{-1}u$ by injectivity of $W$. Hence $W$ maps
$\ker(A_1-x)$ bijectively onto $\ker(A_2-x)$, and the dimensions
agree.
\end{proof}

The isomorphism $\upsilon_S$ of Definition~\ref{def:ccm23-semilocal} is
a theorem of this abstract kind together with a compatibility with the
amplification maps; the entire-function realization and norm of
$\mathcal B^S_\lambda$ depend on $S$ and $\lambda$. By
Proposition~\ref{prop:R10}, an identification with the global space of
Section~\ref{sec:debranges} would need an intertwiner of the specified
operators, the same norm, a controlled limit in $S$ and $\lambda$, and
multiplicity data; none of these is supplied by
\cite[Proposition~4.8]{ConnesConsaniMoscovici2024Prolate}.

It is natural to expect that positivity of $\Im(-F'/F)$ obtained high in
the upper half-plane continues down to the real axis, or that an inverse
canonical-system theorem can supply it; in fact positivity can hold on
$\Im z>a$ and fail just below, and a single zero in $\Cp$ forces
negative values next to it.

\begin{definition}\label{def:R11-quartet}
For $T>0$ and $a>0$, the \emph{quartet} is the polynomial
$P_{T,a}(z)=((z-T)^2+a^2)((z+T)^2+a^2)$.
\end{definition}

\begin{proposition}\label{prop:R11}
Let $T>0$ and $a>0$, and write $P=P_{T,a}$
(Definition~\ref{def:R11-quartet}).
\begin{enumerate}[label=(\roman*)]
\item\label{item:R11-shape} $P$ has real coefficients, $P(-z)=P(z)$, and
  its zeros $\pm T\pm ia$ are not real.
\item\label{item:R11-high} $\Im(-P'(z)/P(z))>0$ for every $z$ with
  $\Im z>a$.
\item\label{item:R11-value} If $T=2$ and $a=\tfrac14$, then
  \[
  \Im\Bigl(-\frac{P'(2+\tfrac i8)}{P(2+\tfrac i8)}\Bigr)
  =-\frac{16892192}{3176475},
  \]
  and consequently the HB companion $E_P$ (Definition~\ref{def:hb})
  satisfies $|E_P(2+\tfrac i8)|<|E_P^\#(2+\tfrac i8)|$.
\item\label{item:R11-pole} Let $G$ be an entire function,
  $G\not\equiv0$, with a zero $z_0\in\Cp$ of order $k\ge1$. Then there
  is $\varepsilon_0\in(0,\Im z_0)$ such that $-G'/G$ is holomorphic at
  $z_0-i\varepsilon$ and $\Im(-G'/G)(z_0-i\varepsilon)<0$ for
  $0<\varepsilon<\varepsilon_0$.
\end{enumerate}
\end{proposition}

\begin{proof}[Proof of \ref{item:R11-shape}]
Both factors are real polynomials, and $z\mapsto-z$ interchanges them.
\end{proof}

\begin{proof}[Proof of \ref{item:R11-high}]
With $r$ running over the four zeros,
$-P'/P=\sum_r1/(r-z)$ and
$\Im\bigl(1/(r-z)\bigr)=(\Im z-\Im r)/|z-r|^2$. If $\Im z>a$, every term
is positive, since $\Im r=\pm a$.
\end{proof}

\begin{proof}[Proof of \ref{item:R11-value}]
At $z=2+\tfrac i8$ the four terms of the preceding proof are, for
$r=2+\tfrac i4$, $2-\tfrac i4$, $-2+\tfrac i4$, $-2-\tfrac i4$:
\[
\frac{-1/8}{1/64}=-8,\quad
\frac{3/8}{9/64}=\frac83,\quad
\frac{-1/8}{16+1/64}=-\frac{8}{1025},\quad
\frac{3/8}{16+9/64}=\frac{24}{1033}.
\]
Their sum is $-16892192/3176475$. Since $P^\#=P$, $E_P^\#=P-iP'$, and
$|x+y|^2-|x-y|^2=4\Re(x\bar y)$ with $x=P$, $y=iP'$ gives
$|E_P|^2-|E_P^\#|^2=4\Im(P\overline{P'})=4|P|^2\Im(-P'/P)$, which is
negative at $2+\tfrac i8$.
\end{proof}

\begin{proof}[Proof of \ref{item:R11-pole}]
Near $z_0$, $-G'/G=-k/(z-z_0)+g(z)$ with $g$ holomorphic on a disc
$|z-z_0|<\delta$ containing no other zero of $G$, and $|g|\le M$ on
$|z-z_0|\le\delta/2$. At $z=z_0-i\varepsilon$,
$-k/(z-z_0)=-ik/\varepsilon$, so
$\Im(-G'/G)\le-k/\varepsilon+M<0$ once $\varepsilon<k/M$. Take
$\varepsilon_0=\min(\delta/2,k/M,\Im z_0/2)$ (with $k/M=\infty$ if
$M=0$); then $z_0-i\varepsilon\in\Cp$ for $0<\varepsilon<\varepsilon_0$.
\end{proof}

By Proposition~\ref{prop:hb}, RH is equivalent to positivity of
$\Im(-F'/F)$ on $\Cp$, and to the HB inequality for $E_F=F+iF'$. The
quartet $P$ has the symmetries of $\Xi$ and is positive in this sense
above its zeros, but not below. An inverse canonical-system theorem whose
input already contains the HB inequality for $\Xi+i\Xi'$ cannot establish
that inequality; positivity, and the relevant endpoint condition, must
be obtained independently.

It is natural to expect that an entire function which is the Fourier
transform of a positive measure has only real zeros; in fact positivity
of the kernel alone does not locate zeros.

\begin{definition}\label{def:R12-kernels}
\begin{enumerate}[label=(\alph*)]
\item The \emph{Fourier transform} of a finite positive Borel measure
  $\nu$ on $\R$ with compact support is the entire function
  $\widehat\nu(z)=\int e^{izt}\,d\nu(t)$ (the convention of
  Definition~\ref{def:weil-form}).
\item For $a>0$, put
  $\mu_a=\cosh(a)\delta_0+\tfrac12(\delta_1+\delta_{-1})$.
\item For $a>0$ and $\sigma>0$, put
  $g_\sigma(t)=(4\pi\sigma)^{-1/2}e^{-t^2/(4\sigma)}$ and
  $k_{a,\sigma}(t)=\cosh(a)g_\sigma(t)
  +\tfrac12(g_\sigma(t-1)+g_\sigma(t+1))$.
\end{enumerate}
\end{definition}

\begin{proposition}\label{prop:R12}
Let $a>0$, and let $\mu_a$, $g_\sigma$, and $k_{a,\sigma}$ be as in
Definition~\ref{def:R12-kernels}.
\begin{enumerate}[label=(\roman*)]
\item\label{item:R12-atoms} $\widehat{\mu_a}(z)=\cos z+\cosh a$ for all
  $z\in\C$, and its zeros are the points $(2m+1)\pi\pm ia$, $m\in\Z$, all
  simple and none real.
\item\label{item:R12-gauss} For every $\sigma>0$, $k_{a,\sigma}$ is
  smooth, integrable, and everywhere positive, and
  $\int e^{izt}k_{a,\sigma}(t)\,dt=e^{-\sigma z^2}(\cos z+\cosh a)$
  for all $z\in\C$; in particular it has the same zeros as
  $\widehat{\mu_a}$.
\end{enumerate}
\end{proposition}

\begin{proof}[Proof of \ref{item:R12-atoms}]
$\widehat{\mu_a}(z)=\cosh a+\tfrac12(e^{iz}+e^{-iz})$. Since
$\cos(\pi+ia)=-\cos(ia)=-\cosh a$, and $\cos z=\cos z_0$ if and only if
$z\in\pm z_0+2\pi\Z$, the zeros are $\pm(\pi+ia)+2\pi\Z$, which is the
displayed set. The derivative $-\sin z$ takes the values
$\pm\sin(ia)=\pm i\sinh a\ne0$ at $\pm(\pi+ia)$, and by periodicity at
all zeros.
\end{proof}

\begin{proof}[Proof of \ref{item:R12-gauss}]
Positivity, smoothness, and integrability are clear. For real $x$,
$\int e^{ixt}g_\sigma(t)\,dt=e^{-\sigma x^2}$ (the Gaussian integral);
both sides are entire in $x$, since $g_\sigma$ has Gaussian decay, so the
identity holds for all $z\in\C$. A translate $g_\sigma(t\mp1)$ has
transform $e^{\pm iz}e^{-\sigma z^2}$, and summing gives
$e^{-\sigma z^2}\widehat{\mu_a}(z)$. The factor $e^{-\sigma z^2}$ has no
zeros.
\end{proof}

For $0<a<\tfrac12$ the zeros lie in the strip $|\Im z|<\tfrac12$, which
corresponds to the critical strip under $z=i(s-\tfrac12)$. These
examples refute the inference from positivity of the kernel alone; they
do not reproduce the further identities satisfied by the theta kernel of
$\zeta$.

It is natural to expect that Weil positivity proved for tests with small
support extends to all supports, that a compact perturbation of a
positive operator stays positive, and that positivity of a shifted form
$Q-\varepsilon_{\min}I$ on a quotient gives positivity of $Q$; in fact
none of these passages is valid without further input.

\begin{definition}\label{def:R14-sources}
Let $f\in C_c(\R)$. Its \emph{correlation} is
$c_f(x)=\int_\R f(y)\overline{f(y-x)}\,dy$. The \emph{window
hypotheses} of \cite[Theorem~1]{ConnesConsani2020WeilPositivity} are:
multiplicative support of the test function in $[2^{-1/2},2^{1/2}]$,
that is additive support in $[-\tfrac12\log2,\tfrac12\log2]$, and
vanishing of its transform at $i/2$ and at $0$. The \emph{lowest
eigenvector hypotheses} of
\cite[Theorem~5.10]{ConnesConsaniMoscovici2025ZetaTriples} (version~1)
are: the lowest eigenvalue $\varepsilon_{\min}$ of the
finite Weil matrix is simple, it has an even eigenvector $v$, and a
boundary functional $\delta_N$ satisfies $\delta_N(v)\ne0$ (normalized
there to $\delta_N(v)=1$); the construction uses the inner product
defined by $Q-\varepsilon_{\min}I$ on a quotient.
\end{definition}

\begin{proposition}\label{prop:R14}
\leavevmode
\begin{enumerate}[label=(\roman*)]
\item\label{item:R14-compact} Let $\cH$ be a Hilbert space and $P$ an
  orthogonal projection of rank one. Then $-2P$ is compact and $I-2P$
  has the eigenvalue $-1$.
\item\label{item:R14-shift} For every $\varepsilon<0$ there is a bounded
  self-adjoint operator $Q$ on a Hilbert space such that
  $Q-\varepsilon I$ is positive semidefinite and $Q$ is not.
\item\label{item:R14-support} Let $\ell>0$ and $f\in C_c(\R)$ with
  $\operatorname{supp}f\subset[-\ell,\ell]$. Then $c_f(x)=0$ for
  $|x|\ge2\ell$. If moreover $f$ is real and $f>0$ on $(-\ell,\ell)$,
  then $c_f(x)>0$ for $|x|<2\ell$.
\item\label{item:R14-primes} Let $\ell>0$. If
  $\ell\le\tfrac12\log2$, then $c_f(\log n)=0$ for every $f$ as in
  \ref{item:R14-support} and every integer $n\ge2$. If
  $\ell>\tfrac12\log2$, then there is $f$ as in \ref{item:R14-support}
  with $c_f(\log n)>0$ for every integer $n$ with $2\le n<e^{2\ell}$.
\end{enumerate}
\end{proposition}

\begin{proof}[Proof of \ref{item:R14-compact}]
$-2P$ has rank one. If $Pv=v\ne0$, then $(I-2P)v=-v$.
\end{proof}

\begin{proof}[Proof of \ref{item:R14-shift}]
Take $\cH=\C$ and $Q=\varepsilon$; then $Q-\varepsilon I=0$ and $Q<0$.
\end{proof}

\begin{proof}[Proof of \ref{item:R14-support}]
The integrand $f(y)\overline{f(y-x)}$ vanishes unless $y\in[-\ell,\ell]$
and $y\in[x-\ell,x+\ell]$. For $|x|\ge2\ell$ this set has measure zero.
For $|x|<2\ell$ and $f$ real and positive on $(-\ell,\ell)$, the
integrand is continuous, nonnegative, and positive on the nonempty open
interval $(-\ell,\ell)\cap(x-\ell,x+\ell)$.
\end{proof}

\begin{proof}[Proof of \ref{item:R14-primes}]
If $\ell\le\tfrac12\log2$ and $n\ge2$, then $\log n\ge2\ell$; apply
\ref{item:R14-support}. If $\ell>\tfrac12\log2$, take $f(t)=\max(0,
\ell-|t|)$, which is real and positive on $(-\ell,\ell)$; for $2\le
n<e^{2\ell}$ we have $0<\log n<2\ell$.
\end{proof}

The values $c_f(\log n)$, $n\ge2$ a prime power, are those at which the
prime-power terms of Weil's explicit formula
(Theorem~\ref{thm:imp-weil-crit}) evaluate the correlation of a test
function. By \ref{item:R14-primes}, the window hypotheses of
Definition~\ref{def:R14-sources} are exactly the range in which no such
term is seen; enlarging the support brings in every prime power below
$e^{2\ell}$. Positivity of the Weil form on test functions supported in
this prime-free window, which then involves only the archimedean term,
is due to Yoshida \cite[Theorem~1, p.~310]{Yoshida1992Hermitian}, and
was reproved by Bombieri \cite[\S12, Theorem~12, pp.~226--227]%
{Bombieri2000Lincei} and Burnol \cite[Theorem, p.~2]{Burnol2000Explicites};
compare \cite[p.~2 and (149), p.~50]{ConnesConsani2020WeilPositivity}.
The positivity conclusion of
\cite[Theorem~1]{ConnesConsani2020WeilPositivity} is proved under those
hypotheses, and its proof uses the absence of prime-power terms, which
fails for larger supports; extending the conclusion to larger supports
requires new input. Items \ref{item:R14-compact}
and \ref{item:R14-shift} apply to the
construction of \cite[Theorem~5.10]{ConnesConsaniMoscovici2025ZetaTriples}
(version~1): the form $Q-\varepsilon_{\min}I$ is positive on its
quotient even when $\varepsilon_{\min}<0$, and nothing about positivity
of the Weil form itself follows. That construction is carried here with
its lowest eigenvector hypotheses; no assertion is made about later
versions.

The short support in Definition~\ref{def:R14-sources} cannot simply be
removed even for the archimedean contribution alone. In the centered
convention of \cite[(152)--(153), Appendix~B]%
{ConnesConsani2020WeilPositivity}, put
\[
\widehat g(t)=\int_0^\infty g(u)u^{-it}\,\frac{du}{u},\qquad
m(t)=\Re\frac{\Gamma'}{\Gamma}
  \left(\frac14+\frac{it}{2}\right)-\log\pi.
\]
Then $W_\infty(g\star g^*)=(2\pi)^{-1}
\int_\R m(t)|\widehat g(t)|^2\,dt$, where
$g^*(u)=\overline{g(u^{-1})}$ and $W_\infty=-W_\R$ in that source.

\begin{proposition}\label{prop:arch-window-limit}
There are smooth compactly supported $g$ on $\R_+^\times$ with
$\widehat g(0)=\widehat g(i/2)=\widehat g(-i/2)=0$ and
$W_\infty(g\star g^*)<0$.
\end{proposition}

\begin{proof}
Choose nonzero $\phi\in C_c^\infty((-1,1))$ with $\int\phi=0$ and set
\[
a_L(x)=L^{-1/2}\phi(x/L),\qquad
g_L(e^x)=(1-4\partial_x^2)a_L(x).
\]
Thus $\operatorname{supp}g_L\subset[e^{-L},e^L]$ and
$\widehat g_L(t)=(1+4t^2)\sqrt L\,\widehat\phi(Lt)$, also for
complex $t$. The three asserted zeros follow. Substituting $y=Lt$,
\[
W_\infty(g_L\star g_L^*)=\frac1{2\pi}\int_\R
m(y/L)(1+4y^2/L^2)^2|\widehat\phi(y)|^2\,dy
\longrightarrow m(0)\|\phi\|_2^2.
\]
The convergence follows from the logarithmic growth of $m$ and rapid
decay of $\widehat\phi$. The digamma value at $1/4$ gives
$m(0)=-\gamma-\pi/2-3\log2-\log\pi<0$.
\end{proof}

This is a failure of positivity for the single archimedean term on a
larger test class. The full Weil form also contains prime terms; this
proposition neither contradicts the short-window theorem nor gives an
off-line zero of $\zeta$.

It is natural to expect that convergence of trial functions, a small
residual, or a small $L^2$ error on growing supports identifies the
eigentransforms required by a spectral approach; in fact each of these
is compatible with failure of the intended conclusion.

\begin{definition}\label{def:Fminus}
For $g\in L^2(\R)$ with compact support, its \emph{negative-sign Fourier
transform} is the entire function
$\mathcal F_-g(z)=\int_\R g(t)e^{-izt}\,dt$.
\end{definition}

\begin{definition}\label{def:R16-examples}
\begin{enumerate}[label=(\alph*)]
\item For $\delta>0$, let $A_\delta=\diag(0,\delta,1)$ on $\C^3$,
  $k=(1,1,0)$, and $P_\delta$ the orthogonal projection onto
  $\ker A_\delta$.
\item For $b>0$ and $\ell\ge1$, let
  $f_\ell=e^{-b\ell}(1_{[\ell-1,\ell]}+1_{[-\ell,-\ell+1]})$.
\item Let $A$ be the operator $Ae_n=n^2e_n$ on $\ell^2(\N)$, $n\ge1$,
  with domain $\{\sum c_ne_n:\sum n^4|c_n|^2<\infty\}$.
\end{enumerate}
\end{definition}

\begin{proposition}\label{prop:R16}
Let $A_\delta$, $k$, $P_\delta$, $f_\ell$, and $A$ be as in
Definition~\ref{def:R16-examples}.
\begin{enumerate}[label=(\roman*)]
\item\label{item:R16-gap} For every $\delta>0$,
  $\|A_\delta k\|=\delta$ and $\|k-P_\delta k\|=1$.
\item\label{item:R16-strip} Let $b>0$ and $\ell\ge1$. Then
  $f_\ell$ is even, $\|f_\ell\|_2=\sqrt2e^{-b\ell}$, and
  \[
  \mathcal F_-f_\ell(ib)=\frac{1-e^{-b}}b+e^{-2b\ell}\frac{e^b-1}b
  \longrightarrow\frac{1-e^{-b}}b>0\quad(\ell\to\infty).
  \]
  Moreover, for every $z$ with $\Im z=b$,
  $\sup\{|\mathcal F_-g(z)|:\operatorname{supp}g\subset[-\ell,\ell],
  \ \|g\|_2\le1\}=\sqrt{\sinh(2b\ell)/b}$.
\item\label{item:R16-unbounded} The vectors $x_n=e_n/n$ satisfy
  $\|x_n\|\to0$ and $\|Ax_n\|=n$.
\end{enumerate}
\end{proposition}

\begin{proof}[Proof of \ref{item:R16-gap}]
$A_\delta k=(0,\delta,0)$. Since $\delta>0$, $\ker A_\delta=\C(1,0,0)$,
so $P_\delta k=(1,0,0)$ and $k-P_\delta k=(0,1,0)$.
\end{proof}

\begin{proof}[Proof of \ref{item:R16-strip}]
Evenness and the norm are immediate. Since $e^{-i(ib)t}=e^{bt}$,
\[
\mathcal F_-f_\ell(ib)=e^{-b\ell}\Bigl(\int_{\ell-1}^{\ell}e^{bt}dt
+\int_{-\ell}^{-\ell+1}e^{bt}dt\Bigr)
=\frac{1-e^{-b}}b+e^{-2b\ell}\frac{e^{b}-1}b.
\]
For $\Im z=b$, $|e^{-izt}|=e^{bt}$, so by the Cauchy--Schwarz inequality
$|\mathcal F_-g(z)|\le\|g\|_2(\int_{-\ell}^{\ell}e^{2bt}dt)^{1/2}$, with
equality for $g$ proportional to $\overline{e^{-izt}}$ on $[-\ell,\ell]$;
and $\int_{-\ell}^{\ell}e^{2bt}dt=\sinh(2b\ell)/b$.
\end{proof}

\begin{proof}[Proof of \ref{item:R16-unbounded}]
$\|x_n\|=1/n$ and $Ax_n=ne_n$.
\end{proof}

In \ref{item:R16-gap} the residual tends to zero while the distance to
the target line does not; the ratio of residual to spectral gap, here
exactly one, is the relevant quantity. In \ref{item:R16-strip},
unweighted $L^2$ convergence on growing supports does not control
transforms off the real axis; the evaluation norm grows like
$e^{b\ell}$. In \ref{item:R16-unbounded}, approximation of vectors does
not control an unbounded operator without graph-norm estimates. A
locally uniform nonzero limit of entire functions with only real zeros
has only real zeros (Hurwitz); finite families of trial functions with
real zeros do not by themselves supply such a limit.

\subsection{Coefficients and the local--global passage}
\label{subsec:neg-coefficients}

It is natural to expect that Deninger's obstruction to real coefficients
is either irrelevant to the axioms or rules out every candidate; in fact
it is a theorem with precise hypotheses, and a separate algebraic
obstruction already appears over finite fields.

\begin{remark}\label{rem:R5-import}
By Theorem~\ref{thm:imp-deninger-nogo}, imported from
\cite[Theorem~2.13]{Deninger2022NoWeil}, no functorial theory with real
coefficients on Arakelov compactifications of spectra of rings of
integers satisfies all of the properties 2.1--2.12 of
\cite{Deninger2022NoWeil}. The proof there uses a quaternionic Galois
extension with a symplectic irreducible representation of root number
$-1$, whose central $L$-order is odd, whereas a symplectic irreducible
occurs with even multiplicity in a complexified real representation. It
excludes the naive untwisted real leafwise candidate only insofar as
that candidate satisfies all of the properties 2.1--2.12. It does not
exclude complex theories, quaternionic coefficient targets, the real
Weil--Arakelov complexes of
\cite[\S3]{FlachMorin2020DeningerWeilArakelov}, or a self-adjoint
operator attached to $\zeta$ alone; Deninger himself notes that the
analogy requires a twisted complex version
\cite[p.~6]{Deninger2022NoWeil}.
\end{remark}

A separate algebraic obstruction appears over finite fields.

\begin{proposition}[Serre]\label{prop:R5}
Let $\mathbb H$ be Hamilton's quaternion algebra over $\R$. There is no
unital ring antihomomorphism $\mathbb H\to M_2(\R)$.
\end{proposition}

\begin{proof}
Let $\varphi$ be such a map and put $\psi(x)=\varphi(x)^T$; then $\psi$
is a unital ring homomorphism, since
$\psi(xy)=(\varphi(y)\varphi(x))^T=\psi(x)\psi(y)$. Put $X=\psi(i)$ and
$Y=\psi(j)$. Additivity and $\psi(1)=I$ give $\psi(-1)=-I$, so
$X^2=Y^2=-I$ and $XY=\psi(ij)=\psi(-ji)=-YX$. The $\R$-algebra $\mathbb H$
has basis $1,i,j,ij$ with multiplication table determined by
$i^2=j^2=-1$ and $ij=-ji$; the matrices $I,X,Y,XY$ satisfy the same
relations, hence the same table. So the $\R$-linear map
$\chi:\mathbb H\to M_2(\R)$, $\chi(c_0+c_1i+c_2j+c_3ij)=c_0I+c_1X+c_2Y
+c_3XY$, is a unital $\R$-algebra homomorphism. Its kernel is a
two-sided ideal not containing $1$; as $\mathbb H$ is a division ring,
$\chi$ is injective, hence bijective since both sides have dimension
four. Then $M_2(\R)$ would be a division ring, but
$\bigl(\begin{smallmatrix}0&1\\0&0\end{smallmatrix}\bigr)^2=0$.
\end{proof}

Proposition~\ref{prop:R5} is the algebraic core of Serre's observation
that a Weil cohomology cannot have coefficients in $\R$ (or $\Q$, or
$\Q_p$) \cite[(2.2) and fn.~25, p.~27]{Milne2016RHFF},
\cite[\S1, p.~1]{Deninger2022NoWeil}: for a supersingular
elliptic curve over $\overline{\F}_p$ the endomorphism algebra tensored
with $\R$ is $\mathbb H$ (Deuring; see
\cite[Theorem~V.3.1, pp.~144--145]{Silverman2009AEC}), and a
contravariant functorial two-dimensional real $H^1$ would give a unital
antihomomorphism $\mathbb H\to M_2(\R)$. This is why the choice of
coefficient category is substantive.

It is natural to expect that reduced leafwise cohomology of a foliation
is represented by leafwise harmonic forms, as in Hodge theory, and that
dynamical traces pass to it; in fact this fails without geometric
hypotheses.

\begin{definition}\label{def:reduced-leafwise}
Let $\mathcal F$ be a smooth foliation of a closed manifold $M$ with
leafwise de Rham complex $(\mathcal A^\bullet_{\mathcal F}(M),
d_{\mathcal F})$ of smooth leafwise forms. The \emph{reduced leafwise
cohomology} is $\overline H^j_{\mathcal F}(M)=\ker d_{\mathcal F}/
\overline{\operatorname{im}d_{\mathcal F}}$, the closure being taken in
the smooth Fr\'echet topology. Given a leafwise metric, a leafwise form
$\omega$ is \emph{harmonic} if $d_{\mathcal F}\omega=0=d_{\mathcal
F}^*\omega$.
\end{definition}

\begin{remark}\label{rem:R6}
By Theorem~\ref{thm:imp-ds}, imported from
\cite[Proposition~1.2 and~(12), p.~214]%
{DeningerSinghof2001Counterexample}, there exist a closed
three-dimensional Heisenberg nilmanifold $M$, a smooth one-dimensional
foliation $\mathcal F$ of $M$, and a leafwise metric such that
$\overline H^1_{\mathcal F}(M)$
(Definition~\ref{def:reduced-leafwise}) is infinite-dimensional while
the space of smooth leafwise harmonic $1$-forms is one-dimensional.
\end{remark}

In the source, \S2 also
obstructs a general passage of dynamical traces to reduced cohomology.
The positive Hodge-type results for foliations assume a Riemannian
foliation with a bundle-like metric. A leafwise candidate for
\ref{ax:AW4}--\ref{ax:AW5} must therefore specify its topology and
geometric hypotheses.

It is natural to expect that local determinant formulas at the finite
and infinite places glue to a global cohomology realizing $\xi$; in fact
the simplest gluing, the direct sum of local spectral spaces, keeps the
local divisors, which are disjoint from the zeros of $\xi$.

\begin{definition}\label{def:R7-local}
For a prime $p$, the \emph{local spectral core at $p$} is
$V_p=\bigoplus_{n\in\Z}\C e_{p,n}$ with $\Theta_pe_{p,n}=(2\pi in/\log
p)e_{p,n}$; it is the logarithmic suspension
(Definition~\ref{def:log-susp}) of $V=\C$, $\mathrm F=1$, with $q=p$.
For a set $\mathcal Q$ of primes let
$V_{\mathcal Q}=\bigoplus_{p\in\mathcal Q}V_p$ (algebraic direct sum)
with $\Theta_{\mathcal Q}=\bigoplus_p\Theta_p$.
\end{definition}

\begin{proposition}\label{prop:R7}
Let $\mathcal Q$ be a nonempty set of primes and let
$(V_{\mathcal Q},\Theta_{\mathcal Q})$ be as in
Definition~\ref{def:R7-local}.
\begin{enumerate}[label=(\roman*)]
\item\label{item:R7-spec} For each $p$, the zeros of $1-p^{-s}$ are
  exactly the eigenvalues of $\Theta_p$, each simple. The eigenvalues of
  $\Theta_{\mathcal Q}$ form $\bigcup_{p\in\mathcal Q}(2\pi i/\log p)\Z
  \subset i\R$; the eigenvalue $0$ has multiplicity $\#\mathcal Q$ and
  every nonzero eigenvalue is simple.
\item\label{item:R7-core} If $\mathcal Q$ is infinite,
  $(V_{\mathcal Q},\Theta_{\mathcal Q})$ is not a spectral core.
\item\label{item:R7-xi} No eigenvalue of $\Theta_{\mathcal Q}$ is a zero
  of $\xi$.
\item\label{item:R7-euler} Let $s\in\C$ with $p^s\ne1$ for every prime
  $p$. The series $\sum_p|(1-p^{-s})^{-1}-1|$ converges if and only if
  $\Re s>1$.
\end{enumerate}
\end{proposition}

\begin{proof}[Proof of \ref{item:R7-spec}]
The first assertion is
Proposition~\ref{prop:suspension}\ref{item:susp-zeros} with $\alpha=1$.
If $2\pi in/\log p=2\pi im/\log q$ with primes $p\ne q$, then $n\log
q=m\log p$, so $q^n=p^m$, and unique factorization forces $n=m=0$.
\end{proof}

\begin{proof}[Proof of \ref{item:R7-core}]
The generalized eigenspace at $0$ contains the linearly independent
vectors $e_{p,0}$, $p\in\mathcal Q$, so it is infinite-dimensional.
\end{proof}

\begin{proof}[Proof of \ref{item:R7-xi}]
The eigenvalues lie on $\Re s=0$, while the zeros of $\xi$ are the
nontrivial zeros of $\zeta$, which lie in $0<\Re s<1$
(Theorem~\ref{thm:imp-riemann}).
\end{proof}

\begin{proof}[Proof of \ref{item:R7-euler}]
Write $\sigma=\Re s$ and $a_p=(1-p^{-s})^{-1}-1=1/(p^s-1)$, so
$|p^s|=p^\sigma$. If $\sigma>1$, then
$|a_p|\le1/(p^\sigma-1)\le2p^{-\sigma}$ and $\sum_pp^{-\sigma}$
converges. If $0<\sigma\le1$, then $|a_p|\ge1/(p^\sigma+1)\ge
p^{-\sigma}/2\ge1/(2p)$, and $\sum_p1/p$ diverges (Euler). If
$\sigma\le0$, then $|a_p|\ge1/(p^\sigma+1)\ge\tfrac12$.
\end{proof}

Proposition~\ref{prop:R7} concerns only the direct-sum construction. The
general statement that local determinant formulas do not glue by
themselves is a scope remark: Hesselholt's formula
\cite[Theorem~A]{Hesselholt2016THHZeta} concerns smooth proper schemes
over one finite field, and the archimedean cyclic formula
\cite[Theorem~1.1]{ConnesConsani2014CyclicSerre} concerns gamma factors
of a smooth projective variety at one infinite place. Their categories,
gradings, coefficient rings, and operators differ, and neither supplies
maps identifying local objects, a descent datum, or a global positive
pairing. By \ref{item:R7-euler}, the global divisor of $\zeta$ arises
only after analytic continuation of a product which converges absolutely
only for $\Re s>1$; interchanging infinite determinant products with
analytic continuation requires a theorem. A restricted product, a
descent complex, or a quotient might be meaningful, but none is specified
by the analogy.

\begin{remark}\label{rem:N-composition}
No composition of the comparison theorems
\ref{thm:imp-h17}, \ref{thm:imp-bms1}, \ref{thm:imp-bms2},
\ref{thm:imp-bs}, \ref{thm:imp-cc14}, \ref{thm:imp-a20},
\ref{thm:imp-a19}, and \ref{thm:imp-sb26} produces the comparison asked
for in Question~\ref{q:full-comparison}. Their sources are,
respectively: smooth proper schemes over a finite field; proper smooth
formal schemes over $\cO_C$; $p$-adic completions of smooth
$\cO_C$-algebras, and quasisyntomic rings; smooth $p$-adic formal
schemes over $A/I$ for a bounded prism $(A,I)$; smooth projective
varieties over a number field, at one infinite place; smooth schemes
over the residue field of a mixed-characteristic valued field;
differential modules over a differential field of characteristic zero;
and arc-motives over an algebraically closed nondiscrete Banach field.
A
composite of comparison maps has the source of its first map, so none
takes $\Xbar_\Q$ with a flow as input, and none outputs a Hilbert space
with an operator whose spectrum is the divisor of $\xi$. New
constructions are needed on both ends.
\end{remark}

It is natural to expect that Ayoub's conservativity theorem is available
as a bridge among realizations; in fact the author's own paper list
records a gap in it.

\begin{remark}\label{rem:R8}
The author's list of papers \ayoubpage{} identifies a gap in
\cite[Theorem~8.3.10]{Ayoub2018FoliatedConservativity} and treats the
conservativity conjecture as open pending a correction. Nothing in these
notes uses that statement. This does not affect the results imported as
Theorems~\ref{thm:imp-a20} and~\ref{thm:imp-a19}.
\end{remark}

\subsection{Findings about specific versions of the literature}
\label{subsec:neg-literature}

It is natural to expect that de Branges's positivity conditions for the
space attached to $\xi$ hold, as would follow from a successful
execution of his programme; in fact Conrey and Li exhibit explicit
failures.

\begin{definition}\label{def:conrey-li}
For an entire function $E$ satisfying the HB inequality, the
\emph{de Branges space} $\mathcal H(E)$ is the space of entire $G$ with
$G/E,\,G^\#/E\in H^2(\Cp)$, normed by $\|G\|^2=\int_\R|G/E|^2\,dx$. The
\emph{de Branges positivity conditions} are the shift-positivity
conditions for $\mathcal H(E)$ with $E(z)=\xi(1-iz)$, and for the
associated space of reciprocal functions, as formulated in
\cite[\S3.1, (3.1)--(3.4), pp.~6--7]{ConreyLi2000Positivity}.
\end{definition}

\begin{remark}\label{rem:R13}
Conrey and Li show, by high-precision numerical evaluation, that the de
Branges positivity conditions of Definition~\ref{def:conrey-li} do not
all hold (Theorem~\ref{thm:imp-conrey-li}). Their witnesses in
\cite[\S3.1]{ConreyLi2000Positivity} are numerical values, not
interval-certified bounds, and the conclusion is recorded here at that
strength. For the condition on the reciprocal space they also report a
proof by Sarnak that involves no numerical calculation
\cite[Remark, p.~9]{ConreyLi2000Positivity}.
\end{remark}

In the normalization
$\xi=\tfrac12s(s-1)\Lambda$ of Definition~\ref{def:zeta}, and with
$\rho=\tfrac12+111.029535543\ldots i$ (the $34$th zero), they amount to
\[
-\Re\{\xi'(\rho)\xi(1+\rho)\}\approx-1.34727512680\cdot10^{-69},
\qquad
\Re\{\xi(1+282i)/\xi(2+282i)\}\approx-0.000131957.
\]
Conrey and Li normalize $\xi=s(s-1)\Lambda$, which multiplies the first
value by four, giving $-5.38910050718\cdot10^{-69}$; there is no complex
conjugation in the product. An evaluation at $65$ significant digits
with the \texttt{mpmath} library, recorded in the research record of
this repository, reproduces both values and signs; it is a
high-precision evaluation, not an interval-certified one. These
witnesses refute the specific positivity conditions, not de Branges
theory, not the identity of Suzuki discussed in
Section~\ref{sec:debranges}, and not all uses of canonical systems.

It is natural to expect that the theta kernel displayed in version~1 of
the zeta spectral triples paper transforms to $F$, and that a Gaussian
tail can be dropped when trial functions are truncated; in fact the
displayed kernel transforms to $F/4$, and the tail is positive.

\begin{definition}\label{def:ccm-kernel}
The \emph{displayed kernel} is
$k_0(u)=\tfrac\pi2u^2(2\pi u^2-3)e^{-\pi u^2}$ for $u>0$, and its
\emph{periodization} is the function
$\mathcal Ek_0(u)=u^{1/2}\sum_{n\ge1}k_0(nu)$. The
\emph{Mellin transform} of $k_0$ is
$M_{k_0}(s)=\int_0^\infty k_0(u)u^{s-1}\,du$. For $\lambda>0$ the
\emph{truncation} is $k_{0,\lambda}=k_0\cdot1_{(0,\lambda]}$.
\end{definition}

\begin{proposition}\label{prop:R15}
Let $k_0$, $\mathcal Ek_0$, $M_{k_0}$, and $k_{0,\lambda}$ be as
in Definition~\ref{def:ccm-kernel}.
\begin{enumerate}[label=(\roman*)]
\item\label{item:R15-mellin} For $\Re s>0$ the integral $M_{k_0}(s)$
  converges absolutely and
  $M_{k_0}(s)=\tfrac18s(s-1)\pi^{-s/2}\Gamma(s/2)$.
\item\label{item:R15-xi} For $\Re s>1$,
  $\int_0^\infty\mathcal Ek_0(u)\,u^{s-1/2}\,\frac{du}u$ converges
  absolutely and equals $\zeta(s)M_{k_0}(s)=\tfrac14\xi(s)$.
\item\label{item:R15-tail} Let $\lambda\ge1$. For every $u>0$,
  $\sum_{n\ge1}k_0(nu)-\sum_{n\ge1}k_{0,\lambda}(nu)
  =\sum_{n>\lambda/u}k_0(nu)>0$.
\end{enumerate}
\end{proposition}

\begin{proof}[Proof of \ref{item:R15-mellin}]
For $\Re\beta>0$, the substitution $v=\pi u^2$, under which
$u=(v/\pi)^{1/2}$ and $du/u=dv/(2v)$, gives
\[
\int_0^\infty u^{\beta}e^{-\pi u^2}\frac{du}u
=\frac12\int_0^\infty\Bigl(\frac v\pi\Bigr)^{\beta/2}e^{-v}\frac{dv}v
=\frac12\pi^{-\beta/2}\Gamma(\beta/2),
\]
the integrals converging absolutely because
$|u^\beta|=u^{\Re\beta}$ and $\Re\beta>0$. Now
$k_0(u)u^{s-1}=\bigl(\pi^2u^{s+4}-\tfrac{3\pi}2u^{s+2}\bigr)
e^{-\pi u^2}u^{-1}$, and $\Re(s+2)>0$, so
\[
M_{k_0}(s)=\frac{\pi^2}2\pi^{-\frac s2-2}\Gamma\bigl(\tfrac s2+2\bigr)
-\frac{3\pi}4\pi^{-\frac s2-1}\Gamma\bigl(\tfrac s2+1\bigr)
=\frac12\pi^{-s/2}\Bigl[\Gamma\bigl(\tfrac s2+2\bigr)
-\tfrac32\Gamma\bigl(\tfrac s2+1\bigr)\Bigr].
\]
Using $\Gamma(\tfrac s2+2)=(\tfrac s2+1)\tfrac s2\Gamma(\tfrac s2)$ and
$\Gamma(\tfrac s2+1)=\tfrac s2\Gamma(\tfrac s2)$, the bracket is
$\tfrac s2(\tfrac s2-\tfrac12)\Gamma(\tfrac s2)=\tfrac14s(s-1)
\Gamma(\tfrac s2)$.
\end{proof}

\begin{proof}[Proof of \ref{item:R15-xi}]
Let $\sigma=\Re s>1$. The integrand is
$\sum_nk_0(nu)u^{s-1}$, and by Tonelli's theorem and the substitution
$v=nu$,
$\sum_n\int_0^\infty|k_0(nu)|u^{\sigma-1}du=\sum_nn^{-\sigma}
\int_0^\infty|k_0(v)|v^{\sigma-1}dv<\infty$, the last integral being finite
by the proof of \ref{item:R15-mellin}. By Fubini's theorem the integral
equals $\sum_nn^{-s}M_{k_0}(s)=\zeta(s)M_{k_0}(s)$, and by
\ref{item:R15-mellin} this is
$\tfrac18s(s-1)\Lambda(s)=\tfrac14\xi(s)$ (Definition~\ref{def:zeta}).
\end{proof}

\begin{proof}[Proof of \ref{item:R15-tail}]
The difference is the sum over $n$ with $nu>\lambda$. For $v>0$,
$k_0(v)>0$ if and only if $v>(3/2\pi)^{1/2}$, and $(3/2\pi)^{1/2}<1\le
\lambda$; so every term with $nu>\lambda$ is positive, and there are
infinitely many. The series converges by the Gaussian decay of $k_0$.
\end{proof}

The kernel and its periodization are classical: summed over $n$, it is,
up to normalization, the kernel of Riemann's second integral for $\xi$
\cite[p.~4]{Riemann1859Wilkins}, and in the variable $u=e^t$ it is
P\'olya's kernel $\Phi$ \cite[(7.2)--(7.3), p.~220]{DeBruijn1950Roots}.
With $s=\tfrac12-iz$, \ref{item:R15-xi} gives the transform $F(z)/4$
(Definition~\ref{def:zeta}). In
\arxivlink{ConnesConsaniMoscovici2025ZetaTriples}{2511.22755v1},
displays (7.1)--(7.2) present this kernel and its periodization; with the
multiplicative Fourier transform used there, the transform is $F/4$
rather than $F$. Multiplying the kernel by four repairs the
normalization and leaves the zeros unchanged. The finding concerns that
version and those conventions only; no claim is made about later
versions. Item \ref{item:R15-tail} shows that truncating the support of
the kernel leaves a positive, nonzero tail, which a convergence argument
for truncated trial functions must bound explicitly (a Gaussian bound
suffices for this purpose); the detailed prolate estimate is not
redeveloped here.

It is natural to expect that the printed formulas of Suzuki's comparison
can be combined literally; in fact the printed even extension to
negative time annihilates every even test function.

\begin{definition}\label{def:bump}
The \emph{standard bump} is the function $\psi_0$ on $\R$ with
$\psi_0(t)=\exp(-1/(1-t^2))$ for $|t|<1$ and $\psi_0(t)=0$ otherwise.
\end{definition}

\begin{proposition}\label{prop:R17}
Let $\Omega\subset\C$ be a set with $\overline\Omega=\Omega$, and let
$(S_t)_{t\in\R}$ be a family of functions $\Omega\to\C$ such that
$S_{-t}=S_t$ for every $t$ and $t\mapsto S_t(z)$ is locally integrable
on $\R$ for every $z\in\Omega$. Let $S_t^\#$ be the reflection of
$S_t$ (Definition~\ref{def:hb}). Then:
\begin{enumerate}[label=(\roman*)]
\item\label{item:R17-zero} for every even $\psi\in C_c^\infty(\R)$
  and every $z\in\Omega$,
  \[
  \int_\R S_t^\#(z)\,i\psi'(t)\,dt=0;
  \]
\item\label{item:R17-bump} the standard bump $\psi_0$ of
  Definition~\ref{def:bump} is a nonzero even element of
  $C_c^\infty(\R)$.
\end{enumerate}
\end{proposition}

\begin{proof}[Proof of \ref{item:R17-zero}]
Fix $z\in\Omega$ and put $G(t)=S_t^\#(z)=\overline{S_t(\bar z)}$,
which is defined since $\bar z\in\Omega$ and is locally integrable;
then $G(-t)=\overline{S_{-t}(\bar z)}=G(t)$. Since $\psi$ is even,
$\psi'$ is odd, and the substitution $t\mapsto-t$ gives
$\int_{-\infty}^0G\psi'=\int_0^\infty G(u)\psi'(-u)\,du
=-\int_0^\infty G\psi'$; the two halves cancel.
\end{proof}

\begin{proof}[Proof of \ref{item:R17-bump}]
The function $t\mapsto\exp(-1/(1-t^2))$ is smooth on $(-1,1)$ and all
its derivatives tend to $0$ as $|t|\to1^-$, so $\psi_0$ is smooth; it
is supported in $[-1,1]$, even, and $\psi_0(0)=e^{-1}\ne0$.
\end{proof}

\begin{definition}\label{def:suzuki-lit}
In \cite{Suzuki2023WeilHilbert} (version~3), with $E=E_F=F+iF'$ and
$\vartheta=E^\#/E$, functions $S_t=\tfrac i2(1+\vartheta^\#)P_t$ and
$P_t$ are defined for $t\ge0$ by (1.5)--(1.6) on p.~3, and $P_t$ is
expressed as a series over the zeros of $F$ in (3.2) on p.~7
(Proposition~3.1 of the source); the source extends them to $t<0$
evenly, $S_{-t}=S_t$ after (1.6) on p.~3 and $P_{-t}=P_t$ after (3.2)
on p.~7.
The \emph{literal transform} of (1.7)--(1.8) is
$T_{\mathrm{lit}}\psi=\int_\R S_t^\#\,i\psi'(t)\,dt$ for
$\psi\in C_c^\infty(\R)$ (denoted $\widehat P_{D\psi}$ there), and
(3.9) on p.~10 sets $\|\psi\|_0=\pi^{-1/2}\|T_{\mathrm{lit}}\psi\|_{
L^2(\R)}$, which Lemma~3.2 on p.~10 asserts to be a norm.
\end{definition}

\begin{remark}\label{rem:R17-source}
Let $\Omega$ be the complement in $\C$ of the zeros of $F$, $E$, and
$E^\#$; then $\overline\Omega=\Omega$. For $z\in\Omega$ and $t\ge0$,
the printed $P_t(z)$ equals the series (3.2) of the source,
$\sum_\gamma m_\gamma(e^{-i\gamma t}-1)/(\gamma(z-\gamma))$, over the
distinct zeros $\gamma$ of $F$ with multiplicities $m_\gamma$. Since
$|\Im\gamma|<\tfrac12$, one has $|e^{-i\gamma t}-1|\le2e^{|t|/2}$, and
$|z-\gamma|\ge|\gamma|/2$ for all but finitely many $\gamma$; so the
terms are bounded by $4e^{|t|/2}m_\gamma|\gamma|^{-2}$ for all but
finitely many $\gamma$, which is summable by \eqref{eq:sum-r2}. The
series therefore converges uniformly for $t$ in compact sets, and
$t\mapsto S_t(z)$ is continuous on $[0,\infty)$; after the even
extension it is locally integrable on $\R$. Applied to the printed
definitions on this $\Omega$, Proposition~\ref{prop:R17} gives
$T_{\mathrm{lit}}\psi=0$ on $\Omega$ for every even
$\psi\in C_c^\infty(\R)$, in particular for $\psi_0$. Hence the literal
expression (3.9) of
Definition~\ref{def:suzuki-lit} vanishes on a nonzero test function: it
is a degenerate seminorm, not the norm asserted in the literal
Lemma~3.2. This uses no hypothesis on the zeros of $\zeta$. It does not
refute \cite[Theorem~1.1]{Suzuki2023WeilHilbert}, a Hilbert-space
statement proved independently of this transform. The all-real
expansion (3.6)--(3.8) of the source does not use the even extension.
The author has corrected the even extension in a later paper: the
factor $i$ in (1.5) is to be replaced by $-i$, and the extension to
$t<0$ by $P_t(z)=\overline{P_{-t}(\bar z)}$
\cite[p.~27]{Suzuki2026WeilScrew}; this corrected definition is the
reconstruction of Section~\ref{subsec:suzuki}. The finding above
concerns the printed text of version~3 of \cite{Suzuki2023WeilHilbert},
which retains the even extension.
\end{remark}

\begin{remark}\label{rem:R17-rh}
There is also an incompatibility under RH. Suppose RH holds. Take a
nonnegative even $\psi\in C_c^\infty(\R)$, not identically zero, with
support so close to $0$ that its cosine transform is positive at the
ordinate of one zero of $\zeta$. In the spectral expression of the Weil
form $Q_W$ (Definition~\ref{def:weil-form}), every term is then a
nonnegative squared modulus and that term is positive, so
$Q_W(\psi)>0$. Since $T_{\mathrm{lit}}\psi=0$, the literal transform
cannot satisfy the norm identity $\|T\psi\|^2=\pi Q_W(\psi)$ at $\psi$.
This observation assumes RH and is not a refutation of an equivalence
with RH.
\end{remark}

\subsection{Register of corrected expectations}
\label{subsec:neg-register}

The following table lists the distinct expectations corrected above,
each mapped to the numbered statements that correct it. Repeated
mentions of one obstruction count once; positive approximation
statements are not counted as refutations.

\begin{longtable}{@{}l>{\raggedright\arraybackslash}p{0.56\textwidth}
>{\raggedright\arraybackslash}p{0.26\textwidth}@{}}
\toprule
Entry & Expectation corrected & Corrected by\\
\midrule
\endfirsthead
\toprule
Entry & Expectation corrected & Corrected by\\
\midrule
\endhead
\bottomrule
\endfoot
O01 & a formal arithmetic cohomology is already a positive realization of
  the zeros & Prop.~\ref{prop:R1}; Sections~\ref{sec:axioms},
  \ref{sec:purity}\\
O02 & skew-symmetry controls the whole operator spectrum &
  Prop.~\ref{prop:R4}\\
O03 & duality and a quaternionic star force purity & Prop.~\ref{prop:R2}\\
O04 & a separately chosen positive star also commutes with the flow &
  Prop.~\ref{prop:R2-star}\\
O05 & integer-twist cones see summands with nonreal spectrum &
  Prop.~\ref{prop:R3}\\
O06 & acyclic blocks preserve determinant and special-value data &
  Prop.~\ref{prop:R3}\\
O07 & HB (logarithmic) positivity extends downwards automatically &
  Prop.~\ref{prop:R11}\\
O08 & the HB test or a scalar realization supplies simplicity or
  multiplicities & Prop.~\ref{prop:R9}\\
O09 & inverse canonical-system theory proves its own positivity
  hypothesis & Prop.~\ref{prop:R11}; Section~\ref{sec:debranges}\\
O10 & the displayed version-1 kernel transforms to $F$ with no factor
  four & Prop.~\ref{prop:R15}\\
O11 & a positive Fourier kernel forces real zeros & Prop.~\ref{prop:R12}\\
O12 & compactly supported trials allow omission of the Gaussian tail &
  Prop.~\ref{prop:R15}\\
O13 & convergence of trial transforms identifies ground eigentransforms
  & Prop.~\ref{prop:R16}\\
O14 & restricted-support positivity already covers all primes and
  supports & Prop.~\ref{prop:R14}\\
O15 & a compact error preserves positivity & Prop.~\ref{prop:R14}\\
O16 & shifted quotient positivity proves the original positivity;
  hypotheses can be dropped & Prop.~\ref{prop:R14}\\
O17 & $L^2$ convergence on growing supports implies convergence in a
  complex strip & Prop.~\ref{prop:R16}\\
O18 & a small residual identifies an eigenline without a gap bound &
  Prop.~\ref{prop:R16}\\
O19 & prolate approximation controls an unbounded nonlocal arithmetic
  operator & Prop.~\ref{prop:R16}\\
O20 & the real-coefficient no-go is either irrelevant or universal &
  Rem.~\ref{rem:R5-import}; Prop.~\ref{prop:R5}\\
O21 & arbitrary smooth leafwise cohomology has harmonic representatives
  and admits passage of traces & Rem.~\ref{rem:R6}\\
O22 & the finite-field, local, and motivic results give a global
  realization for RH & Prop.~\ref{prop:R7};
  Rem.~\ref{rem:N-composition}; Sections~\ref{sec:candidates},
  \ref{sec:imported}\\
O23 & positivity of heights on selected classes proves full spectral
  positivity & Prop.~\ref{prop:R1}\\
\end{longtable}

The findings recorded in Remarks~\ref{rem:R8} and~\ref{rem:R13} and in
Propositions~\ref{prop:R10} and~\ref{prop:R17} are corrections to
possible uses of more recent sources and correspond to no entry of the
table; Proposition~\ref{prop:R9} refines O08 to a statement about
cyclic operators. In total, the seventeen findings R1--R17 are carried
as numbered statements: fourteen of them (R1--R5, R7, R9--R12, and
R14--R17) by propositions, in fifteen environments because R2 is split
into Propositions~\ref{prop:R2} and~\ref{prop:R2-star}, and three (R6,
R8, and R13) by Remarks~\ref{rem:R6}, \ref{rem:R8}, and~\ref{rem:R13};
the imported part of R5 is recorded in Remark~\ref{rem:R5-import}.
Every entry O01--O23 of the table is corrected by at least one numbered
statement.

The second table lists the distinct expectations corrected in the
statements about the main theorem, its calibration, and the class of
zeta functions (Sections~\ref{sec:finite-field}, \ref{sec:axioms},
and~\ref{sec:class}). Entries that sharpen an entry of the first table
say so.

\begin{longtable}{@{}l>{\raggedright\arraybackslash}p{0.56\textwidth}
>{\raggedright\arraybackslash}p{0.26\textwidth}@{}}
\toprule
Entry & Expectation corrected & Corrected by\\
\midrule
\endfirsthead
\toprule
Entry & Expectation corrected & Corrected by\\
\midrule
\endhead
\bottomrule
\endfoot
O24 & all zeros of the zeta function of a higher-dimensional arithmetic
  scheme lie on one line & Prop.~\ref{prop:P1Z};
  Thm.~\ref{thm:A-schemes}\\
O25 & the total zeta function determines the spectra in each degree,
  and degrees may cancel whatever the positivity &
  Prop.~\ref{prop:no-positivity-cancel}; Thm.~\ref{thm:A-schemes}\\
O26 & the correct spectrum with perfect duality is already a
  nontrivial arithmetic statement (sharpens O01) &
  Props.~\ref{prop:trivial-model}, \ref{prop:content}\\
O27 & either half of \ref{ax:AW4} (a positive form, or a
  flow-compatible one) forces purity (sharpens O03, O04) &
  Prop.~\ref{prop:content}\\
O28 & Galois descent forces every isotypic divisor and adds content
  beyond zero location & Props.~\ref{prop:descent},
  \ref{prop:class-converse}\\
O29 & a real-coefficient theory is compatible with descent (sharpens
  O20) & Prop.~\ref{prop:real-coeff}\\
O30 & a cohomology for a single analytic $L$-function carries content
  beyond its Riemann hypothesis & Prop.~\ref{prop:trivial-model}\\
O31 & finite-dimensional Weil formalism with duality, rational zeta
  function, functional equation, and integral counts forces purity &
  Prop.~\ref{prop:formal-curve}\\
O32 & a nondegenerate Frobenius-compatible form suffices &
  Prop.~\ref{prop:formal-curve}\\
O33 & a nondegenerate Frobenius-compatible pairing on the algebraic
  classes of $X$ itself controls $H^j$ & Prop.~\ref{prop:formal-curve}\\
O34 & Weil's positivity is a canonical positive Frobenius polarization
  on $H^1$ & Props.~\ref{prop:weil-curves}, \ref{prop:ff-N7}\\
O35 & the suspension may use one logarithm of each eigenvalue &
  Prop.~\ref{prop:one-branch}\\
O36 & the determinant of suspended Frobenius is known only up to an
  unspecified constant & Thm.~\ref{thm:ff-det};
  Prop.~\ref{prop:ff-cut}\\
O37 & the finite-field zeta function satisfies the axioms verbatim, with
  the arithmetic unit & Prop.~\ref{prop:unit}; Thm.~\ref{thm:B}\\
O38 & Deligne's theorem makes the positivity axiom a theorem for every
  smooth projective variety over $\F_q$ &
  Thm.~\ref{thm:B}\ref{item:B-aw4}; Rem.~\ref{rem:B-status}\\
O39 & the regularized determinant over the zeros of $\zeta$ is
  $\xi/2\pi^2$, with normalization term one &
  Thm.~\ref{thm:imp-deninger-product}; Thm.~\ref{thm:det-K} (O57 for
  general $K$)\\
\end{longtable}

Every entry O24--O38 is corrected by a numbered proposition or theorem
proved here or imported with a locator (the carrier of O34 uses the
imported semisimplicity and purity statements), and O39 by an imported
theorem;
its analogue for general
number fields is O57 below. The findings
recorded in Propositions~\ref{prop:N-jordan}, \ref{prop:N-lines}, and
\ref{prop:R2-frob} are further corrected expectations of the same
kind.

The third table lists further corrected expectations, about the
multiplicative form of Weil's positivity, the parity of central orders,
and the K\"unneth conjecture of Section~\ref{sec:kunneth}.

\begin{longtable}{@{}l>{\raggedright\arraybackslash}p{0.56\textwidth}
>{\raggedright\arraybackslash}p{0.26\textwidth}@{}}
\toprule
Entry & Expectation corrected & Corrected by\\
\midrule
\endfirsthead
\toprule
Entry & Expectation corrected & Corrected by\\
\midrule
\endhead
\bottomrule
\endfoot
O40 & Weil's Rosati positivity is literally an instance of the additive
  algebra form, and the positivity axiom on the suspension of a curve
  can only be derived from the Riemann hypothesis for the curve &
  Props.~\ref{prop:aw4-EF}, \ref{prop:rosati-EF};
  Thm.~\ref{thm:B}\ref{item:B-curves}\\
O41 & self-duality and a simple pole at $s=1$ suffice for the collapse
  of a cohomology for one $L$-function (sharpens O30) &
  Prop.~\ref{prop:N-central-order}\\
O42 & a K\"unneth formula holds for fibre products over $\Spec\Z$ &
  Prop.~\ref{prop:kun-over-Z}\\
O43 & a derived K\"unneth formula over a base holds without a
  hypothesis on the map (it holds, for every cochain enhancement,
  when the derived pushforward of the
  constant sheaf lies in the thick subcategory generated by the
  constant sheaf, and fails for the double cover $E\to\mathbb{P}^1$ over
  the cochains as over the cohomology ring) &
  Props.~\ref{prop:kun-EM}, \ref{prop:EM-cochain}, \ref{prop:EM},
  \ref{prop:EM-x}; Rem.~\ref{rem:N-EM-scope}\\
O44 & the tensor product of the cohomologies of two arithmetic curves is
  a spectral core, on which \ref{ax:AW2} and AW-det make sense &
  Thm.~\ref{thm:kun-square}\\
O45 & the sums $\rho+\rho'$ of zeros form a discrete spectrum &
  Thm.~\ref{thm:kun-square}\ref{item:ks-acc} (under a hypothesis on
  the zeros)\\
O46 & Kurokawa's cone is the cohomology of the arithmetic surface &
  Prop.~\ref{prop:kun-cone}\\
O47 & compact resolvent (\ref{ax:AW6}) survives tensor products, and
  geometric compactness repairs the product &
  Props.~\ref{prop:kun-diag}, \ref{prop:kun-AW6}\\
O48 & the suspension of a product over $\F_q$ is the tensor product over
  $\C$ of the suspensions & Prop.~\ref{prop:kun-Fq-C};
  Thm.~\ref{thm:kun-Fq}\\
O49 & the absolute product of closed orbits is a well-behaved spectral
  object & Prop.~\ref{prop:kun-orbits}\\
O50 & the trace formula of the product is a Lefschetz formula over
  isolated closed orbits & Prop.~\ref{prop:kun-torus}\\
O51 & archimedean factors of tensor products are obtained by adding
  spectra & Prop.~\ref{prop:kun-arch}\\
O52 & the real-coefficient obstruction is an accident of degree one
  (sharpens O20, O29) & Rem.~\ref{rem:kun-real}\\
O53 & positivity on the arithmetic surface is a condition unrelated to
  positivity on the curve, and the bare K\"unneth conjecture carries
  content beyond the Riemann hypothesis & Thm.~\ref{thm:kun-hodge};
  Prop.~\ref{prop:kun-collapse}\\
O54 & for a given star, the Hodge index clause on the arithmetic
  surface is equivalent to the algebra form AW4-E of positivity &
  Prop.~\ref{prop:kun-hodge-strict}\\
O55 & the divisor identity for an open object decides the connecting
  maps of its localization sequence & Prop.~\ref{prop:kun-open}\\
O56 & the counting bound for a tensor sum of semibounded spectra holds
  with the shifts $c_2$ and $c_1$, or with $c_1+c_2$ for arbitrary signs
  & Prop.~\ref{prop:kun-diag}\ref{item:kd-semi}, \ref{item:kd-sharp}\\
\end{longtable}

Every entry O40--O56 except O52 is corrected by a numbered statement
proved here; O52 is carried by a remark with an imported statement, and
O45 by a statement proved under a stated hypothesis on the zeros.

The fourth table lists the corrected expectations about the
normalization of the determinant over a number field
(Section~\ref{subsec:det-normalization}) and about the trace-formula
axiom (Section~\ref{subsec:trace-axiom}).

\begin{longtable}{@{}l>{\raggedright\arraybackslash}p{0.56\textwidth}
>{\raggedright\arraybackslash}p{0.26\textwidth}@{}}
\toprule
Entry & Expectation corrected & Corrected by\\
\midrule
\endfirsthead
\toprule
Entry & Expectation corrected & Corrected by\\
\midrule
\endhead
\bottomrule
\endfoot
O57 & for every number field the regularized determinant of the zeros
  reproduces $\Lambda_K$ with normalization term one (sharpens O39) &
  Thm.~\ref{thm:det-K}\\
O58 & zeta-regularizability of $\divisor\xi_K$ follows from its
  counting function, main term, and error $O(\log T)$ &
  Prop.~\ref{prop:N-count}\\
O59 & the normalization over a number field cannot be determined because
  the half-plane splitting of the spectral zeta function needs a finer
  regularization & Thm.~\ref{thm:det-K}; Prop.~\ref{prop:Zpm}\\
O60 & the half-plane pieces $Z_\pm$ of the spectral zeta function are
  regular at $w=0$ & Prop.~\ref{prop:Zpm}\\
O61 & for the zeros of $\zeta_K$ the $O(\log T)$ fluctuation of the count
  contributes to the residues of $Z_\pm$ at $w=0$ &
  Prop.~\ref{prop:Zpm} (for general spectra it can:
  Prop.~\ref{prop:N-count})\\
O62 & the determinant of $\divisor\xi_K$ requires the general theory of
  regularized products & Thm.~\ref{thm:det-K}\\
O63 & the trace formula on a bounded window $(0,L)$ of times determines
  the degree-one spectrum & Thm.~\ref{thm:window}\\
O64 & finitely many Lefschetz identities determine the characteristic
  polynomial of Frobenius & Prop.~\ref{prop:ff-phantom}\\
O65 & a two-sided window $(-L,L)$ gives no more than the one-sided
  window $(0,L)$ & Prop.~\ref{prop:window-two}\\
O66 & a faithful trace-formula axiom must be one-sided, the two-sided
  form being stronger than Lefschetz & Thms.~\ref{thm:ff-lefschetz},
  \ref{thm:aw6-div}\\
O67 & \ref{ax:AW5} is a purely analytic condition, so that
  \ref{ax:AW5} with AW6$'$ implying \ref{ax:AW2} is the analogue of
  ``Lefschetz implies rationality'' & Prop.~\ref{prop:aw5-contains};
  Thm.~\ref{thm:aw6-hilb}\\
O68 & a class of test functions determines the spectrum if and only if
  it separates points & Thm.~\ref{thm:window};
  Prop.~\ref{prop:exp-unique}\\
O69 & the trace formula integrates to the determinant identity, as over
  $\F_q$ & Thms.~\ref{thm:aw6-div}, \ref{thm:det-K}\\
O70 & the clause of \ref{ax:AW6} that the completion loses no spectrum
  is independent of \ref{ax:AW5} & Prop.~\ref{prop:aw5-contains}\\
O71 & an exponential sum over a discrete spectrum determines its
  coefficients only if it vanishes on all positive times &
  Prop.~\ref{prop:exp-unique}\ref{item:eu-tail};
  Thm.~\ref{thm:window}\\
O72 & Schwartz functions of the spectral variable are admissible test
  functions & Prop.~\ref{prop:schwartz}\\
O73 & replacing \ref{ax:AW2} by a trace-formula axiom weakens the
  arithmetic axioms, as the Lefschetz formula makes rationality a
  theorem & Thm.~\ref{thm:aw6-div}\\
O74 & over the \'etale cochain algebras the derived K\"unneth formula
  for $E\times_{\mathbb{P}^1}E$ is repaired, the failure over the
  cohomology ring being due to the coefficient algebra &
  Props.~\ref{prop:EM-bar}, \ref{prop:EM-cochain}\\
O75 & the failure of the derived K\"unneth formula for
  $E\times_{\mathbb{P}^1}E$ is a defect of the base $\mathbb{P}^1$ &
  Props.~\ref{prop:EM-product}, \ref{prop:EM-x}\\
\end{longtable}

Every entry O57--O75 is corrected by a numbered statement proved here.
The following seven expectations arose in the comparison with BSD in
Subsection~\ref{subsec:prior-sc-bsd}. The examples concern stated
formal deductions, not elliptic curves violating BSD.

\begin{longtable}{@{}l>{\raggedright\arraybackslash}p{0.56\textwidth}
>{\raggedright\arraybackslash}p{0.26\textwidth}@{}}
\toprule
Entry & Expectation corrected & Corrected by\\
\midrule
\endfirsthead
\toprule
Entry & Expectation corrected & Corrected by\\
\midrule
\endhead
\bottomrule
\endfoot
O76 & the geometric divisor class of $E$ counts its rational points &
  Prop.~\ref{prop:bsd-picard-gap}\\
O77 & a positive height or spectral form makes the central
  Abel--Jacobi map onto & Prop.~\ref{prop:bsd-positive-gap}\\
O78 & a spectral divisor and positive metric fix the special-value
  coefficient and regulator lattice &
  Prop.~\ref{prop:bsd-coefficient-gap};
  Prop.~\ref{prop:no-positivity-cancel} for alternating cancellation\\
O79 & determinant order always equals central eigenspace dimension &
  Prop.~\ref{prop:bsd-jordan-gap}\\
O80 & duality and purity alone supply the needed positive form &
  Prop.~\ref{prop:formal-curve}\ref{item:fc-impure}--%
  \ref{item:fc-forms}; Rem.~\ref{rem:B-status}\\
O81 & finiteness of every prime-primary subgroup proves finiteness of
  the whole group & Prop.~\ref{prop:bsd-primary-gap}\\
O82 & failure of a direct implication proves logical nonimplication of
  all geometric standard conjectures and BSD &
  Rem.~\ref{rem:bsd-implication-status}\\
\end{longtable}
The statement that a trace formula determines only the alternating
divisor of a higher-dimensional object is entry O25; for the trace
formula see the paragraph after Theorem~\ref{thm:aw6-div}.

\section{Candidate theories}\label{sec:candidates}

The two constructions below test different parts of the same proposed
theory.  Flach--Morin cohomology tests its integer twists; CCM adelic
cyclic homology tests the middle zero spectrum over a number field.
Neither comparison is an additional cohomology axiom.

\begin{definition}[Comparison probes]\label{def:comparison-probes}
Let $\Xbar\in\Schbar^{\mathrm{proj}}$ be connected, let $D_X=H(\Xbar)$
for a Deninger cohomology theory $H$
(Definition~\ref{def:deninger-theory}), and let $K$ be a number field.
\begin{enumerate}[label=(\alph*)]
\item Put $R=\C[u]$.  A \emph{derived enhancement} of $D_X$ is a
  complex $D(X)$ of $R$-modules whose cohomology is $H^\bullet(X)$,
  with $u$ acting as $\Theta$.  Write $D_c(X)$ for a compact-support
  enhancement when the input includes support data.
\item For $n\in\Z$ the \emph{integer-fibre probe} is
  \[
  P_n(E)=R\operatorname{Hom}_{R}(R/(u-n),E).
  \]
  Its connecting endomorphism of degree one is denoted $\beta_n$.
\item Put $C_K=\mathbb A_K^\times/K^\times$.
  Let $N_K=H^1_{\mathrm{ad}}(K)$ be CCM's relative cyclic-homology
  representation, formed using the quotient by the closure of the
  restriction map's range (Definition~\ref{def:ccm-trace-form}).
  Its \emph{trivial norm-one sector} is
  $N_{K,1}=N_K^{C_{K,1}}$, where
  $C_{K,1}=\ker(|\cdot|:C_K\to\R_+^\times)$.
  Let $T_K$ be the infinitesimal generator of the induced
  $C_K/C_{K,1}\simeq\R_+^\times$ action $\vartheta_m$ on smooth
  vectors, parameterized by norm $e^t$.
\item If the continuous dual of $N_{K,1}$ has finite-dimensional
  generalized eigenspaces
  $E_\rho=\bigcup_{r\ge1}\ker(T_K'-\rho)^r$ with discrete support,
  where $T_K'$ is the transpose, its \emph{semisimple spectral probe}
  is
  \[
  \operatorname{Core}^{\mathrm{ss}}(N_{K,1})
  =\bigoplus_\rho E_\rho^\vee,
  \qquad \Theta|_{E_\rho^\vee}=\rho\,\mathrm{id}.
  \]
\end{enumerate}
\end{definition}

The two-term free resolution of $R/(u-n)$ gives
$P_n(E)\simeq\operatorname{Fib}(\Theta-n:E\to E)$.
When $E$ lies in degree zero, its two terms lie in degrees zero and
one.  The fibre is taken in complexes of $\C$-vector spaces: over
$\F_q$, $\Theta-n$ need not be $\mathbb 1_q$-linear.
The sector $N_{K,1}$ is not fixed by all of $C_K$.
With the Mellin convention of
\cite[Theorem~4.16 and (4.41)]{ConnesConsaniMarcolli2009Weil},
$T_K$ has the unshifted zero parameter $\rho$; the older
half-density action $W$ has parameter $\rho-1/2$.
The finite dual retains $\dim E_\rho$, while the displayed action
discards only nilpotent parts.  The finite-block and trace
properties of this extraction remain to be established for the
CCM topology.

\begin{conjecture}[Common realization of two probes]
\label{conj:fm-ccm-common}
The single assignment of Conjectures~\ref{conj:arithmetic-realization}
and~\ref{conj:arithmetic-positivity} admits a multiplicative derived
enhancement, with compact-support objects and comparison maps, such
that the following hold.
\begin{enumerate}[label=(\roman*)]
\item For each connected $\Xbar\in\Schbar^{\mathrm{proj}}$ covered by
  the hypotheses of
  \cite[Conjectures~2.3 and~4.1, and Remark~4.15]%
  {FlachMorin2020DeningerWeilArakelov}, the Flach--Morin
  compact-support complex $R\Gamma_{\mathrm{ar},c}
  (X,\widetilde\C(n))$ at each $n\in\Z$ is naturally equivalent to
  $P_n(D_c(X))$.  Likewise without supports.  The equivalences carry
  cup product with $\theta$ to $\beta_n$; over $\F_q$ this operator
  corresponds to $\log(q)\cup e_{\F_q}$, with the generator convention
  of \cite[Proposition~3.8]{FlachMorin2020DeningerWeilArakelov}.
\item For each number field $K$, the extraction
  $\operatorname{Core}^{\mathrm{ss}}(N_{K,1})$ is defined and naturally
  isomorphic, with its generator and spectral multiplicities, to
  $H^1D(\Xbar_K)$.  This comparison matches the integrated norm-action
  traces for $C_{K,1}$-invariant functions in $\mathbf S(C_K)$,
  as in \cite[Theorem~4.16]{ConnesConsaniMarcolli2009Weil}.
  It is natural under field isomorphisms.  A pairing on the extracted
  sector, still to be constructed, is carried to the pairing
  $H^1_\rho\times H^1_{1-\rho}\to\C$ of \ref{ax:AW3}.
\end{enumerate}
\end{conjecture}

This is a conjecture of these notes about comparison maps, not a
conjecture attributed in full to either source.  The first clause
refines Flach--Morin's proposed long exact sequence and derived-fibre
remark; the second asks for structure beyond CCM's proved spectral
trace formula.  It does not identify the raw cyclic module with the
derived cohomology complex or equate the CCM trace form with the star
form of \ref{ax:AW4}.  Extension maps between fields, integral
determinant lattices, and the nontrivial Hecke-character sectors require
additional constructions.

For $X=\Spec\F_q$ the unit $\mathbb 1_q$ gives a quick check.  Its
$k$th Fourier mode has eigenvalue $2\pi ik/\log q$.  Thus $P_0$ has
one copy of $\C$ in each of degrees zero and one, while $P_n$ is
acyclic for every nonzero integer $n$.  This is the finite-field
fixed-twist pattern of Theorem~\ref{thm:imp-fm20}, subject to its
finite-generation hypothesis and cup-product normalization.  CCM's
global-field construction does not apply to this zero-dimensional
point.

The two probes are incomplete in different ways.  The family of
integer fibres alone cannot recover noninteger eigenvalues
(Proposition~\ref{prop:R3}); the CCM object here reaches degree one
only.  A scalar positive trace-form completion can lose spectral
multiplicity (Proposition~\ref{prop:R9}), which is why the proposed
CCM extraction keeps the entire generalized block before
semisimplifying.  After the proposed identification, \ref{ax:AW4}
forces the target generator to be semisimple
(Theorem~\ref{thm:main-general}).  The table records the
proved portions and the remaining obligations.


Table~\ref{tab:theory-status} separates constructed cohomology theories
from conjectures about their coefficients, comparison maps, or positivity.
The final column concerns the global arithmetic realization of
Conjectures~\ref{conj:arithmetic-realization}
and~\ref{conj:arithmetic-positivity}. It records what does not follow
from the cited results; it is not a claim that no other approach could
supply the required construction. Section~\ref{sec:imported} states
the comparison theorems with their hypotheses.

{\small
\setlength{\tabcolsep}{4pt}
\begin{longtable}{@{}>{\raggedright\arraybackslash}p{0.20\textwidth}
  >{\raggedright\arraybackslash}p{0.36\textwidth}
  >{\raggedright\arraybackslash}p{0.38\textwidth}@{}}
\caption{Existing theories and their arithmetic obligations.}
\label{tab:theory-status}\\
\toprule
Framework & Constructed objects or proved results &
Remaining requirement for the arithmetic realization\\
\midrule
\endfirsthead
\toprule
Framework & Constructed objects or proved results &
Remaining requirement for the arithmetic realization\\
\midrule
\endhead
Deninger's dynamical cohomology &
The formal implication from a positive compatible star to purity;
geometric dynamical analogues under their stated hypotheses
\cite[pp.~4--5]{Deninger2010HilbertPolya} &
A global arithmetic cohomology with the full zeta spectrum,
duality, and positive star remains conjectural
\cite[Conj.~2.3]{FlachMorin2020DeningerWeilArakelov}.\\
Flach--Morin cohomology (Weil--Arakelov) &
Real complexes for proper regular arithmetic schemes are defined
\cite[\S3]{FlachMorin2020DeningerWeilArakelov} &
Their fixed-twist comparison with dynamical cohomology is
Conjecture~4.1; Theorem~4.3 proves it over finite fields under
finite generation, without supplying the positive star.\\
Ayoub's new Weil cohomology &
A theory on smooth residue-field varieties over a common
coefficient algebra $A_K$, with rigid and $\ell$-adic
specializations \cite[Thm.~1.5]{Ayoub2020NewWeil} &
The domain property and integration injectivity are
Conjectures~\ref{conj:ayoub-period-domain}
and~\ref{conj:ayoub-period-integration};
no global arithmetic flow or positive form is supplied.\\
Ayoub's foliated topology &
The Poincar\'e lemma and de Rham comparison for differential
modules \cite[Thm.~5.7.4, Cor.~5.7.7]{Ayoub2019FoliatedDGalois} &
This topology concerns differential equations. A comparison
with arithmetic dynamical cohomology is additional;
the separate conservativity proof has the author's recorded
gap \cite[\href{https://user.math.uzh.ch/ayoub/}
{conservativity manuscript notice}]{AyoubHomepage}.\\
Connes--Consani archimedean cyclic homology &
Archimedean cyclic homology realizes Serre's local gamma
factors by regularized determinants
\cite[Thm.~1.1]{ConnesConsani2014CyclicSerre} &
A global zero-spectrum identification and compatible positive
form do not follow from this local formula.\\
CCM adelic cyclic homology &
The relative $H^1$ of the cyclic-module restriction-map cokernel
realizes Hecke $L$-function zero spectra
\cite[Definition~4.10 and Proposition~4.13, pp.~21--23]%
{ConnesConsaniMarcolli2009Weil} &
A full graded theory on arithmetic schemes, and a positive star,
are not supplied; the adelic trace-form positivity is RH-equivalent
\cite[Corollary~6.3, p.~29]{ConnesConsaniMarcolli2009Weil}.\\
Arithmetic and scaling sites &
Geometric sites and Riemann--Roch on periodic orbits
\cite[Thms.~3.9, 5.17]{ConnesConsani2017ScalingGeometry};
Riemann--Roch for arithmetic divisors
\cite[Thm.~1.1]{ConnesConsani2023RiemannRochSpecZ} &
The required positivity on the square remains a target of
the Riemann--Roch strategy
\cite[\S3]{ConnesConsani2018RiemannRoch};
the one-dimensional formulas do not prove it.\\
Topological periodic cyclic homology (TP) &
Hesselholt's finite-field determinant formula for the full
even and odd sums \cite[Thm.~A]{Hesselholt2016THHZeta} &
A global arithmetic object, comparison across places, and
positive global form are additional requirements.\\
Prismatic and $A_{\inf}$ cohomology &
Local integral comparison theorems
\cite[Thm.~1.8]{BhattMorrowScholze2018Integral},
\cite[Thm.~1.8]{BhattScholze2022Prisms} &
These comparisons do not construct a complex global flow
whose spectrum is the zeta divisor.\\
Kottwitz and Berkovich settings &
The Kottwitz-valued Weil theory is conjectural
\cite[Conj.~9.5]{Scholze2018PadicGeometry}; Berkovich motives
have pointwise comparison theorems
\cite[Prop.~1.5]{Scholze2026BerkovichMotives} &
Neither statement supplies the global zero spectrum or its
positive completion.\\
\bottomrule
\end{longtable}
}

The distinction is between these existing constructions and the
additional arithmetic realization and positivity assertions, not
between having or lacking cohomology in every sense. In particular,
Deninger's real-coefficient obstruction concerns the specified
functorial Galois multiplicities
\cite[Thm.~2.13]{Deninger2022NoWeil}; it does not contradict the
existence of the real Weil--Arakelov complexes in the table.

Several pieces of the dynamical cohomology $H^\bullet_{\mathrm{dyn}}$ of
\cite[Conj.~2.3]{FlachMorin2020DeningerWeilArakelov} already have
rigorous, unconditional constructions; Table~\ref{tab:dyn-realizations}
lists them. None supplies the positivity that would prove RH. In each
case where a positive form is available, its existence is equivalent
to RH.

{\small
\setlength{\tabcolsep}{4pt}
\begin{longtable}{@{}>{\raggedright\arraybackslash}p{0.20\textwidth}
  >{\raggedright\arraybackslash}p{0.40\textwidth}
  >{\raggedright\arraybackslash}p{0.34\textwidth}@{}}
\caption{Rigorous partial realizations of dynamical cohomology.}
\label{tab:dyn-realizations}\\
\toprule
Piece of $H^\bullet_{\mathrm{dyn}}$ & Rigorous construction &
What it does not supply\\
\midrule
\endfirsthead
\toprule
Piece of $H^\bullet_{\mathrm{dyn}}$ & Rigorous construction &
What it does not supply\\
\midrule
\endhead
Smooth projective varieties over $\F_q$ &
Deninger's spaces $D(H^i(X_{\overline{\F}_q},\Q_\ell)\otimes\C)$ with
their derivation $\Theta$ \cite{Deninger1992Local},
\cite[\S4.2, (7)]{FlachMorin2020DeningerWeilArakelov}; when Frobenius
acts semisimply they satisfy
\cite[Conj.~2.3]{FlachMorin2020DeningerWeilArakelov} except the
$*$-operator~(iv) &
The positive $*$-operator; positivity here is the Riemann hypothesis
over $\F_q$, which is imported, not produced.\\
Degree one for $\overline{\Spec\Z}$ &
The relative $H^1$ of the adelic cyclic module realizes Hecke
$L$-function zero spectra
\cite[Def.~4.10, Prop.~4.13]{ConnesConsaniMarcolli2009Weil}; an
operator on a nuclear Fr\'echet space has the zeros of $\xi$ as
eigenvalues, with algebraic multiplicity equal to the zero order
\cite[Thm.~4.1, Cor.~4.2]{Meyer2005SpectralZeta} &
Products, duality, a star, and the other degrees. Positivity of the
adelic trace form is equivalent to RH
\cite[Cor.~6.3]{ConnesConsaniMarcolli2009Weil}. A zero of order $m>1$
gives a Jordan block of length~$m$, which a positive equivariant form
excludes, so positivity would also force simple zeros.\\
The parts $\Theta=n$, $n\in\Z$ &
Weil--Arakelov complexes, defined unconditionally from motivic and
Deligne cohomology
\cite[\S3]{FlachMorin2020DeningerWeilArakelov}; their comparison with
$H^\bullet_{\mathrm{dyn}}$ is
\cite[Conj.~4.1]{FlachMorin2020DeningerWeilArakelov}, proved over
$\F_q$ under finite generation
\cite[Thm.~4.3]{FlachMorin2020DeningerWeilArakelov} &
Noninteger eigenvalues, and the comparison over $\Z$.\\
Archimedean factor &
Archimedean cyclic homology realizes Serre's local gamma factors by
regularized determinants \cite[Thm.~1.1]{ConnesConsani2014CyclicSerre} &
The nontrivial zeros.\\
The full theory on $\overline{\Spec\Z}$ &
None: cohomology in every degree with cup product, duality and a
positive $*$-operator is conjectural
\cite[Conj.~2.3]{FlachMorin2020DeningerWeilArakelov} &
Everything listed in the other rows at once.\\
\bottomrule
\end{longtable}
}

The tautological constructions show the same circularity. The trivial
model of Definition~\ref{def:trivial-model} exists for every divisor,
but a positive datum with that spectrum exists exactly when the divisor
lies on $\Re s=\tfrac12$ (Proposition~\ref{prop:trivial-model},
\ref{item:tm-iff}); the Lean statements recorded in
Subsection~\ref{sec:lean} check this equivalence. A rigorous candidate
in the useful sense would be a full theory constructed without
reference to the zeros, with positivity proved rather than imposed.

Scholze's real-exponent analytic ring offers another possible input.
For $0<r<1$, let $R_{>r}=\Z((T^\R))_{>r}$ be the union over
$r<r'<1$
of series $\sum_xa_xT^x$ with $\sum_x|a_x|(r')^x<\infty$.
Lecture~XIV of \cite[Theorem~14.11]{Scholze2026AnalyticGeometry}
equips it with an analytic ring structure and the semigroup action
$\sigma_t(T^x)=T^{tx}$ for $t\ge1$. At primes these are Frobenius
lifts. A fixed-radius action need not be invertible.

\begin{proposition}\label{prop:scholze-fixed-radius}
On $R_{>1/4}$, the endomorphism $\sigma_2$ is not surjective.
\end{proposition}

\begin{proof}
The series $b=\sum_{n\ge0}3^nT^n$ belongs to $R_{>1/4}$: its
weighted sum converges whenever $1/4<r'<1/3$. A preimage under
$\sigma_2$ must be $\sum_{n\ge0}3^nT^{n/2}$, whose weighted sum
$\sum_{n\ge0}(3\sqrt{r'})^n$ diverges for every $r'>1/4$.
Coefficient uniqueness rules out another preimage.
\end{proof}

The source also gives canonical isomorphisms between rings at
different radii. The proposition only blocks treating this literal
fixed-radius semigroup as the two-sided flow required above; a
radius-varying or enlarged construction needs its own comparison,
generator, and trace.

No composition of the comparison arrows of
Theorems~\ref{thm:imp-h17}, \ref{thm:imp-bms1}, \ref{thm:imp-bms2},
\ref{thm:imp-bs}, \ref{thm:imp-cc14}, \ref{thm:imp-a20},
\ref{thm:imp-a19}, and~\ref{thm:imp-sb26} produces the object of
Theorem~\ref{thm:rh}: their domains do not match the proposed
arithmetic spectral source without new constructions
(Remark~\ref{rem:N-composition}; cf.\ Proposition~\ref{prop:R7}).

\section{Imported results}\label{sec:imported}

This section states, with their hypotheses and exact locators, the
results from the literature on which Sections~\ref{sec:negative},
\ref{sec:candidates}, and~\ref{sec:debranges} rely. They are not proved
here; the standing convention of these notes is that advanced
comparison theorems are stated precisely and not re-proved. Page numbers
refer to the version cited in the bibliography. None of these results
constructs the object of Theorem~\ref{thm:rh}.

\subsection{Coefficients and leafwise cohomology}

Let $k$ be a number field. Deninger lists properties 2.1--2.12 expected
of a cohomology $H^\bullet(\Spec\overline{\cO}_k,\C)$ with an
endomorphism $\theta$ \cite[pp.~2--5]{Deninger2022NoWeil}; property 2.12
asks that, for $k/\Q$ Galois with group $G$, each irreducible
representation of $G$ occur in the central generalized eigenspace of
$\theta$ on $H^1$ with multiplicity equal to the order of its Artin
$L$-function at $s=\tfrac12$.

\begin{theorem}[{\cite[Theorem~2.13, p.~5]{Deninger2022NoWeil}}]
\label{thm:imp-deninger-nogo}
Let $k$ be a number field. There is no functorial real-valued
cohomology theory, written
$H^\bullet(\Spec\overline{\cO}_k,\R)$, with an endomorphism
$\theta_\R$ such that
$H^\bullet(\Spec\overline{\cO}_k,\R)\otimes_\R\C$ satisfies property
2.12 of \cite{Deninger2022NoWeil}.
\end{theorem}

The source then observes \cite[pp.~5--6]{Deninger2022NoWeil} that
complex leafwise cohomology always has the real structure given by the
cohomology of real forms, so that ``there has to be a twisted version of
leafwise cohomology with complex coefficients that does not have a
natural real structure''.

Reduced leafwise cohomology $\overline H^j_{\mathcal F}$ and leafwise
harmonic forms are as in Definition~\ref{def:reduced-leafwise}.

\begin{theorem}[{\cite[Proposition~1.2 and~(12), p.~214]%
{DeningerSinghof2001Counterexample}}]\label{thm:imp-ds}
There is a one-dimensional foliation $\mathcal F$ with dense leaves on
a three-dimensional Heisenberg nilmanifold, with a leafwise metric,
such that the reduced smooth leafwise cohomology
$\overline H^1_{\mathcal F}$ is infinite-dimensional while the space of
smooth leafwise harmonic $1$-forms is one-dimensional.
\end{theorem}

In the same example the passage from the dynamical side of a trace
formula to reduced leafwise cohomology is not allowed
\cite[\S2, pp.~215--217]{DeningerSinghof2001Counterexample}. The
positive theorem quoted there as Theorem~0.2 \cite[p.~210]%
{DeningerSinghof2001Counterexample}, due to \'Alvarez L\'opez and
Kordyukov, assumes a Riemannian foliation with a bundle-like metric.

\subsection{Finite fields and \texorpdfstring{$p$}{p}-adic comparisons}

Deligne's theorem is stated in Section~\ref{sec:classical}
(Theorem~\ref{thm:imp-deligne}).

\begin{theorem}[{\cite[Theorem~A, p.~2]{Hesselholt2016THHZeta}}]
\label{thm:imp-h17}
Let $k=\F_q$, $W=W(k)$, $\iota:W\to\C$ a ring homomorphism, and
$f:X\to\Spec k$ smooth and proper. Let $\Theta$ be the graded
derivation of periodic topological cyclic homology
$\mathrm{TP}_*(X)\otimes_{W,\iota}\C$ with $q^\Theta=\mathrm{Fr}_q^*$
normalized as in the source. Then
\[
\zeta(X,s)=\frac{\detreg\bigl(s\cdot\mathrm{id}-\Theta\mid
  \mathrm{TP}_{\mathrm{odd}}(X)\otimes_{W,\iota}\C\bigr)}
 {\detreg\bigl(s\cdot\mathrm{id}-\Theta\mid
  \mathrm{TP}_{\mathrm{even}}(X)\otimes_{W,\iota}\C\bigr)}.
\]
\end{theorem}

It is natural to expect Theorem~\ref{thm:imp-h17} to identify each
cohomological degree with a weight-graded piece carrying a metric; in
fact it is an identity of regularized determinants over the full even
and odd graded sums.

\begin{remark}\label{rem:N-h17}
Theorem~\ref{thm:imp-h17} is proved by comparing spectral sequences
\cite[pp.~2--3]{Hesselholt2016THHZeta}; it identifies regularized
determinants over the full even and odd graded sums, not individual
weight-graded pieces with metrics, and the choice $\iota:W\to\C$ is
algebraic, not a continuous embedding supplying a positive metric. It
concerns one finite field and says nothing about $\Spec\Z$.
\end{remark}

Hesselholt's formula uses $\mathrm{TP}$, the topological analogue
of periodic cyclic homology formed as the circle Tate construction
on $\mathrm{THH}$, rather than $\mathrm{TC}$
\cite[Theorem~A and p.~1]{Hesselholt2016THHZeta}. The
$\mathrm{TC}^-$ and $\mathrm{TP}$ filtrations of
Theorem~\ref{thm:imp-bms2} provide $p$-adic comparison data.
Connes--Consani's archimedean determinant instead uses ordinary
cyclic homology \cite[Theorem~1.1]{ConnesConsani2014CyclicSerre}.
These parallel determinant constructions do not by themselves
identify either theory with CCM's adelic cyclic quotient.

\begin{theorem}[{\cite[Theorem~1.8, pp.~5--6]%
{BhattMorrowScholze2018Integral}}]\label{thm:imp-bms1}
Let $C/\Q_p$ be complete and algebraically closed with ring of integers
$\cO_C$ and residue field $k$, let $\mathfrak X/\cO_C$ be a proper smooth
formal scheme, $A_{\inf}=W(\cO_C^\flat)$, and $\mu=[\epsilon]-1$ for a
compatible system of $p$-power roots of unity. There is a perfect
$A_{\inf}$-complex $R\Gamma_{A_{\inf}}(\mathfrak X)$ with a Frobenius and
canonical quasi-isomorphisms
$R\Gamma_{A_{\inf}}(\mathfrak X)\otimes^L_{A_{\inf}}W(k)\simeq
R\Gamma_{\mathrm{crys}}(\mathfrak X_k/W(k))$,
$R\Gamma_{A_{\inf}}(\mathfrak X)\otimes^L_{A_{\inf},\theta}\cO_C\simeq
R\Gamma_{\mathrm{dR}}(\mathfrak X/\cO_C)$, and
$R\Gamma_{A_{\inf}}(\mathfrak X)\otimes_{A_{\inf}}A_{\inf}[1/\mu]\simeq
R\Gamma_{\text{\'et}}(\mathfrak X_C,\Z_p)\otimes_{\Z_p}A_{\inf}[1/\mu]$.
\end{theorem}

\begin{theorem}[{\cite[Theorems~1.8, 1.10, and~1.12(4), pp.~5--7]%
{BhattMorrowScholze2019THH}}]\label{thm:imp-bms2}
If $A$ is the $p$-adic completion of a smooth $\cO_C$-algebra, there is
a Frobenius-equivariant identification
\[
A\Omega_A\simeq R\Gamma_{\mathrm{syn}}\bigl(A,
\pi_0\mathrm{TC}^-(-;\Z_p)\bigr);
\]
if $A$ is smooth over a perfect field $k$ of characteristic $p$,
the right-hand side computes the crystalline cohomology
$R\Gamma_{\mathrm{crys}}(A/W(k))$. If $A$ is
quasisyntomic, $\mathrm{TC}^-(A;\Z_p)$ and $\mathrm{TP}(A;\Z_p)$ carry
complete exhaustive multiplicative filtrations with
\[
\operatorname{gr}^n\mathrm{TC}^-(A;\Z_p)\simeq
N^{\ge n}\widehat\Delta_A\{n\}[2n],\qquad
\operatorname{gr}^n\mathrm{TP}(A;\Z_p)\simeq\widehat\Delta_A\{n\}[2n].
\]
\end{theorem}

\begin{theorem}[{\cite[Theorem~1.8, p.~4]{BhattScholze2022Prisms}}]
\label{thm:imp-bs}
Let $(A,I)$ be a bounded prism and $\mathfrak X$ a smooth $p$-adic formal
scheme over $A/I$, with prismatic complex
$R\Gamma_\Delta(\mathfrak X/A)$. There are canonical comparisons with
crystalline cohomology when $I=(p)$, with de Rham cohomology of
$\mathfrak X/(A/I)$ after Frobenius base change, with Hodge--Tate
cohomology for affine $\mathfrak X$, and, for perfect $A$, with \'etale
cohomology:
$R\Gamma_{\text{\'et}}(\mathfrak X_\eta,\Z/p^n)\simeq
\bigl(R\Gamma_\Delta(\mathfrak X/A)/p^n[1/I]\bigr)^{\varphi=1}$, where
the fixed points are derived.
\end{theorem}

The Breuil--Kisin and $A_{\inf}$ cases are specializations
\cite[Example~1.9, p.~5]{BhattScholze2022Prisms}. The tensor products
and fixed points in Theorems~\ref{thm:imp-bms1}--\ref{thm:imp-bs} are
derived, and an associated-graded comparison is not a splitting.

\begin{theorem}[{\cite[Theorem~1.1 and~(2), p.~2; (4), p.~3]%
{ConnesConsani2014CyclicSerre}}]\label{thm:imp-cc14}
Let $K$ be a number field, $X/K$ smooth projective of dimension $d$, and
$\nu$ an archimedean place. With the archimedean cyclic homology
$HC^{\mathrm{ar}}_*(X_\nu)$ and the operator $\Theta$ defined in the
source through the $\lambda$-operations, one has, for real $s$ away from
the singularities,
\[
\prod_{w=0}^{2d}L_\nu(H^w(X),s)^{(-1)^{w+1}}
=\frac{\detreg\bigl(\tfrac1{2\pi}(s-\Theta)\mid
   HC^{\mathrm{ar}}_{\mathrm{even}}(X_\nu)\bigr)}
 {\detreg\bigl(\tfrac1{2\pi}(s-\Theta)\mid
   HC^{\mathrm{ar}}_{\mathrm{odd}}(X_\nu)\bigr)}.
\]
\end{theorem}

The orientation (even over odd) is the one printed and is opposite to
that of Theorem~\ref{thm:imp-h17}. At a complex place,
$HC^{\mathrm{ar}}$ is the cokernel in an exact sequence
$0\to HP^{\mathrm{real},\lambda}_{*+2}(X_\nu)\to HC_*(X_\nu)\to
HC^{\mathrm{ar}}_*(X_\nu)\to0$ \cite[(4), p.~3]%
{ConnesConsani2014CyclicSerre}. The formula realizes local gamma factors,
not global zeros.

\begin{theorem}[{\cite[Construction~1.3, Th\'eor\`eme~1.5, pp.~1749--%
1750]{Ayoub2020NewWeil}}]\label{thm:imp-a20}
Let $K$ be a valued field of rank one and mixed characteristic $(0,p)$
with residue field $k$ and completion $\widehat K$. There is a Weil
cohomology $\Gamma_{\mathrm{new}}(-/K)$ on smooth $k$-schemes with
coefficients in a $K$-algebra $A_K$, with
$\Gamma_{\mathrm{new}}\otimes^L_{A_K}\widehat K\simeq\Gamma_{\mathrm{rig}}$
and, after choosing $K\hookrightarrow\C$ and $\overline\Q_\ell\simeq\C$
for $\ell\ne p$,
$\Gamma_{\mathrm{new}}\otimes^L_{A_K}\overline\Q_\ell\simeq
\Gamma_\ell\otimes_{\Q_\ell}\overline\Q_\ell$.
\end{theorem}

For a rank-one valued field $K$ of mixed characteristic $(0,p)$,
write $A_K$ for Ayoub's algebra of abstract $p$-adic periods and
$\widehat K$ for the completion of $K$. Write
$I_K:A_K\to\widehat K$ for the canonical $K$-algebra homomorphism
defined by integration in
\cite[Construction~1.3, (1-1)--(1-2), pp.~1749--1750]
{Ayoub2020NewWeil}.

\begin{conjecture}[{Ayoub;
\cite[Remarque~1.6, p.~1750]{Ayoub2020NewWeil}}]
\label{conj:ayoub-period-domain}
Let $K$ be a rank-one valued field of mixed characteristic $(0,p)$
and let $A_K$ be its algebra of abstract $p$-adic periods.
Then $A_K$ is an integral domain.
\end{conjecture}

\begin{conjecture}[{Ayoub;
\cite[Remarque~1.6, p.~1750]{Ayoub2020NewWeil}}]
\label{conj:ayoub-period-integration}
Let $K$ be a number field equipped with a rank-one nonarchimedean
valuation of residue characteristic $p$, let $\widehat K$ be its
completion, and let $A_K$ be its algebra of abstract $p$-adic
periods. Then the integration homomorphism
$I_K:A_K\to\widehat K$ is injective.
\end{conjecture}


\begin{theorem}[{\cite[Th\'eor\`eme~5.7.4 and Corollaire~5.7.7,
pp.~187--188]{Ayoub2019FoliatedDGalois}}]\label{thm:imp-a19}
Let $(K,\Delta)$ be a differential field of characteristic zero.
\begin{enumerate}[label=(\roman*)]
\item (Th\'eor\`eme~5.7.4.) For an ind-holonomic differential module
  $M$, the inclusion of horizontal sections into the de~Rham complex of
  $M$ is an $\mathrm{fttf}$-local equivalence.
\item (Corollaire~5.7.7.) For a bounded-below complex $M^\bullet$ of
  ind-holonomic differential modules,
  $H^i_{\mathrm{dR}}(K;M^\bullet)\simeq H^i_{\mathrm{fttf}}(K;
  (M^\bullet)^\delta)$.
\end{enumerate}
\end{theorem}

The source is an author manuscript; its ``foliated'' topology concerns
algebraic differential equations, not a foliation on an arithmetic
dynamical space.

\begin{conjecture}[{\cite[Conjecture~9.5, pp.~24--25]%
{Scholze2018PadicGeometry}}]\label{conj:imp-s18}
There is a Weil cohomology theory for varieties over $\overline{\F}_p$
with values in Kottwitz's category $\mathrm{Kt}_\Q$ whose images in
$\mathrm{Kt}_{\Q_p}$, $\mathrm{Kt}_{\Q_\ell}$, and $\mathrm{Kt}_\R$ are
crystalline cohomology, $\ell$-adic cohomology, and a Weil cohomology
with values in $\mathrm{Kt}_\R$, respectively.
\end{conjecture}

The source explains that the conjecture follows from the standard
conjectures, the Tate conjecture, and the Hodge conjecture for CM abelian
varieties; it is not an unconditional construction, and its base is
$\overline{\F}_p$, not $\Spec\Z$.

\begin{theorem}[{\cite[Proposition~1.5, p.~4]%
{Scholze2026BerkovichMotives}}]\label{thm:imp-sb26}
Let $C$ be an algebraically closed nondiscrete Banach field and
$X=M_{\mathrm{arc}}(C)$. If $C=\C$, then $D(\Z)\simeq
D^{\mathrm{eff}}_{\mathrm{mot}}(X)$; if $C$ is nonarchimedean and $n$
is invertible in its residue field, then $D(\Z/n)\simeq
D^{\mathrm{eff}}_{\mathrm{mot}}(X,\Z/n)$; if the residue characteristic
is $p$, then $p$ is invertible in $D^{\mathrm{eff}}_{\mathrm{mot}}(X)$;
and in mixed characteristic, tilting gives
$D^{\mathrm{eff}}_{\mathrm{mot}}(C)\simeq
D^{\mathrm{eff}}_{\mathrm{mot}}(C^\flat)$.
\end{theorem}

\subsection{Flach--Morin cohomology (Weil--Arakelov)}

\begin{conjecture}[{\cite[Conjecture~4.1, p.~527]%
{FlachMorin2020DeningerWeilArakelov}}]\label{conj:imp-fm20}
Let $X$ be as in the source and $n\in\Z$. There is a long exact
sequence
\[
\cdots\to H^i_{\mathrm{ar},c}(X,\widetilde\C(n))\to
H^i_{\mathrm{dyn},c}(X)\xrightarrow{\Theta-n}H^i_{\mathrm{dyn},c}(X)\to
H^{i+1}_{\mathrm{ar},c}(X,\widetilde\C(n))\to\cdots
\]
whose composite $H^i_{\mathrm{ar},c}\to H^{i+1}_{\mathrm{ar},c}$ is cup
product with the class $\theta$ of the source, and likewise without
supports.
\end{conjecture}

\begin{theorem}[{\cite[Theorem~4.3, p.~527]%
{FlachMorin2020DeningerWeilArakelov}}]\label{thm:imp-fm20}
Let $X$ be a smooth projective connected variety over $\F_q$ and
$n\in\Z$. If $H^i(X_W,\Z(n))$ is finitely generated for all $i\in\Z$,
then Conjecture~\ref{conj:imp-fm20} holds for $X$ and $n$, with the
dynamical cohomology constructed from $\ell$-adic cohomology.
\end{theorem}

It is natural to read Theorem~\ref{thm:imp-fm20} as a finite-field
verification of the whole dynamical package; in fact it is the
fixed-twist long exact sequence only.

\begin{remark}\label{rem:N-fm20}
Theorem~\ref{thm:imp-fm20} gives the fixed-twist long exact sequence of
Conjecture~\ref{conj:imp-fm20} only. It does not supply the positive
star, and the comparison of cup-product operators carries a
normalization:
\cite[Proposition~3.8, p.~526]{FlachMorin2020DeningerWeilArakelov}
relates the Weil--Arakelov operator to $\log(q)\cup e_{\F_q}$. A derived
enhancement $R\Gamma_{\mathrm{ar},c}\simeq\operatorname{fib}(\Theta-n)$
is a further statement. Deninger's determinant formula
\cite[Conjecture~2, p.~3]{Deninger2010HilbertPolya} and the broader
conjecture \cite[Conjecture~2.3, p.~522]%
{FlachMorin2020DeningerWeilArakelov} remain conjectures.
\end{remark}

\subsection{CCM adelic cyclic homology}

\begin{definition}[Adelic trace form]\label{def:ccm-trace-form}
Let $K$ be a global field and $C_K=\mathbb A_K^\times/K^\times$.
Let $\mathbf S(C_K)$ be the convolution test algebra of
\cite[(4.20), p.~19]{ConnesConsaniMarcolli2009Weil}: for a number
field it is $\bigcap_{\beta\in\R}|\cdot|^\beta\mathcal S(C_K)$,
and for a function field it consists of compactly supported
Schwartz functions. Let $H^1_{\mathrm{ad}}(K)$ denote the source's
$H^1(\mathbb A_K/K^\times,C_K)$, obtained by applying cyclic
homology to the restriction-map cokernel. The cokernel is taken
modulo the closure of the range
\cite[Definition~4.10 and~(4.28), pp.~21--22]%
{ConnesConsaniMarcolli2009Weil}.
Write $\vartheta_m$ for its multiplicative action. For
$f,g\in\mathbf S(C_K)$ set
\[
\langle f,g\rangle_{\mathrm{ad}}=
\operatorname{Tr}\bigl(\vartheta_m(f\star g)
\mid H^1_{\mathrm{ad}}(K)\bigr),\qquad
f^\sharp(x)=|x|^{-1}\overline{f(x^{-1})},
\]
and $Q_K(f)=\langle f,f^\sharp\rangle_{\mathrm{ad}}$.
\end{definition}

\begin{theorem}[{Connes--Consani--Marcolli
\cite[Corollary~6.3, p.~29]{ConnesConsaniMarcolli2009Weil}}]
\label{thm:imp-ccm-positivity}
Let $K$ be a global field and let $Q_K$ be as in
Definition~\ref{def:ccm-trace-form}. Then $Q_K(f)\ge0$ for every
$f\in\mathbf S(C_K)$ if and only if every $L$-function with
Gr\"ossencharakter on $K$ satisfies the Riemann hypothesis.
\end{theorem}

\begin{proof}
Put $f^*(x)=\overline{f(x^{-1})}$ and
$\Delta^z f(x)=|x|^z f(x)$. The Weil positivity criterion, in the
normalization of the source, is
\(\operatorname{Tr}(W(f\star f^*))\ge0\) for every test function,
where $W(f)=\vartheta_m(\Delta^{-1/2}f)$
\cite[Proposition~6.2, p.~29]{ConnesConsaniMarcolli2009Weil}.
For $h=\Delta^{-1/2}f$, one has
$\Delta^{-1/2}f^*=\Delta^{-1}h^*=h^\sharp$; hence that trace is
$Q_K(h)$. Since $\Delta^{-1/2}$ is an automorphism of the test
algebra, the two positivity statements are equivalent.
\end{proof}

\begin{theorem}[Adelic trace positivity implies RH]
\label{thm:ccm-positivity-rh}
Let $Q_{\Q}$ be the trace form on $\mathbf S(C_{\Q})$ in
Definition~\ref{def:ccm-trace-form}. If $Q_{\Q}(f)\ge0$ for every
$f\in\mathbf S(C_{\Q})$, then every nontrivial zero of $\zeta(s)$
has real part $1/2$.
\end{theorem}

\begin{proof}
Apply Theorem~\ref{thm:imp-ccm-positivity} with $K=\Q$.
Its hypothesis is the stated nonnegativity for the full test algebra,
so every Gr\"ossencharakter $L$-function of $\Q$ satisfies RH.
The trivial Gr\"ossencharakter has $L(s,1)=\zeta(s)$, giving the
claim.
\end{proof}

The hypothesis concerns every test function and every Hecke-character
sector. Classical RH for $\zeta$ alone is not asserted to imply it;
the full equivalence is Theorem~\ref{thm:imp-ccm-positivity}.
This is a criterion, not an established positive form: the source
proves that the trace pairing vanishes on the range of its restriction
map \cite[Proposition~6.4, pp.~29--30]%
{ConnesConsaniMarcolli2009Weil}, but does not prove $Q_K\ge0$ on all
test functions. Its form is semidefinite if the criterion holds and
has a radical; identifying it with the positive polarization
\eqref{eq:inf-pol} would require an additional comparison of spaces
and scaling operators. The source also leaves the analogue of
principal divisors open \cite[p.~4]{ConnesConsaniMarcolli2009Weil}.

\begin{remark}[Adelic realization and the full axioms]
\label{rem:ccm-deninger-scope}
Let $K$ be a global field. The $C_K$-action on
$H^1_{\mathrm{ad}}(K)$ realizes the zero spectra of all
Gr\"ossencharakter $L$-functions, and its integrated action has the
trace formula of \cite[Proposition~4.13 and Theorem~4.16, pp.~23--24]%
{ConnesConsaniMarcolli2009Weil}. We call this an \emph{adelic
spectral realization}. For a number field $K$, its Hecke-character
decomposition \cite[Proposition~4.12, p.~23]%
{ConnesConsaniMarcolli2009Weil} contains more sectors than the single
function $Z_{\Xbar_K}=\Lambda_K$ of
Definition~\ref{def:compactified-schemes}. Even after
restricting to the trivial-character sector, the cited construction
does not give the cup product, endpoints, duality, and positive star
required by Definition~\ref{def:weil-cohomology}. In particular,
Theorem~\ref{thm:imp-ccm-positivity} makes positivity of its trace
form an RH-equivalent condition, not an established property. Thus
the adelic realization is a concrete spectral candidate within the
broader program, without being an instance of the full definition.
\end{remark}

\begin{remark}[The sense in which CCM has cohomology]
\label{rem:ccm-relative-cohomology}
Let $K$ be a global field. CCM define a cokernel in cyclic modules
and then apply cyclic homology to obtain
$H^1(\mathbb A_K/K^\times,C_K)$. They explicitly interpret it as a
first relative cohomology group for the inclusion of $C_K$ in the
adelic class space \cite[Definition~4.10 and (4.28), pp.~21--22]%
{ConnesConsaniMarcolli2009Weil}. We call this the \emph{CCM adelic
cyclic homology construction}. The cited paper does not construct
the graded functor with products, endpoints, and duality of
Definition~\ref{def:deninger-theory}, or a polarization; hence this
relative $H^1$
alone is not a Deninger cohomology theory in the sense of these notes.
\end{remark}

\begin{theorem}[{\cite[Theorems~3.9, p.~14, and~5.17, p.~32]%
{ConnesConsani2017ScalingGeometry}}]\label{thm:imp-cc16}
The points of the scaling topos $[0,\infty)\rtimes\N^\times$ are
canonically identified with $\mathcal A(\R_+^{\max})$. For a prime $p$
and a divisor $D$ on the periodic orbit $C_p$, with the continuous
dimension $\operatorname{Dim}_\R$ of the source,
$\operatorname{Dim}_\R H^0(D)-\operatorname{Dim}_\R H^0(-D)=\deg D$.
\end{theorem}

This Riemann--Roch equality concerns divisors on a periodic orbit; it is
not an intersection-positivity theorem on an arithmetic surface.

\begin{theorem}[{\cite[Proposition~4.8, p.~24]%
{ConnesConsaniMoscovici2024Prolate}}]\label{thm:imp-ccm23}
Let $S\ni\infty$ be a finite set of places of $\Q$ and $\lambda>0$.
The map $\upsilon_S$ of the source is an isomorphism of hilbertian
spaces, in the source's terminology (a linear homeomorphism of
topological vector spaces whose topology, but not a preferred inner
product, comes from a Hilbert-space structure), from the semilocal Sonin space
$\mathcal S_\lambda(X_S,\alpha)$ onto the space $\mathcal B_\lambda$,
compatible with the semilocal amplification map.
\end{theorem}

Giving $\mathcal B_\lambda$ the $S$-dependent induced norm defines the
de~Branges space $\mathcal B^S_\lambda$ of \cite[\S4.8.3]%
{ConnesConsaniMoscovici2024Prolate}. An isomorphism of abstract
hilbertian spaces does not identify norms, embeddings, or operators on
entire functions (Proposition~\ref{prop:R10}).

\subsection{De Branges side}

Let $E=E_F=F+iF'$ (Definition~\ref{def:hb}). Since
$F'(z)=-i\xi'(\tfrac12-iz)$, one has
$E(z)=\xi(\tfrac12-iz)+\xi'(\tfrac12-iz)$, which is the function
denoted $E_\xi$ in \cite[(1.3), p.~2]{Suzuki2023WeilHilbert}. Let
$\vartheta=E^\#/E$ (denoted $\Theta_\xi$ there), let
$\mathcal K(\vartheta)=H^2(\Cp)\ominus\vartheta H^2(\Cp)$ be the model
space, and let $\mathcal H(E)$ be the de Branges space of
Definition~\ref{def:conrey-li}. When $Q_W$
(Definition~\ref{def:weil-form}) is positive definite, let
$\mathcal H_W$ be the Hilbert-space completion of $C_c^\infty(\R)$ for
the norm $\|\psi\|_W=Q_W(\psi)^{1/2}$, also written $\|\cdot\|_W$ on
$\mathcal H_W$, and $[\psi]\in\mathcal H_W$ the class of $\psi$. Let
$\mathsf F$ be the Fourier transform of Definition~\ref{def:weil-form},
extended to $L^2(\R)$, and for $G\in L^2(\R)$ put
$\psi_G=\mathsf F^{-1}G$.

\begin{theorem}[{\cite[Theorem~1.1, p.~2]{Suzuki2023WeilHilbert}}]
\label{thm:imp-sz23}
Assume RH. Then $\mathcal H_W$ is defined, and the map
$G\mapsto[\psi_G]$ of the source is an isomorphism of Hilbert spaces
$\mathcal K(\vartheta)\to\mathcal H_W$, with
$\|EG\|^2_{\mathcal H(E)}=\|G\|^2_{\mathcal K(\vartheta)}
=\pi\|[\psi_G]\|_W^2$ for $G\in\mathcal K(\vartheta)$.
\end{theorem}

In the source, every class in $\mathcal H_W$ has a representative in
$L^2(\R)$ \cite[Theorem~5.5]{Suzuki2023WeilHilbert}, $[\psi]$ denotes
the class represented by $\psi\in L^2(\R)$, and the source's Weil form
$\langle\psi,\psi\rangle_W$ coincides with $Q_W(\psi)$ for
$\psi\in C_c^\infty(\R)$.

The self-adjoint operator of \cite[\S6, p.~18]{Suzuki2023WeilHilbert}
has each distinct zero as an eigenvalue of multiplicity one, whatever its
order as a zero. The source's RH-equivalent norm criterion
\cite[Theorem~1.4, p.~4]{Suzuki2023WeilHilbert} is not imported: with
the printed definitions its norm expression degenerates
(Proposition~\ref{prop:R17}).

For $\omega>0$ and $s=\tfrac12-iz$, the source
\cite[(1.1), p.~1]{Suzuki2012Canonical} puts
$A^\omega(z)=\tfrac12(\xi(s+\omega)+\xi(s-\omega))$ and
$B^\omega(z)=\tfrac i2(\xi(s+\omega)-\xi(s-\omega))$, so that
$\xi(\tfrac12+\omega-iz)=A^\omega(z)-iB^\omega(z)$; these are real
entire functions, and they are the real and imaginary parts of
$\xi(\tfrac12+\omega-iz)$, up to sign, only for real $z$. For
$\omega>1$ and $a>0$, let $\mathsf H_{\omega,a}$ be the operators of
\cite[(2.6), p.~8]{Suzuki2012Canonical}; when
$1\pm\mathsf H_{\omega,a}$ are invertible, put
$m_\omega(a)=\det(1+\mathsf H_{\omega,a})/\det(1-\mathsf H_{\omega,a})$
(Fredholm determinants).

\begin{theorem}[{\cite[Theorem~2.3, p.~8]{Suzuki2012Canonical}}]
\label{thm:imp-sz12}
Let $\omega>1$. Then the operators $\mathsf H_{\omega,a}$ are
self-adjoint of Hilbert--Schmidt type, $1\pm\mathsf H_{\omega,a}$ are
invertible for all $a>0$, and the canonical system with Hamiltonian
$\diag(m_\omega(a)^{-2},m_\omega(a)^2)$ of the source has an explicit
solution $(A_a,B_a)$ with $(A_1,B_1)=(A^\omega,B^\omega)$.
\end{theorem}

\begin{theorem}[{\cite[Theorem~2.4, p.~9]{Suzuki2016Hamiltonians}}]
\label{thm:imp-sz16}
For $\omega>0$ and $\nu\in\Z_{>0}$ with $\nu\omega>1$ let
$K^{\omega,\nu}[t]$ and $J^{\omega,\nu}(t;z,z)$ be the Hankel-type
operators and endpoint kernels of the source for $L=\zeta$. RH holds if
and only if there is a sequence $(\omega_n,\nu_n)$ with $\omega_n$
strictly decreasing to $0$, $\nu_n\omega_n>1$,
$\det(1\pm K^{\omega_n,\nu_n}[t])\ne0$ for all $t\ge0$, and
$\lim_{t\to\infty}J^{\omega_n,\nu_n}(t;z,z)=0$ for all $z\in\Cp$.
\end{theorem}


For $f\in C_c^\infty(\R_+^*)$ write $\widetilde f(s)=\int_0^\infty
f(x)x^{s-1}\,dx$, $f^\sharp(x)=x^{-1}f(x^{-1})$, and $\bar f$ for the
complex conjugate; the convolution is $(f*g)(u)=\int_0^\infty
f(v)g(u/v)\,dv/v$. For a prime $p$ put
$W_p(f)=(\log p)\sum_{m\ge1}\bigl(f(p^m)+f^\sharp(p^m)\bigr)$, and put
$W_\R(f)=(\log4\pi+\gamma_E)f(1)+\int_1^\infty\bigl(f(x)+f^\sharp(x)
-\tfrac2xf(1)\bigr)\frac{dx}{x-x^{-1}}$, with $\gamma_E$ Euler's
constant.

\begin{theorem}[{Weil \cite{Weil1952Explicit},
\cite[\S16, p.~16]{Weil1972Explicit};
\cite[(148), p.~50, and Proposition~C.1,
p.~51]{ConnesConsani2020WeilPositivity}}]\label{thm:imp-weil-crit}
For $f\in C_c^\infty(\R_+^*)$ one has the explicit formula
$\sum_\rho\widetilde f(\rho)=\int_0^\infty f(x)\,dx+\int_0^\infty
f^\sharp(x)\,dx-\sum_vW_v(f)$, the sum on the left running over the
nontrivial zeros of $\zeta$, counted with multiplicity, and $v$ over
the places of $\Q$. Moreover, if $\mathcal P\subset\C$ is a finite set
containing $\{0,1\}$ and disjoint from the nontrivial zeros, then RH
holds if and only if $\sum_vW_v(g*\bar g^\sharp)\le0$ for every
$g\in C_c^\infty(\R_+^*)$ with $\widetilde g(z)=0$ for all
$z\in\mathcal P$.
\end{theorem}

The source states the explicit formula (148) without proof, following
Bombieri, and its proof of Proposition~C.1 refers to a proposition of
Yoshida \cite[Proposition~1 and Lemma~1, p.~285]{Yoshida1992Hermitian},
with a sketched refinement; the criterion is also recorded in
\cite[\S V, pp.~8--9]{Bombieri2000Clay}, and in a strengthened form in
\cite[Theorem~1, p.~191]{Bombieri2000Lincei}. It is due to Weil
\cite{Weil1952Explicit}, as reported in \cite[\S V, p.~9]{Bombieri2000Clay};
in \cite[\S16, p.~16]{Weil1972Explicit} Weil extends it to
Artin--Hecke $L$-functions, where positivity is equivalent to the
Riemann hypothesis together with Artin's conjecture.
Proposition~\ref{prop:weil-crit-notes} translates it
into the conventions of Definition~\ref{def:weil-form}.

Let $\xi_{\mathrm{CL}}(s)=s(s-1)\pi^{-s/2}\Gamma(s/2)\zeta(s)=2\xi(s)$,
the normalization of \cite{ConreyLi2000Positivity}.

\begin{theorem}[{\cite[\S3.1, (3.2) p.~6 and~(3.4) p.~7]%
{ConreyLi2000Positivity}}]\label{thm:imp-conrey-li}
For the $34$th nontrivial zero $\rho=\tfrac12+i\,111.0295355\ldots$ of
$\zeta$, $-\Re\{\xi_{\mathrm{CL}}'(\rho)\xi_{\mathrm{CL}}(1+\rho)\}
=-5.389100507\ldots\times10^{-69}<0$, and
$\Re\{\xi_{\mathrm{CL}}(1+282i)/\xi_{\mathrm{CL}}(2+282i)\}
=-0.000131957\ldots<0$.
\end{theorem}

These are high-precision numerical evaluations, not interval-certified
bounds. In the normalization $\xi=\tfrac12\xi_{\mathrm{CL}}$ the first
quantity is one quarter of the displayed value, about
$-1.3472751268\times10^{-69}$; the second is a ratio and is independent
of the normalization.

\section{De Branges spaces and the Weil form}\label{sec:debranges}

This section records what the analytic side can and cannot contribute
to the positivity required in Theorem~\ref{thm:rh}. Throughout, $F$ is
the function of Definition~\ref{def:zeta}(d), and we use the following
facts from Theorem~\ref{thm:imp-riemann}: $F$ is a real entire function
($F(\bar z)=\overline{F(z)}$), even, of order one; every zero of $F$ lies
in the open strip $\mathbb S=\{z\in\C:|\Im z|<\tfrac12\}$; and the
number $N(R)$ of zeros of $F$ in $|z|\le R$, counted with multiplicity,
satisfies $N(R)=O(R\log R)$ as $R\to\infty$, so that \eqref{eq:sum-r2}
holds. By Proposition~\ref{prop:F0}, $F(0)\ne0$. Moreover, $F$ has at
least one zero: an even entire function of order one without zeros is
$e^{A+Bz}$ with $B=0$, hence constant, whereas $\xi(s)\to\infty$ as
$s\to+\infty$ along the reals.

The Weil form $Q_W$ and its spectral expression are introduced in
Section~\ref{subsec:weil-form} (Definition~\ref{def:weil-form} and
Proposition~\ref{prop:weil-form}). By
Proposition~\ref{prop:weil-form}\ref{item:wf-rh}, RH implies $Q_W\ge0$.
The converse, that nonnegativity of $Q_W$ on a suitable class of tests
forces all zeros to be real, is Weil's positivity criterion
(Theorem~\ref{thm:imp-weil-crit}, in the present conventions
Proposition~\ref{prop:weil-crit-notes}); it is used in
Proposition~\ref{prop:weil-norm-converse}.

\subsection{The Hermite--Biehler criterion}\label{subsec:hb}

Since $F$ is even, $F(0)\ne0$, and a zero $r$ has the same order as
$-r\ne r$, the zeros listed with multiplicity split into pairs
$\{r,-r\}$. Fix a multiset $R_+$ containing one member of each pair.

\begin{proposition}\label{prop:hb}
Let $F$ be as in Definition~\ref{def:zeta}, let $\mathfrak m=-F'/F$,
and let $E=E_F=F+iF'$ (Definition~\ref{def:hb}), so that
$E^\#=F-iF'$. The following are equivalent.
\begin{enumerate}[label=(\roman*)]
\item\label{item:hb-rh} Every zero of $F$ is real (the Riemann
  hypothesis).
\item\label{item:hb-herglotz} $\mathfrak m$ is holomorphic on $\Cp$
  and $\Im\mathfrak m(z)>0$ for every $z\in\Cp$.
\item\label{item:hb-weak} $\Im\mathfrak m(z)\ge0$ for every $z\in\Cp$
  with
  $F(z)\ne0$.
\item\label{item:hb-hb} $E$ satisfies the HB inequality of
  Definition~\ref{def:hb}.
\end{enumerate}
\end{proposition}

For this $E$ the equivalence of the Riemann hypothesis with the HB
inequality is Lagarias's theorem
\cite[Theorem~1, p.~9]{Lagarias2006HilbertEntire}, whose proof passes
through the Herglotz condition \cite[(10), p.~9]{Lagarias2006HilbertEntire},
which Lagarias takes from \cite{Lagarias1999Positivity};
Proposition~\ref{prop:hb} is reproduced
to keep multiplicities explicit.
Common real zeros of $E$ and $E^\#$ are allowed in \ref{item:hb-hb}.
The proof uses two identities. First, the product and partial-fraction
expansions
\begin{equation}\label{eq:hb-product}
F(z)=F(0)\prod_{r\in R_+}\Bigl(1-\frac{z^2}{r^2}\Bigr),\qquad
\mathfrak m(z)=\sum_{r\in R_+}\Bigl(\frac1{r-z}+\frac1{-r-z}\Bigr),
\end{equation}
the product converging absolutely and locally uniformly on $\C$ and the
series absolutely and locally uniformly on $\C$ minus the zeros of $F$.
Second, the identity of entire functions
\begin{equation}\label{eq:hb-modulus}
|F+iF'|^2-|F-iF'|^2=4\Im\bigl(F\,\overline{F'}\bigr)
=4|F|^2\Im\Bigl(-\frac{F'}{F}\Bigr)\quad\text{where }F\ne0.
\end{equation}

\begin{proof}[Proof of \eqref{eq:hb-product} and \eqref{eq:hb-modulus}]
Since $F(\bar z)=\overline{F(z)}$, one has $F^\#=F$ and $(F')^\#=F'$,
so $E^\#=F-iF'$. By Hadamard's factorization theorem for the order-one
function $F$, $F(z)=e^{A+Bz}\prod_n(1-z/r_n)e^{z/r_n}$, and by
\eqref{eq:sum-r2} this product converges absolutely and locally
uniformly. Absolute convergence allows grouping $r$ with $-r$, and
$(1-z/r)e^{z/r}(1+z/r)e^{-z/r}=1-z^2/r^2$; thus $F=e^{A+Bz}G$ with $G$
the even product in \eqref{eq:hb-product}. Evenness of $F$ and $G$ gives
$e^{2Bz}=1$ off the zeros of $G$, so $B=0$, and $e^A=F(0)$. Taking the
logarithmic derivative of the locally uniformly convergent product
termwise and using $\frac{d}{dz}\log(1-z^2/r^2)=-2z/(r^2-z^2)$ and
$\frac1{r-z}-\frac1{r+z}=\frac{2z}{r^2-z^2}$ gives the series for
$\mathfrak m$;
for $|r|\ge2|z|$ its terms are bounded by $8|z|/(3|r|^2)$, so
\eqref{eq:sum-r2} gives locally uniform absolute convergence. For
\eqref{eq:hb-modulus}, $|a+b|^2-|a-b|^2=4\Re(a\bar b)$ with $a=F$,
$b=iF'$ gives $4\Re(-iF\overline{F'})=4\Im(F\overline{F'})$, and where
$F\ne0$ one has $\Im(F\overline{F'})=|F|^2\Im(\overline{F'}/\overline
F)=|F|^2\Im(-F'/F)$.
\end{proof}

\begin{proof}[Proof of
\ref{item:hb-rh}$\Rightarrow$\ref{item:hb-herglotz}]
If all zeros are real then $F$ has no zero in $\Cp$ and $\mathfrak m$
is holomorphic there. For real $r$ and $z=x+iy$ with $y>0$,
$\Im\frac1{r-z}=\frac{y}{(x-r)^2+y^2}>0$. Taking imaginary parts
in \eqref{eq:hb-product},
\[
\Im\mathfrak m(z)=\sum_{\gamma\in\Gamma}m_\gamma
\frac{y}{(x-\gamma)^2+y^2},
\]
a convergent series of positive terms; it is nonempty because $F$ has
a zero.
\end{proof}

\begin{proof}[Proof of
\ref{item:hb-herglotz}$\Rightarrow$\ref{item:hb-weak}]
Immediate.
\end{proof}

\begin{proof}[Proof of \ref{item:hb-weak}$\Rightarrow$\ref{item:hb-rh}]
Suppose $F(z_0)=0$ with $z_0\in\Cp$, of order $k\ge1$. Write
$F=(z-z_0)^kg$ near $z_0$ with $g(z_0)\ne0$, so that
$\mathfrak m(z)=-k/(z-z_0)-g'(z)/g(z)$ and $|g'/g|\le M$ on a disc
$|z-z_0|\le\eta$ with $0<\eta<\Im z_0$ on which $z_0$ is the only zero
of $F$. For $z=z_0-i\epsilon$, $0<\epsilon<\eta$, we have $z\in\Cp$,
$F(z)\ne0$, and $-k/(z-z_0)=-ik/\epsilon$, so
$\Im\mathfrak m(z)\le-k/\epsilon+M<0$ once $\epsilon<k/M$. This
contradicts
\ref{item:hb-weak}; hence $F$ has no zero in $\Cp$. If $F(z_0)=0$ with
$\Im z_0<0$ then $F(\bar z_0)=\overline{F(z_0)}=0$ with $\bar
z_0\in\Cp$, which is excluded. So every zero is real.
\end{proof}

\begin{proof}[Proof of
\ref{item:hb-herglotz}$\Rightarrow$\ref{item:hb-hb}]
At a zero of $F$ of order $k$, $\mathfrak m$ has a simple pole with
residue $-k$; so holomorphy of $\mathfrak m$ on $\Cp$ means $F\ne0$ on
$\Cp$. By \eqref{eq:hb-modulus}, $|E|^2-|E^\#|^2=4|F|^2\Im\mathfrak m>0$
on $\Cp$.
\end{proof}

\begin{proof}[Proof of \ref{item:hb-hb}$\Rightarrow$\ref{item:hb-rh}]
If $F(z_0)=0$ with $z_0\in\Cp$, then $E(z_0)=iF'(z_0)$ and
$E^\#(z_0)=-iF'(z_0)$ have equal moduli (both zero if the zero is
multiple), contradicting the strict HB inequality at $z_0$. Hence $F$
has no zero in $\Cp$, and by reality none in the lower half-plane.
\end{proof}

\begin{remark}\label{rem:hb}
Let $r$ be a real zero of $F$ of order $k$, and write
$F=(z-r)^kg$ with $g(r)\ne0$. Then
$E=(z-r)^{k-1}\bigl[(z-r)g+i(kg+(z-r)g')\bigr]$, and the bracket
equals $ikg(r)\ne0$ at $z=r$; likewise for $E^\#$. So a real zero of
order $k$ produces a common real zero of $E$ and $E^\#$ of order exactly
$k-1$. Proposition~\ref{prop:hb} is therefore an equivalence with RH
itself, not with RH together with simplicity of the zeros; a convention
forbidding real zeros of the structure function would, for this $E$,
add simplicity \cite[\S1]{Lagarias2006HilbertEntire}. Discarding the
common factor loses divisor data (Proposition~\ref{prop:R9}). Under RH
the expression for $\Im\mathfrak m$ in the proof is the Poisson
integral of the
measure $\sum_\gamma m_\gamma\delta_\gamma$: multiplicities enter as
weights, not as dimensions of eigenspaces. In the variable
$s=\tfrac12-iz$ one has $F'(z)=-i\xi'(s)$, so
$\mathfrak m(z)=i\xi'(s)/\xi(s)$ and
$\Im\mathfrak m(z)=\Re(\xi'(s)/\xi(s))$; since $\Im z>0$ exactly when
$\Re s>\tfrac12$, condition~\ref{item:hb-herglotz} reads
$\Re(\xi'/\xi)>0$ on $\Re s>\tfrac12$
\cite[(10), p.~9]{Lagarias2006HilbertEntire}.
\end{remark}

\begin{proposition}\label{prop:hb-high}
Let $F$ be as in Definition~\ref{def:zeta} and $\mathfrak m=-F'/F$.
Then $F$ has no zero in the closed half-plane $\{\Im z\ge\tfrac12\}$,
and $\Im\mathfrak m(z)>0$ for every $z$ with $\Im z\ge\tfrac12$.
\end{proposition}

\begin{proof}
All zeros lie in $\mathbb S$, so $F\ne0$ on $\Im z\ge\tfrac12$. For a zero $r$
and $z=x+iy$,
$\Im\frac1{r-z}=\frac{y-\Im r}{|r-z|^2}$, and likewise
$\Im\frac1{-r-z}=\frac{y+\Im r}{|r+z|^2}$. If $y\ge\tfrac12$ then
$y\mp\Im r>0$ because $|\Im r|<\tfrac12$. So every term of the series
\eqref{eq:hb-product} has positive imaginary part, and so does the
convergent sum.
\end{proof}

Proposition~\ref{prop:hb-high} is unconditional, but its argument uses
the location of the zeros term by term and gives nothing on
$0<\Im z<\tfrac12$. A zero $z_0$ with $0<\Im z_0<\tfrac12$ makes
$\Im\mathfrak m$ negative just below $z_0$ (proof of
\ref{item:hb-weak}$\Rightarrow$\ref{item:hb-rh}), while $\Im\mathfrak m$
remains
positive above $\tfrac12$; an explicit even real polynomial exhibiting
both behaviours is given in Proposition~\ref{prop:R11}. Positivity high
in $\Cp$ therefore does not propagate down to the real axis.

\subsection{Canonical systems}\label{subsec:canonical}

Let $J_0=\bigl(\begin{smallmatrix}0&-1\\1&0\end{smallmatrix}\bigr)$
and let $H:[0,\infty)\to M_2(\R)$ be locally integrable with $H(t)$
symmetric and positive semidefinite for almost every $t$. For $z\in\C$
consider
\begin{equation}\label{eq:canonical}
\partial_tY(t,z)=zJ_0H(t)Y(t,z),\qquad Y(0,z)=(1,0)^T.
\end{equation}
The Picard series $Y=\sum_{n\ge0}z^nY_n$, with $Y_0=(1,0)^T$ and
$Y_n(t)=\int_0^tJ_0H(u)Y_{n-1}(u)\,du$, satisfies
$|Y_n(t)|\le(\int_0^t\|H\|)^n/n!$; hence it defines a locally
absolutely continuous solution in $t$ which is entire in $z$, and
Gronwall's inequality shows it is the unique locally absolutely
continuous solution, also backwards from any initial time. Write
$Y=(A,B)^T$ and $E_t(z)=A(t,z)-iB(t,z)$.

\begin{proposition}\label{prop:canonical}
Let $H$, $J_0$, and $Y$ be as in \eqref{eq:canonical}, and let
$t\ge0$ and $z,w\in\C$. Then:
\begin{enumerate}[label=(\roman*)]
\item\label{item:can-real} $Y(t,\bar z)=\overline{Y(t,z)}$; hence
  $E_t^\#=A+iB$;
\item\label{item:can-lagrange} for almost every $t$,
  \begin{equation}\label{eq:lagrange}
  \partial_t\bigl(Y(t,w)^*J_0Y(t,z)\bigr)
  =(\bar w-z)\,Y(t,w)^*H(t)Y(t,z);
  \end{equation}
\item\label{item:can-hb}
  \begin{equation}\label{eq:can-hb}
  |E_t(z)|^2-|E_t^\#(z)|^2
  =4\Im z\int_0^tY(u,z)^*H(u)Y(u,z)\,du;
  \end{equation}
\item\label{item:can-pos} if $H(u)$ is positive definite for almost
  every $u$ and $t>0$, then $E_t$ satisfies the HB inequality.
\end{enumerate}
\end{proposition}

\begin{proof}[Proof of \ref{item:can-real}]
Since $H$ and $J_0$ are real, $\overline{Y(\cdot,z)}$ solves
\eqref{eq:canonical} with $\bar z$ in place of $z$ and the same initial
value; uniqueness gives the first claim. Then
$E_t^\#(z)=\overline{A(t,\bar z)}+i\,\overline{B(t,\bar z)}
=A(t,z)+iB(t,z)$.
\end{proof}

\begin{proof}[Proof of \ref{item:can-lagrange}]
Since $J_0^*=J_0^T=-J_0$, $J_0^2=-I$, and $H^*=H$, the adjoint of
$\partial_tY(t,w)=wJ_0HY(t,w)$ is
$\partial_tY(t,w)^*=-\bar wY(t,w)^*HJ_0$. Products of locally
absolutely continuous functions are locally absolutely continuous, and
almost everywhere
\[
\partial_t(Y_w^*J_0Y_z)=-\bar wY_w^*HJ_0^2Y_z+zY_w^*J_0^2HY_z
=(\bar w-z)Y_w^*HY_z,
\]
where $Y_w=Y(t,w)$ and $Y_z=Y(t,z)$.
\end{proof}

\begin{proof}[Proof of \ref{item:can-hb}]
For $Y=(A,B)^T$ one computes $Y^*J_0Y=-\bar AB+\bar BA=2i\Im(A\bar B)$,
which vanishes at $t=0$. Integrating \eqref{eq:lagrange} with $w=z$,
$2i\Im(A\bar B)(t,z)=(\bar z-z)\int_0^tY^*HY\,du$, that is,
$\Im(A\bar B)=-\Im z\int_0^tY^*HY\,du$. By \ref{item:can-real},
$|E_t|^2-|E_t^\#|^2=|A-iB|^2-|A+iB|^2=4\Re(A\,\overline{-iB})
=-4\Im(A\bar B)$.
\end{proof}

\begin{proof}[Proof of \ref{item:can-pos}]
If $Y(u_0,z)=0$ for some $u_0$, uniqueness from $u_0$ gives
$Y(\cdot,z)\equiv0$, contradicting $Y(0,z)=(1,0)^T$. Hence
$Y(u,z)^*H(u)Y(u,z)>0$ for almost every $u$, the integral in
\eqref{eq:can-hb} is positive for $t>0$, and \eqref{eq:can-hb} gives
$|E_t(z)|>|E_t^\#(z)|$ for $z\in\Cp$.
\end{proof}

For $H=I$ one has $Y(t,z)=(\cos zt,\sin zt)^T$ and $E_t(z)=e^{-izt}$;
with $y=\Im z$, both sides of \eqref{eq:can-hb} equal
$2\sinh(2ty)$, since $|\cos w|^2+|\sin w|^2=\cosh(2\Im w)$.

Proposition~\ref{prop:canonical} shows that a positive Hamiltonian on
an initial interval produces HB functions $E_t$. To reach
$\Xi+i\Xi'$, or a shifted $\xi$, one must in addition prove the
endpoint limit, nonvanishing, and normalization, and the relevant domain
properties. Conversely, an inverse theorem for canonical systems whose
input is an HB function cannot be applied to $\Xi+i\Xi'$ without first
proving Proposition~\ref{prop:hb}\ref{item:hb-hb}
(Proposition~\ref{prop:R11}).

It is natural to expect that a canonical system with positive
Hamiltonian on an initial interval, whose solutions are entire, already
forces the HB inequality for the function it is meant to produce at
the endpoint; in fact positivity on an initial interval imposes no sign
on a prescribed terminal function.

\begin{remark}\label{rem:N-short-interval}
Let $H$ be as in \eqref{eq:canonical} on $[0,t_0]$ and let $t_1>t_0$.
Integrating \eqref{eq:lagrange} from $t_0$ to $t_1$ gives
\[
|E_{t_1}(z)|^2-|E_{t_1}^\#(z)|^2
=|E_{t_0}(z)|^2-|E_{t_0}^\#(z)|^2
+4\Im z\int_{t_0}^{t_1}Y^*HY\,du,
\]
so the sign at $t_1$ is controlled only if $H\ge0$ on all of
$[t_0,t_1]$; at an infinite endpoint, the terminal boundary term
$\lim_{t\to\infty}(|E_t|^2-|E_t^\#|^2)$ must itself be identified and
controlled. Positivity of $H$ on $[0,t_0]$ says nothing about the
Hamiltonian, if any, that continues the system to an endpoint with a
prescribed terminal function such as $\Xi+i\Xi'$.
\end{remark}

Suzuki constructs, for the shifted functions $\xi(\tfrac12+\omega-iz)$
with $\omega>1$, canonical systems with positive diagonal Hamiltonians
whose solutions match the shifted function
(Theorem~\ref{thm:imp-sz12}); his systems are written in a
multiplicative variable $a$ with operator $-a\partial_a$, so comparison
with \eqref{eq:canonical} requires a change of variable and a reversal
of time. For powers $\xi(\tfrac12+\omega-iz)^\nu$ with $\nu\omega>1$ he
obtains a positive Hamiltonian on some initial interval unconditionally,
and shows that RH is equivalent to the existence of shifts
$\omega_n\downarrow0$ and integers $\nu_n$ with $\nu_n\omega_n>1$ for
which the relevant determinants do not vanish for all $t\ge0$ and an
endpoint kernel tends to zero (Theorem~\ref{thm:imp-sz16}).

It is natural to expect that these results, pushed to small shifts or
large exponents, reach $\xi$ itself; in fact the checked range stops
short of it.

\begin{remark}\label{rem:N-suzuki-range}
Theorem~\ref{thm:imp-sz12} covers only shifts $\omega>1$, where the
shifted function has no zeros in $\Cp$ unconditionally; it does not
give the corresponding system for every $\omega>0$. In
Theorem~\ref{thm:imp-sz16}, increasing $\nu$ improves the regularity of
the kernel but does not prove RH: the content of the criterion is the
global invertibility $\det(1\pm K^{\omega_n,\nu_n}[t])\ne0$, for all
$t\ge0$, of the kernel operators of that theorem, together with the
endpoint condition, and the unconditional part is the existence of the
system on an initial interval.
\end{remark}

\subsection{Norm identities and the converse}\label{subsec:norm}

Let $\mathcal T_W=\{\psi\in C_c^\infty(\R):\widehat\psi(\tfrac i2)
=\widehat\psi(-\tfrac i2)=0\}$.

\begin{proposition}[Weil]\label{prop:weil-crit-notes}
Let $F$ and $Q_W$ be as in Definition~\ref{def:weil-form}. Then every
zero of $F$ is real if and only if $Q_W(\psi)\ge0$ for every
$\psi\in\mathcal T_W$.
\end{proposition}

\begin{proof}
For $\psi\in C_c^\infty(\R)$ put $g(x)=x^{-1/2}\psi(\log x)$; this is a
bijection $C_c^\infty(\R)\to C_c^\infty(\R_+^*)$. With the notation of
Theorem~\ref{thm:imp-weil-crit} and $x=e^t$,
$\widetilde g(\tfrac12-iz)=\int_\R\psi(t)e^{-izt}\,dt=\widehat\psi(-z)$.
Hence $\widetilde g(0)=\widehat\psi(\tfrac i2)$ and
$\widetilde g(1)=\widehat\psi(-\tfrac i2)$, so $g$ satisfies the
vanishing conditions for $\mathcal P=\{0,1\}$ exactly when
$\psi\in\mathcal T_W$; and $\{0,1\}$ contains no nontrivial zero. Let
$f=g*\bar g^\sharp$. The Mellin transform turns $*$ into the product,
and the substitution $x=1/y$ gives $\widetilde{\bar g^\sharp}(s)=
\overline{\widetilde g(1-\bar s)}$; thus $\widetilde f(s)=\widetilde
g(s)\overline{\widetilde g(1-\bar s)}$. In particular
$\int_0^\infty f\,dx=\widetilde f(1)=0$ and
$\int_0^\infty f^\sharp\,dx=\widetilde f(0)=0$ when
$\psi\in\mathcal T_W$. A nontrivial zero $\rho$ corresponds to the zero
$z$ of $F$ with $\rho=\tfrac12-iz$, with the same multiplicity, and then
$1-\bar\rho=\tfrac12-i\bar z$. The explicit formula of
Theorem~\ref{thm:imp-weil-crit} therefore gives
\[
-\sum_vW_v(g*\bar g^\sharp)=\sum_{z}\widehat\psi(-z)
\overline{\widehat\psi(-\bar z)}=\sum_z\widehat\psi(z)
\overline{\widehat\psi(\bar z)}=Q_W(\psi),
\]
the sums running over the zeros of $F$ with multiplicity; the middle
equality uses that the zeros of the even function $F$ are symmetric
under $z\mapsto-z$ with multiplicity, and the rearrangement is justified
by absolute convergence (Proposition~\ref{prop:weil-form}). The
criterion of Theorem~\ref{thm:imp-weil-crit} is thus exactly the
statement.
\end{proof}

\begin{definition}\label{def:norm-realization}
Let $Q_W$ be as in Definition~\ref{def:weil-form} and let
$\mathcal T_0\subset C_c^\infty(\R)$ be a linear subspace. A
\emph{norm realization} of $Q_W$ on $\mathcal T_0$ is a pair
$(\mathcal K,\mathcal T)$ of a Hilbert space $\mathcal K$ and a linear
map $\mathcal T:\mathcal T_0\to\mathcal K$ with
$\|\mathcal T\psi\|^2=\pi Q_W(\psi)$ for every $\psi\in\mathcal T_0$.
\end{definition}

\begin{proposition}\label{prop:weil-norm-converse}
Let $F$ and $Q_W$ be as in Definition~\ref{def:weil-form}, let
$\mathcal T_W$ be as in Proposition~\ref{prop:weil-crit-notes}, and
suppose that $Q_W$ has a norm realization $(\mathcal K,\mathcal T)$ on
a subspace $\mathcal T_0\supset\mathcal T_W$. Then every zero of $F$ is
real.
\end{proposition}

\begin{proof}
For $\psi\in\mathcal T_W\subset\mathcal T_0$,
$Q_W(\psi)=\pi^{-1}\|\mathcal T\psi\|^2\ge0$.
Proposition~\ref{prop:weil-crit-notes} applies.
\end{proof}

Linearity of $\mathcal T$ is not used in the proof. A norm realization
on all of $C_c^\infty(\R)$ restricts to $\mathcal T_W$; one defined
only on a smaller class, for instance on even tests, suffices only if
that class contains $\mathcal T_W$. The hypothesis cannot be weakened
to a norm identity on a single support window (Proposition~\ref{prop:R14}).
Under RH a norm realization on $C_c^\infty(\R)$ exists trivially:
$\psi\mapsto(\sqrt{\pi m_\gamma}\,\widehat\psi(\gamma))_\gamma$ lands in
$\ell^2(\Gamma)$ by Proposition~\ref{prop:weil-form} and has the
required norm. As with Proposition~\ref{prop:converse}, the content lies
in constructing $\mathcal T$ and $\mathcal K$ independently of the
location of the zeros.

\subsection{Suzuki's comparison}\label{subsec:suzuki}

Let $E=E_F=F+iF'$, $\vartheta=E^\#/E$, $\mathcal K(\vartheta)$,
$\mathcal H_W$, $\|\cdot\|_W$, and $\psi_G=\mathsf F^{-1}G$ be as in
Section~\ref{sec:imported}, before Theorem~\ref{thm:imp-sz23}.
Under RH, $E$ satisfies the HB inequality
(Proposition~\ref{prop:hb}), so $E$ has no zero in $\Cp$ and
$|\vartheta|<1$ there. Theorem~\ref{thm:imp-sz23}
\cite[Theorem~1.1]{Suzuki2023WeilHilbert} states that, assuming RH, the
Hilbert space $\mathcal H_W$ obtained by completing $C_c^\infty(\R)$
for $Q_W$ is identified with $\mathcal K(\vartheta)$ through the
inverse Fourier transform, with
$\|G\|_{\mathcal K(\vartheta)}^2=\pi\|[\psi_G]\|_W^2$. This is an
identification of Hilbert spaces under RH, not a construction of a
positive space from which RH would follow.

The same source proposes an explicit transform realizing the norm
identity for every test. With the printed definitions, which extend its
time-dependent kernels evenly to negative $t$, the transform vanishes on
every even test, so the norm it defines is degenerate
(Proposition~\ref{prop:R17}). The following reconstruction defines the
zero series for all real $t$ without even extension. Since
$m_\gamma=m_{\bar\gamma}$, it agrees for $t<0$ with the rule
$P_t(z)=\overline{P_{-t}(\bar z)}$ of the author's correction
\cite[p.~27]{Suzuki2026WeilScrew}, which also replaces the factor $i$
in the source's (1.5) by $-i$ and so multiplies $S_t$ by $-1$.

For $t\in\R$ and $z\in\C$ away from the zeros of $F$ put
\begin{equation}\label{eq:rec-P}
P_t^{\mathrm{rec}}(z)=\sum_{\gamma\in\Gamma}m_\gamma
\frac{e^{-i\gamma t}-1}{\gamma(z-\gamma)},\qquad
S_t^{\mathrm{rec}}(z)=\frac{i\bigl(1+\vartheta^\#(z)\bigr)}{2}\,
P_t^{\mathrm{rec}}(z).
\end{equation}
The series converges absolutely, uniformly for $t$ in a compact set and
$z$ in a compact set disjoint from $\Gamma$. Indeed, since
$|\Im\gamma|<\tfrac12$, one has $|e^{-i\gamma t}-1|\le2e^{|t|/2}$; for
$z$ in a compact set, $|z-\gamma|\ge|\gamma|/2$ for all but finitely
many $\gamma$; so all but finitely many terms are bounded by
$4e^{|t|/2}m_\gamma|\gamma|^{-2}$, and
$\sum_\gamma m_\gamma|\gamma|^{-2}<\infty$ by \eqref{eq:sum-r2}. This is
the justification given after \cite[(3.2), p.~7]{Suzuki2023WeilHilbert}
for $t\ge0$.
For $\psi\in C_c^\infty(\R)$ define
\begin{equation}\label{eq:rec-T}
\mathcal T_{\mathrm{rec}}\psi
=\int_\R(S_t^{\mathrm{rec}})^\#\,i\psi'(t)\,dt .
\end{equation}
Since $\overline{\vartheta^\#(\bar z)}=\vartheta(z)$ and
$\int_\R(e^{i\bar\gamma t}-1)\,i\psi'(t)\,dt
=\widehat{i\psi'}(\bar\gamma)=\bar\gamma\widehat\psi(\bar\gamma)$,
exchanging sum and integral (justified by the same majorant, since
$\psi'$ has compact support) and reindexing by
$\gamma\mapsto\bar\gamma$ gives, pointwise off $\Gamma$ and wherever
$\vartheta$ is defined,
\begin{equation}\label{eq:rec-expansion}
\mathcal T_{\mathrm{rec}}\psi
=-\sum_{\gamma\in\Gamma}\sqrt{\pi m_\gamma}\,\widehat\psi(\gamma)
f_\gamma,\qquad
f_\gamma(z)=\sqrt{\frac{m_\gamma}{\pi}}\,
\frac{i\bigl(1+\vartheta(z)\bigr)}{2(z-\gamma)} .
\end{equation}
The functions $f_\gamma$ are those of \cite[(3.5), p.~10]%
{Suzuki2023WeilHilbert}, and the overall sign of
\eqref{eq:rec-expansion} is opposite to that of \cite[(3.8), p.~10]%
{Suzuki2023WeilHilbert}, which does not affect norms. Under RH, the
$f_\gamma$ form an orthonormal basis of $\mathcal K(\vartheta)$
\cite[Proposition~4.1, p.~10]{Suzuki2023WeilHilbert}, an imported
statement used only in this paragraph; it makes the series
\eqref{eq:rec-expansion} converge in $L^2(\R)$ with
$\|\mathcal T_{\mathrm{rec}}\psi\|_2^2=\pi Q_W(\psi)$. Without RH,
neither the $L^2$ boundary behaviour of $\mathcal T_{\mathrm{rec}}\psi$
nor a prime-side formula for negative $t$ is available.

\begin{question}\label{q:reconstruction}
Let $\mathcal T_{\mathrm{rec}}$ be defined by \eqref{eq:rec-P} and
\eqref{eq:rec-T}. Without assuming RH, is $\mathcal T_{\mathrm{rec}}\psi$
the boundary value of a function whose restriction to $\R$ lies in
$L^2(\R)$, for every $\psi\in C_c^\infty(\R)$, and does
$\|\mathcal T_{\mathrm{rec}}\psi\|_{L^2(\R)}^2=\pi Q_W(\psi)$ hold for
every such $\psi$?
\end{question}

An affirmative answer to Question~\ref{q:reconstruction} would imply RH
by Proposition~\ref{prop:weil-norm-converse}, provided the Hilbert space
is $L^2(\R)$ with its own (positive) norm.
By the agreement just noted, Question~\ref{q:reconstruction} concerns
the corrected definition of the source; the unconditional norm identity
is also stated as a target in \cite[pp.~26--27]{Suzuki2026WeilScrew}.

\emph{Multiplicities.} Under RH, Proposition~\ref{prop:weil-form} gives
one evaluation coordinate $\widehat\psi(\gamma)$ per distinct zero,
carrying weight $m_\gamma$. The self-adjoint operator specified in
\cite[\S6]{Suzuki2023WeilHilbert} has multiplicity one at each distinct
zero. The ordinary trace of $f$ of an operator on this scalar completion
is $\sum_\gamma f(\gamma)$, whereas the divisor identity of
\ref{ax:AW2}, the determinant identity \eqref{eq:aw-det}, and the trace in
\ref{ax:AW6} require
$\sum_\gamma m_\gamma f(\gamma)$. A commuting weight operator, a
weighted trace, or a Hilbert space with fibres $\C^{m_\gamma}$ would
supply the difference; these are distinct constructions, each needing
independent arithmetic meaning. A cyclic self-adjoint operator cannot
have a two-dimensional eigenspace (Proposition~\ref{prop:R9}). For zero
location alone, an independently established scalar norm identity
suffices (Proposition~\ref{prop:weil-norm-converse}).

\emph{Semilocal spaces.} For a finite set $S$ of places containing
$\infty$ and a cutoff $\lambda>0$, Connes, Consani, and Moscovici
construct from Sonin spaces an isomorphism of hilbertian spaces (a
linear homeomorphism; no preferred inner product is transported) onto
a space $\mathcal B_\lambda$ of entire functions, which with its
$S$-dependent norm is a de Branges space $\mathcal B_\lambda^S$
(Theorem~\ref{thm:imp-ccm23}). This comparison involves arithmetic
places, entire functions, and positivity unconditionally. It does not
identify $\mathcal B_\lambda^S$ with $\mathcal K(\vartheta)$ or
$\mathcal H_W$, nor prove convergence as $S$ and $\lambda$ grow; an
abstract isomorphism of Hilbert spaces is too weak to transport spectra
(Proposition~\ref{prop:R10}).

\subsection{Spectral approximation}\label{subsec:gap}

For $\ell>0$ we regard $L^2([-\ell,\ell])$ as the subspace of
$L^2(\R)$ of functions vanishing outside $[-\ell,\ell]$; its elements
are integrable with compact support, so their Fourier transforms
(Definition~\ref{def:weil-form}) are entire.

\begin{definition}\label{def:residual-gap}
Let $\ell>0$.
\begin{enumerate}[label=(\alph*)]
\item A \emph{gap datum} on $[-\ell,\ell]$ is a tuple
  $(\mathcal E,A,v,\mu)$ with $\mathcal E\subset L^2([-\ell,\ell])$ a
  finite-dimensional subspace, $A$ a self-adjoint operator on
  $\mathcal E$, $v\in\mathcal E$ a unit eigenvector of $A$, and
  $\mu\in\R$, such that
  $d=\operatorname{dist}\bigl(\mu,\spec(A|_{v^\perp})\bigr)>0$, where
  $v^\perp$ is the orthogonal complement of $v$ in $\mathcal E$ (and
  $d=+\infty$ if $v^\perp=0$).
\item For $k\in\mathcal E$, its \emph{residual} is
  $r=\|(A-\mu)k\|_2$ and its \emph{projection} is $w=Pk$, where $P$ is
  the orthogonal projection onto $\C v$.
\item For $b\ge0$ put $C(\ell,b)=\sqrt{\sinh(2b\ell)/b}$ if $b>0$ and
  $C(\ell,0)=\sqrt{2\ell}$.
\end{enumerate}
\end{definition}

The subspace $v^\perp$ is $A$-invariant because $v$ is an eigenvector
and $A$ is self-adjoint, so $\spec(A|_{v^\perp})$ is defined.

\begin{proposition}\label{prop:gap}
Let $\ell>0$, let $(\mathcal E,A,v,\mu)$ be a gap datum on
$[-\ell,\ell]$ with gap $d$, and let $k\in\mathcal E$ have residual $r$
and projection $w$. Then
\begin{equation}\label{eq:gap-l2}
\|k-w\|_2\le r/d,
\end{equation}
and for every $b\ge0$,
\begin{equation}\label{eq:gap-strip}
\sup_{|\Im z|\le b}\bigl|\widehat k(z)-\widehat w(z)\bigr|
\le C(\ell,b)\,r/d .
\end{equation}
\end{proposition}

\begin{proof}
Choose an orthonormal eigenbasis $v=e_0,e_1,\dots,e_n$ of $A$ with
eigenvalues $\lambda_0,\dots,\lambda_n$, and write
$k=\sum_ic_ie_i$. Then $w=c_0e_0$ and
$r^2=\sum_i|\lambda_i-\mu|^2|c_i|^2\ge d^2\sum_{i\ge1}|c_i|^2
=d^2\|k-w\|_2^2$, whatever $c_0$ and $\lambda_0$ are; this is
\eqref{eq:gap-l2}. For $g\in L^2([-\ell,\ell])$ and $z=x+iy$, the
Cauchy--Schwarz inequality gives
\[
|\widehat g(z)|^2\le\|g\|_2^2\int_{-\ell}^{\ell}e^{-2yt}\,dt
=\|g\|_2^2\,\frac{\sinh(2y\ell)}{y}
\]
for $y\ne0$, and $\|g\|_2^2\cdot2\ell$ for $y=0$. The function
$y\mapsto\sinh(2y\ell)/y$ is even, with a power series in $y^2$ having
positive coefficients and value $2\ell$ at $y=0$, so on $|y|\le b$ it is
at most $C(\ell,b)^2$. Applying this to $g=k-w$ and using
\eqref{eq:gap-l2} proves \eqref{eq:gap-strip}.
\end{proof}

We use Hurwitz's theorem in the following form: if $f_j\to f$
uniformly on a closed disc $D$ inside a domain where all $f_j$ and $f$
are holomorphic, and $f$ has a zero in the interior of $D$ and none on
$\partial D$, then $f_j$ has a zero in $D$ for all large $j$. (By
Rouch\'e's theorem, once $\sup_{\partial D}|f_j-f|<\min_{\partial
D}|f|$.)

\begin{proposition}\label{prop:rouche-route}
Let $F$ be as in Definition~\ref{def:zeta} and
$\mathbb S=\{|\Im z|<\tfrac12\}$.
For each $j\ge1$ let $\ell_j>0$, let $(\mathcal E_j,A_j,v_j,\mu_j)$ be a
gap datum on $[-\ell_j,\ell_j]$ with gap $d_j$, and let
$k_j\in\mathcal E_j$ have residual $r_j$ and projection $w_j$. Assume:
\begin{enumerate}[label=(\alph*)]
\item\label{item:rr-real} every zero of $\widehat v_j$ is real, for
  every $j$;
\item\label{item:rr-conv} $\widehat k_j\to F$ uniformly on compact
  subsets of $\mathbb S$;
\item\label{item:rr-gap} $C(\ell_j,b)\,r_j/d_j\to0$ as $j\to\infty$,
  for every $0<b<\tfrac12$.
\end{enumerate}
Then every zero of $F$ is real.
\end{proposition}

\begin{proof}
Let $L\subset\mathbb S$ be compact; then $L\subset\{|\Im z|\le b\}$
for some $0<b<\tfrac12$. By Proposition~\ref{prop:gap} and
\ref{item:rr-gap}, $\widehat k_j-\widehat w_j\to0$ uniformly on $L$,
so by \ref{item:rr-conv}, $\widehat w_j\to F$ uniformly on $L$. Since
$F(0)>0$ (Proposition~\ref{prop:F0}), $\widehat w_j(0)\ne0$ for large
$j$; as $w_j=\langle k_j,v_j\rangle v_j$, this forces
$\langle k_j,v_j\rangle\ne0$, and then $\widehat w_j$ and $\widehat v_j$
have the same zeros, all real by \ref{item:rr-real}. Suppose $F(z_0)=0$
with $z_0\notin\R$. Then $z_0\in\mathbb S$. Choose $\eta$ with
$0<\eta<\min(|\Im z_0|,\tfrac12-|\Im z_0|)$ such that $F$ has no zero on
$|z-z_0|=\eta$ (zeros of $F\not\equiv0$ are isolated). The closed disc
$D=\{|z-z_0|\le\eta\}$ lies in $\mathbb S$ and does not meet $\R$. By Hurwitz's
theorem, $\widehat w_j$ has a zero in $D$ for all large $j$, a nonreal
zero, contradicting what was shown. Hence every zero of $F$, all of
which lie in $\mathbb S$, is real.
\end{proof}

The argument uses neither $\ell_j\to\infty$, which is the regime in
which the hypotheses are expected to be attainable, nor any assumption
on the multiplicities of the real zeros of $F$; and \ref{item:rr-conv}
may be weakened to $\widehat k_j\to cF$ with a constant $c\ne0$, which
absorbs normalization factors such as the factor $4$ of
Proposition~\ref{prop:R15}. Since $C(\ell,b)\sim e^{b\ell}/\sqrt{2b}$ as
$\ell\to\infty$, condition~\ref{item:rr-gap} asks, in terms of
$\lambda_j=e^{\ell_j}$, that $\lambda_j^br_j/d_j\to0$ for every
$b<\tfrac12$; the estimate $r_j/d_j=O(\lambda_j^{-1/2})$ suffices.
Neither a small residual without a gap, nor unweighted $L^2$
convergence on growing supports, nor approximation in norm without
graph-norm control of an unbounded operator gives
\ref{item:rr-conv}--\ref{item:rr-gap} (Proposition~\ref{prop:R16}).

Hypothesis~\ref{item:rr-real} is where the arithmetic enters. In the
first arXiv version of \cite[Theorem~5.10]
{ConnesConsaniMoscovici2025ZetaTriples}, the compression of the Weil
quadratic form to periodic Fourier modes on an interval is a Hermitian
matrix, and, if its lowest eigenvalue is simple, a corresponding
eigenvector $v$ is even, and a boundary functional does not vanish on
$v$, then $\widehat v$ has only real zeros. The positivity used there is
that of $Q-\varepsilon_{\min}I$ on a quotient, which does not give
positivity of the Weil form itself (Proposition~\ref{prop:R14}). What
remains open along this route is a bound on $r_j/d_j$ satisfying
\ref{item:rr-gap}, together with verification of the simplicity,
evenness, and boundary hypotheses along an unbounded sequence of
intervals.

\subsection{Two comparison targets}\label{subsec:targets}

The discussion above isolates two targets of different strength.

\begin{question}\label{q:norm-target}
Let $Q_W$ be as in Definition~\ref{def:weil-form} and $\mathcal T_W$
as in Proposition~\ref{prop:weil-crit-notes}. Construct from arithmetic or
adelic data, without using the location of the zeros of $F$, a norm
realization $(\mathcal K,\mathcal T)$ of $Q_W$
(Definition~\ref{def:norm-realization}) on a subspace containing
$\mathcal T_W$, preferably on all of $C_c^\infty(\R)$, in which the
positivity of the norm of $\mathcal K$ is independent of RH.
\end{question}

By Proposition~\ref{prop:weil-norm-converse} an answer to
Question~\ref{q:norm-target} implies RH; polarization then determines
the sesquilinear Weil pairing. The transform $\mathcal T_{\mathrm{rec}}$
of Question~\ref{q:reconstruction} is one candidate; the literal
transform of Proposition~\ref{prop:R17} is not. This target concerns
zero location only and need not resolve multiplicities.

\begin{question}\label{q:full-comparison}
Let $\Xbar_\Q$ carry a Weil datum satisfying
\ref{ax:AW1}--\ref{ax:AW4} for $(\Lambda,\mathbb 1)$ that also
satisfies
\ref{ax:AW5}, with completion $\cH^1$ of $(V^1,h_1)$. Is
there an analytic model Hilbert space $\mathcal M$ carrying the
operator $M_{\mathrm{an}}$ of multiplication by $z$, and a unitary
$U:\cH^1\to\mathcal M$ matching the cup-product and star structures,
such that on complete domains
\begin{equation}\label{eq:full-comparison}
U\bigl[-i(\Theta_1-\tfrac12)\bigr]U^{-1}=-M_{\mathrm{an}},
\end{equation}
and such that $U$ preserves spectral multiplicities, trace
distributions, and the normalization of the determinant in
\eqref{eq:aw-det}?
\end{question}

By \ref{ax:AW5}, $-i(\Theta_1-\tfrac12)$ is the self-adjoint operator
$T_1$; the minus sign in \eqref{eq:full-comparison} reflects the
correspondence $\rho=\tfrac12+i\tau\leftrightarrow z=-\tau$ of
Definition~\ref{def:zeta}. If some zero of $F$ is multiple, the scalar
model $\mathcal K(\vartheta)$ of Theorem~\ref{thm:imp-sz23} cannot serve
as $\mathcal M$, since its operator has simple spectrum
(Proposition~\ref{prop:R9}); one must then prove simplicity separately,
weaken the conclusion to zero location
(Question~\ref{q:norm-target}), or construct and justify additional
multiplicity data. An isomorphism of uncompleted vector spaces, a
comparison valid only after assuming RH, or an agreement of
integer-twist cones (Proposition~\ref{prop:R3}) does not give
\eqref{eq:full-comparison}.

\section{Open obligations}\label{sec:open}

The following questions are what remains between the conditional
Theorem~\ref{thm:rh} and an unconditional statement, together with one
question on the trace-formula axiom (Question~\ref{q:trace-axiom}).
They are listed in a useful order of attack; none is resolved in these
notes.

\begin{question}\label{q:obl-object}
Define a functor on $\Schbar_1$
(Definition~\ref{def:compactified-schemes}), or on
the single object $\Xbar_\Q$, together with its complex coefficient
structure, spectral core, and completion, without presupposing its
spectrum. Does it pass the Galois-functoriality test of
Remark~\ref{rem:R5-import}, that is, is it genuinely complex or
quaternionic rather than the complexification of a real theory?
\end{question}

\begin{question}\label{q:obl-trace}
Is there an arithmetic object whose trace formula or norm identity is
exactly the explicit formula, with prime terms, gamma terms, and
normalization derived from the object? In particular, is there an
independently defined transform satisfying the norm identity of
Question~\ref{q:norm-target} for every test function
(Proposition~\ref{prop:weil-norm-converse})?
\end{question}

The trace-formula axiom for $\Xbar_K$ is
Definition~\ref{def:aw6prime}; by Theorem~\ref{thm:aw6-div} it
determines the degree-one spectrum on
every tail class of test functions, by Theorem~\ref{thm:window} not on
a one-sided bounded window, and by Proposition~\ref{prop:window-two} a
two-sided bounded window detects every spectrum added on the critical
line. What remains is the following.

\begin{question}\label{q:trace-axiom}
Let $K$ be a number field such that every zero of $\xi_K$ has real part
$\tfrac12$, let $L>0$, and let $\nu=\sum_\rho\ord_\rho(\xi_K)
\,\delta_{\Im\rho}$. Is there a measure $\mu=\sum_\gamma m(\gamma)
\delta_\gamma\ne\nu$ on $\R$, with integer weights $m(\gamma)\ge0$,
$m(-\gamma)=m(\gamma)$, $m(0)$ even, and polynomially bounded counting
function, such that the Fourier transform of $\mu-\nu$ vanishes on
$(-L,L)$ in the sense of distributions? Equivalently: is there a Weil
datum on $\Xbar_K$ satisfying \ref{ax:AW1}, \ref{ax:AW3},
\ref{ax:AW4}, \ref{ax:AW5}, and AW6$'$ on $\mathcal S_{(-L,L)}$ whose
degree-one spectrum is not $\divisor\xi_K$?
\end{question}

By Proposition~\ref{prop:window-two} such a $\mu$ cannot satisfy
$\mu\ge\nu$: it must omit some zeros. The trivial model with spectrum
$\{\tfrac12+i\gamma\}$, weighted by $m$, gives the equivalence, as in
the proof of Theorem~\ref{thm:window}. By analogy with the extension
problem for positive definite functions on an interval, which is in
general not uniquely solvable, one expects the answer to be yes; no
such $\mu$ is constructed here.

\begin{question}\label{q:obl-positivity}
Can the positive compatible form of \ref{ax:AW4}, or the global
invertibility and endpoint condition of Theorem~\ref{thm:imp-sz16}, be
established? Algebraic duality, a positive Hamiltonian on an initial
interval, and semilocal compression do not suffice
(Propositions~\ref{prop:R2}, \ref{prop:R11}, and~\ref{prop:R14}).
\end{question}

\begin{question}\label{q:obl-multiplicity}
Is the goal RH alone, RH together with simplicity of the zeros, or the
full determinant identity \eqref{eq:aw-det} with zero orders? For the
last, construct the multiplicity-preserving object required by
Question~\ref{q:full-comparison}; a scalar de Branges model cannot
supply it when a zero is multiple (Proposition~\ref{prop:R9}).
\end{question}

\begin{question}\label{q:obl-descent}
Formulate an actual all-place descent: objects, overlap maps, topology,
and derived gluing linking the finite-place realizations of
Theorems~\ref{thm:imp-h17}, \ref{thm:imp-bms1}, \ref{thm:imp-bms2},
and~\ref{thm:imp-bs} with the archimedean realization of
Theorem~\ref{thm:imp-cc14}. Determinant formulas alone are not descent
data (Remark~\ref{rem:N-composition}; cf.\
Proposition~\ref{prop:R7}).
\end{question}

\begin{question}\label{q:obl-special-values}
Which additional information (analytically controlled dependence on the
twist, determinant metrics, a resolvent-level comparison) allows the
fixed-twist comparisons of Conjecture~\ref{conj:imp-fm20} to recover a
global generator? The cones alone cannot (Proposition~\ref{prop:R3}).
\end{question}

The remaining open items elsewhere in the notes are
Question~\ref{q:reconstruction} (a repaired transform in
Suzuki's framework) and Questions~\ref{q:norm-target}
and~\ref{q:full-comparison} (the two comparison targets). The class of
zeta functions is fixed in
Section~\ref{sec:class}; the self-contained account of Weil's positivity
for curves is Section~\ref{subsec:weil-curves}. Beyond the class
$\cC_{\mathrm{Ded}}$, the obligations above have a class-level form: a
functor on $\Schbar_1$ with Galois descent (Definition~\ref{def:descent}),
whose construction would also yield Artin's conjecture for rational
characters (Proposition~\ref{prop:descent}).
The K\"unneth conjecture (Conjecture~\ref{conj:kunneth}) is a further
open item: it asks for the arithmetic surface on which the Hodge index
clause of Theorem~\ref{thm:kun-hodge} could be proved.

\section{Software and AI assistance}
\label{sec:ai}

\subsection{Methods: typesetting}
The manuscript was assembled from AI research reports and a Lean
formalization, then revised into a survey of the previous work described
in Section~\ref{sec:prior}. The drafting passes of 25--26~September~2026
used Anthropic's Claude Opus~5.5, model identifier
\nolinkurl{claude-opus-5-5[1m]}, in Claude Code~2.1.282
\cite[\href{https://claude.com/product/claude-code}{product description}]%
{AnthropicClaudeCode,AnthropicClaudeOpus}, with the \nolinkurl{using-latexpowers} router and
\nolinkurl{harvest-dumps} workflow. They selected and organized the reports
into definitions, statements, and proofs, checked exact constants and
source locators, and revised the text after independent agent reviews.
The starting reports were AI-generated; the research questions and the
request to prepare these notes came from Taylor Dupuy. The human author
has only thumbed through the manuscript, as stated in the abstract.

\subsection{Methods: agent management}
The repository has draft contracts describing its agent entry point,
file layout, canonical \TeX{} files, writing conventions, prior decisions,
and AI disclosure. Those proposals remain unadopted. The recorded
synthesis used \nolinkurl{using-latexpowers} and \nolinkurl{math-workflow} to
organize the writing, \nolinkurl{create-exposition} and
\nolinkurl{rapid-prototype} to prepare its plan, \nolinkurl{fill-in-prototype}
for drafting, and \nolinkurl{contradiction-check} to compare conflicting
claims. A disagreement about the determinant normalization was recorded
and investigated rather than silently resolved by choosing one report.
The models identify the origin of the drafts; they are neither authors
nor mathematical authorities.

This manuscript has not been checked by a human. Readers use it at their
own risk and are responsible for independently verifying any claim,
proof, computation, or citation they rely on.


The workflow instructions named here are credited to Mathcity and its
mirror of the author's agent skills
\cite[\href{https://github.com/tdupu/mathcity/tree/0205e0adfb8963a843c28616fcedceafc7ca5680/subdomains/latex}
{LaTeX workflows}]{Dupuy2026mathcity},
\cite[\href{https://github.com/tdupu/mathcity/blob/c87fb990708ba1e8493133c6a8f17877004444c1/docs/agent-skills-mirror.md}
{workflow mirror}]{Dupuy2026skills}.
These are later reference snapshots, not evidence of the versions
executed in September. The exact historical versions and public
retrievability of all corresponding pinned files remain incompletely
verified; the procedural descriptions here do not require those files.

\subsection{Methods: record keeping}
The project retains provenance and token records and the dated research,
review, and verification packages supporting them. The provenance record
states the prompts, requested and observed models, harnesses, work
performed, derivations, searches, and negative findings;
\nolinkurl{update-ai-usage} is its writer. The token record, maintained by
\nolinkurl{update-tokens}, distinguishes measured counts, estimates,
unavailable counts, and prices. The available counts are tabulated in
Appendix~\ref{app:tokens}.
Missing runtime or token information is
reported as missing. In particular, the requested model for the report
of 11~September~2026 was \nolinkurl{gpt-6-astra}, but the actual runtime model
and effort were not exposed; that request does not establish what ran.
Some citation helpers in Claude Code were recorded only as
\nolinkurl{sonnet}; their exact model version was not retained.

\subsection{Methods: natural language proofs}
Research questions were sent through \nolinkurl{frontier-dump}, with the
report distinguishing conditional arguments, counterexamples,
computations, and open obligations. The report of 19~September~2026 used
OpenAI's \nolinkurl{gpt-6-astra} in the Codex desktop harness at reasoning
effort \nolinkurl{ultra}
\cite[\href{https://openai.com/codex/}{product description}]{OpenAICodex,OpenAIGPT6Astra};
the earlier report of 11~September requested the same model and effort,
with the runtime gap just stated. The reports on zeta-function axioms,
finite-field calibration, the arithmetic K\"unneth conjecture,
determinant normalization, and the trace-formula axiom on
25~September~2026, and on the Eilenberg--Moore comparison on
26~September~2026, used Claude Fable~5.1, model identifier
\nolinkurl{claude-fable-5-1}, in Claude Code~2.1.282
\cite[\href{https://claude.com/product/claude-code}{product description}]%
{AnthropicClaudeCode,AnthropicClaudeFable}. Their reasoning setting was not exposed.
The exposition re-derived the proofs and corrected claims in those
reports; a research report was an input, never mathematical warrant.
The four-question prompt, supplied proposal, research report, and caller
validation for the September search are retained in
\nolinkurl{ai/2026-09-19-arithmetic-weil-cohomology/}.
The preceding investigation is retained in
\nolinkurl{ai/2026-09-11-prove-riemann-hypothesis/}; its workflow was then
named \nolinkurl{astra-dump}. The source-search records of
26~September are in the two \nolinkurl{provenance-a} and
\nolinkurl{provenance-b} packages under \nolinkurl{ai/}.
These records distinguish supplied ideas, generated candidates,
source checks, and later corrections. They support the account in
Section~\ref{subsec:search-findings}; they are not measurements of a
constructed arithmetic cohomology or an experiment establishing RH.


\subsection{Methods: adversarial reviews}
\begingroup\emergencystretch=1em
The recorded drafting process used independent agents with
\nolinkurl{fp-finder-latex} and \nolinkurl{critical-review} for outline,
presentation, and integrated-manuscript reviews, followed by
\nolinkurl{referee-report} on successive completed drafts and revisions
through \nolinkurl{revise}. Some reviewers shared the writer's model family.
The review records distinguish those passes from self-checks, list
findings and closure tests, and record what was corrected or deferred.
These are simulated reviews by agents; they establish no human review or
acceptance of the mathematics.
\par\endgroup

\subsection{Methods: AI assisted computations}
The research and review packages contain rerunnable checks with exact
rational arithmetic and numerical diagnostics. Recorded systems include
Python~3.12.14 and~3.13.7
\cite[\href{https://www.python.org/about/}{project description}]
{Python2026python}, mpmath~1.3.0
\cite[\href{https://mpmath.readthedocs.io/en/1.3.0/}{documentation}]
{Mpmath2026mpmath}, SymPy~1.13.2
\cite[\href{https://doi.org/10.7717/peerj-cs.103}{software description}]
{Meurer2017sympy}, and NumPy~2.3.2
\cite[pp.~357--362]{Harris2020array}; several runs used the Python
runtime from SageMath~10.8
\cite[\href{https://www.sagemath.org/}{project website}]{Sage2026sagemath}. They checked
finite-dimensional witnesses, finite-field point counts, determinant
constants, explicit-formula identities, and signs in the de Branges
examples. These numerical diagnostics are not interval certificates and
do not support an existence claim for the conjectural cohomology. Source
checks opened the cited papers when obtainable; an inaccessible original
is cited together with an opened source reporting it. The manuscript was
built with \nolinkurl{latexmk}
\cite[\href{https://ctan.org/pkg/latexmk}{package description}]
{Collins2026latexmk}; the historical build version is not recorded.

\subsection{Methods: AI statement}
This statement was revised by an OpenAI GPT-6 agent in the Codex desktop
harness on 14~October~2026
\cite[\href{https://openai.com/codex/}{product description}]{OpenAICodex}.
Its exact runtime model identifier and reasoning effort were not exposed
to this drafting agent. The \nolinkurl{write-materials-and-methods} and
\nolinkurl{latex-ai-statement} instructions were consulted to prepare this
revision; the final disclosure was also checked against the recorded trail.
The structure follows the software statement of the author's PCF manuscript.
The requested models were \nolinkurl{gpt-6-astra} at reasoning effort
\nolinkurl{high} for the source and mathematical review,
\nolinkurl{gpt-6.1-sol} at effort \nolinkurl{high} for the disclosure comparison,
and \nolinkurl{gpt-6-luna} at effort \nolinkurl{medium} for the independent
\nolinkurl{check-paragraph} reviewers, all in Codex desktop.
Their exact runtime identifiers were not independently reported.
The continuation of this revision reconstructed the literature-search
purpose from the retained reports and made the source category a
parameter of the common spectral formalism. It requested
\nolinkurl{gpt-6-astra} at \nolinkurl{high} for source and mathematical
review, \nolinkurl{gpt-6.1-sol} at \nolinkurl{high} for archival
provenance and paragraph coordination, and \nolinkurl{gpt-6-luna} at
\nolinkurl{medium} for fresh reviews of changed paragraphs. Actual
runtime identifiers were not independently reported. This was an
archival reconstruction and manuscript revision, not a new
\nolinkurl{frontier-dump} investigation. The paragraph verdicts are
retained separately from the manuscript.

Human review and acceptance are not established by these agent passes.
Historical workflow versions were not recorded for several passes;
public, version-pinned sources for those usages have not been established.

On 15~October~2026 an OpenAI \nolinkurl{gpt-6-sol} agent at
\nolinkurl{xhigh} effort in the Codex desktop harness
\cite[\href{https://developers.openai.com/api/docs/models/gpt-6-sol}
{model description}]{OpenAIGPT6Sol}
drafted Definition~\ref{def:comparison-probes},
Conjecture~\ref{conj:fm-ccm-common}, and the revised opening
\cite[\href{https://openai.com/codex/}{product description}]%
{OpenAICodex}. These are observed harness labels, not a provider
response identity. The inputs were the existing manuscript and the
AI-generated research report in
\nolinkurl{ai/2026-10-15-single-deninger-formalism/}.
The agent checked the relevant Flach--Morin and CCM source statements
and made two revisions after a \nolinkurl{gpt-6-luna} agent review;
another \nolinkurl{gpt-6-luna} agent audited the writing advice and
a third checked the revised opening
\cite[\href{https://developers.openai.com/api/docs/models/gpt-6-luna}
{model description}]{OpenAIGPT6Luna}.
The sources motivate the two
separate comparisons but do not state their common-realization
conjecture. The bounded literature check found no existing combined
comparison; it was not exhaustive. These agent reviews are not human
verification, and neither comparison has been proved.

On 17~October~2026, a Codex desktop \nolinkurl{gpt-6-sol} agent at
\nolinkurl{xhigh} effort used \nolinkurl{harvest-dumps} to add
Subsection~\ref{subsec:prior-sc-bsd} and the corresponding negative
register entries. Its selected input was the AI-generated research
package \nolinkurl{ai/2026-10-17-standard-conjectures-bsd/}, prepared
by an observed \nolinkurl{gpt-6-astra} agent at \nolinkurl{ultra}
effort through \nolinkurl{frontier-dump}
\cite[\href{https://openai.com/codex/}{product description}]
{OpenAICodex},
\cite[\href{https://developers.openai.com/api/docs/models/gpt-6-sol}
{model description}]{OpenAIGPT6Sol}, and
\cite[\href{https://developers.openai.com/api/docs/models/gpt-6-astra}
{model description}]{OpenAIGPT6Astra}.
The workflow files came from a local \nolinkurl{agent-skills}
checkout at revision \nolinkurl{8ed41877};
\nolinkurl{harvest-dumps} had uncommitted local changes. A public
\cite[\href{https://github.com/tdupu/mathcity/blob/%
c87fb990708ba1e8493133c6a8f17877004444c1/%
docs/agent-skills-mirror.md}{workflow reference snapshot}]
{Dupuy2026skills} describes the workflow family, but the exact
executed files have no verified public snapshot.
The package
distinguished Grothendieck's geometric standard conjectures from
Tate's divisor conjecture and from arithmetic Abel--Jacobi and
special-value comparisons. The conditional rank argument and the
function-field implication are expositions of Deninger's proposal and
the cited arithmetic theorems. Independent Codex agents reviewed the
outline, candidate mathematics and presentation; these reviews and the
source audit are retained in the dated \nolinkurl{scratch/} and
\nolinkurl{ai/} packages. No human verification or mathematical
acceptance is established by these agent checks.

\subsection{Proofreading and adversarial review}
The Opus drafting passes and Opus or Fable reviewer passes of
25--26~September~2026 used the harnesses stated above. Reviews caught
errors in determinant normalization, Galois descent, tensor-product
hypotheses, the strength of the Hodge-index clause, the scope of a
compact-support argument, and citation locators. The later
Eilenberg--Moore harvest replaced a weight argument by a degree argument
and supplied a direct proof where a cochain enhancement had previously
been assumed without proof. No human referee has reviewed the manuscript.

\subsection{Literature review}
Source searches were run in the research packages and source retrieval
through \nolinkurl{track-down-reference} checked pinpoint locators. On
26~September~2026 two Claude Fable~5.1 agents in Claude Code, with delegated
subagents, searched for prior sources of the theorem-class statements.
They used arXiv, general web search, publisher and bibliographic databases,
and a local corpus of the cited literature. The Semantic Scholar index
was rate-limited and could not be queried systematically. The resulting
prior-work survey, provenance table, and attributions were written by
Claude Opus~5.5 in Claude Code. Source verification by these agents is not
verification by the human author. Some earlier packages had recorded
unperformed literature searches; those are gaps, not zero-hit searches.

\subsection{AI findings and procedure}
The inputs were the author's research questions, the published work
surveyed in Section~\ref{sec:prior}, and the earlier AI-generated reports.
The negative findings are retained in Section~\ref{sec:negative} and the
register of Section~\ref{subsec:neg-register}; they delimit the claims
made by the axioms and candidate constructions. The principal corrections
and contributions from the drafting and research passes are:
\begin{itemize}
\item The descent argument was restricted to what permutation characters
actually determine (Proposition~\ref{prop:descent}); a claimed determinant
normalization was replaced by the computation of
Theorem~\ref{thm:det-K}.
\item The tensor-square analysis acquired the zero-counting hypothesis
in Theorem~\ref{thm:kun-square}\ref{item:ks-acc}, the sign condition in
Proposition~\ref{prop:kun-diag}\ref{item:kd-semi}, and the one-way
implication distinguished in Proposition~\ref{prop:kun-hodge-strict}.
The unsupported compact-support injectivity claim was not used in
Proposition~\ref{prop:kun-open}.
\item The trace-formula report's prime-sum divergence was placed in the
spectral normalization in Proposition~\ref{prop:schwartz}; the derivation
after Theorem~\ref{thm:imp-explicit} uses a polynomial bound on horizontal
segments in place of the cited estimate.
\item The exposition gathers Deninger's conditional positivity argument,
the finite-field model, and the counterexamples, with their prior sources
and the method comparisons recorded in Subsection~\ref{subsec:provenance}.
\item The later cochain comparison was harvested into
Section~\ref{sec:kunneth}, with the corrections described under
``Proofreading and adversarial review.''
\end{itemize}
These drafting and research contributions were produced by the Opus,
Fable, and Codex passes dated and identified above. Their proofs and the
cited mathematical sources, rather than a model's assertions, supply
the arguments.

\subsection{Provenance}
\label{subsec:provenance}
Section~\ref{sec:prior} and Table~\ref{tab:provenance} give the sources and
relations for the statements retained in these notes. In particular,
the implication from a positive star compatible with the flow to the
Riemann hypothesis is Deninger's formal argument after Serre, rather
than a new discovery of the agents. The table distinguishes imported,
classical, adapted, and standard statements from those for which the
26~September search found no prior appearance. The last class asserts
only that search outcome, with the limits just stated; it asserts no
novelty. Some individual earlier proof-writing and Lean tasks recorded
no separate literature search, a limitation that the later statement
survey does not turn into an exhaustive novelty check.

Table~\ref{tab:provenance} lists the theorem-class statements of
Sections~\ref{sec:conventions}--\ref{sec:open} with the place where
each appears in the literature or where its idea comes from. The
classes are: \emph{imported}, a result stated with its source and
locator; \emph{classical}, a result in the literature, cited at the
statement; \emph{adapted}, a variant or reformulation of a result in
the literature; \emph{standard}, textbook material or folklore; and
\emph{not found}, a statement for which no prior appearance was found
in the literature we searched. The last class is not a claim of
novelty: the search used arXiv, general web search, and a local corpus
of the cited literature, the Semantic Scholar index could not be
queried systematically, and the statements concerned are mostly
explicit examples and elementary equivalences of a kind that may well
appear in print. Definitions are listed only where they carry an
attribution.

{\small
\begin{longtable}{@{}>{\raggedright\arraybackslash}p{0.2\textwidth}
>{\raggedright\arraybackslash}p{0.6\textwidth}
>{\raggedright\arraybackslash}p{0.12\textwidth}@{}}
\caption{Provenance of the statements of these notes.}
\label{tab:provenance}\\
\toprule
Statement & Where it appears, or where the idea comes from & Class\\
\midrule
\endfirsthead
\toprule
Statement & Where it appears, or where the idea comes from & Class\\
\midrule
\endhead
\bottomrule
\endfoot
\multicolumn{3}{@{}l}{\emph{Section~\ref{sec:conventions}}}\\
Def.~\ref{def:inf-pol}, Prop.~\ref{prop:purity} &
\cite[pp.~392--393]{Serre1960Analogues};
\cite[(7.11)--(7.13), pp.~32--33]{Deninger1994Motivic};
\cite[p.~169]{Deninger1998Analogies}; \cite[2.7, p.~3]{Deninger2022NoWeil}
& classical\\
Def.~\ref{def:spectral-core} &
\cite[Definition~2.1(i), p.~521]{FlachMorin2020DeningerWeilArakelov};
\cite[\S1]{Deninger1994Motivic}; the accumulation clause added &
adapted\\
Thms.~\ref{thm:imp-riemann}, \ref{thm:imp-hecke},
\ref{thm:imp-deligne} & sources cited at the statements & imported\\
Props.~\ref{prop:F0}, \ref{prop:xiK} & consequences of Hecke's theorem
and of $\eta(s)>0$ on $(0,1)$ & standard\\
Def.~\ref{def:weil-form}, Prop.~\ref{prop:weil-form} & Weil's
functional \cite[\S16]{Weil1972Explicit};
\cite[\S V, p.~9]{Bombieri2000Clay};
\cite[Theorem~1, p.~191]{Bombieri2000Lincei} & classical\\
\multicolumn{3}{@{}l}{\emph{Section~\ref{sec:finite-field}}}\\
Def.~\ref{def:frob-pol} & hypothesis of
\cite[Th\'eor\`eme~1]{Serre1960Analogues} & adapted\\
Props.~\ref{prop:weil-linear}, \ref{prop:algebra-form} &
\cite[p.~393]{Serre1960Analogues};
\cite[Lemma~4.6, p.~385]{Kleiman1968Cycles};
\cite[Proposition~3.10]{Milne2016RHFF} & classical\\
Props.~\ref{prop:weil-converse}, \ref{prop:N-jordan} & linear algebra
& standard\\
Prop.~\ref{prop:formal-curve} & that positivity is needed:
\cite[p.~360]{Kleiman1968Cycles}; the example itself & not found\\
Thms.~\ref{thm:imp-weil-coh}, \ref{thm:imp-rosati} & sources cited at
the statements & imported\\
Prop.~\ref{prop:weil-curves} & Weil \cite{Weil1941RHFunctionFields};
\cite[Th\'eor\`eme~2, p.~394]{Serre1960Analogues};
\cite[Lemma~4.6 and Theorem~4.7]{Kleiman1968Cycles};
\cite[Theorem~1.27, p.~21]{Milne2016RHFF} & classical\\
Prop.~\ref{prop:ff-N7} & \cite[p.~3]{Deninger2010HilbertPolya} &
standard\\
Defs.~\ref{def:suspension}, \ref{def:log-susp},
Prop.~\ref{prop:suspension}\ref{item:susp-zeros}--\ref{item:susp-core}
& \cite[\S2, Corollary~(2.6) and~(2.8), and \S4, p.~16]%
{Deninger1994Motivic}; \cite[Remark~7.2]{Hesselholt2016THHZeta} &
classical\\
Prop.~\ref{prop:suspension}\ref{item:susp-pol}--\ref{item:susp-line},
Prop.~\ref{prop:susp-converse} & \cite[p.~4]{Deninger2010HilbertPolya}
(integrated and infinitesimal star laws) & adapted\\
Prop.~\ref{prop:aw4-EF}\ref{item:EF-pure} &
\cite[Lemma~4.6]{Kleiman1968Cycles} & classical\\
Prop.~\ref{prop:aw4-EF}\ref{item:EF-sheet}--\ref{item:EF-add} &
ingredients \cite[Lemma~4.6]{Kleiman1968Cycles} & not found\\
Prop.~\ref{prop:rosati-EF} & Weil; \cite[\S14]{MilneAV2008};
\cite[Lemma~4.6]{Kleiman1968Cycles} & adapted\\
\multicolumn{3}{@{}l}{\emph{Section~\ref{sec:axioms}}}\\
Defs.~\ref{def:regdet}, \ref{def:principal-cut} &
\cite[(1.1)--(1.7), pp.~3--6]{Deninger1994Motivic};
\cite[Definition~2.1]{FlachMorin2020DeningerWeilArakelov} & imported\\
Def.~\ref{def:unit} & \cite[pp.~167--168]{Deninger1998Analogies} &
classical\\
Defs.~\ref{def:aw}, \ref{def:aw4-forms}, \ref{def:aw-det} &
\cite[Conjecture~2, p.~3, and pp.~4--6]{Deninger2010HilbertPolya};
\cite[Conjecture~2.3, p.~522]{FlachMorin2020DeningerWeilArakelov};
\cite[p.~168, p.~172]{Deninger1998Analogies};
\cite[(7.13), (7.19)]{Deninger1994Motivic};
\cite[Lemma~4.6]{Kleiman1968Cycles} & adapted\\
Prop.~\ref{prop:weighted-algebra-obstruction} &
derivation, trace, and positivity in the former weighted internal
axioms; proved here & adapted\\
Defs.~\ref{def:test-classes}, \ref{def:geom-side}, \ref{def:aw6prime}
& Weil, as reported in \cite[\S V]{Bombieri2000Clay};
\cite[(4) and (10), pp.~170--172]{Deninger1998Analogies};
\cite[(17)--(19), pp.~15--16]{Deninger2001NumberTheory} & adapted\\
\multicolumn{3}{@{}l}{\emph{Section~\ref{sec:purity}}}\\
Prop.~\ref{prop:polarization-identity} &
\cite[p.~5]{Deninger2010HilbertPolya};
\cite[(7.12)--(7.13)]{Deninger1994Motivic} & classical\\
Prop.~\ref{prop:aw4-forms}\ref{item:aw4e-pure}--\ref{item:aw4v-e} &
\cite[Lemma~4.6]{Kleiman1968Cycles} in additive form;
\cite[(7.13)]{Deninger1994Motivic} & classical\\
Prop.~\ref{prop:aw4-forms}\ref{item:aw4v-star} & linear algebra &
not found\\
Thm.~\ref{thm:rh} & Deninger's formal argument
\cite[(7.11)--(7.13)]{Deninger1994Motivic},
\cite[pp.~168--169]{Deninger1998Analogies},
\cite[pp.~4--5]{Deninger2010HilbertPolya},
\cite[2.7]{Deninger2022NoWeil},
\cite[Conjecture~2.3(iv)--(v)]{FlachMorin2020DeningerWeilArakelov},
after \cite{Serre1960Analogues} & classical\\
Prop.~\ref{prop:hilbert} & spectral theorem;
\cite[p.~33]{Deninger1994Motivic} & standard\\
Prop.~\ref{prop:converse}, Rem.~\ref{rem:reformulation} &
Hilbert--P\'olya: \cite[p.~3]{Connes1998TraceFormula},
\cite[p.~147]{Deninger1992Local}, \cite[p.~5]{Sarnak2004Clay} &
standard\\
Prop.~\ref{prop:R1} & indefinite forms;
\cite[pp.~7--9]{Deninger2010HilbertPolya} & standard\\
Props.~\ref{prop:R2}, \ref{prop:R2-star} & mechanism
\cite[2.7]{Deninger2022NoWeil}; the example itself & not found\\
\multicolumn{3}{@{}l}{\emph{Section~\ref{sec:class}}}\\
Thms.~\ref{thm:imp-artin}, \ref{thm:imp-fs} & sources cited at the
statements & imported\\
Thm.~\ref{thm:A} & (i), (iv) as Thm.~\ref{thm:rh}; (ii), (iii) Artin
formalism & classical\\
Props.~\ref{prop:trivial-model}, \ref{prop:content} &
\cite[2.6--2.9, pp.~3--4]{Deninger2022NoWeil};
\cite[Proposition~3.1]{Deninger2025BSDDedekind} & adapted\\
Prop.~\ref{prop:N-central-order} & \cite[2.9 and
Theorem~2.13]{Deninger2022NoWeil} & adapted\\
Def.~\ref{def:descent} & \cite[p.~4]{Deninger2022NoWeil};
\cite[\S\S2.8--2.10]{Deninger2025BSDDedekind} & adapted\\
Prop.~\ref{prop:perm-span} & Artin;
\cite[\S13.1, Theorems~29--30]{Serre1977LinearReps} & classical\\
Prop.~\ref{prop:descent} & ingredients
\cite[\S13.1, p.~104]{Serre1977LinearReps},
\cite[2.10--2.12]{Deninger2022NoWeil} & not found\\
Prop.~\ref{prop:class-converse} &
\cite[Theorem~2.15]{Deninger2025BSDDedekind} & adapted\\
Prop.~\ref{prop:real-coeff} & \cite[Theorem~2.13]{Deninger2022NoWeil}
& imported\\
Prop.~\ref{prop:root-parity} & \cite[2.7--2.8]{Deninger2022NoWeil};
\cite[Proposition~3.1]{Deninger2025BSDDedekind} & adapted\\
Prop.~\ref{prop:susp-product}, Def.~\ref{def:susp-datum} &
\cite[\S4, pp.~16--17]{Deninger1994Motivic} & classical\\
Thms.~\ref{thm:imp-hurwitz}, \ref{thm:imp-stirling} & sources cited
at the statements & imported\\
Thm.~\ref{thm:ff-det}, Prop.~\ref{prop:ff-cut} &
\cite[(2.7)--(2.10), (4.1)]{Deninger1994Motivic};
\cite[Proposition~7.1]{Hesselholt2016THHZeta} & classical\\
Props.~\ref{prop:one-branch}, \ref{prop:ff-phantom} & direct
computation & standard\\
Prop.~\ref{prop:unit} & \cite[pp.~167--168, Proposition~3.1]%
{Deninger1998Analogies} & classical\\
Thm.~\ref{thm:B} & (i)--(iii) \cite[\S4, (4.1)]{Deninger1994Motivic},
\cite[Theorem~3.3, p.~172]{Deninger1998Analogies}; (iv)--(vi) Weil,
\cite[Lemma~4.6]{Kleiman1968Cycles}, \cite[\S3]{Milne2016RHFF} &
classical\\
Prop.~\ref{prop:P1Z} & \cite[Conjecture~1(c), p.~2]%
{Deninger2010HilbertPolya} & classical\\
Thm.~\ref{thm:A-schemes} & (i) \cite[p.~172]{Deninger1998Analogies},
\cite[p.~5]{Deninger2010HilbertPolya}; (ii)--(iii) stated here &
adapted\\
Prop.~\ref{prop:no-positivity-cancel} & the example itself & not
found\\
Rem.~\ref{rem:A-motivic} & \cite[pp.~6--7]{Deninger2010HilbertPolya};
\cite[(7.11)--(7.13)]{Deninger1994Motivic} & adapted\\
Thm.~\ref{thm:imp-deninger-product} &
\cite[Theorem~3.3 and Remark~3.4(1)]{Deninger1992Local};
Soul\'e; \cite[p.~24]{Kurokawa1988Parabolic} & imported\\
Thm.~\ref{thm:det-K} & \cite[Theorem~7.18(ii)]{Deninger1994Motivic}
(Soul\'e, Illies); \cite[p.~24]{Kurokawa1988Parabolic}; the constant
relative to $\Lambda_K$ not found written out & classical\\
Thm.~\ref{thm:det-K-rf}, Prop.~\ref{prop:Zpm} &
\cite[Lemma~(1.7), Theorem~(7.20), and (7.19.1)]{Deninger1994Motivic};
the explicit residues not found & adapted\\
Prop.~\ref{prop:N-count} & compare
\cite[Remark~(3.4)(2)]{Deninger1992Local} & not found\\
Props.~\ref{prop:xiK-asymp}, \ref{prop:beta}, \ref{prop:exp-unique},
\ref{prop:aw5-contains}, \ref{prop:schwartz} & Stirling, beta
integral, uniqueness of Dirichlet series, Stone's theorem,
$\sum_p1/p=\infty$ & standard\\
Thm.~\ref{thm:imp-explicit} & Weil, as reported in
\cite[\S V]{Bombieri2000Clay};
\cite[Theorem~4.1]{Deninger2001NumberTheory};
\cite{Barner1981Explicit} & imported\\
Thm.~\ref{thm:aw6-div} & trace formula as an expectation:
\cite[(10)]{Deninger1998Analogies},
\cite[(19)]{Deninger2001NumberTheory}; the equivalence & not found\\
Thm.~\ref{thm:ff-lefschetz} & (i) Grothendieck; (ii)
\cite[(2), p.~3]{Deninger2002Explicit},
\cite[(10)]{Deninger1998Analogies}; (iii) & classical\\
Thm.~\ref{thm:aw6-hilb} & \cite[p.~33]{Deninger1994Motivic},
\cite[p.~169]{Deninger1998Analogies} & adapted\\
Prop.~\ref{prop:aw6-weil-pos} & Weil;
\cite[\S V, p.~9]{Bombieri2000Clay} & classical\\
Thm.~\ref{thm:window}, Prop.~\ref{prop:window-two} & the examples
themselves & not found\\
\multicolumn{3}{@{}l}{\emph{Section~\ref{sec:kunneth}}}\\
Props.~\ref{prop:tensor-eigen}, \ref{prop:B-support},
\ref{prop:kun-torus}, \ref{prop:kun-open}, \ref{prop:EM-product},
\ref{prop:kun-tensor} & linear algebra, Kronecker sums, Euler
characteristics, the K\"unneth formula over a field & standard\\
Def.~\ref{def:surface-datum} & \cite[\S1.2, p.~363, and
p.~360]{Kleiman1968Cycles} & adapted\\
Def.~\ref{def:tensor-datum} & \cite[\S1.6, p.~133]{Manin1995Lectures}
& adapted\\
Def.~\ref{def:kurokawa-cone} & \cite[p.~219]{Kurokawa1992MultipleZeta};
\cite[\S5.2, p.~158]{Manin1995Lectures} & adapted\\
Prop.~\ref{prop:kun-over-Z} & $\Spec\Z\times_{\Spec\Z}\Spec\Z=\Spec\Z$:
\cite[p.~121]{Manin1995Lectures}; \cite[p.~41]{Deninger1994Motivic}
& adapted\\
Props.~\ref{prop:kun-EM}, \ref{prop:EM-cochain}, \ref{prop:EM-x} &
ingredients \cite[pp.~22--23, 89--90, 105]{MilneLEC2013}; the examples
themselves & not found\\
Prop.~\ref{prop:EM-bar}, Rem.~\ref{rem:N-EM-scope} &
\cite[Theorem~7.6, p.~241]{McCleary2001UsersGuide} & adapted\\
Thm.~\ref{thm:kun-square}, Props.~\ref{prop:kun-AW6},
\ref{prop:kun-Fq-C} & \cite[p.~132]{Manin1995Lectures};
\cite[(7.32), p.~42]{Deninger1994Motivic} & adapted\\
Prop.~\ref{prop:kun-cone} & the cone of
\cite[p.~219]{Kurokawa1992MultipleZeta}; the statement & not found\\
Prop.~\ref{prop:kun-diag} & tensor sums of self-adjoint operators &
standard\\
Prop.~\ref{prop:kun-orbits} & \cite[pp.~167--168, p.~181]%
{Deninger1998Analogies}; \cite[p.~220]{Kurokawa1992MultipleZeta} &
adapted\\
Prop.~\ref{prop:kun-arch} & \cite[\S3, (25)]{Serre1970Facteurs};
\cite[VII.12, p.~535]{Neukirch1999ANT} & classical\\
Rem.~\ref{rem:kun-real} & \cite[p.~184]{Deninger1998Analogies}; Serre
\cite[fn.~25, p.~27]{Milne2016RHFF} & imported\\
Conj.~\ref{conj:kunneth} & \cite[(7.32)]{Deninger1994Motivic};
\cite[\S3]{ConnesConsani2018RiemannRoch};
\cite[\S6]{Scholze2018PadicGeometry}; the form stated & not found\\
Thm.~\ref{thm:kun-Fq} & sources cited in the proof & adapted\\
Def.~\ref{def:EM-setting}, Props.~\ref{prop:EM},
\ref{prop:EM-criterion} & \cite[Definition~1.1.7, Proposition~1.1.9,
Corollary~1.1.10]{Lurie2011DAGXIII}; Dwyer & adapted\\
Prop.~\ref{prop:EM-fibre} & \stackstag{0F0C}, \stackstag{0F0G} &
imported\\
Prop.~\ref{prop:kun-susp-tor} & bar resolution
\cite[X.2, Theorem~2.1]{MacLane1963Homology},
\cite[8.6.12]{Weibel1994Homological} & adapted\\
Prop.~\ref{prop:kun-tate} & \cite[Proposition~2.2.3]{Kahn2018Fonctions}
& classical\\
Thm.~\ref{thm:kun-hodge} & \cite[3.11]{Kleiman1968Cycles};
\cite[\S21, Theorem~1]{Mumford1974AV};
\cite[Th\'eor\`eme~2]{Serre1960Analogues} & adapted\\
Props.~\ref{prop:kun-hodge-strict}, \ref{prop:kun-collapse} &
compare \cite[pp.~4--5]{Deninger2010HilbertPolya} & not found\\
\multicolumn{3}{@{}l}{\emph{Section~\ref{sec:negative}}}\\
Props.~\ref{prop:R2-frob}, \ref{prop:R11}, \ref{prop:R16} & compare
\cite[Lemma~2.1]{Lagarias2005Differenced},
\cite[Theorem~2.3]{Suzuki2012Canonical},
\cite[Theorem~3.9, p.~63]{Saad2011Eigenvalue} & not found\\
Prop.~\ref{prop:R3} & \cite[\S1, pp.~246--247]{Deninger1991Gamma};
\cite[(2), p.~215]{QuineHeydariSong1993Zeta} & adapted\\
Props.~\ref{prop:N-lines}, \ref{prop:R4}, \ref{prop:R10},
\ref{prop:R7} & Hasse's bound; closed symmetric operators; linear
algebra; \cite[(1.3), (2.3)]{Deninger1992Local} & standard\\
Prop.~\ref{prop:R9} & \cite[p.~4 and Theorem~1(ii)]%
{Lagarias2006HilbertEntire} & adapted\\
Prop.~\ref{prop:R12} & \cite[\S7(b), p.~223]{DeBruijn1950Roots} &
standard\\
Prop.~\ref{prop:R14} & \cite[Theorem~1, p.~310]{Yoshida1992Hermitian};
\cite[Theorem~12]{Bombieri2000Lincei};
\cite{Burnol2000Explicites} & standard\\
Prop.~\ref{prop:R5} & Serre \cite[(2.2), p.~27]{Milne2016RHFF};
\cite[\S1, p.~1]{Deninger2022NoWeil} & classical\\
Rems.~\ref{rem:R5-import}, \ref{rem:R6}, \ref{rem:N-composition},
\ref{rem:R8}, \ref{rem:R13} & sources cited at the remarks &
imported\\
Prop.~\ref{prop:R15} & \cite[p.~4]{Riemann1859Wilkins};
\cite[(7.2)--(7.3), p.~220]{DeBruijn1950Roots} & classical\\
Prop.~\ref{prop:R17}, Rems.~\ref{rem:R17-source}, \ref{rem:R17-rh} &
about version~3 of \cite{Suzuki2023WeilHilbert}; corrected in
\cite[p.~27]{Suzuki2026WeilScrew} & not found\\
\multicolumn{3}{@{}l}{\emph{Section~\ref{sec:imported}}}\\
all statements & sources cited at the statements & imported\\
\multicolumn{3}{@{}l}{\emph{Section~\ref{sec:debranges}}}\\
Prop.~\ref{prop:hb}, Rem.~\ref{rem:hb} &
\cite[Theorem~1 and (10), p.~9]{Lagarias2006HilbertEntire} &
classical\\
Prop.~\ref{prop:hb-high} & \cite[Lemma~2.1]{Lagarias2005Differenced}
& standard\\
Prop.~\ref{prop:canonical} & de Branges's theorem;
\cite[(12)--(13) and Proposition~4]{Romanov2014Canonical} &
classical\\
Prop.~\ref{prop:weil-crit-notes} & Weil;
\cite[Theorem~1]{Bombieri2000Lincei};
\cite[Proposition~C.1]{ConnesConsani2020WeilPositivity} & classical\\
Prop.~\ref{prop:weil-norm-converse} & the ``if'' direction of
\cite[Theorems~1.4 and~4.4]{Suzuki2023WeilHilbert} & adapted\\
Question~\ref{q:reconstruction} & the definition of
\cite[p.~27]{Suzuki2026WeilScrew} & adapted\\
Props.~\ref{prop:gap}, \ref{prop:rouche-route} & residual bounds
\cite[Theorem~3.9]{Saad2011Eigenvalue}; Hurwitz's theorem & standard\\
\end{longtable}
}

The statements for which no prior appearance was found in the
literature we searched are Propositions~\ref{prop:formal-curve},
\ref{prop:aw4-EF}\ref{item:EF-sheet}--\ref{item:EF-add},
\ref{prop:aw4-forms}\ref{item:aw4v-star}, \ref{prop:R2},
\ref{prop:R2-star}, \ref{prop:descent},
\ref{prop:no-positivity-cancel}, \ref{prop:N-count},
\ref{prop:window-two}, \ref{prop:kun-EM}, \ref{prop:EM-cochain},
\ref{prop:kun-cone}, \ref{prop:EM-x}, \ref{prop:kun-hodge-strict},
\ref{prop:kun-collapse}, \ref{prop:R2-frob}, \ref{prop:R11},
\ref{prop:R16}, and~\ref{prop:R17}, Theorems~\ref{thm:aw6-div}
and~\ref{thm:window}, Conjecture~\ref{conj:kunneth} in the form
stated, and Remarks~\ref{rem:R17-source} and~\ref{rem:R17-rh}.

\subsection{Due diligence}
The finite-field package tested the axioms against their classical
model; the K\"unneth package tested stronger product expectations; the
determinant package computed the normalization; the trace-formula
package tested how much of the spectrum its test-function class detects;
and the cochain package investigated the Eilenberg--Moore comparison.
Their corrections and counterexamples are retained with the positive
statements. The recorded computations and agent reviews leave the
geometric and analytic existence obligations open.

On 15~October~2026, an OpenAI \nolinkurl{gpt-6-astra} agent at reasoning
effort \nolinkurl{max}, using the Codex desktop harness and the
\nolinkurl{frontier-dump} workflow,
\cite[\href{https://openai.com/codex/}{product description}]%
{OpenAICodex,OpenAIGPT6Astra}, searched for a counterexample to
Flach--Morin's Conjecture~\ref{conj:imp-fm20}. It checked the
finite-field point and the published finite-field theorem and found
no verified counterexample to the full conjecture in that bounded
search. Proposition~\ref{prop:R3} records the narrower failure of
integer-twist cone data to determine a zero spectrum; it is not a
counterexample to the conjecture. The search and its source checks
are retained in
\nolinkurl{ai/2026-10-15-flach-morin-cyclic-assembly/}.

The same research pass attempted to rule out any assembly of
finite-prime and archimedean local theories. It established the
specific direct-sum obstruction in Proposition~\ref{prop:R7}, but
found no universal impossibility theorem for constructions with
global relations, quotients, or other comparison maps. The class of
all possible assemblies was not specified, so the unsuccessful
search is not evidence that an assembly exists. This revision was
integrated by an OpenAI \nolinkurl{gpt-6-sol} agent at reasoning effort
\nolinkurl{xhigh} in Codex desktop on the same date
\cite[\href{https://openai.com/codex/}{product description}]{OpenAICodex}.
The mathematical warrants are the displayed arguments and the cited
primary papers, not the agents' judgments.

A second \nolinkurl{frontier-dump} pass on 15~October~2026 used an
OpenAI \nolinkurl{gpt-6-astra} agent at \nolinkurl{max} effort in the
Codex desktop harness \cite{OpenAICodex,OpenAIGPT6Astra}. It tested
the proposed adelic positivity proof, local-to-global comparison
routes, and Scholze's analytic-ring example. It obtained neither
global positivity nor an arithmetic spectral realization. The
compact-support argument in Proposition~\ref{prop:arch-window-limit}
and the fixed-radius argument in
Proposition~\ref{prop:scholze-fixed-radius} are the narrower negative
results incorporated here. Its primary-source checks and finite
verification are recorded in
\nolinkurl{ai/2026-10-15-adelic-positivity-prismatic-scholze/}.
The present manuscript integration was made by an OpenAI
\nolinkurl{gpt-6-sol} agent at \nolinkurl{xhigh} effort in Codex
desktop. The mathematical warrants remain the proofs and the cited
sources.

A further \nolinkurl{frontier-dump} pass on 15~October~2026 used an
OpenAI \nolinkurl{gpt-6-astra} agent at \nolinkurl{max} effort in
Codex desktop
\cite[\href{https://openai.com/codex/}{product description}]%
{OpenAICodex,OpenAIGPT6Astra}. It checked two function-field trace
examples and searched for a comparison between Flach--Morin's
integer-twist complexes and the Connes--Consani--Marcolli adelic
cyclic quotient. The examples are calculations in known
function-field cases; they do not establish adelic positivity over
$\Q$. The searched sources supplied a common Frobenius spectrum in
the function-field setting, but no verified canonical map between
the two constructions. The methods differ: one uses motivic
regulators and fixed integer twists, while the other uses an adelic
convolution action on cyclic homology. The bounded search and
calculations are retained in
\nolinkurl{ai/2026-10-15-ccm-two-examples-flach-morin/}.
Theorem~\ref{thm:ccm-positivity-rh} is the direct specialization of
\cite[Corollary~6.3, p.~29]{ConnesConsaniMarcolli2009Weil} to the
trivial character of $\Q$; it proves no new positivity statement.
The manuscript integration was made by an OpenAI
\nolinkurl{gpt-6-sol} agent at \nolinkurl{xhigh} effort in Codex
desktop. Its warrant is the cited criterion and the displayed proof.

The terminology and scope remarks
\ref{rem:deninger-terminology} and~\ref{rem:ccm-deninger-scope}
were drafted on 15~October~2026 by an OpenAI
\nolinkurl{gpt-6-sol} agent at \nolinkurl{xhigh} effort in Codex
desktop \cite[\href{https://openai.com/codex/}{product description}]%
{OpenAICodex}. The pass compared Flach--Morin's conjectural
dynamical groups and fixed-twist Weil--Arakelov complexes with
Connes--Consani--Marcolli's cyclic trace realization. A bounded
search of those papers and targeted uniqueness queries found no
theorem characterizing a global Deninger cohomology theory uniquely;
this is a search limit, not a nonexistence result. The source checks
and query record are in
\nolinkurl{scratch/2026-10-15-deninger-terminology/report.md}.
The mathematical warrants are the cited papers and AW1--AW2, not
the drafting agent.

The later naming pass on 15~October~2026 was drafted by an OpenAI
\nolinkurl{gpt-6-sol} agent at \nolinkurl{xhigh} effort in Codex
desktop \cite[\href{https://openai.com/codex/}{product description}]%
{OpenAICodex}. It checked Flach--Morin's name and fixed-twist
comparison, CCM's relative $H^1$ construction, Hesselholt's
$\mathrm{TP}$ determinant, and Connes--Consani's ordinary cyclic
archimedean determinant against the cited primary papers. The
methods have different inputs: motivic regulators, an adelic cyclic
quotient, a circle Tate construction on $\mathrm{THH}$, and
archimedean cyclic homology. The table and remarks record the
correction that Hesselholt's formula uses $\mathrm{TP}$ rather than
$\mathrm{TC}$. This pass found no direct identification among the
adelic quotient, $\mathrm{TP}$, and $\mathrm{TC}$ in those sources;
it does not rule out a future comparison. The source audit is in
\nolinkurl{scratch/2026-10-15-fm-ccm-names/report.md}.

\subsection{Formalization}\label{sec:lean}
The Lean formalization was written on 25--26~September~2026 by Claude
Opus~5.5 in Claude Code~2.1.282. It uses Lean~4.34.1
\cite[pp.~625--635]{Moura2021lean} and mathlib at tag
\nolinkurl{v4.34.1}
\cite[\href{https://github.com/leanprover-community/mathlib4}
{project repository}]{Mathlib2026library}; the recorded checks build the project and inspect the
axioms of its declarations. The table below states claim by claim
what is encoded, what is assumed, and what is outside the formalization.
The finite-dimensional reductions and their proof checking do not
construct the arithmetic cohomology or discharge the analytic inputs.
Machine checking is not human review of the formalization choices.

The \texttt{lean/} directory of this repository contains a Lean~4
formalization (Lean~4.34.1 with Mathlib~v4.34.1) of the statements
listed below, by claim identifier; the tree contains no
\texttt{axiom} declaration and no \texttt{sorry}, so every declaration
in it depends only on the standard axioms \texttt{propext},
\texttt{Classical.choice}, and \texttt{Quot.sound}; the script
\texttt{lean/scripts/PrintAxioms.lean} prints the axioms of $202$
declarations, among them every theorem named in the table below
(definitions such as \texttt{completedXi} and \texttt{dirichletDatum}
are not among the $202$).
 The formal
statements are conditional in the same way as the text: they take the
axioms as hypotheses and never assert the Riemann hypothesis.

The central formal fact is an equivalence about the structure
\[
\text{\texttt{WeilRH.ArithmeticWeilCohomologyQMult}},
\]
which records a complex inner product space (the positive form $h_1$ of
\ref{ax:AW4}), a linear map
$\Theta$ that is locally finite, identity \eqref{eq:inf-pol} with
$j=1$, and the order condition of \ref{ax:AW2}: every eigenspace
$\ker(\Theta-\rho)$ is finite-dimensional of dimension equal to the
order of $\xi$ at $\rho$. The theorem
\texttt{WeilRH.\allowbreak{}nonempty\_\allowbreak{}%
arithmeticWeil\allowbreak{}CohomologyQMult\_\allowbreak{}iff} states
that this structure is inhabited if and only if the Mathlib
statement \texttt{Riemann\allowbreak{}Hypothesis} holds. This is the
formal counterpart of Proposition~\ref{prop:converse} with
multiplicities; it encodes no cup product, trace, or star, so it
formalizes the spectral part of
Proposition~\ref{prop:content}\ref{item:ct-iff} for $K=\Q$ only. The
data encoded are equivalent to the Riemann hypothesis, and the content of a
proof along these lines lies in constructing them from independent
arithmetic input.

The formalization builds on, and in places overlaps with, earlier and
contemporaneous work. The statement \texttt{RiemannHypothesis} is
Mathlib's \ghlink{Mathlib2026RiemannHypothesis}{leanprover-community/mathlib4},
introduced by D.~Loeffler, and the functional equations of the completed
zeta function used here come from the same development. The definition
of $\xi$ as
\texttt{completedXi} coincides with \texttt{Complex.riemannXi} of the
PrimeNumberTheoremAnd project
\ghlink{Cipollina2026CompletedXi}{AlexKontorovich/PrimeNumberTheoremAnd}
and with \texttt{riemannXi} of Bulka's formalization of Li's criterion
\ghlink{Bulka2026LiCriterionLean}%
{nicholasbulka/li-criterion-rh-equivalence-lean}. An independent,
contemporaneous Lean formalization of the Hilbert--P\'olya reduction
\ghlink{Csillag2026HilbertPolyaLean}%
{frostynumeral/hilbert-polya-lean-verification} proves that a
symmetric operator whose eigenvalues realize the zeros implies
\texttt{RiemannHypothesis}, and that such a realization exists if and
only if \texttt{RiemannHypothesis} holds; the structure above differs
in taking identity \eqref{eq:inf-pol} rather than symmetry as its
hypothesis, in encoding local finiteness, and in tracking
multiplicities, and the caveat that the existence of the structure is a
reformulation of the Riemann hypothesis applies to both. The thick
closure \texttt{EM.ThickClosure} of the table below is a local
packaging of Mathlib's \texttt{ObjectProperty.triangEnvelope}
\ghlink{Mathlib2026TriangEnvelope}%
{leanprover-community/mathlib4/pull/37046}, built on J.~Riou's
triangulated subcategories, and \texttt{EM.thickClosure\_le} is the
analogue of \texttt{triangEnvelope\_le\_iff}.

{\small
\setlength{\tabcolsep}{4pt}
\begin{longtable}{@{}>{\raggedright\arraybackslash}p{0.2\textwidth}
>{\raggedright\arraybackslash}p{0.36\textwidth}
>{\raggedright\arraybackslash}p{0.36\textwidth}@{}}
\toprule
Claim (statement) & Lean declarations & Fidelity\\
\midrule
\endhead
\bottomrule
\endfoot
C-PUR (Prop.~\ref{prop:purity}) &
\texttt{re\_\allowbreak{}eq\_\allowbreak{}half\_\allowbreak{}of\_\allowbreak{}%
hasEigenvalue},
\texttt{isSemisimple\_\allowbreak{}of\_\allowbreak{}additive},
\texttt{maxGenEigenspace\_\allowbreak{}eq\_\allowbreak{}%
eigenspace\_\allowbreak{}of\_\allowbreak{}additive} &
faithful for parts~\ref{item:pur-real} and~\ref{item:pur-ss}; the
eigenvalue part holds in any complex inner product space,
semisimplicity in finite dimension; part~\ref{item:pur-orth} has no
declaration of its own\\
C-WEIL (Prop.~\ref{prop:weil-linear}) &
\texttt{norm\_\allowbreak{}eq\_\allowbreak{}rpow\_\allowbreak{}%
of\_\allowbreak{}weil}, \texttt{isSemisimple\_\allowbreak{}of\_\allowbreak{}%
weil} &
faithful to the modulus and diagonalizability clauses, with $q>0$ in
place of $q>1$; the $h$-orthonormal basis and the orthogonality of
distinct eigenspaces are not stated; no $\ell$-adic cohomology,
geometric polarization, or embedding $\iota$\\
C-RH (Thm.~\ref{thm:rh}, $K=\Q$) &
\texttt{ArithmeticWeilCohomologyQ.\allowbreak{}riemannHypothesis},
\texttt{ArithmeticWeilCohomologyQMult.\allowbreak{}riemannHypothesis} &
\ref{ax:AW2} enters only as a divisor condition (set-level, or with
zero orders); regularized determinants, degrees $0$ and $2$,
\ref{ax:AW3}, the star, and cup product are not encoded\\
C-CONV (Prop.~\ref{prop:converse}) &
\texttt{nonempty\_\allowbreak{}arithmeticWeilCohomologyQ\_\allowbreak{}iff},
\texttt{nonempty\_\allowbreak{}arithmeticWeilCohomologyQMult\_\allowbreak{}%
iff} &
faithful; the second keeps multiplicities; the models are built from
the zeros and are not arithmetic\\
$\xi$ and its zeros (Def.~\ref{def:zeta}) &
\texttt{completedXi}, \texttt{completedXi\_\allowbreak{}eq\_\allowbreak{}%
zero\_\allowbreak{}iff},
\texttt{analyticOrderNatAt\_\allowbreak{}completedXi\_\allowbreak{}%
pos\_\allowbreak{}iff} &
faithful; $\xi$ is defined as an entire function and its zeros are
exactly the nontrivial zeros for every $\rho$\\
flow form of \eqref{eq:inf-pol} & \texttt{P\_\allowbreak{}iff\_\allowbreak{}%
flow},
\texttt{P\_\allowbreak{}iff\_\allowbreak{}flow\_\allowbreak{}of\_\allowbreak{}%
finiteDimensional} &
faithful for bounded generators and finite-dimensional blocks;
unbounded generators and \ref{ax:AW5} not covered\\
Q1-L4.2 (Prop.~\ref{prop:polarization-identity}) &
\texttt{WeightedStarDatum.\allowbreak{}derivation\_\allowbreak{}law},
\texttt{WeightedStarDatum.\allowbreak{}re\_\allowbreak{}eq\_\allowbreak{}%
of\_\allowbreak{}hasEigenvalue},
\texttt{WeightedStarDatum.\allowbreak{}isSemisimple} &
faithful to the weighted-slice algebraic identity, specialized to
$w=0$ in Proposition~\ref{prop:polarization-identity}; it does not
encode the full unital datum. Semisimplicity is finite-dimensional;
perfectness of
the pairing, graded commutativity, other
degrees, and \ref{ax:AW2} not encoded\\
Q2-T1$'$ (Prop.~\ref{prop:algebra-form}, Prop.~\ref{prop:aw4-forms}) &
\texttt{PositiveTrace.\allowbreak{}norm\_\allowbreak{}eq\_\allowbreak{}%
rpow\_\allowbreak{}of\_\allowbreak{}mulLeft},
\texttt{PositiveTrace.\allowbreak{}isSemisimple\_\allowbreak{}mulLeft},
\texttt{PositiveTrace.\allowbreak{}minpoly\_\allowbreak{}squarefree},
\texttt{PositiveTrace.\allowbreak{}re\_\allowbreak{}eq\_\allowbreak{}%
half\_\allowbreak{}of\_\allowbreak{}mulLeft} &
faithful over $k=\C$ for any finite-dimensional complex $*$-algebra
with positive trace, including the additive form used in AW4-E; Lean's
\texttt{PositiveTrace.}$\tau$\texttt{\_pos} asks $0<\Re\tau(aa^\dagger)$
and derives realness, which is equivalent to
Definition~\ref{def:pos-involution} given
$\tau(a^\dagger)=\overline{\tau(a)}$; the cases $k=\Q,\R$,
part~\ref{item:af-rep}, and the packaging of AW4-E as a subalgebra of
$\operatorname{End}(V^j)$ are not encoded\\
C-A (Thm.~\ref{thm:A}) &
\texttt{grh\_\allowbreak{}dedekind\_\allowbreak{}of\_\allowbreak{}awc},
\texttt{riemannHypothesis\_\allowbreak{}of\_\allowbreak{}awc\_\allowbreak{}%
dedekind\_\allowbreak{}rat},
\texttt{dirichletDatum}, \texttt{grh\_\allowbreak{}dirichlet\_\allowbreak{}%
of\_\allowbreak{}awc} &
part~\ref{item:A-dedekind} in class form, with the encoding of row
C-RH (set-level spectrum; no cup product, trace, star, or
multiplicities), which asks less of each object than
\ref{ax:AW1}--\ref{ax:AW4}; Mathlib has no continuation
of $\zeta_K$, so $\Lambda_K$ is supplied as data carrying Hecke's
theorem, proved unique, and constructed only for $K=\Q$, where the
statement recovers \texttt{RiemannHypothesis}; for other $K$ the
hypothesis cannot be instantiated inside this development, since no
continuation is constructed in Lean, although one exists
mathematically. The
hypothesis quantifies over all fields $K$ in the universe \texttt{Type}
and over every continuation datum, which is harmless, since the
conclusion does not depend on
the continuation (\texttt{dedekindDatum\_\allowbreak{}grh\_\allowbreak{}iff}).
Dirichlet
$L$-functions are stated directly, not derived from
part~\ref{item:A-dedekind}; parts~\ref{item:A-hecke}--\ref{item:A-artin}
(class field theory, Brauer induction), descent, and multiplicities are
not encoded\\
Q1 Cor.~5.2 (Prop.~\ref{prop:content}) &
\texttt{WeightedStarDatum.\allowbreak{}lineModel},
\texttt{R2.\allowbreak{}counterexample} &
witnesses only: a positive one-dimensional model, and the indefinite
form satisfying \eqref{eq:inf-pol} with an off-line eigenvalue; the
latter is not repackaged as a Weil datum, and
parts~\ref{item:ct-free}--\ref{item:ct-iff} are not formalized\\
C-B (Thm.~\ref{thm:B}\ref{item:B-aw4},
Props.~\ref{prop:suspension}, \ref{prop:susp-converse}) &
\texttt{suspension\_\allowbreak{}P}, \texttt{flow\_\allowbreak{}suspension},
\texttt{weil\_\allowbreak{}of\_\allowbreak{}P\_\allowbreak{}of\_\allowbreak{}%
flow}, \texttt{suspension\_\allowbreak{}P\_\allowbreak{}iff\_\allowbreak{}%
weil} &
for a given positive form and, one $n\in\Z$ at a time, the sheet
$\alpha\mapsto(\operatorname{Log}\alpha+2\pi in)/\log q$ with the same
$n$ on every eigenspace (Def.~\ref{def:aw4-EF}; $\operatorname{Log}$
the principal branch): the form is a positive Frobenius polarization
of weight $j$ if and only if it satisfies \eqref{eq:inf-pol} for the
generator of the sheet and the flow of that generator at time $\log q$
is $\mathrm F$; the AW4-E leg of
Thm.~\ref{thm:B}\ref{item:B-aw4}, the existence direction of
Prop.~\ref{prop:weil-converse}\ref{item:wc-exist}, the full direct sum,
the nilpotent part, the regularized determinant
(Thm.~\ref{thm:ff-det}), the suspended duality and star, and
$\ell$-adic cohomology are not encoded\\
R2 (Props.~\ref{prop:R2}, \ref{prop:R2-star}) &
\texttt{R2.\allowbreak{}counterexample}, \texttt{R2.\allowbreak{}%
Jplus\_\allowbreak{}not\_\allowbreak{}commute} &
faithful to the matrices, for all real $a\ne0$ and $T$, including
that the indefinite form satisfies \eqref{eq:inf-pol}; the signature
$(2,2)$ of part~\ref{item:R2-form} enters only through $H^2=I$ and
$\tr H=0$, and part~\ref{item:R2-off} only as one eigenvalue off the
line; the graded
algebra and Prop.~\ref{prop:R2-star}\ref{item:R2s-none} not encoded\\
R3 (Prop.~\ref{prop:R3}) & \texttt{R2.\allowbreak{}det\_\allowbreak{}eq},
\texttt{R2.\allowbreak{}isUnit\_\allowbreak{}sub\_\allowbreak{}intCast} &
the polynomial and the invertibility of integer twists, for $T\ne0$
(the text assumes $T>0$); ``$p(s)>0$ for real $s$'' is formalized as
$p(s)\ne0$ for real $s$ (\texttt{R2.\allowbreak{}p\_\allowbreak{}real\_%
\allowbreak{}ne\_\allowbreak{}zero}); determinant lines not encoded\\
C-K6 (Thm.~\ref{thm:kun-hodge}\ref{item:kh-iso}, \ref{item:kh-iff},
\ref{item:kh-pure}) &
\texttt{posDef\_\allowbreak{}kronecker\_\allowbreak{}iff},
\texttt{SurfaceDatum.\allowbreak{}h\_\allowbreak{}eq\_\allowbreak{}%
zero\_\allowbreak{}of\_\allowbreak{}re\_\allowbreak{}ne\_\allowbreak{}half},
\texttt{SurfaceDatum.\allowbreak{}hodgeIndex\_\allowbreak{}iff\_\allowbreak{}%
pureDefinite},
\texttt{SurfaceDatum.\allowbreak{}re\_\allowbreak{}eq\_\allowbreak{}half\_%
\allowbreak{}of\_\allowbreak{}hodgeIndex},
\texttt{SurfaceDatum.\allowbreak{}isSemisimple\_\allowbreak{}of\_%
\allowbreak{}hodgeIndex} &
a finite-dimensional, blockwise Gram-matrix model of $N_S$ (the Hodge
index clause as positivity of the Kronecker products of the Gram
matrices of $h_1$ on $V_\rho$ and $V_{1-\rho}$); the three
\texttt{SurfaceDatum} statements about the Hodge index clause assume
$V^1$ finite-dimensional (\texttt{FiniteDimensional}), which loses
nothing, because $h_\otimes$ is positive definite on the algebraic
direct sum $N_K$ if and only if it is so on each block: off-line
blocks are $h_\otimes$-isotropic by part~\ref{item:kh-iso}, and
distinct on-line blocks are $h_\otimes$-orthogonal, as in the proof of
part~\ref{item:kh-iff}; \texttt{h\_eq\_zero\_of\_re\_ne\_half} holds
without it, and \texttt{posDef\_kronecker\_iff} concerns finite
matrices on nonempty index sets; of part~\ref{item:kh-pure} only
diagonalizability and $\Re\rho=\tfrac12$ are formalized, the clauses on
AW4-V, AW4-E, and on $\xi_K$ having no declaration; the form $B$
is assumed nondegenerate in the
first variable only (\texttt{SurfaceDatum.\allowbreak{}B\_nondeg}),
which is weaker than \ref{ax:AW3} and suffices; no completion, no
diagonal classes, and no K\"unneth map $\kappa$; the signature formula
of part~\ref{item:kh-sign} is not formalized, and the Riemann
hypothesis for $\zeta_K$ is not derived, since that needs
\ref{ax:AW2}\\
R12 (Prop.~\ref{prop:R12}) & \texttt{R12.\allowbreak{}fourier\_\allowbreak{}eq},
\texttt{R12.\allowbreak{}cos\_\allowbreak{}add\_\allowbreak{}%
cosh\_\allowbreak{}eq\_\allowbreak{}zero\_\allowbreak{}plus},
\texttt{R12.\allowbreak{}cos\_\allowbreak{}add\_\allowbreak{}%
cosh\_\allowbreak{}eq\_\allowbreak{}zero\_\allowbreak{}minus} &
the atomic measure: its Fourier transform, and that $\pi\pm ia$ are
zeros; that the zeros are exactly the points of
part~\ref{item:R12-atoms}, simple, and non-real, is not formalized; the
Gaussian variant not encoded\\
Def.~\ref{def:aw4-EF}, Prop.~\ref{prop:aw4-EF} &
\texttt{AW4EFDatum}, \texttt{SatisfiesAW4EF}, \texttt{IsSheet},
\texttt{AW4EFDatum.\allowbreak{}sheetElem};
\ref{item:EF-pure}: \texttt{AW4EFDatum.\allowbreak{}isUnit},
\texttt{.mul\_\allowbreak{}inv}, \texttt{.isSemisimple},
\texttt{.norm\_\allowbreak{}eq};
\ref{item:EF-sheet}: \texttt{.star\_\allowbreak{}sheetElem},
\texttt{.}$\iota$\texttt{\_\allowbreak{}sheetElem\_\allowbreak{}apply};
\ref{item:EF-add}: \texttt{.truncation},
\texttt{AW4EAddDatum.\allowbreak{}re\_\allowbreak{}eq\_\allowbreak{}half};
\ref{item:EF-dual}: \texttt{AW4EFDatum.\allowbreak{}dual} &
the subalgebra $A$ is encoded as an abstract algebra with an injective
unital map to $\operatorname{End}_\C(V)$, an equivalent formulation;
the suspension $(\widehat V,\Theta)$, the sheet spaces, and the sheet
maps are not built, so \ref{item:EF-sheet} is encoded as the action of
$\theta_\sigma$ by $\sigma(\alpha)$ on $\ker(\mathrm F-\alpha)$, and
\ref{item:EF-add} for $\widehat V'$ the sum of the copies of $V$
indexed by finitely many sheets, without the claim that such spaces
exhaust $\widehat V$\\
Prop.~\ref{prop:rosati-EF} &
\texttt{rosati\_\allowbreak{}aw4EF},
\texttt{rosati\_\allowbreak{}gram\_\allowbreak{}posDef},
\texttt{rosati\_\allowbreak{}trace\_\allowbreak{}pos} &
part~\ref{item:rEF-EF} for an abstract algebra with a positive trace
involution acting faithfully on $V$ and an element $\pi$ with
$\pi\pi^\dagger=q$, all three taken as hypotheses; of
part~\ref{item:rEF-inv} only the positivity step is formalized: a
rational symmetric Gram matrix positive on $\Q^n$ gives a positive
definite Hermitian form on $\C^n$; the Jacobian,
$\operatorname{End}^0(J)$, $V_\ell J$, part~\ref{item:rEF-inj}, and
Theorem~\ref{thm:imp-rosati} are not encoded\\
Prop.~\ref{prop:ff-phantom} &
\texttt{sum\_\allowbreak{}phantomSet\_\allowbreak{}pow\_\allowbreak{}%
eq\_\allowbreak{}zero},
\texttt{sum\_\allowbreak{}phantomSet\_\allowbreak{}pow\_\allowbreak{}self},
\texttt{norm\_\allowbreak{}of\_\allowbreak{}mem\_\allowbreak{}phantomSet},
\texttt{div\_\allowbreak{}mem\_\allowbreak{}phantomSet},
\texttt{spectrum\_\allowbreak{}suspension\_\allowbreak{}phantom} &
both parts; in part~\ref{item:ph-susp} the spectrum is encoded as the
union of the progressions $\Sigma_q(\alpha)$, $\alpha\in S_N$, which is
the index set of Definition~\ref{def:log-susp}\\
Thm.~\ref{thm:ff-lefschetz} &
\texttt{Lefschetz.\allowbreak{}altCount\_\allowbreak{}eq\_\allowbreak{}%
of\_\allowbreak{}lefschetzNumber\_\allowbreak{}eq},
\texttt{Lefschetz.\allowbreak{}zetaRat\_\allowbreak{}eq\_\allowbreak{}%
of\_\allowbreak{}lefschetzNumber\_\allowbreak{}eq} &
a finite analogue of part~\ref{item:fl-converse}: for finitely many
eigenvalues with multiplicities in each degree, the numbers $N_k$ for
all $k\ge1$ determine the alternating multiplicities at every
$\alpha\ne0$ and the rational function
$\prod_j\det(1-\mathrm FT\mid V^j)^{(-1)^{j+1}}$; $N_k$ is the
alternating sum of power sums of eigenvalues, not a trace of a linear
map; parts~\ref{item:fl-counts} and~\ref{item:fl-formula} and the
infinite suspended spectrum are not formalized\\
Props.~\ref{prop:EM}, \ref{prop:EM-criterion} &
\texttt{EM.\allowbreak{}ThickClosure},
\texttt{EM.\allowbreak{}thickClosure\_\allowbreak{}le},
\texttt{EM.\allowbreak{}isIso\_\allowbreak{}app\_\allowbreak{}of\_%
\allowbreak{}mem\_\allowbreak{}thickClosure},
\texttt{EM.\allowbreak{}isIso\_\allowbreak{}comp\_\allowbreak{}of\_%
\allowbreak{}mem\_\allowbreak{}thickClosure} &
the thick-subcategory induction in any pretriangulated category: a
natural transformation of triangulated functors compatible with the
shift that is an isomorphism on a class of objects is one on its thick
closure, the smallest strictly full triangulated subcategory closed
under direct summands containing it; \'etale derived categories,
cochain enhancements, $\mu_S$, Proposition~\ref{prop:EM-fibre}, and the
filtration and truncation steps of Proposition~\ref{prop:EM-criterion}
are inputs, not formalized\\
\end{longtable}
}

All declaration names lie in the namespace \texttt{WeilRH}. The Lean
conventions differ from the text in one point that is proved harmless:
Mathlib's inner product is linear in the second argument, and the
formal statements use the form $h(u,v)=\langle v,u\rangle$, for which
identity \eqref{eq:inf-pol} with real $j$ is equivalent in either
convention. The remaining formal obligations are the regularized
determinants and AW-det, the continuation of $\Lambda_K$ in Mathlib,
parts~\ref{item:A-hecke}--\ref{item:A-artin} of Theorem~\ref{thm:A},
the Hilbert realization \ref{ax:AW5}, Proposition~\ref{prop:R4}, the
full suspension of Section~\ref{subsec:suspension}, and the
determinant and trace-formula statements of
Sections~\ref{subsec:det-normalization} and~\ref{subsec:trace-axiom}.
 The file
\texttt{lean/README.md} records, for each declaration, the
correspondence with the statement it formalizes, and the script
\texttt{lean/scripts/PrintTypes.lean} prints its full type.

\appendix
\section{Token accounting}\label{app:tokens}

Table~\ref{tab:tokens} reports the retained counting windows and
estimates, not a complete bill or a total for the project. The measured
19~September window records 653,852 uncached input tokens, 16,242,944
cache-read tokens, and 107,121 output tokens, with no recorded cache
writes. The historical rate-card check on 19~September~2026 used
USD~10, 1, and~50 per million tokens in those three classes for the
Standard scenario and 2.5 times those rates for the Fast scenario.
The resulting valuations are USD~28.137514 and~70.343785, respectively.
They are alternative token-charge scenarios, not amounts to add, and
exclude search and other tool fees; the actual service tier and billing
agreement are not recorded.

All remaining quantities below are unsplit harness aggregates or
context-budget estimates. They exclude unreported primary-agent use
and, where noted in the records, later review rounds. They are left
unpriced here because their input, output, and cache classes cannot be
recovered. Approximation signs distinguish budget estimates from
retained aggregate counters; neither is a complete project total.
The model identifiers and harness versions are given in
Section~\ref{sec:ai}.

{\small
\begin{longtable}{@{}p{0.32\textwidth}p{0.24\textwidth}
 p{0.22\textwidth}p{0.13\textwidth}@{}}
\caption{Retained AI usage; incomplete windows and estimates.}
\label{tab:tokens}\\
\toprule
Task & Model and harness & Tokens & USD valuation\\
\midrule
\endfirsthead
\toprule
Task & Model and harness & Tokens & USD valuation\\
\midrule
\endhead
\bottomrule
\endfoot
19 September arithmetic cohomology research & GPT-6 Astra, Codex
desktop & 17,003,917; class split above & 28.137514 Standard;
70.343785 Fast\\
Finite-field calibration & Fable 5.1, Claude Code & $\approx380,000$
unsplit & unpriced\\
Arithmetic K\"unneth research & Fable 5.1, Claude Code &
$\approx480,000$ unsplit & unpriced\\
Trace-formula research & Fable 5.1, Claude Code & $\approx400,000$
unsplit & unpriced\\
First Lean reduction & Opus 5.5, Claude Code & $\approx170,000$ plus
71,191 & unpriced\\
Lean multiplicities & Opus 5.5, Claude Code & $\approx165,000$ &
unpriced\\
First synthesis, eight descendants & Opus 5.5, Claude Code &
1,996,417 & unpriced\\
Lean class calibration and surface & Opus 5.5, Claude Code &
$\approx265,000$ and $\approx215,000$ & unpriced\\
Second synthesis, five descendants & Opus/Fable/Sonnet, Claude Code &
1,288,264 & unpriced\\
First revision, four descendants & Opus/Fable/Sonnet, Claude Code &
1,173,021 & unpriced\\
Determinant source-retrieval descendant & Fable 5.1, Claude Code &
180,421 & unpriced\\
Second referee pass, two descendants & Fable 5.1, Claude Code &
646,091 & unpriced\\
Third referee pass, one descendant & Fable 5.1, Claude Code &
136,153 & unpriced\\
Second revision, three descendants & Opus/Fable, Claude Code &
801,948 & unpriced\\
Fourth referee pass, two descendants & Fable 5.1, Claude Code &
418,967 & unpriced\\
Fourth revision, reviewer & Opus 5.5, Claude Code & 218,300 &
unpriced\\
Eilenberg--Moore source retrieval & Fable 5.1, Claude Code &
209,948 & unpriced\\
Fifth revision, reviewer & Fable 5.1, Claude Code & 192,601 &
unpriced\\
Other primary-agent work and 11 September research & Models as
disclosed in Section~\ref{sec:ai} & not recorded & not recorded\\
14 October revision and reviews & GPT-6 family; requested Astra,
Sol, and Luna; Codex desktop & not exposed & not computed\\
\end{longtable}
}

\bibliographystyle{alphaurl}
\bibliography{notes}
\end{document}
